% Options for packages loaded elsewhere
\PassOptionsToPackage{unicode}{hyperref}
\PassOptionsToPackage{hyphens}{url}
\documentclass[
  ukrainian,
]{article}
\usepackage{xcolor}
\usepackage[margin=2.5cm]{geometry}
\usepackage{amsmath,amssymb}
\setcounter{secnumdepth}{-\maxdimen} % remove section numbering
\usepackage{iftex}
\ifPDFTeX
  \usepackage[T1]{fontenc}
  \usepackage[utf8]{inputenc}
  \usepackage{textcomp} % provide euro and other symbols
\else % if luatex or xetex
  \usepackage{unicode-math} % this also loads fontspec
  \defaultfontfeatures{Scale=MatchLowercase}
  \defaultfontfeatures[\rmfamily]{Ligatures=TeX,Scale=1}
\fi
\usepackage{lmodern}
\ifPDFTeX\else
\setmainfont{DejaVu Serif}
\setsansfont{DejaVu Sans}
\setmonofont{DejaVu Sans Mono}
\setmathfont{STIX Two Math}
\fi
\ifPDFTeX\else
  % xetex/luatex font selection
\fi
% Use upquote if available, for straight quotes in verbatim environments
\IfFileExists{upquote.sty}{\usepackage{upquote}}{}
\IfFileExists{microtype.sty}{% use microtype if available
  \usepackage[]{microtype}
  \UseMicrotypeSet[protrusion]{basicmath} % disable protrusion for tt fonts
}{}
\makeatletter
\@ifundefined{KOMAClassName}{% if non-KOMA class
  \IfFileExists{parskip.sty}{%
    \usepackage{parskip}
  }{% else
    \setlength{\parindent}{0pt}
    \setlength{\parskip}{6pt plus 2pt minus 1pt}}
}{% if KOMA class
  \KOMAoptions{parskip=half}}
\makeatother
\usepackage{longtable,booktabs,array}
\newcounter{none} % for unnumbered tables
\usepackage{calc} % for calculating minipage widths
% Correct order of tables after \paragraph or \subparagraph
\usepackage{etoolbox}
\makeatletter
\patchcmd\longtable{\par}{\if@noskipsec\mbox{}\fi\par}{}{}
\makeatother
% Allow footnotes in longtable head/foot
\IfFileExists{footnotehyper.sty}{\usepackage{footnotehyper}}{\usepackage{footnote}}
\makesavenoteenv{longtable}
\usepackage{graphicx}
\graphicspath{{/Users/kalinicenkovitalijmikolajovic/.codex/worktrees/5c15/Script_for_translate_book/Output_Final/}}
\setkeys{Gin}{width=0.66\linewidth,height=0.70\textheight,keepaspectratio}
\makeatletter
\newsavebox\pandoc@box
\newcommand*\pandocbounded[1]{% scales image to fit in text height/width
  \sbox\pandoc@box{#1}%
  \Gscale@div\@tempa{\textheight}{\dimexpr\ht\pandoc@box+\dp\pandoc@box\relax}%
  \Gscale@div\@tempb{\linewidth}{\wd\pandoc@box}%
  \ifdim\@tempb\p@<\@tempa\p@\let\@tempa\@tempb\fi% select the smaller of both
  \ifdim\@tempa\p@<\p@\scalebox{\@tempa}{\usebox\pandoc@box}%
  \else\usebox{\pandoc@box}%
  \fi%
}
% Set default figure placement to htbp
\def\fps@figure{htbp}
\makeatother
\ifLuaTeX
\usepackage[bidi=basic,shorthands=off,provide=*]{babel}
\else
\usepackage[bidi=default,shorthands=off,provide=*]{babel}
\fi
\ifLuaTeX
  \usepackage{selnolig} % disable illegal ligatures
\fi
\setlength{\emergencystretch}{3em} % prevent overfull lines
\providecommand{\tightlist}{%
  \setlength{\itemsep}{0pt}\setlength{\parskip}{0pt}}
\usepackage{bookmark}
\IfFileExists{xurl.sty}{\usepackage{xurl}}{} % add URL line breaks if available
\urlstyle{same}
\hypersetup{
  pdflang={uk},
  hidelinks,
  pdfcreator={LaTeX via pandoc}}

\author{}
\date{}

\usepackage{float}
\usepackage[section]{placeins}
\raggedbottom
\clubpenalty=10000
\widowpenalty=10000
\displaywidowpenalty=10000
\brokenpenalty=10000
\emergencystretch=2em
\allowdisplaybreaks[2]
\setlength{\textfloatsep}{12pt plus 2pt minus 4pt}
\setlength{\floatsep}{10pt plus 2pt minus 2pt}
\setlength{\intextsep}{10pt plus 2pt minus 2pt}
\makeatletter
\def\fps@figure{htbp}
\makeatother
\begin{document}

\section{С. Фарлоу}\label{ux441.-ux444ux430ux440ux43bux43eux443}

\section{РІВНЯННЯ З ПОХІДНИМИ ПОХІДНИМИ ДЛЯ НАУКОВЦІВ І
ІНЖЕНЕРІВ}\label{ux440ux456ux432ux43dux44fux43dux43dux44f-ux437-ux43fux43eux445ux456ux434ux43dux438ux43cux438-ux43fux43eux445ux456ux434ux43dux438ux43cux438-ux434ux43bux44f-ux43dux430ux443ux43aux43eux432ux446ux456ux432-ux456-ux456ux43dux436ux435ux43dux435ux440ux456ux432}

\section{Фарлоу С.}\label{ux444ux430ux440ux43bux43eux443-ux441.}

24 Рівняння з частими похідними для науковців та інженерів. Москва,
видавництво Мир, 1985, 384 с.

King's American Mathematician, який є підручником з теорії
диференціальних рівнянь у похідних похідних. Він характеризується
компактністю, ясністю та ясністю подачі, а також неформальним підходом
до подачі матеріалу. Вона містить багато ілюстрацій, графіків і діаграм,
але замість суворих доказів часто даються роздуми, засновані на інтуїції
або аналогіях.

Для інженерів і нематематиків --- біологів, диміків, а також студентів
університету.
\texttt{\textbackslash{}Phi\ =\ \textbackslash{}frac\{1702050000-380\}\{041(01)-85\}\ 26-85,\ \textbackslash{}text\{\ q.\ 1\}}
BBK 22.161.6 \# З РЕДАКТОРА ПЕРЕКЛАДУ

Книга, запропонована для читачів, слід розглядати як підручник для
курсів: «Диференціальні рівняння в похідних похідних» та «Рівняння
математичної фізики» для студентів технічних університетів. Вона також
буде корисною для інженерів різних спеціальностей. У книзі представлено
традиційні методи розв'язання диференціальних рівнянь у похідних
похідних, включаючи деякі чисельні методи. Водночас особлива увага
приділяється методу розділення змінних і методу інтегральних
перетворень, і цю частину книги можна використовувати на практичних
заняттях в університетах

Орієнтація книги на широку «нематематичну» аудиторію визначила стиль
викладання розглянутих питань: інтуїтивний підхід і велика увага до
фізичного значення не лише самих рівнянь, а й меж і початкових умов для
різних задач. Вважається, що читачі знайомі лише з основами
диференціального та інтегрального числення та елементами теорії
звичайних диференціальних рівнянь.

Фахівці можуть не задовольнитися інтуїтивним описом певних понять і
методів автором, але більш суворий виклад неминуче призведе до
ускладнення книги і звузить коло її читачів, що, очевидно, не було
наміром автора.

Більшим запереченням може бути відсутність загальних результатів щодо
задач для диференціальних рівнянь у похідних з змінними коефіцієнтами,
включно з рівняннями вищого порядку. Однак читач має чітко розуміти, що
ця книга є вступним курсом з теорії диференціальних рівнянь у часткових
похідних. Питання лінійної теорії диференціальних рівнянь, які автор не
включив, можна знайти в літературі, наведеній наприкінці книги.

Матеріал книги подається автором з великою методологічною майстерністю.
Однак не все було однаково успішним, і винний читач знайде це,
наприклад, у вивченні принципу Дюамеля, сферичних заповітів та деяких
інших уривків у книзі.

Без сумніву, запропонована книга С. Фарлоу знайде читача не лише серед
студентів і інженерів, а й серед викладачів, які можуть багато чого
навчитися з конкретного способу презентації та розробки спеціальних
математичних курсів.

\section{ПЕРЕДМОВА}\label{ux43fux435ux440ux435ux434ux43cux43eux432ux430}

За останні кілька років спостерігається значне зростання кількості
учнів, які починають вивчати диференціальні рівняння з похідними у своїх
одинадцятих класах. Крім того, багато з них спеціалізуються на сферах,
де інтуїція важливіша за математичну строгість. Тому, пишучи цю книгу, я
прагнув створити текст, який стимулює інтуїтивне мислення читачів, не
втрачаючи при цьому значної математичної строгості. Можливо, на одному
полюсі, подати предмет на високому рівні «епсилон-дельта», але тоді, як
це зазвичай буває, студенти не знатимуть, що робити з отриманою
інформацією. Інша крайність --- повне ігнорування математичних тонкощів,
але тоді ні учні, ні вчителі не знатимуть, як діяти. Я намагався знайти
розумний компроміс між цими крайнощами.

Ця книга виникла з курсу лекцій, які я читав протягом останніх п'яти
років. Розташування матеріалу не є загальноприйнятим --- вони майже
незалежні один від одного

Друг лекцій.

Основна увага приділяється двом найважливішим аналітичним методам:
розділенню змінних та інтегральним перетворенням. Також представлені
деякі нестандартні розділи, такі як метод Монте-Карло, варіаційне
числення, теорія керування, теорія потенціалів, інтегральні рівняння. Це
пов'язано з тим, що більшості студентів доведеться стикатися з цими
питаннями під час навчання, і якщо вони не ознайомляться з ними зараз,
навряд чи коли-небудь систематично їх вивчатимуть.

Цю книгу можна використовувати як навчальний посібник для одно- або
двосеместрового курсу як для початківців, так і для старших студентів.
Читачі повинні знати лише диференціальне та інтегральне числення та
звичайні диференціальні рівняння. Матеріал більшості лекцій подається в
одному-двох уроках. Типовий односеместровий курс може включати лекції
1-13, 15-17, 19-20, 22-23, 25-27, 30-32, 37-39. Усі 47 лекцій легко
читати протягом двох семестрів. Водночас часу буде достатньо для
вирішення проблем.

Автор вдячний Вайлі за пропозицію написати книгу, а також рецензентам
--- професорам Крісу Рорресу та М. Куршиду Алі, які зробили низку
корисних коментарів до тексту книги. Будь-які пропозиції щодо подальшого
покращення цього підручника будуть вдячні як від учнів, так і від
викладачів. Я також вдячний Дороті, Сюзан, Александер і Дезі Фарлоу.

Стенлі Дж. Фарлоу

\section{ЗМІСТ}\label{ux437ux43cux456ux441ux442}

\subsection{Частина 1.
Вступ}\label{ux447ux430ux441ux442ux438ux43dux430-1.-ux432ux441ux442ux443ux43f}

\begin{enumerate}
\def\labelenumi{\arabic{enumi}.}
\tightlist
\item
  Вступ до теорії диференціальних рівнянь у частинних похідних
\end{enumerate}

\subsection{Частина 2. Дифузійні
задачі}\label{ux447ux430ux441ux442ux438ux43dux430-2.-ux434ux438ux444ux443ux437ux456ux439ux43dux456-ux437ux430ux434ux430ux447ux456}

\begin{enumerate}
\def\labelenumi{\arabic{enumi}.}
\setcounter{enumi}{1}
\tightlist
\item
  Задачі типу дифузії (параболічні рівняння)
\item
  Граничні умови в задачах дифузії
\item
  Виведення рівняння теплопровідності
\item
  Розділення змінних
\item
  Перетворення неоднорідних граничних умов в однорідні
\item
  Розв'язання складніших задач методом розділення змінних
\item
  Перетворення складних рівнянь у простішу форму
\item
  Розв'язання неоднорідних УЧП методом власних функцій
\item
  Інтегральні перетворення (синусне і косинусне)
\item
  Ряди і перетворення Фур'є
\item
  Перетворення Фур'є та його застосування до УЧП
\item
  Перетворення Лапласа
\item
  Принцип Дюамеля
\item
  Конвективний член \texttt{u\_x} в дифузійній задачі
\end{enumerate}

\subsection{Частина 3. Хвильові
задачі}\label{ux447ux430ux441ux442ux438ux43dux430-3.-ux445ux432ux438ux43bux44cux43eux432ux456-ux437ux430ux434ux430ux447ux456}

\begin{enumerate}
\def\labelenumi{\arabic{enumi}.}
\setcounter{enumi}{15}
\tightlist
\item
  Одномірне хвильове рівняння (гіперболічні рівняння)
\item
  Формула Д'Аламбера
\item
  Формула Д'Аламбера (продовження)
\item
  Хвильове рівняння і граничні умови
\item
  Коливання обмеженої струни (стоячі хвилі)
\item
  Коливання балки (рівняння четвертого порядку)
\item
  Перехід до безрозмірних змінних
\item
  Класифікація рівнянь з ЧП (канонічна форма гіперболічного рівняння)
\item
  Хвильове рівняння у вільному просторі (2D і 3D)
\item
  Скінченні перетворення Фур'є (синус- і косинус-перетворення)
\item
  Принцип суперпозиції --- основа теорії лінійних систем
\item
  Рівняння першого порядку (метод характеристик)
\item
  Нелінійні рівняння першого порядку (закони збереження)
\item
  Системи рівнянь з частинними похідними
\item
  Коливання мембрани (хвильове рівняння в полярних координатах)
\end{enumerate}

\subsection{Частина 4. Еліптичні
задачі}\label{ux447ux430ux441ux442ux438ux43dux430-4.-ux435ux43bux456ux43fux442ux438ux447ux43dux456-ux437ux430ux434ux430ux447ux456}

\begin{enumerate}
\def\labelenumi{\arabic{enumi}.}
\setcounter{enumi}{30}
\tightlist
\item
  Лапласіан (інтуїтивний опис)
\item
  Загальні властивості крайових задач
\item
  Внутрішня задача Діріхле для круга
\item
  Задача Діріхле в кільці
\item
  Рівняння Лапласа в сферичних координатах (сферичні гармоніки)
\item
  Неоднорідна задача Діріхле (функція Гріна)
\end{enumerate}

\subsection{Частина 5. Чисельні та наближені
методи}\label{ux447ux430ux441ux442ux438ux43dux430-5.-ux447ux438ux441ux435ux43bux44cux43dux456-ux442ux430-ux43dux430ux431ux43bux438ux436ux435ux43dux456-ux43cux435ux442ux43eux434ux438}

\begin{enumerate}
\def\labelenumi{\arabic{enumi}.}
\setcounter{enumi}{36}
\tightlist
\item
  Чисельні розв'язки (еліптичні задачі)
\item
  Явні різницеві схеми
\item
  Неявні різницеві схеми (схема Кранка--Ніколсона)
\item
  Порівняння аналітичних і чисельних розв'язків
\item
  Класифікація рівнянь (параболічні та еліптичні рівняння)
\item
  Метод Монте-Карло (вступ)
\item
  Розв'язання УЧП методом Монте-Карло
\item
  Варіаційне числення (рівняння Ейлера--Лагранжа)
\item
  Варіаційні методи розв'язання УЧП / метод Рітца
\item
  Розв'язання УЧП методами теорії збурень
\item
  Розв'язання УЧП методом конформних відображень
\end{enumerate}

\section{Лекція 1. ВСТУП ДО ТЕОРІЇ ДИФЕРЕНЦІАЛЬНИХ РІВНЯНЬ У ПОХІДНИХ
ПОХІДНИХ}\label{ux43bux435ux43aux446ux456ux44f-1.-ux432ux441ux442ux443ux43f-ux434ux43e-ux442ux435ux43eux440ux456ux457-ux434ux438ux444ux435ux440ux435ux43dux446ux456ux430ux43bux44cux43dux438ux445-ux440ux456ux432ux43dux44fux43dux44c-ux443-ux43fux43eux445ux456ux434ux43dux438ux445-ux43fux43eux445ux456ux434ux43dux438ux445}

МЕТА ЛЕКЦІЇ: Показати, що таке диференціальні рівняння в похідних
похідних, де і як вони виникають і як їх розв'язують. Коротко розгляньте
методи класифікації рівнянь і надайте список базових понять, які будуть
детально вивчені на наступних лекціях.

Більшість фізичних явищ у таких сферах, як гідродинаміка, електрика та
магнетизм, механіка, оптика та теплопередача, загалом можна описати за
допомогою рівнянь у похідних похідних (PNP). Більшість рівнянь у
математичній фізиці є диференціальними рівняннями з похідними похідними.
Правда, з деякими спрощеними припущеннями ці рівняння зводяться до
звичайних диференціальних рівнянь, але повний опис таких систем неминуче
призводить до використання диференціальних рівнянь у похідних похідних.

\section{Що таке часткові диференціальні
рівняння?}\label{ux449ux43e-ux442ux430ux43aux435-ux447ux430ux441ux442ux43aux43eux432ux456-ux434ux438ux444ux435ux440ux435ux43dux446ux456ux430ux43bux44cux43dux456-ux440ux456ux432ux43dux44fux43dux43dux44f}

Рівняння з частковими похідними --- це рівняння, яке містить
диференціали у похідних похідних. На відміну від звичайних
диференціальних рівнянь (TAC), де невідома функція залежить лише від
однієї змінної, у рівняннях часткових похідних вона залежить від кількох
змінних (наприклад, температура \texttt{u(x,\ t)} залежить від
координати x і часу t).

Тепер наведемо деякі з найважливіших диференціальних рівнянь у похідних
похідних. Для спрощення запису ми використаємо наступні нотації

\begin{verbatim}
u_t \equiv \frac{\partial u}{\partial t}
\end{verbatim}

\texttt{,\ \textquotesingle{}u\_x\ \textbackslash{}equiv\ \textbackslash{}frac\{\textbackslash{}partial\ u\}\{\textbackslash{}partial\ x\}\textquotesingle{}\ ,\ \textquotesingle{}u\_\{xx\}\ \textbackslash{}equiv\ \textbackslash{}frac\{\textbackslash{}partial\^{}2\ u\}\{\textbackslash{}partial\ x\^{}2\}\textquotesingle{}\ ,\ ...}

\section{Деякі диференціальні рівняння в похідних
похідних}\label{ux434ux435ux44fux43aux456-ux434ux438ux444ux435ux440ux435ux43dux446ux456ux430ux43bux44cux43dux456-ux440ux456ux432ux43dux44fux43dux43dux44f-ux432-ux43fux43eux445ux456ux434ux43dux438ux445-ux43fux43eux445ux456ux434ux43dux438ux445}

\begin{verbatim}
u_t = u_{xx}
\end{verbatim}

(одновимірне рівняння теплопровідності),
\texttt{u\_t\ =\ u\_\{xx\}\ +\ u\_\{yy\}} (двовимірне рівняння
теплопровідності),
\texttt{u\_\{rr\}\ +\ \textbackslash{}frac\{1\}\{r\}u\_r\ +\ \textbackslash{}frac\{1\}\{r\^{}2\}u\_\{00\}\ =\ 0}
(рівняння Лапласа за полярними координатами),
\texttt{u\_\{tt\}\ =\ u\_\{xx\}\ +\ u\_\{yy\}\ +\ u\_\{zz\}} (тривимірне
хвильове рівняння),
\texttt{u\_\{tt\}\ =\ u\_\{xx\}\ +\ zu\_t\ +\ \textbackslash{}beta\ u}
(телеграфне рівняння).

\section{Примітки щодо наведених
прикладів}\label{ux43fux440ux438ux43cux456ux442ux43aux438-ux449ux43eux434ux43e-ux43dux430ux432ux435ux434ux435ux43dux438ux445-ux43fux440ux438ux43aux43bux430ux434ux456ux432}

У всіх наведених вище прикладах функція невідома і залежить від більше
ніж однієї змінної. Змінна і (яку ми диференціюємо) називається залежною
змінною. Змінні, за якими відбувається диференціювання, називаються
незалежними змінними. Наприклад, у рівнянні.

\begin{verbatim}
u_t = u_{xx}
\end{verbatim}

Залежна змінна \texttt{u(x,\ t)} є функцією двох незалежних змінних x і
t; та в рівнянні

\begin{verbatim}
u_t = u_{rr} + \frac{1}{r} u_r + \frac{1}{r^2} u_{00}
\end{verbatim}

змінна \texttt{u(r,\ \textbackslash{}theta,\ t)} залежить від
\texttt{r,\ \textbackslash{}theta,\ t} .

\section{Чому необхідно вивчати диференціальні рівняння частинних
похідних?}\label{ux447ux43eux43cux443-ux43dux435ux43eux431ux445ux456ux434ux43dux43e-ux432ux438ux432ux447ux430ux442ux438-ux434ux438ux444ux435ux440ux435ux43dux446ux456ux430ux43bux44cux43dux456-ux440ux456ux432ux43dux44fux43dux43dux44f-ux447ux430ux441ux442ux438ux43dux43dux438ux445-ux43fux43eux445ux456ux434ux43dux438ux445}

Більшість фізичних законів природи можна сформулювати через рівняння з
похідними похідними. Прикладами є рівняння Максвелла, закон
теплопередачі Ньютона, рівняння Нав'є-Стокса, рівняння руху Ньютона та
рівняння Шредінгера в кватичній механіці. У всіх цих рівняннях фізичні
явища описуються мовою просторових і часових похідних. Похідні
з'являються в рівняннях, оскільки вони описують найважливіші фізичні
величини (такі як швидкість, прискорення, сила, тертя, потік, струм
тощо). Таким чином, виникають диференціальні рівняння в частинних
похідних, які містять невідому функцію, яку потрібно визначити.

Мета цієї книги --- показати читачеві:

\begin{enumerate}
\def\labelenumi{\arabic{enumi}.}
\item
  Як сформулювати фізичну задачу у формі диференціального рівняння у
  похідних похідних (побудова математичної моделі).
\item
  Як розв'язати диференціальне рівняння в похідних (з урахуванням
  початкових і граничних умов).
\end{enumerate}

Перш ніж перейти до побудови математичних моделей, давайте коротко
розглянемо методи розв'язання рівнянь у похідних похідних.

\section{Як розв'язати диференціальні рівняння в похідних
похідних?}\label{ux44fux43a-ux440ux43eux437ux432ux44fux437ux430ux442ux438-ux434ux438ux444ux435ux440ux435ux43dux446ux456ux430ux43bux44cux43dux456-ux440ux456ux432ux43dux44fux43dux43dux44f-ux432-ux43fux43eux445ux456ux434ux43dux438ux445-ux43fux43eux445ux456ux434ux43dux438ux445}

Гарне питання. Виявляється, існує цілий арсенал методів, придатних для
практичного використання. Найважливішими є ті, у яких диференціальні
рівняння з похідними похідними зводяться до звичайних диференціальних
рівнянь. Перелічимо десять методів розв'язання диференціальних рівнянь у
похідних похідних:

\begin{enumerate}
\def\labelenumi{\arabic{enumi}.}
\item
  Метод розділення змінних. Рівняння з факторами, отриманими з n
  незалежними змінними, зводиться до n звичайних диференціальних
  рівняння .
\item
  Метод інтегральних перетворень. Рівняння в похідних з n незалежними
  змінними зводиться до рівняння з похідними у (n-1) незалежних змінних;
  Відповідно, рівняння з частковими похідними з двома незалежними
  змінними можна звести до звичайного диференціального рівняння ).
\item
  Метод перетворення координат. Початкове диференціальне рівняння в
  похідних похідних зводиться до звичайного рівняння з частковими
  похідними або до іншого, простішого рівняння у похідних похідних
  шляхом відповідного перетворення координат (наприклад, обертання осей
  координат тощо).
\item
  Перетворення залежної змінної. Початкове диференціальне рівняння в
  похідних похідних перетворюється у таке рівняння з похідними похідними
  для іншої невідомої функції, яку легше розв'язати, ніж початкові 3).
\item
  Чисельні методи. Початкове диференціальне рівняння в похідних похідних
  зводиться до системи різницей, яку розв'язують методом ітерацій на
  комп'ютері. У багатьох випадках це єдиний спосіб розв'язати
  диференціальне рівняння з похідними похідними. Окрім різницьких
  методів для розв'язання рівнянь з частковими про-
\end{enumerate}

\begin{enumerate}
\def\labelenumi{\alph{enumi})}
\setcounter{enumi}{4}
\tightlist
\item
  Ефективніше використовувати загальні перетворення незалежних і
  залежних змінних, включаючи похідні. Ед.
\end{enumerate}

4) Це стосується повного розділення змінних. При частковому розділенні
змінних одне диференціальне рівняння в похідних похідних зводиться до
кількох рівнянь у часткових похідних з меншою кількістю незалежних
змінних. --- Ед. Ед.

2) Це відбувається у випадку одномірного інтегрального перетворення. У
випадку k-вимірного інтегрального перетворення LCP з n незалежними
змінними зводиться до рівняння часткових з n---k незалежними змінними.
--- Ед. Ед.

Існують також інші чисельні методи, зокрема ті, що базуються на
наближенні розв'язків за поліноміальними поверхнями (наближення за
сплайнами).

\begin{enumerate}
\def\labelenumi{\arabic{enumi}.}
\setcounter{enumi}{5}
\tightlist
\item
  Метод теорії збурень. Початкова нелінійна задача зводиться до
  послідовності лінійних задач, апрокенн-
\end{enumerate}

нелінійна задача 1).

\begin{itemize}
\tightlist
\item
  \begin{enumerate}
  \def\labelenumi{\arabic{enumi}.}
  \setcounter{enumi}{6}
  \tightlist
  \item
    Метод функцій Гріна. Початкові та граничні умови замінюються
    системою найпростіших джерел, і задача розв'язується для кожного
    найпростішого джерела. Повне розв'язання початкової задачі
    отримується шляхом підсумовування розв'язків елементарних джерел.
  \end{enumerate}
\item
  \begin{enumerate}
  \def\labelenumi{\arabic{enumi}.}
  \setcounter{enumi}{7}
  \tightlist
  \item
    Метод інтегральних рівнянь. Рівняння з частковими похідними
    зводяться до інтегрального рівняння (рівняння, в якому невідома
    функція знаходиться під знаком інтеграла). Існує багато різних
    методів розв'язання інтегральних рівнянь.
  \end{enumerate}
\item
  \begin{enumerate}
  \def\labelenumi{\arabic{enumi}.}
  \setcounter{enumi}{8}
  \tightlist
  \item
    Варіаційні методи. Замість рівняння з частковими похідними
    диференціальними розв'язують задачу мінімізації. Виявляється, що
    функція, яка дає мінімум певному виразу (наприклад, загальній
    енергії системи), одночасно є розв'язком початкового рівняння
    похідниху похідних 2.
  \end{enumerate}
\item
  \begin{enumerate}
  \def\labelenumi{\arabic{enumi}.}
  \setcounter{enumi}{9}
  \tightlist
  \item
    Метод розкладу за власними функціями. Розв'язок рівняння часткових
    похідних шукається як послідовність власних функцій. Ці власні
    функції знаходяться як розв'язки так званої задачі власних значень,
    що відповідає початковій задачі для рівняння часткових похідних.
  \end{enumerate}
\end{itemize}

\section{Типи рівнянь у похідних
похідних}\label{ux442ux438ux43fux438-ux440ux456ux432ux43dux44fux43dux44c-ux443-ux43fux43eux445ux456ux434ux43dux438ux445-ux43fux43eux445ux456ux434ux43dux438ux445}

Рівняння в похідних похідних можна класифікувати за багатьма
характеристиками. Класифікація рівнянь важлива, оскільки, як виявилося,
кожен клас має свою загальну теорію та методи розв'язання рівнянь.

Тут ми наведемо шість основних методів класифікації

рівняння.

\begin{enumerate}
\def\labelenumi{\arabic{enumi}.}
\tightlist
\item
  Порядок рівняння. Порядок рівняння --- це найвищий порядок часткових
  похідних, включених до рівняння. Наприклад,
\end{enumerate}

Примітка. Ед.

1) Точніше, це метод лінеаризації нелінійних задач. Загальна теорія
збурень містить лінійні та нелінійні теорії збурень, кожна з яких, у
свою чергу, є регулярною та особливою. --- Ед. Ед.

Ці методи застосовні до спеціальних задач, що зумовлено існуванням
відповідних виразів --- функціоналів для цих задач.

Застосовувався раніше

\begin{verbatim}
u_t = u_{xx}
\end{verbatim}

(рівняння другого порядку), \texttt{u\_t\ =\ u\_x} (рівняння першого
порядку), \texttt{u\_1\ =\ uu\_\{xxx\}\ +\ \textbackslash{}sin\ x}
(рівняння трьох порядків).

\begin{enumerate}
\def\labelenumi{\arabic{enumi}.}
\setcounter{enumi}{1}
\tightlist
\item
  Кількість змінних. Кількість змінних --- це кількість незалежних
  змінних. Наприклад,
\end{enumerate}

\begin{verbatim}
u_t = u_{xx}
\end{verbatim}

\texttt{(рівняння\ з\ двома\ змінними\ \textquotesingle{}x\textquotesingle{}\ і\ \textquotesingle{}t\textquotesingle{}\ ),\ \textquotesingle{}u\_t\ =\ u\_\{rr\}\ +\ \textbackslash{}frac\{1\}\{r\}\ u\_r\ +\ \textbackslash{}frac\{1\}\{r\^{}2\}\ u\_\{\textbackslash{}theta\textbackslash{}theta\}\textquotesingle{}\ (рівняння\ з\ трьома\ змінними\ \textquotesingle{}r\textquotesingle{}\ ,}

\begin{enumerate}
\def\labelenumi{\arabic{enumi}.}
\setcounter{enumi}{2}
\tightlist
\item
  Лінійність. Рівняння з частковими похідними можуть бути лінійними та
  нелінійними. У лінійних рівняннях залежна змінна та всі її часткові
  похідні включаються лінійно, зокрема, вони не множаться одна на одну,
  не множаться в квадраті тощо). Точніше, лінійне рівняння другого
  порядку з двома незалежними змінними є рівнянням вигляду
\end{enumerate}

\begin{verbatim}
(1.1) Au_{xx} + Bu_{xy} + Cu_{yy} + Du_x + Eu_y + Fu = G
\end{verbatim}

де A, B, C, D, E, F і G --- це константи або задані функції
пояснювальних змінних x і y.

\begin{verbatim}
u_{tt} = e^{-t}u_{xx} + \sin t
\end{verbatim}

(лінійне рівняння), \texttt{uu\_\{xx\}\ +\ u\_t\ =\ 0} (нелінійне
рівняння), \texttt{u\_\{xx\}\ +\ yu\_\{yy\}\ =\ 0} (лінійне рівняння),
\texttt{xu\_x\ +\ yu\_y\ +\ u\^{}2\ =\ 0} (нелінійне рівняння).

\begin{itemize}
\tightlist
\item
  \begin{enumerate}
  \def\labelenumi{\arabic{enumi}.}
  \setcounter{enumi}{3}
  \tightlist
  \item
    Однорідність. Рівняння (1.1) називають однорідним, якщо права
    частина \texttt{G(x,\ y)} ідентична нулю для всіх x і y. Якщо
    \texttt{G(x,\ y)} не ідентичне нулю, то рівняння називають
    неоднорідним.
  \end{enumerate}
\item
  \begin{enumerate}
  \def\labelenumi{\arabic{enumi}.}
  \setcounter{enumi}{4}
  \tightlist
  \item
    Типи шансів. Якщо коефіцієнти A, B, C, D, E і F рівняння (1.1) є
    сталими, то це рівняння називається рівнянням зі сталими
    коефіцієнтами (інакше рівнянням зі змінними коефіцієнтами).
  \end{enumerate}
\item
  \begin{enumerate}
  \def\labelenumi{\arabic{enumi}.}
  \setcounter{enumi}{5}
  \tightlist
  \item
    Три основні типи лінійних рівнянь. Усі лінійні рівняння з частковими
    похідними другого порядку виду (1.1) належать до одного з трьох
    типів: а) параболічних, б) гіперболічних, в) еліптичних.
  \end{enumerate}
\end{itemize}

Параболічний тип. Рівняння параболічного типу описують процеси
теплопровідності та дифузії і визначаються умовою
\texttt{B\^{}2\ -\ 4AC\ =\ 0} .

Гіперболічний тип. Рівняння гіперболічного типу описів

осцилюючі системи та хвильові рухи визначаються умовою
\texttt{B\^{}2-4AC\textgreater{}0} .

Еліпічний тип. Рівняння еліптичного типу описують усталені процеси і
визначаються умовним \texttt{B\^{}2-4AC\textless{}0} . Приклади

\begin{enumerate}
\def\labelenumi{\alph{enumi})}
\tightlist
\item
\end{enumerate}

\begin{verbatim}
u_i = u_{xx}
\end{verbatim}

, \texttt{B\^{}2\ -\ 4AC\ =\ 0} (параболічна), (гіперболічна),
(гіперболічна), (гіперболічна \texttt{y\ \textgreater{}\ 0} ),
(гіперболічна), (еліптична), (еліптична), параболічна на
\texttt{y\ =\ 0} , гіперболічна на \texttt{y\ \textless{}\ 0} .

(У випадку змінних коефіцієнтів тип рівняння може змінюватися від точки
до точки.)

\section{ЗАУВАЖЕННЯ}\label{ux437ux430ux443ux432ux430ux436ux435ux43dux43dux44f}

\begin{itemize}
\tightlist
\item
  \begin{enumerate}
  \def\labelenumi{\arabic{enumi}.}
  \tightlist
  \item
    Загалом, значення \texttt{B\^{}2-4AC} є функцією незалежних змінних.
    Відповідно, базовий тип рівняння може змінюватися в межах визначення
    рівняння, хоча це не є обов'язковим.
  \end{enumerate}
\item
  \begin{enumerate}
  \def\labelenumi{\arabic{enumi}.}
  \setcounter{enumi}{1}
  \tightlist
  \item
    У рівнянні (1.1) незалежними змінними є x і y. У багатьох задачах
    однією з двох незалежних змінних є час, і рівняння (1.1) можна
    записати через x і t.
  \end{enumerate}
\end{itemize}

\pandocbounded{\includegraphics[keepaspectratio]{images/_page_10_Figure_9.jpeg}}

PHC 1.1. Діаграма класифікації диференціальних рівнянь у похідних
похідних.

\begin{enumerate}
\def\labelenumi{\arabic{enumi}.}
\setcounter{enumi}{2}
\tightlist
\item
  Рис. 1.1. Представлено діаграму класифікації для диференціальних
  рівнянь у похідних похідних.
\end{enumerate}

\section{ЗАВДАННЯ}\label{ux437ux430ux432ux434ux430ux43dux43dux44f}

\begin{itemize}
\tightlist
\item
  \begin{enumerate}
  \def\labelenumi{\arabic{enumi}.}
  \tightlist
  \item
    Класифікуємо наступні рівняння відповідно до всіх характеристик,
    зазначених у днаграмі на рис. 1.1:
  \end{enumerate}
\item
  \begin{enumerate}
  \def\labelenumi{\alph{enumi})}
  \tightlist
  \item
    \texttt{u\_t\ =\ u\_\{xx\}\ +\ 2u\_x\ +\ u} ,
  \end{enumerate}
\end{itemize}

\texttt{6)\ \textquotesingle{}u\_t\ =\ u\_\{xx\}\ +\ e\^{}\{-t\}\textquotesingle{}\ ,\ B)\ \textquotesingle{}u\_\{xx\}\ +\ 3u\_\{xy\}\ +\ u\_\{yy\}\ =\ \textbackslash{}sin\ x\textquotesingle{}\ ,\ \textquotesingle{}u\_\{tt\}\ =\ uu\_\{xxx\}\ +\ e\^{}\{-t\}\textquotesingle{}\ .}

\texttt{-\ 2.\ Скільки\ розв\textquotesingle{}язків\ \textquotesingle{}u\_1\ =\ u\_\{xx\}\textquotesingle{}\ рівняння\ існує?\ Спробуйте\ знайти\ рішення,\ як\ \textquotesingle{}u(x,\ t)\ =\ e\^{}\{a\textbackslash{}lambda\ +\ bt\}\textquotesingle{}\ .}
- 3. Якщо функції \texttt{u\_1(x,\ t)} і \texttt{u\_2(x,\ t)}
задовольняють рівняння (1.1), чи задовольняє сума цих функцій його?
Доведи це. - 4. Ймовірно, найпростіше диференціальне рівняння в похідних
---

\begin{verbatim}
\frac{\partial u(x, y)}{\partial x} = 0.
\end{verbatim}

Чи можете ви розв'язати це рівняння? (Знайдіть усі функції, що
задовольняють це рівняння.)

\begin{enumerate}
\def\labelenumi{\arabic{enumi}.}
\setcounter{enumi}{4}
\tightlist
\item
  Що можна сказати про рівняння
\end{enumerate}

\begin{verbatim}
\frac{\partial^2 u(x, y)}{\partial x \partial y} = 0
\end{verbatim}

Чи можете ви знайти всі розв'язки цього рівняння? Скільки їх? Порівняйте
з кількістю розв'язків звичайного диференціального рівняння

\begin{verbatim}
\frac{d^2y}{dx^2} = 0.
\end{verbatim}

\section{Лекція 2. ЗАДАЧІ ТИПУ ДИФУЗІЇ (ПАРАБОЛІЧНІ
РІВНЯННЯ)}\label{ux43bux435ux43aux446ux456ux44f-2.-ux437ux430ux434ux430ux447ux456-ux442ux438ux43fux443-ux434ux438ux444ux443ux437ux456ux457-ux43fux430ux440ux430ux431ux43eux43bux456ux447ux43dux456-ux440ux456ux432ux43dux44fux43dux43dux44f}

МЕТА ЛЕКЦІЇ: Показати, як параболічні рівняння використовуються для
розв'язання задач теплопровідності та дифузії. Наведено кілька прикладів
параболічних рівнянь, а також пояснюється фізичне значення різних членів
\texttt{(u\_1,\ u\_2,\ u\_3,\ u\_4,\ u\_6)} , включених у ці рівняння.

На конкретному прикладі розглядається основна ідея зміщеної задачі для
параболічного рівняння. Одна з головних цілей цієї лекції --- надати
читачеві інтуїтивне уявлення

Розуміння параболічних проблем.

Почнемо з розгляньмо просту фізичну задачу і те, як її можна описати
математичною моделлю з рівнянням у похідних похідних.

Потім ми ускладнимо задачу і покажемо, як нові фізичні умови призводять
до появи нових рівнянь у похідних похідних. У цій лекції диференціальні
рівняння не виводяться і не розв'язуються. Про це ми поговоримо пізніше.

\section{Простий експеримент з
теплопровідністю}\label{ux43fux440ux43eux441ux442ux438ux439-ux435ux43aux441ux43fux435ux440ux438ux43cux435ux43dux442-ux437-ux442ux435ux43fux43bux43eux43fux440ux43eux432ux456ux434ux43dux456ux441ux442ux44e}

Припустимо, що ми провели простий експеримент, який розбили на такі
кроки:

КРОК 1. Візьміть достатньо довгий (скажімо,
\texttt{L=2\textbackslash{},\textbackslash{}mathrm\{m\}} ) стрижень
(наприклад, мідний) діаметром
\texttt{2\textbackslash{},\textbackslash{}mathrm\{cm\}} , бічна поверхня
(але не кінці) якої термічно ізольована. Замість стрижня можна
використовувати мідну трубку, покриту теплоізоляцією зовні та зсередини.

Іншими словами, тепло може проходити лише через кінці стрижня і не може
проходити через його бічну поверхню.

КРОК 2. Тепер помістимо цей стрижень у пристрій із фіксованою
температурою \texttt{T\_\{\textbackslash{}mathfrak\{o\}\}} (градуси
Цельсія) на достатньо довгий час, щоб температура всередині стрижня
стала

Так само, як і температура пристрою. Для простоти припустимо, що
температура пристрою \texttt{T\_0\ =\ 10\^{}\{\textbackslash{}circ\}} C.

КРОК 3. Давайте витягнемо стрижень з нашого пристрою у певний момент
часу, що зручно взяти як нуль t = 0, і приєднаємо до нього дві термопари
з кінців. Завдання цих елементів --- підтримувати фіксовані температури
\texttt{T\_1} і \texttt{T\_2} на кінцях (наприклад,
\texttt{T\_1=0\^{}\{\textbackslash{}circ\}} C і
\texttt{T\_2=50\^{}\{\textbackslash{}circ\}} C). Іншими словами, два
термостати постійно контролюють температуру на кінцях стрижня, і якщо
температура відхиляється від встановлених значень \texttt{T\_1} і
\texttt{T\_2} , активуються потужні нагрівальні або охолоджувальні
елементи (пристрої) для відповідного регулювання температури. Наш
експеримент ілюструє RPS, 2 1,

\pandocbounded{\includegraphics[keepaspectratio]{images/_page_13_Figure_4.jpeg}}

Рис. 2.1. Схема експерименту: a --- це термопара, яка підтримує
встановлену температуру \texttt{T\_1} на лівому кінці стрижня;
\texttt{\textbackslash{}delta} --- це термопара, яка підтримує
встановлену температуру \texttt{T\_2} на правому кінці стрижня.

КРОК 4. Тепер ми будемо слідкувати за температурним профілем у стрижні
за допомогою пристрою візуалізації. Чому ми хочемо провести такий
експеримент? Це інше питання, про яке ми поговоримо пізніше.

Це кінець опису нашого експерименту.

Головна мета цієї лекції --- показати, як цю фізичну проблему та її
модифікації можна змоделювати мовою параболічних рівнянь.

\section{Математична модель
теплопровідності}\label{ux43cux430ux442ux435ux43cux430ux442ux438ux447ux43dux430-ux43cux43eux434ux435ux43bux44c-ux442ux435ux43fux43bux43eux43fux440ux43eux432ux456ux434ux43dux43eux441ux442ux456}

Опис нашої фізичної проблеми вимагає трьох типів відносин.

\begin{itemize}
\tightlist
\item
  \begin{enumerate}
  \def\labelenumi{\arabic{enumi}.}
  \tightlist
  \item
    Рівняння з частковими похідними, які описують фізичне явище
    теплопровідності.
  \end{enumerate}
\item
  \begin{enumerate}
  \def\labelenumi{\arabic{enumi}.}
  \setcounter{enumi}{1}
  \tightlist
  \item
    Граничні умови, що описують процеси теплопередачі на межах.
  \end{enumerate}
\item
  \begin{enumerate}
  \def\labelenumi{\arabic{enumi}.}
  \setcounter{enumi}{2}
  \tightlist
  \item
    Початкові твердження, які описують стан системи на початку процесу.
  \end{enumerate}
\end{itemize}

\section{Рівняння
теплопровідності}\label{ux440ux456ux432ux43dux44fux43dux43dux44f-ux442ux435ux43fux43bux43eux43fux440ux43eux432ux456ux434ux43dux43eux441ux442ux456}

Базове одновимірне рівняння теплопровідності записується у вилці

(2.1)

\begin{verbatim}
u_t = \alpha^2 u_{xx}, \quad 0 < x < L, \quad 0 < t < \infty.
\end{verbatim}

Це рівняння пов'язує значення \texttt{u\_t} --- швидкість зміни
температури з часом (виміряна в градусах/с)

та величину \texttt{u\_\{xx\}} --- це увігнутий профілю температури
\texttt{u(x,\ t)} (який є мірою різниці між температурою в заданій точці
та температурою в сусідніх точках).

\pandocbounded{\includegraphics[keepaspectratio]{images/_page_14_Figure_14.jpeg}}

Рис. 2.2. Стрілки показують зміну температури відповідно до рівняння
\texttt{u\_t\ =\ \textbackslash{}alpha\^{}2\ u\_\{xx\}} .

\texttt{Це\ рівняння\ буде\ виведено\ з\ базового\ рівняння\ збереження\ тепла\ в\ наступних\ лекціях,\ але\ наразі\ ми\ розглянемо\ його\ заданим.\ Це\ рівняння\ каже,\ що\ температура\ \textquotesingle{}u\textbackslash{}left(x,\textbackslash{}\ t\textbackslash{}right)\textquotesingle{}\ (у\ певній\ точці\ t\ у\ певній\ точці\ x\ стрижня)\ зростає\ \textquotesingle{}(u\_t\textgreater{}0)\textquotesingle{}\ або\ зменшується\ \textquotesingle{}(u\_t\textless{}0)\textquotesingle{}\ залежно\ від\ того,\ чи\ є\ друга\ похідна\ \textquotesingle{}u\_\{xx\}\textquotesingle{}\ позитивною\ чи\ негативною.\ Фіг.\ Рисунок\ 2.2\ ілюструє\ зміну\ температури\ в\ різних\ точках\ на\ елементі.}

Давайте подивимось, як значення \texttt{u\_\{xx\}} можна інтерпретувати
мовою теплопровідності. Припустимо, ми наблизимо значення
\texttt{u\_\{xx\}} фінальних постів

\begin{verbatim}
u_{xx}(x, t) \cong \frac{1}{\Delta x^2} [u(x + \Delta x, t) - 2u(x, t) + u(x - \Delta x, t)].
\end{verbatim}

Це співвідношення можна переписати у вигляді

\begin{verbatim}
u_{\lambda x}(x, t) \simeq \frac{2}{\Delta x^2} \left[ \frac{u(x + \Delta x, t) + u(x - \Delta x, t)}{2} - u(x, t) \right].
\end{verbatim}

Тепер ми можемо дати наступну інтерпретацію значення \texttt{u\_\{xx\}}
.

\begin{enumerate}
\def\labelenumi{\arabic{enumi}.}
\item
  Якщо температура \texttt{u(x,\ t)} менша за середнє значення
  температури у двох сусідніх точках, то
  \texttt{u\_\{xx\}\ \textgreater{}\ 0} (тут загальний тепловий потік по
  осі x додатний).
\item
  Якщо температура \texttt{u(x,\ t)} дорівнює середньому значенню
  температур у двох сусідніх точках, то \texttt{u\_\{xx\}\ =\ 0} (тут
  загальний тепловий потік по осі x дорівнює нулю).
\end{enumerate}

\texttt{3.\ Якщо\ температура\ \textquotesingle{}u\textbackslash{}left(x,\textbackslash{}\ t\textbackslash{}right)\textquotesingle{}\ вища\ за\ середню\ температуру\ у\ двох\ сусідніх\ точках,\ то\ \textquotesingle{}u\_\{xx\}\textless{}0\textquotesingle{}\ (тут\ —\ повний\ потік}

тепла вздовж осі x є від'ємною).

\texttt{Усе\ це\ ілюструється\ на\ рис.\ 2.2.\ Іншими\ словами,\ якщо\ температура\ в\ точці\ x\ більша\ за\ середню\ температуру\ у\ двох\ сусідніх\ точках\ \textquotesingle{}x\ -\ \textbackslash{}Delta\ x\textquotesingle{}\ і\ \textquotesingle{}x\ +\ \textbackslash{}Delta\ x\textquotesingle{}\ ,\ то\ температура\ в\ точці\ x\ зменшиться.\ Відповідно,\ точна\ швидкість\ зниження\ температури,\ а\ саме\ величина\ \textquotesingle{}u\_t\textquotesingle{}\ ,\ пропорційна\ цій\ різниці.\ Коефіцієнт\ пропорційності\ \textquotesingle{}\textbackslash{}alpha\^{}2\textquotesingle{}\ визначається\ властивостями\ матеріалу.\ Ми\ розглянемо\ цю\ константу\ детальніше\ у\ майбутніх\ лекціях.}

\section{Граничні умови
(ГУ)}\label{ux433ux440ux430ux43dux438ux447ux43dux456-ux443ux43cux43eux432ux438-ux433ux443}

Усі фізичні проблеми характеризуються наявністю певних меж. Отже, щоб
рухатися вперед, ми повинні включити межі в математичну модель, щоб
адекватно описати фізичну проблему. У нашому експерименті граничні умови
(ГІ) дуже прості. Оскільки температура на кінцях x=0 і x=L завжди
залишається сталою і дорівнює \texttt{T\_1} та \texttt{T\_2} , можливо
записати

(2.2)

\begin{verbatim}
\left\{ \begin{array}{l} u\left(0,\ t\right) = T_{t}, \\ u\left(L,\ t\right) = T_{2}, \end{array} \right. \quad 0 < t < \infty,
\end{verbatim}

\section{Початкові умови
(ПУ)}\label{ux43fux43eux447ux430ux442ux43aux43eux432ux456-ux443ux43cux43eux432ux438-ux43fux443}

Усі фізичні процеси мають починатися у певний момент часу (зазвичай
приймається як нуль t=0), тому ми повинні встановити фізичні умови саме
в цю точку. Відтоді, як ми почали контролювати температуру з моменту,
коли паличка мала постійну температуру
\texttt{T\_\{\textbackslash{}mathrm\{\ \}o\}} , ми

(2.3)

\begin{verbatim}
u(x, 0) = T_0, \quad 0 \le x \le L
\end{verbatim}

Тепер ми створили математичну модель експерименту. Система співвідношень
(2.1), (2.2) та (2.3) називається задачею зміщення для рівняння
теплопровідності і зазвичай записується як

(2.4)

\begin{verbatim}
\begin{aligned} & (\text{YUII}) \quad u_t = \alpha^2 u_{xx}, & 0 < x < L, \ 0 < t < \infty, \\ & \left( \Gamma \text{Y} \right) \quad \begin{cases} u \ (0, \ t) = T_1, \\ u \ (L, \ t) = T_2, \end{cases} & 0 < t < \infty, \\ & (\text{HY}) \quad u \ (x, \ 0) = T_0, \quad 0 \leqslant x \leqslant L. \end{aligned}
\end{verbatim}

Цікаво, що це зовсім не очевидно: чи існує одна функція
\texttt{u(x,\ t)}, яка задовольняє умови (2.4), і чи описує ця функція
температуру в стрижні? Отже, наше безпосереднє завдання --- знайти цю
функцію \texttt{u(x,\ t)} з (2.4).

Перш ніж завершити цю лекцію, давайте розглянемо деякі варіації основної
проблеми. Почнемо з невеликих змін рівняння теплопровідності.

\section{Деякі рівняння дифузійного
типу}\label{ux434ux435ux44fux43aux456-ux440ux456ux432ux43dux44fux43dux43dux44f-ux434ux438ux444ux443ux437ux456ux439ux43dux43eux433ux43e-ux442ux438ux43fux443}

Теплопередача через бічну поверхню пропорційна різниці температур

Рівняння

\begin{verbatim}
u_t = \alpha^2 u_{xx} - \beta (u - u_0)
\end{verbatim}

Описує теплопровідність елемента, враховуючи не лише дифузію теплових
\texttt{\textbackslash{}alpha\^{}2\ u\_\{xx\}} вздовж елемента, а й
теплопередачу через бічну поверхню елемента. Відтік тепла (
\texttt{\textbackslash{}beta\ \textgreater{}\ 0} ) або його приплив (
\texttt{\textbackslash{}beta\ \textless{}\ 0} ) пропорційний різниці між
температурою u стрижня та температурою навколишнього середовища
\texttt{u\_0} ( \texttt{\textbackslash{}beta} --- це сталий коефіцієнт
пропорційності). Якщо коефіцієнт \texttt{\textbackslash{}beta} великий
порівняно з \texttt{\textbackslash{}alpha\^{}2} , то теплопотік уздовж
стрижня буде малим порівняно з

Протікання через бічну поверхню, і тому тепло буде проходити через неї
відповідно до закону
\texttt{u\_t\ =\ -\textbackslash{}beta\ (u\ -\ u\_a)} .

У хімії, де величина відіграє роль концентрації, рівняння

\begin{verbatim}
u_1 = \alpha^2 u_{xx} - \beta (u - u_0)
\end{verbatim}

означає, що швидкість зміни \texttt{u\_i} кількості речовини залежить як
від дифузії \texttt{\textbackslash{}alpha\^{}2\ u\_\{xx\}} (у напрямку
осі x), так і від появи ( \texttt{\textbackslash{}beta\ \textless{}\ 0}
) або розпаду ( \texttt{\textbackslash{}beta\ \textgreater{}\ 0} )
речовини в хімічній реакції, і пропорційна різниці між двома
концентраціями u і \texttt{u\_\{\textbackslash{}theta\}} .

\section{Внутрішнє джерело
тепла}\label{ux432ux43dux443ux442ux440ux456ux448ux43dux454-ux434ux436ux435ux440ux435ux43bux43e-ux442ux435ux43fux43bux430}

Неоднородне рівняння

\begin{verbatim}
u_t = \alpha^2 u_{xx} + f(x, t)
\end{verbatim}

відповідає випадку, коли всередині елемента є джерело тепла (розташоване
вздовж усього елемента і діє в будь-який момент t). Наприклад, може бути
так, що дріт з електричним струмом проходить через стрижень і опір є
постійним джерелом тепла \texttt{f(x,t)} = K.

\section{Рівняння конвективної
дифузії}\label{ux440ux456ux432ux43dux44fux43dux43dux44f-ux43aux43eux43dux432ux435ux43aux442ux438ux432ux43dux43eux457-ux434ux438ux444ux443ux437ux456ux457}

\texttt{Припустимо,\ що\ домішка\ поширюється\ вздовж\ потоку,\ що\ рухається\ зі\ швидкістю\ v.\ Очевидно,\ що\ концентрація\ \textquotesingle{}u\textbackslash{}left(x,\textbackslash{},t\textbackslash{}right)\textquotesingle{}\ домішок\ змінюється\ залежно\ від\ x\ (позитивний\ напрямок\ осі\ x\ обирається\ вздовж\ потоку)\ і\ часу\ t.\ Швидкість\ зміни\ концентрації\ \textquotesingle{}u\_t\textquotesingle{}\ визначається\ рівнянням\ конвективної\ дифузії}

\begin{verbatim}
u_t = \alpha^2 u_{xx} - v u_x.
\end{verbatim}

Термін \texttt{\textbackslash{}alpha\^{}2\ u\_\{xx\}} описує внесок
дифузії, а \texttt{vu\_x} --- конвективний компонент. Що важливіше,
дифузія або конвекція залежить від співвідношення
\texttt{\textbackslash{}alpha\^{}2} до v. Димові частинки піднімаються
конвективно через рух гарячого повітря і одночасно розсіюють через
вихорові рухи повітря. Окрім цих модифікацій рівняння теплопровідності,
можливо змінювати граничні умови на кінцях стрижня відповідно до
фізичної ситуації. Деякі з цих змін ми обговоримо у третій лекції.

\section{ЗАУВАЖЕННЯ}\label{ux437ux430ux443ux432ux430ux436ux435ux43dux43dux44f-1}

Рівняння теплопровідності
\texttt{u\_t\ =\ \textbackslash{}alpha\^{}2(x)\ u\_\{xx\}} з змінним
коефіцієнтом \texttt{\textbackslash{}alpha(x)} відповідає задачі, в якій
дифузія залежить від x (тобто матеріал є гетерогенним). Наприклад, мідь
і сталь---

\pandocbounded{\includegraphics[keepaspectratio]{images/_page_18_Picture_4.jpeg}}

PHC. 2.3.

Пластини контактують, як показано на рис. 2.3. Припустимо, що
температура \texttt{u(0,t)=0\^{}\{\textbackslash{}circ\}} C
підтримується на лівому боці міді, а температура
\texttt{u(L,t)=20\^{}\{\textbackslash{}circ\}} C --- з правого боку
сталі, тоді рівняння, що описує теплопровідність у цій системі, виглядає
так:

\begin{verbatim}
u_t = \alpha^2(x) \, u_{_{\! A \, X}}, \quad 0 < x < L,
\end{verbatim}

де

\begin{verbatim}
\alpha(x) = \begin{cases}
\alpha_1, & \text{коефіцієнт дифузії в міді}, \quad 0 < x < L/2, \\
\alpha_2, & \text{коефіцієнт дифузії в сталі}, \quad L/2 < x < L.
\end{cases}
\end{verbatim}

\section{ЗАВДАННЯ}\label{ux437ux430ux432ux434ux430ux43dux43dux44f-1}

\begin{enumerate}
\def\labelenumi{\arabic{enumi}.}
\tightlist
\item
  Якщо початкова температура в стрижні
\end{enumerate}

\begin{verbatim}
u(x, 0) = \sin \pi x, \quad 0 \leqslant x \leqslant 1,
\end{verbatim}

а граничні умови такі:

\begin{verbatim}
u(0, t) = 0, \qquad u(1, t) = 0.
\end{verbatim}

Тоді як поводитиметься температура в стрижні \texttt{U(x,\ t)} при
\texttt{t\ \textgreater{}\ 0} ? 2. Припустимо, що в стрижні є постійне
внутрішнє джерело тепла, так що рівняння теплопровідності в стрижні
записується як

\begin{verbatim}
u_t = \alpha^2 u_{xx} + 1, \quad 0 < x < 1.
\end{verbatim}

Припустимо, що ми зафіксували температуру на межі з значеннями
\texttt{u(0,\ t)} = 0 і \texttt{u(1,\ t)} = 1. Якою буде стаціонарна
температура у стрижні? Іншими словами, який c-незалежний розподіл
\texttt{U(x,\ t)} перешкоджує \texttt{u(x,\ t)}?

\begin{enumerate}
\def\labelenumi{\arabic{enumi}.}
\setcounter{enumi}{2}
\tightlist
\item
  Припустимо, що металевий стрижень втрачає тепло через бічну поверхню
  згідно з рівнянням
\end{enumerate}

\begin{verbatim}
u_1 = \alpha^2 u_{xx} - \beta u
\end{verbatim}

\texttt{0\ \textless{}\ x\ \textless{}\ 1} .

І припустимо, що ми підтримуємо кінці стрижня при температурах
\texttt{U(0,\ T)} = 1, \texttt{U(1,\ T)} = 1. Знайдіть стаціонарний
розподіл температури по стрижню. Побудуйте графік цього розподілу. У
яких моментах відбувається тепловий потік?

\texttt{4.\ Припустимо,\ що\ термоізольований\ металевий\ стрижень\ довжини\ L\ =\ 1\ має\ початкову\ температуру\ \textquotesingle{}\textbackslash{}sin(3\textbackslash{}pi\ x)\textquotesingle{}\ ,\ а\ фіксовану\ температуру\ —\ 0\ і\ \textquotesingle{}10\textbackslash{},\^{}\{\textbackslash{}circ\}\textquotesingle{}\ C\ на\ лівому\ та\ правому\ кінцях.}

\section{Лекція 3. ГРАНИЧНІ УМОВИ В ЗАДАЧАХ ДИФУЗІЙНОГО
ТИПУ}\label{ux43bux435ux43aux446ux456ux44f-3.-ux433ux440ux430ux43dux438ux447ux43dux456-ux443ux43cux43eux432ux438-ux432-ux437ux430ux434ux430ux447ux430ux445-ux434ux438ux444ux443ux437ux456ux439ux43dux43eux433ux43e-ux442ux438ux43fux443}

МЕТА ЛЕКЦІЇ: Показати, як виникають різні граничні умови у задачах
теплопровідності та дифузії, а також ввести важливе поняття теплового
потоку.

Розглядаються три найважливіші типи граничних умов:

\begin{itemize}
\tightlist
\item
  \begin{enumerate}
  \def\labelenumi{\arabic{enumi}.}
  \tightlist
  \item
    u = \texttt{g(t)} (температура задається на межі),
    \texttt{-\ \textquotesingle{}2.\ \textbackslash{}\ \textbackslash{}frac\{\textbackslash{}partial\ u\}\{\textbackslash{}partial\ n\}\ +\ \textbackslash{}lambda\ u\ =\ g\textbackslash{}left(t\textbackslash{}right)\textquotesingle{}\ (встановлення\ температури\ навколишнього\ середовища);}
  \end{enumerate}
\end{itemize}

n --- вектор зовнішньої нормалі до межі),
\texttt{3.\ \textbackslash{}frac\{\textbackslash{}partial\ u\}\{\textbackslash{}partial\ n\}\ =\ g(t)}
(визначений тепловий потік через межу).

У задачах теплопровідності зазвичай використовують три основні типи
граничних умов. Ця лекція наводить приклади того, як різні фізичні умови
експерименту призводять до різних типів граничних умов.

\section{Граничні умови першого виду (температура задається на
межі)}\label{ux433ux440ux430ux43dux438ux447ux43dux456-ux443ux43cux43eux432ux438-ux43fux435ux440ux448ux43eux433ux43e-ux432ux438ux434ux443-ux442ux435ux43cux43fux435ux440ux430ux442ux443ux440ux430-ux437ux430ux434ux430ux454ux442ux44cux441ux44f-ux43dux430-ux43cux435ux436ux456}

Розглянемо тепловий потік в одномірному стрижні, показаний на рис. 3 1,
і припустимо, що температура кінців змінюється відповідно до законів
\texttt{g\_1(t)} і \texttt{g\_2(t)} .

\pandocbounded{\includegraphics[keepaspectratio]{images/_page_20_Picture_4.jpeg}}

Рис. 3,.1. Температура буде відновлена на межі.

Як згадувалося в попередній лекції, кожен пристрій, що підтримує
встановлену температуру на кінцях, складається з термостата та
нагрівального елемента для підтримки відповідних теплових умов. Проблеми
з граничними умовами першого виду дуже поширені. У деяких випадках
складність полягає у пошуку контролерів граничних умов, тобто температур
\texttt{g\_1(t)} та \texttt{g\_2(t)} меж, які призводять до певної зміни
температури всередині стрижня. Наприклад, у металургії часто необхідно
обирати граничні умови контролю, щоб температура металу всередині печі
змінювалася з часом, а температурний градієнт був невеликим.

\pandocbounded{\includegraphics[keepaspectratio]{images/_page_20_Picture_7.jpeg}}

Фіг. 3.2. Зниження температури на межі.

Граничні умови першого виду також виникають у багатовимірних задачах
теплопровідності. Як приклад можна навести цікаву задачу визначення
температурного поля всередині кругового диска з радіусом \emph{R} (див.
рис. 3.2), коли гранична температура задається у полярних координатах
формулою
\texttt{u(R,\ \textbackslash{}theta,\ t)\ =\ \textbackslash{}cos\ t\ \textbackslash{}sin\ \textbackslash{}theta}
.

Звісно, щоб розв'язати задачу, потрібно знати початкову температуру, але
в цьому випадку початкові умови перестають впливати на розв'язок через
дуже короткий час, і отримана температура всередині диска визначається
лише значеннями температури на межі.

\section{Граничні умови другого виду (задана температура навколишнього
середовища)}\label{ux433ux440ux430ux43dux438ux447ux43dux456-ux443ux43cux43eux432ux438-ux434ux440ux443ux433ux43eux433ux43e-ux432ux438ux434ux443-ux437ux430ux434ux430ux43dux430-ux442ux435ux43cux43fux435ux440ux430ux442ux443ux440ux430-ux43dux430ux432ux43aux43eux43bux438ux448ux43dux44cux43eux433ux43e-ux441ux435ux440ux435ux434ux43eux432ux438ux449ux430}

Припустимо, що ми знову розглядаємо термічно ізольований мідний
стержень, але тепер відмовляємося від заданого температурного режиму на
кінцях стрижня. Натомість давайте наблизимо кінці стрижня до контакту з
двома середовищами. Нехай температура одного з них \texttt{g\_1(t)} , а
іншого \texttt{g\_2(t)} . Іншими словами, припустимо, що лівий кінець
стрижня поміщений у контейнер із рідиною, температура якої змінюється
відповідно до закону \texttt{g\_1(t)}. А правий кінець поміщається в
іншу рідину з температурою \texttt{g\_2(t)} (рис. 3, 3)

\pandocbounded{\includegraphics[keepaspectratio]{images/_page_21_Picture_6.jpeg}}

Рис. 3.3. Конвективна теплопередача через межі, a --- це рідина при
температурі \texttt{g\_1(t)} , \texttt{\textbackslash{}delta} --- рідина
при температурі \texttt{g\_2(t)} .

Задаючи граничні умови такого типу, ми не можемо вважати, що граничні
температури стрижня збігаються з температурами рідини. Але ми знаємо
закон Ньютона: якщо температура одного з кінців стрижня менша за
температуру відповідної рідини, тепло потече в стрижень зі швидкістю,
пропорційною різницям температур. Іншими словами, для одномірного
стрижня з межами x=0 і x=1 закон теплопередачі Ньютона формулюється так:

(3.1)

\texttt{вихідний\ тепловий\ потік\ на\ \textquotesingle{}x=0\textquotesingle{}\ дорівнює\ \textquotesingle{}h\textbackslash{}left{[}u\textbackslash{}left(0,\textbackslash{},t\textbackslash{}right)-g\_\{1\}\textbackslash{}left(t\textbackslash{}right)\textbackslash{}right{]}\textquotesingle{}\ ,\ вихідний\ тепловий\ потік\ при\ \textquotesingle{}x=L\textquotesingle{}\ дорівнює\ \textquotesingle{}h\textbackslash{}left{[}u\textbackslash{}left(L,\textbackslash{},t\textbackslash{}right)-g\_\{2\}\textbackslash{}left(t\textbackslash{}right)\textbackslash{}right{]}\textquotesingle{}\ ,}

де h --- коефіцієнт теплопередачі, який показує, скільки калорій
проходить через межу за одну секунду при різниці температур у один
градус. Вихідний потік дорівнює кількості кало-

Пройшов крізь палочний спис за одну секунду. Зверніть увагу, що вихідний
тепловий потік буде позитивним на кінці стрижня, де температура стрижня
вища за температуру навколишнього середовища. Рівняння (3.1) разом із
відомим законом теплопровідності Фур'є тепер можна використовувати для
отримання Рис

\pandocbounded{\includegraphics[keepaspectratio]{images/_page_22_Figure_3.jpeg}}

. 3.4. Ілюстрація закону Фур'є:
\texttt{\textbackslash{}frac\{\textbackslash{}partial\ u\}\{\textbackslash{}partial\ n\}\ \textgreater{}\ 0}
означає, що тепло надходить у цю область;
\texttt{\textbackslash{}frac\{\textbackslash{}partial\ u\}\{\textbackslash{}partial\ n\}\ \textless{}\ 0}
означає, що тепло виходить із цієї зони; n --- напрямок зовнішньої
нормалі;
\texttt{\textbackslash{}frac\{\textbackslash{}partial\ u\}\{\textbackslash{}partial\ n\}\ -}
Зміна температури в напрямку N;
\texttt{\textbackslash{}frac\{\textbackslash{}partial\ u\}\{\textbackslash{}partial\ n\}\ =\ -\textbackslash{}frac\{\textbackslash{}partial\ u\}\{\textbackslash{}partial\ (-n)\}}
є похідною у напрямку внутрішньої нормальної --- це похідна, взята
\texttt{\textbackslash{}mathbf\{e\}} протилежному знаку протилежного
знака пормалі.

граничні умови. Закон Фур'є дає друге представлення (перше міститься у
(3.1)) для вихідного теплового потоку. Прирівнюючи ці два вирази до
потоку, ми отримуємо граничні умови, які шукаємо. Сформулюємо закон
Фур'є (встановлений експериментально):

(3.2) Тепловий потік, що проходить через межу площі, пропорційний
нормальній похідній температури у напрямку внутрішньої норми.

Цей закон стверджує, що якщо температура швидко підвищується у напрямку
зовнішньої температури до межі D (рисунок 3.4), то потік буде текти з
навколишнього середовища до області D.

У нашій одновимірній задачі закон Фур'є має вигляд:

(3.3)

\begin{verbatim}
\begin{cases}
\text{вихідний тепловий потік при } x=0 \text{ дорівнює } k \frac{\partial u}{\partial x}, \\
\text{вихідний тепловий потік при } x=L \text{ дорівнює } -k \frac{\partial u}{\partial x}.
\end{cases}
\end{verbatim}

де k --- теплопровідність матеріалу, яка є мірою того, наскільки добре
матеріал проводить тепло. Погано провідні матеріали мають коефіцієнт
близький до нуля в одиницях GHS, тоді як мідь та алюміній мають значення
значно вище нуля.

\pandocbounded{\includegraphics[keepaspectratio]{images/_page_23_Figure_5.jpeg}}

Фіг. 3.5. Ще одна ілюстрація закону Чотирьох.

Закон Фур'є (3.3) фактично описує теплопровідність всюди всередині
стрижня, а не лише на його межі; Наприклад, (див. рис. 3.5).

(3.4) Потік, що проходить через точку

x

\texttt{зліва\ направо\ \textquotesingle{}=-k\textbackslash{}frac\{\textbackslash{}partial\ n\}\{\textbackslash{}partial\ x\}\textquotesingle{}\ Закон\ Фур\textquotesingle{}є\ (3.4)\ стверджує:\ якщо\ \textquotesingle{}u\_x\ \textless{}\ 0\textquotesingle{}\ ,\ то\ в\ точці\ \textquotesingle{}x\_0\textquotesingle{}\ тепло\ тече\ зліва\ направо,\ якщо\ \textquotesingle{}u\_x\ \textgreater{}\ 0\textquotesingle{}\ ,\ то\ в\ точці\ \textquotesingle{}x\_0\textquotesingle{}\ тепло\ тече\ справа\ наліво.\ (Тепло\ завжди\ переходить\ від\ вищих\ температур\ до\ нижчих.)}

Нарешті, якщо використовувати вирази (3.1) і (3.3) для теплового потоку,
то отримуємо бажані граничні умови для задачі, описаної на рис. 3.3, у
чисто математичній формі

\begin{verbatim}
(\Gamma Y) \quad \begin{cases} \frac{\partial u\left(0,\,t\right)}{\partial x} = \frac{h}{k}\left[u\left(0,\,t\right) - g_{1}\left(t\right)\right], \\ \frac{\partial u\left(L,\,t\right)}{\partial x} = -\frac{h}{k}\left[u\left(L,\,t\right) - g_{2}\left(t\right)\right]. \end{cases} \quad 0 < t < \infty,
\end{verbatim}

Дуже часто константа h/k позначається як \texttt{\textbackslash{}lambda}
, а граничні умови теплового потоку на межі записуються у вигляді

(3.5)

\begin{verbatim}
u_{x}(0, t) = \lambda [u(0, t) - g_{1}(t)], u_{x}(L, t) = -\lambda [u(L, t) - g_{2}(t)].
\end{verbatim}

Для випадку вищих розмірностей отримують подібні граничні умови.
Наприклад, якщо межа кругового диска осяяна рухомою рідиною з
температурою \texttt{g(\textbackslash{}theta,t)} , то граничні умови
записуються як

\begin{verbatim}
\frac{\partial u}{\partial r}(R, \theta, t) = -\frac{h}{k} [u(R, \theta, t) - g(\theta, t)].
\end{verbatim}

Тут
\texttt{\textbackslash{}frac\{\textbackslash{}partial\ u\}\{\textbackslash{}partial\ r\}(R,\ \textbackslash{}theta,\ t)}
позначає зовнішню нормальну похідну (у позитивному напрямку r-осі)
значення u, обчислену в точці \texttt{(R,\ \textbackslash{}theta)} межі.
Ми назвемо такі граничні умови лінійними (оскільки вони лінійні в U і
\texttt{u\_r} ), але неоднорідними, оскільки функція
\texttt{g(\textbackslash{}theta,\ t)} знаходиться праворуч.

\section{Граничні умови третього виду (враховуючи потік; зокрема,
враховується випадок термічно ізольованих
меж)}\label{ux433ux440ux430ux43dux438ux447ux43dux456-ux443ux43cux43eux432ux438-ux442ux440ux435ux442ux44cux43eux433ux43e-ux432ux438ux434ux443-ux432ux440ux430ux445ux43eux432ux443ux44eux447ux438-ux43fux43eux442ux456ux43a-ux437ux43eux43aux440ux435ux43cux430-ux432ux440ux430ux445ux43eux432ux443ux454ux442ux44cux441ux44f-ux432ux438ux43fux430ux434ux43eux43a-ux442ux435ux440ux43cux456ux447ux43dux43e-ux456ux437ux43eux43bux44cux43eux432ux430ux43dux438ux445-ux43cux435ux436}

Жоден потік не проходить через термоізольовані межі, тому нормальна
похідна (зовнішня чи внутрішня) має зрівнятися з нулем на межі, оскільки
потік пропорційний нормальній похідній. У випадку одномірного елемента з
ізольованими кінцями x=0 і x=L, граничні умови мають вигляд

\begin{verbatim}
u_{\lambda}(0, t) = 0,
\end{verbatim}

\texttt{\textquotesingle{}u\_\{\textbackslash{}lambda\}(L,\ t)\ =\ 0.\textquotesingle{}\ \textquotesingle{}0\ \textless{}\ t\ \textless{}\ \textbackslash{}infty,\textquotesingle{}\ У\ двовимірній\ області\ ізоляція\ межі\ означає,\ що\ нормальна\ похідна\ температури\ на\ межі\ перетворюється\ на\ нуль.\ Припустимо,\ наприклад,\ що\ круглий\ диск\ термічно\ ізольований\ уздовж\ межі.\ Тоді\ граничну\ умову\ слід\ записати\ так:\ \textquotesingle{}u\_r(R,\ \textbackslash{}theta,\ t)\ =\ 0\textquotesingle{}\ для\ всіх\ \textquotesingle{}0\ \textbackslash{}le\ \textbackslash{}theta\ \textless{}\ 2\textbackslash{}pi\textquotesingle{}\ і\ всіх\ \textquotesingle{}0\ \textless{}\ t\ \textless{}\ \textbackslash{}infty\textquotesingle{}\ .}

З іншого боку, якщо вказати кількість тепла, що проходить через межу
нашого диска, то граничні умови набувають вигляду

\begin{verbatim}
u_r(R, \theta, t) = f(\theta, t),
\end{verbatim}

де \texttt{f(0,\ t)} --- це кількість тепла, що переходить у круглий
диск із зовнішнього джерела.

Тепер проілюструємо різні типи граничних умов.

\section{Типові граничні умови для задачі одномірної
теплопровідності}\label{ux442ux438ux43fux43eux432ux456-ux433ux440ux430ux43dux438ux447ux43dux456-ux443ux43cux43eux432ux438-ux434ux43bux44f-ux437ux430ux434ux430ux447ux456-ux43eux434ux43dux43eux43cux456ux440ux43dux43eux457-ux442ux435ux43fux43bux43eux43fux440ux43eux432ux456ux434ux43dux43eux441ux442ux456}

Припустимо, що у нас є мідний стрижень довжиною 200 см, бікова поверхня
якого термічно ізольована, а початкова температура становить 0 °C.
Припустимо також, що верхній кінець стрижня x=0 термічно ізольований, а
нижній кінець x=200 змивається рухомою водою, яка має постійну
температуру \texttt{g\_2(t)=20} °C (див. рис. 3.6).

\pandocbounded{\includegraphics[keepaspectratio]{images/_page_25_Picture_6.jpeg}}

Фіг. 3.6. Змішане завдання.

Математична модель цієї задачі описується наступними чотирма
відношеннями:

\begin{verbatim}
(3.6)
\begin{aligned}
(\text{УЧП})\quad & u_t = \alpha^2 u_{xx}, && 0 < x < 200,\ 0 < t < \infty, \\
(\text{ГУ})\quad & u_x(0,t) = 0, && 0 < t < \infty, \\
& u_x(200,t) = -\frac{h}{k}\left[u(200,t) - 20\right], && 0 < t < \infty, \\
(\text{НУ})\quad & u(x,0) = 0^\circ\mathrm{C}, && 0 \leqslant x \leqslant 200.
\end{aligned}

де \alpha^2 = 1,16 <sup>см 2</sup>/с — теплова дифузія міді;

k = 0,93 кал/см·с·°C — теплопровідність міді;

h — коефіцієнт теплопередачі. Визначення цього співвідношення — велика проблема. Він вимірюється швидкістю виходу тепла від другого кінця стрижня до води. Це залежить від швидкості потоку води, стану поверхні, яку слід змивати тощо.

Читач може спробувати провести експеримент, щоб визначити ділення значення цього коефіцієнта.

# ЗАУВАЖЕННЯ

1. Типова задача внутрішньої теплопровідності для квадрата показана на рис. 3.7.

У цій задачі початкова температура u\left(x,\,y,\,t\right) при t = 0 задається всередині квадрата, і рівняння з похідними з частковими має вигляд Р и с

![](images/_page_26_Figure_9.jpeg)

. 3.7. Типові граничні умови для задачі дифузії в квадраті: a — кількість тепла, f_1(t) проходить через межу; b — термоізольована межа; a — це навколишня температура Равії 10^{\circ} C; z — температура підтримується рівною 10^{\circ} C.

а граничні умови будуть виконані на 0 < t < \infty і визначатимуть температуру в будь-який момент. Якою б не була ця температура, вона повинна відповідати граничним умовам, зазначеним на рис. 3.7.

2. Зверніть увагу, що гранична умова
\end{verbatim}

\texttt{u\_r(R,\ \textbackslash{}theta,\ t)\ =\ -\textbackslash{}frac\{h\}\{k\}\ {[}u(R,\ \textbackslash{}theta,\ t)\ -\ g(\textbackslash{}theta,\ t){]}}

\begin{verbatim}
На колі не означає, що гранична температура дорівнює g(\theta,t) , але якщо коефіцієнт теплопередачі h великий, то ці умови фактично еквівалентні вимогі g(\theta,t) температури на межі.

# ЗАВДАННЯ

1. Намалюйте контури графіків для розв'язання змішаної задачі (3.6) для різних моментів часу. Чи відповідають ваші графи визначеним граничним умовам? Яка стаціонарна температура стрижня? Очевидно, ваша відповідь базується на інтуїції?

2. Як би ви інтерпретували наступне змішане завдання?
\end{verbatim}

\texttt{\textbackslash{}begin\{aligned\}}
\texttt{(\textbackslash{}text\{УЧП\})\textbackslash{}quad\ \&\ u\_t\ =\ \textbackslash{}alpha\^{}2\ u\_\{xx\},\ \&\&\ 0\ \textless{}\ x\ \textless{}\ 1,\textbackslash{}\ 0\ \textless{}\ t\ \textless{}\ \textbackslash{}infty,\ \textbackslash{}\textbackslash{}}
\texttt{(\textbackslash{}text\{ГУ\})\textbackslash{}quad\ \&\ u(0,t)\ =\ 0,\textbackslash{}quad\ u\_x(1,t)\ =\ 1,\ \&\&\ 0\ \textless{}\ t\ \textless{}\ \textbackslash{}infty,\ \textbackslash{}\textbackslash{}}
\texttt{(\textbackslash{}text\{НУ\})\textbackslash{}quad\ \&\ u(x,0)\ =\ \textbackslash{}sin(\textbackslash{}pi\ x),\ \&\&\ 0\ \textbackslash{}leqslant\ x\ \textbackslash{}leqslant\ 1.}
\texttt{\textbackslash{}end\{aligned\}}

\begin{verbatim}
Чи можете ви графічно показати рішення цієї проблеми в різні моменти часу? Чи буде це рішення, як правило, стаціонарним? Це очевидно?

3. Яку фізичну інтерпретацію ви можете дати проблемі
\end{verbatim}

\texttt{\textbackslash{}begin\{aligned\}}
\texttt{(\textbackslash{}text\{УЧП\})\textbackslash{}quad\ \&\ u\_t\ =\ \textbackslash{}alpha\^{}2\ u\_\{xx\},\ \&\&\ 0\ \textless{}\ x\ \textless{}\ 1,\textbackslash{}\ 0\ \textless{}\ t\ \textless{}\ \textbackslash{}infty,\ \textbackslash{}\textbackslash{}}
\texttt{(\textbackslash{}text\{ГУ\})\textbackslash{}quad\ \&\ u\_x(0,t)\ =\ 0,\textbackslash{}quad\ u\_x(1,t)\ =\ 0,\ \&\&\ 0\ \textless{}\ t\ \textless{}\ \textbackslash{}infty,\ \textbackslash{}\textbackslash{}}
\texttt{(\textbackslash{}text\{НУ\})\textbackslash{}quad\ \&\ u(x,0)\ =\ \textbackslash{}sin(\textbackslash{}pi\ x),\ \&\&\ 0\ \textbackslash{}leqslant\ x\ \textbackslash{}leqslant\ 1.}
\texttt{\textbackslash{}end\{aligned\}}

\begin{verbatim}
Чи можете ви графічно показати рішення цієї проблеми в різні моменти часу?

Що можна сказати про стабільну температуру?

4. Припустимо, що бічна поверхня металевого стрижня не ізольована і має початкову температуру 20°C, але температура одного кінця стає 50°C. Потім стрижень промивають рідиною з температурою 30°C. Як виглядає змішана задача, що відповідає цьому випадку?

# Лекція 4. ВИВЕДЕННЯ РІВНЯННЯ ТЕПЛОПРОВІДНОСТІ

МЕТА ЛЕКЦІЇ: Показати, як можна отримати одномірне рівняння теплопровідності
\end{verbatim}

\texttt{u\_t\ =\ \textbackslash{}alpha\^{}2\ u\_\{xx\}\ +\ f(x,\ t),}

\begin{verbatim}
спираючись лише на фундаментальний закон збереження кількості тепла. Розглядається залежність швидкості теплопередачі від основних термофізичних параметрів: теплопровідності, теплоємності та густини матеріалу. Розглядаються різні модифікації базового рівняння теплопровідності.

У кожній галузі знань існує набір базових принципів, правильність яких очевидна. Усі інші твердження мають бути виведені з цих базових принципів. Звісно, те, що очевидно одній людині, може здатися сумнівним для іншої. Історія будь-якої науки полягає у постійному пошуку більш фундаментальних принципів, що лежать в основі цієї науки. Основні принципи мають бути настільки універсальними, щоб ніхто не сумнівався.

? — Припущення A — Припущення B — Припущення C — Припущення D — ?
Рис. 4.1. Аксіоматичний метод.

Наприклад, одна людина може вважати, що всі факти певної науки можна вивести з одного базового припущення, яке ми позначимо як B. На основі припущення B вона або вона може довести теорему C, що призведе до теореми D, а яка, у свою чергу, призведе до доведення багатьох інших теорем (див. рисунок 4.1). Ось як виглядає прогрес у науці — дослідник може отримувати все більше нових результатів. Найвидатніші фізики, хіміки та біологи або саме на цьому шляху.

З іншого боку, замість доведення нових теорем можна спробувати, якщо можливо, знайти новий принцип (назвемо його A), який є більш фундаментальним за B. Припущення A природно вважається більш фундаментальним, ніж B.

Якщо B можна вивести з A. Отже, якщо можна знайти новий принцип A, тоді межа наших знань зміщується до сфери більш фундаментальних ідей про певну науку. У теорії теплопровідності основним принципом є закон збереження енергії (теплова енергія). Усі інші твердження випливають із цього базового принципу (рис. 4.2).

Закон збереження енергії
\end{verbatim}

\texttt{u\_t\ =\ \textbackslash{}alpha\^{}2\ u\_\{xx\}\ +\ f(x,t)}

\begin{verbatim}
Інші властивості теплопровідності Рис. 4.2. Енергозбереження: наріжний камінь теплопровідності.

Звісно, ми могли б забути цю лекцію і використати рівняння теплопровідності як фундаментальне поняття (дехто може так подумати, якщо це очевидно), але тоді ми обдуримо більш серйозних студентів, адже закон збереження вважається найфундаментальнішим принципом у науці. Науковці, моделюючи свої задачі, часто записують відношення, що виражають закони збереження, а потім переписують їх у вигляді рівнянь у похідних похідних.

Повернемося до головної мети нашої лекції — вивести рівняння теплопровідності з рівняння збереження енергії.

# Виведення рівняння теплопровідності

Розглянемо однорідний стрижень довжини L, щодо якого зробимо такі припущення:

1. Стрижень виготовлений з одного однорідного провідного матеріалу.

![](images/_page_29_Figure_10.jpeg)

Фіг. 4.3. Тонкий тепловий стержень.

2. Бікова поверхня стрижня термічно ізольована (тепло може поширюватися лише вздовж осі x).

3. Стрижень тонка, тобто температура всіх точок у кожному поперечному перерізі стрижня є сталою. Якщо розглядати частину вудки на сегменті [x, x+\Delta x] і використовувати

закон збереження кількості тепла, тоді можна написати:

(4.1) Загальна зміна кількості тепла на сегменті [x, x + \Delta x] = дорівнює загальній кількості тепла, що пройшло через межі + Нульова кількість тепла, утвореного всередині сегмента [x, x + \Delta x] .

Оскільки загальна кількість тепла (у калоріях) всередині сегмента \{x, x + \Delta x\} в будь-який момент часу t обчислюється формулою:

Загальна кількість тепла всередині сегмента [x, x+\Delta x] = \int\limits_{-\infty}^{x+\Delta x} coAu (s, t) ds,

де c — питома теплоємність матеріалу (показує здатність матеріалу накопичувати тепло),

Щільність матеріалу,

A — це площа поперечного перерізу стрижня.

**va**con збереження енергії (4.1) можна надати у такій математичній формі:

(4.2)
\end{verbatim}

\texttt{\textbackslash{}frac\{d\}\{dt\}\ \textbackslash{}int\_\{x\}\^{}\{x+\textbackslash{}Delta\ x\}\ c\textbackslash{}rho\ Au(s,\ t)\ ds\ =\ c\textbackslash{}rho\ A\ \textbackslash{}int\_\{x\}\^{}\{x+\textbackslash{}Delta\ x\}\ u\_\{t\}(s,\ t)\ ds\ =\ eA\ \textbackslash{}left{[}\ u\_\{x\}(x+\textbackslash{}Delta\ x,\ t)\ -\ u\_\{x\}(x,\ t)\ \textbackslash{}right{]}\ +\ A\ \textbackslash{}int\_\{x+\textbackslash{}Delta\ x\}\^{}\{x+\textbackslash{}Delta\ x\}\ f(s,\ t)\ ds.}

\begin{verbatim}
**де** k — теплопровідність матеріалу (вказує на здатність матеріалу проводити тепло), f(x, t) — об'ємна потужність зовнішнього джерела тепла (калорії/см·с).

Задача полягає в тому, щоб записати рівняння (4.2) у формі, яка не містить інтегралів. Щоб розв'язати цю задачу, нагадаємо читачеві теорему про середнє значення з ходу інтегральної нечисельності.

# Теорема про значення

Якщо функція f(x) є неперервною на сегменті [a, b], то існує принаймні одна точка \xi \in [a, b] така, що
\end{verbatim}

\texttt{\textbackslash{}int\_\{a\}\^{}\{b\}\ f(x)\ dx\ =\ f(\textbackslash{}xi)\ (b-a).}

\begin{verbatim}
Застосовуючи цей результат до рівняння (4.2), ми отримуємо наступний результат
\end{verbatim}

\texttt{coA\textbackslash{},u\_t(\textbackslash{}xi,\ t)\textbackslash{},\textbackslash{}Delta\ x\ =\ kA\textbackslash{}left{[}u\_x(x\ +\ \textbackslash{}Delta\ x,\ t)\ -\ u\_x(x,\ t)\textbackslash{}right{]}\ +\ Af(\textbackslash{}xi,\ t)\textbackslash{},\textbackslash{}Delta\ x,}

\begin{verbatim}
де x < \xi < x + \Delta x, тобто
\end{verbatim}

\texttt{u\_t(\textbackslash{}xi,\ t)\ =\ \textbackslash{}frac\{k\}\{c\textbackslash{}rho\}\textbackslash{}left{[}\textbackslash{}frac\{u\_x(x\ +\ \textbackslash{}Delta\ x,\ t)\ -\ u\_x(x,\ t)\}\{\textbackslash{}Delta\ x\}\textbackslash{}right{]}\ +\ \textbackslash{}frac\{1\}\{c\textbackslash{}rho\}\ f(\textbackslash{}xi,\ t).}

\begin{verbatim}
Спрямовуючи \Delta x до 0, отримуємо потрібне рівняння

(4.3)
\end{verbatim}

\texttt{u\_t(x,\ t)\ =\ \textbackslash{}alpha\^{}2\ u\_\{xx\}(x,\ t)\ +\ F(x,\ t).}

\begin{verbatim}
де
\end{verbatim}

\texttt{\textbackslash{}alpha\^{}\{2\}=\textbackslash{}frac\{k\}\{c\textbackslash{}rho\}}

\begin{verbatim}
— коефіцієнт теплової дифузії, F\left(x,\ t\right)=\frac{1}{c\rho}f\left(x,\ t\right) — густина джерел тепла.

Ми могли б зупинитися на цьому. Однак ми також розглянемо випадок, коли бічна поверхня стрижня не термічно ізольована. Припустимо, що кількість теплового потоку через бічну поверхню стрижня в цьому випадку пропорційна різниці між температурою стрижня u(x,t) і температурою навколишнього середовища, яка залишається сталою і дорівнює нулю. У цьому випадку закон збереження кількості тепла веде до рівняння
\end{verbatim}

\texttt{(4.4)\ u\_t\ =\ \textbackslash{}alpha\^{}2\ u\_\{xx\}\ -\ \textbackslash{}beta\ u\ +\ F(x,\ t),}

\begin{verbatim}
де \beta — коефіцієнт пропорційності потоку через бічну поверхню.

# ЗАУВАЖЕННЯ

1. Як зазначалося вище, константа k називається теплопровідністю речовини і чисельно дорівнює кількості тепла (у калоріях), що проходить за секунду через пластину товщиною 1 см і площею поперечного перерізу 1 см³ при різниці температур на протилежних гранях пластини рівною 1°C. Числові значення теплопровідності для різних матеріалів можна знайти в довідниках з фізики або хімії. Типові значення теплопровідності варіюються від 1 (для міді) до значень, близьких до нуля для хороших теплоізоляторів. Якщо стрижень зроблений із однорідного матеріалу, то значення k не залежить від x. Для деяких матеріалів значення k залежить від температури, і тому рівняння теплопровідності в цьому випадку стає нелінійним:
\end{verbatim}

\texttt{u\_\{t\}\ =\ \textbackslash{}frac\{1\}\{c\textbackslash{}rho\}\ \textbackslash{}frac\{\textbackslash{}partial\}\{\textbackslash{}partial\ x\}\textbackslash{}left(k(u)\textbackslash{},u\_\{x\}\textbackslash{}right).}

\begin{verbatim}
Однак найчастіше теплопровідність k дуже слабо впливає на температуру u, і цю нелінійність можна ігнорувати.

2. Значення c зазвичай називають питомою теплоємністю речовини. Він характеризує здатність речовини зберігати теплову енергію. Наприклад, запечена картопля має високу теплоємність, оскільки може зберігати велику кількість тепла на одиницю маси (саме тому вона зберігає високу температуру так довго). Чисельна теплоємність дорівнює кількості калорій тепла, які потрібно передати 1 граму речовини, щоб підвищити її температуру на 1°C. Для більшості розглянутих задач значення c можна вважати сталим, не записуючи ні x, ні i. Значення питомої теплоємності для конкретних матеріалів можна знайти у фізичних довідниках.

3. Одиниці вимірювання основних величин, включених у рівняння

теплопровідність (у системі GHS) наступна:

i-температура (градуси Цельсія), u_t — швидкість зміни температури (°C/s), u_x — нахил температурної кривої (°C/см), u_{xx} — опуклість температурної кривої (°C/см²),

c = щільність тепла (калор/г \cdot °C)

k — теплопровідність (кал/см \cdot при \cdot °C), \rho — щільність (г/<sup>см3</sup>), \alpha^2 — це теплова дифузія (см<sup>²</sup>/с).

4. Слід зазначити, що теплова дифузія матеріалу \alpha^2 = \frac{k}{\epsilon \rho} прямо пропорційна теплопровідності матеріалу і обернено пропорційна густині \rho та питому теплоємності \epsilon Ми сподіваємося, що цей зв'язок відповідає інтуїтивним ідеям читача.

# ЗАВДАННЯ

1. Підставити розмірність кожної з u,\ u_t,\ \dots величин у рівняння
\end{verbatim}

\texttt{u\_t\ =\ \textbackslash{}alpha\^{}2\ u\_\{xx\}\ -\ \textbackslash{}beta\ u}

\begin{verbatim}
і переконайтеся, що всі терміни мають однаковий розмір {}^{\circ}\mathrm{C}/c .

2. Підставити розмірність кожної з величин у рівняння
\end{verbatim}

\texttt{u\_t\ =\ \textbackslash{}alpha\^{}2\ u\_\{xx\}\ -\ v\ u\_x,}

\begin{verbatim}
де \upsilon вимірюється в одиницях швидкості, і переконайтеся, що всі терми мають однакову розмірність.

3. Отримаємо рівняння теплопровідності
\end{verbatim}

\texttt{u\_t\ =\ \textbackslash{}frac\{1\}\{c\textbackslash{}rho\}\textbackslash{}frac\{\textbackslash{}partial\}\{\textbackslash{}partial\ x\}\textbackslash{}left{[}h(x)u\_x\textbackslash{}right{]}\ +\ f(x,t).}

\begin{verbatim}
для випадку, коли коефіцієнт теплопровідності залежить від координати x.

4. Припустимо, що u(x,t) — концентрація речовини в потоці, який рухається зі швидкістю v. Якщо концентрація змінюється внаслідок дифузії та конвекції, то відповідне рівняння має вигляд
\end{verbatim}

\texttt{u\_t\ =\ \textbackslash{}alpha\^{}2\ u\_\{xx\}\ -\ v\ u\_x.}

\begin{verbatim}
Виходячи з того, що в будь-який момент часу в певному [x, x + \Delta x] сегменті матерії не з'являється і не зникає.

ПРИМІТКА. Закон збереження матерії виглядає так:

Зміни маси речовини на сегменті [x, x + \Delta x] = = Зміна внаслідок дифузії через межі + Зміна внаслідок переносу через межі.

# Лекція 5. ЗМІННЕ РОЗДІЛЕННЯ

МЕТА ЛЕКЦІЇ: Познайомити читача з потужним методом розділення змінних і показати, як цей метод можна використати для розв'язання добре відомої проблеми дифузії. Оскільки цей метод створює труднощі для деяких учнів (головним чином через громіздкі алгебранічні перетворення), ми супроводжуємо презентацію інтуїтивними міркуваннями.

Основна ідея методу полягає в тому, щоб розкласти початкову умову на найпростіші компоненти, знайти відгук системи на кожен найпростіший компонент і підсумувати всі відповіді. Таким чином можна знайти відповідь на довільну початкову умову.

Існуючий метод поділу змінних певною мірою зменшує можливість такої інтерпретації методу, але все ж не виключає його.

Метод розділення змінних є одним із найвідоміших методів розв'язання зсувних задач і використовується, якщо:

- 1. Рівняння є лінійним і однорідним (не обов'язково з постійними коефіцієнтами).
 - 2. Граничні умови визначаються як
\end{verbatim}

\texttt{\textbackslash{}alpha\ u\_x(0,\ t)\ +\ \textbackslash{}beta\ u(0,\ t)\ =\ 0,\ \textbackslash{}qquad\ \textbackslash{}gamma\ u\_x(1,\ t)\ +\ \textbackslash{}delta\ u(1,\ t)\ =\ 0.}

\begin{verbatim}
де \alpha, \beta, \gamma і \delta — константи (граничні умови, задані в такій формі, називаються лінійними однорідними граничними умовами).

Метод був створений за часів Фур'є (зазвичай відомий як метод Фур'є) і, здається, є найпопулярнішим (де це застосовно).

![](images/_page_34_Figure_8.jpeg)

Рис. 5.1. Схематичне представлення задачі дифузії.

Замість того, щоб вивчати метод у загальному випадку, давайте спочатку розглянемо конкретну задачу (загальний випадок обговоримо пізніше). Розглянемо змішану задачу типу дифузії: щоб знайти розв'язок
\end{verbatim}

\texttt{\textbackslash{}begin\{aligned\}}
\texttt{(\textbackslash{}text\{УЧП\})\textbackslash{}quad\ \&\ u\_t\ =\ \textbackslash{}alpha\^{}2\ u\_\{xx\},\ \&\&\ 0\ \textless{}\ x\ \textless{}\ 1,\textbackslash{}quad\ 0\ \textless{}\ t\ \textless{}\ \textbackslash{}infty,\ \textbackslash{}\textbackslash{}}
\texttt{(\textbackslash{}text\{ГУ\})\textbackslash{}quad\ \&\ u(0,t)=0,\textbackslash{}quad\ u(1,t)=0,\ \&\&\ 0\ \textless{}\ t\ \textless{}\ \textbackslash{}infty,\ \textbackslash{}\textbackslash{}}
\texttt{(\textbackslash{}text\{НУ\})\textbackslash{}quad\ \&\ u(x,0)=\textbackslash{}varphi(x),\ \&\&\ 0\ \textbackslash{}le\ x\ \textbackslash{}le\ 1.}
\texttt{\textbackslash{}end\{aligned\}}

\begin{verbatim}
Перш ніж почати розділяти змінні, давайте дамо фізичну інтерпретацію нашої проблеми. Отже, існує стрижень скінченної довжини, кінці якого підтримуються при сталій температурі, рівній нулю (насправді кінці можна підтримувати при значно вищій температурі, значення якої вважається початковою точкою). Додаткові дані про завдання подаються як початкова умова. Наша мета — знайти розподіл температури u\left(x,t\right) в наступний час.

# Загальні принципи методу розділення змінних

Для найпростішого диференціального рівняння в похідних похідних розділення змінних є пошуком розв'язків у вигляді
\end{verbatim}

u(x, t) = X(x)T(t),

\begin{verbatim}
де X(x) — це функція, що залежить лише від змінної x, а T(t) — функція, що залежить лише від t. Це розв'язок у певному сенсі є найпростішим, оскільки температура u(x, t), представлена у такій формі, зберігатиме «форму» профілю у різний час часу t (див. рисунок 5.2).

![](images/_page_35_Figure_6.jpeg)

РИСУНОК 5.2 Графік функції X(x) T(t) у різних точках часу t.

Загальна ідея полягає в тому, щоб знайти нескінченну кількість таких часткових розв'язків, які задовольняють граничні умови
\end{verbatim}

\texttt{u(0,t)=0,\ \textbackslash{}qquad\ u(1,t)=0.}

\begin{verbatim}
Для цього підставляємо розв'язки (5.1) у ці граничні умови. У результаті отримуємо
\end{verbatim}

\texttt{\textbackslash{}begin\{aligned\}}
\texttt{u(0,t)\ \&=\ B\ e\^{}\{-\textbackslash{}lambda\^{}2\ \textbackslash{}alpha\^{}2\ t\}\ =\ 0\ \textbackslash{}Rightarrow\ B\ =\ 0,\ \textbackslash{}\textbackslash{}}
\texttt{u(1,t)\ \&=\ A\ e\^{}\{-\textbackslash{}lambda\^{}2\ \textbackslash{}alpha\^{}2\ t\}\textbackslash{}sin\ \textbackslash{}lambda\ =\ 0\ \textbackslash{}Rightarrow\ \textbackslash{}sin\ \textbackslash{}lambda\ =\ 0.}
\texttt{\textbackslash{}end\{aligned\}}

\begin{verbatim}
Отже,
\end{verbatim}

\texttt{\textbackslash{}lambda\_n\ =\ n\textbackslash{}pi,\ \textbackslash{}qquad\ n\ =\ 1,\ 2,\ \textbackslash{}ldots}

\begin{verbatim}
Зверніть увагу, що друга гранична умова могла б виконуватися і при A=0 , але тоді розв'язок (5.1) був би тривіально нульовим.

![](images/_page_38_Figure_3.jpeg)

Рис. 5.3. Фундаментальне рішення n_n(x, t) = A_n e^{-(\pi n \alpha)^{n/t}} \sin{(\pi n x)} .

Тепер, коли ми зробили другий крок, маємо нескінченний набір ознак

(5.2)
\end{verbatim}

\texttt{u\_n(x,\ t)\ =\ A\_n\ e\^{}\{-(n\textbackslash{}pi\textbackslash{}alpha)\^{}2\ t\}\ \textbackslash{}sin(n\textbackslash{}pi\ x),\ \textbackslash{}quad\ n\ =\ 1,\ 2\ \textbackslash{}dots}

\begin{verbatim}
кожна з яких задовольняє рівняння з частковими похідними та граничними умовами). Отримані функції — це будівельні блоки, з яких ми будемо створювати рішення нашої проблеми. Це рішення буде певною сумою

<sup>2)</sup> Зверніть увагу, що функції u_n і u_{-n} відрізняються лише знаком.

цих найпростіших функцій. Конкретний тип суми залежить від початкового стану. Фіг. Рисунок 5.3 показує графіки кількох таких фундаментальних функцій

КРОК 3. (Пошук розв'язку, який задовольняє рівняння, межу та початкові умови.)

Останній (і, ймовірно, найцікавіший з математичної точки зору) крок — знайти таку суму фундаментальних розв'язків
\end{verbatim}

\texttt{u(x,\ t)\ =\ \textbackslash{}sum\_\{n=1\}\^{}\{\textbackslash{}infty\}\ A\_n\ e\^{}\{-(n\textbackslash{}pi\textbackslash{}alpha)\^{}n\ t\}\ \textbackslash{}sin(n\textbackslash{}pi\ x),}

\begin{verbatim}
тобто при виборі таких коефіцієнтів A_n , що функція задовольнятиме початкову умову
\end{verbatim}

\texttt{u\textbackslash{}left(\ x,\textbackslash{}\ 0\textbackslash{}right)\ =\textbackslash{}varphi\textbackslash{}left(\ x\textbackslash{}right)\ .}

\begin{verbatim}
Підставляючи суму в початковому стані, отримуємо

(5.3)
\end{verbatim}

\texttt{q(x)\ =\ \textbackslash{}sum\_\{n=1\}\^{}\{\textbackslash{}infty\}\ A\_n\ \textbackslash{}sin(n\textbackslash{}pi\ x).}

\begin{verbatim}
Це рівняння підводить нас до цікавого питання: чи можна розкласти початкову температуру \phi(x) на ряд елементарних функцій вигляду?
\end{verbatim}

\texttt{A\_1\ \textbackslash{}sin(\textbackslash{}pi\ x)\ +\ A\_2\ \textbackslash{}sin(2\textbackslash{}pi\ x)\ +\ A\_3\ \textbackslash{}sin(3\textbackslash{}pi\ x)\ +\ \textbackslash{}dots}

\begin{verbatim}
Позитивну відповідь на це питання дав французький математик Жозеф Фур'є. Виявилося, що для достатньо хороших функцій така декомпозиція можлива 
^1 ). Тоді виникає нове питання: знайти коефіцієнти розкладу A_n ?

Насправді це дуже просто. Для цього використаємо властивість системи функцій
\end{verbatim}

\texttt{\textbackslash{}\{\textbackslash{}sin(n\textbackslash{}pi\ x)\textbackslash{}\}\_\{n=1\}\^{}\{\textbackslash{}infty\}.}

\begin{verbatim}
відомий як ортогональність.

У нашому випадку (див. завдання 2) ці функції задовольняють умови
\end{verbatim}

\texttt{\textbackslash{}int\_\{0\}\^{}\{1\}\ \textbackslash{}sin(m\textbackslash{}pi\ x)\ \textbackslash{}sin(n\textbackslash{}pi\ x)\ dx\ =\ \textbackslash{}begin\{cases\}\ 0,\ \&\ m\ \textbackslash{}neq\ n,\ \textbackslash{}\textbackslash{}\ 1/2,\ \&\ m\ =\ n.\ \textbackslash{}end\{cases\}}

\begin{verbatim}
Умови розкладу функцій у рядах Фур'є наведені в лекції 11. Для глибшого ознайомлення з цим питанням можна звернутися до численних освітніх літератур (див., наприклад: Г. П. Толстов, Серія Фур'є. — Москва: Наука, 1980). — Ед. Ед.

Ця властивість чітко показана на рис. 5.4, яка показує графіки деяких функцій із цієї системи.

Отже, ми хочемо знайти коефіцієнти у розкладі
\end{verbatim}

\texttt{\textbackslash{}varphi(x)\ =\ A\_1\ \textbackslash{}sin(\textbackslash{}pi\ x)\ +\ A\_2\ \textbackslash{}sin(2\textbackslash{}pi\ x)\ +\ A\_3\ \textbackslash{}sin(3\textbackslash{}pi\ x)\ +\ \textbackslash{}dots}

\begin{verbatim}
![](images/_page_40_Figure_5.jpeg)

РИСУНОК 5.4. Послідовність ортогональних функцій

Помножте обидві сторони цього співвідношення на \sin{(m\pi x)} (m — довільне ціле число) і інтегруйте від нуля до одиниці. У результаті ми отримуємо
\end{verbatim}

\texttt{\textbackslash{}int\_\{0\}\^{}\{1\}\ \textbackslash{}varphi(x)\ \textbackslash{}sin(m\textbackslash{}pi\ x)\ \textbackslash{},\ dx\ =\ A\_\{m\}\ \textbackslash{}int\_\{0\}\^{}\{1\}\ \textbackslash{}sin\^{}\{2\}(m\textbackslash{}pi\ x)\ \textbackslash{},\ dx\ =\ \textbackslash{}frac\{A\_\{m\}\}\{2\}}

\begin{verbatim}
(усі інші члени зрівнялися з нулем через ортогональність). Розв'язуючи рівняння відносно A_m , отримуємо
\end{verbatim}

\texttt{A\_m\ =\ 2\ \textbackslash{}int\_0\^{}1\ \textbackslash{}varphi(x)\ \textbackslash{}sin(m\textbackslash{}pi\ x)\ dx.}

\begin{verbatim}
Отже, отримано, що розв'язок записується у вигляді

(5.4)
\end{verbatim}

\texttt{u(x,\ t)\ =\ \textbackslash{}sum\_\{n=1\}\^{}\{\textbackslash{}infty\}\ A\_n\ e\^{}\{-(n\textbackslash{}pi\textbackslash{}alpha)\^{}2\ t\}\ \textbackslash{}sin(\textbackslash{}pi\ n\ x),}

\begin{verbatim}
де коефіцієнти A_n визначаються формулами
\end{verbatim}

\texttt{A\_n\ =\ 2\textbackslash{}int\_0\^{}1\ \textbackslash{}varphi(x)\ \textbackslash{}sin(\textbackslash{}pi\ nx)\ dx}

\begin{verbatim}
Ви можете переконатися, що отримане рішення задовольняє всі умови початкової задачі. На цьому завершується третій крок.

Багато студентів розчаровані, коли бачать, наскільки складним виглядає рішення, і майже не намагаються дивитися на нього знову (виглядає дуже погано). Однак це рішення не здаватиметься таким складним, якщо ви витратите час на його аналіз. Насправді, складніша форма розв'язку робить його більш інформативним. Нижче наведено кілька нотаток, які допоможуть вам інтерпретувати це рішення.

# ЗАУВАЖЕННЯ

Зверніть увагу, що єдина різниця між синусоїдальним розкладом функції φ(x) і розв'язком (5.4) — це наявність часового фактора e<sup>-(nπα)<sup>3</sup>1</sup> у кожному члені ряду. Отже, якщо початкова умова має просту декомпозицію, наприклад, вигляду
\end{verbatim}

\texttt{\textbackslash{}varphi(x)\ =\ \textbackslash{}sin(\textbackslash{}pi\ x)\ +\ \textbackslash{}frac\{1\}\{2\}\textbackslash{}sin(3\textbackslash{}pi\ x),}

\begin{verbatim}
тоді рішення можна негайно записати
\end{verbatim}

\texttt{u(x,\ t)\ =\ e\^{}\{-(\textbackslash{}pi\textbackslash{}alpha)\^{}2\ t\}\ \textbackslash{}sin(\textbackslash{}pi\ x)\ +\ \textbackslash{}frac\{1\}\{2\}\ e\^{}\{-(3\textbackslash{}pi\textbackslash{}alpha)\^{}2\ t\}\ \textbackslash{}sin(3\textbackslash{}pi\ x).}

\begin{verbatim}
У цьому випадку очевидно, що якщо ми розкладемо \phi(x) на ряд Фур'є вздовж синусів, то отримуємо
\end{verbatim}

A\_1 = 1,

\begin{verbatim}
A_2 = 0, A_3 = 1/2, A_4 = A_5 = \dots = 0. - 2. Розв'язок (5.4) можна інтерпретувати так: ми представляємо початкову температуру \varphi(x) як суму найпростіших функцій виду A_n \sin(\pi nx) ; кожна така функція породжує відповідь A_n e^{-(\pi n\alpha)^2 t} \sin(\pi nx) . Сумуючи всі такі відповіді, отримуємо розв'язок, який відповідає початковій умові u(x, 0) = \varphi(x) .
- 3. Кожен член у декомпозиції
\end{verbatim}

\texttt{u(x,\ t)\ =\ A\_1\ e\^{}\{-(\textbackslash{}pi\textbackslash{}alpha)\^{}2\ t\}\ \textbackslash{}sin(\textbackslash{}pi\ x)\ +\ A\_2\ e\^{}\{-(2\textbackslash{}pi\textbackslash{}alpha)\^{}2\ t\}\ \textbackslash{}sin(2\textbackslash{}pi\ x)\ +\ \textbackslash{}cdots}

\begin{verbatim}
є функцією від x і t. Зверніть увагу, що внесок членів з великими числами в t>0 дуже малий через множник e^{-(n\pi\alpha)^2 t} . Відповідно, після закінчення досить великого Р и с.

![](images/_page_42_Figure_2.jpeg)

5.5. У задачах дифузії вищі порядки терміни швидко затухають; \phi(x) — початкова температура. Слід зазначити, що по-перше, у розчині зникають високі частоти — по-перше, зникли найбільші коливання міді.

Повний розв'язок приблизно збігається з першим членом, який є півхвилі синусоїди, що з часом згасає (рис. 1). 5.5.)

# ЗАВДАННЯ

1. Показати, що функції
\end{verbatim}

\texttt{u(x,\ t)\ =\ e\^{}\{-\textbackslash{}lambda\^{}2\ \textbackslash{}alpha\textquotesingle{}\ t\}\ \textbackslash{}left{[}\ A\ \textbackslash{}sin(\textbackslash{}lambda\ x)\ +\ B\ \textbackslash{}cos(\textbackslash{}lambda\ x)\ \textbackslash{}right{]}}

\begin{verbatim}
задовольняють рівняння u_l = \alpha^2 u_{xx} для довільних значень A, B і \lambda .

2. Покажи це
\end{verbatim}

\texttt{\textbackslash{}int\_0\^{}1\ \textbackslash{}sin(m\textbackslash{}pi\ x)\ \textbackslash{}sin(n\textbackslash{}pi\ x)\ dx\ =\ \textbackslash{}begin\{cases\}\ 0,\ \&\ m\ \textbackslash{}neq\ n,\ \textbackslash{}\textbackslash{}\ 1/2,\ \&\ m\ =\ n.\ \textbackslash{}end\{cases\}}

\begin{verbatim}
ПРИМІТКА. Використання ідентичності
\end{verbatim}

\texttt{\textbackslash{}sin(mx)\textbackslash{}sin(nx)\ =\ \textbackslash{}frac\{1\}\{2\}{[}\textbackslash{}cos(m-n)x\ -\ \textbackslash{}cos(m+n)x{]}.}

\begin{verbatim}
3. Знайти розклад у рядах Фур'є синусами функції \varphi(x) = 1 на сегменті [0, 1]. Побудуйте перші три терміни розкладу.

Використовуючи результати завдання 3, знайдіть розв'язок наступної змішаної задачі:

(PЧА)
\end{verbatim}

u\_t = u\_\{xx\}

\begin{verbatim}
0 < x < 1 , 
(ГУ)
\end{verbatim}

\texttt{\textbackslash{}begin\{cases\}\ u(0,\ t)\ =\ 0,\ \textbackslash{}\textbackslash{}\ u(1,\ t)\ =\ 0,\ \textbackslash{}end\{cases\}}

\begin{verbatim}
0 < t < \infty , 
(ПУ) u(x, 0) = 1, \quad 0 \le x \le 1 .

(Зверніть увагу, що ця проблема фізично не має значення, оскільки припускається, що температура на кінцях стрижня миттєво знижується з одного до нуля \phi(x) .

5. Знайти розв'язок проблеми 4, якщо початкова умова встановлена у вигляді
\end{verbatim}

\texttt{u(x,\ 0)\ =\ \textbackslash{}sin(2\textbackslash{}pi\ x)\ +\ \textbackslash{}frac\{1\}\{3\}\textbackslash{}sin(4\textbackslash{}pi\ x)\ +\ \textbackslash{}frac\{1\}\{5\}\textbackslash{}sin(6\textbackslash{}pi\ x).}

\begin{verbatim}
6. Розв'яжіть задачу 4, якщо початкова умова u\left( x,\ 0\right) =x-x^{2},\quad 0< x<1. 
# Лекція 6. ПЕРЕТВОРЕННЯ НЕОДНОРІДНИХ ГРАНИЧНИХ УМОВ У ОДНОРІДНІ ГРАНИЧНІ УМОВИ

МЕТА ЛЕКЦІЇ: Показати, як змішану задачу з неоднорідними граничними умовами можна звести до задачі з однорідними граничними умовами.
\end{verbatim}

\texttt{\textbackslash{}begin\{aligned\}}
\texttt{(\textbackslash{}text\{УЧП\})\textbackslash{}quad\ \&\ u\_t\ -\ \textbackslash{}alpha\^{}2\ u\_\{xx\}\ =\ f(x,t),\ \textbackslash{}\textbackslash{}}
\texttt{(\textbackslash{}text\{ГУ\})\textbackslash{}quad\ \&\ \textbackslash{}alpha\_1\ u\_x(0,t)\ +\ \textbackslash{}beta\_1\ u(0,t)\ =\ g\_1(t),\ \textbackslash{}\textbackslash{}}
\texttt{\&\ \textbackslash{}alpha\_2\ u\_x(L,t)\ +\ \textbackslash{}beta\_2\ u(L,t)\ =\ g\_2(t),\ \textbackslash{}\textbackslash{}}
\texttt{(\textbackslash{}text\{НУ\})\textbackslash{}quad\ \&\ u(x,0)\ =\ \textbackslash{}varphi(x).}
\texttt{\textbackslash{}end\{aligned\}}

\begin{verbatim}
Після відповідної заміни невідомої функції ця задача часто зводиться до задачі вигляду
\end{verbatim}

\texttt{\textbackslash{}begin\{aligned\}}
\texttt{U\_t\ -\ \textbackslash{}alpha\^{}2\ U\_\{xx\}\ \&=\ F(x,t),\ \textbackslash{}\textbackslash{}}
\texttt{\textbackslash{}alpha\_1\ U\_x(0,t)\ +\ \textbackslash{}beta\_1\ U(0,t)\ \&=\ 0,\ \textbackslash{}\textbackslash{}}
\texttt{\textbackslash{}alpha\_2\ U\_x(L,t)\ +\ \textbackslash{}beta\_2\ U(L,t)\ \&=\ 0,\ \textbackslash{}\textbackslash{}}
\texttt{U(x,0)\ \&=\ \textbackslash{}overline\{\textbackslash{}varphi\}(x),}
\texttt{\textbackslash{}end\{aligned\}}

\begin{verbatim}
тобто до задачі з нульовими граничними умовами. Якщо F(x,t)=0, нову задачу можна розв'язувати методом розділення змінних; якщо ж F(x,t)\neq 0, доводиться використовувати інші методи, зокрема інтегральні перетворення або розклади за власними функціями.

Щоб метод розділення змінних був застосовним, граничні умови мають бути лінійними та однорідними, тобто мати вигляд

(6.1)
\end{verbatim}

\texttt{\textbackslash{}begin\{aligned\}}
\texttt{\textbackslash{}alpha\_1\ u\_x(0,t)\ +\ \textbackslash{}beta\_1\ u(0,t)\ \&=\ 0,\ \textbackslash{}\textbackslash{}}
\texttt{\textbackslash{}alpha\_2\ u\_x(L,t)\ +\ \textbackslash{}beta\_2\ u(L,t)\ \&=\ 0.}
\texttt{\textbackslash{}end\{aligned\}}

\begin{verbatim}
Мета цієї лекції — пояснити, як задачі з неоднорідними граничними умовами форми

(6.2)
\end{verbatim}

\texttt{\textbackslash{}begin\{aligned\}}
\texttt{(\textbackslash{}text\{УЧП\})\textbackslash{}quad\ \&\ u\_t\ =\ \textbackslash{}alpha\^{}2\ u\_\{xx\},\ \textbackslash{}\textbackslash{}}
\texttt{(\textbackslash{}text\{ГУ\})\textbackslash{}quad\ \&\ \textbackslash{}alpha\_1\ u\_x(0,t)\ +\ \textbackslash{}beta\_1\ u(0,t)\ =\ g\_1(t),\ \textbackslash{}\textbackslash{}}
\texttt{\&\ \textbackslash{}alpha\_2\ u\_x(L,t)\ +\ \textbackslash{}beta\_2\ u(L,t)\ =\ g\_2(t),\ \textbackslash{}\textbackslash{}}
\texttt{(\textbackslash{}text\{НУ\})\textbackslash{}quad\ \&\ u(x,0)\ =\ \textbackslash{}varphi(x),}
\texttt{\textbackslash{}end\{aligned\}}

\begin{verbatim}
можна звести до задач з однорідними граничними умовами.

Почнемо з найпростішого прикладу, коли кінці стрижня підтримуються при сталих температурах.

# Перетворення неоднорідних граничних умов у однорідні граничні умови

Розглянемо проблему поширення тепла в термоізольованому стрижні, кінці якого підтримуються при сталих температурах k_1 і k_2, тобто
\end{verbatim}

\texttt{\textbackslash{}begin\{aligned\}}
\texttt{(\textbackslash{}text\{УЧП\})\textbackslash{}quad\ \&\ u\_t\ =\ \textbackslash{}alpha\^{}2\ u\_\{xx\},\ \&\&\ 0\ \textless{}\ x\ \textless{}\ L,\textbackslash{}\ 0\ \textless{}\ t\ \textless{}\ \textbackslash{}infty,\ \textbackslash{}\textbackslash{}}
\texttt{(6.3)\textbackslash{}\ (\textbackslash{}text\{ГУ\})\textbackslash{}quad\ \&\ u(0,t)\ =\ k\_1,\textbackslash{}quad\ u(L,t)\ =\ k\_2,\ \&\&\ 0\ \textless{}\ t\ \textless{}\ \textbackslash{}infty,\ \textbackslash{}\textbackslash{}}
\texttt{(\textbackslash{}text\{НУ\})\textbackslash{}quad\ \&\ u(x,0)\ =\ \textbackslash{}varphi(x),\ \&\&\ 0\ \textbackslash{}leqslant\ x\ \textbackslash{}leqslant\ L.}
\texttt{\textbackslash{}end\{aligned\}}

\begin{verbatim}
Складність цієї задачі полягає в тому, що неоднорідні граничні умови не дозволяють безпосередньо застосувати метод розділення змінних. Однак очевидно, що при t \to \infty розв'язок має наближатися до стаціонарного лінійного профілю температури. Тому природно шукати розв'язок у вигляді суми стаціонарної та перехідної частин:
\end{verbatim}

\texttt{u(x,t)=\textbackslash{}left{[}k\_1+\textbackslash{}frac\{x\}\{L\}(k\_2-k\_1)\textbackslash{}right{]}+U(x,t).}

\begin{verbatim}
У цьому випадку достатньо знайти перехідну температуру U(x,t).

![](images/_page_45_Figure_7.jpeg)

Рис. 6.1. Розв'язання задачі (6.3) у різні моменти часу:
а) початкова температура u(x,0)=\varphi(x);
б) стаціонарна температура u(x,\infty)=k_1+\frac{x}{L}(k_2-k_1);
в) перехідна температура.

Після цієї заміни задача (6.3) зводиться до нової задачі для функції U(x,t):

(6.4)
\end{verbatim}

\texttt{\textbackslash{}begin\{aligned\}}
\texttt{U\_t\ \&=\ \textbackslash{}alpha\^{}2\ U\_\{xx\},\ \&\&\ 0\ \textless{}\ x\ \textless{}\ L,\ \textbackslash{}\textbackslash{}}
\texttt{U(0,t)\ \&=\ 0,\textbackslash{}quad\ U(L,t)\ =\ 0,\ \&\&\ 0\ \textless{}\ t\ \textless{}\ \textbackslash{}infty,\ \textbackslash{}\textbackslash{}}
\texttt{U(x,0)\ \&=\ \textbackslash{}varphi(x)\ -\ \textbackslash{}left{[}k\_1\ +\ \textbackslash{}frac\{x\}\{L\}(k\_2\ -\ k\_1)\textbackslash{}right{]}\ =\ \textbackslash{}overline\{\textbackslash{}varphi\}(x).}
\texttt{\textbackslash{}end\{aligned\}}

\begin{verbatim}
де \overline{\varphi}(x) — нова, але відома початкова умова.

Це вже задача з однорідними граничними умовами, отже її можна розв'язувати методом розділення змінних. Розв'язок має вигляд

(6.5)
\end{verbatim}

\texttt{U(x,t)=\textbackslash{}sum\_\{n=1\}\^{}\{\textbackslash{}infty\}\ a\_n\ e\^{}\{-(n\textbackslash{}pi\textbackslash{}alpha/L)\^{}2\ t\}\textbackslash{}sin\textbackslash{}!\textbackslash{}left(\textbackslash{}frac\{n\textbackslash{}pi\ x\}\{L\}\textbackslash{}right),}

\begin{verbatim}
де
\end{verbatim}

\texttt{a\_n=\textbackslash{}frac\{2\}\{L\}\textbackslash{}int\_0\^{}L\ \textbackslash{}overline\{\textbackslash{}varphi\}(\textbackslash{}xi)\textbackslash{}sin\textbackslash{}!\textbackslash{}left(\textbackslash{}frac\{n\textbackslash{}pi\ \textbackslash{}xi\}\{L\}\textbackslash{}right)\textbackslash{},d\textbackslash{}xi.}

\begin{verbatim}
Це все для сердечників із фіксованою температурою на кінцях. А як щодо більш реалістичних ПУ з правою стороною, залежною від часу? Основні ідеї тут ті ж, що й у попередній задачі, але вони набагато складніші.

# Перетворення граничних умов, залежних від часу, у нульові

Розглянемо типову задачу
\end{verbatim}

\texttt{\textbackslash{}begin\{aligned\}}
\texttt{(\textbackslash{}text\{PPP\})\textbackslash{}quad\ \&\ u\_t\ =\ \textbackslash{}alpha\^{}2\ u\_\{xx\},\ \&\&\ 0\ \textless{}\ x\ \textless{}\ L,\textbackslash{}quad\ 0\ \textless{}\ t\ \textless{}\ \textbackslash{}infty,\ \textbackslash{}\textbackslash{}}
\texttt{(\textbackslash{}text\{ГУ\})\textbackslash{}quad\ \&\ u(0,t)=g\_1(t),\textbackslash{}quad\ u\_x(L,t)+h\textbackslash{},u(L,t)=g\_2(t),\ \&\&\ 0\ \textless{}\ t\ \textless{}\ \textbackslash{}infty,\ \textbackslash{}\textbackslash{}}
\texttt{(\textbackslash{}text\{НУ\})\textbackslash{}quad\ \&\ u(x,0)=\textbackslash{}varphi(x),\ \&\&\ 0\ \textbackslash{}le\ x\ \textbackslash{}le\ L.}
\texttt{\textbackslash{}end\{aligned\}}

\begin{verbatim}
Щоб перетворити ці граничні умови на нульові умови, ми, після деяких спроб і помилок, зупинилися на такій формі розв'язку:

(6.7)
\end{verbatim}

\texttt{u(x,\ t)=A(t)\textbackslash{}left(1-\textbackslash{}frac\{x\}\{L\}\textbackslash{}right)+B(t)\textbackslash{}frac\{x\}\{L\}+U(x,\ t).}

\begin{verbatim}
де функції A\left(t\right) і B\left(t\right) обираються так, що «квазі-стаціонарна» частина розв'язку (6.7)

(6.8)
\end{verbatim}

\texttt{S(x,\ t)=A(t)\textbackslash{}left(1-\textbackslash{}frac\{x\}\{L\}\textbackslash{}right)+B(t)\textbackslash{}frac\{x\}\{L\}.}

\begin{verbatim}
задовольняє граничні умови початкової задачі. У цьому випадку функція U\left(x,\ t\right) задовольнятиме однорідні граничні умови.

Підставляючи функцію S(x, t) у граничні умови
\end{verbatim}

\texttt{S(0,\ t)=g\_1(t),\ \textbackslash{}qquad\ S\_x(L,\ t)+hS(L,\ t)=g\_2(t).}

\begin{verbatim}
одержуємо два рівняння для A(t) і B(t). У результаті маємо

(6.9)
\end{verbatim}

\texttt{A(t)=g\_1(t),\ \textbackslash{}qquad\ B(t)=\textbackslash{}frac\{g\_1(t)+Lg\_2(t)\}\{1+Lh\}.}

\begin{verbatim}
Отже,
\end{verbatim}

\texttt{u(x,\ t)=g\_1(t)\textbackslash{}left(1-\textbackslash{}frac\{x\}\{L\}\textbackslash{}right)+\textbackslash{}frac\{g\_1(t)+Lg\_2(t)\}\{1+Lh\}\textbackslash{}frac\{x\}\{L\}+U(x,\ t).}

\begin{verbatim}
Якщо підставити ці вирази для u(x,t) у початкову задачу, то отримаємо нову задачу для невідомої функції U(x,t):
\end{verbatim}

\texttt{\textbackslash{}begin\{aligned\}}
\texttt{(\textbackslash{}text\{УЧП\})\textbackslash{}quad\ \&\ U\_t\ =\ \textbackslash{}alpha\^{}2\ U\_\{xx\}\ -\ S\_t(x,t),\ \&\&\ 0\ \textless{}\ x\ \textless{}\ L,\textbackslash{}quad\ 0\ \textless{}\ t\ \textless{}\ \textbackslash{}infty,\ \textbackslash{}\textbackslash{}}
\texttt{(\textbackslash{}text\{ГУ\})\textbackslash{}quad\ \&\ U(0,t)\ =\ 0,\textbackslash{}quad\ U\_x(L,t)\ +\ hU(L,t)\ =\ 0,\ \&\&\ 0\ \textless{}\ t\ \textless{}\ \textbackslash{}infty,\ \textbackslash{}\textbackslash{}}
\texttt{(\textbackslash{}text\{НУ\})\textbackslash{}quad\ \&\ U(x,0)\ =\ \textbackslash{}varphi(x)\ -\ S(x,0),\ \&\&\ 0\ \textbackslash{}leqslant\ x\ \textbackslash{}leqslant\ L.}
\texttt{\textbackslash{}end\{aligned\}}

\begin{verbatim}
Тепер ми стикаємося з новою проблемою однорідних граничних умов (на жаль, рівняння стало гетерогенним). Ми не можемо розв'язати цю задачу методом розділення змінних, але якщо читач погоджується чекати, то за кілька лекцій він ознайомиться з розв'язком цієї задачі методом цілісних перетворень і розкладання за власними функціями.

# ЗАУВАЖЕННЯ .

1. Як зазначалося раніше, мета цієї лекції — навчити читача, як перетворювати задачі з неунірольними граничними умовами на задачі з нульовими граничними умовами. Після цього може статися, що нове диференціальне рівняння в похідних похідних стане однорідним. Якщо це станеться, ми можемо вважати себе щасливцями, оскільки в цьому випадку проблему можна розв'язати методом розділення змінних (див. перший приклад з цієї лекції). Якщо перетворене рівняння виявиться гетерогенним, тоді нову задачу потрібно розв'язати іншим методом.

2. Лінійні неоднорідні граничні умови найзагальної форми
\end{verbatim}

\texttt{\textbackslash{}alpha\_1\ u\_x(0,t)\ +\ \textbackslash{}beta\_1\ u(0,t)\ =\ g\_1(t),\ \textbackslash{}qquad\ \textbackslash{}alpha\_2\ u\_x(L,t)\ +\ \textbackslash{}beta\_2\ u(L,t)\ =\ g\_2(t).}

\begin{verbatim}
Ви також можете перетворити ці умови на нульові граничні умови, використовуючи метод, описаний у другому прикладі. Звісно, нове рівняння при цьому, як правило, буде гетерогенним.

- 3. Деякі методи розв'язку не містять вимог щодо унікальності граничних умов і, отже, не потребують попередніх перетворень. Пізніше, коли ми ознайомимося з перетворенням Лапласа, побачимо, що нульові граничні умови є необов'язковими (у випадку нульових граничних умов розв'язок дещо простіший, ніж в інших випадках).
- 4. Для граничних умов типу
\end{verbatim}

\texttt{u(0,\ t)=g\_1(t),\ \textbackslash{}qquad\ u(L,\ t)=g\_2(t).}

\begin{verbatim}
Метод, описаний у другому прикладі, призводить до наступної форми розв'язання задачі:
\end{verbatim}

\texttt{u(x,\ t)=g\_1(t)+\textbackslash{}frac\{x\}\{L\}\textbackslash{}left{[}g\_2(t)-g\_1(t)\textbackslash{}right{]}+U(x,\ t).}

\begin{verbatim}
# ЗАВДАННЯ

1. Знайти розв'язок наступної змішаної задачі:
\end{verbatim}

\texttt{\textbackslash{}begin\{aligned\}}
\texttt{(\textbackslash{}text\{РЧП\})\textbackslash{}quad\ \&\ u\_t\ =\ \textbackslash{}alpha\^{}2\ u\_\{xx\}\ +\ x\textbackslash{}cos\ t,\ \&\&\ 0\ \textless{}\ x\ \textless{}\ 1,\textbackslash{}quad\ 0\ \textless{}\ t\ \textless{}\ \textbackslash{}infty,\ \textbackslash{}\textbackslash{}}
\texttt{(\textbackslash{}text\{ГУ\})\textbackslash{}quad\ \&\ u(0,t)\ =\ 1,\textbackslash{}quad\ u\_x(1,t)\ +\ hu(1,t)\ =\ 1,\ \&\&\ 0\ \textless{}\ t\ \textless{}\ \textbackslash{}infty,\ \textbackslash{}\textbackslash{}}
\texttt{(\textbackslash{}text\{ПУ\})\textbackslash{}quad\ \&\ u(x,0)\ =\ \textbackslash{}sin(\textbackslash{}pi\ x),\ \&\&\ 0\ \textbackslash{}leqslant\ x\ \textbackslash{}leqslant\ 1.}
\texttt{\textbackslash{}end\{aligned\}}

\begin{verbatim}
попередньо перетворюючи його на задачу з нульовими граничними умовами. Чи відповідає рішення, яке ви отримуєте, вашим задумам щодо цієї проблеми?

2. Перетворіть у задачу з нульовими граничними умовами і розв'яжіть задачу
\end{verbatim}

\texttt{\textbackslash{}begin\{aligned\}}
\texttt{(\textbackslash{}text\{РЧП\})\textbackslash{}quad\ \&\ u\_t\ =\ u\_\{xx\},\ \&\&\ 0\ \textless{}\ x\ \textless{}\ 1,\textbackslash{}quad\ 0\ \textless{}\ t\ \textless{}\ \textbackslash{}infty,\ \textbackslash{}\textbackslash{}}
\texttt{(\textbackslash{}text\{ГУ\})\textbackslash{}quad\ \&\ u(0,t)=0,\textbackslash{}quad\ u(1,t)=1,\ \&\&\ 0\ \textless{}\ t\ \textless{}\ \textbackslash{}infty,\ \textbackslash{}\textbackslash{}}
\texttt{(\textbackslash{}text\{ПУ\})\textbackslash{}quad\ \&\ u(x,0)=x\^{}2,\ \&\&\ 0\ \textbackslash{}le\ x\ \textbackslash{}le\ 1.}
\texttt{\textbackslash{}end\{aligned\}}

\begin{verbatim}
Як виглядатиме розв'язок у різні моменти часу? Чи відповідає він вашим інтуїтивним уявленням? Що таке стаціонарний розв'язок і як виглядає перехідна частина?

3. Завдання
\end{verbatim}

\texttt{\textbackslash{}begin\{aligned\}}
\texttt{(\textbackslash{}text\{УЧП\})\textbackslash{}quad\ \&\ u\_t\ =\ u\_\{xx\},\ \&\&\ 0\ \textless{}\ x\ \textless{}\ 1,\textbackslash{}\ 0\ \textless{}\ t\ \textless{}\ \textbackslash{}infty,\ \textbackslash{}\textbackslash{}}
\texttt{(\textbackslash{}text\{ГУ\})\textbackslash{}quad\ \&\ u\_x(0,t)\ =\ 0,\textbackslash{}quad\ u\_x(1,t)\ +\ h\ u(1,t)\ =\ 1,\ \&\&\ 0\ \textless{}\ t\ \textless{}\ \textbackslash{}infty,\ \textbackslash{}\textbackslash{}}
\texttt{(\textbackslash{}text\{НУ\})\textbackslash{}quad\ \&\ u(x,0)\ =\ \textbackslash{}sin(\textbackslash{}pi\ x),\ \&\&\ 0\ \textbackslash{}leqslant\ x\ \textbackslash{}leqslant\ 1.}
\texttt{\textbackslash{}end\{aligned\}}

\begin{verbatim}
Перетворимо на задачу з нульовими граничними умовами. Чи буде нове рівняння однорідним?
# Лекція 7. РОЗВ'ЯЗАННЯ СКЛАДНІШИХ ЗАДАЧ ШЛЯХОМ РОЗДІЛЕННЯ ЗМІННИХ

МЕТА ЛЕКЦІЇ: Показати, як більш складні проблеми теплопровідності можна розв'язати методом розділення змінних. У цій лекції ми головним чином розглядаємо такі задачі, розв'язання яких дозволить читачеві краще ознайомитися з методом ділення змінних. Сподіваємося, що читач зможе екстраполювати викладені тут ідеї для власних завдань.

У тій же лекції читач ознайомиться з задачею за її власними функціями, або, як її ще називають, проблемою Штурма-Ліувіля. Обговорюються деякі властивості проблеми Штурма-Ліувіля.

Як згадувалося вище, мета цієї лекції — розв'язати змішану задачу шляхом розділення змінних. Можливо, саме таким методом читач вирішить свою проблему. Водночас ми сподіваємося, що він зможе перенести метод, розглянутий у цій лекції, для власних завдань.

Почнемо з одномірної задачі теплопровідності, в якій одна з граничних умов містить похідну.

# Задача теплопровідності з похідною в граничній умові

Розглянемо пристрій, показаний на рис. 7.1. Припустимо, що верхній кінець стрижня підтримується при нульовій температурі, тобто u(0, t) = 0, а нижній кінець занурений у воду, температура якої також стала і дорівнює нулю (у цьому випадку — нуль)

![](images/_page_50_Picture_4.jpeg)

Рис. 7.1. Схематичне представлення зсувної задачі.

![](images/_page_50_Figure_6.jpeg)

Рис. 7.2. Природа шахраїв, як бути задоволеними
\end{verbatim}

\texttt{(\textbackslash{}text\{ГУ\})\textbackslash{}quad}
\texttt{\textbackslash{}begin\{cases\}} u(0,t)=0,\textbackslash{}
u\_x(1,t)=-h,u(1,t). \texttt{\textbackslash{}end\{cases\}}

\begin{verbatim}
відповідає певній еталонній температурі). Теплообмін між стрижнем і водою відбувається згідно із законом Ньютона, який записується у x=1 як наступна гранична умова:
\end{verbatim}

u\_x(1, t) = -hu(1, t).

\begin{verbatim}
Тепер припустимо, що початкова температура стрижня дорівнює u(x, 0) = x, але одразу при t > 0 включаємо граничні умови. Щоб визначити температуру в стрижні, необхідно розв'язати проблему зміщення
\end{verbatim}

\texttt{\textbackslash{}begin\{aligned\}}
\texttt{(\textbackslash{}text\{УЧП\})\textbackslash{}quad\ \&\ u\_t\ =\ \textbackslash{}alpha\^{}2\ u\_\{xx\},\ \&\&\ 0\ \textless{}\ x\ \textless{}\ 1,\textbackslash{}\ 0\ \textless{}\ t\ \textless{}\ \textbackslash{}infty,\ \textbackslash{}\textbackslash{}}
\texttt{(7.1)\textbackslash{}\ (\textbackslash{}text\{ГУ\})\textbackslash{}quad\ \&\ u(0,t)\ =\ 0,\textbackslash{}quad\ u\_x(1,t)\ +\ h\ u(1,t)\ =\ 0,\ \&\&\ 0\ \textless{}\ t\ \textless{}\ \textbackslash{}infty,\ \textbackslash{}\textbackslash{}}
\texttt{(\textbackslash{}text\{НУ\})\textbackslash{}quad\ \&\ u(x,0)\ =\ x,\ \&\&\ 0\ \textbackslash{}leq\ x\ \textbackslash{}leq\ 1.}
\texttt{\textbackslash{}end\{aligned\}}

\begin{verbatim}
Щоб використати метод розділення змінних, розбиємо процес розв'язання на кілька етапів.

**КРОК** 1. (Перетворення диференціального рівняння в похідних у два звичайних диференціальних рівняння)

Підставляючи u(x,t)=X(x)T(t) у рівняння (7.1), отримуємо
\end{verbatim}

\texttt{XT\textquotesingle{}=\textbackslash{}alpha\^{}2X\textquotesingle{}\textquotesingle{}T.}

\begin{verbatim}
Поділивши обидві частини на \alpha^2XT, маємо
\end{verbatim}

\texttt{\textbackslash{}frac\{T\textquotesingle{}\}\{\textbackslash{}alpha\^{}2\ T\}\ =\ \textbackslash{}frac\{X\textquotesingle{}\textquotesingle{}\}\{X\}\ .}

\begin{verbatim}
Оскільки ліва частина останнього рівняння залежить лише від часу, а права — лише від x (і оскільки x і t незалежні), обидві сторони рівняння мають бути константами. Прирівнюючи обидві частини до однієї сталої \mu , отримуємо

(7.2)
\end{verbatim}

\texttt{T\textquotesingle{}\ -\ \textbackslash{}mu\ \textbackslash{}alpha\^{}\{2\}\ T\ =\ 0,\ X\textquotesingle{}\textquotesingle{}\ -\ \textbackslash{}mu\ X\ =\ 0.}

\begin{verbatim}
Процес розщеплення змінних завершено.

КРОК 2. (Визначення константи відокремлення)

По-перше, константа і не повинна бути додатною, інакше функція T\left( t\right) буде зростати експоненційно, що призведе до нескінченного підвищення температури u=XT. Очевидно, що таке рішення не має фізичного значення.

По-друге, припустимо, що \mu = 0 . Тоді
\end{verbatim}

X'\,' = 0

\begin{verbatim}
і, отже.
\end{verbatim}

\texttt{X(x)\ =\ A\ +\ \textbackslash{}hat\{B\}x.}

\begin{verbatim}
Однак, оскільки граничні умови вимагають, що
\end{verbatim}

\texttt{u(0,\ t)\ =\ X(0)\ T(t)\ =\ 0,\ \textbackslash{}qquad\ u(1,t)=\textbackslash{}bigl(X\textquotesingle{}(1)+hX(1)\textbackslash{}bigr)T(t)=0.}

\begin{verbatim}
Ми одразу дійшли висновку, що
\end{verbatim}

\texttt{X(0)=0\textbackslash{}Rightarrow\ A=0,\textbackslash{}qquad\ X\textquotesingle{}(1)+hX(1)=0\textbackslash{}Rightarrow\ B=0.}

\begin{verbatim}
Це означає, що u(x, t) \equiv 0 . Іншими словами, лише нульове розв'язання відповідає нульовій константі розділення. Ми відкинемо його, оскільки нас цікавлять ненульові розв'язки.

Якщо, нарешті, \mu < 0 , можна поставити \mu = -\lambda^2 і переписати рівняння (7.2) у вигляді
\end{verbatim}

\texttt{T\textquotesingle{}\ +\ \textbackslash{}lambda\^{}2\ \textbackslash{}alpha\^{}2\ T\ =\ 0,\ \textbackslash{}qquad\ X\textquotesingle{}\textquotesingle{}\ +\ \textbackslash{}lambda\^{}2\ X\ =\ 0.}

\begin{verbatim}
<sup>1)</sup> Розв'язок X(x) відповідної задачі з граничним значенням буде лише тривіальним, X(x) = 0. 
Загальні розв'язки цих рівнянь легко знайти:
\end{verbatim}

\texttt{T(t)\ =\ Ae\^{}\{-(\textbackslash{}lambda\ \textbackslash{}alpha)\^{}2\ t\},}
\texttt{\textbackslash{}qquad\ X(x)\ =\ B\textbackslash{}sin(\textbackslash{}lambda\ x)\ +\ C\textbackslash{}cos(\textbackslash{}lambda\ x).}

\begin{verbatim}
Отже, будь-яка функція провідника

(7.3)
\end{verbatim}

\texttt{u(x,\ t)\ =\ e\^{}\{-(\textbackslash{}lambda\ \textbackslash{}alpha)\^{}2\ t\}\ {[}A\ \textbackslash{}sin(\textbackslash{}lambda\ x)\ +\ B\ \textbackslash{}cos(\textbackslash{}lambda\ x){]}}

\begin{verbatim}
при будь-яких значеннях \lambda A і B задовольняє рівняння (читач може перевірити це самостійно простою заміною).

Тепер нам потрібно з'ясувати, які з цих функцій задовольняють граничним умовам

(7.4)
\end{verbatim}

u(0, t) = 0, u\_x(1, t) + hu(1, t) = 0.

\begin{verbatim}
Підстановка розв'язків виду (7.3) у граничні умови (7.4) дозволяє знайти відношення, яким мають задовольняти значення \lambda , \Lambda та B:
\end{verbatim}

\texttt{Be\^{}\{-(\textbackslash{}lambda\textbackslash{}alpha)\^{}2\ t\}\ =\ 0\ \textbackslash{}Rightarrow\ B\ =\ 0,}

\begin{verbatim}
A\lambda e^{-(\lambda\alpha)^2 t} \cos \lambda + hAe^{-(\lambda\alpha)^2 t} \sin \lambda = 0. Після простих алгебранхічних перетворень останнє рівняння зводиться до вигляду

\tan \lambda = -\lambda/h
Іншими словами, щоб визначити \lambda , потрібно знайти точки перетину графів функцій tg \lambda і — \lambda/h (рисунок 7.3).

Числові значення значень \lambda_1 , \lambda_2 , ... можна знайти за допомогою комп'ютера, якщо задано значення h Р и с Р

![](images/_page_52_Figure_15.jpeg)

ис 7.3. Перетин графів функцій tg(\lambda) і -\lambda/h .

викликати власні значення задачі граничного значення

(7.5)
\end{verbatim}

\texttt{X\textquotesingle{}\textquotesingle{}+\textbackslash{}lambda\^{}\{2\}X=0,\textbackslash{}qquad}
\texttt{\textbackslash{}begin\{cases\}} X(0)=0,\textbackslash{}
X'(1)+hX(1)=0. \texttt{\textbackslash{}end\{cases\}}

\begin{verbatim}
Іншими словами, з цими значеннями параметра \lambda існують нетривіальні розв'язки задачі (7.5). Власні значення задачі (7.5) визначаються як корені рівняння
\end{verbatim}

\texttt{\textbackslash{}tan\textbackslash{}lambda\ =\ -\textbackslash{}frac\{\textbackslash{}lambda\}\{h\}.}

\begin{verbatim}
Їхні числові значення знаходилися на комп'ютері для випадку h=1. Таблиця 7.1 показує результати розрахунків перших кількох власних значень.

| n | \lambda_B | rt . | \lambda_n |
|-------------|----------------------|--------|----------------|
| 1<br>2<br>3 | 2,02<br>4,91<br>7,98 | 4<br>5 | 11,08<br>14,20 |

Розв'язки задачі (7.5), що відповідають власним значенням \lambda_n , називаються власними функціями X_n(x) . Для розглянутого завдання власні функції виглядають так:
\end{verbatim}

\texttt{X\_n(x)\ =\ \textbackslash{}sin(\textbackslash{}lambda\_n\ x)}

\begin{verbatim}
зображені на рис. 7.4.

![](images/_page_53_Figure_10.jpeg)

Рис. 7.4. Власні функції X_n(x) завдання (7.5) при h=1; Усі власні функції задовольняють граничні умови.

КРОК 3  (Пошук фундаментальних рішень)

Тепер у нас є нескінченний набір функцій (фундаментальних розв'язків) у нашому розпорядженні
\end{verbatim}

\texttt{u\_n(x,\ t)\ =\ X\_n(x)\ T\_n(t)\ =\ e\^{}\{-(\textbackslash{}lambda\_n\textbackslash{}alpha)\^{}2\ t\}\ \textbackslash{}sin(\textbackslash{}lambda\_n\ x),}

\begin{verbatim}
кожна з яких задовольняє рівняння та граничні умови. Залишився ще один крок — загалом
\end{verbatim}

\texttt{u(x,\ t)\ =\ \textbackslash{}sum\_\{n=1\}\^{}\{\textbackslash{}infty\}\ a\_n\ X\_n(x)\ T\_n(t)}
\texttt{=\ \textbackslash{}sum\_\{n=1\}\^{}\{\textbackslash{}infty\}\ a\_n\ e\^{}\{-(\textbackslash{}lambda\_n\ \textbackslash{}alpha)\^{}2\ t\}\textbackslash{}sin(\textbackslash{}lambda\_n\ x)}

\begin{verbatim}
Вибрати коефіцієнти a_n так, щоб функція u(x, t) задовольняла початкову умову u(x, 0) = x (у цьому випадку функція u(x, t) задовольнятиме рівняння та граничні умови, оскільки, по-перше, кожен член суми їх задовольняє, а по-друге, і рівняння, і граничні умови є лінійними та однорідними).

Отже, ймовірність a_n слід визначати з цієї умови

(7 6)
\end{verbatim}

\texttt{u(x,\ 0)\ =\ x\ =\ \textbackslash{}sum\_\{n=1\}\^{}\{\textbackslash{}infty\}\ a\_n\ \textbackslash{}sin(\textbackslash{}lambda\_n\ x).}

\begin{verbatim}
Ми маємо зробити останній крок.

КРОК 4. (Розклад початкової умови на ряд за власними функціями)

Щоб знайти коефіцієнти a_n у розкладі (7,6) на власні функції, помножте обидві частини цього співвідношення на \sin(\lambda_m x) і інтегруйте від 0 до 1. У результаті ми отримуємо
\end{verbatim}

\texttt{\textbackslash{}int\_\{0\}\^{}\{1\}\ x\ \textbackslash{}sin(\textbackslash{}lambda\_m\ x)\textbackslash{},dx\ =\ \textbackslash{}sum\_\{n=1\}\^{}\{\textbackslash{}infty\}\ a\_\{n\}\ \textbackslash{}int\_\{0\}\^{}\{1\}\ \textbackslash{}sin(\textbackslash{}lambda\_\{n\}\ \textbackslash{}xi)\ \textbackslash{}sin(\textbackslash{}lambda\_\{m\}\ \textbackslash{}xi)\ d\textbackslash{}xi}
\texttt{=\ a\_\{m\}\ \textbackslash{}int\_\{0\}\^{}\{1\}\ \textbackslash{}sin\^{}\{2\}(\textbackslash{}lambda\_\{m\}\ \textbackslash{}xi)\ d\textbackslash{}xi}
\texttt{=\ a\_\{m\}\ \textbackslash{}left(\ \textbackslash{}frac\{\textbackslash{}lambda\_\{m\}\ -\ \textbackslash{}sin\ \textbackslash{}lambda\_\{m\}\ \textbackslash{}cos\ \textbackslash{}lambda\_\{m\}\}\{2\textbackslash{}lambda\_\{m\}\}\ \textbackslash{}right).}

\begin{verbatim}
Розв'язавши останнє відношення відносно a_m і замінивши m на n, отримуємо бажаний результат:

(7.7)
\end{verbatim}

\texttt{a\_n\ =\ \textbackslash{}frac\{2\textbackslash{}lambda\_n\}\{(\textbackslash{}lambda\_n\ -\ \textbackslash{}sin\ \textbackslash{}lambda\_n\ \textbackslash{}cos\ \textbackslash{}lambda\_n)\}\ \textbackslash{}int\_0\^{}1\ x\ \textbackslash{}sin(\textbackslash{}lambda\_n\ x)\textbackslash{},dx.}

\begin{verbatim}
Таблиця 7.2.
# Шанси a_n у формулі (7.8)

| n | an | n | a_n |
|-------------|-----------------------|---------------|----------------|
| 1<br>2<br>3 | 0,24<br>0,22-0,03 <br>| <b>4</b><br>5 | -0.11-0.09 <br>|

Отже, розв'язок задачі (7.1) записується у вигляді

(7.8)
\end{verbatim}

\texttt{u(x,\ t)\ =\ \textbackslash{}sum\_\{n=1\}\^{}\{\textbackslash{}infty\}\ a\_n\ e\^{}\{-(\textbackslash{}lambda\_n\ \textbackslash{}alpha)\^{}2\ t\}\ \textbackslash{}sin(\textbackslash{}lambda\_n\ x),}

\begin{verbatim}
де коефіцієнти a_n задаються формулами (7.7). Перші п'ять коефіцієнтів a_n нашої задачі були знайдені чисельно. Їхні значення наведені в Таблиці 1. 7.2.

![](images/_page_55_Figure_8.jpeg)

Рис. 7.5. Графічне представлення функції (7.8); межа умовно задовольняється у всі моменти часу.

Відповідно, перші три члени розв'язання змішаної задачі

(7.9)
\end{verbatim}

\texttt{\textbackslash{}begin\{aligned\}}
\texttt{(\textbackslash{}text\{УЧП\})\textbackslash{}quad\ \&\ u\_t=u\_\{xx\},\ \&\&\ 0\textless{}x\textless{}1,\textbackslash{}\ 0\textless{}t\textless{}\textbackslash{}infty,\textbackslash{}\textbackslash{}}
\texttt{(\textbackslash{}text\{ГУ\})\textbackslash{}quad\ \&\ u(0,t)=0,\textbackslash{}quad\ u\_x(1,t)+h\textbackslash{},u(1,t)=0,\ \&\&\ 0\textless{}t\textless{}\textbackslash{}infty,\textbackslash{}\textbackslash{}}
\texttt{(\textbackslash{}text\{НУ\})\textbackslash{}quad\ \&\ u(x,0)=x,\ \&\&\ 0\textbackslash{}le\ x\textbackslash{}le1.}
\texttt{\textbackslash{}end\{aligned\}}

\begin{verbatim}
мають вигляд
\end{verbatim}

\texttt{u(x,\ t)\ =\ =\ 0.24e\^{}\{-4t\}\ \textbackslash{}sin(2.02x)\ +\ 0.22e\^{}\{-24t\}\ \textbackslash{}sin(4.9x)\ -\ 0.03e\^{}\{-48.3t\}\ \textbackslash{}sin(7.98x)\ +\ \textbackslash{}dots}

\begin{verbatim}
На рисунку 7-5 показані графіки цього розв'язку для кількох точок часу. Рекомендуємо читачеві відповісти на такі питання: Чи відповідає тип рішення, яке ви отримуєте, з вашими інтуїтивними ідеями? Чи виконуються граничні умови?

# ЗАУВАЖЕННЯ

Задача власних значень (7.5) є особливим випадком більш загальної задачі, яку зазвичай називають проблемою Штурма-Ліувіля:

(7.10)
\end{verbatim}

\texttt{\textbackslash{}begin\{aligned\}}
\texttt{(\textbackslash{}text\{ОДУ\})\textbackslash{}quad\ \&\ {[}p(x)y\textquotesingle{}{]}\textquotesingle{}\ -\ q(x)y\ +\ \textbackslash{}lambda\ r(x)y\ =\ 0,\textbackslash{}qquad\ 0\textless{}x\textless{}1,\textbackslash{}\textbackslash{}}
\texttt{(\textbackslash{}text\{ГУ\})\textbackslash{}quad\ \&\ \textbackslash{}alpha\_1\ y(0)\ +\ \textbackslash{}beta\_1\ y\textquotesingle{}(0)=0,\textbackslash{}\textbackslash{}}
\texttt{\&\ \textbackslash{}alpha\_2\ y(1)\ +\ \textbackslash{}beta\_2\ y\textquotesingle{}(1)=0.}
\texttt{\textbackslash{}end\{aligned\}}

\begin{verbatim}
Якщо розв'язати задачу з однорідними граничними умовами методом розділення змінних, то функції X(x) завжди знаходяться як розв'язки деякої задачі Штурма–Ліувіля. Очевидно, що задача власних значень (7.5) є особливим випадком задачі (7.10).

Перелічимо основні властивості задачі (7.10), які були встановлені Штурмом і Ліувілем за певних природних обмежень.

функції p(x), q(x) та r(x).

1. Існує нескінченна послідовність власних значень, що задовольняють нерівності.

\lambda_1 < \lambda_2 < \lambda_3 < \ldots < \lambda_n < \ldots \to \infty
2. Кожне власне значення \lambda_n відповідає одній (точній до константного множника) ненульовій власній функції y_n(x) .
3. Якщо y_n(x) і y_m(x) дві різні власні функції, що відповідають різним власним значенням \lambda_n і \lambda_m , то вони є ортогональними з вагою r(x) на сегменті [0, 1], тобто вони задовольняють це відношення
\end{verbatim}

\texttt{\textbackslash{}int\_\{0\}\^{}\{1\}\ r(x)\ y\_\{n\}(x)\ y\_\{m\}(x)\ dx\ =\ 0.}

\begin{verbatim}
Більше деталей про проблему Штурма-Ліувіля можна знайти в літературі, наведеній наприкінці книги.
<sup>1)</sup> Див., наприклад, Тіхоног А. Н., Самарський А. А. Рівняння математичної фізики. — Москва: Наука, 1972. — ред. Ед.

# ЗАВДАННЯ
1. Розв'язати наступну задачу для рівняння теплопровідності методом розкладу
(PЧП)
\end{verbatim}

\texttt{u\_1\ =\ u\_\{xx\},\ \textbackslash{}quad\ 0\ \textless{}\ x\ \textless{}\ 1,\ \textbackslash{}quad\ 0\ \textless{}\ t\ \textless{}\ \textbackslash{}infty,}

\begin{verbatim}
(ГУ)
\end{verbatim}

\texttt{\textbackslash{}begin\{cases\}\ u(0,\ t)\ =\ 0,\ \textbackslash{}\textbackslash{}\ u\_x(1,\ t)\ =\ 0,\ \textbackslash{}\textbackslash{}\ u(x,\ 0)\ =\ x,\ \textbackslash{}quad\ 0\ \textbackslash{}le\ x\ \textbackslash{}le\ 1.\ \textbackslash{}end\{cases\}}

\begin{verbatim}
Чи відповідає рішення, яке ви знаходите, з інтонативними представленнями? Що таке стаціонарне рішення?
2. Знайдіть власні цінності та функції проблеми Ассолта-Ліувіля
(ОДУ)
\end{verbatim}

\texttt{X\textquotesingle{}\ +\ \textbackslash{}lambda\ X\ =\ 0,\ \textbackslash{}quad\ 0\ \textless{}\ x\ \textless{}\ 1,}

\begin{verbatim}
(ГУ)
\end{verbatim}

\texttt{\textbackslash{}begin\{cases\}\ X(0)\ =\ 0,\ \textbackslash{}\textbackslash{}\ X\textquotesingle{}(1)\ =\ 0.\ \textbackslash{}end\{cases\}}

\begin{verbatim}
Задайте специфічні вирази для функцій p(x), q(x) та r(x) із загальної задачі Штурма–Ліувіля для цього випадку.
3. Розв'язати задачу теплопровідності для планки з термоізольованими межами
(РЧП)
\end{verbatim}

\texttt{u\_t\ =\ u\_\{xx\},\ \textbackslash{}quad\ 0\ \textless{}\ x\ \textless{}\ 1,\ \textbackslash{}quad\ 0\ \textless{}\ t\ \textless{}\ \textbackslash{}infty,}

\begin{verbatim}
(ГУ)
\end{verbatim}

\texttt{\textbackslash{}begin\{cases\}\ u\_x\ (0,\ t)\ =\ 0,\ \textbackslash{}\textbackslash{}\ u\_x\ (1,\ t)\ =\ 0,\ \textbackslash{}end\{cases\}\ \textbackslash{}quad\ 0\ \textless{}\ t\ \textless{}\ \textbackslash{}infty,}

\begin{verbatim}
(ГУ) u(x, 0) = x, \quad 0 \le x \le. 
1.Чи відповідає рішення, яке ви отримуєте, з вашою інтерпретацією цієї проблеми? Що таке стаціонарне рішення? Яке його фізичне значення?
4. Знайдіть власні цінності та функції завдань
(ОДА)
\end{verbatim}

\texttt{X\textquotesingle{}\textquotesingle{}\ +\ \textbackslash{}lambda\ X\ =\ 0,\ \textbackslash{}quad\ 0\ \textless{}\ x\ \textless{}\ 1,}

\begin{verbatim}
(ГУ)
\end{verbatim}

\texttt{\textbackslash{}begin\{cases\}\ X\textquotesingle{}(0)\ =\ 0,\ \textbackslash{}\textbackslash{}\ X\textquotesingle{}(1)\ =\ 0.\ \textbackslash{}end\{cases\}}

\begin{verbatim}
# Лекція 8. ПЕРЕТВОРЕННЯ СКЛАДНИХ РІВНЯНЬ У ПРОСТІШУ ФОРМУ МЕТА ЛЕКЦІЇ: Показати, як рівняння теплопровідності з додатковими членами можна звести до стандартного рівняння для нової функції w(x,t).

У цій лекції ми розглядаємо два типові приклади:

u_t = \alpha^2 u_{xx} - \beta u,
та
\end{verbatim}

\texttt{u\_t\ =\ \textbackslash{}alpha\^{}2\ u\_\{xx\}\ -\ v\ u\_x,\ \textbackslash{}qquad\ v=\textbackslash{}text\{const\}.}

Обидва рівняння можна перетворити до найпростішого рівняння
теплопровідності

\begin{verbatim}
w_t = \alpha^2 w_{xx}
\end{verbatim}

за допомогою відповідних замін:

\begin{verbatim}
\begin{aligned}
u(x,t)&=e^{-\beta t}w(x,t),\\
u(x,t)&=\exp\!\left[\frac{v}{2\alpha^2}\!\left(x-\frac{vt}{2}\right)\right]w(x,t).
\end{aligned}
\end{verbatim}

Розв'язавши просте рівняння для \texttt{w(x,t)}, ми повертаємося до
початкової функції \texttt{u(x,t)} за цими формулами. Те саме
перетворення, звісно, треба послідовно застосовувати і до початкових, і
до граничних умов.

Після прочитання двох останніх лекцій читач може створити враження, що
метод розділення змінних застосовується \emph{лише} до рівняння

\begin{verbatim}
u_t = \alpha^2 u_{xx}.
\end{verbatim}

Насправді рівняння теплопотоку є найпростішим, але не єдиним
параболічним рівнянням, яке можна розв'язати методом розділення змінних.
Як уже згадувалося, якщо рівняння є лінійним і однорідним, його можна
розв'язати, поділивши змінні. Наприклад, двовимірна задача
теплопровідності в колі описується рівнянням

\begin{verbatim}
u_t = \alpha^2 \left[ u_{rr} + \frac{1}{r} u_r + \frac{1}{r^2} u_{\theta\theta} \right],
\end{verbatim}

І хоча це рівняння з \emph{змінними коефіцієнтами}, його можна звести до
трьох звичайних диференціальних рівнянь. У цій лекції ми покажемо
читачеві, що іноді, замість того, щоб безпосередньо розв'язувати
диференціальне рівняння в похідних, краще перетворити його на простіше
рівняння, яке можна розв'язати методом розділення змінних або іншим
методом. Наступний приклад ілюструє цей підхід.

\section{Перетворення проблеми з перенесенням бічної поверхні у проблему
з термічно ізольованою бічною поверхнею Розглянемо наступне
завдання:}\label{ux43fux435ux440ux435ux442ux432ux43eux440ux435ux43dux43dux44f-ux43fux440ux43eux431ux43bux435ux43cux438-ux437-ux43fux435ux440ux435ux43dux435ux441ux435ux43dux43dux44fux43c-ux431ux456ux447ux43dux43eux457-ux43fux43eux432ux435ux440ux445ux43dux456-ux443-ux43fux440ux43eux431ux43bux435ux43cux443-ux437-ux442ux435ux440ux43cux456ux447ux43dux43e-ux456ux437ux43eux43bux44cux43eux432ux430ux43dux43eux44e-ux431ux456ux447ux43dux43eux44e-ux43fux43eux432ux435ux440ux445ux43dux435ux44e-ux440ux43eux437ux433ux43bux44fux43dux435ux43cux43e-ux43dux430ux441ux442ux443ux43fux43dux435-ux437ux430ux432ux434ux430ux43dux43dux44f}

(8.1)

\begin{verbatim}
\begin{aligned}
u_{t} &= \alpha^{2}u_{xx} - \beta u, \quad 0 < x < 1, \quad 0 < t < \infty,\\
u(0, t) &= 0, \quad u(1, t) = 0, \quad 0 < t < \infty,\\
u(x, 0) &= \varphi(x), \quad 0 \le x \le 1.
\end{aligned}
\end{verbatim}

де термін \texttt{-\textbackslash{}beta\ u} описує потік тепла через
бічну поверхню стрижня (рис. 8.1).

\pandocbounded{\includegraphics[keepaspectratio]{images/_page_59_Figure_7.jpeg}}

Рис. 8.1. Задача теплопровідності, описана рівнянням
\texttt{u\_t=\textbackslash{}alpha\^{}2u\_\{xx\}-\textbackslash{}beta\ u};
член \texttt{-\textbackslash{}beta\ u} відповідає втратам тепла через
бічну поверхню стрижня.
\texttt{Наша\ мета\ —\ ввести\ нову\ температурну\ \textquotesingle{}w\textbackslash{}left(x,\textbackslash{}\ t\textbackslash{}right)\textquotesingle{}\ замість\ \textquotesingle{}u\textbackslash{}left(x,\textbackslash{}\ t\textbackslash{}right)\textquotesingle{}\ ,\ щоб\ рівняння\ для\ \textquotesingle{}w\textbackslash{}left(x,\textbackslash{}\ t\textbackslash{}right)\textquotesingle{}\ стало\ простішим\ за\ початкове\ рівняння}

\begin{verbatim}
u_t = \alpha^2 u_{xx} - \beta u.
\end{verbatim}

Це досить загальний метод розв'язання диференціальних рівнянь у
частинних похідних, але вибір відповідного перетворення базується на
інтуїтивних уявленнях про поведінку розв'язків початкового рівняння.
Наприклад, у задачі (8.1) температура \texttt{u(x,\ t)} у кожній точці
\texttt{x\_0} змінюється внаслідок двох факторів:

\begin{itemize}
\tightlist
\item
  \begin{enumerate}
  \def\labelenumi{\arabic{enumi})}
  \tightlist
  \item
    дифузія тепла вздовж стрижня (через компонентний
    \texttt{\textbackslash{}alpha\^{}2\ u\_\{xx\}} ),
  \end{enumerate}
\item
  \begin{enumerate}
  \def\labelenumi{\arabic{enumi})}
  \setcounter{enumi}{1}
  \tightlist
  \item
    теплопередача через бічну поверхню стрижня (через член
    \texttt{-\textbackslash{}beta\ u}).
  \end{enumerate}
\end{itemize}

Найважливішим елементом нашого міркування є наступне: якщо не було
дифузії вздовж стрижня ( \texttt{\textbackslash{}alpha\ =\ 0} ),

тоді температура в кожній точці стрижня експоненційно знижується до нуля
згідно із законом

\begin{verbatim}
u(x_0,t)=u(x_0,0)e^{-\beta t}.
\end{verbatim}

Виходячи з цього спостереження, запитаємо, чи можна розв'язок проблеми
(8.1) представити як утворення двох факторів

\begin{verbatim}
(8.2)\quad u(x,t)=e^{-\beta t}w(x,t).
\end{verbatim}

де \texttt{w(x,t)} описує розподіл температури, що виникає лише через
дифузію. Після підстановки цього виразу у задачу (8.1) отримуємо
перетворену задачу без члена \texttt{-\textbackslash{}beta\ u}. (8.4)

\begin{verbatim}
\begin{aligned}
w(x,t)&=\sum_{n=1}^{\infty} a_n e^{-(n\pi\alpha)^2 t}\sin(n\pi x),\\
a_n&=2\int_0^1 \varphi(\xi)\sin(n\pi \xi)\,d\xi.
\end{aligned}
\end{verbatim}

і, отже, розв'язок початкової задачі (8.1) має вигляд

\begin{verbatim}
u(x, t) = e^{-\beta t}w(x, t).
\end{verbatim}

Нижче, у розділі «Нотатки», наведено ще один приклад проблеми,
розв'язаної цим методом.

\section{ПРИМІТКИ}\label{ux43fux440ux438ux43cux456ux442ux43aux438}

\begin{enumerate}
\def\labelenumi{\arabic{enumi}.}
\tightlist
\item
  Вирішити проблему
\end{enumerate}

\begin{verbatim}
\begin{aligned}
u_t &= u_{xx} - u, && 0 < x < 1,\ 0 < t < \infty, \\
u(0,t) &= 0,\quad u(1,t) = 0, && 0 < t < \infty, \\
u(x,0) &= \sin(\pi x) + 0.5\sin(3\pi x), && 0 \leqslant x \leqslant 1.
\end{aligned}
\end{verbatim}

За методом, описаним вище, ми - a) тимчасово відкинути конвективний член
u у початковому рівнянні. - b) записати розв'язок змішаної задачі без
конвективного члена

\begin{verbatim}
u(x, t) = e^{-\pi^2 t} \sin(\pi x) + 0.5e^{-(3\pi)^2 t} \sin(3\pi x),
\end{verbatim}

\begin{enumerate}
\def\labelenumi{\alph{enumi})}
\setcounter{enumi}{2}
\tightlist
\item
  помножимо це розв'язання на конвективний фактор
  \texttt{e\^{}\{-\textbackslash{}beta\ t\}=e\^{}\{-t\}} отримаємо
  остаточне розв'язання початкової задачі
\end{enumerate}

\begin{verbatim}
u(x, t) = e^{-t} \left[ e^{-\pi^2 t} \sin(\pi x) + 0.5e^{-(3\pi)^2 t} \sin(3\pi x) \right].
\end{verbatim}

\begin{enumerate}
\def\labelenumi{\arabic{enumi}.}
\setcounter{enumi}{1}
\tightlist
\item
  Рівняння конвективної дифузії
\end{enumerate}

\begin{verbatim}
u_t = \alpha^2 u_{xx} - v u_x
\end{verbatim}

(коефіцієнт і є константою) також можна перетворити у вигляд

\begin{verbatim}
w_t = \alpha^2 w_{xx}.
\end{verbatim}

Відповідне перетворення має вигляд

\begin{verbatim}
u(x, t) = e^{v(x-vt/2)/2\alpha^2}w(x, t).
\end{verbatim}

Це перетворення підкреслює експоненціальний фактор, який повністю
зумовлений рухом середовища. Зверніть увагу, що цей множник описує
хвилю, що поширюється зліва направо зі швидкістю
\texttt{\textbackslash{}upsilon/2} . Читач матиме можливість
скористатися цією трансформацією.

\section{ЗАВДАННЯ}\label{ux437ux430ux432ux434ux430ux43dux43dux44f-2}

\begin{enumerate}
\def\labelenumi{\arabic{enumi}.}
\tightlist
\item
  Розв'язати задачу конвективної дифузії
\end{enumerate}

\begin{verbatim}
\begin{aligned}
(\text{РЧП})\quad & u_t = u_{xx} - u_x, && 0 < x < 1,\ 0 < t < \infty, \\
(\text{ГУ})\quad & u(0,t) = 0,\quad u(1,t) = 0, && 0 < t < \infty, \\
(\text{ПУ})\quad & u(x,0) = e^{x/2}, && 0 \leqslant x \leqslant 1.
\end{aligned}
\end{verbatim}

методом перетворення до простішого рівняння. На що схожий розв'язок?
Можна дати таку інтерпретацію цієї задачі: \texttt{u(x,t)} ---
концентрація деякої речовини в рухомому середовищі, яке рухається зліва
направо зі швидкістю \texttt{v\ =\ 1}; початковий розподіл концентрації
задається функцією \texttt{e\^{}\{x/2\}}, а на межах концентрація
підтримується нульовою. Чи відповідає ваш розв'язок такій інтерпретації?
2. Розв'язати задачу

\begin{verbatim}
\begin{aligned}
(\text{РЧП})\quad & u_t = u_{xx} - u + x, && 0 < x < 1,\ 0 < t < \infty, \\
(\text{ГУ})\quad & u(0,t) = 0,\quad u(1,t) = 1, && 0 < t < \infty, \\
(\text{ГУ})\quad & u(x,0) = 0, && 0 \leqslant x \leqslant 1.
\end{aligned}
\end{verbatim}

\begin{itemize}
\tightlist
\item
  а) перетворивши неоднорідні граничні умови на однорідні;
\item
  б) перетворивши вихідне рівняння на нове рівняння без члена
  \texttt{-u};
\item
  в) розв'язавши задачу, що виникає після цих перетворень.
\end{itemize}

\begin{enumerate}
\def\labelenumi{\arabic{enumi}.}
\setcounter{enumi}{2}
\tightlist
\item
  Розв'язати задачу
\end{enumerate}

\begin{verbatim}
\begin{aligned}
(\text{РЧП})\quad & u_t = u_{xx} - u, && 0 < x < 1,\ 0 < t < \infty, \\
(\text{ГУ})\quad & u(0,t) = 0,\quad u(1,t) = 0, && 0 < t < \infty, \\
(\text{ПУ})\quad & u(x,0) = \sin(\pi x), && 0 \leqslant x \leqslant 1.
\end{aligned}
\end{verbatim}

застосовуючи метод розділення змінних безпосередньо до вихідного
рівняння, не виконуючи жодних попередніх перетворень. Чи збігається
отриманий розв'язок з розв'язком, знайденим раніше після заміни?

\begin{verbatim}
u(x,t) = e^{-t}w(x,t)
\end{verbatim}

\section{Лекція 9. РОЗВ'ЯЗАННЯ НЕОДНОРІДНИХ УЧП МЕТОДОМ РОЗКЛАДУ ЗА
ВЛАСНИМИ ФУНКЦІЯМИ ЦІЛЬ ЛЕКЦІЇ: Навчитися розв'язувати змішані задачі
вигляду}\label{ux43bux435ux43aux446ux456ux44f-9.-ux440ux43eux437ux432ux44fux437ux430ux43dux43dux44f-ux43dux435ux43eux434ux43dux43eux440ux456ux434ux43dux438ux445-ux443ux447ux43f-ux43cux435ux442ux43eux434ux43eux43c-ux440ux43eux437ux43aux43bux430ux434ux443-ux437ux430-ux432ux43bux430ux441ux43dux438ux43cux438-ux444ux443ux43dux43aux446ux456ux44fux43cux438-ux446ux456ux43bux44c-ux43bux435ux43aux446ux456ux457-ux43dux430ux432ux447ux438ux442ux438ux441ux44f-ux440ux43eux437ux432ux44fux437ux443ux432ux430ux442ux438-ux437ux43cux456ux448ux430ux43dux456-ux437ux430ux434ux430ux447ux456-ux432ux438ux433ux43bux44fux434ux443}

\begin{verbatim}
\begin{aligned}
(\text{РЧП})\quad & u_t = \alpha^2 u_{xx} + f(x,t), && 0 < x < 1,\ 0 < t < \infty, \\
(\text{ГУ})\quad & \alpha_1 u_x(0,t) + \beta_1 u(0,t) = 0, && 0 < t < \infty, \\
& \alpha_2 u_x(1,t) + \beta_2 u(1,t) = 0, && 0 < t < \infty, \\
(\text{ПУ})\quad & u(x,0) = \varphi(x), && 0 \leqslant x \leqslant 1.
\end{aligned}
\end{verbatim}

Розв'язок неоднорідного УЧП можна шукати у вигляді ряду

\begin{verbatim}
u(x,t) = \sum_{n=1}^{\infty} T_n(t) X_n(x).
\end{verbatim}

де \texttt{X\_n(x)} --- власні функції відповідної однорідної задачі

\begin{verbatim}
\begin{aligned}
(\text{УЧП})\quad & u_t = \alpha^2 u_{xx}, && 0 < x < 1,\ 0 < t < \infty, \\
(\text{ГУ})\quad & \alpha_1 u_x(0,t) + \beta_1 u(0,t) = 0, && 0 < t < \infty, \\
& \alpha_2 u_x(1,t) + \beta_2 u(1,t) = 0, && 0 < t < \infty, \\
(\text{НУ})\quad & u(x,0) = \varphi(x), && 0 \leqslant x \leqslant 1.
\end{aligned}
\end{verbatim}

А функції \texttt{T\_n(t)} визначаються як розв'язки деяких звичайних
диференціальних рівнянь. На лекції 6 ми ознайомилися з тим, як
перетворювати неоднорідні граничні умови на однорідні. На жаль, рівняння
залишається гетерогенним, і ми стоїмо перед завданням

\begin{verbatim}
\begin{aligned}
u_{t} &= \alpha^{2} u_{xx} + f(x, t), \quad 0 < x < 1, \quad 0 < t < \infty,\\
\alpha_{1} u_{x}(0, t) + \beta_{1} u(0, t) &= 0, \\
\alpha_{2} u_{x}(1, t) + \beta_{2} u(1, t) &= 0, \\
u(x, 0) &= \varphi(x), \quad 0 \le x \le 1.
\end{aligned}
\end{verbatim}

Мета нашої лекції --- опанувати метод розв'язання таких задач. Цей метод
аналогічний методу варіації довільних констант. теорія звичайних
диференціальних рівнянь називається методом розкладу за власними
функціями. Ідея методу надзвичайно проста. Відомо, що розв'язок задачі
(9.1) при \texttt{f(x,\ t)} = 0 (так звана відповідна однорідна задача)
записується як ряд

\begin{verbatim}
u\left(x,\ t\right)=\sum_{n=1}^{\infty}a_{n}e^{-\left(\lambda_{n}\alpha\right)^{2}}/X_{n}\left(x\right),
\end{verbatim}

де \texttt{X\_n(x)} --- власні функції задачі Штурма-Ліувіля, (9.2)

\begin{verbatim}
X'' + \lambda^2 X = 0, \\ \alpha_1 X'(0) + \beta_1 X(0) = 0, \\ \alpha_2 X'(1) + \beta_2 X(1) = 0.
\end{verbatim}

Запитаймо себе, чи можна також шукати розв'язок гетерогенної проблеми
(9.1) у вигляді серії власних функцій, але більш загальної форми

\begin{verbatim}
u(x, t) = \sum_{n=1}^{\infty} T_n(t) X_n(x).
\end{verbatim}

Існують певні фізичні підстави для такого припущення. Насправді терміни
\texttt{f(x,t)} описують джерело тепла, розташоване всередині стрижня.
Це означає, що залежність рішення від часу не буде виражена показниками.

\begin{verbatim}
e^{-(\lambda_n\alpha)^2 t}
\end{verbatim}

Це стосувалося однорідного завдання, але з деякими іншими функціями.

Щоб показати цей метод у дії і не ускладнювати подання деталями, ми
застосуємо його до розв'язання однієї дуже простої задачі.

\section{Построзв'язання методом демаркації за власними
функціями,}\label{ux43fux43eux441ux442ux440ux43eux437ux432ux44fux437ux430ux43dux43dux44f-ux43cux435ux442ux43eux434ux43eux43c-ux434ux435ux43cux430ux440ux43aux430ux446ux456ux457-ux437ux430-ux432ux43bux430ux441ux43dux438ux43cux438-ux444ux443ux43dux43aux446ux456ux44fux43cux438}

Розглянемо гетерогенну проблему

(9.3) (ГУ)

\begin{verbatim}
\begin{aligned} u_t &= \alpha^2 u_{xx} + f(x, t), & 0 < x < 1, & 0 < t < \infty, \\ u(0, t) &= 0, & 0 < t < \infty, \\ u(1, t) &= 0, & 0 < t < \infty, \\ u(x, 0) &= \varphi(x), & 0 \leqslant x \leqslant 1. \end{aligned}
\end{verbatim}

Процес розв'язання задачі подається у вигляді послідовності наступних
кроків. КРОК 1. Основна ідея методу полягає в розкладанні густини
джерела \texttt{f(x,\ t)} на ряд власними функціями

\begin{verbatim}
f(x, t) = f_1(t) X_1(x) + f_2(t) X_2(x) + \dots + f_n(t) X_n(x) + \dots
\end{verbatim}

\texttt{а\ також\ визначити\ реакції\ системи\ \textquotesingle{}u\_n\textbackslash{}left(x,\textbackslash{},t\textbackslash{}right)\textquotesingle{}\ на\ вплив\ кожного\ компонента\ \textquotesingle{}f\_n\textbackslash{}left(t\textbackslash{}right)X\_n\textbackslash{}left(x\textbackslash{}right)\textquotesingle{}\ .\ Якщо\ знайти\ всі\ відповіді,\ рішення\ проблеми\ виглядатиме\ так:}

\begin{verbatim}
u(x, t) = \sum_{n=1}^{\infty} u_n(x, t).
\end{verbatim}

Головна складність цього методу --- розклад густини джерела на
компоненти \texttt{f\_n(t)\textbackslash{},X\_n(x)} . Виявляється, що
\texttt{X\_n(x)} факторами в цій задачі є власними векторами системи
Штурма--Ліувіля, які виникають при розв'язанні відповідної однорідної
задачі методом розділення змінних, а саме

\begin{verbatim}
\begin{array}{c} u_t = \alpha^2 u_{xx} \quad \text{(помітим, що } f(x, \ t) = 0), \\ u(0, \ t) = 0, \\ u(1, \ t) = 0, \\ u(x, \ 0) = \varphi(x). \end{array}
\end{verbatim}

Завдання Штурма-Ліувіля в цьому випадку виглядає так:

\begin{verbatim}
X'' + \lambda^2 X = 0,
\end{verbatim}

\texttt{X(0)\ =\ 0,} \texttt{X(1)\ =\ 0,} а його розв'язки є функціями

\begin{verbatim}
X_n(x) = \sin(n\pi x), \quad n = 1, 2, 3 \dots
\end{verbatim}

Отже, розклад на густину джерела представлений як ряд

\begin{verbatim}
(9.5) \ f(x,t) = f_1(t) \sin(\pi x) + f_2(t) \sin(2\pi x) + \dots + f_n(t) \sin(n\pi x) + \dots
\end{verbatim}

\texttt{\textquotesingle{}\textbackslash{}mathcal\{U\}\textquotesingle{}\ нарешті,\ щоб\ знайти\ функції\ \textquotesingle{}f\_n(t)\textquotesingle{}\ ,\ помножте\ обидві\ сторони\ останнього\ відношення\ на\ \textquotesingle{}\textbackslash{}sin(m\textbackslash{}pi\ x)\textquotesingle{}\ і\ інтегруйте\ від\ 0\ до\ 1\ за\ координатою\ x.}

У результаті ми отримуємо

\begin{verbatim}
\int_{0}^{1} f(x, t) \sin(m\pi x)\,dx = \sum_{n=1}^{\infty} f_{n}(t) \int_{0}^{1} \sin(m\pi x) \sin(n\pi x)\,dx = \frac{1}{2} f_{m}(t)
\end{verbatim}

або (після заміни t на p)

(9.6)

\begin{verbatim}
f_{n}(t) = 2 \int_{0}^{1} f(x, t) \sin(n\pi x)\,dx.
\end{verbatim}

Останнє співвідношення встановлює співвідношення між коефіцієнтами
\texttt{f\_n(t)} і часткою джерела тепла \texttt{f(x,\ t)}.

КРОК 2. (Пошук \texttt{u\_n(x,\ t)\ =\ X\_n(x)\ T\_n(t)} реакції на
вхідний вплив \texttt{f\_n(t)\ X\_n(x)} .)

Замінимо густину джерела на його розклад у рядок

\begin{verbatim}
f(x, t) = \sum_{n=1}^{\infty} f_n(t) \sin(n\pi x)
\end{verbatim}

і намагатися знайти індивідуальні відповіді, тобто функції,
\texttt{T\_\{\textbackslash{}mathbf\{n\}\}(t)} у

\begin{verbatim}
u(x, t) = \sum_{n=1}^{\infty} T_n(t) \sin(n\pi x).
\end{verbatim}

Оскільки розв'язок задачі має вигляд

\begin{verbatim}
u(x, t) = \sum_{n=1}^{\infty} T_n(t) \sin(n\pi x),
\end{verbatim}

потім підставимо цей вираз у початкову задачу

\begin{verbatim}
u_{t} = \alpha^{2} u_{xx} + \sum_{n=1}^{\infty} f_{n}(t) \sin(n\pi x),
\quad u(0, t) = 0, \quad u(1, t) = 0, \quad u(x, 0) = \varphi(x),
\end{verbatim}

У результаті ми отримаємо

\begin{verbatim}
\sum_{n=1}^{\infty} T'_n(t) \sin(n\pi x) = -\alpha^2 \sum_{n=1}^{\infty} (n\pi)^2 T_n(t) \sin(n\pi x) + + \sum_{n=1}^{\infty} f_n(t) \sin(\pi n\lambda),
\end{verbatim}

(9.7)

\begin{verbatim}
\sum_{n=1}^{\infty} T_n(t) \sin 0 = 0
\end{verbatim}

\texttt{(задовольняється\ ідентично),\ \textquotesingle{}\textbackslash{}sum\_\{n=1\}\^{}\{\textbackslash{}infty\}\ T\_n(t)\ \textbackslash{}sin\ (n\textbackslash{}pi)\ =\ 0\textquotesingle{}\ (задовольняється\ ідентично),\ \textquotesingle{}\textbackslash{}sum\_\{n=1\}\^{}\{\textbackslash{}infty\}\ T\_n(0)\ \textbackslash{}sin\ (n\textbackslash{}pi\ x)\ =\ \textbackslash{}phi(x)\textquotesingle{}\ .}

Переписування рівняння та початкової умови

(РЧП)

\begin{verbatim}
\sum_{n=1}^{\infty} [T'_n + (n\pi\alpha)^2 T_n - f_n(t)] \sin(n\pi x) = 0,
\end{verbatim}

(ПУ)

\begin{verbatim}
\sum_{n=1}^{\infty} T_n(0) \sin(n\pi x) = \varphi(x),
\end{verbatim}

можна побачити, що функції \texttt{T\_n(\textbackslash{}ell)} є
розв'язками задачі Коші

(9.8)

\begin{verbatim}
\begin{aligned}
T_n' + (n\pi\alpha)^2 T_n &= f_n(t), \\
T_n(0) &= 2 \int_0^1 \varphi(\xi)\sin(n\pi \xi)\,d\xi = a_n.
\end{aligned}
\end{verbatim}

Усі ці TAC легко розв'язати, а їхні рішення записуються у вигляді

(9.9)

\begin{verbatim}
T_n(t) = a_n e^{-(n\pi\alpha)^2 t} + \int_0^t e^{-(n\pi\alpha)^2 (t-\tau)} f_n(\tau)\,d\tau.
\end{verbatim}

З того місця, де легко отримати повне розв'язання проблеми (9.3):

(9.10)

\begin{verbatim}
\begin{aligned}
u(x, t) &= \sum_{n=1}^{\infty} T_n(t)\sin(n\pi x) \\
&= \sum_{n=1}^{\infty} a_n e^{-(n\pi\alpha)^2 t}\sin(n\pi x) \\
&\quad + \sum_{n=1}^{\infty} \sin(n\pi x)\int_0^t e^{-(n\pi\alpha)^2 (t-\tau)} f_n(\tau)\,d\tau.
\end{aligned}
\end{verbatim}

Перехідна частина (обумовлена стаціонарною частиною (умовлена початковою
умовою) правою частиною \texttt{f(x,\ t)}).

Розв'язок (9.10) показує, що температурна реакція складається з двох
частин: перша частина зумовлена початковими умовами, а друга ---
густиною джерела \texttt{f(x,t)}. Зверніть увагу, що термін «стаціонарна
частина» не є найкращим описом другої частини температурної реакції,
оскільки він не означає, що ця частина не залежить від часу
(стаціонарний стан може бути періодичним, якщо \texttt{f(x,\ t)} є
періодичною функцією часу t).

Проблема повністю вирішена. Перш ніж перейти до наступної теми, давайте
покажемо, як використовувати описаний вище метод для розв'язання

Конкретне завдання.

\section{Розв'язання конкретної задачі методом декомпозиції за власними
функціями}\label{ux440ux43eux437ux432ux44fux437ux430ux43dux43dux44f-ux43aux43eux43dux43aux440ux435ux442ux43dux43eux457-ux437ux430ux434ux430ux447ux456-ux43cux435ux442ux43eux434ux43eux43c-ux434ux435ux43aux43eux43cux43fux43eux437ux438ux446ux456ux457-ux437ux430-ux432ux43bux430ux441ux43dux438ux43cux438-ux444ux443ux43dux43aux446ux456ux44fux43cux438}

Розглянемо просту задачу

(9.11)

\begin{verbatim}
\begin{aligned}
(\text{УЧП})\quad & u_t = \alpha^2 u_{xx} + \sin(3\pi x), && 0 < x < 1,\quad 0 < t < \infty, \\
(\text{ГУ})\quad & u(0,t)=0,\quad u(1,t)=0, && 0 < t < \infty, \\
(\text{НУ})\quad & u(x,0)=\sin(\pi x), && 0 \le x \le 1.
\end{aligned}
\end{verbatim}

Наша мета --- обчислити коефіцієнти \texttt{T\_n(t)} при розкладі

\begin{verbatim}
u\left(x,\,t\right)=\sum_{n=1}^{\infty}T_n(t)\sin(n\pi x)
\end{verbatim}

(Власні функції цієї задачі \texttt{X\_n(x)} збігатися з власними
функціями попереднього). Якщо підставити цю декомпозицію на початкову
задачу, отримаємо наступну задачу Коші для визначення функцій
\texttt{T\_n(t)} :

\begin{verbatim}
\begin{aligned}
T_n' + (n\pi\alpha)^2 T_n &= f_n(t), \\
f_n(t) &= 2 \int_0^1 \sin(3\pi x)\sin(n\pi x)\,dx
= \begin{cases}
1, & n = 3, \\
0, & n \ne 3,
\end{cases} \\
T_n(0) &= 2 \int_0^1 \sin(\pi \xi)\sin(n\pi \xi)\,d\xi
= \begin{cases}
1, & n = 1, \\
0, & n \ne 1.
\end{cases}
\end{aligned}
\end{verbatim}

Перепишемо ці рівняння послідовно для
\texttt{n=1,2,\textbackslash{}dots}.

\begin{verbatim}
\begin{aligned}
(n=1)\quad & T_1' + (\pi\alpha)^2 T_1 = 0, && T_1(0) = 1, && T_1(t) = e^{-(\pi\alpha)^2 t}, \\
(n=2)\quad & T_2' + (2\pi\alpha)^2 T_2 = 0, && T_2(0) = 0, && T_2(t) = 0, \\
(n=3)\quad & T_3' + (3\pi\alpha)^2 T_3 = 4, && T_3(0) = 0, && T_3(t) = \frac{1}{(3\pi\alpha)^2}\left[1 - e^{-(3\pi\alpha)^2 t}\right], \\
(n\geqslant 4)\quad & T_n' + (n\pi\alpha)^2 T_n = 0, && T_n(0) = 0, && T_n(t) = 0.
\end{aligned}
\end{verbatim}

Отже, розв'язок задачі можна записати у вигляді

\begin{verbatim}
u(x, t) = e^{-(\pi\alpha)^2 t} \sin(\pi x) + \frac{1}{(3\pi\alpha)^2} [1 - e^{-(3\pi\alpha)^2 t}] \sin(3\pi x).
\end{verbatim}

(9.12)

\section{ЗАУВАЖЕННЯ}\label{ux437ux430ux443ux432ux430ux436ux435ux43dux43dux44f-2}

\begin{enumerate}
\def\labelenumi{\arabic{enumi}.}
\item
  Метод розкладу власних функцій є одним із найпотужніших методів
  розв'язання неоднорідних диференціальних рівнянь у похідних похідних.
  Пізніше, коли ми ознайомимося з інтегральними перетвореннями,
  побачимо, що існують й інші методи розв'язання задач такого типу.
\item
  Власні функції \texttt{X\_n(x)} , за допомогою якої виконується
  розклад, залежать від задачі, яку потрібно розв'язати, точніше від
  рівнянь і граничних умов. Читача просять розв'язати задачу 4
  наприкінці лекції, щоб перевірити, чи дійсно він навчився знаходити
  власні функції \texttt{X\_n(x)} .
\item
  Якщо читач пам'ятає основні факти з теорії звичайних диференціальних
  рівнянь, то він знає, що розв'язки, що відповідають неоднорідним
  членам вигляду
\end{enumerate}

\begin{verbatim}
P_n(x) e^{\alpha x} \cdot
\begin{cases} \sin{(\beta x)}, \\ \cos{(\beta x)}, \end{cases}
\end{verbatim}

можна знайти методом невизначених коефіцієнтів. Те саме і тут. Задача
(9.11) може бути розв'язана таким самим способом.

\section{ЗАДАЧІ}\label{ux437ux430ux434ux430ux447ux456}

\section{1. Рішення
задачі}\label{ux440ux456ux448ux435ux43dux43dux44f-ux437ux430ux434ux430ux447ux456}

\begin{verbatim}
\begin{aligned}
(\text{УЧП}) \quad & u_t = u_{xx} + \sin(3\pi x), && 0 < x < 1, \\
(\Gamma \text{У}) \quad & u(0,t) = 0,\quad u(1,t) = 0, && 0 < t < \infty, \\
(\text{НУ}) \quad & u(x,0) = \sin(\pi x), && 0 \leqslant x \leqslant 1.
\end{aligned}
\end{verbatim}

Отже, розв'язок задається формулою (9.12). Чи відповідає це рішення
вашим інтуїтивним уявленням? Що саме означає отриманий розв'язок? 2.
Вирішіть задачу

\begin{verbatim}
\begin{aligned}
(\text{РЧП})\quad & u_t = u_{xx} + \sin(\pi x) + \sin(2\pi x), && 0 < x < 1,\ 0 < t < \infty,\\
(\text{ГУ})\quad & u(0,t)=0,\quad u(1,t)=0, && 0 < t < \infty,\\
(\text{ПУ})\quad & u(x,0)=0, && 0 \le x \le 1.
\end{aligned}
\end{verbatim}

\section{3. Розв'яжіть
задачу}\label{ux440ux43eux437ux432ux44fux436ux456ux442ux44c-ux437ux430ux434ux430ux447ux443}

(PЧП)

\begin{verbatim}
u_t = u_{xx} + \sin(\pi x), \quad 0 < x < 1, \quad 0 < t < \infty,
\end{verbatim}

(гУ)

\begin{verbatim}
\begin{cases} u(0, t) = 0, \\ u(1, t) = 0, \\ u(x, 0) = 1, \quad 0 < x < 1, \end{cases}
\end{verbatim}

методом розкладу за власними функціями. 4. Найдіть рішення задачі

\begin{verbatim}
\begin{aligned}
(\text{УЧП})\quad & u_t = u_{xx} + \sin(\lambda_1 x), && 0 < x < 1,\ 0 < t < \infty, \\
(\text{ГУ})\quad & u(0,t) = 0,\quad u_x(1,t) + u(1,t) = 0, && 0 < t < \infty, \\
(\text{НУ})\quad & u(x,0) = 0, && 0 \leqslant x \leqslant 1.
\end{aligned}
\end{verbatim}

\texttt{Методом\ розкладу\ за\ власними\ функціями,\ якщо\ \textquotesingle{}\textbackslash{}lambda\_1\textquotesingle{}\ —\ перший\ корінь\ рівняння\ \textquotesingle{}\textbackslash{}operatorname\{tg\}\textbackslash{}lambda\ =\ -\textbackslash{}lambda\textquotesingle{}.\ Які\ власні\ функції\ цієї\ задачі?}

\section{5. Розв'яжіть
задачу}\label{ux440ux43eux437ux432ux44fux436ux456ux442ux44c-ux437ux430ux434ux430ux447ux443-1}

\begin{verbatim}
\begin{aligned}
(\text{РЧП})\quad & u_t = u_{xx}, && 0 < x < 1,\ 0 < t < \infty, \\
(\text{ГУ})\quad & u(0,t) = 0,\quad u(1,t) = \cos t, && 0 < t < \infty, \\
(\text{ПУ})\quad & u(x,0) = \varphi(x), && 0 \leqslant x \leqslant 1.
\end{aligned}
\end{verbatim}

Спочатку перетворивши граничні умови в нуль. Отримана задача
розв'язується методом розкладу за власними функціями.

\section{Лекція 10. ІНТЕГРАЛЬНІ ПЕРЕТВОРЕННЯ (СИНУСНІ ТА КОСИНУСНІ
ПЕРЕТВОРЕННЯ)}\label{ux43bux435ux43aux446ux456ux44f-10.-ux456ux43dux442ux435ux433ux440ux430ux43bux44cux43dux456-ux43fux435ux440ux435ux442ux432ux43eux440ux435ux43dux43dux44f-ux441ux438ux43dux443ux441ux43dux456-ux442ux430-ux43aux43eux441ux438ux43dux443ux441ux43dux456-ux43fux435ux440ux435ux442ux432ux43eux440ux435ux43dux43dux44f}

\textbf{МЕТА ЛЕКЦІЇ:} Ввести поняття інтегрального перетворення та
показати, як його застосування дозволяє перейти від ЧПР з \texttt{n}
змінними до нового рівняння з \texttt{n\ -\ 1} змінною.

Лекція демонструє, що інтегральні перетворення можна інтерпретувати як
розклад вхідного впливу на елементарні (частотний розклад), визначення
реакції системи на елементарні впливи та подальше підсумовування всіх
отриманих результатів.

В інтегральних перетвореннях операція диференціювання переходить в
операцію множення, і, відповідно, деякі частинні похідні переходять в
алгебраїчні вирази.

Для розв'язання задач дифузії на необмежених інтервалах вводять синусні
та косинусні перетворення. Розв'язок, отриманий у цьому випадку, цікавий
тим, що містить функцію, яка є доповненням до ймовірнісного інтеграла.

Інтегральним перетворенням зазвичай називають перетворення, яке кожній
функції \texttt{f(t)} ставить у відповідність нову функцію \texttt{F(s)}
згідно з формулою

\begin{verbatim}
F(s) = \int_{A}^{B} K(s, t) f(t) dt.
\end{verbatim}

Зверніть увагу, що початкова функція залежить від \texttt{t}, а отримана
функція --- від \texttt{s}. Функція \texttt{K(s,t)} називається ядром
перетворення. Це головний елемент, який відрізняє одне перетворення від
іншого. Зазвичай ядро обирають так, щоб перетворення мало певні задані
властивості. Межі інтегрування також залежать від типу перетворення.

Основною ідеєю, що лежить в основі інтегральних перетворень, є
можливість позбутися часткових похідних за однією з незалежних змінних.
Це означає, що перетворене рівняння матиме на одну змінну менше, ніж
початкове. Наприклад, якщо застосувати інтегральне перетворення до
рівняння, то

щоб позбутися похідної за часом, отримуємо звичайне диференціальне
рівняння з незалежною змінною \texttt{x}.

\begin{verbatim}
u_{xx} + u_{cy} + u_{zz} = 0
\end{verbatim}

застосувати перетворення Фур'є до змінної x, тоді похідна
\texttt{u\_\{xx\}} буде усунена, і в новому рівнянні залишаться лише дві
незалежні змінні y та z. Тоді отримуємо звичайне диференціальне рівняння
з однією незалежною змінною. Іншими словами, інтегральні перетворення
спрощують початкове завдання. Потім розв'язується перетворення задачі,
і, застосовуючи обернене перетворення до її розв'язку, отримується
розв'язок початкового перетворення. Схема описаного алгоритму показана
на рис. 10.1

\pandocbounded{\includegraphics[keepaspectratio]{images/_page_72_Figure_5.jpeg}}

Рис. 10.1. Загальна схема застосування інтегральних перетворень.

Як видно з рис. 10.1, кожному прямому інтегральному перетворенню
відповідає обернене перетворення, яке відновлює початкову функцію з
перетвореної.

Пряме та обернене перетворення зазвичай називають парою перетворень. У
таблиці 10.1 наведено кілька пар перетворень, які ми будемо
використовувати під час розв'язування рівнянь.

Слід зазначити, що в цих перетвореннях використовуються два типи
позначень. Наприклад, у випадку перетворення Лапласа запис
\texttt{\textbackslash{}mathcal\{L\}{[}f{]}} означає, що функцію
\texttt{f} перетворено, тоді як позначення \texttt{F(s)} вказує, що
результатом є функція, яка залежить від \texttt{s}. У цій лекції ми не
прагнемо розглянути всі пари перетворень. Ми зосередимося лише на
синусному та косинусному перетвореннях Фур'є (пари 1 і 2). Пізніше ми
розглянемо також інші пари. Ознайомлюючись із перетвореннями, ми
з'ясуємо їхні переваги та недоліки.

\subsection{Таблиця 10.1. Деякі пари інтегральних
перетворень}\label{ux442ux430ux431ux43bux438ux446ux44f-10.1.-ux434ux435ux44fux43aux456-ux43fux430ux440ux438-ux456ux43dux442ux435ux433ux440ux430ux43bux44cux43dux438ux445-ux43fux435ux440ux435ux442ux432ux43eux440ux435ux43dux44c}

\textbf{1. Синусне перетворення Фур'є:}

\begin{verbatim}
\mathcal{F}_s[f](\omega) = \frac{2}{\pi}\int_0^\infty f(t)\sin(\omega t)\,dt
\mathcal{F}_s^{-1}[F](t) = \int_0^\infty F(\omega)\sin(\omega t)\,d\omega
\end{verbatim}

\textbf{2. Косинусне перетворення Фур'є:}

\begin{verbatim}
\mathcal{F}_c[f](\omega) = \frac{2}{\pi}\int_0^\infty f(t)\cos(\omega t)\,dt
\mathcal{F}_c^{-1}[F](t) = \int_0^\infty F(\omega)\cos(\omega t)\,d\omega
\end{verbatim}

\textbf{3. Перетворення Фур'є:}

\begin{verbatim}
\mathcal{F}[f](\omega) = \frac{1}{\sqrt{2\pi}}\int_{-\infty}^{\infty} f(x)e^{-i\omega x}\,dx
\mathcal{F}^{-1}[F](x) = \frac{1}{\sqrt{2\pi}}\int_{-\infty}^{\infty} F(\omega)e^{i\omega x}\,d\omega
\end{verbatim}

\textbf{4. Скінченне синусне перетворення:}

\begin{verbatim}
S_n = \frac{2}{L}\int_0^L f(x)\sin\left(\frac{n\pi x}{L}\right)\,dx
f(x) = \sum_{n=1}^{\infty} S_n \sin\left(\frac{n\pi x}{L}\right)
\end{verbatim}

\textbf{5. Скінченне косинусне перетворення:}

\begin{verbatim}
C_n = \frac{2}{L}\int_0^L f(x)\cos\left(\frac{n\pi x}{L}\right)\,dx
f(x) = \frac{C_0}{2} + \sum_{n=1}^{\infty} C_n \cos\left(\frac{n\pi x}{L}\right)
\end{verbatim}

\textbf{6. Перетворення Лапласа:}

\begin{verbatim}
\mathcal{L}[f](s) = \int_0^\infty e^{-st}f(t)\,dt
\mathcal{L}^{-1}[F](t) = \frac{1}{2\pi i}\int_{\gamma-i\infty}^{\gamma+i\infty} e^{st}F(s)\,ds
\end{verbatim}

\textbf{7. Перетворення Ханкеля:}

\begin{verbatim}
H\{f\}(\xi) = F_n(\xi) = \int_0^\infty r J_n(\xi r) f(r)\,dr
H^{-1}[F_n](r) = \int_0^\infty \xi J_n(\xi r) F_n(\xi)\,d\xi
\end{verbatim}

Однак перед тим, як перейти до вивчення інтегральних перетворень, буде
корисно зрозуміти, що зазвичай називають спектром функції (або
спектральним розкладом функції).

\section{Спектральна функція Інтегральні перетворення та спектр функції
тісно пов'язані. Насправді інтегральне перетворення можна інтерпретувати
як розклад функції на певний спектр компонент. Цей спектр залежить від
обраного перетворення, але кожна функція має певний спектр відносно
цього перетворення Р и с
Р}\label{ux441ux43fux435ux43aux442ux440ux430ux43bux44cux43dux430-ux444ux443ux43dux43aux446ux456ux44f-ux456ux43dux442ux435ux433ux440ux430ux43bux44cux43dux456-ux43fux435ux440ux435ux442ux432ux43eux440ux435ux43dux43dux44f-ux442ux430-ux441ux43fux435ux43aux442ux440-ux444ux443ux43dux43aux446ux456ux457-ux442ux456ux441ux43dux43e-ux43fux43eux432ux44fux437ux430ux43dux456.-ux43dux430ux441ux43fux440ux430ux432ux434ux456-ux456ux43dux442ux435ux433ux440ux430ux43bux44cux43dux435-ux43fux435ux440ux435ux442ux432ux43eux440ux435ux43dux43dux44f-ux43cux43eux436ux43dux430-ux456ux43dux442ux435ux440ux43fux440ux435ux442ux443ux432ux430ux442ux438-ux44fux43a-ux440ux43eux437ux43aux43bux430ux434-ux444ux443ux43dux43aux446ux456ux457-ux43dux430-ux43fux435ux432ux43dux438ux439-ux441ux43fux435ux43aux442ux440-ux43aux43eux43cux43fux43eux43dux435ux43dux442.-ux446ux435ux439-ux441ux43fux435ux43aux442ux440-ux437ux430ux43bux435ux436ux438ux442ux44c-ux432ux456ux434-ux43eux431ux440ux430ux43dux43eux433ux43e-ux43fux435ux440ux435ux442ux432ux43eux440ux435ux43dux43dux44f-ux430ux43bux435-ux43aux43eux436ux43dux430-ux444ux443ux43dux43aux446ux456ux44f-ux43cux430ux454-ux43fux435ux432ux43dux438ux439-ux441ux43fux435ux43aux442ux440-ux432ux456ux434ux43dux43eux441ux43dux43e-ux446ux44cux43eux433ux43e-ux43fux435ux440ux435ux442ux432ux43eux440ux435ux43dux43dux44f-ux440-ux438-ux441-ux440}

\pandocbounded{\includegraphics[keepaspectratio]{images/_page_74_Picture_7.jpeg}}

ис 10.2 Розклад періодичної функції на ряд синусами та косинусами;
Послідовність квадратних імпульсів апроксимується синусами та косинусами
Наприклад, розглянемо розклад періодичної функції в ряд за синусами та
косинусами (ряд Фур'є):

\begin{verbatim}
f(x) = \sum_{n=0}^{\infty} \left[ A_n \cos(nx) + B_n \sin(nx) \right]
\end{verbatim}

(див. рис. 10.2).

Тут коефіцієнти \texttt{A\_n} і \texttt{B\_n} визначають внесок
\texttt{\textbackslash{}cos(nx)} та \texttt{\textbackslash{}sin(nx)} у
функцію \texttt{f(x)}, тоді як величина

\begin{verbatim}
\sqrt{A_n^2 + B_n^2}
\end{verbatim}

(яку називають спектром функції) показує амплітуду складової функції
\texttt{f(x)}, що відповідає частоті \texttt{n}.

1) Ряди Фур'є детальніше розглядаються у лекції 11.

Наприклад, якщо функція \texttt{f(x)} має вигляд

\begin{verbatim}
f(x) = 1 + \sin x + \frac{1}{5}\sin(3x) + \cos x + \frac{1}{2}\cos(2x) + \frac{1}{4}\cos(4x)
\end{verbatim}

тоді її спектр (дискретний) виглядатиме так, як показано на рис. 10.3.
Зчитуючи значення значення
\texttt{\textbackslash{}sqrt\{A\_n\^{}2\ +\ B\_n\^{}2\}} з графіка,
можна визначити амплітудне значення компоненти функції \texttt{f(x)},
яка має частоту n.

\pandocbounded{\includegraphics[keepaspectratio]{images/_page_75_Figure_8.jpeg}}

Рис. 10 3. Дискретний спектр функції \texttt{f(r)}. Періодичні функції
розкладаються в нескінченні ряди (вони мають дискретний спектр), тоді як
неперіодичні функції мають неперервний спектр. Якщо функція визначена
лише на скінченному інтервалі, її можна періодично продовжити на всю
вісь і розкласти в ряд Фур'є, який представлятиме цю функцію в межах
області визначення.

Хоча неперіодичну функцію не можна подати у вигляді нескінченного ряду
синусів і косинусів, можна записати неперервний аналог ряду Фур'є:
\texttt{f(x)\ =\ \textbackslash{}int\_\{-\textbackslash{}infty\}\^{}\{+\textbackslash{}infty\}\ \textbackslash{}left{[}C(\textbackslash{}omega)\textbackslash{}cos(\textbackslash{}omega\ x)\ +\ S(\textbackslash{}omega)\textbackslash{}sin(\textbackslash{}omega\ x)\textbackslash{}right{]}\textbackslash{},d\textbackslash{}omega}
де функції \texttt{S(\textbackslash{}omega)} і
\texttt{C(\textbackslash{}omega)} слугують мірою спектральної густини
синуса і косинуса у функції \texttt{f(x)} а величина

\begin{verbatim}
\sqrt{S^2(\omega) + C^2(\omega)}
\end{verbatim}

слугує мірою загального внеску частоти \texttt{\textbackslash{}omega} у
функцію \texttt{f(x)} і називається неперервним спектром функції
\texttt{f(x)}. Виходячи з наведеного вище інтуїтивного визначення
спектра функції, давайте проаналізуємо інтегральні перетворення «гвинт
за гвинтом». Спершу запишемо кілька властивостей цих перетворень, які ми
будемо широко використовувати.

\subsection{Синусне та косинусне перетворення
похідних}\label{ux441ux438ux43dux443ux441ux43dux435-ux442ux430-ux43aux43eux441ux438ux43dux443ux441ux43dux435-ux43fux435ux440ux435ux442ux432ux43eux440ux435ux43dux43dux44f-ux43fux43eux445ux456ux434ux43dux438ux445}

Для подальших застосувань нам знадобляться такі властивості:

\begin{verbatim}
\mathcal{F}_s[f'] = -\omega \mathcal{F}_c[f],
\qquad
\mathcal{F}_s[f''] = \frac{2}{\pi}\omega f(0) - \omega^2 \mathcal{F}_s[f]
\mathcal{F}_c[f'] = -\frac{2}{\pi}f(0) + \omega \mathcal{F}_s[f],
\qquad
\mathcal{F}_c[f''] = -\frac{2}{\pi}f'(0) - \omega^2 \mathcal{F}_c[f]
\end{verbatim}

Деякі інші властивості прямих і обернених синусних та косинусних
перетворень наведено в таблиці наприкінці книги. Тепер покажемо, як
синусне перетворення можна використати для розв'язання важливої змішаної
задачі.

\section{Розв'язання задачі дифузії на
напівпрямій}\label{ux440ux43eux437ux432ux44fux437ux430ux43dux43dux44f-ux437ux430ux434ux430ux447ux456-ux434ux438ux444ux443ux437ux456ux457-ux43dux430-ux43dux430ux43fux456ux432ux43fux440ux44fux43cux456ux439}

Задача, яку ми хочемо розв'язати, формулюється так:

\textbf{(Задача Коші):}

\begin{verbatim}
\begin{aligned}
(\text{РЧП})\quad & u_t = \alpha^2 u_{xx}, && 0 < x < \infty,\quad 0 < t < \infty, \\
(\text{ГУ})\quad & u(0,t) = A, && 0 < t < \infty, \\
(\text{ПУ})\quad & u(x,0) = 0, && 0 \le x < \infty.
\end{aligned}
\end{verbatim}

\pandocbounded{\includegraphics[keepaspectratio]{images/_page_76_Figure_14.jpeg}}

\emph{Рис. 10.4. Задача дифузії в напівнескінченній області:
\texttt{u(0,t)=A}, \texttt{u(x,0)=0}.}

\subsection{КРОК 1}\label{ux43aux440ux43eux43a-1}

Застосуємо синусне перетворення Фур'є за змінною \texttt{x} до обох
частин рівняння. Тоді отримаємо:

\begin{verbatim}
\mathcal{F}_s[u_t] = \alpha^2 \mathcal{F}_s[u_{xx}]
\end{verbatim}

Оскільки перетворення виконується за змінною \texttt{x}, змінна
\texttt{t} розглядається як параметр. Тому:

\begin{verbatim}
\mathcal{F}_s[u_t]  
= \frac{2}{\pi}\int_{0}^{\infty} u_t(x,t)\sin(\omega x)\,dx  
= \frac{d}{dt}\mathcal{F}_s[u]  
= \frac{dU}{dt},
\end{verbatim}

де
\texttt{U\ =\ U(\textbackslash{}omega,\ t)\ =\ \textbackslash{}mathcal\{F\}\_s{[}u{]}}.

Для другого доданка, використовуючи властивості синусного перетворення
та граничну умову \texttt{u(0,t)=A}, маємо:

\texttt{\textbackslash{}mathcal\{F\}\_s{[}u\_\{xx\}{]}\ =\ \textbackslash{}frac\{2\}\{\textbackslash{}pi\}\textbackslash{}omega\ u(0,t)\ -\ \textbackslash{}omega\^{}2\ \textbackslash{}mathcal\{F\}\_s{[}u{]}}
\texttt{=\ \textbackslash{}frac\{2A\textbackslash{}omega\}\{\textbackslash{}pi\}\ -\ \textbackslash{}omega\^{}2\ U(\textbackslash{}omega,t)}
Підставляючи ці вирази у рівняння
\texttt{u\_t\ =\ \textbackslash{}alpha\^{}2\ u\_\{xx\}}, отримуємо
звичайне диференціальне рівняння:
\texttt{\textbackslash{}frac\{dU\}\{dt\}\ =\ \textbackslash{}alpha\^{}2\textbackslash{}left{[}-\textbackslash{}omega\^{}2\ U\ +\ \textbackslash{}frac\{2A\textbackslash{}omega\}\{\textbackslash{}pi\}\textbackslash{}right{]}.}

Початкова умова для \texttt{U} випливає з умови \texttt{u(x,0)=0}:
\texttt{\textbackslash{}mathcal\{F\}\_s{[}u(x,0){]}\ =\ U(\textbackslash{}omega,0)\ =\ 0.}

Отже, замість початкової змішаної задачі ми отримали задачу Коші для
звичайного диференціального рівняння:

\textbf{(10.2)}

\begin{verbatim}
\begin{aligned}
\frac{dU}{dt} + \omega^2 \alpha^2 U &= \frac{2A\omega \alpha^2}{\pi}, \\
U(\omega,0) &= 0.
\end{aligned}
\end{verbatim}

КРОК 2. Щоб розв'язати отриману задачу Коші, можна використати будь-який
елементарний метод розв'язання звичайних диференціальних рівнянь. У
результаті маємо

\begin{verbatim}
U(\omega,t)=\frac{2A}{\pi\omega}\left(1-e^{-\omega^2\alpha^2 t}\right).
\end{verbatim}

Отже, перед нами результат застосування синусного перетворення до
шуканої функції \texttt{u(x,t)}. Останній крок полягає в застосуванні
оберненого перетворення:

\begin{verbatim}
u(x,t)=\mathcal{F}_s^{-1}[U].
\end{verbatim}

КРОК 3. Скориставшись таблицями перетворень, отримуємо

\begin{verbatim}
u(x,t)=A\operatorname{erfc}\!\left(\frac{x}{2\alpha\sqrt{t}}\right),
\end{verbatim}

де \texttt{\textbackslash{}operatorname\{erfc\}(x)} --- доповнююча
функція помилки:

\begin{verbatim}
\operatorname{erfc}(x)=\frac{2}{\sqrt{\pi}}\int_x^\infty e^{-s^2}\,ds.
\end{verbatim}

Графік цієї функції показано на рис. 10.5. Вона не виражається через
елементарні функції, але широко табульована в довідниках.

\pandocbounded{\includegraphics[keepaspectratio]{images/_page_79_Figure_4.jpeg}}

Рис. 10.5. Графіки функцій
\texttt{\textbackslash{}operatorname\{erf\}(x)} та
\texttt{\textbackslash{}operatorname\{erfc\}(x)}.

\section{Інтерпретація рішення
Розв'язок}\label{ux456ux43dux442ux435ux440ux43fux440ux435ux442ux430ux446ux456ux44f-ux440ux456ux448ux435ux43dux43dux44f-ux440ux43eux437ux432ux44fux437ux43eux43a}

\begin{verbatim}
u(x,t)=A\operatorname{erfc}\!\left(\frac{x}{2\alpha\sqrt{t}}\right)
\end{verbatim}

дає багато їжі для роздумів. На рисунку 10.6 показані графіки
ймовірнісної комплементарної функції на різних значеннях часу t.

\pandocbounded{\includegraphics[keepaspectratio]{images/_page_79_Figure_10.jpeg}}

Рис. 10.6. Розподіл температури в напівнескінченному стрижні, на кінці
якого підтримується постійна температура. Вони відрізняються один від
одного лише характеристичною шкалою вздовж осі x.

\section{ЗАВДАННЯ - 1. Довести ідентичності
(10.1).}\label{ux437ux430ux432ux434ux430ux43dux43dux44f---1.-ux434ux43eux432ux435ux441ux442ux438-ux456ux434ux435ux43dux442ux438ux447ux43dux43eux441ux442ux456-10.1.}

\begin{itemize}
\tightlist
\item
  \begin{enumerate}
  \def\labelenumi{\arabic{enumi}.}
  \setcounter{enumi}{1}
  \tightlist
  \item
    Розв'яжіть задачу Коші (10.2).
  \end{enumerate}
\item
  \begin{enumerate}
  \def\labelenumi{\arabic{enumi}.}
  \setcounter{enumi}{2}
  \tightlist
  \item
    Використайте синусоїдальне або косинусне перетворення для
    розв'язання задачі
  \end{enumerate}
\end{itemize}

\begin{verbatim}
\begin{aligned}
(\text{УЧП})\quad & u_t = \alpha^2 u_{xx}, && 0 < x < \infty,\ 0 < t < \infty, \\
(\text{ГУ})\quad & u_x(0,t) = 0, && 0 < t < \infty, \\
(\text{НУ})\quad & u(x,0) = H(1-x), && 0 \leqslant x < \infty.
\end{aligned}
\end{verbatim}

де \texttt{H(x)} --- функція Хевісайду:

\begin{verbatim}
H(x)=
\begin{cases}
0, & x<0,\\
1, & x\ge 0.
\end{cases}
\end{verbatim}

Вхідна до початкової умови функція зображена на рисунку.

\pandocbounded{\includegraphics[keepaspectratio]{images/_page_80_Picture_7.jpeg}}

Як виглядатиме графік розв'язку в різні моменти часу? 5. Знайдіть
прообраз константи при синусному перетворенні, тобто знайдіть
\texttt{f(x)}, якщо
\texttt{F(\textbackslash{}omega)=c=\textbackslash{}mathrm\{const\}}.

\section{\texorpdfstring{Лекція 11. РЯДИ І ПЕРЕТВОРЕННЯ ФУР'Є МЕТА
ЛЕКЦІЇ: Ввести поняття ряду Фур'є і показати, як періодичну функцію
\texttt{f(x)} можна подати рядом за синусами і косинусами. Якщо ж
функція \texttt{f(x)} є неперіодичною і визначена на всій дійсній осі,
то ряд Фур'є природно переходить у перетворення
Фур'є.}{Лекція 11. РЯДИ І ПЕРЕТВОРЕННЯ ФУР'Є МЕТА ЛЕКЦІЇ: Ввести поняття ряду Фур'є і показати, як періодичну функцію f(x) можна подати рядом за синусами і косинусами. Якщо ж функція f(x) є неперіодичною і визначена на всій дійсній осі, то ряд Фур'є природно переходить у перетворення Фур'є.}}\label{ux43bux435ux43aux446ux456ux44f-11.-ux440ux44fux434ux438-ux456-ux43fux435ux440ux435ux442ux432ux43eux440ux435ux43dux43dux44f-ux444ux443ux440ux454-ux43cux435ux442ux430-ux43bux435ux43aux446ux456ux457-ux432ux432ux435ux441ux442ux438-ux43fux43eux43dux44fux442ux442ux44f-ux440ux44fux434ux443-ux444ux443ux440ux454-ux456-ux43fux43eux43aux430ux437ux430ux442ux438-ux44fux43a-ux43fux435ux440ux456ux43eux434ux438ux447ux43dux443-ux444ux443ux43dux43aux446ux456ux44e-fx-ux43cux43eux436ux43dux430-ux43fux43eux434ux430ux442ux438-ux440ux44fux434ux43eux43c-ux437ux430-ux441ux438ux43dux443ux441ux430ux43cux438-ux456-ux43aux43eux441ux438ux43dux443ux441ux430ux43cux438.-ux44fux43aux449ux43e-ux436-ux444ux443ux43dux43aux446ux456ux44f-fx-ux454-ux43dux435ux43fux435ux440ux456ux43eux434ux438ux447ux43dux43eux44e-ux456-ux432ux438ux437ux43dux430ux447ux435ux43dux430-ux43dux430-ux432ux441ux456ux439-ux434ux456ux439ux441ux43dux456ux439-ux43eux441ux456-ux442ux43e-ux440ux44fux434-ux444ux443ux440ux454-ux43fux440ux438ux440ux43eux434ux43dux43e-ux43fux435ux440ux435ux445ux43eux434ux438ux442ux44c-ux443-ux43fux435ux440ux435ux442ux432ux43eux440ux435ux43dux43dux44f-ux444ux443ux440ux454.}

Важлива роль рядів Фур'є в теорії рівнянь з частинними похідними
пов'язана з тим, що періодичні функції, а також функції, визначені на
скінченному відрізку, можна розкладати в ряди за синусами і косинусами.
Це часто дозволяє звести початкову задачу до задач для окремих гармонік.
Наприклад, пилкоподібну функцію

~~~
`f(x)=x, \qquad -L < x < L, \qquad f(x+2L)=f(x),`
~~~

можна подати рядом Фур'є (11.1)

~~~
`f(x) = \frac{a_0}{2} + \sum_{n=1}^{\infty} \left[a_n \cos\!\left(\frac{n\pi x}{L}\right) + b_n \sin\!\left(\frac{n\pi x}{L}\right)\right].`
~~~

Коефіцієнти Фур'є \texttt{a\_n} і \texttt{b\_n} знаходяться за формулами
Ейлера. Для цієї функції маємо

~~~
`a_n = \frac{1}{L}\int_{-L}^{L} f(x)\cos\!\left(\frac{n\pi x}{L}\right)\,dx = 0, \qquad n=0,1,2,\dots,`
~~~

~~~
`b_n = \frac{1}{L}\int_{-L}^{L} f(x)\sin\!\left(\frac{n\pi x}{L}\right)\,dx`
`= -(-1)^n \frac{2L}{n\pi}, \qquad n=1,2,\dots.`
~~~

\pandocbounded{\includegraphics[keepaspectratio]{images/_page_82_Figure_7.jpeg}}

Рис. 11.1. Розклад пилкоподібної функції в ряд Фур'є. Тому розклад цієї
функції набирає вигляду (11.3)

~~~
`f(x) = \frac{2L}{\pi}\left[\sin\!\left(\frac{\pi x}{L}\right) - \frac{1}{2}\sin\!\left(\frac{2\pi x}{L}\right) + \frac{1}{3}\sin\!\left(\frac{3\pi x}{L}\right) - \dots \right].`
~~~

Кожен доданок цього ряду називається гармонікою, а всі гармоніки мають
частоти, кратні основній частоті. У теорії рядів Фур'є важливу роль
відіграє теорема Діріхле: якщо обмежена періодична функція на кожному
періоді має лише скінченне число максимумів, мінімумів і точок розриву,
то її ряд Фур'є збігається до самої функції в точках неперервності і до
півсуми односторонніх границь у точках розриву. Перед переходом до
перетворення Фур'є корисно ввести поняття дискретного частотного спектра
періодичної функції. Якщо

~~~
`c_n = \sqrt{a_n^2 + b_n^2}, \qquad n=0,1,2,\dots,`
~~~

то числа \texttt{c\_n} показують внесок відповідних частотних компонент
у функцію \texttt{f(x)}. Для пилкоподібної функції

~~~
`c_n = \frac{2L}{n\pi}, \qquad n=1,2,\dots, \qquad c_0 = 0.`
~~~

\pandocbounded{\includegraphics[keepaspectratio]{images/_page_84_Figure_8.jpeg}}

Рис. 11.2. Дискретний частотний спектр пилкоподібної функції: a ---
основна частота (період \texttt{2L}), b --- \texttt{n}-та гармоніка
(період \texttt{2L/n}). Послідовність
\texttt{\textbackslash{}\{c\_n\textbackslash{}\}} нагадує спектральний
розклад білого світла на кольори. Введемо тепер поняття перетворення
Фур'є. Для неперіодичних функцій, визначених на всій дійсній осі, роль
ряду Фур'є відіграє інтеграл Фур'є (11.4)

~~~
`f(x) = \int_{0}^{\infty} a(\xi)\cos(\xi x)\,d\xi + \int_{0}^{\infty} b(\xi)\sin(\xi x)\,d\xi,`
~~~

де (11.5)

~~~
`a(\xi) = \frac{1}{\pi}\int_{0}^{\infty} f(x)\cos(\xi x)\,dx, \qquad b(\xi) = \frac{1}{\pi}\int_{0}^{\infty} f(x)\sin(\xi x)\,dx.`
~~~

За аналогією з рядом Фур'є визначимо неперервний частотний спектр

~~~
`C(\xi) = \sqrt{a^2(\xi) + b^2(\xi)}.`
~~~

Частотний спектр є мірою внеску частоти \texttt{\textbackslash{}xi} у
функцію \texttt{f(x)}.

\pandocbounded{\includegraphics[keepaspectratio]{images/_page_85_Figure_7.jpeg}}

Рис. 11.3. Частотні спектри деяких функцій. Тепер перейдемо до
перетворення Фур'є.

\section{\texorpdfstring{Преобразование Фурье Як уже зазначалося раніше,
неперіодичні функції, задані на
\texttt{(-\textbackslash{}infty,+\textbackslash{}infty)}, не
розкладаються у звичайні ряди Фур'є. Проте для деяких таких функцій
можна побудувати аналог ряду Фур'є у вигляді інтеграла Фур'є, тобто
розкладу за неперервним спектром
частот:}{Преобразование Фурье Як уже зазначалося раніше, неперіодичні функції, задані на (-\textbackslash infty,+\textbackslash infty), не розкладаються у звичайні ряди Фур'є. Проте для деяких таких функцій можна побудувати аналог ряду Фур'є у вигляді інтеграла Фур'є, тобто розкладу за неперервним спектром частот:}}\label{ux43fux440ux435ux43eux431ux440ux430ux437ux43eux432ux430ux43dux438ux435-ux444ux443ux440ux44cux435-ux44fux43a-ux443ux436ux435-ux437ux430ux437ux43dux430ux447ux430ux43bux43eux441ux44f-ux440ux430ux43dux456ux448ux435-ux43dux435ux43fux435ux440ux456ux43eux434ux438ux447ux43dux456-ux444ux443ux43dux43aux446ux456ux457-ux437ux430ux434ux430ux43dux456-ux43dux430--inftyinfty-ux43dux435-ux440ux43eux437ux43aux43bux430ux434ux430ux44eux442ux44cux441ux44f-ux443-ux437ux432ux438ux447ux430ux439ux43dux456-ux440ux44fux434ux438-ux444ux443ux440ux454.-ux43fux440ux43eux442ux435-ux434ux43bux44f-ux434ux435ux44fux43aux438ux445-ux442ux430ux43aux438ux445-ux444ux443ux43dux43aux446ux456ux439-ux43cux43eux436ux43dux430-ux43fux43eux431ux443ux434ux443ux432ux430ux442ux438-ux430ux43dux430ux43bux43eux433-ux440ux44fux434ux443-ux444ux443ux440ux454-ux443-ux432ux438ux433ux43bux44fux434ux456-ux456ux43dux442ux435ux433ux440ux430ux43bux430-ux444ux443ux440ux454-ux442ux43eux431ux442ux43e-ux440ux43eux437ux43aux43bux430ux434ux443-ux437ux430-ux43dux435ux43fux435ux440ux435ux440ux432ux43dux438ux43c-ux441ux43fux435ux43aux442ux440ux43eux43c-ux447ux430ux441ux442ux43eux442}

(11.4)

~~~
`f(x) = \int_{0}^{\infty} a(\xi)\cos(\xi x)\,d\xi + \int_{0}^{\infty} b(\xi)\sin(\xi x)\,d\xi,`
~~~

де (11.5)

~~~
`a(\xi) = \frac{1}{\pi}\int_{0}^{\infty} f(x)\cos(\xi x)\,dx, \qquad b(\xi) = \frac{1}{\pi}\int_{0}^{\infty} f(x)\sin(\xi x)\,dx.`
~~~

З формули (11.4) видно, що інтеграл Фур'є розкладає функцію
\texttt{f(x)} у спектр за всіма частотами
\texttt{0\textless{}\textbackslash{}xi\textless{}\textbackslash{}infty},
а не лише за частотами, кратними основній, як це було для періодичних
функцій. За аналогією з рядами Фур'є визначимо частотний спектр формулою

~~~
`C(\xi) = \sqrt{a^2(\xi) + b^2(\xi)}.`
~~~

Частотний спектр є мірою внеску частоти \texttt{\textbackslash{}xi} у
функцію \texttt{f(x)}. На рис. 11.3 зображено кілька функцій та їхні
спектри. Функції з різкими зламами мають широкий частотний спектр, тоді
як для простої гармоніки
\texttt{f(x)=\textbackslash{}sin(\textbackslash{}xi\_0\ x)} спектр
зосереджується поблизу однієї частоти
\texttt{\textbackslash{}xi=\textbackslash{}xi\_0}. Тепер ми готові
перейти до експоненціальної форми перетворення Фур'є. Скориставшись
формулою Ейлера

~~~
`e^{i\theta} = \cos\theta + i\sin\theta,`
~~~

після нескладних перетворень приходимо до інтегрального представлення
Фур'є (11.6)

~~~
`f(x) = \frac{1}{\sqrt{2\pi}} \int_{-\infty}^{+\infty}`
`\left[\frac{1}{\sqrt{2\pi}} \int_{-\infty}^{+\infty} f(s)e^{-i\xi s}\,ds \right]`
`e^{i\xi x}\,d\xi.`
~~~

Звідси отримуємо пару інтегральних перетворень (11.7)

~~~
`\mathcal{F}[f] = F(\xi) = \frac{1}{\sqrt{2\pi}} \int_{-\infty}^{+\infty} f(x)e^{-i\xi x}\,dx`
~~~

(пряме перетворення Фур'є) та

~~~
`\mathcal{F}^{-1}[F] = f(x) = \frac{1}{\sqrt{2\pi}} \int_{-\infty}^{+\infty} F(\xi)e^{i\xi x}\,d\xi`
~~~

(обернене перетворення Фур'є). У наступній лекції ми вивчатимемо
властивості цих перетворень.

\section{ЗАУВАЖЕННЯ}\label{ux437ux430ux443ux432ux430ux436ux435ux43dux43dux44f-3}

\begin{enumerate}
\def\labelenumi{\arabic{enumi}.}
\tightlist
\item
  Образ Фур'є \texttt{F(\textbackslash{}xi)} функції \texttt{f(x)} може
  бути комплекснозначною функцією. Наприклад, для функції
\end{enumerate}

~~~
f(x)=
`\begin{cases}`
`0, & x \le 0,\\`
e^{-x}, & x > 0,
`\end{cases}`
~~~

образ Фур'є має вигляд

~~~
`F(\xi) = \frac{1}{\sqrt{2\pi}} \frac{1-i\xi}{1+\xi^2}.`
~~~

\begin{enumerate}
\def\labelenumi{\arabic{enumi}.}
\setcounter{enumi}{1}
\tightlist
\item
  Абсолютне значення образу Фур'є називається частотним спектром функції
  \texttt{f(x)}. Наприклад, якщо позначити
  \texttt{C(\textbackslash{}xi)=\textbar{}F(\textbackslash{}xi)\textbar{}},
  то маємо
\end{enumerate}

~~~
`C(\xi) = \sqrt{\frac{1}{2\pi(1+\xi^2)}}.`
~~~

\texttt{3.\ Не\ всі\ функції\ допускають\ перетворення\ Фур\textquotesingle{}є,\ тобто\ інтеграл\ (11.7)\ може\ не\ існувати.\ Наприклад,\ функції\ \textquotesingle{}f(x)=c\textquotesingle{},\ \textquotesingle{}\textbackslash{}sin\ x\textquotesingle{}\ і\ \textquotesingle{}e\^{}x\textquotesingle{}\ не\ мають\ класичного\ образу\ Фур\textquotesingle{}є.\ У\ прикладних\ задачах\ перетворення\ Фур\textquotesingle{}є\ зазвичай\ застосовують\ до\ температурних,\ хвильових\ та\ інших\ фізичних\ полів,\ які\ прямують\ до\ нуля\ при\ \textquotesingle{}\textbar{}x\textbar{}\textbackslash{}to\textbackslash{}infty\textquotesingle{}.}

\section{ЗАВДАННЯ}\label{ux437ux430ux432ux434ux430ux43dux43dux44f-3}

\begin{enumerate}
\def\labelenumi{\arabic{enumi}.}
\tightlist
\item
  Розкласти у ряд Фур'є за синусами функцію
\end{enumerate}

~~~
f(x)=
`\begin{cases}`
-1, & -1 < x < 0,\\
`1, & 0 \le x < 1,`
`\end{cases}`
`\qquad f(x+2)=f(x).`
~~~

Побудуйте суми перших 2, 3 і 4 членів ряду, щоб простежити його
збіжність до \texttt{f(x)}. Побудуйте також графіки частотного спектра
цієї функції. 2. Показати, що якщо почленно диференціювати ряд Фур'є
(11.3) для пилкоподібної функції, то отримаємо послідовність, яка не
збігається до похідної цієї функції. 3. Побудувати частотний спектр
таких періодичних функцій:
\texttt{-\ a)\ \textquotesingle{}f(x)=\textbackslash{}sin\ x\textquotesingle{},}
\texttt{-\ б)\ \textquotesingle{}f(x)=\textbackslash{}sin\ x+\textbackslash{}cos\ 2x\textquotesingle{},}
\texttt{-\ в)\ \textquotesingle{}f(x)=\textbackslash{}sin\ x+\textbackslash{}cos\ x+0.5\textbackslash{}sin\ 3x\textquotesingle{}.}
4. Знайти образ Фур'є \texttt{F(\textbackslash{}xi)} та частотний спектр
\texttt{C(\textbackslash{}xi)=\textbar{}F(\textbackslash{}xi)\textbar{}}
для функції

~~~
f(x)=
`\begin{cases}`
1, & -1 < x < 1,\\
`0, & \text{в інших точках.}`
`\end{cases}`
~~~

\texttt{5.\ Показати,\ що\ абсолютне\ значення\ функції\ \textquotesingle{}F(\textbackslash{}xi)=\textbackslash{}dfrac\{1\}\{1+i\textbackslash{}xi\}\textquotesingle{}\ дорівнює}

~~~
`\left|F(\xi)\right| = \sqrt{\frac{1}{1+\xi^2}}.`
~~~

\texttt{ПРИМІТКА.\ Спочатку\ помножте\ чисельник\ і\ знаменник\ на\ \textquotesingle{}1-i\textbackslash{}xi\textquotesingle{},\ а\ потім\ позбудьтеся\ уявної\ одиниці\ \textquotesingle{}i\textquotesingle{}\ у\ знаменнику.}
\texttt{6.\ Перевірити\ ортогональність\ функцій\ \textquotesingle{}\textbackslash{}\{\textbackslash{}sin(n\textbackslash{}pi\ x/L)\textbackslash{}\}\textquotesingle{}\ та\ \textquotesingle{}\textbackslash{}\{\textbackslash{}cos(n\textbackslash{}pi\ x/L)\textbackslash{}\}\textquotesingle{}\ на\ відрізку\ \textquotesingle{}{[}-L,L{]}\textquotesingle{}:}

~~~
`\int_{-L}^{L} \sin\!\left(\frac{m\pi x}{L}\right)\sin\!\left(\frac{n\pi x}{L}\right)\,dx =`
`\begin{cases}`
`0, & m \neq n,\\`
L, & m = n,
`\end{cases}`
~~~

~~~
`\int_{-L}^{L} \cos\!\left(\frac{m\pi x}{L}\right)\cos\!\left(\frac{n\pi x}{L}\right)\,dx =`
`\begin{cases}`
`0, & m \neq n,\\`
L, & m = n,
`\end{cases}`
~~~

~~~
`\int_{-L}^{L} \sin\!\left(\frac{m\pi x}{L}\right)\cos\!\left(\frac{n\pi x}{L}\right)\,dx = 0`
`\quad \text{для всіх } m,n=1,2,\dots`
~~~

\section{Лекція 12. ПЕРЕТВОРЕННЯ ФУР'Є ТА ЙОГО ЗАСТОСУВАННЯ ДО
РОЗВ'ЯЗАННЯ ДИФЕРЕНЦІАЛЬНИХ РІВНЯНЬ У
ПОХІДНИХ}\label{ux43bux435ux43aux446ux456ux44f-12.-ux43fux435ux440ux435ux442ux432ux43eux440ux435ux43dux43dux44f-ux444ux443ux440ux454-ux442ux430-ux439ux43eux433ux43e-ux437ux430ux441ux442ux43eux441ux443ux432ux430ux43dux43dux44f-ux434ux43e-ux440ux43eux437ux432ux44fux437ux430ux43dux43dux44f-ux434ux438ux444ux435ux440ux435ux43dux446ux456ux430ux43bux44cux43dux438ux445-ux440ux456ux432ux43dux44fux43dux44c-ux443-ux43fux43eux445ux456ux434ux43dux438ux445}

\subsection{МЕТА ЛЕКЦІЇ Проілюструвати основні властивості перетворення
Фур'є та показати, як ці властивості застосовуються до розв'язання
рівнянь з частинними похідними. Зокрема, ми побачимо, як під дією
перетворення Фур'є операція диференціювання замінюється множенням, а
диференціальні рівняння переходять в
алгебраїчні.}\label{ux43cux435ux442ux430-ux43bux435ux43aux446ux456ux457-ux43fux440ux43eux456ux43bux44eux441ux442ux440ux443ux432ux430ux442ux438-ux43eux441ux43dux43eux432ux43dux456-ux432ux43bux430ux441ux442ux438ux432ux43eux441ux442ux456-ux43fux435ux440ux435ux442ux432ux43eux440ux435ux43dux43dux44f-ux444ux443ux440ux454-ux442ux430-ux43fux43eux43aux430ux437ux430ux442ux438-ux44fux43a-ux446ux456-ux432ux43bux430ux441ux442ux438ux432ux43eux441ux442ux456-ux437ux430ux441ux442ux43eux441ux43eux432ux443ux44eux442ux44cux441ux44f-ux434ux43e-ux440ux43eux437ux432ux44fux437ux430ux43dux43dux44f-ux440ux456ux432ux43dux44fux43dux44c-ux437-ux447ux430ux441ux442ux438ux43dux43dux438ux43cux438-ux43fux43eux445ux456ux434ux43dux438ux43cux438.-ux437ux43eux43aux440ux435ux43cux430-ux43cux438-ux43fux43eux431ux430ux447ux438ux43cux43e-ux44fux43a-ux43fux456ux434-ux434ux456ux454ux44e-ux43fux435ux440ux435ux442ux432ux43eux440ux435ux43dux43dux44f-ux444ux443ux440ux454-ux43eux43fux435ux440ux430ux446ux456ux44f-ux434ux438ux444ux435ux440ux435ux43dux446ux456ux44eux432ux430ux43dux43dux44f-ux437ux430ux43cux456ux43dux44eux454ux442ux44cux441ux44f-ux43cux43dux43eux436ux435ux43dux43dux44fux43c-ux430-ux434ux438ux444ux435ux440ux435ux43dux446ux456ux430ux43bux44cux43dux456-ux440ux456ux432ux43dux44fux43dux43dux44f-ux43fux435ux440ux435ux445ux43eux434ux44fux442ux44c-ux432-ux430ux43bux433ux435ux431ux440ux430ux457ux447ux43dux456.}

\textbf{(12.1)}

\begin{verbatim}
\mathcal{F}[f] = F(\xi) = \frac{1}{\sqrt{2\pi}} \int_{-\infty}^{+\infty} f(x)e^{-i\xi x}\,dx.
\end{verbatim}

Отже, якщо функція \texttt{f(x)} визначена на всій дійсній осі, то
підставляючи її у формулу (12.1), ми отримуємо нову функцію
\texttt{F(\textbackslash{}xi)}, визначену для
\texttt{-\textbackslash{}infty\ \textless{}\ \textbackslash{}xi\ \textless{}\ \textbackslash{}infty}.

Таблиця 1. 12.1 показує деякі функції та їхні фур'є-образи. Більш
розгорнута таблиця наведена в додатку. З наведеної таблиці видно, що
образ Фур'є може бути комплекснозначною функцією аргументу
\texttt{\textbackslash{}xi}. У першому прикладі функція
\texttt{F(\textbackslash{}xi)} містить уявну одиницю \texttt{i}, тобто є
комплексною функцією. Аргумент \texttt{\textbackslash{}xi} --- дійсний,
але значення функції --- комплексні.

Перетворення Фур'є дуже корисне, оскільки воно замінює диференціювання
множенням і перетворює диференціальні рівняння на алгебраїчні.

\subsection{Таблиця 12.1. Приклади перетворень
Фур'є}\label{ux442ux430ux431ux43bux438ux446ux44f-12.1.-ux43fux440ux438ux43aux43bux430ux434ux438-ux43fux435ux440ux435ux442ux432ux43eux440ux435ux43dux44c-ux444ux443ux440ux454}

\begin{enumerate}
\def\labelenumi{\arabic{enumi}.}
\tightlist
\item
  Для непарної функції
\end{enumerate}

\begin{verbatim}
f(x)=
\begin{cases}
e^{-x}, & x\ge 0,\\
-e^{x}, & x<0,
\end{cases}
\qquad F(\xi)=-i\sqrt{\frac{2}{\pi}}\frac{\xi}{1+\xi^2}.
\end{verbatim}

\begin{enumerate}
\def\labelenumi{\arabic{enumi}.}
\setcounter{enumi}{1}
\tightlist
\item
  Для прямокутного імпульсу
\end{enumerate}

\begin{verbatim}
f(x)=
\begin{cases}
1, & |x|<1,\\
0, & |x|\ge 1,
\end{cases}
\qquad F(\xi)=\sqrt{\frac{2}{\pi}}\frac{\sin\xi}{\xi}.
\end{verbatim}

\begin{enumerate}
\def\labelenumi{\arabic{enumi}.}
\setcounter{enumi}{2}
\tightlist
\item
  Для гауссової функції
\end{enumerate}

\begin{verbatim}
f(x)=e^{-x^2},
\qquad F(\xi)=\frac{1}{\sqrt{2}}e^{-\xi^2/4}.
\end{verbatim}

\section{Основні властивості перетворення
Фур'є}\label{ux43eux441ux43dux43eux432ux43dux456-ux432ux43bux430ux441ux442ux438ux432ux43eux441ux442ux456-ux43fux435ux440ux435ux442ux432ux43eux440ux435ux43dux43dux44f-ux444ux443ux440ux454}

\subsection{Властивість 1 (Пара перетворень
Фур'є)}\label{ux432ux43bux430ux441ux442ux438ux432ux456ux441ux442ux44c-1-ux43fux430ux440ux430-ux43fux435ux440ux435ux442ux432ux43eux440ux435ux43dux44c-ux444ux443ux440ux454}

\begin{verbatim}
\mathcal{F}[f] = F(\xi) = \frac{1}{\sqrt{2\pi}} \int_{-\infty}^{+\infty} f(x)e^{-i\xi x}\,dx.
\end{verbatim}

Обернене перетворення:

\begin{verbatim}
\mathcal{F}^{-1}[F] = f(x) = \frac{1}{\sqrt{2\pi}} \int_{-\infty}^{+\infty} F(\xi)e^{i\xi x}\,d\xi.
\end{verbatim}

\begin{figure}
\centering
\pandocbounded{\includegraphics[keepaspectratio,alt={Рис. 12.1}]{images/_page_90_Figure_7.jpeg}}
\caption{Рис. 12.1}
\end{figure}

Рис. 12.1. Графіки функції та її образу Фур'є.

Приклад:

\begin{verbatim}
e^{-|x|} \xrightarrow{\mathcal{F}} \sqrt{\frac{2}{\pi}}\frac{1}{1+\xi^2} \xrightarrow{\mathcal{F}^{-1}} e^{-|x|}.
\end{verbatim}

\subsection{Властивість 2
(Лінійність)}\label{ux432ux43bux430ux441ux442ux438ux432ux456ux441ux442ux44c-2-ux43bux456ux43dux456ux439ux43dux456ux441ux442ux44c}

\begin{verbatim}
\mathcal{F}[af+bg] = a\mathcal{F}[f] + b\mathcal{F}[g].
\end{verbatim}

Наприклад:

\begin{verbatim}
\mathcal{F}\left[e^{-|x|} + 3e^{-x^2}\right]
= \mathcal{F}\left[e^{-|x|}\right] + 3\mathcal{F}\left[e^{-x^2}\right].
\end{verbatim}

\subsection{Властивість 3 (Перетворення
похідних)}\label{ux432ux43bux430ux441ux442ux438ux432ux456ux441ux442ux44c-3-ux43fux435ux440ux435ux442ux432ux43eux440ux435ux43dux43dux44f-ux43fux43eux445ux456ux434ux43dux438ux445}

Для похідних за \texttt{x}:

\begin{verbatim}
\mathcal{F}[u_x] = i\xi \mathcal{F}[u],
\mathcal{F}[u_{xx}] = -\xi^2 \mathcal{F}[u].
\end{verbatim}

Для похідних за \texttt{t}:

\begin{verbatim}
\mathcal{F}[u_t] = \frac{\partial}{\partial t}\mathcal{F}[u],
\mathcal{F}[u_{tt}] = \frac{\partial^2}{\partial t^2}\mathcal{F}[u].
\end{verbatim}

\subsection{Властивість 4
(Згортка)}\label{ux432ux43bux430ux441ux442ux438ux432ux456ux441ux442ux44c-4-ux437ux433ux43eux440ux442ux43aux430}

\begin{verbatim}
\mathcal{F}[f(x)g(x)] \neq \mathcal{F}[f]\mathcal{F}[g].
\end{verbatim}

\textbf{(12.2)}

\begin{verbatim}
\mathcal{F}[f*g] = \mathcal{F}[f]\mathcal{F}[g].
\end{verbatim}

\textbf{(12.3)}

\begin{verbatim}
(f*g)(x) = \frac{1}{\sqrt{2\pi}} \int_{-\infty}^{+\infty} f(x-\xi)g(\xi)\,d\xi.
\end{verbatim}

\subsection{Приклад
згортки}\label{ux43fux440ux438ux43aux43bux430ux434-ux437ux433ux43eux440ux442ux43aux438}

Нехай:

\begin{verbatim}
f(x) = x, \qquad g(x) = e^{-x^2}.
\end{verbatim}

Тоді:

\begin{verbatim}
(f*g)(x) = \frac{1}{\sqrt{2\pi}} \int_{-\infty}^{+\infty} (x-\xi)e^{-\xi^2}\,d\xi = \frac{x}{\sqrt{2}}.
\end{verbatim}

Використано:

\begin{verbatim}
\int_{-\infty}^{+\infty} e^{-\xi^2}\,d\xi = \sqrt{\pi}.
\end{verbatim}

\subsection{Інтерпретація}\label{ux456ux43dux442ux435ux440ux43fux440ux435ux442ux430ux446ux456ux44f}

\textbf{(12.4)}

\begin{verbatim}
\mathcal{F}^{-1}\left\{\mathcal{F}[f]\mathcal{F}[g]\right\}.
\end{verbatim}

Із властивості згортки:

\begin{verbatim}
f * g = \mathcal{F}^{-1}\left\{\mathcal{F}[f]\mathcal{F}[g]\right\}.
\end{verbatim}

Тепер ми готові розв'язувати рівняння з частинними похідними за
допомогою перетворення Фур'є. \# Розв'язання проблеми Коші Розглянемо
проблему поширення тепла в \emph{нескінченному} стрижні, якщо задано
початкову температуру \texttt{u(x,\ 0)\ =\ \textbackslash{}varphi(x)} .
Цю задачу зазвичай називають задачою Коші для рівняння теплопровідності
або задачою початкових умов

\begin{verbatim}
\begin{aligned}
(\text{РЧП})\quad & u_t=\alpha^2u_{xx}, && -\infty<x<\infty,\ 0<t<\infty,\\
(\text{НУ})\quad & u(x,0)=q(x), && -\infty<x<\infty.
\end{aligned}
\end{verbatim}

Розв'язання задачі можна поділити на три основні етапи. КРОК 1
(Конверсія завдань). Оскільки просторова змінна x змінюється між
\texttt{-\textbackslash{}infty} і \texttt{+\textbackslash{}infty} , ми
піддамо рівняння та початкову умову (12,6) перетворенню Фур'є на змінній
\texttt{x\^{}\{(1)\}} (x --- інтегрована змінна у перетворенні). Після
цього ми отримуємо

\begin{verbatim}
\mathcal{F}[u_t] = \alpha^2 \mathcal{F}[u_{xx}],\qquad
\mathcal{F}[u(x,0)] = \mathcal{F}[\varphi(x)].
\end{verbatim}

Використовуючи властивості перетворення Фур'є, отримуємо задачу Коші

\begin{verbatim}
\begin{cases}
\dfrac{dU}{dt}=-\alpha^2\xi^2U,\\
U(0)=\Phi(\xi).
\end{cases}
\end{verbatim}

де
\texttt{U=U(\textbackslash{}xi,t)=\textbackslash{}mathcal\{F\}{[}u(x,t){]}}
та
\texttt{\textbackslash{}Phi(\textbackslash{}xi)=\textbackslash{}mathcal\{F\}{[}\textbackslash{}varphi(x){]}}.
\#\# КРОК 2 (Розв'язання перетвореної задачі)

Оскільки нова змінна відіграє роль константи в задачі (12.7), її
розв'язок легко знаходиться:

\textbf{(12.8)}

\begin{verbatim}
U(\xi,t)=\Phi(\xi)e^{-\alpha^2\xi^2 t}.
\end{verbatim}

\subsection{КРОК 3 (Знаходження оберненого
перетворення)}\label{ux43aux440ux43eux43a-3-ux437ux43dux430ux445ux43eux434ux436ux435ux43dux43dux44f-ux43eux431ux435ux440ux43dux435ux43dux43eux433ux43e-ux43fux435ux440ux435ux442ux432ux43eux440ux435ux43dux43dux44f}

Шуканий розв'язок \texttt{u(x,t)} знаходиться за формулою:

\begin{verbatim}
u(x, t) = \mathcal{F}^{-1}\!\left[U(\xi, t)\right]
= \mathcal{F}^{-1}\!\left[\Phi(\xi)e^{-\alpha^2 \xi^2 t}\right].
\end{verbatim}

Припускається, що існує перетворення Фур'є цих функцій.

Тут застосовується теорема про згортку. Використовуючи її, отримуємо:

\textbf{(12.9)}

\begin{verbatim}
\begin{aligned}
\end{verbatim}

\texttt{u(x,\ t)}

\begin{verbatim}
&= \mathcal{F}^{-1}\!\left[\Phi(\xi)e^{-\alpha^{2}\xi^{2} t}\right] \\
&= \mathcal{F}^{-1}\!\left[\Phi(\xi)\right] * \mathcal{F}^{-1}\!\left[e^{-\alpha^{2}\xi^{2} t}\right] \\
&= \Phi(x) * \left[\frac{1}{2\alpha \sqrt{\pi t}} e^{-x^{2}/4\alpha^{2} t}\right] \\
&= \frac{1}{2\alpha \sqrt{\pi t}} \int_{-\infty}^{+\infty} \Phi(\xi)\,e^{-(x-\xi)^{2}/4\alpha^{2} t}\,d\xi.
\end{aligned}
\end{verbatim}

\subsection{Інтерпретація
результату}\label{ux456ux43dux442ux435ux440ux43fux440ux435ux442ux430ux446ux456ux44f-ux440ux435ux437ux443ux43bux44cux442ux430ux442ux443}

Формула (12.9) дає розв'язок задачі. Підінтегральна функція складається
з двох множників:

\begin{enumerate}
\def\labelenumi{\arabic{enumi}.}
\tightlist
\item
  Початкова температура \texttt{q(x)}\\
\item
  Функція\\
  \texttt{G(x,\ t)\ =\ \textbackslash{}frac\{1\}\{2\textbackslash{}alpha\ \textbackslash{}sqrt\{\textbackslash{}pi\ t\}\}\ e\^{}\{-(x-\textbackslash{}xi)\^{}2/4\textbackslash{}alpha\^{}2\ t\},}
  яка називається \textbf{функцією Гріна} (або функцією джерела).
\end{enumerate}

Функція \texttt{G(x,t)} описує відгук системи на початковий
температурний одиничний імпульс.

\begin{figure}
\centering
\pandocbounded{\includegraphics[keepaspectratio,alt={Рис. 12.2}]{images/_page_94_Figure_9.jpeg}}
\caption{Рис. 12.2}
\end{figure}

Рис. 12.2. Відгук \texttt{G(x,t)} на температурний імпульс у точці
\texttt{x=\textbackslash{}xi}.

Графік \texttt{G(x,t)} подібний до нормального розподілу, причому час
\texttt{t} відіграє роль дисперсії. При малих значеннях \texttt{t}
функція вузька і висока, а при збільшенні \texttt{t} --- розширюється і
«розмивається».

Ідея інтерпретації:\\
початкова температура \texttt{u(x,0)=\textbackslash{}phi(x)} може
розглядатися як неперервна сукупність точкових імпульсів. Кожен імпульс
дає внесок \texttt{G(x,t)}, а загальний розв'язок отримується як
інтеграл усіх внесків.

Це є проявом \textbf{принципу суперпозиції}.

\subsection{ЗАУВАЖЕННЯ}\label{ux437ux430ux443ux432ux430ux436ux435ux43dux43dux44f-4}

Головний недолік перетворення Фур'є полягає в тому, що воно не
застосовується до всіх функцій.

Наприклад:

\begin{verbatim}
f(x)=\text{const}, \qquad f(x)=e^x, \qquad f(x)=\sin x
\end{verbatim}

не мають фур'є-образу, оскільки інтеграл

\begin{verbatim}
\mathcal{F}[f] = \frac{1}{\sqrt{2\pi}} \int_{-\infty}^{+\infty} f(x)e^{-i\xi x}\,dx
\end{verbatim}

розходиться.

\texttt{Перетворення\ Фур\textquotesingle{}є\ застосовується\ лише\ до\ функцій,\ які\ достатньо\ швидко\ прямують\ до\ нуля\ при\ \textquotesingle{}\textbar{}x\textbar{}\textbackslash{}to\textbackslash{}infty\textquotesingle{}.}

Крім того, не можна застосовувати перетворення Фур'є за часом для задач
із початковими умовами, оскільки змінна часу обмежена:

\begin{verbatim}
0 < t < \infty.
\end{verbatim}

\section{ЗАДАЧІ}\label{ux437ux430ux434ux430ux447ux456-1}

\begin{enumerate}
\def\labelenumi{\arabic{enumi}.}
\tightlist
\item
  Найдите фурье-образ функции
\end{enumerate}

\begin{verbatim}
f(x)=
\begin{cases}
0, & x \le 0, \\
e^{-x}, & x > 0.
\end{cases}
\end{verbatim}

використовуючи таблиці, наведені в додатку. - 2. На 10м переконайтеся,
що прямі та обернені перетворення Фур'є є лінійними перетвореннями. - 3.
Використовуючи перетворення Фур'є, розв'яжете задачу Коші

\begin{verbatim}
\begin{aligned}
(\text{РЧП})\quad & u_t = \alpha^2 u_{xx}, && -\infty < x < \infty,\\
(\text{ПУ})\quad & u(x,0)=e^{-x^2}, && -\infty < x < \infty.
\end{aligned}
\end{verbatim}

\begin{enumerate}
\def\labelenumi{\arabic{enumi}.}
\setcounter{enumi}{3}
\tightlist
\item
  Перевірте справедливість співвідношень
\end{enumerate}

\begin{verbatim}
\begin{aligned}
\mathcal{F}\left[u_{x}\right] &= i\xi \mathcal{F}\left[u\right], \\
\mathcal{F}\left[u_{xx}\right] &= -\xi^{2} \mathcal{F}\left[u\right].
\end{aligned}
\end{verbatim}

ВКАЗІВКА. Скористайтеся інтегруванням частинами. 5. Перевірити, чи можна
записати згортку двох функцій f і g у двох еквівалентних формах

\texttt{(f\ *\ g)(x)\ =\ \textbackslash{}frac\{1\}\{\textbackslash{}sqrt\{2\textbackslash{}pi\}\}\ \textbackslash{}int\_\{-\textbackslash{}infty\}\^{}\{+\textbackslash{}infty\}\ f(\textbackslash{}xi)\textbackslash{},\ g(x\ -\ \textbackslash{}xi)\textbackslash{},\ d\textbackslash{}xi}
або

\begin{verbatim}
(f * g)(x) = \frac{1}{\sqrt{2\pi}} \int_{-\infty}^{+\infty} f(x-\xi)\, g(\xi)\, d\xi.
# Лекція 13. ПЕРЕТВОРЕННЯ ЛАПЛАСА МЕТА ЛЕКЦІЇ: Ввести перетворення Лапласа, ознайомити читача з формулою оберненого перетворення та показати, як ці інструменти застосовуються до диференціальних рівнянь.

Перетворення Лапласа має важливу перевагу над перетворенням Фур'є: завдяки згасаючому множнику e^{-st} воно застосовується до значно ширшого класу функцій. Крім того, перетворення Лапласа природно працює на півосі 0 \le t < \infty, що особливо зручно для задач з початковими умовами.

# Властивості перетворення Лапласа Властивість 1 (Пряме і обернене перетворення)
Перетворення Лапласа та формула його інверсії мають вигляд
\end{verbatim}

\texttt{\textbackslash{}mathscr\{L\}{[}f{]}\ =\ F(s)\ =\ \textbackslash{}int\_\{0\}\^{}\{\textbackslash{}infty\}\ f(t)\ \textbackslash{},\ e\^{}\{-st\}\ \textbackslash{},\ dt}
\texttt{\textbackslash{}quad\ \textbackslash{}text\{(перетворення\ Лапласа).\}}

\begin{verbatim}
(13.2)
\end{verbatim}

\texttt{\textbackslash{}mathscr\{L\}\^{}\{-1\}{[}F{]}\ =\ f(t)\ =\ \textbackslash{}frac\{1\}\{2\textbackslash{}pi\ i\}\ \textbackslash{}int\_\{c-i\textbackslash{}infty\}\^{}\{c+i\textbackslash{}infty\}\ F(s)\ \textbackslash{},\ e\^{}\{st\}\ \textbackslash{},\ ds}
\texttt{\textbackslash{}quad\ \textbackslash{}text\{(формула\ інверсії).\}}

\begin{verbatim}
Перетворення Лапласа має велику перевагу над перетворенням Фур'є: множник e^{-st} забезпечує збіжність для ширшого класу функцій. Більш точний опис цього класу дає наступна теорема.

# Достатні умови існування перетворення Лапласа Нехай
1. для будь-якого додатного числа A функція f(t) є кусково-неперервною на інтервалі 0 \le t \le A;
2. існують такі константи M, a і T, що |f(t)| \le Me^{at} для всіх t > T.
Тоді перетворення Лапласа
\end{verbatim}

\texttt{\textbackslash{}mathscr\{L\}{[}f{]}\ =\ F(s)\ =\ \textbackslash{}int\_0\^{}\textbackslash{}infty\ f(t)e\^{}\{-st\}\textbackslash{},dt}

\begin{verbatim}
існує для всіх s > a.
Наведемо приклади деяких функцій та їхніх перетворень Лапласа.
У визначенні перетворення Лапласа зазвичай вважають, що s — дійсна змінна. Однак часто корисно розширити це визначення на комплексні значення s, а обернене перетворення записати як
\end{verbatim}

\texttt{\textbackslash{}mathscr\{L\}\^{}\{-1\}{[}F{]}\ =\ \textbackslash{}frac\{1\}\{2\textbackslash{}pi\ i\}\ \textbackslash{}int\_\{\textbackslash{}epsilon\ -\ i\textbackslash{}infty\}\^{}\{\textbackslash{}epsilon\ +\ i\textbackslash{}infty\}\ F(s)e\^{}\{st\}\textbackslash{},ds.}

\begin{verbatim}
Використання цієї формули передбачає знайомство з технікою інтегрування вздовж контурів у комплексній площині та теорією дедукцій. Ми не будемо відволікатися на вивчення цих питань тут, і коли йдеться про перетворення Лапласа, ми використаємо таблицю, наведену в додатку.

# Властивість 2 (Перетворення часткових похідних)
Нехай u(x,t) — функція двох змінних, а
\end{verbatim}

\texttt{U(x,s)=\textbackslash{}mathcal\{L\}\textbackslash{}!\textbackslash{}left{[}u(x,t)\textbackslash{}right{]}.}

\begin{verbatim}
Оскільки перетворення Лапласа виконується за змінною t, маємо формули
\end{verbatim}

\texttt{\textbackslash{}begin\{aligned\}}
\texttt{\textbackslash{}mathcal\{L\}\textbackslash{}!\textbackslash{}left{[}u\_t\textbackslash{}right{]}}
\texttt{\&=\textbackslash{}int\_0\^{}\textbackslash{}infty\ u\_t(x,t)e\^{}\{-st\}\textbackslash{},dt}
=sU(x,s)-u(x,0),\textbackslash{}
\texttt{\textbackslash{}mathcal\{L\}\textbackslash{}!\textbackslash{}left{[}u\_\{tt\}\textbackslash{}right{]}}
\texttt{\&=\textbackslash{}int\_0\^{}\textbackslash{}infty\ u\_\{tt\}(x,t)e\^{}\{-st\}\textbackslash{},dt}
=s\^{}2U(x,s)-su(x,0)-u\_t(x,0),\textbackslash{}
\texttt{\textbackslash{}mathcal\{L\}\textbackslash{}!\textbackslash{}left{[}u\_x\textbackslash{}right{]}}
\texttt{\&=\textbackslash{}int\_0\^{}\textbackslash{}infty\ u\_x(x,t)e\^{}\{-st\}\textbackslash{},dt}
\texttt{=\textbackslash{}frac\{\textbackslash{}partial\ U\}\{\textbackslash{}partial\ x\}(x,s),\textbackslash{}\textbackslash{}}
\texttt{\textbackslash{}mathcal\{L\}\textbackslash{}!\textbackslash{}left{[}u\_\{xx\}\textbackslash{}right{]}}
\texttt{\&=\textbackslash{}int\_0\^{}\textbackslash{}infty\ u\_\{xx\}(x,t)e\^{}\{-st\}\textbackslash{},dt}
\texttt{=\textbackslash{}frac\{\textbackslash{}partial\^{}2\ U\}\{\textbackslash{}partial\ x\^{}2\}(x,s).}
\texttt{\textbackslash{}end\{aligned\}}

\begin{verbatim}
Нижче наведено кілька простих прикладів функцій та їхніх образів Лапласа.

1. f(t)=1, 0<t<\infty.
\end{verbatim}

\texttt{F(s)=\textbackslash{}int\_0\^{}\textbackslash{}infty\ e\^{}\{-st\}\textbackslash{},dt=\textbackslash{}frac\{1\}\{s\}.}

\begin{verbatim}
![](images/_page_99_Figure_3.jpeg)

![](images/_page_99_Figure_5.jpeg)

2. f(t)=e^{2t}, 0<t<\infty.
\end{verbatim}

\texttt{F(s)=\textbackslash{}frac\{1\}\{s-2\},\ \textbackslash{}qquad\ s\textgreater{}2.}

\begin{verbatim}
![](images/_page_99_Figure_7.jpeg)

![](images/_page_99_Figure_9.jpeg)

3. f(t)=\sin(\omega t).
\end{verbatim}

\texttt{F(s)=\textbackslash{}frac\{\textbackslash{}omega\}\{s\^{}2+\textbackslash{}omega\^{}2\}.}

\begin{verbatim}
![](images/_page_99_Figure_12.jpeg)

![](images/_page_99_Figure_13.jpeg)

4. f(t)=e^{t^2} не має перетворення Лапласа.
Рис. 13.1. Графіки деяких функцій та їхніх зображень за Лапласом.
Правила перетворення просторових похідних випливають із теореми про диференціювання інтеграла за параметром:
\end{verbatim}

\texttt{\textbackslash{}frac\{d\}\{dx\}\textbackslash{}int\_\{a\}\^{}\{b\}f\textbackslash{}left(x,\textbackslash{}\ y\textbackslash{}right)dy=\textbackslash{}int\_\{a\}\^{}\{b\}\textbackslash{}frac\{df\}\{dx\}\textbackslash{}left(x,\textbackslash{}\ y\textbackslash{}right)dy,}

\begin{verbatim}
а формули для перетворення похідних u_t і u_{tt} можна отримати інтегруванням частинами.

# Властивість 3 (Теорема про згортку)

Згортка відіграє тут ту ж роль, що і в теорії перетворення Фур'є, але вводиться дещо інакше.

# Визначення скінченної згортки

Скінчена згортка двох функцій f і g визначається формулами

(f * g)(t)
\end{verbatim}

\texttt{=\ \textbackslash{}int\_0\^{}t\ f(\textbackslash{}tau)g(t-\textbackslash{}tau)\textbackslash{},d\textbackslash{}tau}
\texttt{=\ \textbackslash{}int\_0\^{}t\ f(t-\textbackslash{}tau)g(\textbackslash{}tau)\textbackslash{},d\textbackslash{}tau.}

\begin{verbatim}
(Можна довести, що ці інтеграли рівні). Тут інтегрування здійснюється вздовж скінченного інтервалу від 0 до t, а не вздовж нескінченної ( \infty , +\infty ), як це було раніше з визначенням нескінченної згортки.
Ось приклад скінченної згортки двох функцій. Нехай
\end{verbatim}

\texttt{f(t)=t,\textbackslash{}qquad\ g(t)=t.}

\begin{verbatim}
тоді
\end{verbatim}

\texttt{(f\ *\ g)(t)=\textbackslash{}int\_0\^{}t\ \textbackslash{}tau(t-\textbackslash{}tau)\textbackslash{},d\textbackslash{}tau=\textbackslash{}frac\{t\^{}3\}\{6\}.}

\begin{verbatim}
Як і у випадку нескінченної згортки, її найважливіша властивість визначається формулою
(13.4)
\end{verbatim}

\texttt{\textbackslash{}mathcal\{L\}{[}f*g{]}=\textbackslash{}mathcal\{L\}{[}f{]}\textbackslash{}mathcal\{L\}{[}g{]}.}

\begin{verbatim}
або еквівалентна формула
(13.5)
\end{verbatim}

\texttt{\textbackslash{}mathcal\{L\}\^{}\{-1\}\textbackslash{}!\textbackslash{}left(\textbackslash{}mathcal\{L\}{[}f{]}\textbackslash{}mathcal\{L\}{[}g{]}\textbackslash{}right)=f*g.}

\begin{verbatim}
Ці формули дозволяють застосовувати перетворення Лапласа у випадках, коли образ можна подати як добуток \mathcal{L}[f] і \mathcal{L}[g].
Функції F і G легко знайти. Після знаходження функцій f і g можна знайти їхню згортку. Наприклад,
\end{verbatim}

\texttt{\textbackslash{}mathcal\{L\}\^{}\{-1\}\textbackslash{}!\textbackslash{}left(\textbackslash{}frac\{1\}\{s(s\^{}2+1)\}\textbackslash{}right)=1-\textbackslash{}cos\ t.}

\begin{verbatim}
Тут використано пари
\end{verbatim}

\texttt{\textbackslash{}mathcal\{L\}\^{}\{-1\}\textbackslash{}!\textbackslash{}left(\textbackslash{}frac\{1\}\{s\}\textbackslash{}right)=1,\textbackslash{}qquad}
\texttt{\textbackslash{}mathcal\{L\}\^{}\{-1\}\textbackslash{}!\textbackslash{}left(\textbackslash{}frac\{1\}\{s\^{}2+1\}\textbackslash{}right)=\textbackslash{}sin\ t.}

\begin{verbatim}
Тепер ми готові до важливого змішаного виклику.

# Теплопровідність у напівнескінченному середовищі Розглянемо глибокий резервуар з рідиною, і нехай бічна поверхня резервуара буде теплоізольованою. Нехай u_0 буде початковою температурою рідини, а температура повітря над рідиною —

![](images/_page_101_Picture_7.jpeg)

Рис. 13.2. Схематичне представлення проблеми теплопровідності:  
a — тепло тече всередину, якщо u_x(0, t) < 0, і назовні, якщо u_x(0, t) > 0;  
b — резервуар настільки глибокий, що гранична умова внизу не впливає на розв’язок при значеннях x.

кістка дорівнює нулю (орієнтир температур) (див. рис. 13.2).  
Наше завдання — знайти температурне поле в рідині на різних глибинах і в різний час.  
Іншими словами, нам потрібно розв’язати проблему

(13.6)
\end{verbatim}

\texttt{\textbackslash{}begin\{aligned\}}
\texttt{(\textbackslash{}text\{УЧП\})\textbackslash{}quad\ \&\ u\_t\ =\ u\_\{xx\},\ \&\&\ 0\ \textless{}\ x\ \textless{}\ \textbackslash{}infty,\textbackslash{};\ 0\ \textless{}\ t\ \textless{}\ \textbackslash{}infty,\ \textbackslash{}\textbackslash{}}
\texttt{(\textbackslash{}Gamma\ \textbackslash{}text\{У\})\textbackslash{}quad\ \&\ u\_x(0,t)\ -\ u(0,t)\ =\ 0,\ \&\&\ 0\ \textless{}\ t\ \textless{}\ \textbackslash{}infty,\ \textbackslash{}\textbackslash{}}
\texttt{(\textbackslash{}text\{НУ\})\textbackslash{}quad\ \&\ u(x,0)\ =\ u\_0,\ \&\&\ 0\ \textless{}\ x\ \textless{}\ \textbackslash{}infty.}
\texttt{\textbackslash{}end\{aligned\}}

\begin{verbatim}
Щоб розвинути це завдання, ми можемо відтворити Лапласа у мінливу t. Відновити повагу, щоб можна було зупинити перетворення Лапласа на зміну x, змінити з 0 на \infty . Обробка розмноження

Ми отримуємо звичайне диференціальне рівняння за змінною x.

**(13.7)**
\end{verbatim}

\texttt{sU(x)\ -\ u\_0\ =\ \textbackslash{}frac\{d\^{}2U\}\{dx\^{}2\},\ \textbackslash{}quad\ 0\ \textless{}\ x\ \textless{}\ \textbackslash{}infty.}

\begin{verbatim}
Гранична умова:
\end{verbatim}

\texttt{\textbackslash{}frac\{dU\}\{dx\}(0)\ =\ U(0).}

\begin{verbatim}
(Ми перетворили лише рівняння та граничну умову, а не початкову умову.)

Маємо звичайне диференціальне рівняння другого порядку з однією граничною умовою при x=0. Насправді існує ще одна гранична умова, яка визначається фізичними міркуваннями:
\end{verbatim}

\texttt{U(x)\ \textbackslash{}to\ 0\ \textbackslash{}quad\ \textbackslash{}text\{при\}\ \textbackslash{}quad\ x\ \textbackslash{}to\ +\textbackslash{}infty.}

\begin{verbatim}
Зверніть увагу, що для спрощення позначення параметр s опущено в аргументах функції U(x).

Щоб розв’язати задачу (13.7), запишемо загальний розв’язок (сума однорідного і часткового):
\end{verbatim}

\texttt{U(x)\ =\ c\_1\ e\^{}\{\textbackslash{}sqrt\{s\}x\}\ +\ c\_2\ e\^{}\{-\textbackslash{}sqrt\{s\}x\}\ +\ \textbackslash{}frac\{u\_0\}\{s\}.}

\begin{verbatim}
Підставляючи цей вираз у граничні умови, визначаємо сталі c_1 і c_2.

Оскільки при x \to +\infty температура не повинна зростати нескінченно, маємо:
\end{verbatim}

c\_1 = 0.

\begin{verbatim}
Після використання граничної умови при x=0 отримуємо остаточний вираз для U(x):

**(13.8)**
\end{verbatim}

\texttt{U(x)\ =\ \textbackslash{}frac\{u\_0\}\{s\}\ \textbackslash{}left(\ \textbackslash{}frac\{e\^{}\{-\textbackslash{}sqrt\{s\}x\}\}\{1\ +\ \textbackslash{}sqrt\{s\}\}\ \textbackslash{}right)\ +\ \textbackslash{}frac\{1\}\{s\}.}

\begin{verbatim}
Залишається зробити останній крок. Щоб визначити температурне поле, потрібно обчислити:
\end{verbatim}

\texttt{u(x,\ t)\ =\ \textbackslash{}mathcal\{L\}\^{}\{-1\}{[}U(x,\ s){]}.}

\begin{verbatim}
Ми повертаємося до запису U(x, s) із явним аргументом s.

Використовуючи таблиці оберненого перетворення Лапласа, отримуємо:

**(13.9)**
\end{verbatim}

\texttt{u(x,\ t)\ =\ u\_0\ -\ u\_0\ \textbackslash{}left{[}\ \textbackslash{}operatorname\{erfc\}\textbackslash{}!\textbackslash{}left(\ \textbackslash{}frac\{x\}\{2\textbackslash{}sqrt\{t\}\}\ \textbackslash{}right)\ +\ \textbackslash{}operatorname\{erfc\}\textbackslash{}!\textbackslash{}left(\ \textbackslash{}frac\{x\}\{2\textbackslash{}sqrt\{t\}\}\ +\ \textbackslash{}sqrt\{t\}\ \textbackslash{}right)\ e\^{}\{x+t\}\ \textbackslash{}right{]}.}

\begin{verbatim}
де
\end{verbatim}

\texttt{\textbackslash{}operatorname\{erfc\}(x)\ =\ \textbackslash{}frac\{2\}\{\textbackslash{}sqrt\{\textbackslash{}pi\}\}\ \textbackslash{}int\_\{x\}\^{}\{\textbackslash{}infty\}\ e\^{}\{-\textbackslash{}xi\^{}2\}\ d\textbackslash{}xi}

\begin{verbatim}
— функція, спряжена до інтеграла ймовірностей.
Графіки цієї функції показані на рис. 13.3.

![Рис. 13.3](images/_page_102_Figure_7.jpeg)

Рис. 13.3. Графік ймовірнісного інтеграла та його доповнюючої функції.

Аналіз цього рішення показує розподіл температури всередині напівнескінченного середовища в різні моменти часу (див. рис. 13.4).

![Рис. 13.4](images/_page_103_Figure_7.jpeg)

# ПРИМІТКИ

1. Перетворення Лапласа можна звести до диференціального рівняння з частковими похідними або використовувати метод розділення змінних.

2. Для розв’язання задач також використовуються перетворення Ганкеля та Мелліна. Наприклад:
\end{verbatim}

\texttt{\textbackslash{}mathscr\{L\}{[}y\textquotesingle{}{]}\ =\ s\textbackslash{}mathscr\{L\}{[}y{]}\ -\ y(0),}
\texttt{H{[}y{]}\ =\ \textbackslash{}int\_\{0\}\^{}\{\textbackslash{}infty\}\ r\ J\_0(\textbackslash{}xi\ r)\ y(r)\textbackslash{},dr,}
\texttt{H\textbackslash{}left{[}y\textquotesingle{}\textquotesingle{}(r)+\textbackslash{}frac\{1\}\{r\}y\textquotesingle{}(r)\textbackslash{}right{]}\ =\ -\textbackslash{}xi\^{}2\ H{[}y{]}.}

\begin{verbatim}
# ЗАВДАННЯ

1. Перевірити формулу:
\end{verbatim}

\texttt{\textbackslash{}mathscr\{L\}{[}u\_t(x,t){]}\ =\ sU(x,s)\ -\ u(x,0).}

\begin{verbatim}
2. Розв’язати задачу:
\end{verbatim}

\texttt{\textbackslash{}begin\{aligned\}}
\texttt{u\_t\ \&=\ \textbackslash{}alpha\^{}2\ u\_\{xx\},\ \textbackslash{}\textbackslash{}}
\texttt{u(x,0)\ \&=\ \textbackslash{}sin\ x.}
\texttt{\textbackslash{}end\{aligned\}}

\begin{verbatim}
3. Розв’язати задачу:
\end{verbatim}

\texttt{\textbackslash{}begin\{aligned\}} u\_t \&= u\_\{xx\},
\textbackslash{}
\texttt{u(0,t)\ \&=\ \textbackslash{}sin\ t,\ \textbackslash{}\textbackslash{}}
u(x,0) \&= 0. \texttt{\textbackslash{}end\{aligned\}}

\begin{verbatim}
4. Розв’язати задачу:
\end{verbatim}

\texttt{\textbackslash{}frac\{d\^{}2U\}\{dx\^{}2\}\ -\ sU\ =\ A,\ \textbackslash{}quad\ 0\ \textless{}\ x\ \textless{}\ 1,}
\texttt{\textbackslash{}begin\{cases\}}
\texttt{\textbackslash{}frac\{dU\}\{dx\}(0)\ =\ 0,\ \textbackslash{}\textbackslash{}}
U(1) = 0. \texttt{\textbackslash{}end\{cases\}}

\begin{verbatim}
# Лекція 14. ПРИНЦИП ДЮАМЕЛЯ

МЕТА ЛЕКЦІЇ: Показати, як переклад Лапласа можна використовувати для ідентифікації характеристик диференціальних рівнів. Зокрема, за допомогою алгебранічного перетворення образів Лапласа, розвиток рівня часткового оберненого можна відкинути принципом погляду Дюамеля. Інтерпретація цього принципу в теорії найвищих диференціальних рівнів світу використовується для розробки певних завдань РЧП.

технології.

Давайте подивимось, як трансформація Лапласа здійснюється спробою розвитку часткових диференціальних ліній, що дозволяє нам розкрити фізичну сутність різних розвитків. За допомогою трансформації Лапласа ми вже розробили принцип Дюамеля. Перше, що потрібно зробити — звернутися до принципу Дюамеля, ми бачимо проблему, яка часто є проблемою

# Теплопровідність у стрижні, кінці якого підтримуються при заданій температурі

Дуже часто необхідно знайти температуру всередині області, якщо відомо, що на межі області вона змінюється з часом відповідно до певного закону (граничні умови, залежні від часу).

На лівому кінці підтримується умова u(0,t)=0, а на правому кінці температура задається функцією f(t).

Рис. 14.1. Граничні умови, що залежать від часу.

Наприклад, розглянемо термоізольований стержень, правий кінець якого підтримується у f(t) (див. рисунок 14.1). Потрібно вирішити цю проблему

(14.1)
\end{verbatim}

\texttt{u\_t\ =\ u\_\{xx\},\ \textbackslash{}quad\ 0\ \textless{}\ x\ \textless{}\ 1,\ \textbackslash{}quad\ 0\ \textless{}\ t\ \textless{}\ \textbackslash{}infty.}

\begin{verbatim}
Граничні умови:
\end{verbatim}

\texttt{\textbackslash{}begin\{cases\}} u(0, t) = 0, \textbackslash{}
u(1, t) = f(t), \texttt{\textbackslash{}end\{cases\}}
\texttt{\textbackslash{}quad\ 0\ \textless{}\ t\ \textless{}\ \textbackslash{}infty.}

\begin{verbatim}
Початкова умова:
\end{verbatim}

\texttt{u(x,\ 0)\ =\ 0,\ \textbackslash{}quad\ 0\ \textbackslash{}le\ x\ \textbackslash{}le\ 1.}

\begin{verbatim}
Сподіваємося, що розв'язання цієї задачі (14.1) буде легким, оскільки ми знаємо розв'язок простішої задачі, де температура в кінці є сталою:

Насправді, якщо обидві задачі ((14.1) і (14.2)) розв'язуються одночасно за допомогою перетворення Лапласа, то отримаємо визначний результат: розв'язок задачі (14.1) виражається через розв'язання задачі (14.2).

Отже, розв'язуючи обидві задачі одночасно, ми отримуємо

# Легке завдання (14.2) (Постійні солдати)

# Застосування перетворення Лапласа до задачі (14.2)
\end{verbatim}

\texttt{\textbackslash{}frac\{d\^{}2\ W\}\{dx\^{}2\}\ -\ s\ W(x)\ =\ 0,}
\texttt{W(0)\ =\ 0,\ \textbackslash{}quad\ W(1)\ =\ \textbackslash{}frac\{1\}\{s\}.}

\begin{verbatim}
## Розв’язання ОДУ
\end{verbatim}

\texttt{W(x,\ s)\ =\ \textbackslash{}frac\{1\}\{s\}\ \textbackslash{}cdot\ \textbackslash{}frac\{\textbackslash{}sinh(x\ \textbackslash{}sqrt\{s\})\}\{\textbackslash{}sinh(\textbackslash{}sqrt\{s\})\}.}

\begin{verbatim}
## Пошук оберненого перетворення
\end{verbatim}

\texttt{w(x,\ t)\ =\ x\ +\ \textbackslash{}frac\{2\}\{\textbackslash{}pi\}\ \textbackslash{}sum\_\{n=1\}\^{}\{\textbackslash{}infty\}\ \textbackslash{}frac\{(-1)\^{}n\}\{n\}\ e\^{}\{-(n\textbackslash{}pi)\^{}2\ t\}\ \textbackslash{}sin(n\textbackslash{}pi\ x).}

\begin{verbatim}
(це розв’язок задачі з постійними граничними умовами)

# Важке завдання (14.1) (граничні умови, що залежать від часу)

## Застосування перетворення Лапласа до задачі (14.1)
\end{verbatim}

\texttt{\textbackslash{}frac\{d\^{}2\ U\}\{dx\^{}2\}\ -\ s\ U(x)\ =\ 0,}
\texttt{U(0)\ =\ 0,\ \textbackslash{}quad\ U(1)\ =\ F(s).}
\texttt{U(x,\ s)\ =\ F(s)\textbackslash{}cdot\ \textbackslash{}frac\{\textbackslash{}sinh(x\ \textbackslash{}sqrt\{s\})\}\{\textbackslash{}sinh(\textbackslash{}sqrt\{s\})\}.}

\begin{verbatim}
## Представлення через оператор Лапласа

Запишемо:
\end{verbatim}

\texttt{U(x,\ s)\ =\ F(s)\ \textbackslash{}cdot\ \textbackslash{}left(\ \textbackslash{}mathcal\{L\}{[}w\_t{]}\ \textbackslash{}right).}

\begin{verbatim}
Використовуючи співвідношення:
\end{verbatim}

\texttt{\textbackslash{}mathcal\{L\}{[}w\_t{]}\ =\ s\ W\ -\ w(x,\ 0),}

\begin{verbatim}
отримуємо зв’язок для подальшого оберненого перетворення.

## Обернене перетворення Лапласа
\end{verbatim}

\texttt{u(x,\ t)\ =\ \textbackslash{}mathcal\{L\}\^{}\{-1\}\textbackslash{}\{F(s)\textbackslash{}cdot\ \textbackslash{}mathcal\{L\}{[}w\_t{]}\textbackslash{}\}.}

\begin{verbatim}
За теоремою про згортку:
\end{verbatim}

\texttt{u(x,\ t)\ =\ \textbackslash{}mathcal\{L\}\^{}\{-1\}{[}F(s){]}\ *\ \textbackslash{}mathcal\{L\}\^{}\{-1\}{[}w\_t{]}.}

\begin{verbatim}
Отже,
\end{verbatim}

u(x, t) = f(t) * w\_t(t),
\texttt{u(x,\ t)\ =\ \textbackslash{}int\_0\^{}t\ w\_t(x,\ t-\textbackslash{}tau)\textbackslash{},\ f(\textbackslash{}tau)\textbackslash{},\ d\textbackslash{}tau.}

\begin{verbatim}
Після інтегрування частинами:
\end{verbatim}

\texttt{u(x,\ t)\ =\ \textbackslash{}int\_\{0\}\^{}\{t\}\ w(x,\ t-\textbackslash{}tau)\textbackslash{},\ f\textquotesingle{}(\textbackslash{}tau)\textbackslash{},\ d\textbackslash{}tau\ +\ f(0)\textbackslash{},\ w(x,\ t).}

\begin{verbatim}
## Висновок

Розв’язок задачі з граничними умовами, що змінюються в часі, можна виразити через розв’язок задачі з постійними граничними умовами.
(14.3)

Плутанина у наступних формулах. Записати диференціальне рівняння у частинних похідних у нових координатах (\xi, \eta), ми скористаємося наступними формулами переходу.

Наприклад, якщо \frac{dy}{dx} = \frac{B}{2A} = 3, то y - 3x = c, а отже  
\xi(x, y) = y - 3x. Ця функція задовольняє умову  
\frac{\xi_x}{\xi_y} = -3 і, відповідно, занулює коефіцієнт \overline{A}.

Отже, половина роботи виконана: ми знайшли нову координату  
\xi = \xi(x, y), яка забезпечує рівність \overline{A} = 0.  
Залишається знайти \eta(x, y) так, щоб коефіцієнт \overline{B} також дорівнював нулю.

Виявляється, що якщо вибрати B^2 - 4AC = 0 і \xi так, щоб \overline{A} = 0, то коефіцієнт \overline{B} автоматично стане нульовим.

Коефіцієнт \overline{B} задається виразом:
\end{verbatim}

\texttt{\textbackslash{}overline\{B\}\^{}2\ -\ 4\textbackslash{}overline\{A\}\textbackslash{}overline\{C\}\ =\ (B\^{}2\ -\ 4AC)\textbackslash{},(\textbackslash{}xi\_x\ \textbackslash{}eta\_y\ -\ \textbackslash{}xi\_y\ \textbackslash{}eta\_x)\^{}2}
\texttt{\textbackslash{}overline\{B\}\ =\ 2\textbackslash{}left(\ \textbackslash{}overline\{A\}\textbackslash{},\textbackslash{}xi\_x\ +\ \textbackslash{}overline\{C\}\textbackslash{},\textbackslash{}xi\_y\ \textbackslash{}right)}
\texttt{\textbackslash{}left(\ \textbackslash{}overline\{A\}\textbackslash{},\textbackslash{}eta\_x\ +\ \textbackslash{}overline\{C\}\textbackslash{},\textbackslash{}eta\_y\ \textbackslash{}right).}

\begin{verbatim}
Функціональна діаграма робить ці формули зрозумілішими. Вона особливо корисна при обчисленні частинних похідних функцій \xi та \eta.

## Подальші перетворення

Підставляємо знайдені вирази u_t, u_x та u_{xx} у вихідне рівняння:
\end{verbatim}

\texttt{\textbackslash{}overline\{B\}\ =\ 2\textbackslash{}sqrt\{\textbackslash{}overline\{A\}\}\ \textbackslash{}left{[}\ \textbackslash{}sqrt\{\textbackslash{}overline\{A\}\}\textbackslash{},\ \textbackslash{}eta\_x\ +\ \textbackslash{}sqrt\{\textbackslash{}overline\{C\}\}\textbackslash{},\ \textbackslash{}eta\_y\ \textbackslash{}right{]}.}

\begin{verbatim}
Звідси отримуємо:
\end{verbatim}

u\_\{xx\} + 2u\_\{xy\} + u\_\{yy\} = 0.

\begin{verbatim}
## Нова задача

Отже, нова задача у змінних \xi та \eta має вигляд:
\end{verbatim}

\texttt{\textbackslash{}frac\{dy\}\{dx\}\ =\ -\textbackslash{}frac\{\textbackslash{}xi\_x\}\{\textbackslash{}xi\_y\}\ =\ \textbackslash{}frac\{B\}\{2A\}\ =\ 1.}

\begin{verbatim}
(ПУ)
\end{verbatim}

\texttt{u(\textbackslash{}xi,\ 0)\ =\ 1\ -\ H(\textbackslash{}xi),\ \textbackslash{}quad\ -\textbackslash{}infty\ \textless{}\ \textbackslash{}xi\ \textless{}\ \textbackslash{}infty.}

\begin{verbatim}
Оскільки \xi = x при t = 0, граничні умови збігаються зі старими.

## Розв’язок

Після перетворення:
\end{verbatim}

\texttt{u(\textbackslash{}xi,\ \textbackslash{}eta)\ =\ \textbackslash{}eta\ f(\textbackslash{}xi)\ +\ g(\textbackslash{}xi),}

\begin{verbatim}
переходимо до змінних x і y:
\end{verbatim}

u(x, y) = y f(y - x) + g(y - x).

\begin{verbatim}
Для конкретного випадку:
\end{verbatim}

\texttt{u(x,\ y)\ =\ y\ \textbackslash{}sin(y\ -\ x)\ +\ (y\ -\ x)\^{}2.}

\begin{verbatim}
## Параметри рівняння
\end{verbatim}

A u\_\{xx\} + B u\_\{xy\} + C u\_\{yy\} + D u\_x + E u\_y + F u = G
\texttt{u\_t\ =\ u\_\{xx\}\ -\ 2u\_x,\ \textbackslash{}quad\ -\textbackslash{}infty\ \textless{}\ x\ \textless{}\ \textbackslash{}infty,\ \textbackslash{}quad\ 0\ \textless{}\ t\ \textless{}\ \textbackslash{}infty,}
\texttt{u(x,\ 0)\ =\ \textbackslash{}sin\ x,\ \textbackslash{}quad\ -\textbackslash{}infty\ \textless{}\ x\ \textless{}\ \textbackslash{}infty.}

\begin{verbatim}
## Інші форми запису
\end{verbatim}

\texttt{\textbackslash{}frac\{dy\}\{dx\}\ =\ \textbackslash{}frac\{B\ -\ \textbackslash{}sqrt\{B\^{}2\ -\ 4AC\}\}\{2A\}.}
\texttt{\textbackslash{}xi(x,\ y)\ =\ \textbackslash{}text\{const\},\ \textbackslash{}quad\ \textbackslash{}eta(x,\ y)\ =\ \textbackslash{}text\{const\}.}
\texttt{\textbackslash{}frac\{dy\}\{dx\}\ =\ -\textbackslash{}sqrt\{-4x\^{}2\}\ =\ -2ix.}

\begin{verbatim}
## Додаткові зауваження

Перехід до рухомої системи координат дозволяє усунути член V u_x у рівнянні конвективної дифузії.

Після розв’язання нової задачі:
\end{verbatim}

\texttt{\textbackslash{}text\{(узагальнене\ рівняння\ у\ нових\ змінних)\}}

\begin{verbatim}
де \phi позначає праві частини рівнянь у загальному вигляді.

Замість доведення складних перетворень, розглянемо простий приклад.

# Зведення рівняння y^2 u_{xx} + x^2 u_{yy} = 0 до канонічної форми

Розглянемо рівняння:
\end{verbatim}

y\^{}2 u\_\{xx\} + x\^{}2 u\_\{yy\} = 0,

\begin{verbatim}
де A = y^2, B = 0, C = x^2, D = E = F = G = 0.

Дискримінант:
\end{verbatim}

B\^{}2 - 4AC = -4x\^{}2 y\^{}2,

\begin{verbatim}
і ми зведемо це рівняння до канонічної форми в першому квадранті x > 0, y > 0.

## КРОК 1. Перше перетворення

Характеристичне рівняння:
\end{verbatim}

\texttt{\textbackslash{}frac\{dy\}\{dx\}\ =\ \textbackslash{}frac\{B\ +\ \textbackslash{}sqrt\{B\^{}2\ -\ 4AC\}\}\{2A\}\ =\ \textbackslash{}frac\{\textbackslash{}sqrt\{-4x\^{}2\ y\^{}2\}\}\{2y\^{}2\}\ =\ i\ \textbackslash{}frac\{x\}\{y\}.}

\begin{verbatim}
Інтегруючи, отримуємо:
\end{verbatim}

\texttt{y\^{}2\ +\ i\ x\^{}2\ =\ \textbackslash{}text\{const\}.}

\begin{verbatim}
Вводимо змінні:
\end{verbatim}

\texttt{\textbackslash{}xi\ =\ y\^{}2\ +\ i\ x\^{}2,\ \textbackslash{}quad\ \textbackslash{}eta\ =\ y\^{}2\ -\ i\ x\^{}2.}

\begin{verbatim}
(можна міняти місцями \xi та \eta)

## КРОК 2. Друга трансформація

Після другої заміни отримуємо:
\end{verbatim}

\texttt{\textbackslash{}begin\{cases\}}
\texttt{u\_\{tt\}\ =\ \textbackslash{}alpha\^{}2\ u\_\{xx\}\ +\ F(x,t),\ \textbackslash{}\textbackslash{}}
\texttt{u\_\{tt\}\ =\ \textbackslash{}alpha\^{}2\ u\_\{xx\}\ -\ \textbackslash{}beta\ u\_t\ -\ \textbackslash{}lambda\ u\ +\ F(x,t).}
\texttt{\textbackslash{}end\{cases\}}

\begin{verbatim}
(де \beta = x^2 — уявна частина відповідних змінних)

## КРОК 3. Отримання нового рівняння

Канонічна форма:
\end{verbatim}

\texttt{\textbackslash{}overline\{A\}u\_\{\textbackslash{}xi\textbackslash{}xi\}\ +\ \textbackslash{}overline\{B\}u\_\{\textbackslash{}xi\textbackslash{}eta\}\ +\ \textbackslash{}overline\{C\}u\_\{\textbackslash{}eta\textbackslash{}eta\}\ +\ \textbackslash{}overline\{D\}u\_\{\textbackslash{}xi\}\ +\ \textbackslash{}overline\{E\}u\_\{\textbackslash{}eta\}\ +\ \textbackslash{}overline\{F\}u\ =\ \textbackslash{}overline\{G\}.}

\begin{verbatim}
## Хвильове рівняння (гіперболічний тип)

Одновимірне хвильове рівняння:
\end{verbatim}

u\_\{tt\} = c\^{}2 u\_\{xx\}.

\begin{verbatim}
## Вібрації струни

Сила натягу:
\end{verbatim}

\texttt{\textbackslash{}text\{Сила\}\ \textbackslash{}approx\ T\ \textbackslash{}left{[}\ u\_x(x\ +\ \textbackslash{}Delta\ x,\ t)\ -\ u\_x(x,\ t)\ \textbackslash{}right{]}.}

\begin{verbatim}
Закон Ньютона приводить до:
\end{verbatim}

\texttt{\textbackslash{}rho\ u\_\{tt\}\ =\ T\ u\_\{xx\}\ +\ F(x,t)\ -\ \textbackslash{}beta\ u\_t\ -\ \textbackslash{}gamma\ u.}

\begin{verbatim}
Поділивши на \rho, отримуємо телеграфне рівняння.

## Примітки

1. Класифікація рівнянь:
   - параболічні,
   - гіперболічні,
   - еліптичні.

2. Це важливо, бо кожен тип описує різні фізичні процеси.

# ЗАВДАННЯ

1. Які з цих параболічних і еліптичних рівнянь записані у канонічній формі?

a) u_t = u_{xx} + hu,

б) u_{xy} + u_{xx} + 3u = \sin x,

в) u_{xx} + 2u_{yy} = 0,

г) u_{xx} = \sin^2 x?

2. Перетворіть параболічне рівняння u_{xx} + 2u_{xy} + u_{yy} + u = 2 до канонічної форми.

3. Перетворіть еліптичне рівняння u_{xx} + 2u_{yy} + x^2 u_x = e^{-x^2/2} до канонічної форми.

# МЕТОД МОНТЕ-КАРЛО (ВСТУП)

МЕТА ЛЕКЦІЇ: Пояснити основні ідеї, що лежать в основі методу Монте-Карло, і показати, як цей метод можна використовувати для розв’язання різних задач. Виявляється, що за допомогою випадкових експериментів можна змоделювати (загалом, на комп’ютері) наближені розв’язки реальних задач. Як простий приклад можна вказати обчислення інтеграла
\end{verbatim}

\texttt{\textbackslash{}int\_a\^{}b\ f(x)\textbackslash{},dx}

\begin{verbatim}
на основі інтуїтивних міркувань.

# Інтуїтивна інтерпретація хвильового рівняння

Читач може запитати: чому рівняння типу (16.2) описує вібрації струни? Щоб зрозуміти це, необхідно розглянути фізичний зміст величин.

![](images/_page_121_Picture_9.jpeg)

Рівняння (16.2) можна інтерпретувати так: прискорення струни, викликане натягом, у кожній точці пропорційне вигину струни:
\end{verbatim}

\texttt{u\_\{tt\}\ =\ \textbackslash{}alpha\^{}2\ u\_\{xx\},\ \textbackslash{}quad\ \textbackslash{}alpha\^{}2\ =\ \textbackslash{}frac\{T\}\{\textbackslash{}rho\}}

\begin{verbatim}
(див. рис. 13.3).

# Пояснення

1. Хвильове рівняння
\end{verbatim}

\texttt{u\_\{tt\}\ =\ \textbackslash{}alpha\^{}2\ u\_\{xx\}}

\begin{verbatim}
також описує поздовжні та крутильні коливання стрижня.

У поздовжніх коливаннях зміщення паралельні осі стрижня, і поздовжнє зміщення позначається як u(x,t). Такі коливання виникають, наприклад, коли ударяють по кінцю стрижня.

# Ймовірнісна інтерпретація (Метод Монте-Карло)

Приклад: кидаємо дротик у одиничний квадрат  
\{(x, y): 0 < x < 1,\; 0 < y < 1\}.

Після приблизно 100 кидків відносна частка точок, що лежать нижче кривої y = x^2, дає наближення інтеграла.

Узагальнення такої гри передбачає генерацію випадкових чисел.

# Загальна ідея методу

Методи Монте-Карло — це група методів розв’язання детермінованих задач, у яких використовуються випадкові величини.

Загальна схема:

Ймовірнісна гра → випадкова модель  
Детермінована задача → шуканий результат  
Наближення: \hat{p} \approx p

# Обчислення інтеграла

Припустимо, потрібно обчислити інтеграл:
\end{verbatim}

\texttt{\textbackslash{}int\_a\^{}b\ f(x)\textbackslash{},dx}

\begin{verbatim}
Для цього вводимо прямокутну область:
\end{verbatim}

\texttt{R\ =\ \textbackslash{}\{(x,y):\ a\ \textbackslash{}le\ x\ \textbackslash{}le\ b,\textbackslash{};\ 0\ \textbackslash{}le\ y\ \textbackslash{}le\ \textbackslash{}max\ f(x)\textbackslash{}\}}

\begin{verbatim}
(див. рис. 42.2)

# Фізичний зміст хвиль

Модуль Юнга k характеризує пружність матеріалу. Чим більший k, тим вищі частоти коливань.

Звукові хвилі — це поздовжні хвилі.

# Узагальнене хвильове рівняння

Якщо густина \rho(x) залежить від координати, хвильове рівняння має вигляд:
\end{verbatim}

\texttt{\textbackslash{}rho(x)\textbackslash{},u\_\{tt\}\ =\ \textbackslash{}frac\{\textbackslash{}partial\}\{\textbackslash{}partial\ x\}\textbackslash{}left(\ T\ \textbackslash{}frac\{\textbackslash{}partial\ u\}\{\textbackslash{}partial\ x\}\ \textbackslash{}right)}

\begin{verbatim}
# Розв’язки хвильового рівняння

Типові розв’язки мають вигляд хвиль, що рухаються:
\end{verbatim}

u(x,t) = e\^{}\{-(x - ct)\^{}2\}

\begin{verbatim}
Хвилі можуть рухатися у різних напрямках зі швидкістю c.

# Варіаційний підхід
\end{verbatim}

\texttt{\textbackslash{}int\_a\^{}b\ \textbackslash{}left(\ \textbackslash{}frac\{\textbackslash{}partial\ F\}\{\textbackslash{}partial\ u\}\ \textbackslash{}eta(x)\ +\ \textbackslash{}frac\{\textbackslash{}partial\ F\}\{\textbackslash{}partial\ u\textquotesingle{}\}\ \textbackslash{}eta\textquotesingle{}(x)\ \textbackslash{}right)\ dx\ =\ 0}

\begin{verbatim}
Звідси отримуємо:
\end{verbatim}

\texttt{\textbackslash{}frac\{\textbackslash{}partial\ F\}\{\textbackslash{}partial\ u\}\ -\ \textbackslash{}frac\{d\}\{dx\}\textbackslash{}left(\ \textbackslash{}frac\{\textbackslash{}partial\ F\}\{\textbackslash{}partial\ u\textquotesingle{}\}\ \textbackslash{}right)\ =\ 0}

\begin{verbatim}
# Висновок

Рівняння хвильового типу описує процеси, які мають вигляд хвиль, що поширюються в просторі.

КРОК 4 (Урахування початкових умов).

Нагадаємо читачеві загальний метод розв'язання задачі Коші в теорії звичайних диференціальних рівнянь. Спочатку знаходять загальний розв'язок рівняння, а потім підставляють його в початкові умови, щоб визначити конкретні значення довільних сталих. У нас ситуація аналогічна. Для розв'язання задачі Коші для хвильового рівняння ми підставимо загальний розв'язок
\end{verbatim}

\texttt{\textbackslash{}frac\{\textbackslash{}partial\ F\}\{\textbackslash{}partial\ y\}\ -\ \textbackslash{}frac\{d\}\{dx\}\ \textbackslash{}left{[}\ \textbackslash{}frac\{\textbackslash{}partial\ F\}\{\textbackslash{}partial\ y\textquotesingle{}\}\ \textbackslash{}right{]}\ =\ 0}

\begin{verbatim}
(який містить дві довільні функції) у початкові умови
\end{verbatim}

\texttt{J{[}y{]}\ =\ \textbackslash{}int\_\{0\}\^{}\{1\}\ {[}y\^{}\{2\}\ +\ y\textquotesingle{}\^{}\{2\}{]}\ dx.}

\begin{verbatim}
u_1(x, 0) = g(x),
\end{verbatim}

щоб знайти конкретні вирази для довільних функцій φ і ψ. Після
підстановки отримуємо

(17.5)

\begin{verbatim}
F_y - \frac{d}{dx} F_{y'} = 0,
\end{verbatim}

Якщо тепер проінтегрувати друге рівняння (17.5), отримаємо алгебраїчну
систему двох рівнянь відносно невідомих функцій
\texttt{\textbackslash{}varphi(x)} і \texttt{\textbackslash{}psi(x)}.
Дійсно, після інтегрування в межах від \texttt{x\_0} до \texttt{x} маємо

(17.6)

\begin{verbatim}
\begin{aligned}
F_y &= 2y, \\
F_{y'} &= 2y'.
\end{aligned}
\end{verbatim}

Тепер розв'язуємо рівняння (17.6) разом із першим рівнянням (17.5) і
отримуємо такі вирази для функцій \texttt{\textbackslash{}varphi(x)} і
\texttt{\textbackslash{}psi(x)}:

(17.7)

\begin{verbatim}
2y - \frac{d}{dx}(2y') = 0,
y'' - y = 0.
\end{verbatim}

Отже, загальний розв'язок задачі (17.1) задається формулою

(17.8)

\begin{verbatim}
\overline{y}(x) = 0.42e^x - 0.42e^{-x}
\end{verbatim}

Мета досягнута. Ми отримали загальний розв'язок задачі (17.1). Розв'язок
у вигляді (17.8) зазвичай називають формулою Д'Аламбера. Читача
запрошують самостійно простежити за межами інтегрування, щоб
переконатися, що вони змінюються з \texttt{x-ct} на \texttt{x+ct}.
Проблему повністю розв'язано.

Перед завершенням лекції наведемо кілька прикладів застосування формули
Д'Аламбера.

\section{Приклади застосування формули
Д'Аламбера}\label{ux43fux440ux438ux43aux43bux430ux434ux438-ux437ux430ux441ux442ux43eux441ux443ux432ux430ux43dux43dux44f-ux444ux43eux440ux43cux443ux43bux438-ux434ux430ux43bux430ux43cux431ux435ux440ux430}

\section{1. Рух синусоїдальних
хвиль}\label{ux440ux443ux445-ux441ux438ux43dux443ux441ux43eux457ux434ux430ux43bux44cux43dux438ux445-ux445ux432ux438ux43bux44c}

Розглянемо початкові умови вигляду

\begin{verbatim}
J[y] = \int_a^b y \sqrt{1 + y'^2} \, dx,
u_1(x, 0) = 0.
\end{verbatim}

За формулою Д'Аламбера отримуємо розв'язок

\begin{verbatim}
\overline{y}(x) = \alpha \operatorname{ch}\left(\frac{x - \beta}{\alpha}\right),
\end{verbatim}

Його можна інтерпретувати так: початкове зміщення струни
\texttt{u(x,\ 0)\ =\ \textbackslash{}sin\ x} розкладається на дві
однакові частини

\begin{verbatim}
\overline{y}(x) = 0.42e^x - 0.42e^{-x}
\end{verbatim}

та \texttt{\textbackslash{}frac\{\textbackslash{}sin\ x\}\{2\}}.

Кожна з цих частин поширюється зі швидкістю \texttt{c} у вигляді хвиль.
Одна хвиля рухається зліва направо, інша --- у протилежному напрямку.
Читача запрошують замислитися, як виглядає отримана хвиля в цьому
випадку.

\section{2. Розповсюдження найпростішого прямокутного
імпульсу}\label{ux440ux43eux437ux43fux43eux432ux441ux44eux434ux436ux435ux43dux43dux44f-ux43dux430ux439ux43fux440ux43eux441ux442ux456ux448ux43eux433ux43e-ux43fux440ux44fux43cux43eux43aux443ux442ux43dux43eux433ux43e-ux456ux43cux43fux443ux43bux44cux441ux443}

У цьому випадку початкові умови задаються так:

\begin{verbatim}
J[y] = \int_0^1 (y^2 + y'^2)\,dx,
u_t(x, 0) = 0.
\end{verbatim}

Оскільки початкове зміщення розкладається на дві півхвилі, що рухаються
у протилежних напрямках, отриманий рух виглядає так, як показано на рис.
17.2.

\pandocbounded{\includegraphics[keepaspectratio]{images/_page_129_Figure_2.jpeg}}

Фіг. 17.2. Початкове збурення розділяється на дві хвилі, що рухаються.

\section{3. Задана початкова
швидкість}\label{ux437ux430ux434ux430ux43dux430-ux43fux43eux447ux430ux442ux43aux43eux432ux430-ux448ux432ux438ux434ux43aux456ux441ux442ux44c}

Тепер припустимо, що в початковий момент струна перебуває у стані
рівноваги. Надамо їй початкову швидкість форми
\texttt{\textbackslash{}sin\ x}, тобто

\begin{verbatim}
\int\limits_{t_1}^{t_2} J[u] = \iint\limits_D F(x, y, u, u_x, u_y)\,dx\,dy,
\end{verbatim}

Розв'язок знаходимо за формулою Д'Аламбера:

\begin{verbatim}
F_u - \frac{\partial}{\partial x}F_{ux} - \frac{\partial}{\partial y}F_{uy} = 0.
\end{verbatim}

Це рішення представляє суму двох хвиль, що рухаються у вигляді
косинусоїдальних хвиль. Читач може самостійно оцінити, чи є цей
результат фізично зрозумілим.

На цьому лекція 17 завершується. У наступній лекції буде показано, як
формула Д'Аламбера інтерпретується в площині змінних \texttt{x} і
\texttt{t}. \# ЗАУВАЖЕННЯ

\begin{enumerate}
\def\labelenumi{\arabic{enumi}.}
\item
  Загальний розв'язок диференціального рівняння в частинних похідних
  другого порядку містить дві довільні функції, тоді як загальний
  розв'язок звичайного диференціального рівняння містить дві довільні
  сталі. Отже, рівняння з частинними похідними має більше розв'язків,
  ніж звичайне диференціальне рівняння.
\item
  Заміна змінних для зведення рівняння з частинними похідними до
  простішого вигляду є одним із загальних методів теорії. Нові
  координати \texttt{(\textbackslash{}xi,\ \textbackslash{}eta)}
  прийнято називати канонічними координатами. Ми будемо звертатися до
  них і надалі, особливо при вивченні гіперболічних рівнянь.
\item
  Метод знаходження загального розв'язку рівняння з частинними похідними
  та підстановки цього розв'язку в початкові умови не є загальним
  методом розв'язання таких рівнянь. Розв'язок, розглянутий у цій
  лекції, є лише прикладом використання цього методу. Зазвичай неможливо
  знайти загальний розв'язок рівняння з частинними похідними, а навіть
  якщо це вдається, підстановка його в граничні умови є досить складною
  задачею.
\end{enumerate}

\section{ЗАВДАННЯ}\label{ux437ux430ux432ux434ux430ux43dux43dux44f-4}

\begin{enumerate}
\def\labelenumi{\arabic{enumi}.}
\item
  Перевірте, що розв'язок Д'Аламбера (17.8) задовольняє початкові умови
  задачі (17.1).
\item
  Підставте (17.7) у загальний розв'язок
  \texttt{u(x,\ t)\ =\ \textbackslash{}psi(x\ -\ ct)\ +\ \textbackslash{}psi(x\ +\ ct)}
  і отримайте формулу Д'Аламбера.
\item
  Знайдіть розв'язок задачі Коші:
\end{enumerate}

\begin{verbatim}
J[y] = \int_{0}^{1} \sqrt{1 + y'^{2}}\,dx.
\end{verbatim}

Як виглядає цей розв'язок у різні моменти часу?

\begin{enumerate}
\def\labelenumi{\arabic{enumi}.}
\setcounter{enumi}{3}
\tightlist
\item
  Знайдіть розв'язок задачі Коші:
\end{enumerate}

(УЧП)

\begin{verbatim}
\int_{t_{1}}^{t_{2}} (KE - PE)\,dt = \frac{1}{2} \int_{t_{1}}^{t_{2}} (m\dot{y}^{2} - ky^{2})\,dt
\end{verbatim}

(НУ)

\begin{verbatim}
J[y] = \int_{0}^{\pi/2} (y'^{2} - y^{2})\,dx
\end{verbatim}

Побудуйте графік розв'язку в різні моменти часу.

\begin{enumerate}
\def\labelenumi{\arabic{enumi}.}
\setcounter{enumi}{4}
\tightlist
\item
  Розв'язуючи спільно рівняння (17.6) і перше рівняння з (17.5),
  отримайте формули (17.7).
\end{enumerate}

\section{\texorpdfstring{Лекція 15. КОНВЕКТИВНИЙ ЧЛЕН \texttt{u\_x} У
ДИФУЗІЙНІЙ
ЗАДАЧІ}{Лекція 15. КОНВЕКТИВНИЙ ЧЛЕН u\_x У ДИФУЗІЙНІЙ ЗАДАЧІ}}\label{ux43bux435ux43aux446ux456ux44f-15.-ux43aux43eux43dux432ux435ux43aux442ux438ux432ux43dux438ux439-ux447ux43bux435ux43d-u_x-ux443-ux434ux438ux444ux443ux437ux456ux439ux43dux456ux439-ux437ux430ux434ux430ux447ux456}

ЦЕЛЬ ЛЕКЦИИ: Показать, каким образом слагаемое \texttt{u\_x} в
диффузионной задаче

\begin{verbatim}
u_t = Du_{xx} - Vu_x
\end{verbatim}

Диффузионный член Конвективный член

описывает явление конвекции. Приведенное выше уравнение конвективной
диффузии возникает во многих задачах и описывает как диффузию, так и
конвекцию. Роль каждого из этих механизмов определяется отношением
коэффициентов D и V.

Поскольку конвекция вещества обусловлена его движением вместе с
окружающей средой, то естественно перейти в систему координат, которая
движется вместе со средой. В результате такого преобразования
конвективный член будет исключен в получившееся уравнение можно решить
одним из известных методов. Преобразовав полученное решение к
неподвижной системе отсчета, получим решение исходной задачи.

До сих пор мы имели дело с теплопроводностью (или некоторым видом
диффузии) в одномерных системах. Рассмотрим теперь задачу о
распределении концентрации некоторого вещества, выходящего из земли и
попадающего в восходящие потоки воздуха. Будем учитывать как диффузию
вещества, так и конвективный перенос воздухом, который движется со
скоростью V. Ясно, что в данном случае конвективный механизм переноса
вещества играет более важную роль, чем диффузионный (роль каждого
механизма определяется отношением коэффициентов диффузии и скорости
воздуха.) Диффузия---это проникновение вещества через воздух, а
конвекция---это перемещение вещества вместе с движущимся воздухом. Наша
задача---вывести уравнение конвективной диффузии

\begin{verbatim}
u_t = Du_{xx} - Vu_x
\end{verbatim}

Чтобы убедиться в правильности уравнения конвективной диффузии, мы
воспользуемся двумя основными фактами.

\begin{enumerate}
\def\labelenumi{\arabic{enumi}.}
\tightlist
\item
  Поток, обусловленный конвекцией.
\end{enumerate}

Конвективный поток вещества слева направо через поперечное сечение в
данной точке равен Vu(x, t), где V--- скорость движения среды (см/с), а
\texttt{u(x,\ t)}--- концентрация вещества (г/см) (рис. 15.1).

\pandocbounded{\includegraphics[keepaspectratio]{images/_page_112_Picture_4.jpeg}}

\begin{verbatim}
\begin{array}{c} \vdots \\ \vdots \\ \vdots \\ \vdots \\ \vdots \\ \vdots \\ \vdots \\ \vdots \\ \vdots \\ u_i + Vu_x = 0 \\ (\textit{SpaBhenue nepenoca nep8ozo nops@ka}) \\ \vdots \\ u(x,0) = 0 \end{array}
\end{verbatim}

\texttt{РИС.\ 15.1.\ Количество\ вещества,\ переносимого\ в\ одну\ секунду\ через\ поперечное\ сечение\ в\ точке\ x\ за\ счет\ конвекции,\ равно\ \ \textquotesingle{}Vu\textbackslash{}left(\ v,\textbackslash{}\ t\ \textbackslash{}right).\textquotesingle{}}

РИС. 15.2. Задача чисто конвективного переноса.

\begin{enumerate}
\def\labelenumi{\arabic{enumi}.}
\setcounter{enumi}{1}
\tightlist
\item
  Поток, обусловленный диффузией.
\end{enumerate}

\texttt{\{\textbackslash{}it\ Диффузионный\ поток\}} вещества слева
направо через поперечное сечение в данной точке равен ---
\texttt{Du\_x(x,\ t)} , где D---коэффициент диффузии.

Если подставить эти два выражения в закон сохранения из лекции 3, то
получим основное уравнение с частными производными

\begin{verbatim}
u_t = Du_{xx} - Vu_x.
\end{verbatim}

Чтобы разобраться в том, как влияет конвективный член на характер
решения, рассмотрим сначала только конвективный механизм переноса.
Типичный пример---задача о сбросе некоторого вещества в реку (скорость
течения реки V). Пусть x--- расстояние вниз по течению от места сброса.
Если вещество не диффундирует в воде, то концентрация \texttt{u(x,\ t)}
вещества является решением следующей задачи (рис. 15.2):

(УЧП)

\begin{verbatim}
u_t = -Vu_x
\end{verbatim}

\begin{verbatim}
,  0 < x < \infty ,  0 < t < \infty ,  (\Gamma Y)   u(0, t) = P \leftarrow  Постоянная концентрация в месте сброса,

(НУ)  u(x, 0) = 0 \leftarrow Чистая река в начальный момент.

Прежде чем решать задачу, следует задуматься над тем, каким должно получиться решение. Очевидно, что загрязнение реки

(количество вещества в воде) будет сначала равно нулю, а затем, поддерживаемое постоянным в точке x=0, будет смещаться вниз по течению со скоростью V. Чтобы убедиться в этом, давайте решим задачу (15.1). Поскольку уравнение (15.1) линейное и граничное условие тоже линейное, можно надеяться на решение этой задачи методом разделения переменных или методом интегральных преобразований. Однако поскольку координата x здесь не ограничена, то метод разделения переменных неприменим. Используем преобразование Лапласа по времени t.

#### Решение уравнения переноса методом преобразования Лапласа

С помощью преобразования Лапласа
\end{verbatim}

\texttt{U(x)\ =\ \textbackslash{}int\_\{0\}\^{}\{\textbackslash{}infty\}\ u(x,\ t)\ e\^{}\{-st\}\ dt}

\begin{verbatim}
задачу конвективного переноса (15.1) можно свести к виду
\end{verbatim}

\texttt{sU(x)\ =\ -V\ \textbackslash{}frac\{dU\}\{dx\},\ \textbackslash{}quad\ 0\ \textless{}\ x\ \textless{}\ \textbackslash{}infty,}
U(0) = P/s.

\begin{verbatim}
Решая эту очень простую задачу Коши, получаем  U(x) = Pe^{-(sx/V)}/s .

По таблицам обратного преобразования находим решение задачи  u(x, t) = \mathcal{L}^{-1}[U] = PH(t-x/V), 

где Н (ξ) -- ступенчатая функция Хевисайда
\end{verbatim}

\texttt{H(\textbackslash{}xi)\ =\ \textbackslash{}begin\{cases\}\ 0,\ \&\ \textbackslash{}xi\ \textless{}\ 0,\ \textbackslash{}\textbackslash{}\ 1,\ \&\ \textbackslash{}xi\ \textbackslash{}geqslant\ 0.\ \textbackslash{}end\{cases\}}

\begin{verbatim}
Следовательно, решению нашей задачи можно придать такую форму:
\end{verbatim}

\texttt{u(x,\ t)\ =\ \textbackslash{}begin\{cases\}\ 0,\ \&t\ \textless{}\ x/V,\ \textbackslash{}\textbackslash{}\ P,\ \&\ t\ \textbackslash{}ge\ x/V.\ \textbackslash{}end\{cases\}}

\begin{verbatim}
Конечно, все это очень просто! Не сложнее, чем сбрасывать что-нибудь на ленту транспортера, а затем наблюдать за его движением. Однако все становится очень интересным, если примесь диффундирует в среде. Чтобы разобраться в том, что про-исходит с движущейся волной при наличии диффузии, решим

следующую задачу

(15.2)
\end{verbatim}

\texttt{\textbackslash{}begin\{aligned\}\ (3\ \textbackslash{}text\{\ HI\})\ \textbackslash{}quad\ u\_t\ \&=\ Du\_\{xx\}\ -\ Vu\_x,\ \textbackslash{}quad\ -\textbackslash{}infty\ \textless{}\ x\ \textless{}\ \textbackslash{}infty,\ \textbackslash{}\textbackslash{}\ (HY)\ \textbackslash{}quad\ u(x,\ 0)\ \textbackslash{}quad\ 1\ -\ H(x),\ \textbackslash{}quad\ -\textbackslash{}infty\ \textless{}\ x\ \textless{}\ \textbackslash{}infty.\ \textbackslash{}end\{aligned\}}

\begin{verbatim}
где, как обычно, H(x) — функция Хевисайда Начальное распределение концентрации изображено на рис 15 4.

![](images/_page_114_Figure_4.jpeg)

РИС. 15.3 Конвективная волна a—передиції фронт примеси,  \delta —на одну единицу вперед волнового фронта

Обратим внимание на то, что в задаче (15.2) мы отодвинули границу в —  \infty  (теперь решается задача Коши), так что она не оказывает теперь влияние на оценку решение. Чтобы решить задачу (15.2), можно воспользоваться преобразованием Лапласа Р И

![](images/_page_114_Figure_7.jpeg)

С. 15 4 Начальные условия в задаче конвективной диффузии.

ло переменной t или преобразованием Фурье по переменной x. Однако в этом случае гораздо интереснее поступить совсем по-другому Введем новую систему координат, которая движется вдоль старой со скоростью V. Другими словами, вместо системы координат, привязанной к берегу реки, мы рассмотрим систему координат, которая движется со скоростью фронта примеси (конечно, при налични диффузии фронт будет «смазан»). С точки зрения математики это означает, что мы заменяем пространственную координату x на новую  \xi = x - Vt .

Теперь ясно, что

когда  \xi = 0 , мы находимся па фронте распространяющейся примеси.

когда  \xi = 1 , мы находимся на одну единицу длины впереди фронта.

когда  \xi = -1 , мы находимся на одну единицу длины позади фронта.

Наша задача — преобразовать исходную задачу с НУ
\end{verbatim}

\texttt{\textbackslash{}begin\{array\}\{lll\}\ (\textbackslash{}mathrm\{Y\}\ +\ \textbackslash{}Pi)\ \&\ u\_t\ =\ Du\_\{xx\}\ -\ Vu\_x,\ \&\ -\textbackslash{}infty\ \textless{}\ x\ \textless{}\ \textbackslash{}infty,\ \textbackslash{}\textbackslash{}\ (\textbackslash{}mathrm\{H\}\ \textbackslash{}mathrm\{Y\})\ \&\ u(x,\ 0)\ =\ 1\ -\ H(x),\ \&\ -\textbackslash{}infty\ \textless{}\ x\ \textless{}\ \textbackslash{}infty,\ \textbackslash{}end\{array\}}

\begin{verbatim}
в новую задачу в движущейся системе координат, решить новую задачу, а затем полученное решение преобразовать к старым координатам (x, t). Проведем сначала замену переменных (замену независимых переменных). Вместо старых координат (x, t) вводим новые  (\xi, \tau)  по формулам
\end{verbatim}

\texttt{\textbackslash{}xi\ =\ x\ -\ Vt,}
\texttt{\textbackslash{}tau\ =\ t.}

\begin{verbatim}
Надо заметить, что буквами  \tau  и t обозначена одна и та же ведичина. Такие обозначения введены для того, чтобы избежать

путаницы в последующих формулах. Чтобы записать уравнение с частными производными в новых координатах (, т), мы воспользуемся следующими формулами перехода;
\end{verbatim}

\texttt{u\_\{t\}\ =\ u\_\{\textbackslash{}xi\}\ \textbackslash{},\ \textbackslash{}xi\_\{t\}\ +\ u\_\{\textbackslash{}tau\}\ \textbackslash{},\ \textbackslash{}tau\_\{t\}\ =\ -V\ u\_\{\textbackslash{}xi\}\ +\ u\_\{\textbackslash{}tau\},}
\texttt{u\_\{x\}\ =\ u\_\{\textbackslash{}xi\}\ \textbackslash{},\ \textbackslash{}xi\_\{x\}\ =\ u\_\{\textbackslash{}xi\},}
\texttt{u\_\{xx\}\ =\ (u\_\{\textbackslash{}xi\})\_\{x\}\ =\ u\_\{\textbackslash{}xi\textbackslash{}xi\}\ \textbackslash{}xi\_\{x\}\ =\ u\_\{\textbackslash{}xi\textbackslash{}xi\}.}

\begin{verbatim}
Функциональная диаграмма на рис. 15.5 делает эти формулы совсем очевидными. Эта диаграмма особенно полезна при вычислении частных производных функций и и и по переменным Р И

![](images/_page_115_Picture_14.jpeg)

С 15.5 Диаграмма, иллюстрирующая функциональную зависимость переменных

 \xi  и  \tau , поскольку  \xi  зависит и от x, и от t. Переменная  \tau  зависит только от t. Это все, что касается преобразований. Теперь подставим найденные выражения  u_t ,  u_x  и  u_{vx}  в уравнение (15 2) и получим
\end{verbatim}

\texttt{-Vu\_\{\textbackslash{}xi\}+u\_\{\textbackslash{}tau\}=Du\_\{\textbackslash{}xi\textbackslash{}xi\}-Vu\_\{\textbackslash{}xi\},}

\begin{verbatim}
откуда
\end{verbatim}

\texttt{u\_\{\textbackslash{}tau\}\ =\ Du\_\{\textbackslash{}xi\textbackslash{}xi\}}

.

Следовательно, новая задача с начальным условием в персменных
\texttt{\textbackslash{}xi} и \texttt{\textbackslash{}tau} имеет вид

(УЧП)

\begin{verbatim}
u_{\tau} = Du_{\xi\xi}, \quad -\infty < \xi < \infty,
\end{verbatim}

(НУ)
\texttt{u(\textbackslash{}xi,\ 0)\ =\ 1\ -\ H(\textbackslash{}xi),\ \textbackslash{}quad\ -\textbackslash{}infty\ \textless{}\ \textbackslash{}xi\ \textless{}\ \textbackslash{}infty.}

Так как \texttt{\textbackslash{}xi\ =\ x} при t = 0, то новые граничные
условия совпадают со старыми. Эта задача была решена в лекции 12 с
помощью интегрального преобразования Фурье, и ее решение записывается в
виде

\begin{verbatim}
u(\xi, \tau) = \frac{1}{2\sqrt{D\pi\tau}} \int_{-\infty}^{\infty} \varphi(\beta) e^{-(\xi-\beta)^2/4D\tau} d\beta,
\end{verbatim}

где \texttt{\textbackslash{}phi(\textbackslash{}beta)} --- начальное
условие. Следовательно, в нашем случае

\begin{verbatim}
u(\xi, \tau) = \frac{1}{2\sqrt{D\pi\tau}} \int_{-\infty}^{0} e^{-(\xi-\beta)^2/4D\tau} d\beta.
\end{verbatim}

После подстановки

\begin{verbatim}
\overline{\beta} = \frac{\xi - \beta}{2\sqrt{D\tau}}, \quad d\overline{\beta} = \frac{-1}{2\sqrt{D\tau}}d\beta
\end{verbatim}

получаем окончательный результат

(15.3)

\begin{verbatim}
u(\xi, \tau) = \frac{1}{2} \left[ \frac{2}{\sqrt{\pi}} \int_{\frac{\xi}{\sqrt{D\tau}}}^{\infty} e^{-\overline{p}^{2}} d\overline{\beta} \right] = \begin{cases} \frac{1}{2} \left[ 1 + \operatorname{erf} \left( \frac{-\xi}{2\sqrt{D\tau}} \right) \right], & \xi < 0, \\ \frac{1}{2} \operatorname{erfc} \left( \frac{\xi}{2\sqrt{D\tau}} \right), & 0 \leq \xi. \end{cases}
\end{verbatim}

Графики этого решения для различных моментов времени изображены на рис.
15.6. Осталось сделать последний шаг --- записать решение нашей задачи в
координатах x и t. Подставляя

\begin{verbatim}
\xi = x - Vt
\end{verbatim}

\begin{verbatim}
,  \tau = t 

в формулы (15.3), получаем

(15.4)
\end{verbatim}

\texttt{u(x,\ t)\ =\ \textbackslash{}begin\{cases\}\ \textbackslash{}frac\{1\}\{2\}\ \textbackslash{}left{[}\ 1\ +\ \textbackslash{}operatorname\{erf\}\textbackslash{}left(\textbackslash{}frac\{Vt\ -\ x\}\{2\textbackslash{}sqrt\{Dt\}\}\textbackslash{}right)\ \textbackslash{}right{]},\ \&\ Vt\ \textgreater{}\ x,\ \textbackslash{}\textbackslash{}\ \textbackslash{}frac\{1\}\{2\}\ \textbackslash{}operatorname\{erfc\}\textbackslash{}left(\textbackslash{}frac\{x\ -\ Vt\}\{2\textbackslash{}sqrt\{Dt\}\}\textbackslash{}right),\ \&\ Vt\ \textbackslash{}leq\ x.\ \textbackslash{}end\{cases\}}

\begin{verbatim}
Перед нами решение задачи конвективной диффузии (15.2). Его очень легко интерпретировать, если представить, что мы движемся относительно графика, изображенного на рис. 15.6.

![](images/_page_117_Figure_3.jpeg)

РИС. 15.6. Диффузня из области высокой концентрации в область низкой концентрации. Чем больше коэффициент диффузии, тем быстрее установится стационарное значение.

Другими словами, в зависимости от относительной величины D (коэффициента диффузии) и V (скорости потока) решение движется слева направо со скоростью V и в то же самое время передний фронт расплывается со скоростью, определяемой величиной D.

![](images/_page_117_Figure_6.jpeg)

РИС. 15.7. Решение задачи конвективнои диффузии. Вещество одновременно движется и диффундирует.

(На рис. 15.7 показано, как расплывается передний фронт концентрации.)

#### ЗАМЕЧАНИЯ

Преобразование координат является важным методом решения уравнений с частными производными. Выбрав подходящую систему координат, можно существенно упростить уравнение.

#### ЗАДАЧИ

1. Решите задачу Коши:
\end{verbatim}

\texttt{\textbackslash{}begin\{array\}\{ll\}\ u\_t\ =\ u\_\{xx\}\ -\ 2u\_x,\ \&\ -\textbackslash{}infty\ \textless{}\ x\ \textless{}\ \textbackslash{}infty,\ \&\ 0\ \textless{}\ t\ \textless{}\ \textbackslash{}infty,\ \textbackslash{}\textbackslash{}\ u(x,\ 0)\ =\ \textbackslash{}sin\ x,\ \&\ -\textbackslash{}infty\ \textless{}\ x\ \textless{}\ \textbackslash{}infty.\ \textbackslash{}end\{array\}}

\begin{verbatim}
2. Найдите решение задачи Коши для уравнения конвективной диффузии:
\end{verbatim}

\texttt{\textbackslash{}begin\{array\}\{ll\}\ u\_t\ =\ u\_\{xx\}\ -\ 2u\_x,\ \&\ -\textbackslash{}infty\ \textless{}\ x\ \textless{}\ \textbackslash{}infty,\ \textbackslash{}\ 0\ \textless{}\ t\ \textless{}\ \textbackslash{}infty,\ \textbackslash{}\textbackslash{}\ u(x,\ \textbackslash{}\ 0)\ =\ e\^{}x\ \textbackslash{}sin\ x,\ \&\ -\textbackslash{}infty\ \textless{}\ x\ \textless{}\ \textbackslash{}infty.\ \textbackslash{}end\{array\}}

\begin{verbatim}
Воспользуйтесь преобразованием из лекции 8.

3. Найдите решение задачи переноса:
\end{verbatim}

\texttt{\textbackslash{}begin\{array\}\{ll\}\ (\textbackslash{}text\{YHI\})\ \&\ u\_t\ =\ -2u\_x,\ \textbackslash{}\textbackslash{}\ (\textbackslash{}text\{HY\})\ \&\ u\_t(x,\ 0)\ =\ e\^{}\{-x\^{}2\}.\ \textbackslash{}end\{array\}\ \textbackslash{}\textbackslash{}\ -\textbackslash{}infty\ \textless{}\ x\ \textless{}\ \textbackslash{}infty,\ \textbackslash{}quad\ 0\ \textless{}\ t\ \textless{}\ \textbackslash{}infty,}

\begin{verbatim}
Проверьте ваш ответ.

4. Решите задачу
\end{verbatim}

\texttt{u\_t\ =\ Du\_\{xx\}\ -\ Vu\_x,\ \textbackslash{}quad\ -\textbackslash{}infty\ \textless{}\ x\ \textless{}\ \textbackslash{}infty,\ \textbackslash{}quad\ 0\ \textless{}\ t\ \textless{}\ \textbackslash{}infty,}
\texttt{u(x,\ 0)\ =\ e\^{}\{-x\^{}2\},\ \textbackslash{}quad\ -\textbackslash{}infty\ \textless{}\ x\ \textless{}\ \textbackslash{}infty.}

\begin{verbatim}
Проверьте ваше решение. Как выглядит решение в различные моменты времени?

УКАЗАНИЕ. Переход к движущейся системе отсчета позволяет сразу отбросить слагаемое  Vu_x  в уравнении конвективной диффузии. После решения новой задачи
\end{verbatim}

\texttt{\textbackslash{}begin\{array\}\{ll\}\ (\textbackslash{}text\{HV\})\ \&\ u\_\{\textbackslash{}tau\}\ =\ Du\_\{\textbackslash{}xi\textbackslash{}xi\},\ \&\ -\textbackslash{}infty\ \textless{}\ \textbackslash{}xi\ \textless{}\ \textbackslash{}infty,\ \&\ 0\ \textless{}\ \textbackslash{}tau\ \textless{}\ \textbackslash{}infty,\ \textbackslash{}\textbackslash{}\ (\textbackslash{}text\{HV\})\ \&\ u(\textbackslash{}xi,\ 0)\ =\ e\^{}\{-\textbackslash{}xi\^{}2\},\ \&\ -\textbackslash{}infty\ \textless{}\ \textbackslash{}xi\ \textless{}\ \textbackslash{}infty,\ \textbackslash{}end\{array\}}

\begin{verbatim}
сделайте замену переменных  \xi = x - Vt ,  \tau = t . В нашем случае можно в явном виде вычислить интеграл
\end{verbatim}

\texttt{u(\textbackslash{}xi,\ \textbackslash{}tau)\ =\ \textbackslash{}frac\{1\}\{2\textbackslash{}sqrt\{D\textbackslash{}pi\textbackslash{}tau\}\}\ \textbackslash{}int\_\{-\textbackslash{}infty\}\^{}\{+\textbackslash{}infty\}\ e\^{}\{-\textbackslash{}beta\^{}2\}\ e\^{}\{-(\textbackslash{}xi-\textbackslash{}beta)\^{}2/4D\textbackslash{}tau\}\ d\textbackslash{}beta.}

\begin{verbatim}
Такой интеграл уже встречался в лекции 12. Может быть, читателю будет удобнее обратиться к таблицам преобразования Фурье.

# Лекція 16. ОДНОВИМІРНЕ ХВИЛЬОВЕ РІВНЯННЯ (ГІПЕРБОЛІЧНІ РІВНЯННЯ)

ЦЕЛЬ ЛЕКЦИИ: Показать, что колебания сгруны описываются уравнением
\end{verbatim}

\texttt{u\_\{tt\}\ =\ \textbackslash{}alpha\^{}2\ u\_\{xx\}}

\begin{verbatim}
и что это уравнение — следствие законов Ньютона. Рассматриваются и некоторые другие виды волнового уравнения:
\end{verbatim}

\texttt{\textbackslash{}begin\{split\}\ u\_\{tt\}\ \&=\ \textbackslash{}alpha\^{}2\ u\_\{xx\}\ +\ F\textbackslash{}left(x,\ \textbackslash{}\ t\textbackslash{}right),\ \textbackslash{}\textbackslash{}\ u\_\{tt\}\ \&=\ \textbackslash{}alpha\^{}2\ u\_\{xx\}\ -\ \textbackslash{}beta\ u\_t\ -\ \textbackslash{}lambda\ u\ +\ F\textbackslash{}left(x,\ \textbackslash{}\ t\textbackslash{}right).\ \textbackslash{}end\{split\}}

\begin{verbatim}
До сих пор мы занкмались физическими процессами, которые описываются одномерными параболическими уравнениями (диффузионные задачи). Теперь мы приступим к изучению еще одного класса уравнений с частными производными—гиперболических уравнений. Начнем с одномерного волнового уравнения, которое описывает, в частности, поперечные колебания струны.

## Колебания струны

Рассмотрим малые колебания струны с закрепленными концами. Предположим, что струна туго натянута, сделана из однородного материала и колеблется в одной плоскости (рис. 16.1).

![](images/_page_119_Figure_10.jpeg)

РИС. 16.1. Поперечные колебания струны: a — смещение струны от положения равновесия.

Действием сил тяжести на струну будем пренебрегать. Для построения математической модели колебаний струны рассмотрим все силы, действующие на небольшой участок струны (рис. 16.2). Оказывается, что волновое уравнение—это не что иное, как закон движения Ньютона (изменение импульса равно сумме действующих сил), примененный к струне. Глядя на рис. 16.1, мы можем представить те силы, которые действуют на струну в направлений, перпендикулярном оси х.

![](images/_page_120_Figure_2.jpeg)

РИС. 16.2. Малый участок  (x, x + \Delta x)  колеблющейся струны.

1. Сила, обусловленная натяжением струны ( \alpha^2 u_{xx} ). Величина поперечной компоненты силы натяжения находится по формуле

Сила = 

Т \сін \theta_2 - Т \сін \theta_1 =

\simeq T [u_x(x + \Delta x, t) - u_x(x, t)]. 2. Внешняя сила F (x, t).

Внешняя сила F(x, t), приложенная к струне, может произвольным образом зависеть от x и t. Некоторые примеры внешних сил приводятся ниже:

а) гравитационная сила F(x, t) = -mg;

б) импульсы, распределенные вдоль струны и действующие на нее в различные моменты времени:

в) силы, возникающие при действии звуковой волны на кожу барабана (позже мы подробно изучим колебания мембраны, которые описываются волновым уравнением).

3. Силы трения, действующие на струну (— \beta u_t ). Если струна колеблется в среде, то возникает сила сопротивления, которая пропорциональна скорости  u_t .

4. Возвращающая сила (- үи).

Эта сила направлена противоположно смещению струны. Если смещение u положительно, то сила отрицательна.

Если применить теперь уравнения движения Ньютона к небольшому участку струны, то получим
\end{verbatim}

\texttt{\textbackslash{}Delta\ x\ \textbackslash{}rho\ u\_\{tt\}\ =\ T\ \textbackslash{}left{[}\ u\_x\ \textbackslash{}left(\ x\ +\ \textbackslash{}Delta\ x,\ \textbackslash{}\ t\ \textbackslash{}right)\ -\ u\_x\ \textbackslash{}left(\ x,\ \textbackslash{}\ t\ \textbackslash{}right)\ \textbackslash{}right{]}\ +\ \textbackslash{}Delta\ x\ F\ \textbackslash{}left(\ x,\ \textbackslash{}\ t\ \textbackslash{}right)\ -\ \textbackslash{}Delta\ x\ \textbackslash{}beta\ u\_t\ \textbackslash{}left(\ x,\ \textbackslash{}\ t\ \textbackslash{}right)\ -\ \textbackslash{}Delta\ x\ \textbackslash{}gamma\ u\ \textbackslash{}left(\ x,\ \textbackslash{}\ t\ \textbackslash{}right)}

\begin{verbatim}
где  \rho  — линейная плотность струны. Разделим обе части этого сротношения на  \Delta x  и устремим  \Delta x  — 0. В результате получим хорошо известное телеграфное уравнение

(16.1)
\end{verbatim}

\texttt{u\_\{tt\}\ =\ \textbackslash{}alpha\^{}\{2\}\ u\_\{xx\}\ -\ \textbackslash{}beta\ u\_\{t\}\ -\ \textbackslash{}gamma\ u\ +\ F(x,\ t).}

\begin{verbatim}
В этом уравнении следовало бы разделить  \alpha ,  \beta ,  \gamma  и F(x,t) на  \rho , но мы ради простоты сохраним для этих величин прежние обозначения. Искомое уравнение получено. Приведем интерпретацию простейшего волнового уравнения
\end{verbatim}

\texttt{(16.2)\ u\_\{tt\}\ =\ \textbackslash{}alpha\^{}2\ u\_\{xx\},}

\begin{verbatim}
основанную на интуитивных соображениях.

### Интуитивная интерпретация волнового уравнения

У читателя может возникнуть вопрос: почему уравнение вида (16.2) должно описывать нечто такое, что похоже на колебания скрипичной струны? Чтобы разобраться в этом, необходимо понять,

![](images/_page_121_Picture_9.jpeg)

РИС. 16 3. Интерпретация волнового уравнения  u_{tt} = \alpha^2 u_{xx} 

что величина  u_{tt}  является вертикальным ускорением струны в точке x. Следовательно, уравнение (16.2) можно интерпретировать следующим образом: ускорение струны, обусловленное натяжением, в каждой точке тем больше, чем больше вогнутость струны  u_{xx}  в данной точке (коэффициентом пропорциональности служит константа  \alpha^2 = T/\rho ) (см. рис. 133).

#### ЗАМЕЧАНИЯ

1. Волновое уравнение  u_{tt} = \alpha^2 u_{xx}  описывает также продольные и крутильные колебания стержня. В продольных колебаниях смещения параллельны стержню, а через  u\left(x,\,t\right)  обозначается величина продольного смещения относительно равновесного положения. Такие колебания возникают, например, при ударе молотком по торцу стержня. Аналогично описываются крутиль-

ные колебания
\end{verbatim}

u\_\{tt\} = ku\_\{xx\}

\begin{verbatim}
Здесь *k* — модуль Юнга, характеризующий упругость материала. Материалы, обладающие большим модулем Юнга, колеблются на более высоких частотах. Звуковые волны представляют собой продольные волны.

2. Если линейная плотность струны  \rho(x)  зависит от координаты, то волновое урагнение записывается в виде
\end{verbatim}

\texttt{u\_\{tt\}\ =\ \textbackslash{}frac\{\textbackslash{}partial\}\{\textbackslash{}partial\ x\}\ {[}\textbackslash{}alpha\^{}2(x)\ u\_x{]},}

\begin{verbatim}
т. е. получается УЧП с переменными коэффициентами.

3. Поскольку волновое уравнение  u_{tt} = \alpha^2 u_{xx}  содержит производную по времени второго порядка, для получения единственного решения при t>0 необходимо задать  \partial sa  начальных условия:
\end{verbatim}

u(x, 0) = f(x)

\begin{verbatim}
(начальное смещение струны),  u_t(x, 0) = g(x)  (начальная скорость струны).

Этим волновое уравнение отличается от уравнения теплопроводности, для которого необходимо задать только одно начальное условие.

4. С помощью волнового уравнения можно описать распределение электрического тока в проводе. Из законов Кирхгофа получается следующая система двух уравнений с частными производными первого порядка:

(16.3)
\end{verbatim}

i\_x + Cv\_t + Gv = 0, \textbackslash{} v\_x + Li\_t + Ri = 0,

\begin{verbatim}
где

ж — координата вдоль провода,

*t* — время,

i(x, t) — распределение тока вдоль провода,

 v\left(x,\ t\right)  — распределение потенциала вдоль провода,

С -- емкость провода на единицу длины,

G — утечка на единицу длины,

Р — сопротивление на единицу длины,

-- индуктивность провода на единицу длины.

Уравнения (16.3) принято называть системой телеграфных уравнений. В такой форме они обычно не используются. Для преобразования этой системы продифференцируем первое уравнение по x, второе продифференцируем по t, умножим обе части на C и вычтем из первого. В результате получаем
\end{verbatim}

i\_\{xx\} + Gv\_x - CLi\_\{tt\} - CRi\_t = 0.

\begin{verbatim}
Воспользуемся теперь вторым уравнением из (16.3)
\end{verbatim}

v\_\{x\} = -Li\_\{t\} - Ri,

\begin{verbatim}
окончательно получаем

(16.4)
\end{verbatim}

i\_\{xx\} = CLi\_\{tt\} + (CR + GL)i\_t + GRi.

\begin{verbatim}
Это уравнение для тока, известное под названием *телеграфного уравнения*, является гиперболическим уравнением второго порядка (если конечно ии C, ни L не равны нулю, в противном случае уравнение становится параболическим).

Напряжение описывается точно таким же уравнением:

(16.5)
\end{verbatim}

v\_\{xx\} = CLv\_\{tt\} + (CR + GL)v\_t + GRv.

\begin{verbatim}
Если G = R = 0, то уравнения (16.4)—(16.5) упрощаются:

(16.6)
\end{verbatim}

\texttt{v\_\{tt\}\ =\ \textbackslash{}alpha\^{}2\ v\_\{xx\},\ \textbackslash{}\textbackslash{}\ i\_\{tt\}\ =\ \textbackslash{}alpha\^{}2\ i\_\{xx\},\ \textbackslash{}qquad\ \textbackslash{}alpha\^{}2\ =\ 1/CL,}

\begin{verbatim}
#### ЗАДАЧИ

- 1. Получите уравнение (16.5) для *и* из системы двух уравнений первого порядка (16.3)
- 2. Зная физический смысл каждого члена волнового уравнения, что вы можете сказать о поведении во времени решения следующей задачи:

(YYII)
\end{verbatim}

u\_\{tt\} = u\_\{xx\} - u\_t

\texttt{,\ \ \textquotesingle{}0\ \textless{}\ x\ \textless{}\ 1\textquotesingle{}\ ,\ \ \textquotesingle{}0\ \textless{}\ t\ \textless{}\ \textbackslash{}infty\textquotesingle{}\ ,}
(ГУ)

\begin{verbatim}
\begin{cases} u(0, t) = 0, & 0 < t < \infty, \\ u(1, t) = 0, & 0 < t < \infty, \end{cases}
\end{verbatim}

(HУ)

\begin{verbatim}
\begin{cases} u(x, 0) = \sin(\pi x), & 0 \le x \le 1, \\ u_t(x, 0) = 0. & 0 \le x \le 1, \end{cases}
\end{verbatim}

Как выглядит в различные моменты времени решение следующей задачи

\begin{verbatim}
\begin{aligned} & (\text{UCH}\Pi) & u_{tt} = u_{xx}, & 0 < x < 1, \\ & (\text{GU}) & \begin{cases} u(0, t) = 0, & 0 < t < \infty, \\ u(1, t) = \sin t, & 0 < t < \infty, \end{cases} \\ & (\text{NU}) & \begin{cases} u(x, 0) = 0, & 0 \le x \le 1, \\ u_t(x, 0) = 0 \end{cases}
\end{verbatim}

Дайте физическую интерпретацию этой задачи.

\begin{enumerate}
\def\labelenumi{\arabic{enumi}.}
\setcounter{enumi}{3}
\tightlist
\item
  Для уравнения \texttt{u\_\{tt\}\ =\ u\_\{xx\}} найдите все решения
  вида
\end{enumerate}

\begin{verbatim}
u(x, t) = e^{ax+bt}
\end{verbatim}

Будет ли решением сумма двух решений такого вида?

\section{Лекція 17. ФОРМУЛА
Д'АЛАМБЕРА}\label{ux43bux435ux43aux446ux456ux44f-17.-ux444ux43eux440ux43cux443ux43bux430-ux434ux430ux43bux430ux43cux431ux435ux440ux430}

ЦЕЛЬ ЛЕКЦИИ: Построить решение задачи Коши для волнового уравнения

\begin{verbatim}
(\text{UBP}) \qquad u_{tt} = c^2 u_{xx}, \quad -\infty < x < \infty, \quad 0 < t < \infty,
(\text{HY}) \quad \begin{cases} u(x, 0) = f(x), \\ u_t(x, 0) = g(x), \end{cases} \quad -\infty < x < \infty.
\end{verbatim}

Эта задача описывает движение бесконечной струны с заданными начальными
условиями. Она была решена в 1750 году французским математиком
Даламбером. Формула Даламбера

\begin{verbatim}
u(x, t) = \frac{1}{2} [f(x - ct) + f(x + ct)] + \frac{1}{2c} \int_{x-ct}^{x+ct} g(\xi) d\xi
\end{verbatim}

позволяет легко найти решение, если известны начальные условия. Кроме
того, она позволяет дать интересную физическую интерпретацию решения на
языке бегущих волн.

Как, наверное, помнит читатель, в параболническом случае мы сначала
решали диффузионные задачи на конечном отрезке (методом разделения
переменных), а уже затем перешли к неограниченным промежуткам (---
\texttt{\textbackslash{}infty\ \textless{}\ x\ \textless{}\ \textbackslash{}infty}
), где применяли преобразование Фурье. В гиперболическом случае мы
поступим наоборот. Начнем с решения одномерного волнового уравнения на
всей действительной оси, т. е. с задачи Коши

(17.1)

\begin{verbatim}
u_{tt} = c^{2}u_{xx}, \qquad -\infty < x < \infty, \quad 0 < t < \infty,
(17.1) \quad (HY) \quad \begin{cases} u(x, 0) = f(x), \\ u_{t}(x, 0) = g(x), \end{cases} \quad -\infty < x < \infty,
\end{verbatim}

содержащей только начальные условия. Эту задачу можно было бы решить
либо с помощью преобразования Фурье (по переменной t), либо с помощью
преобразования Лапласа (по переменной t). Однако мы воспользуемся другим
методом (методом канонических координат), который, как мы надеемся,
заинтересует читателя. Этот метод базируется на тех же идеях, что и
метод перехода к движущейся системе отсчета, с которым мы познакомились
в лекции 15. Итак, приступим к решению задачи (17.1).

\paragraph{Решение одномерного волнового уравнения. Формула
Даламбера}\label{ux440ux435ux448ux435ux43dux438ux435-ux43eux434ux43dux43eux43cux435ux440ux43dux43eux433ux43e-ux432ux43eux43bux43dux43eux432ux43eux433ux43e-ux443ux440ux430ux432ux43dux435ux43dux438ux44f.-ux444ux43eux440ux43cux443ux43bux430-ux434ux430ux43bux430ux43cux431ux435ux440ux430}

Решение задачи (17.1) разобьем на несколько шагов:

ШАГ 1 (Замена координат (x, t) новыми каноническими координатами
\texttt{(\textbackslash{}xi,\ \textbackslash{}eta)} ).

Для решения задачи (17.1) воспользуемся тем, что если заменить две
независимые переменные x и t новыми пространственновременными
координатами

\begin{verbatim}
\xi = x + ct,
\eta = x - ct,
\end{verbatim}

то уравнение

\begin{verbatim}
u_{ti} = c^2 u_{xx}
\end{verbatim}

преобразуется в уравнение

\begin{verbatim}
u_{\xi n} = 0
\end{verbatim}

\begin{verbatim}
.

В этом легко убедиться, если воспользоваться формулами преобразования координат:

(17.2)
\end{verbatim}

\texttt{u\_\{x\}\ =\ u\_\{\textbackslash{}xi\}\ +\ u\_\{\textbackslash{}eta\},\ u\_\{t\}\ =\ c\ (u\_\{\textbackslash{}xi\}\ -\ u\_\{\textbackslash{}eta\}),\ u\_\{xx\}\ =\ u\_\{\textbackslash{}xi\textbackslash{}xi\}\ +\ 2u\_\{\textbackslash{}xi\textbackslash{}eta\}\ +\ u\_\{\textbackslash{}eta\textbackslash{}eta\},\ u\_\{tt\}\ =\ c\^{}\{2\}\ (u\_\{\textbackslash{}xi\textbackslash{}xi\}\ -\ 2u\_\{\textbackslash{}xi\textbackslash{}eta\}\ +\ u\_\{\textbackslash{}eta\textbackslash{}eta\}).}

\begin{verbatim}
Подставляя эти выражения для производных в волновое уравнение, получим
\end{verbatim}

\texttt{u\_\{\textbackslash{}xi\textbackslash{}eta\}\ =\ 0}

Первый шаг сделан.

ШАГ 2 (Решение преобразованного уравнения).

Преобразованное уравнение легко решить двумя последовательными
интегрированиями (сначала по переменной ξ, а затем по η). После
интегрирования по ξ получаем

\begin{verbatim}
u_{\eta}(\xi, \eta) = \varphi_{1}(\eta)
\end{verbatim}

--- произвольная функция от \texttt{\textbackslash{}eta} .

Интегрирование последнего соотношения по \textbf{η} приводит к общему
решению

\begin{verbatim}
u(\xi, \eta) = \varphi(\eta) + \psi(\xi),
\end{verbatim}

где \texttt{\textbackslash{}phi(\textbackslash{}eta)} --- первообразная
функция \texttt{\textbackslash{}phi\_1(\textbackslash{}eta)} , а
\texttt{\textbackslash{}psi(\textbackslash{}xi)} --- произвольная
функция. Итак, общее решение уравнения

\begin{verbatim}
u_{\xi\eta} = 0
\end{verbatim}

записывается в виде

(17.3)

\begin{verbatim}
u(\xi, \eta) = \varphi(\eta) + \psi(\xi),
\end{verbatim}

где \texttt{\textbackslash{}phi(\textbackslash{}eta)} и
\texttt{\textbackslash{}psi(\textbackslash{}xi)} --- произвольные
функции своих аргументов. Папример, читатель может проверить, что
функции

\begin{verbatim}
u(\xi, \eta) = \sin \eta + \xi^2,
\end{verbatim}

\texttt{u(\textbackslash{}xi,\ \textbackslash{}eta)\ =\ 1/\textbackslash{}eta\ +\ tg\ \textbackslash{}xi,}
\texttt{u(\textbackslash{}xi,\ \textbackslash{}eta)\ =\ \textbackslash{}eta\^{}2\ +\ e}

являются решениями уравнения
\texttt{u\_\{\textbackslash{}mathrm\{\ \}nh\}=0} . Второй шаг сделан.

\texttt{\textquotesingle{}\textbackslash{}text\{IIIA\}\textbackslash{}Gamma\textquotesingle{}\ \ 3\ (Возвращение\ к\ старым\ координатам\ x\ и\ t).}

Для нахождения общего решения, т. е. всех решений уравнения
\texttt{u\_\{tt\}\ =\ c\^{}2\ u\_\{xx\}} , подставим

\begin{verbatim}
\xi = x - ct,
n = x + ct
\end{verbatim}

в решение

\begin{verbatim}
u(\xi, \eta) = \varphi(\eta) + \psi(\xi).
\end{verbatim}

В результате получаем

(17.4)

\begin{verbatim}
u(x, t) = \varphi(x-ct) + \psi(x+ct).
\end{verbatim}

Это общее решение волнового уравнения. С физической точки эрения оно
интересно тем, что представляет сумму двух бегущих волн Р И

\pandocbounded{\includegraphics[keepaspectratio]{images/_page_126_Figure_14.jpeg}}

С. 17.1. Волна \texttt{u(x,\ t)\ =\ e\^{}\{-(x-ct)\^{}2\}} движется
слева направо

произвольной формы, движущихся в прогивоположных направлениях со
скоростью c.~Например, функции

\begin{verbatim}
u(x, t) = \sin(x-ct)
\end{verbatim}

(волна, Лвижущаяся слева направо),

\begin{verbatim}
u(x, t) = (x + ct)^2
\end{verbatim}

(волна, движущаяся справа налево),

\begin{verbatim}
u(x, t) = \sin(x - ct) + (x + ct)^2
\end{verbatim}

(две волны, движущиеся в противоположных направлениях)

являются типичными решениями волиового уравнения. На рисунке 17.1
изображена простейшая бегущая волна.

ШАГ 4 (Учет начальных условий).

Напомним читателю общий метод решения задачи Коши в теории обыкновенных
дифференциальных уравнений. Сначала находят общее решение уравнения, а
затем оно подставляется в начальные условия для того, чтобы найти
конкретные значения произвольных постоянных. У нас ситуация аналогична.
Для решения задачи Коши для волнового уравнения мы подставим общее
решение

\begin{verbatim}
u(x, t) = \varphi(x - ct) + \psi(x + ct)
\end{verbatim}

(содержащее две произвольные функции) в начальные условия

\begin{verbatim}
u(x, 0) = f(x),
\end{verbatim}

\texttt{u\_t(x,\ 0)\ =\ g(x),} чтобы найти конкретные выражения для
произвольных функций ф и ф. После подстановки получаем

(17.5)

\begin{verbatim}
\varphi(x) + \psi(x) = f(x), \\ -c\varphi'(x) + c\psi'(x) = g(x).
\end{verbatim}

Если теперь проинтегрировать второе уравнение (17.5), то получится
алгебранческая система двух уравнений относительно неизвестных функций
\texttt{\textbackslash{}varphi(x)} и \texttt{\textbackslash{}psi(x)} . В
самом деле, после интегрирования в пределах от \texttt{x\_0} до x
получаем

(17.6)

\begin{verbatim}
-c\varphi(x) + c\psi(x) = \int_{x_0}^{x} g(\xi) d\xi + K.
\end{verbatim}

Решаем теперь уравнение (17.6) совместно с первым уравнением (17.5) и
получаем следующие выражения для функций
\texttt{\textbackslash{}varphi(x)} и \texttt{\textbackslash{}psi(x)} :

(17.7)

\begin{verbatim}
\psi(x) = \frac{1}{2} f(x) + \frac{1}{2c} \int_{x_0}^{x} g(\xi) d\xi + K/2c,
\varphi(x) = \frac{1}{2} f(x) - \frac{1}{2c} \int_{x_0}^{x} g(\xi) d\xi - K/2c.
\end{verbatim}

Следовательно, общее решение задачи (17.1) дается формулой

(17.8)

\begin{verbatim}
u(x, t) = \frac{1}{2} \left[ f(x-ct) + f(x+ct) \right] + \frac{1}{2c} \int_{x-ct}^{x+ct} \mathbf{g}(\xi) d\xi.
\end{verbatim}

Цель достигнута. Мы получили общее решение вадачи (17.1). Решение вида
(17.8) принято называть формулой Даламбера. Читателю предлагается
самостоятельно проследить за расстановкой пределов интегрирования, чтобы
убедиться в том, что они меняются от x-ct до x+ct. Задача полностью
решена.

Прежде чем завершить лекцию, приведем несколько примеров, показывающих,
как применять формулу Даламбера в конкретных валачах.

\subsection{Примеры применения формулы
Даламбера}\label{ux43fux440ux438ux43cux435ux440ux44b-ux43fux440ux438ux43cux435ux43dux435ux43dux438ux44f-ux444ux43eux440ux43cux443ux43bux44b-ux434ux430ux43bux430ux43cux431ux435ux440ux430}

\subsection{1. Движение синусондальной
волны}\label{ux434ux432ux438ux436ux435ux43dux438ux435-ux441ux438ux43dux443ux441ux43eux43dux434ux430ux43bux44cux43dux43eux439-ux432ux43eux43bux43dux44b}

Рассмотрим начальные условия вида

\begin{verbatim}
u(x, 0) = \sin x,
\end{verbatim}

\texttt{u\_t(x,\ 0)\ =\ 0.} По формуле Даламбера получаем решение

\begin{verbatim}
u(x, t) = \frac{1}{2} [\sin(x-ct) + \sin(x+ct)].
\end{verbatim}

Его можно интерпретировать следующим образом. Начальное смещение струны
\texttt{u(x,\ 0)\ =\ \textbackslash{}sin\ x} делится на две одинаковые
части

\begin{verbatim}
\frac{\sin x}{2}
\end{verbatim}

\texttt{H\ \ \textquotesingle{}\textbackslash{}frac\{\textbackslash{}sin\ x\}\{2\}\textquotesingle{}\ .}

Каждая из частей распространяется со скоростью с в виде бегущих волн.
Одна из волн перемещается слева направо, а вторая---в противоположном
направлении. Читателю предлагается подумать над тем, как выглядит в этом
случае результирующая волна.

\section{2. Распространение простейшего прямоугольного
импульса}\label{ux440ux430ux441ux43fux440ux43eux441ux442ux440ux430ux43dux435ux43dux438ux435-ux43fux440ux43eux441ux442ux435ux439ux448ux435ux433ux43e-ux43fux440ux44fux43cux43eux443ux433ux43eux43bux44cux43dux43eux433ux43e-ux438ux43cux43fux443ux43bux44cux441ux430}

В этом случае начальные условия задаются следующим образом:

\begin{verbatim}
u(x, 0) = \begin{cases} 1, & -1 < x < 1, \\ 0, & \end{cases}
\end{verbatim}

\texttt{u\_t(x,\ 0)\ =\ 0} (в остальных точках).

Поскольку начальное смещение расшепляется на две полуволны, движущиеся в
противоположных направлениях, результирующее движение будет таким, как
показано на рис. 17.2.

\pandocbounded{\includegraphics[keepaspectratio]{images/_page_129_Figure_2.jpeg}}

РИС, 17.2. Начальное возмущение распадается на две бегущие волны.

\paragraph{3. Задана начальная
скорость}\label{ux437ux430ux434ux430ux43dux430-ux43dux430ux447ux430ux43bux44cux43dux430ux44f-ux441ux43aux43eux440ux43eux441ux442ux44c}

Предположим теперь, что в начальный момент струна находится в положении
равновесия. Придадим струне начальную скорость (как в фортепиано) вида
\texttt{\textbackslash{}sin\ x} , т. е.

\begin{verbatim}
u(x, 0) = 0,
u_t(x, 0) = \sin x.
\end{verbatim}

Решение находим по формуле Даламбера

\begin{verbatim}
u(x, t) = \frac{1}{2c} \int_{x-ct}^{x+ct} \sin \xi \, d\xi = \frac{1}{2c} \left[ \cos (x+ct) - \cos (x-ct) \right].
\end{verbatim}

Оно представляет сумму двух бегущих косинусоидальных воли. Читатель
должен спросить себя, кажется ли ему это решение разумным.

На этом лекция 17 заканчивается. В следующей лекции мы покажем, как
формула Даламбера помогает дать интерпретацию в плоскости переменных x и
t.

\paragraph{ЗАМЕЧАНИЯ}\label{ux437ux430ux43cux435ux447ux430ux43dux438ux44f}

\begin{itemize}
\tightlist
\item
  \begin{enumerate}
  \def\labelenumi{\arabic{enumi}.}
  \tightlist
  \item
    Общее решение уравнения с частными производными второго порядка
    содержит две произвольные функции, а общее решение обыкновенного
    дифференциального уравнения содержит две произвольные постоянные.
    Значит, у уравнения с частными производными больше решений, чем у
    обыкновенного дифференциального уравнения.
  \end{enumerate}
\item
  2 Замена переменных для приведения уравнения с частными производными к
  более простому виду является одним из общих методов теории. Новые
  координаты (ξ, η) принято называть каноническими координатами. Мы
  будем обращаться к ним и в дальнейшем, особенно при изучении
  гиперболических уравнений.
\item
  3 Метод нахождения общего решения уравнения с частными производными и
  подстановки этого решения в начальные условия не является общим
  методом решения уравнений с частными производными. Решение, которое
  рассматривается в данной лекции,--- единственный пример использования
  этого метода. Обычно невозможно найти общее решение уравнения с
  частными производными, но даже если удается это сделать, оказывается
  очень сложным подставить это решение в краевые условия.
\end{itemize}

\paragraph{ЗАДАЧН}\label{ux437ux430ux434ux430ux447ux43d}

\begin{enumerate}
\def\labelenumi{\arabic{enumi}.}
\item
  Проверьте, что решение Даламбера (17-8) удовлетворяет начальным
  условиям задачи (17.1).
\item
  Подставьте (17.7) в общее решение
  \texttt{u(x,\ t)\ =\ \textbackslash{}psi(x\ -\ ct)\ +\ \textbackslash{}psi(x\ +\ ct)}
  и получите формулу Даламбера.
\item
  Найдите решение задачи Коши:
\end{enumerate}

\begin{verbatim}
(\text{UBP}) \qquad u_{tt} = u_{xx}, \quad -\infty < x < \infty, \quad 0 < t < \infty, \\ (\text{HV}) \qquad \left\{ \begin{array}{ll} u(x, \ 0) = e^{-x^2} \\ u_t(x, \ 0) = 0. \end{array} \right. \quad -\infty < x < \infty, \\ \end{array}
\end{verbatim}

Как выглядит это решение в различные моменты времени? 4. Найдите решение
задачи Коши:

(УЧП)

\begin{verbatim}
u_{tt} = u_{xx}, -\infty < x < \infty, 0 < t < \infty,
\end{verbatim}

(НУ)

\begin{verbatim}
\begin{cases} u(x, 0) = 0, \\ u_t(x, 0) = xe^{-\lambda^2}, -\infty < x < \infty, \end{cases}
\end{verbatim}

Постройте график решения в различные моменты времени. 5. Решая совместно
уравнения (17.6) и первое уравнение из (17.5),

получите формулы (17.7).

\section{Лекція 18. ФОРМУЛА Д'АЛАМБЕРА
(ПРОДОВЖЕННЯ)}\label{ux43bux435ux43aux446ux456ux44f-18.-ux444ux43eux440ux43cux443ux43bux430-ux434ux430ux43bux430ux43cux431ux435ux440ux430-ux43fux440ux43eux434ux43eux432ux436ux435ux43dux43dux44f}

МЕТА ЛЕКЦІЇ: показати, як за допомогою формули Д'Аламбера можна знайти
розв'язок хвильового рівняння на напівобмеженій прямій.

(УЧП)

\begin{verbatim}
J[u] = \iint_D F(x, y, u, u_x, u_y)\,dx\,dy,
\end{verbatim}

(ГУ)
\texttt{u(0,\ t)\ =\ 0,\textbackslash{}quad\ 0\ \textless{}\ t\ \textless{}\ \textbackslash{}infty}
(НУ)

\begin{verbatim}
F_u - \frac{\partial}{\partial x}F_{ux} - \frac{\partial}{\partial y}F_{uy} = 0.
\end{verbatim}

Додатково подається інтерпретація формули Д'Аламбера в координатній
площині \texttt{x,\ t}.

На попередній лекції було показано, що якщо задані початкове зміщення
струни \texttt{u(x,0)=f(x)} і початкова швидкість
\texttt{u\_t(x,0)=g(x)}, то

(18.1)

\begin{verbatim}
u(x,t)=\frac{1}{2}\big[f(x-ct)+f(x+ct)\big]+\frac{1}{2c}\int_{x-ct}^{x+ct} g(\xi)\,d\xi
\end{verbatim}

описує зміщення струни в довільний момент часу.

У цій лекції ми розглянемо інтерпретацію цієї формули в площині змінних
\texttt{x} і \texttt{t} (просторово-часова площина) та покажемо, як її
модифікувати для задачі коливань напівнескінченної струни.

\section{Просторово-часова інтерпретація формули
Д'Аламбера}\label{ux43fux440ux43eux441ux442ux43eux440ux43eux432ux43e-ux447ux430ux441ux43eux432ux430-ux456ux43dux442ux435ux440ux43fux440ux435ux442ux430ux446ux456ux44f-ux444ux43eux440ux43cux443ux43bux438-ux434ux430ux43bux430ux43cux431ux435ux440ux430}

Розглянемо задачу Коші:

(18.2)

\begin{verbatim}
u_{tt} = c^2 u_{xx}, \quad -\infty < x < \infty,\quad t > 0
\end{verbatim}

Розв'язок зображається на площині \texttt{(x,t)}, а не в координатах
\texttt{(x,u)}.

\pandocbounded{\includegraphics[keepaspectratio]{images/_page_133_Figure_5.jpeg}}

Рис. 18.2. Розв'язок задачі Коші на площині \texttt{x,t}:\\
a --- передній фронт хвилі, що рухається вліво;\\
b --- задній фронт хвилі;\\
c --- передній фронт хвилі, що рухається вправо.

Тепер розглянемо випадок, коли початкове зміщення дорівнює нулю, а
початкова швидкість довільна.

\section{Випадок 2 (початкове зміщення дорівнює
нулю)}\label{ux432ux438ux43fux430ux434ux43eux43a-2-ux43fux43eux447ux430ux442ux43aux43eux432ux435-ux437ux43cux456ux449ux435ux43dux43dux44f-ux434ux43eux440ux456ux432ux43dux44eux454-ux43dux443ux43bux44e}

Розглянемо початкові умови:

\begin{verbatim}
u(x,0)=0,\quad u_t(x,0)=g(x).
\end{verbatim}

Розв'язок має вигляд формули Д'Аламбера.

Значення розв'язку в точці \texttt{(x\_e,t\_e)} можна інтерпретувати як
інтеграл початкової швидкості по відрізку \texttt{{[}x-ct,\ x+ct{]}}.

\pandocbounded{\includegraphics[keepaspectratio]{images/_page_134_Figure_2.jpeg}}

Рис. 18.3. Інтерпретація задачі з початковою швидкістю.

Перші наближення в методі Рітца мають вигляд:

\begin{verbatim}
u_1(x,y)=a_1xy(1-x)(1-y),
u_3(x,y)=xy(1-x)(1-y)(a_1+a_2x+a_3y),
\end{verbatim}

і так далі.

Відповідний функціонал:

\begin{verbatim}
J[u_n] = \int_0^1 \int_0^1 \left[\left(\sum_{j=1}^n a_j \frac{\partial \varphi_j}{\partial x}\right)^2 + \left(\sum_{j=1}^n a_j \frac{\partial \varphi_j}{\partial y}\right)^2 + 2f \sum_{j=1}^n a_j \varphi_j\right] dx\,dy.
\end{verbatim}

Умови мінімуму:

\begin{verbatim}
\frac{\partial J[u_n]}{\partial a_k}=0,\quad k=1,\dots,n.
\end{verbatim}

Після цього задача зводиться до системи лінійних алгебраїчних рівнянь
для коефіцієнтів \texttt{a\_1,\textbackslash{}dots,a\_n}.

Розв'язавши систему, отримуємо наближений розв'язок задачі.

\section{ЗАУВАЖЕННЯ}\label{ux437ux430ux443ux432ux430ux436ux435ux43dux43dux44f-5}

\begin{enumerate}
\def\labelenumi{\arabic{enumi}.}
\tightlist
\item
  У цій лекції методом Рітца мінімізувався функціонал, що залежить від
  функції двох змінних. Аналогічно можна мінімізувати функціонали інших
  типів, наприклад:
\end{enumerate}

\begin{verbatim}
y(0)=0,\quad y(1)=1.
\end{verbatim}

\begin{enumerate}
\def\labelenumi{\arabic{enumi}.}
\setcounter{enumi}{1}
\tightlist
\item
  Для хвильового рівняння
\end{enumerate}

\begin{verbatim}
u_{tt}=c^2u_{xx}, \quad 0<x<\infty,\quad t>0,
\end{verbatim}

з граничними та початковими умовами:

\begin{verbatim}
u(0,t)=0,\quad u(x,0)=f(x),\quad u_t(x,0)=g(x),
\end{verbatim}

розв'язок будується за допомогою формули Д'Аламбера. \# Розв'язання
задачі для напівнескінченної струни за допомогою формули Д'Аламбера

Наше завдання --- знайти хвильові рухи напівнескінченної струни, лівий
кінець якої закріплений, за заданих початкових умов. Для розв'язання
задачі (18.6) ми використаємо той самий метод, що й при знаходженні
загального розв'язку хвильового рівняння на всій прямій.

Замість загального розв'язку

\begin{verbatim}
u_{xx} + u_{yy} = 1
\end{verbatim}

підставляємо в початкові умови та отримуємо (як у попередній лекції):

(18.7)

\begin{verbatim}
J[y] = \int_{0}^{1} (y^{2} + {y'}^{2})\,dx,\quad y(0) = 0,\; y(1) = 1,
z(x) = (1 - x)\, y(x).
\end{verbatim}

Тепер перед нами виникає задача, якої не було для нескінченної струни.
Розв'язок \texttt{u(x,t)} повинен бути визначений лише в першому
квадранті \texttt{(x\textgreater{}0,\textbackslash{}\ t\textgreater{}0)}
площини змінних \texttt{x} і \texttt{t}. Отже, потрібно вміти
обчислювати функції

\begin{verbatim}
\varphi(x-ct), \quad \psi(x+ct)
\end{verbatim}

для всіх
\texttt{-\textbackslash{}infty\ \textless{}\ x-ct\ \textless{}\ \textbackslash{}infty}
та \texttt{0\ \textless{}\ x+ct\ \textless{}\ \textbackslash{}infty}.

На жаль, перша з формул (18.7) дозволяє обчислювати
\texttt{\textbackslash{}varphi(x-ct)} лише при
\texttt{x-ct\ \textbackslash{}ge\ 0}, оскільки початкові функції
\texttt{f(x)} і \texttt{g(x)} задані тільки для додатних значень
аргументу.

Якщо \texttt{x\ -\ ct\ \textbackslash{}ge\ 0}, то після підстановки
(18.7) в загальний розв'язок отримуємо:

\begin{verbatim}
J[u] = \int_{0}^{1} \int_{0}^{1} (u_{x}^{2} + u_{y}^{2})\,dx\,dy,
u_{xx} + u_{yy} = \sin(\pi x).
\end{verbatim}

Для напівнескінченної струни розв'язок зручно записати у вигляді:

\begin{verbatim}
u(x,t) = \varphi(x-ct) + \psi(x+ct)
= \frac{1}{2}\bigl[f(x-ct) + f(x+ct)\bigr] + \frac{1}{2c} \int_{x-ct}^{x+ct} g(\xi)\,d\xi.
\end{verbatim}

Якщо \texttt{x\ \textbackslash{}ge\ ct}, ця формула збігається зі
звичайною формулою Д'Аламбера для нескінченної струни. Якщо ж
\texttt{x\ \textless{}\ ct}, необхідно враховувати граничну умову
\texttt{u(0,t)=0}, тобто відбиття хвилі від закріпленого кінця.

У результаті для напівнескінченної струни маємо: u(x,t)=
\texttt{\textbackslash{}begin\{cases\}}
\texttt{\textbackslash{}dfrac12\textbackslash{}bigl{[}f(x-ct)+f(x+ct)\textbackslash{}bigr{]}+\textbackslash{}dfrac\{1\}\{2c\}\textbackslash{}displaystyle\textbackslash{}int\_\{x-ct\}\^{}\{x+ct\}\ g(\textbackslash{}xi)\textbackslash{},d\textbackslash{}xi,\ \&\ x\ \textbackslash{}ge\ ct,\ \textbackslash{}\textbackslash{}{[}1.2ex{]}}
\texttt{\textbackslash{}dfrac12\textbackslash{}bigl{[}f(x+ct)-f(ct-x)\textbackslash{}bigr{]}+\textbackslash{}dfrac\{1\}\{2c\}\textbackslash{}displaystyle\textbackslash{}int\_\{ct-x\}\^{}\{x+ct\}\ g(\textbackslash{}xi)\textbackslash{},d\textbackslash{}xi,\ \&\ x\ \textless{}\ ct.}
\texttt{\textbackslash{}end\{cases\}}

Цим завершується побудова розв'язку. Його геометрична інтерпретація
подана на рис. 18.6.

\pandocbounded{\includegraphics[keepaspectratio]{images/_page_137_Figure_9.jpeg}}

Рис. 18.6. Інтерпретація задачі Коші для півобмеженої струни в
\texttt{xt}-площині: \texttt{a} --- збурення відбивається від межі;
\texttt{\textbackslash{}delta} --- збурення поширюється вздовж
характеристики з точки \texttt{(x\_0+ct\_0,\textbackslash{},0)}. \#
ЗАУВАЖЕННЯ 1. Формула для півобмеженої струни повністю узгоджується з
фізичною картиною: при \texttt{x\ \textbackslash{}ge\ ct} хвиля ще не
встигла відбитися від межі, а при \texttt{x\textless{}ct} у розв'язку
з'являється відбита хвиля зі зміненим знаком.

\begin{enumerate}
\def\labelenumi{\arabic{enumi}.}
\setcounter{enumi}{1}
\item
  Якщо замінити граничну умову \texttt{u(0,t)=0} на іншу, то явний
  вигляд розв'язку теж зміниться. Але сама ідея побудови - доозначення
  хвильових складових за допомогою граничної умови - залишається тією
  самою.
\item
  Прямі
\end{enumerate}

\begin{verbatim}
x-ct=\text{const}, \qquad x+ct=\text{const},
\end{verbatim}

називаються характеристиками хвильового рівняння. Саме вздовж цих прямих
поширюється збурення.

\section{ЗАВДАННЯ}\label{ux437ux430ux432ux434ux430ux43dux43dux44f-5}

\begin{enumerate}
\def\labelenumi{\arabic{enumi}.}
\item
  Побудуйте графіки розв'язку для напівнескінченної струни в різні
  моменти часу та поясніть, як хвиля відбивається від закріпленої межі.
\item
  Опишіть алгоритм побудови такого розв'язку:

  \begin{itemize}
  \tightlist
  \item
    продовжити початкові дані на всю вісь,
  \item
    побудувати ліву і праву хвилі,
  \item
    врахувати відбиття від межі \texttt{x=0}.
  \end{itemize}
\item
  Розв'яжіть змішану задачу для напівнескінченної струни, коли початкові
  дані задано графічно.
\end{enumerate}

\pandocbounded{\includegraphics[keepaspectratio]{images/_page_139_Figure_3.jpeg}}

\section{\texorpdfstring{Розв'язання нелінійного рівняння
\texttt{\textbackslash{}Delta\ u\ +\ u\^{}2\ =\ 0} методом
збурень}{Розв'язання нелінійного рівняння \textbackslash Delta u + u\^{}2 = 0 методом збурень}}\label{ux440ux43eux437ux432ux44fux437ux430ux43dux43dux44f-ux43dux435ux43bux456ux43dux456ux439ux43dux43eux433ux43e-ux440ux456ux432ux43dux44fux43dux43dux44f-delta-u-u2-0-ux43cux435ux442ux43eux434ux43eux43c-ux437ux431ux443ux440ux435ux43dux44c}

Розглянемо нелінійну задачу Діріхле як збурення простішої лінійної
задачі. Для цього вводять сімейство задач

\begin{verbatim}
\Delta u + \varepsilon u^2 = 0, \qquad 0 \le \varepsilon \le 1,
\end{verbatim}

і шукають розв'язок у вигляді ряду

\begin{verbatim}
u = u_0 + \varepsilon u_1 + \varepsilon^2 u_2 + \cdots .
\end{verbatim}

Нульове наближення \texttt{u\_0} обирають як розв'язок відповідної
лінійної задачі. Далі, підставляючи ряд у рівняння і прирівнюючи
коефіцієнти при однакових степенях \texttt{\textbackslash{}varepsilon},
послідовно знаходять поправки
\texttt{u\_1,\ u\_2,\ \textbackslash{}ldots}.

(46.4)

Використайте метод, аналогічний розглянутому в лекції. Дайте фізичну
інтерпретацію цієї задачі.

\textbf{4.} Припустимо, що коливання струни описуються рівнянням

\begin{verbatim}
u_{tt} = u_{xx},
\end{verbatim}

початкове зміщення зображене на рисунку, а початкова швидкість дорівнює
нулю. Зобразіть розв'язок цієї задачі в \texttt{x}--\texttt{t} площині.
Зверніть увагу, що початкові умови в цій задачі є розривними.

\pandocbounded{\includegraphics[keepaspectratio]{images/_page_139_Figure_8.jpeg}}

\section{Лекція 19. ХВИЛЬОВЕ РІВНЯННЯ І ГРАНИЧНІ
УМОВИ}\label{ux43bux435ux43aux446ux456ux44f-19.-ux445ux432ux438ux43bux44cux43eux432ux435-ux440ux456ux432ux43dux44fux43dux43dux44f-ux456-ux433ux440ux430ux43dux438ux447ux43dux456-ux443ux43cux43eux432ux438}

\textbf{Мета лекції.} Ознайомити читача з основними типами граничних
умов для хвильового рівняння на скінченному відрізку.

Розглянемо одномірне хвильове рівняння:

\begin{verbatim}
u_{tt} = c^2 u_{xx}, \qquad 0 < x < L,\quad 0 < t < \infty.
\end{verbatim}

Найчастіше виникають три типи граничних умов:

\begin{enumerate}
\def\labelenumi{\arabic{enumi}.}
\tightlist
\item
  \textbf{Умови першого роду}:
\end{enumerate}

\begin{verbatim}
u(0,t)=g_1(t), \qquad u(L,t)=g_2(t).
\end{verbatim}

\begin{enumerate}
\def\labelenumi{\arabic{enumi}.}
\setcounter{enumi}{1}
\tightlist
\item
  \textbf{Умови другого роду}:
\end{enumerate}

\begin{verbatim}
u_x(0,t)=g_1(t), \qquad u_x(L,t)=g_2(t).
\end{verbatim}

\begin{enumerate}
\def\labelenumi{\arabic{enumi}.}
\setcounter{enumi}{2}
\tightlist
\item
  \textbf{Умови третього роду}:
\end{enumerate}

\begin{verbatim}
\alpha_1 u_x(0,t)+\beta_1 u(0,t)=g_1(t), \qquad
\alpha_2 u_x(L,t)+\beta_2 u(L,t)=g_2(t).
\end{verbatim}

У фізичному сенсі це відповідає заданим переміщенням, силам або пружному
закріпленню кінців струни.

Далі розглянемо задачу:

\begin{verbatim}
\begin{aligned}
u_{tt} &= c^2u_{xx}, && 0 < x < L,\quad 0 < t < \infty, \\
u(0,t) &= g_1(t), \quad u(L,t) = g_2(t), && 0 < t < \infty, \\
u(x,0) &= f(x), \quad u_t(x,0) = g(x), && 0 \leqslant x \leqslant L.
\end{aligned}
\end{verbatim}

Щоб розв'язати таку неоднорідну задачу, зручно подати розв'язок як суму
спеціально підібраної функції, що задовольняє граничні умови, і нової
невідомої функції з однорідними граничними умовами. Саме цей підхід
використовується далі.

\texttt{Правильні\ частини\ виду\ \textquotesingle{}r\^{}n\textquotesingle{},\ \textquotesingle{}r\^{}n\ \textbackslash{}cos(n\textbackslash{}theta)\textquotesingle{},\ \textquotesingle{}r\^{}n\ \textbackslash{}sin(n\textbackslash{}theta)\textquotesingle{}\ ведуть\ до\ розв’язків}
\texttt{\textquotesingle{}Ar\^{}\{n+2\}\textquotesingle{},\ \textquotesingle{}Br\^{}\{n+2\}\ \textbackslash{}cos(n\textbackslash{}theta)\textquotesingle{},\ \textquotesingle{}Cr\^{}\{n+2\}\ \textbackslash{}sin(n\textbackslash{}theta)\textquotesingle{}\ відповідно.}

Підставляючи функції \texttt{u\_p(r,\textbackslash{}theta)} у
неоднорідне рівняння, отримуємо:
\texttt{A\ =\ -\textbackslash{}frac\{1\}\{32\},\ \textbackslash{}qquad\ B\ =\ -\textbackslash{}frac\{1\}\{24\}.}

Отже,

\begin{verbatim}
u_x(1,t) = \frac{1}{k} v(t) - (k - \text{модуль Юнга}).
\end{verbatim}

Щоб розв'язати задачу, потрібно підставити загальний розв'язок

\begin{verbatim}
u(r,\theta)=u_h(r,\theta)+u_p(r,\theta)
\end{verbatim}

у граничні умови \texttt{u(1,\textbackslash{}theta)=0} і знайти
коефіцієнти \texttt{a\_n} та \texttt{b\_n}.

Після цього отримуємо:

\begin{verbatim}
u_x(0,t) = \frac{h}{T} u(0,t).
\end{verbatim}

Отже, \texttt{a\_0=1/32}, \texttt{a\_2=1/24}, а всі інші \texttt{a\_n} і
\texttt{b\_n} дорівнюють нулю. Таким чином, розв'язок задачі знайдено:
\texttt{u\_x(L,t)\ =\ -\textbackslash{}frac\{h\}\{T\}\ u(L,t).}

Функція \texttt{u\_1(r,\textbackslash{}theta)} називається збуренням
першого порядку. Додавши \texttt{u\_0} і \texttt{u\_1}, отримуємо новий
підхід до розв'язання задачі (46.2):

(46.5)

\begin{verbatim}
u_x(0,t) - \frac{h}{T}u(0,t) = 0,
\end{verbatim}

Щоб знайти збурення наступного порядку
\texttt{u\_z(r,\textbackslash{}theta)}, потрібно розв'язати задачу:
\texttt{u\_x(L,t)\ +\ \textbackslash{}frac\{h\}\{T\}u(L,t)\ =\ 0.}

разом із системою:

\begin{verbatim}
\begin{cases}
\Delta u_{1} = -2u_{0}u_{1}, \\
u_{2}(1,\theta) = 0,
\end{cases}
\end{verbatim}

де функції \texttt{u\_\{\textbackslash{}phi\}(r,\textbackslash{}theta)}
та \texttt{u\_\{1\}(r,0)} вже відомі.

\texttt{Звісно,\ виконувати\ такі\ алгебраїчні\ перетворення\ без\ комп’ютера\ дуже\ складно.\ На\ щастя,\ вираз\ (46.5)\ є\ майже\ точним\ розв’язком\ задачі.\ Насправді,\ підставивши\ його\ в\ рівняння\ \textquotesingle{}\textbackslash{}Delta\ u\ +\ u\^{}\{2\}\textquotesingle{},\ отримаємо\ вираз,\ який\ майже\ всюди\ дорівнює\ нулю\ в\ області\ \textquotesingle{}0\ \textless{}\ r\ \textless{}\ 1\textquotesingle{}.}

Окрім нелінійних задач, теорія збурень може використовуватися для задач
у областях зі складною межею (якщо нерегулярність незначна). Розглянемо
приклад.

\section{Приклад задачі з порушеною
межею}\label{ux43fux440ux438ux43aux43bux430ux434-ux437ux430ux434ux430ux447ux456-ux437-ux43fux43eux440ux443ux448ux435ux43dux43eux44e-ux43cux435ux436ux435ux44e}

Розглянемо задачу Лапласа в деформованій області. Нехай на межі задано:

(46.6)

\begin{verbatim}
u\left(1 + \frac{1}{4} \sin \theta,\ \theta\right) = \cos \theta,\quad 0 \le \theta \le 2\pi.
\end{verbatim}

Узагальнено:

\begin{verbatim}
u(1+\epsilon\sin\theta,\ \theta) = \cos\theta,\quad 0 \le \epsilon \le \frac{1}{4}.
\end{verbatim}

Таким чином, ми завершуємо розгляд основних граничних умов для
гіперболічних задач. Наступні лекції будуть присвячені методам
розв'язання задач з такими умовами.\# Лекція 47

\section{РОЗВ'ЯЗАННЯ ДИФЕРЕНЦІАЛЬНИХ РІВНЯНЬ У ПОХІДНИХ ПОХІДНИХ ЗА
ДОПОМОГОЮ МЕТОДУ КОНФОРМНИХ
ВІДОБРАЖЕНЬ}\label{ux440ux43eux437ux432ux44fux437ux430ux43dux43dux44f-ux434ux438ux444ux435ux440ux435ux43dux446ux456ux430ux43bux44cux43dux438ux445-ux440ux456ux432ux43dux44fux43dux44c-ux443-ux43fux43eux445ux456ux434ux43dux438ux445-ux43fux43eux445ux456ux434ux43dux438ux445-ux437ux430-ux434ux43eux43fux43eux43cux43eux433ux43eux44e-ux43cux435ux442ux43eux434ux443-ux43aux43eux43dux444ux43eux440ux43cux43dux438ux445-ux432ux456ux434ux43eux431ux440ux430ux436ux435ux43dux44c}

МЕТА ЛЕКЦІЇ: Показати, як деякі двовимірні граничні задачі можна
перетворити на інші, простіші задачі за допомогою конформного
відображення. Припустимо, наприклад, що ми хочемо розв'язати рівняння
Лапласа

\begin{verbatim}
\begin{cases} u(0, t) = 0, \\ u_x(1, t) = -u(1, t), \\ u(x, 0) = x, \\ u_t(x, 0) = 0. \end{cases}
\end{verbatim}

З граничною умовою в області комплексної форми (площина змінних
\texttt{x} і \texttt{y}) цю задачу можна перетворити на нову, в якій
потрібно знайти розв'язок рівняння Лапласа

\begin{verbatim}
\text{у простішій області (площина змінних } \xi \text{ і } \eta \text{)}.
\end{verbatim}

\pandocbounded{\includegraphics[keepaspectratio]{images/_page_147_Picture_5.jpeg}}

\subsection{\texorpdfstring{20.1. Три типові стоячі хвилі:
\texttt{X(x)T(t)}}{20.1. Три типові стоячі хвилі: X(x)T(t)}}\label{ux442ux440ux438-ux442ux438ux43fux43eux432ux456-ux441ux442ux43eux44fux447ux456-ux445ux432ux438ux43bux456-xxtt}

Декілька прикладів стоячих хвиль показано на рис. 20.1. Якщо профілі цих
стоячих хвиль відомі \texttt{X\_n(x)} і природа коливань кожної хвилі
\texttt{T\_n(t)}, то розв'язок задачі коливань гітарної струни є
суперпозицією елементарних коливань \texttt{X\_n(x)T\_n(t)}.

Залишається лише обрати коефіцієнти \texttt{c\_n} так, щоб при
\texttt{t\ =\ 0} початкові умови
\texttt{u(x,0)=f(x),\ \textbackslash{}qquad\ u\_t(x,0)=g(x)} були
виконані.

Задачу коливань гітарної струни розв'яжемо методом розділення змінних.

\subsection{Розв'язання задачі методом розділення
змінних}\label{ux440ux43eux437ux432ux44fux437ux430ux43dux43dux44f-ux437ux430ux434ux430ux447ux456-ux43cux435ux442ux43eux434ux43eux43c-ux440ux43eux437ux434ux456ux43bux435ux43dux43dux44f-ux437ux43cux456ux43dux43dux438ux445}

Щоб розв'язати змішану задачу:

\begin{verbatim}
\begin{cases}
u(0,t)=0,\\
u(L,t)=0,
\end{cases}
\end{verbatim}

розглянемо розв'язок у вигляді добутку:

\begin{verbatim}
u(x,t)=X(x)T(t).
\end{verbatim}

Підстановка дає:

\begin{verbatim}
\sin(\beta L)=0,
\end{verbatim}

звідки

\begin{verbatim}
\beta_n = \frac{n\pi}{L}, \qquad n=1,2,3,\ldots
\end{verbatim}

(випадок \texttt{n=0} дає тривіальний розв'язок).

Отже, маємо послідовність елементарних коливань струни:

\texttt{u\_n(x,t)=\textbackslash{}sin\textbackslash{}!\textbackslash{}left(\textbackslash{}frac\{n\textbackslash{}pi\ x\}\{L\}\textbackslash{}right)}
\texttt{\textbackslash{}left{[}a\_n\ \textbackslash{}sin\textbackslash{}!\textbackslash{}left(\textbackslash{}frac\{n\textbackslash{}pi\ \textbackslash{}alpha\ t\}\{L\}\textbackslash{}right)}
\texttt{+\ b\_n\ \textbackslash{}cos\textbackslash{}!\textbackslash{}left(\textbackslash{}frac\{n\textbackslash{}pi\ \textbackslash{}alpha\ t\}\{L\}\textbackslash{}right)\textbackslash{}right{]}.}
або

\begin{verbatim}
u_n(x,t)=R_n \sin\!\left(\frac{n\pi x}{L}\right)
\cos\!\left(\frac{n\pi \alpha (t-\delta_n)}{L}\right).

Тут a_n, b_n, R_n, \delta_n — довільні сталі.

## Комплексні потенціали та конформні відображення

Розглянемо функцію:
\end{verbatim}

\texttt{\textbackslash{}ln\textbackslash{}left{[}\textbackslash{}frac\{z-1\}\{z+1\}\textbackslash{}right{]}}
\texttt{=\ \textbackslash{}ln\textbackslash{}left\textbar{}\textbackslash{}frac\{z-1\}\{z+1\}\textbackslash{}right\textbar{}}
\texttt{+\ i\ \textbackslash{}arg\textbackslash{}left{[}\textbackslash{}frac\{z-1\}\{z+1\}\textbackslash{}right{]}.}

\begin{verbatim}
Тоді:
\end{verbatim}

\texttt{\textbackslash{}varphi(x,y)=\textbackslash{}frac\{1\}\{\textbackslash{}pi\}\textbackslash{}arctan\textbackslash{}left(\textbackslash{}frac\{2y\}\{x\^{}2+y\^{}2-1\}\textbackslash{}right).}

\begin{verbatim}
## Характеристичне рівняння

Для параболічного випадку характеристика має вигляд:
\end{verbatim}

\texttt{\textbackslash{}frac\{dy\}\{dx\}\ =\ -\textbackslash{}frac\{B\}\{2A\}.}

\begin{verbatim}
Для прикладу рівняння:
\end{verbatim}

u\_\{xx\} + 2u\_\{xy\} + u\_\{yy\} = 0

\begin{verbatim}
характеристичне рівняння:
\end{verbatim}

\texttt{\textbackslash{}frac\{dy\}\{dx\}\ =\ -1.}

\begin{verbatim}
Тому вводимо нові змінні:
\end{verbatim}

\texttt{\textbackslash{}xi\ =\ y\ -\ x,\ \textbackslash{}qquad\ \textbackslash{}eta\ =\ y.}

\begin{verbatim}
У цих координатах рівняння набуває вигляду:
\end{verbatim}

\texttt{u\_\{\textbackslash{}eta\textbackslash{}eta\}\ =\ 0.}

\begin{verbatim}
Загальний розв’язок:
\end{verbatim}

\texttt{u(\textbackslash{}xi,\textbackslash{}eta)\ =\ \textbackslash{}eta\ f(\textbackslash{}xi)\ +\ g(\textbackslash{}xi),}

\begin{verbatim}
де f і g — довільні функції.

## ЗАВДАННЯ

1. Розв’яжіть задачу Коші методом характеристик для лінійного рівняння першого порядку.

2. Зведіть рівняння
\end{verbatim}

u\_\{xx\} + 2u\_\{xy\} + u\_\{yy\} + u = 2

\begin{verbatim}
до канонічної форми.

3. Зведіть рівняння
\end{verbatim}

u\_\{xx\} + 2u\_\{yy\} + x\^{}2 u\_x = e\textsuperscript{\{-x}2/2\}

\begin{verbatim}
до канонічної форми.

4. Для хвильового рівняння запишіть граничні умови першого, другого і третього роду та поясніть їх фізичний зміст.

5. Знайдіть розв’язок задачі переносу методом переходу до рухомої координати.

Щоб завершити перетворення, потрібно так само перенести початкові умови на нову невідому функцію. Після цього задача для хвильового рівняння з неоднорідними граничними умовами зводиться до задачі з однорідними граничними умовами, але вже з модифікованими правими частинами.

На практиці це означає, що складність переноситься з межі в саме рівняння, зате далі можна застосовувати стандартні методи: розділення змінних, розкладання за власними функціями або інтегральні перетворення.

![](images/_page_144_Picture_4.jpeg)

Рис. 19.6. Виникнення неоднорідних граничних умов пружного закріплення: h - коефіцієнт жорсткості пружини.

Для умов пружного закріплення однорідні граничні умови часто записують у вигляді
\end{verbatim}

\texttt{u\_x(0,t)\ -\ h\_1u(0,t)\ =\ 0,\ \textbackslash{}qquad\ u\_x(L,t)\ +\ h\_2u(L,t)\ =\ 0.}

\begin{verbatim}
Саме з такого запису зазвичай починають побудову власних функцій і відповідних спектральних розкладів для подальшого розв'язання задачі.

![](images/_page_150_Picture_5.jpeg)

Фіг. 20.3. Стоячи u_n(x, t) = X_n(x)T_n(t) .

У нашій задачі рівняння і граничні умови є лінійними та однорідними, тому довільна сума елементарних стоячих хвиль також є розв'язком хвильового рівняння і задовольняє граничним умовам. Тепер ми повинні знайти таку суперпозицію стоячих хвиль, щоб вона також задовольняла початкові умови, і тоді отримаємо розв'язок нашої проблеми. Заміна кількості
\end{verbatim}

\texttt{u(x,t)=\textbackslash{}sum\_\{n=1\}\^{}\{\textbackslash{}infty\}\textbackslash{}left(a\_n\textbackslash{}sin\textbackslash{}frac\{n\textbackslash{}pi\textbackslash{}alpha\ t\}\{L\}+b\_n\textbackslash{}cos\textbackslash{}frac\{n\textbackslash{}pi\textbackslash{}alpha\ t\}\{L\}\textbackslash{}right)\textbackslash{}sin\textbackslash{}frac\{n\textbackslash{}pi\ x\}\{L\}.}

\begin{verbatim}
![](images/_page_154_Figure_5.jpeg)

Головна різниця між поперечними коливаннями тонкої балки та поперечними коливаннями струни полягає в тому, що промінь чинить опір згину. Не заглиблюючись у механіку тонких балок, припустимо, що, враховуючи опір згинання (замість рівняння хвилі) веде до рівняння четвертого порядку

(21.1)
\end{verbatim}

\texttt{u\_\{tt\}\ =\ -a\^{}2\ u\_\{xxxx\},\ \textbackslash{}qquad\ 0\ \textless{}\ x\ \textless{}\ L,\textbackslash{};\ 0\ \textless{}\ t\ \textless{}\ \textbackslash{}infty.}

\begin{verbatim}
# Лекція 20. КОЛИВАННЯ ОБМЕЖЕНОЇ СТРУНИ (СТОЯЧІ ХВИЛІ)

ЦЕЛЬ ЛЕКЦИИ: Показать, что решение смешанной задачи для поперечных колебаний струны

~~
\begin{array}{ll} (\text{UCH}\Pi) & u_{tt} = \alpha^2 u_{xx}, & 0 < x < L, & 0 < t < \infty, \\ (\text{GU}) & \begin{cases} u \ (0, t) = 0, & 0 < t < \infty, \\ u \ (L, t) = 0, & 0 < t < \infty, \end{cases} \\ (\text{HY}) & \begin{cases} u \ (x, 0) = f \ (x), & 0 \leqslant x \leqslant L, \\ u_t \ (x, 0) = g \ (x),  & 0 \leqslant x \leqslant L, \end{cases}
~~

можно найти стандартным методом разделения переменных. При этом оказывается, что решение  u\left(x,\,t\right) , представленное в виде ряда \end{array}
\end{verbatim}

\texttt{u(x,\ t)\ =\ \textbackslash{}sum\_\{n=1\}\^{}\{\textbackslash{}infty\}\ X\_n(x)\ T\_n(t),}

\begin{verbatim}
является суперпозицией простейших колебаний. Пространственная часть  X_n(x)  каждого элементарного колебания является решением (собственной функцией) некоторой краевой задачи Штурма—Лиувилля.

Итак, мы научились решать волновое уравнение  u_{tt}=c^2u_{xx}  на всей прямой  -\infty < x < \infty . Формула Даламбера представляет решение в виде суммы двух бегущих волн, распространяющихся в противоположные стороны. Если же рассматривать это же уравнение на ограниченном промежутке 0 < x < L, то бегущих волн уже не станет, так как они будут взаимодействовать с границами. Вместо них возникают другие волны, которые называют стоячими волнами. Посмотрим, например, что случится, если гитарную струну, закрепленную в точках x=0, x=L, привести в движение. Для ответа на этот вопрос нам необходимо решить следующую смешанную задачу гиперболического типа:
\end{verbatim}

\texttt{\textbackslash{}begin\{array\}\{ll\}\ (\textbackslash{}mathrm\{Y\}\textbackslash{}mathrm\{H\}\textbackslash{}Pi)\ \&\ u\_\{tt\}\ =\ \textbackslash{}alpha\^{}2\ u\_\{xx\},\ \&\ 0\ \textless{}\ x\ \textless{}\ L,\ \&\ 0\ \textless{}\ t\ \textless{}\ \textbackslash{}infty,\ \textbackslash{}\textbackslash{}\ .\ (\textbackslash{}mathrm\{\textbackslash{}Gamma\}\textbackslash{}mathrm\{Y\})\ \&\ \textbackslash{}begin\{cases\}\ u\ \textbackslash{}\ (0,\ \textbackslash{}\ t)\ =\ 0,\ \&\ 0\ \textless{}\ t\ \textless{}\ \textbackslash{}infty\ \textbackslash{}\textbackslash{}\ u\ \textbackslash{}\ (L,\ \textbackslash{}\ t)\ =\ 0,\ \&\ 0\ \textless{}\ t\ \textless{}\ \textbackslash{}infty\ \textbackslash{}end\{cases\}\ \textbackslash{}\textbackslash{}\ (\textbackslash{}mathrm\{H\}\textbackslash{}mathrm\{Y\})\ \&\ \textbackslash{}begin\{cases\}\ u\ \textbackslash{}\ (x,\ \textbackslash{}\ 0)\ =\ f\ \textbackslash{}\ (x),\ \&\ 0\ \textbackslash{}leqslant\ x\ \textbackslash{}leqslant\ L,\ \textbackslash{}\textbackslash{}\ u\_t\ \textbackslash{}\ (x,\ \textbackslash{}\ 0)\ =\ g\ \textbackslash{}\ (x).\ \&\ 0\ \textbackslash{}leqslant\ x\ \textbackslash{}leqslant\ L,\ \textbackslash{}end\{cases\}\ \textbackslash{}end\{array\}}

\begin{verbatim}
Оказывается, в этом случае бегущие волны отражаются от границ таким образом, что результирующие колебания становятся не бегущими, а сохраняющими форму в одном положении, т. е.

![](images/_page_147_Picture_5.jpeg)

РИС. 20.1. Три типичных стоячих волны X(x) T(t).

превращаются в стоячие волны Несколько примеров стоячих волн показано на рис. 20.1. Если известны профили этих стоячих волн  X_n(x)  и характер колебаний каждой волны  T_n(t) , то решение задачи о колебаниях гитарной струны представляет собой суперпозицию простейших колебаний  X_n(x)\,T_n(t)  и запи-

сывается в виде
\end{verbatim}

\texttt{u(x,\ t)\ =\ \textbackslash{}sum\_\{n=1\}\^{}\{\textbackslash{}infty\}\ c\_n\ X\_n(x)\ T\_n(t),}

\begin{verbatim}
Единственное, что остается сделать нам,— это подобрать коэффициенты  c_n  так, чтобы при t=0 удовлетворялись начальные условия  u\left(x,\ 0\right)=f\left(x\right)  и  u_t\left(x,\ 0\right)=g\left(x\right) .

Задачу о колебаниях гитарной струны мы решим методом разделения переменных.

#### Решение задачи о колебаниях ограниченной струны методом разделения переменных

Чтобы решить смешанную задачу

(УЧП)
\end{verbatim}

\texttt{u\_\{tt\}\ =\ \textbackslash{}alpha\^{}2\ u\_\{xx\},\ \textbackslash{}quad\ 0\ \textless{}\ x\ \textless{}\ L,\ \textbackslash{}quad\ 0\ \textless{}\ t\ \textless{}\ \textbackslash{}infty,}
\texttt{(\textbackslash{}Gamma\ Y)\ \textbackslash{}begin\{cases\}\ u(0,\ t)\ =\ 0,\ \textbackslash{}\textbackslash{}\ u(L,\ t)\ =\ 0,\ \textbackslash{}end\{cases\}\ \textbackslash{}quad\ 0\ \textless{}\ t\ \textless{}\ \textbackslash{}infty,}
\texttt{(HY)\ \textbackslash{}begin\{cases\}\ u(x,\ 0)\ =\ f(x),\ \textbackslash{}\textbackslash{}\ u\_t(x,\ 0)\ =\ g(x),\ \textbackslash{}end\{cases\}\ \textbackslash{}quad\ 0\ \textbackslash{}le\ x\ \textbackslash{}le\ L,}

\begin{verbatim}
найдем сначала стоячне волны, т. е. решения вида
\end{verbatim}

u(x, t) = X(x) T(t).

\begin{verbatim}
Подставив это выражение в уравнение (20.1) и разделив переменные, получим два обыкновенных дифференциальных уравнения
\end{verbatim}

\texttt{T\textquotesingle{}\textquotesingle{}\ -\ \textbackslash{}alpha\^{}2\ \textbackslash{}lambda\ T\ =\ 0,}
\texttt{X\textquotesingle{}\textquotesingle{}\ -\ \textbackslash{}lambda\ X\ =\ 0,}

\begin{verbatim}
где константа  \lambda  — пока что любое действительное число —  \infty < \lambda < \infty  .

Результаты исследования решений этих двух уравнений при различных значениях  \lambda  приведены на рис 20.2.

Теперь мы должны *опбросить* те решения, которые либо неограничены при  t \to \infty , либо становятся тождественно равными нулю, если потребовать выполнения граничных условий u(0, t) = u(L, t) = 0.

В качестве упражнения читателю предлагается проверить, что нетривиальные и ограниченные решения можно получить только при отрицательных значениях параметра  \lambda . Таким образом, наша задача— найти такие значения констант A, B, C, D и отрица-

тельной константы  \lambda = -\beta^2 , чтобы выражение

(20.2)
\end{verbatim}

\texttt{u(x,\ t)\ =\ {[}C\ \textbackslash{}sin(\textbackslash{}beta\ x)\ +\ D\ \textbackslash{}cos(\textbackslash{}beta\ x){]}{[}A\ \textbackslash{}sin(\textbackslash{}alpha\ \textbackslash{}beta\ t)\ +\ B\ \textbackslash{}cos(\textbackslash{}alpha\ \textbackslash{}beta\ t){]}}

\begin{verbatim}
удовлетворяло граничным условиям. Если это удастся, то мы будем располагать множеством простейших колебаний струны, каждое из которых удовлетворяет уравнению с частными произ- Р И

![](images/_page_149_Figure_5.jpeg)

С 20.2. Построение стоячих воли при раздичных значениях параметра λ.

водными и граничным условиям. Затем останется только найти такую линейную комбинацию этих колебаний, чтобы сумма удовлетворяла начальным условиям при t=0.

Подстановка (20.2) в граничные условия u(0, t) = u(L, t) = 0 дает  u(0, t) = X(0) T(t) = D[A \sin(\alpha\beta t) + B \cos(\alpha\beta t)] = 0 \Rightarrow D = 0 , (20.3)  u(L, t) = X(L) T(t) = C \sin(\beta L) [A \sin(\alpha\beta t) + B \cos(\alpha\beta t)] = 0 \Rightarrow \sin(\beta L) = 0 .

Другими словами, константа разделения  \beta  (мы можем вместо  \lambda  искать константу  \beta ) должна удовлетворять уравнению  \sin{(\beta L)} = 0 , откуда
\end{verbatim}

\texttt{\textbackslash{}beta\_n\ =\ \textbackslash{}frac\{n\textbackslash{}pi\}\{L\},\ \textbackslash{}quad\ n\ =\ 0,\ 1,\ 2,\ \textbackslash{}ldots}

\begin{verbatim}
Заметим, что если во втором уравнении (20.3) положить C=0, то получится тривиальное решение  X\left(x\right)T\left(t\right)\equiv0 . Таким образом, мы нашли последовательность элементарных колебаний струны (будем помечать их индексом n)
\end{verbatim}

\texttt{u\_n(x,\ t)\ =\ X\_n(x)\ T\_n(t)\ =\ \textbackslash{}sin(n\textbackslash{}pi\ x/L)\ {[}a\_n\ \textbackslash{}sin(n\textbackslash{}pi\ at/L)\ +\ b\_n\ \textbackslash{}cos(n\textbackslash{}pi\ at/L){]},}

\begin{verbatim}
(20.4)

или
\end{verbatim}

\texttt{u\_n(x,\ t)\ =\ R\_n\ \textbackslash{}sin(n\textbackslash{}pi\ x/L)\ \textbackslash{}cos(n\textbackslash{}pi\textbackslash{}alpha\ (t\ -\ \textbackslash{}delta\_n)/L{]},}

\begin{verbatim}
где  a_n ,  b_n ,  R_n  и  \delta_n —произвольные постоянные. Заметим, что элементарные колебания удовлетворяют граничным условиям. Читателю должно быть ясно, что каждое элементарное колебание представляет собой стоячую волну. В стоячей волне все точки колеблются с одной и той же частотой, а форма волны подобна волнам, изображенным на рис. 20.3.

![](images/_page_150_Picture_5.jpeg)

РИС. 20.3. Стоячие волны  u_n(x, t) = X_n(x) T_n(t) .

В нашей задаче уравнение и граничные условия линейны и однородны, поэтому произвольная сумма элементарных стоячих волн также является решением волнового уравнения и удовлетворяет граничным условиям. Мы должны теперь найти такую суперпозицию стоячих волн, чтобы она удовлетворяла также и начальным условиям, тогда получим решение нашей задачи. Подстановка суммы
\end{verbatim}

\texttt{u(x,\ t)\ =\ \textbackslash{}sum\_\{n=1\}\^{}\{\textbackslash{}infty\}\ \textbackslash{}sin(n\textbackslash{}pi\ x/L)\ \textbackslash{}left{[}\ a\_n\ \textbackslash{}sin(n\textbackslash{}pi\ \textbackslash{}alpha\ t/L)\ +\ b\_n\ \textbackslash{}cos(n\textbackslash{}pi\ \textbackslash{}alpha\ t/L)\ \textbackslash{}right{]}}

\begin{verbatim}
в начальные условия
\end{verbatim}

\texttt{u(x,\ \textbackslash{}theta)\ =\ f(x),}

\begin{verbatim}
u_t(x, \theta) = g(x) приводит к двум уравнениям
\end{verbatim}

\texttt{\textbackslash{}sum\_\{n=1\}\^{}\{\textbackslash{}infty\}\ b\_n\ \textbackslash{}sin(n\textbackslash{}pi\ x/L)\ =\ f(x),}
\texttt{\textbackslash{}sum\_\{n=1\}\^{}\{\textbackslash{}infty\}\ a\_n\ (n\textbackslash{}pi\ \textbackslash{}alpha/L)\ \textbackslash{}sin(n\textbackslash{}pi\ x/L)\ =\ g(x).}

\begin{verbatim}
Воспользовавшись соотношениями ортогональности
\end{verbatim}

\texttt{\textbackslash{}int\_\{0\}\^{}\{L\}\ \textbackslash{}sin(m\textbackslash{}pi\ x/L)\ \textbackslash{}sin(n\textbackslash{}pi\ x/L)\ dx\ =\ \textbackslash{}begin\{cases\}\ 0,\ \&m\ \textbackslash{}neq\ n,\ \textbackslash{}\textbackslash{}\ L/2,\ \&\ m\ =\ n,\ \textbackslash{}end\{cases\}}

\begin{verbatim}
находим коэффициенты  a_n  и  b_n 

(20.5)
\end{verbatim}

\texttt{a\_\{n\}\ =\ \textbackslash{}frac\{2\}\{n\textbackslash{}pi\textbackslash{}alpha\}\ \textbackslash{}int\_\{0\}\^{}\{L\}\ g(x)\ \textbackslash{}sin(n\textbackslash{}pi\ x/L)\ dx,\ \textbackslash{}\textbackslash{}\ b\_\{n\}\ =\ \textbackslash{}frac\{2\}\{L\}\ \textbackslash{}int\_\{0\}\^{}\{L\}\ f(x)\ \textbackslash{}sin(n\textbackslash{}pi\ x/L)\ dx.}

\begin{verbatim}
Итак, задача решена. Ее решение записывается в виде

(20.6)
\end{verbatim}

\texttt{u(x,\ t)\ =\ \textbackslash{}sum\_\{n=1\}\^{}\{\textbackslash{}infty\}\ \textbackslash{}sin(n\textbackslash{}pi\ x/L)\ \textbackslash{}left{[}\ a\_n\ \textbackslash{}sin(n\textbackslash{}pi\ \textbackslash{}alpha\ t/L)\ +\ b\_n\ \textbackslash{}cos(n\textbackslash{}pi\ \textbackslash{}alpha\ t/L)\ \textbackslash{}right{]},}

\begin{verbatim}
а коэффициенты  a_n  и  b_n  находятся по формулам (20.5). Прежде чем закончить лекцию, сделаем несколько замечаний.

## ЗАМЕЧАНИЯ

1. Если начальная *скорость* струны равна нулю, то решение (20.6) принимает вид
\end{verbatim}

\texttt{u(x,\ t)\ =\ \textbackslash{}sum\_\{n=1\}\^{}\{\textbackslash{}infty\}\ b\_n\ \textbackslash{}sin(n\textbackslash{}pi\ x/L)\ \textbackslash{}cos(n\textbackslash{}pi\textbackslash{}alpha\ t/L)}

\begin{verbatim}
и его можно интерпретировать следующим образом. Предположим, что мы разложили начальное смещение струны
\end{verbatim}

\texttt{u(x,\ \textbackslash{}theta)\ =\ f(x)}

\begin{verbatim}
в ряд по синусам
\end{verbatim}

\texttt{f(x)\ =\ \textbackslash{}sum\_\{n=1\}\^{}\{\textbackslash{}infty\}\ b\_n\ \textbackslash{}sin(n\textbackslash{}pi\ x/L).}

\begin{verbatim}
Тогда каждый член такого разложения должен совершать колебания
\end{verbatim}

\texttt{u\_n(x,\ t)\ =\ b\_n\ \textbackslash{}sin(n\textbackslash{}pi\ x/L)\ \textbackslash{}cos(n\textbackslash{}pi\textbackslash{}alpha\ t/L).}

\begin{verbatim}
Если теперь просуммировать все эти колебания, то получится решение нашей задачи. Например, предположим, что

начальное смещение струны задается функцией  f(x) = \sin(\pi x/L) + 0.5 \sin(3\pi x/L) + 0.25 \sin(5\pi x/L) .

Полный отклик системы на начальное условие такого вида равен сумме откликов на каждое слагаемое, т. е.
\end{verbatim}

\texttt{u(x,\ t)\ =\ \textbackslash{}sin(\textbackslash{}pi\ x/L)\textbackslash{}cos(\textbackslash{}pi\ \textbackslash{}alpha\ t/L)\ +\ 0.5\textbackslash{}sin(3\textbackslash{}pi\ x/L)\textbackslash{}cos(3\textbackslash{}pi\ \textbackslash{}alpha\ t/L)\ +\ 0.25\textbackslash{}sin(5\textbackslash{}pi\ x/L)\textbackslash{}cos(5\textbackslash{}pi\ \textbackslash{}alpha\ t/L).}

\begin{verbatim}
2. Слагаемое с номером n в разложении (20.6) называется **п-й модой колебаний** или **п-й гармоникой**. Используя тригонометрические формулы, можно записать эти гармоники в виде
\end{verbatim}

\texttt{R\_n\ \textbackslash{}sin(n\textbackslash{}pi\ x/L)\ \textbackslash{}cos{[}n\textbackslash{}pi\textbackslash{}alpha(t-\textbackslash{}delta\_n)/L{]},}

\begin{verbatim}
где  R_n  и  \delta_n  — новые произвольные постоянные (амплитуда и фазовый угол). Эта новая форма представления n-й моды оказывается более полезной при анализе колебаний. Частота колебаний од (рад/с) и-й моды находится по формуле
\end{verbatim}

\texttt{\textbackslash{}omega\_n\ =\ \textbackslash{}frac\{n\textbackslash{}pi\textbackslash{}alpha\}\{L\}\ =\ \textbackslash{}frac\{n\textbackslash{}pi\}\{L\}\ \textbackslash{}sqrt\{\textbackslash{}frac\{T\}\{\textbackslash{}rho\}\},}

\begin{verbatim}
где Т — натяжение струны, р — линейная плотность струны. Отметим, также, что частота n-й гармоники в n раз больше основной частоты (n=1). Это свойство частоты характерно не для всех типов колебаний. Приятный звук скривичной или гитарной струны обусловлен кратностью частогы обертонов основной частоте. В барабане частоты гармоник высокого порядка не кратны основной частоте колебаний.

## ЗАДАЧИ

1. Найдите решение задачи (20.1), если начальные условия задаются формулами
\end{verbatim}

\texttt{u(x,\ 0)\ =\ \textbackslash{}sin(\textbackslash{}pi\ x/L)\ +\ 0,5\ \textbackslash{}sin(3\textbackslash{}pi\ x/L),}

\begin{verbatim}
u_t(x, 0) = 0. Постройте графики решения для различных моментов времени. Является ли это решение периодическим по времени? Каков его период?

2. Найдите решение задачи о колебаниях струны (20.1) при начальных условиях вида
\end{verbatim}

u(x, 0) = 0,

\begin{verbatim}
u_t(x, 0) = \sin(3\pi x/L). Как выглядит график решения в различные моменты времени?

- 3. Покажите, что если на рис. 20.2 величина  \lambda \geqslant 0 , то решения X(x) T(t) будут либо не ограничены, либо тождественно равны нулю.
- Найдите решение задачи о колебаниях струны при начальных условиях вида
\end{verbatim}

\texttt{u(x,\ 0)\ =\ \textbackslash{}sin(3\textbackslash{}pi\ x/L),}

\begin{verbatim}
u_t(x, 0) = (3\pi \alpha/L)\sin(3\pi x/L). 5. Гитарная струна длины L=1 оттягивается за среднюю точку на величину h (см. рис.) Начальное положение струны можно задать в виде
\end{verbatim}

\texttt{u(x,\ 0)\ =\ \textbackslash{}begin\{cases\}\ 2hx,\ \&\ 0\ \textbackslash{}le\ x\ \textbackslash{}le\ 0.5,\ \textbackslash{}\textbackslash{}\ 2h(1-x),\ \&\ 0.5\ \textbackslash{}le\ x\ \textbackslash{}le\ 1.\ \textbackslash{}end\{cases\}}

\begin{verbatim}
![](images/_page_153_Picture_8.jpeg)

Какие движения будет совершать струна, если ее отпустить? 6. Решите задачу о затухающих колебаниях струны
\end{verbatim}

\texttt{\textbackslash{}begin\{array\}\{lll\}\ (\textbackslash{}mathrm\{Y\}\ \textbackslash{}Pi\ \textbackslash{}Pi)\ \&\ u\_\{tt\}\ =\ \textbackslash{}alpha\^{}2\ u\_\{xx\}\ -\ \textbackslash{}beta\ u\_t,\ \&\ 0\ \textless{}\ x\ \textless{}\ 1,\ \&\ 0\ \textless{}t\ \textless{}\ \textbackslash{}infty,\ \textbackslash{}\textbackslash{}\ (\textbackslash{}Gamma\ \textbackslash{}mathrm\{Y\})\ \&\ \textbackslash{}begin\{cases\}\ u\ \textbackslash{}\ (0,\ t)\ =\ 0,\ \textbackslash{}\textbackslash{}\ u\ \textbackslash{}\ (1,\ t)\ =\ 0,\ \textbackslash{}end\{cases\}\ \&\ 0\ \textless{}\ t\ \textless{}\ \textbackslash{}infty,\ \textbackslash{}\textbackslash{}\ (\textbackslash{}mathrm\{H\}\ \textbackslash{}mathrm\{Y\})\ \&\ \textbackslash{}begin\{cases\}\ u\ \textbackslash{}\ (x,\ 0)\ =\ f\ \textbackslash{}\ (x),\ \textbackslash{}\ u\_t\ \textbackslash{}\ (x,\ 0)\ =\ 0.\ \textbackslash{}end\{cases\}\ \&\ 0\ \textbackslash{}leqslant\ x\ \textbackslash{}leqslant\ 1,\ \textbackslash{}end\{array\}}

\begin{verbatim}
Представляется ли вам разумным полученное решение? Удовлетворяет ли оно уравнению с частными производными, граничным и начальным условиям?

7. Как вы будете решать неоднородное уравнение с заданными начальными и граничными условиям?
\end{verbatim}

\texttt{\textbackslash{}begin\{array\}\{ll\}\ (\textbackslash{}mathrm\{Y\}\ \textbackslash{}mathrm\{H\}\ \textbackslash{}Pi)\ \&\ u\_\{tt\}\ =\ \textbackslash{}alpha\^{}2\ u\_\{xx\}\ +\ Kx,\ \textbackslash{}quad\ 0\ \textless{}\ x\ \textless{}\ 1,\ \textbackslash{}\ 0\ \textless{}\ t\ \textless{}\ \textbackslash{}infty,\ \textbackslash{}\textbackslash{}\ (\textbackslash{}mathrm\{\textbackslash{}Gamma\}\ \textbackslash{}mathrm\{Y\})\ \&\ \textbackslash{}begin\{cases\}\ u\ \textbackslash{}\ (0,\ \textbackslash{}\ t)\ =\ 0,\ \textbackslash{}\textbackslash{}\ u\ \textbackslash{}\ (1,\ \textbackslash{}\ t)\ =\ 0,\ \textbackslash{}end\{cases\}\ \&\ 0\ \textless{}\ t\ \textless{}\ \textbackslash{}infty,\ \textbackslash{}\textbackslash{}\ (\textbackslash{}mathrm\{H\}\ \textbackslash{}mathrm\{Y\})\ \&\ \textbackslash{}begin\{cases\}\ u\ \textbackslash{}\ (x,\ \textbackslash{}\ 0)\ =\ f\ \textbackslash{}\ (x),\ \textbackslash{}\textbackslash{}\ u\_t\ \textbackslash{}\ (x,\ \textbackslash{}\ 0)\ =\ 0.\ \textbackslash{}end\{cases\}\ \&\ 0\ \textbackslash{}leqslant\ x\ \textbackslash{}leqslant\ 1,\ \textbackslash{}end\{array\}}

\begin{verbatim}
# Лекція 21. ШАРНІРНО ОПЕРТА БАЛКА

Розглянемо малі коливання тонкої балки, кінці якої вільно спираються на дві опори. Це означає, що кінці балки не зміщуються, але можуть вільно повертатися.

Тому на кінцях балки мають виконуватися граничні умови
\end{verbatim}

\texttt{u(0,\ t)\ =\ 0,\ \textbackslash{}qquad\ u(1,\ t)\ =\ 0,}

\begin{verbatim}
а також умови нульового згинального моменту
\end{verbatim}

\texttt{u\_\{xx\}(0,\ t)\ =\ 0,\ \textbackslash{}qquad\ u\_\{xx\}(1,\ t)\ =\ 0.}

\begin{verbatim}
Із теорії тонких балок випливає, що згинальний момент пропорційний u_{xx}, тому для вільно опертих кінців саме ці умови є природними.

![](images/_page_155_Figure_8.jpeg)

Рис. 21.1. Шарнірно оперта балка.

Отже, коливання балки описуються такою змішаною задачею (для простоти відповідний коефіцієнт вважаємо рівним одиниці):

(21.2)
\end{verbatim}

\texttt{\textbackslash{}begin\{aligned\}}
\texttt{(\textbackslash{}text\{УЧП\})\textbackslash{}quad\ \&\ u\_\{tt\}\ =\ -u\_\{xxxx\},\ \&\&\ 0\ \textless{}\ x\ \textless{}\ 1,\textbackslash{}\ 0\ \textless{}\ t\ \textless{}\ \textbackslash{}infty,\ \textbackslash{}\textbackslash{}}
\texttt{(\textbackslash{}text\{ГУ\})\textbackslash{}quad\ \&\ u(0,t)\ =\ 0,\textbackslash{}quad\ u\_\{xx\}(0,t)\ =\ 0,\ \textbackslash{}\textbackslash{}}
\texttt{\&\ u(1,t)\ =\ 0,\textbackslash{}quad\ u\_\{xx\}(1,t)\ =\ 0,\ \&\&\ 0\ \textless{}\ t\ \textless{}\ \textbackslash{}infty,\ \textbackslash{}\textbackslash{}}
\texttt{(\textbackslash{}text\{НУ\})\textbackslash{}quad\ \&\ u(x,0)\ =\ f(x),\ \textbackslash{}quad\ u\_t(x,0)\ =\ g(x),\ \&\&\ 0\ \textbackslash{}leqslant\ x\ \textbackslash{}leqslant\ 1.}
\texttt{\textbackslash{}end\{aligned\}}

\begin{verbatim}
Щоб розв’язати цю задачу, використаємо метод розділення змінних. Будемо шукати періодичні розв’язки (гармонічні флуктуації).

(21.3)
\end{verbatim}

u(x,t) = X(x),T(t).

\begin{verbatim}
Підставляючи у рівняння, отримуємо:
\end{verbatim}

\texttt{\textbackslash{}frac\{T\textquotesingle{}\textquotesingle{}(t)\}\{T(t)\}\ =\ -\textbackslash{}frac\{X\textquotesingle{}\textquotesingle{}\textquotesingle{}\textquotesingle{}(x)\}\{X(x)\}\ =\ -\textbackslash{}omega\^{}2.}

\begin{verbatim}
Звідси маємо дві задачі:
\end{verbatim}

\texttt{\textbackslash{}begin\{cases\}}
\texttt{T\textquotesingle{}\textquotesingle{}\ +\ \textbackslash{}omega\^{}2\ T\ =\ 0,\ \textbackslash{}\textbackslash{}}
\texttt{X\textquotesingle{}\textquotesingle{}\textquotesingle{}\textquotesingle{}\ -\ \textbackslash{}omega\^{}2\ X\ =\ 0.}
\texttt{\textbackslash{}end\{cases\}}

\begin{verbatim}
Загальний розв’язок для просторової частини:
\end{verbatim}

\texttt{X(x)\ =\ C\ +\ Dx\ +\ E\ e\^{}\{\textbackslash{}sqrt\{\textbackslash{}omega\}x\}\ +\ F\ e\^{}\{-\textbackslash{}sqrt\{\textbackslash{}omega\}x\}.}

\begin{verbatim}
Підставляємо граничні умови:

1. u(0,t)=0 \Rightarrow X(0)=0:
C + E + F = 0

2. u_{xx}(0,t)=0 \Rightarrow X''(0)=0:
E\omega + F\omega = 0 \Rightarrow E + F = 0

3. u(1,t)=0 \Rightarrow X(1)=0:
C + D + E e^{\sqrt{\omega}} + F e^{-\sqrt{\omega}} = 0

4. u_{xx}(1,t)=0 \Rightarrow X''(1)=0:
E e^{\sqrt{\omega}} + F e^{-\sqrt{\omega}} = 0

Із перших двох умов:
\end{verbatim}

\texttt{E\ =\ -F,\ \textbackslash{}qquad\ C\ =\ 0.}

\begin{verbatim}
Підставляючи у решту умов, отримуємо систему для визначення констант та власних значень \omega.

Нетривіальний розв’язок існує тоді, коли:
\end{verbatim}

\texttt{\textbackslash{}sin(\textbackslash{}sqrt\{\textbackslash{}omega\})\ =\ 0,}

\begin{verbatim}
звідки:
\end{verbatim}

\texttt{\textbackslash{}omega\_n\ =\ n\^{}2\ \textbackslash{}pi\^{}2,\ \textbackslash{}qquad\ n\ =\ 1,2,3,\textbackslash{}dots}

\begin{verbatim}
## Розклад розв’язку

Загальний розв’язок можна подати у вигляді суперпозиції:
\end{verbatim}

\texttt{u(x,t)\ =\ \textbackslash{}sum\_\{n=1\}\^{}\{\textbackslash{}infty\}\ \textbackslash{}left{[}}
\texttt{A\_n\ \textbackslash{}cos(\textbackslash{}omega\_n\ t)\ +\ B\_n\ \textbackslash{}sin(\textbackslash{}omega\_n\ t)}
\texttt{\textbackslash{}right{]}\ X\_n(x),}

\begin{verbatim}
де X_n(x) — власні функції задачі.

Коефіцієнти A_n, B_n визначаються з початкових умов:
\end{verbatim}

\texttt{u(x,0)\ =\ f(x),\ \textbackslash{}qquad\ u\_t(x,0)\ =\ g(x).}

\begin{verbatim}
# Лекція 22. ПЕРЕХІД ДО БЕЗРОЗМІРНИХ ЗМІННИХ

**Мета лекції:** показати, як крайові, початкові та фізичні задачі можна записати у безрозмірному вигляді.

Перехід до безрозмірних змінних дозволяє привести рівняння до універсального вигляду без явної залежності від фізичних параметрів. Це спрощує аналіз і підкреслює математичну структуру задачі.

Після такого переходу багато різних задач (у фізиці, хімії, біології) зводяться до однієї і тієї ж канонічної форми.

![](images/_page_161_Figure_4.jpeg)

Рис. 22.1. Різні задачі, що приводяться до безрозмірної форми.

## Приведення дифузійної задачі до безрозмірного вигляду

Почнемо зі змішаної задачі для відрізка з початковою умовою на температуру  
`u(x,0)=\sin\left(\frac{\pi x}{L}\right)`  
та граничними температурами на кінцях T_1 і T_2. Іншими словами, ми розв'язуємо задачу

(22.1)
\end{verbatim}

\texttt{\textbackslash{}begin\{cases\}}
\texttt{u\_t\ =\ \textbackslash{}alpha\^{}2\ u\_\{xx\},\ \&\ 0\ \textless{}\ x\ \textless{}\ L,\textbackslash{}\ 0\ \textless{}\ t\ \textless{}\ \textbackslash{}infty,\ \textbackslash{}\textbackslash{}}
\texttt{u(0,t)\ =\ T\_1,\textbackslash{}quad\ u(L,t)\ =\ T\_2,\ \&\ 0\ \textless{}\ t\ \textless{}\ \textbackslash{}infty,\ \textbackslash{}\textbackslash{}}
\texttt{u(x,0)\ =\ \textbackslash{}sin\textbackslash{}left(\textbackslash{}frac\{\textbackslash{}pi\ x\}\{L\}\textbackslash{}right),\ \&\ 0\ \textbackslash{}le\ x\ \textbackslash{}le\ L.}
\texttt{\textbackslash{}end\{cases\}}

\begin{verbatim}
(см. рис. 22.2).

![](images/_page_161_Figure_10.jpeg)

Рис. 22.2. Задача теплопровідності в площині змінних `(x,t)`.

## Перехід до безрозмірних змінних

Введемо безрозмірні змінні:
- U — безрозмірна температура,
- \xi — безрозмірна координата,
- \tau — безрозмірний час.

## Перетворення залежної змінної

Визначимо:
\end{verbatim}

\texttt{U(x,t)\ =\ \textbackslash{}frac\{u(x,t)\ -\ T\_1\}\{T\_2\ -\ T\_1\}.}

\begin{verbatim}
Тоді граничні умови стають:
\end{verbatim}

\texttt{U(0,t)=0,\ \textbackslash{}qquad\ U(L,t)=1.}

\begin{verbatim}
Рівняння:
\end{verbatim}

\texttt{U\_t\ =\ \textbackslash{}alpha\^{}2\ U\_\{xx\}.}

\begin{verbatim}
## Перетворення просторової змінної

Вводимо:
\end{verbatim}

\texttt{\textbackslash{}xi\ =\ \textbackslash{}frac\{x\}\{L\}.}

\begin{verbatim}
Тоді:
\end{verbatim}

\texttt{U\_x\ =\ \textbackslash{}frac\{1\}\{L\}\ U\_\textbackslash{}xi,\ \textbackslash{}qquad\ U\_\{xx\}\ =\ \textbackslash{}frac\{1\}\{L\^{}2\}\ U\_\{\textbackslash{}xi\textbackslash{}xi\}.}

\begin{verbatim}
Отже задача набуває вигляду:
\end{verbatim}

\texttt{\textbackslash{}begin\{cases\}}
\texttt{U\_t\ =\ \textbackslash{}frac\{\textbackslash{}alpha\^{}2\}\{L\^{}2\}\ U\_\{\textbackslash{}xi\textbackslash{}xi\},\ \&\ 0\ \textless{}\ \textbackslash{}xi\ \textless{}\ 1,\ \textbackslash{}\textbackslash{}}
\texttt{U(0,t)=0,\textbackslash{}quad\ U(1,t)=1,\ \textbackslash{}\textbackslash{}}
\texttt{U(\textbackslash{}xi,0)=\textbackslash{}dfrac\{\textbackslash{}sin(\textbackslash{}pi\ \textbackslash{}xi)-T\_1\}\{T\_2-T\_1\}.}
\texttt{\textbackslash{}end\{cases\}}

\begin{verbatim}
## Перетворення часу

Введемо безрозмірний час:
\end{verbatim}

\texttt{\textbackslash{}tau\ =\ \textbackslash{}frac\{\textbackslash{}alpha\^{}2\}\{L\^{}2\}\ t.}

\begin{verbatim}
Тоді:
\end{verbatim}

\texttt{U\_t\ =\ \textbackslash{}frac\{\textbackslash{}alpha\^{}2\}\{L\^{}2\}\ U\_\textbackslash{}tau.}

\begin{verbatim}
Отримуємо остаточну безрозмірну задачу:
\end{verbatim}

\texttt{\textbackslash{}begin\{cases\}}
\texttt{U\_\textbackslash{}tau\ =\ U\_\{\textbackslash{}xi\textbackslash{}xi\},\ \&\ 0\ \textless{}\ \textbackslash{}xi\ \textless{}\ 1,\textbackslash{}\ 0\ \textless{}\ \textbackslash{}tau\ \textless{}\ \textbackslash{}infty,\ \textbackslash{}\textbackslash{}}
\texttt{U(0,\textbackslash{}tau)=0,\textbackslash{}quad\ U(1,\textbackslash{}tau)=1,\ \textbackslash{}\textbackslash{}}
\texttt{U(\textbackslash{}xi,0)=\textbackslash{}dfrac\{\textbackslash{}sin(\textbackslash{}pi\ \textbackslash{}xi)-T\_1\}\{T\_2-T\_1\}.}
\texttt{\textbackslash{}end\{cases\}}

\begin{verbatim}
## Відновлення початкової змінної
\end{verbatim}

\texttt{u(x,t)\ =\ T\_1\ +\ (T\_2\ -\ T\_1)\textbackslash{},\ U\textbackslash{}!\textbackslash{}left(\textbackslash{}frac\{x\}\{L\},\ \textbackslash{}frac\{\textbackslash{}alpha\^{}2\ t\}\{L\^{}2\}\textbackslash{}right).}

\begin{verbatim}
## Висновки

1. Рівняння не містить фізичних параметрів.
2. Граничні умови спрощені.
3. Початкова задача стала універсальною.
4. Розв’язок легко повертається до фізичних змінних.

## ДОДАТКОВО

Для хвильового рівняння:
\end{verbatim}

\texttt{u\_\{tt\}\ =\ \textbackslash{}alpha\^{}2\ u\_\{xx\}}

\begin{verbatim}
використовують заміну:
\end{verbatim}

\texttt{\textbackslash{}xi\ =\ \textbackslash{}frac\{x\}\{L\},\ \textbackslash{}qquad\ \textbackslash{}tau\ =\ \textbackslash{}frac\{\textbackslash{}alpha\}\{L\}\ t,}

\begin{verbatim}
що приводить до:
\end{verbatim}

\texttt{u\_\{\textbackslash{}tau\textbackslash{}tau\}\ =\ u\_\{\textbackslash{}xi\textbackslash{}xi\}.}

\begin{verbatim}
## ЗАВДАННЯ

1. Перетворити задачу (22.5) до безрозмірного вигляду.
2. Записати задачу теплопровідності у безрозмірній формі.
3. Довести усунення параметра \alpha заміною часу.
4. Вибрати змінну \xi для рівняння переносу u_t + v u_x = 0.
5. вона відокремлює хвильові, дифузійні та стаціонарні процеси;
6. саме для канонічних форм існує найбільше готових аналітичних методів;
7. чисельні схеми для гіперболічних задач зазвичай теж будують уже після такого перетворення.

Як тільки розв'язок знайдено в нових координатах, його завжди можна перетворити назад у початкові координати (x,y).

# ЗАВДАННЯ

1. Визначте, чи є наступні рівняння гіперболічними, параболічними чи еліптичними:
   u_{xx}u_{xy}=0; u_{tt}=u_{xx}+u_x+hu; u_{xx}+3u_{yy}=\sin x; u_{xx}+u_{yy}=f(x,y); u_{rr}+\frac{1}{r}u_r+\frac{1}{r^2}u_{\theta\theta}=f(r,\theta).

2. Отримайте співвідношення (23.4), (23.5) та (23.6).

3. Переконайтеся, що рівняння
\end{verbatim}

3u\_\{xx\} + 7u\_\{xy\} + 2u\_\{yy\} = 0

\begin{verbatim}
є гіперболічним при всіх x та y, і знайдіть його характеристики.

4. Зведіть рівняння з попередньої задачі до канонічної форми
\end{verbatim}

\texttt{u\_\{\textbackslash{}xi\textbackslash{}eta\}\ =\ \textbackslash{}Phi(\textbackslash{}xi,\ \textbackslash{}eta,\ u,\ u\_\{\textbackslash{}xi\},\ u\_\{\textbackslash{}eta\}).}

\begin{verbatim}
5. Продовжте задачу 4 і приведіть це рівняння до другої канонічної форми
\end{verbatim}

\texttt{u\_\{\textbackslash{}alpha\textbackslash{}alpha\}\ -\ u\_\{\textbackslash{}beta\textbackslash{}beta\}\ =\ \textbackslash{}Psi(\textbackslash{}alpha,\ \textbackslash{}beta,\ u,\ u\_\{\textbackslash{}alpha\},\ u\_\{\textbackslash{}beta\}).}

\begin{verbatim}
6. Найдите характеристики уравнения
\end{verbatim}

u\_\{xx\} + 4u\_\{xy\} = 0.

\begin{verbatim}
Перетворіть рівняння на нові координати, розв'яжіть його і поверніться до координат x і y, щоб отримати розв'язок початкової здогадки.

# Лекція 23. КЛАСИФІКАЦІЯ РІВНЯНЬ З ЧАСТКОВИМИ ПОХІДНИМИ (КАНОНІЧНА ФОРМА ГІПЕРБОЛІЧНОГО РІВНЯННЯ)

ЦЕЛЬ ЛЕКЦИИ: Показать, что линейное уравнение второго порядка с двумя независимыми переменными
\end{verbatim}

Au\_\{xx\} + Bu\_\{xy\} + Cu\_\{yy\} + Du\_x + Eu\_y + Fu = G

\begin{verbatim}
(A, B, C, D, E, F и G являются функциями x и y или могут быть константами) относится к одному из следующих типов:

- 1) гиперболическому (если  B^2 - 4AC > 0 ),
- 2) параболическому (если  B^2 - 4AC = 0 ),
- 3) эллиптическому (если  B^2 - 4AC < 0 ).

Показать, какие новые переменные  \xi = \xi(x, y)  и  \eta = \eta(x, y)  можно ввести вместо x и y, чтобы упростить исходное уравнение. В новых переменных  \xi  и  \eta  уравнение с частными производными приводится к одному из следующих видов (в вависимости от того, будет ли величина  B^2 - 4AC  положительна, равна нулю или отрицательна):

- 1.
\end{verbatim}

\texttt{\textbackslash{}begin\{cases\}\ u\_\{\textbackslash{}xi\textbackslash{}xi\}\ -\ u\_\{\textbackslash{}eta\textbackslash{}eta\}\ =\ \textbackslash{}Psi(\textbackslash{}xi,\ \textbackslash{}eta,\ u,\ u\_\{\textbackslash{}xi\},\ u\_\{\textbackslash{}eta\})\ \&\ \textbackslash{}text\{(две\ канонические\ формы\ для\ гиперболического\ уравнения),\}\ \textbackslash{}end\{cases\}}

\begin{verbatim}
- 2.  u_{\eta\eta} = \Phi(\xi, \eta, u, u_{\xi}, u_{\eta})  (каноническая форма параболического уравнения),
- 3.  u_{\xi\xi} + u_{\eta\eta} = \Phi(\xi, \eta, u, u_{\xi}, u_{\eta})  (каноническая форма эллиптического уравнения),

где  \Phi  и  \Psi —функции от частных производных первого порядка, зависимой переменной u и новых независимых переменных  \xi  и  \eta . Конкретный вид функций  \Phi  и  \Psi  зависит от исходного уравнения.

У читателя может возникнуть мысль, что глава, посвященная классификации уравнений, должна размещаться в самом начале книги. Вероятно, это правильно, и многие книги действительно начинаются с этого материала. Однако правда и то, что многие студенты не могут с энтузиазмом изучать нечто такое, о чем у них нет никакого представления. Именно по этой причине мы отложили изучение темы, посвященной классификации уравнений с частными производными, до настоящего момента.

Наша задача - классифицировать все уравнения вида

(23 1)
\end{verbatim}

Au\_\{xx\} + Bu\_\{xy\} + Cu\_\{yy\} + Du\_x + Eu\_y + Fu = G,

\begin{verbatim}
где A, B, C, D, E, F и G—вообще говоря, функции [x и y. Согласно предлагаемой схеме классификации уравнение называется

1) гиперболическим в точке  (x_0, y_0) , если  B^2(x_0, y_0) - 4A(x_0, y_0) \times  \times C(x_0, y_0) > 0 2) параболнческим в точке  (x_0, y_0) , если  B^2(x_0, y_0) - 4A(x_0, y_0) \times C(x_0, y_0) = 0 ,

3) эллиптическим в точке  (x_0, y_0) , если  B^2(x_0, y_0) - 4A(x_0, y_0) \times  \times C(x_0, y_0) < 0 и в зависимости от типа его можно привести к одной из канонических форм. Для того чтобы было легче разобраться в классификационной схеме, мы приведем четыре примера гиперболических, параболических и эллиптических уравнений.

## Примеры гиперболических, параболических и эллиптических уравнений 10

1. У равнение теплопроводности  u_t = u_{xx}  — линейное уравнение второго порядка вида (23.1) с коэффициентами
\end{verbatim}

A=1

, \texttt{B=0} , \texttt{C=0} , \texttt{D=0} , \texttt{E=-1} ,
\texttt{F=0} , \texttt{G=0} ,

так что \texttt{B\^{}2-4AC=0} при всех x и t. Следовательно, это
уравнение будет параболическим при всех x и t. В общем уравнении (23 1)
переменная y играет роль времени t. Результат не изменится, если в общем
уравнении поменять местами x и y.

\begin{enumerate}
\def\labelenumi{\arabic{enumi}.}
\setcounter{enumi}{1}
\tightlist
\item
  Волновое уравнение \texttt{u\_\{tt\}\ =\ u\_\{xx\}} также принадлежит
  классу уравнений вида (23.1) с коэффициентами
\end{enumerate}

\begin{verbatim}
A=1
\end{verbatim}

\begin{verbatim}
,  B=0 ,  C=-1 ,  D=E=F=G=0 .

Так как  B^2 - 4AC = 4  для всех x и t, то это уравнение является еиперболическим при всех x и t.

<sup>1)</sup> Следует указать, что приведенные в примерах уравнения являются частными случаями более общих параболических, гиперболических в эллиптических уравнений, а основные их свойства характерны и для общего случая.

3. У равнение Лапласа  u_{xx} + u_{yy} = 0  является эллиптическим при всех x и y, так как  B^2 - 4AC = -4 < 0 .

4. Линейное уравнение  xu_{xx} + u_{yy} = \sin x  с переменными коэффициентами также принадлежит классу уравнений вида (23.1), но для него  B^2-4AC=-4x  и, следовательно, уравнение будет эллиптическим при x > 0,

лараболическим при x=0, гиперболическим при x < 0.

Этот пример иллюстрирует тот факт, что тип уравнения о переменными коэффициентами может изменяться от точки к точке.

Читатель должен обратить внимание на следующее: тип уравнения (23.1) определяется только коэффициентами при вторых првизводных и никак не зависит ни от коэффициентов при первых производных и самой функции, ни от свободного члена.

Обратимся теперь к главной теме нашей лекции — приведению гиперболических уравнений к каноническому виду. Известно, что если в данной области пространства уравнение является гиперболическим, то вместо x и y можно ввести координаты  \xi  и  \eta (характеристические переменные) таким образом, что уравнение примет простейшую форму

(23.2)
\end{verbatim}

\texttt{u\_\{\textbackslash{}xi\textbackslash{}eta\}\ =\ \textbackslash{}Phi(\textbackslash{}xi,\ \textbackslash{}eta,\ u,\ u\_\{\textbackslash{}xi\},\ u\_\{\textbackslash{}eta\}).}

\begin{verbatim}
Это уравнение содержит только одну производную второго порядка  u_{\xi\eta} , а функция  \Phi(\xi,\eta,u,u_{\xi},u_{\eta})  зависит только от новых независимых переменных § и η, зависимой переменной и и первых производных из и ип. Конкретный вид функции Ф зависит, конечно, от исходного уравнения и формул перехода к новым координатам § и п. Научимся отыскивать эти функции.

## Каноническая форма гиперболического уравнения

Рассмотрим уравнение общего вида

(23.3)
\end{verbatim}

Au\_\{xx\} + Bu\_\{xy\} + Cu\_\{yy\} + Du\_x + Eu\_y + Fu = G

\begin{verbatim}
и пусть в интересующей нас области  B^2 - 4AC > 0 . Нам нужно ввести новые координаты 1)
\end{verbatim}

\texttt{\textbackslash{}xi\ =\ \textbackslash{}xi(x,\ y),}
\texttt{\textbackslash{}eta\ =\ \textbackslash{}eta(x,\ y)}

\begin{verbatim}
так, чтобы в уравнении (23.3) осталась только одна частная производная второго порядка  u_{\xi\eta}  (не следует пытаться искать такую замену переменных, с помощью которой гиперболическое уравне-

<sup>1)</sup> Необходимо, чтобы это преобразование координат было локально обратимо и дважды дифференцируемо. — Прим. ред.

ние можно было бы привести и параболической или заливитической канонической форме, это невозможно сделать).

Вычислим сначала частные производные
\end{verbatim}

\texttt{\textbackslash{}begin\{array\}\{c\}\ u\_x\ =\ u\_\textbackslash{}xi\ \textbackslash{}xi\_x\ +\ u\_\textbackslash{}eta\ \textbackslash{}eta\_x,\ \textbackslash{}\textbackslash{}\ u\_y\ =\ u\_\textbackslash{}xi\ \textbackslash{}xi\_y\ +\ u\_\textbackslash{}eta\ \textbackslash{}eta\_y,\ \textbackslash{}\textbackslash{}\ (23.4)\ \textbackslash{}\ u\_\{xx\}\ =\ u\_\{\textbackslash{}xi\textbackslash{}xi\}\ \textbackslash{}xi\_x\^{}2\ +\ 2u\_\{\textbackslash{}xi\textbackslash{}eta\}\ \textbackslash{}xi\_x\ \textbackslash{}eta\_x\ +\ u\_\{\textbackslash{}eta\textbackslash{}eta\}\ \textbackslash{}eta\_x\^{}2\ +\ u\_\{\textbackslash{}xi\}\ \textbackslash{}xi\_\{xx\}\ +\ u\_\{\textbackslash{}eta\}\ \textbackslash{}eta\_\{xx\},\ \textbackslash{}\textbackslash{}\ u\_\{xy\}\ =\ u\_\{\textbackslash{}xi\textbackslash{}xi\}\ \textbackslash{}xi\_x\ \textbackslash{}xi\_y\ +\ u\_\{\textbackslash{}xi\textbackslash{}eta\}\ (\textbackslash{}xi\_x\ \textbackslash{}eta\_y\ +\ \textbackslash{}xi\_y\ \textbackslash{}eta\_x)\ +\ u\_\{\textbackslash{}eta\textbackslash{}eta\}\ \textbackslash{}eta\_x\ \textbackslash{}eta\_y\ +\ u\_\{\textbackslash{}xi\}\ \textbackslash{}xi\_\{xy\}\ +\ u\_\{\textbackslash{}eta\}\ \textbackslash{}eta\_\{xy\},\ \textbackslash{}\textbackslash{}\ u\_\{yy\}\ =\ u\_\{\textbackslash{}xi\textbackslash{}xi\}\ \textbackslash{}xi\_y\^{}2\ +\ 2u\_\{\textbackslash{}xi\textbackslash{}eta\}\ \textbackslash{}xi\_y\ \textbackslash{}eta\_y\ +\ u\_\{\textbackslash{}eta\textbackslash{}eta\}\ \textbackslash{}eta\_y\^{}2\ +\ u\_\{\textbackslash{}xi\}\ \textbackslash{}xi\_\{yy\}\ +\ u\_\{\textbackslash{}eta\}\ \textbackslash{}eta\_\{yy\}.\ \textbackslash{}end\{array\}}

\begin{verbatim}
Подставляя эти соотношения в исходное уравнение (23.3), после простых, но громоздких вычислений получаем

(23.5)
\end{verbatim}

\texttt{\textbackslash{}overline\{A\}u\_\{\textbackslash{}xi\textbackslash{}xi\}\ +\ \textbackslash{}overline\{B\}u\_\{\textbackslash{}xi\textbackslash{}eta\}\ +\ \textbackslash{}overline\{C\}u\_\{\textbackslash{}eta\textbackslash{}eta\}\ +\ \textbackslash{}overline\{D\}u\_\{\textbackslash{}xi\}\ +\ \textbackslash{}overline\{E\}u\_\{\textbackslash{}eta\}\ +\ \textbackslash{}overline\{F\}u\ =\ \textbackslash{}overline\{G\},}
\texttt{\textbackslash{}overline\{A\}\ =\ A\textbackslash{}xi\_\{x\}\^{}\{2\}\ +\ B\textbackslash{}xi\_\{x\}\textbackslash{}xi\_\{y\}\ +\ C\textbackslash{}xi\_\{y\}\^{}\{2\},}
\texttt{\textbackslash{}overline\{B\}\ =\ 2A\textbackslash{}xi\_\{x\}\textbackslash{}eta\_\{x\}\ +\ B(\textbackslash{}xi\_\{x\}\textbackslash{}eta\_\{y\}\ +\ \textbackslash{}xi\_\{y\}\textbackslash{}eta\_\{x\})\ +\ 2C\textbackslash{}xi\_\{y\}\textbackslash{}eta\_\{y\},}
\texttt{\textbackslash{}overline\{C\}\ =\ A\textbackslash{}eta\_\{x\}\^{}\{2\}\ +\ B\textbackslash{}eta\_\{x\}\textbackslash{}eta\_\{y\}\ +\ C\textbackslash{}eta\_\{y\}\^{}\{2\},}
\texttt{\textbackslash{}overline\{D\}\ =\ A\textbackslash{}xi\_\{xx\}\ +\ B\textbackslash{}xi\_\{xy\}\ +\ C\textbackslash{}xi\_\{yy\}\ +\ D\textbackslash{}xi\_\{x\}\ +\ E\textbackslash{}xi\_\{y\},}
\texttt{\textbackslash{}overline\{E\}\ =\ A\textbackslash{}eta\_\{xx\}\ +\ B\textbackslash{}eta\_\{xy\}\ +\ C\textbackslash{}eta\_\{yy\}\ +\ D\textbackslash{}eta\_\{x\}\ +\ E\textbackslash{}eta\_\{y\},}
\texttt{\textbackslash{}overline\{F\}\ =\ F,}
\texttt{\textbackslash{}overline\{G\}\ =\ G.}

\begin{verbatim}
Читателю предоставится возможность проделать эти выкладки при решении задач.

Наш следующий шаг—выбрать функции  \xi = \xi(x, y)  и  \eta = \eta(x, y)  таким образом, чтобы коэффициенты  \overline{A}  и  \overline{C}  обратились в нуль. Это позволит нам привести исходное уравиение к канонической форме. Итак, необходимо потребовать, чтобы
\end{verbatim}

\texttt{\textbackslash{}bar\{A\}\ =\ A\textbackslash{}xi\_x\^{}2\ +\ B\textbackslash{}xi\_x\textbackslash{}xi\_y\ +\ C\textbackslash{}xi\_y\^{}2\ =\ 0,}

\begin{verbatim}
\bar{C} = A\eta_x^2 + B\eta_x\eta_y + C\eta_y^2 = 0. Эти уравнения можно представить в виде
\end{verbatim}

\texttt{\textbackslash{}begin\{split\}\ \&A\textbackslash{}left{[}\textbackslash{}xi\_\{x\}/\textbackslash{}xi\_\{y\}\textbackslash{}right{]}\^{}\{2\}+B\textbackslash{}left{[}\textbackslash{}xi\_\{x\}/\textbackslash{}xi\_\{y\}\textbackslash{}right{]}+C=0,\textbackslash{}\textbackslash{}\ \&A\textbackslash{}left{[}\textbackslash{}eta\_\{x\}/\textbackslash{}eta\_\{y\}\textbackslash{}right{]}\^{}\{2\}+B\textbackslash{}left{[}\textbackslash{}eta\_\{x\}/\textbackslash{}eta\_\{y\}\textbackslash{}right{]}\^{}\{2\}+C=0.\ \textbackslash{}end\{split\}}

\begin{verbatim}
Разрешив их относительно  [\xi_x/\xi_y]  и  [\eta_x/\eta_y] , получаем

(23.7)
\end{verbatim}

\texttt{\textbackslash{}begin\{bmatrix\}\ \textbackslash{}xi\_x/\textbackslash{}xi\_y\ \textbackslash{}end\{bmatrix\}\ =\ \textbackslash{}frac\{-B\ +\ \textbackslash{}sqrt\{B\^{}2\ -\ 4AC\}\}\{2A\}}

\begin{verbatim}
(характеристические
\end{verbatim}

\texttt{{[}\textbackslash{}eta\_x/\textbackslash{}eta\_y{]}\ =\ \textbackslash{}frac\{-B\ -\ \textbackslash{}sqrt\{B\^{}2\ -\ 4AC\}\}\{2A\}}

\begin{verbatim}
уравнения).

Kaждое из квадратных уравнений для  [\xi_x/\xi_y]  и  [\eta_x/\eta_y]  имеет два кория, но мы оставляем только по одному, но так чтобы они были различны.

![](images/_page_170_Figure_3.jpeg)

РИС. 23.1. Характеристики  \xi(x, y) = c  и  \eta(x, y) = c .

Задача свелась к нахождению двух функций  \xi(x, y)  и  \eta(x, y) , таких, чтобы отношения  [\xi_x/\xi_y]  и  [\eta_x/\eta_y]  удовлетворяли уравнению (23.7). Найти такие функции очень просто 1), если обратить внимание на рис. 23.1. Чтобы понять, каким образом можно это сделать, рассмотрим простое уравнение
\end{verbatim}

u\_\{xx\} - 4u\_\{yy\} + u\_x = 0.

\begin{verbatim}
Его характеристики определяются из уравнений
\end{verbatim}

\texttt{\textbackslash{}begin\{aligned\}}
\texttt{\textbackslash{}frac\{dy\}\{dx\}\ =\ -\textbackslash{}frac\{\textbackslash{}xi\_x\}\{\textbackslash{}xi\_y\}\ =\ \textbackslash{}frac\{B\ -\ \textbackslash{}sqrt\{B\^{}2\ -\ 4AC\}\}\{2A\}\ =\ -2,\ \textbackslash{}qquad}
\texttt{\textbackslash{}frac\{dy\}\{dx\}\ =\ -\textbackslash{}frac\{\textbackslash{}eta\_x\}\{\textbackslash{}eta\_y\}\ =\ \textbackslash{}frac\{B\ +\ \textbackslash{}sqrt\{B\^{}2\ -\ 4AC\}\}\{2A\}\ =\ 2.}
\texttt{\textbackslash{}end\{aligned\}}

\begin{verbatim}
Решая эти уравнения относительно у, получаем
\end{verbatim}

y = -2x + c\_1, y = 2x + c\_2.

\begin{verbatim}
Чтобы найти  \xi  и  \eta , разрешим полученные уравнения относительно констант  c_1  и  c_2 , оставляя их в правой части и перенося остальные слагаемые в левую.

Выражения, стоящие в левых частях, примем за искомые функции § и пр
\end{verbatim}

\texttt{\textbackslash{}xi\ =\ y\ +\ 2x\ =\ c\_1,}
\texttt{\textbackslash{}eta\ =\ y\ -\ 2x\ =\ c\_2.}

\begin{verbatim}
<sup>1)</sup> См., например, Тихонов А. Н., Самарский А. А. Уравнения математической физики. — М.: Наука, 1972 и др. издания. — Прим. ред.

Очевидно, введенные таким образом функции  \xi  и  \eta  удовлетворяют уравнениям характеристик. Эти новые координаты изображены на рис. 23.2. На этом мы заканчиваем обсуждение вопроса о том, как найти новые координаты. Осталось последнее — определить новый вид исходного уравнения.

![](images/_page_171_Figure_3.jpeg)

РИС. 23.2. Характеристические координаты для уравнения  u_{xx} - 4u_{yy} + u_x = 0 .

Это совсем просто: чтобы найти каноническую форму уравнения с частными производными, подставим новые координаты  \xi(x,y)  и  \eta(x,y)  в уравнение
\end{verbatim}

\texttt{\textbackslash{}overline\{A\}u\_\{\textbackslash{}xi\textbackslash{}xi\}\ +\ \textbackslash{}overline\{B\}u\_\{\textbackslash{}xi\textbackslash{}eta\}\ +\ \textbackslash{}overline\{C\}u\_\{\textbackslash{}eta\textbackslash{}eta\}\ +\ \textbackslash{}overline\{D\}u\_\{\textbackslash{}xi\}\ +\ \textbackslash{}overline\{E\}u\_\{\textbackslash{}eta\}\ +\ \textbackslash{}overline\{F\}u\ =\ \textbackslash{}overline\{G\}}

,

где \texttt{\textbackslash{}overline\{A\}} ,
\texttt{\textbackslash{}overline\{B\}} ,
\texttt{\textbackslash{}overline\{C\}} ,
\texttt{\textbackslash{}overline\{D\}} ,
\texttt{\textbackslash{}overline\{E\}} ,
\texttt{\textbackslash{}overline\{F\}} и
\texttt{\textbackslash{}overline\{G\}} определены формулами (23.6).

Прежде чем закончить лекцию, давайте на конкретном примере посмотрим,
как «работает» общий метод приведения к каноническому виду.

Приведение гиперболического уравнения
\texttt{y\^{}2u\_\{xx\}-x\^{}2u\_\{yy\}=0} к канонической форме

Рассмотрим уравнение

\begin{verbatim}
y^2u_{xx}-x^2u_{yy}=0
\end{verbatim}

\begin{verbatim}
,  x>0 ,  y>0 ,

которое является гиперболическим в первом квадранте. Найдем новые координаты таким образом, чтобы привести исходное уравнение к канонической форме.

ШАГ 1. Решаем уравнения характеристик
\end{verbatim}

\texttt{\textbackslash{}frac\{dy\}\{dx\}\ =\ \textbackslash{}frac\{B\ +\ \textbackslash{}sqrt\{B\^{}2\ -\ 4AC\}\}\{2A\}\ =\ \textbackslash{}frac\{x\}\{y\}}

\begin{verbatim}
(напомним, что этот шаг
\end{verbatim}

\texttt{\textbackslash{}frac\{dy\}\{dx\}\ =\ \textbackslash{}frac\{B\ -\ \textbackslash{}sqrt\{B\^{}2\ -\ 4AC\}\}\{2A\}\ =\ -\textbackslash{}frac\{x\}\{y\}.}

\begin{verbatim}
обеспечивает  \overline{A} = \overline{C} = 0 ),

![](images/_page_172_Figure_2.jpeg)

РИС, 23.3. Новые карактеристические координаты.

Проинтегрировав эти обыкновенные дифференциальные уравнения методом разделения переменных, получаем два неявных соотнощения (при желании можно было бы явно выразить у через х)
\end{verbatim}

\texttt{y\^{}2\ -\ x\^{}2\ =\ \textbackslash{}text\{const\},}
\texttt{y\^{}2\ +\ x\^{}2\ =\ \textbackslash{}text\{const\}.}

\begin{verbatim}
Следовательно, новые координаты следует вводить по формулам
\end{verbatim}

\texttt{\textbackslash{}xi\ =\ y\^{}2\ -\ x\^{}2}

,\\
\texttt{\textquotesingle{}\textbackslash{}eta\ =\ y\^{}2\ +\ x\^{}2\textquotesingle{}\ .}

Эти новые координаты изображены на рис. 23.3. Теперь следует получить
новое уравнение. Вычислим все коэффициенты

\texttt{\textbackslash{}overline\{A\}\ =\ 0} (так и должно быть; мы
специально выбирали \texttt{\textbackslash{}xi} и
\texttt{\textbackslash{}eta} так, чтобы этот коэффициент обращался в
нуль).

\begin{verbatim}
\overline{B} = 2A\xi_x\eta_x + B(\xi_x\eta_y + \xi_y\eta_x) + 2C\xi_y\eta_y = -16x^2y^2
\end{verbatim}

\begin{verbatim}
,

 \overline{C} = 0  (по той же причине, что и  \overline{A} = 0 ),
\end{verbatim}

\texttt{\textbackslash{}bar\{D\}\ =\ A\textbackslash{}xi\_\{xx\}\ +\ B\textbackslash{}xi\_\{xy\}\ +\ C\textbackslash{}xi\_\{yy\}\ +\ D\textbackslash{}xi\_x\ +\ E\textbackslash{}xi\_y\ =\ -2(x\^{}2\ +\ y\^{}2)}
\texttt{\textbackslash{}overline\{E\}\ =\ A\textbackslash{}eta\_\{xx\}\ +\ B\textbackslash{}eta\_\{xy\}\ +\ C\textbackslash{}eta\_\{yy\}\ +\ D\textbackslash{}eta\_x\ +\ E\textbackslash{}eta\_y\ =\ 2(y\^{}2\ -\ x\^{}2),}
\texttt{\textbackslash{}overline\{F\}\ =\ F\ =\ 0}

,

\begin{verbatim}
\overline{G} = G = 0
\end{verbatim}

\begin{verbatim}
,

подставим их в уравнение
\end{verbatim}

\texttt{\textbackslash{}overline\{A\}u\_\{\textbackslash{}xi\textbackslash{}xi\}\ +\ \textbackslash{}overline\{B\}u\_\{\textbackslash{}xi\textbackslash{}eta\}\ +\ \textbackslash{}overline\{C\}u\_\{\textbackslash{}eta\textbackslash{}eta\}\ +\ \textbackslash{}overline\{D\}u\_\{\textbackslash{}xi\}\ +\ \textbackslash{}overline\{E\}u\_\{\textbackslash{}eta\}\ +\ \textbackslash{}overline\{F\}u\ =\ \textbackslash{}overline\{G\}}

\begin{verbatim}
и получим
\end{verbatim}

\texttt{u\_\{\textbackslash{}xi\textbackslash{}eta\}\ =\ \textbackslash{}frac\{-(x\^{}2\ +\ y\^{}2)\ u\_\{\textbackslash{}xi\}\ +\ (y\^{}2\ -\ x\^{}2)\ u\_\{\textbackslash{}eta\}\}\{8x\^{}2y\^{}2\}.}

\begin{verbatim}
ШАГ 2. Выражая, наконец, x и y через  \xi  и  \eta , окончательно получаем

(23.8)
\end{verbatim}

\texttt{u\_\{\textbackslash{}xi\textbackslash{}eta\}\ =\ \textbackslash{}frac\{\textbackslash{}eta\ u\_\{\textbackslash{}xi\}\ -\ \textbackslash{}xi\ u\_\{\textbackslash{}eta\}\}\{2(\textbackslash{}xi\^{}2\ -\ \textbackslash{}eta\^{}2)\}.}

\begin{verbatim}
Задача решена. Теперь мы знаем, что нужно сделать, чтобы найти новые координаты и привести уравнение к канонической форме.

### ЗАМЕЧАНИЯ

1. На самом деле для гиперболических уравнений существует две канонические формы. Вторую можно получить из первов заменой переменных вида
\end{verbatim}

\texttt{\textbackslash{}alpha\ =\ \textbackslash{}alpha\ (\textbackslash{}xi,\ \textbackslash{}eta)\ =\ \textbackslash{}xi\ +\ \textbackslash{}eta,}

\begin{verbatim}
\beta = \beta (\xi, \eta) = \xi - \eta. Воспользуемся этой заменой для преобразования уравнения (23.8). Так как
\end{verbatim}

\texttt{\textbackslash{}begin\{array\}\{l\}\ u\_\{\textbackslash{}xi\}\ =\ u\_\{\textbackslash{}alpha\}\textbackslash{}alpha\_\{\textbackslash{}xi\}\ +\ u\_\{\textbackslash{}beta\}\textbackslash{}beta\_\{\textbackslash{}xi\}\ =\ u\_\{\textbackslash{}alpha\}\ +\ u\_\{\textbackslash{}beta\}.\ \textbackslash{}\textbackslash{}\ u\_\{\textbackslash{}eta\}\ =\ u\_\{\textbackslash{}alpha\}\textbackslash{}alpha\_\{\textbackslash{}eta\}\ +\ u\_\{\textbackslash{}beta\}\textbackslash{}beta\_\{\textbackslash{}eta\}\ =\ u\_\{\textbackslash{}alpha\}\ -\ u\_\{\textbackslash{}beta\},\ \textbackslash{}\textbackslash{}\ u\_\{\textbackslash{}xi\textbackslash{}eta\}\ =\ u\_\{\textbackslash{}alpha\textbackslash{}alpha\}\textbackslash{}alpha\_\{\textbackslash{}eta\}\ +\ u\_\{\textbackslash{}alpha\textbackslash{}beta\}\textbackslash{}beta\_\{\textbackslash{}eta\}\ +\ u\_\{\textbackslash{}beta\textbackslash{}alpha\}\textbackslash{}alpha\_\{\textbackslash{}eta\}\ +\ u\_\{\textbackslash{}beta\textbackslash{}beta\}\textbackslash{}beta\_\{\textbackslash{}eta\}\ =\ u\_\{\textbackslash{}alpha\textbackslash{}alpha\}\ -\ u\_\{\textbackslash{}beta\textbackslash{}beta\},\ \textbackslash{}end\{array\}}

\begin{verbatim}
TO

(23 9)
\end{verbatim}

\texttt{u\_\{\textbackslash{}alpha\textbackslash{}alpha\}\ -\ u\_\{\textbackslash{}beta\textbackslash{}beta\}\ =\ \textbackslash{}frac\{-\textbackslash{}beta\ u\_\{\textbackslash{}alpha\}\ -\ \textbackslash{}alpha\ u\_\{\textbackslash{}beta\}\}\{2\textbackslash{}alpha\textbackslash{}beta\}.}

\begin{verbatim}
При желании можно выразить  \alpha  и  \beta  через исходные координаты x и y. В результате получаем
\end{verbatim}

\texttt{\textbackslash{}alpha\ =\ \textbackslash{}xi\ +\ \textbackslash{}eta\ =\ (y\^{}2\ -\ x\^{}2)\ +\ (y\^{}2\ +\ x\^{}2)\ =\ 2y\^{}2,}
\texttt{\textbackslash{}beta\ =\ \textbackslash{}xi\ -\ \textbackslash{}eta\ =\ (y\^{}2\ -\ x\^{}2)\ -\ (y\^{}2\ +\ x\^{}2)\ =\ -2x\^{}2.}

\begin{verbatim}
- Читатель может спросить, зачем нужно классифицировать и приводить к канонической форме уравнения с частными производными.
  - а) Разбиение уравнений на классы гиперболических, параболических и эллиптических соответствует разбиению физических процессов на три основных класса: волновых, диффузионных и стационарных. Математические особенности решений уравнений этих трех типов совершенно различны.
  - б) В большинстве работ, посвященных решению гиперболических задач, предполагается, что исходное уравнение записано в канонической форме, т. е. в виде
\end{verbatim}

\texttt{u\_\{\textbackslash{}xi\textbackslash{}xi\}\ -\ u\_\{\textbackslash{}eta\textbackslash{}eta\}\ =\ \textbackslash{}Psi(\textbackslash{}xi,\ \textbackslash{},\ \textbackslash{}eta,\ \textbackslash{},\ u\_\{\textbackslash{}xi\},\ \textbackslash{},\ u,\ u\_\{\textbackslash{}eta\}).}

\begin{verbatim}
Если у нас есть уравнение и мы хотим получить его решение, мы должны преобразовать его к канонической форме и воспользоваться известными результатами.

в) Для численного решения канонического гиперболического уравнения создано много различных программ для ЭВМ.

Функция  \Phi  ( \xi ,  \eta , u,  u_{\xi} ,  u_{\eta} ) вводится в  \Theta ВМ в виде подпрограммы, так что сначала необходимо привести уравнение к канонической форме. После того как решение в новых координатах найдено, его всегда можно преобразовать обратно к исходным координатам.

## ЗАДАЧИ

- 1. Выясните, являются ли приведенные ниже уравнения гиперболическими, параболическими или эллиптическими:
  - a) u_{xx}+u_{xy}=0,
  - б) u_{tt}=u_{xx}+u_x+hu,
  - в) u_{xx}+3u_{yy}=\sin x,
  - г) u_{xx}+u_{yy}=f(x,y),
  - д) u_{rr}+\frac{1}{r}u_r+\frac{1}{r^2}u_{\theta\theta}=f(r,\theta).
- 2. Получите соотношения (23.4), (23.5) и (23.6).
- 3. Убедитесь в том, что уравнение
\end{verbatim}

3u\_\{xx\} + 7u\_\{xy\} + 2u\_\{yy\} = 0

\begin{verbatim}
гиперболического типа при всех x и y и найдите характеристики.
- 4. Приведите уравнение из предыдущей задачи к каноническому виду
\end{verbatim}

\texttt{u\_\{\textbackslash{}xi\textbackslash{}eta\}\ =\ \textbackslash{}Phi(\textbackslash{}xi,\ \textbackslash{}eta,\ u,\ u\_\{\textbackslash{}xi\},\ u\_\{\textbackslash{}eta\}).}

\begin{verbatim}
- 5. Продолжите задачу 4 и приведите уравнение ко второй канонической форме
\end{verbatim}

\texttt{u\_\{\textbackslash{}alpha\textbackslash{}alpha\}\ -\ u\_\{\textbackslash{}beta\textbackslash{}beta\}\ =\ \textbackslash{}Psi(\textbackslash{}alpha,\ \textbackslash{}beta,\ u,\ u\_\{\textbackslash{}alpha\},\ u\_\{\textbackslash{}beta\}).}

\begin{verbatim}
- 6. Найдите характеристики уравнения
\end{verbatim}

u\_\{xx\} + 4u\_\{xy\} = 0.

\begin{verbatim}
Преобразуйте уравнение к новым координатам, решите и, вернувшись к координатам x и y, получите решение исходной задачи.

# Лекція 24. ХВИЛЬОВЕ РІВНЯННЯ У ВІЛЬНОМУ ПРОСТОРІ (ДВОВИМІРНІ ТА ТРИВИМІРНІ ЗАДАЧІ)
МЕТА ЛЕКЦІЇ: Узагальнити формулу д'Аламбера на двовимірний і тривимірний випадки та пояснити геометричний зміст відповідних розв'язків.

Розглянемо задачу Коші для хвильового рівняння у тривимірному просторі:
\end{verbatim}

\texttt{u\_\{tt\}=c\^{}2\textbackslash{}Delta\ u,}
\texttt{\textbackslash{}qquad}
\texttt{(x,y,z)\textbackslash{}in\textbackslash{}mathbb\{R\}\^{}3,}
\texttt{\textbackslash{}qquad\ t\textgreater{}0,}
\texttt{u(x,y,z,0)=\textbackslash{}varphi(x,y,z),}
\texttt{\textbackslash{}qquad\ u\_t(x,y,z,0)=\textbackslash{}psi(x,y,z).}

\begin{verbatim}
Тут оператор Лапласа має вигляд
\end{verbatim}

\texttt{\textbackslash{}Delta\ u=u\_\{xx\}+u\_\{yy\}+u\_\{zz\}.}

\begin{verbatim}
Одновимірні хвилі в просторі називають плоскими, двовимірні — циліндричними, а тривимірні — сферичними. Для сферичних хвиль роль формули д'Аламбера відіграє **формула Пуассона**.

## Формула Пуассона

Розв'язок задачі Коші у \mathbb{R}^3 можна подати у вигляді
\end{verbatim}

\texttt{u(x,y,z,t)=\textbackslash{}frac\{\textbackslash{}partial\}\{\textbackslash{}partial\ t\}\textbackslash{}!\textbackslash{}left{[}t\textbackslash{},\textbackslash{}overline\{\textbackslash{}varphi\}(x,y,z,t)\textbackslash{}right{]}+t\textbackslash{},\textbackslash{}overline\{\textbackslash{}psi\}(x,y,z,t),}

\begin{verbatim}
де \overline{\varphi} та \overline{\psi} позначають середні значення початкових функцій по сфері радіуса ct з центром у точці (x,y,z).

![](images/_page_177_Picture_3.jpeg)

Фіг. 24.1. Середнє значення початкових даних береться по сфері радіуса ct.

Геометричний зміст формули простий: значення розв'язку в точці (x,y,z) у момент часу t залежить лише від тих початкових даних, що лежать на сфері радіуса ct.

![](images/_page_177_Picture_7.jpeg)

Фіг. 24.2. Початкове збурення поширюється сферично в усіх напрямках.

## Принцип Гюйгенса

Якщо початкове збурення локалізоване в обмеженій області, то в тривимірному випадку розв'язок має чітко визначений передній і задній фронт. Це означає, що після проходження хвильового фронту через точку простору вплив збурення зникає. Такий ефект називають **принципом Гюйгенса**.

![](images/_page_179_Picture_3.jpeg)

Фіг. 24.3. Точка починає реагувати лише тоді, коли сфера радіуса ct перетинає область початкового збурення.

У двовимірному випадку аналогічна задача розв'язується методом спуску, але там за хвильовим фронтом лишається «хвіст», тому принцип Гюйгенса вже не виконується.

## Конформні відображення

Тепер перейдемо до одного важливого класу функцій комплексної змінної, які широко використовують у теорії потенціалу.

### Визначення

Відображення w=f(z) називається **конформним** у точці z_0, якщо похідна f'(z_0)\neq 0. Якщо ця умова виконується в кожній точці області D, то відображення є конформним у всій області D.

Прості приклади:

- f(z)=z^2 є конформним всюди, крім точки z=0;
- f(z)=e^z є конформним у всій комплексній площині, бо f'(z)=e^z\neq 0.

Головна причина, чому ці відображення важливі в задачах математичної фізики, полягає в тому, що вони переносять складні області у простіші, не руйнуючи гармонічну структуру. Саме тому рівняння Лапласа в нових координатах часто набуває зручнішого вигляду.

### Приклад: задача Діріхле у верхній півплощині

Нехай потрібно розв’язати задачу:
\end{verbatim}

\texttt{\textbackslash{}varphi\_\{xx\}\ +\ \textbackslash{}varphi\_\{yy\}\ =\ 0,\ \textbackslash{}qquad\ -\textbackslash{}infty\ \textless{}\ x\ \textless{}\ \textbackslash{}infty,\textbackslash{}quad\ y\ \textgreater{}\ 0,}

\begin{verbatim}
із граничною умовою:
\end{verbatim}

\texttt{\textbackslash{}varphi(x,0)\ =}
\texttt{\textbackslash{}begin\{cases\}}
\texttt{1,\ \&\ \textbar{}x\textbar{}\ \textbackslash{}le\ 1,\ \textbackslash{}\textbackslash{}}
0, \& \textbar x\textbar{} \textgreater{} 1.
\texttt{\textbackslash{}end\{cases\}}

\begin{verbatim}
Для переходу до простішої області використовують комплексне відображення:
\end{verbatim}

\texttt{w\ =\ \textbackslash{}ln\textbackslash{}!\textbackslash{}left(\textbackslash{}frac\{z\ -\ 1\}\{z\ +\ 1\}\textbackslash{}right),\ \textbackslash{}qquad\ z\ =\ x\ +\ iy.}

\begin{verbatim}
Це відображення переводить верхню півплощину у смугу, для якої гранична задача має простіший вигляд.

Далі:

1. Розв’язують задачу у новій області (у змінних w),
2. Повертаються до початкових змінних x, y,
3. Отримують шуканий розв’язок `\varphi(x,y)`.

### Ідея методу

- складна область → простіша область (через конформне відображення),
- розв’язання у новій області,
- зворотне відображення до початкової області.

# Лекція 25. СКІНЧЕННІ ПЕРЕТВОРЕННЯ ФУР’Є (СИНУС- І КОСИНУС-ПЕРЕТВОРЕННЯ)

ЦЕЛЬ ЛЕКЦИИ: Ввести два новых интегральных преобразования (конечные спиус- и косинус-преобразования):
\end{verbatim}

\texttt{S\_n\ =\ S{[}f{]}\ =\ \textbackslash{}frac\{2\}\{L\}\ \textbackslash{}int\_0\^{}L\ f(x)\ \textbackslash{}sin(n\textbackslash{}pi\ x/L)\ dx}

\begin{verbatim}
(конечное синус-преобразование),
\end{verbatim}

\texttt{C\_n\ =\ C{[}f{]}\ =\ \textbackslash{}frac\{2\}\{L\}\ \textbackslash{}int\_0\^{}L\ f(x)\ \textbackslash{}cos(n\textbackslash{}pi\ x/L)\ dx}

\begin{verbatim}
(конечное косинус-преобразование),
\end{verbatim}

\texttt{f(x)\ =\ \textbackslash{}sum\_\{n=1\}\^{}\{\textbackslash{}infty\}\ S\_n\ \textbackslash{}sin(n\textbackslash{}pi\ x/L)}

\begin{verbatim}
(обратное синус-преобразование),
\end{verbatim}

\texttt{f(x)\ =\ \textbackslash{}frac\{C\_0\}\{2\}\ +\ \textbackslash{}sum\_\{n=1\}\^{}\{\textbackslash{}infty\}\ C\_n\ \textbackslash{}cos(n\textbackslash{}pi\ x/L)}

\begin{verbatim}
(обратное косинус-преобразование),

и показать, как с помощью этих преобразований решать краевые задачи (в частности, неоднородные). Ранее мы познакомились с регулярными преобразованиями Фурье и Лапласа и увидели, как применение этих преобразований к уравнению с частными производными сводит задачу к решению обыкновениого дифференциального уравнения. Обычное преобразование Фурье подразумевает, что преобразуемая переменная изменяется от  -\infty  до  +\infty , и, следовательно, оно применяется для решения задач в свободном пространстве (без границ). В этой лекции мы покажем, как решать краевые задачи (с границами), преобразуя ограниченные переменные (с таким преобразованием мы встречаемся впервые).

Давайте на время забудем о наших задачах, а займемся самими преобразованнями. Дадим определение преобразований, приведем формулы обращения и основные свойства. Для чего нужны преобразования и как ими пользоваться, мы узнаем несколько позже. Впрочем, если говорить кратко, то суть интегральных преобразований в разложении входящих в задачу функций по

частотам, решении семейства задач для каждой частоты и после-

дующем суммировании всех полученных результатов.

Начнем с функции f(x), определенной на интервале [0, L]. Конечные синус- и косинус-преобразования этой функции определяются формулами

Конечное синус-преобразование  S[f] = S_n = \frac{2}{L} \int_{\sigma}^{L} f(x) \sin(n\pi x/L) dx   n = 1, 2, \dots   C[f] = C_n = \frac{2}{L} \int_{\sigma}^{L} f(x) \cos(n\pi x/L) dx   n = 0, 1, \dots 

Конечное косинус-преобразование

Читатель, наверное, заметил, что эти преобразования не отличаются от формул для определения коэффициентов Фурье при разложении их в ряд по синусам или косинусам. Формулы обращения для этих преобразований—это обычные ряды Фурье по синусам и косинусам.

Обратное синус-преобразование
\end{verbatim}

\texttt{f(x)\ =\ \textbackslash{}sum\_\{n=1\}\^{}\{\textbackslash{}infty\}\ S\_n\ \textbackslash{}sin(n\textbackslash{}pi\ x/L)}
\texttt{f(x)\ =\ \textbackslash{}frac\{C\_0\}\{2\}\ +\ \textbackslash{}sum\_\{n=1\}\^{}\{\textbackslash{}infty\}\ C\_n\ \textbackslash{}cos(n\textbackslash{}pi\ x/L)}

\begin{verbatim}
Обратное косинус-преобразование

Отметим, что в обратном косинус-преобразовании суммирование начинается с n=0, а в синус-преобразовании — с n=1.

## Примеры синус-преобразования

Пусть
\end{verbatim}

\texttt{f(x)\ =\ 1,\ \textbackslash{}quad\ 0\ \textbackslash{}le\ x\ \textbackslash{}le\ 1,}

\begin{verbatim}
тогда
\end{verbatim}

\texttt{S\_n\ =\ S{[}1{]}\ =\ 2\ \textbackslash{}int\_0\^{}1\ \textbackslash{}sin(n\textbackslash{}pi\ x)\ dx\ =\ \textbackslash{}begin\{cases\}\ 0,\ \&\ \textbackslash{}text\{if\ \}\ n\ -\ \textbackslash{}text\{even\}\ \textbackslash{}\textbackslash{}\ \textbackslash{}frac\{4\}\{n\textbackslash{}pi\},\ \&\ \textbackslash{}text\{if\ \}\ n\ -\ \textbackslash{}text\{odd.\}\ \textbackslash{}end\{cases\}}

\begin{verbatim}
График функции f(x) и ее преобразования изображены на рис. 25.1. Обратное преобразование записывается в виде
\end{verbatim}

\texttt{f(x)\ =\ 1\ =\ \textbackslash{}frac\{4\}\{\textbackslash{}pi\}\ \textbackslash{}sum\_\{n=1\}\^{}\{\textbackslash{}infty\}\ \textbackslash{}left{[}\ \textbackslash{}frac\{1\}\{2n-1\}\ \textbackslash{}right{]}\ \textbackslash{}sin(n\textbackslash{}pi\ x).}

\begin{verbatim}
Знаете ли вы, как выглядит график этой функции ене отрезка [0, 1]? Подумайте над этим. Отметим также, что синус-преобразование функции f(x) определено только для положительных целых n. Другими словами, конечные синус- и косинус-преобразования преобразуют функции в числовые последовательности.

![](images/_page_185_Figure_3.jpeg)

РИС. 25.1. Графики функции f(x) = 1 и ее преобразования.

## Свойства преобразований

Прежде чем приступить к решению задач, мы должны получить некоторые свойства этих преобразований.

Если u(x, t) — функция  \partial \theta yx  переменных и мы проводим преобразование по переменной x, то
\end{verbatim}

\texttt{S{[}u{]}\ =\ S\_n(t)\ =\ \textbackslash{}frac\{2\}\{L\}\ \textbackslash{}int\_0\^{}L\ u(x,\ t)\ \textbackslash{}sin(\textbackslash{}pi\ nx/L)\ dx,}
\texttt{C{[}u{]}\ =\ C\_n(t)\ =\ \textbackslash{}frac\{2\}\{L\}\ \textbackslash{}int\_0\^{}L\ u(x,\ t)\ \textbackslash{}cos(\textbackslash{}pi\ nx/L)\ dx.}

\begin{verbatim}
(Заметим, что преобразование проводилось по переменной x, и полученная последовательность зависит только от времени t.)

А как быть с производными? Ниже приводятся несколько полезных формул
\end{verbatim}

\texttt{\textbackslash{}begin\{split\}\ S\textbackslash{}left{[}u\_\{t\}\textbackslash{}right{]}\ \&=\ \textbackslash{}frac\{dS\textbackslash{}left{[}u\textbackslash{}right{]}\}\{dt\},\ \textbackslash{}quad\ S\textbackslash{}left{[}u\_\{tt\}\textbackslash{}right{]}\ =\ \textbackslash{}frac\{d\^{}\{2\}S\textbackslash{}left{[}u\textbackslash{}right{]}\}\{dt\^{}\{2\}\},\ \textbackslash{}\textbackslash{}\ S\textbackslash{}left{[}u\_\{xx\}\textbackslash{}right{]}\ \&=\ -\textbackslash{}left{[}n\textbackslash{}pi/L\textbackslash{}right{]}\^{}\{2\}S\textbackslash{}left{[}u\textbackslash{}right{]}\ +\ \textbackslash{}frac\{2n\textbackslash{}pi\}\{L\^{}\{2\}\}\textbackslash{}left{[}u\textbackslash{}left(0,\textbackslash{}\ t\textbackslash{}right)\ +\ (-1)\^{}\{n+2\}u\textbackslash{}left(L,\textbackslash{}\ \ t\textbackslash{}right)\textbackslash{}right{]},\ \textbackslash{}\textbackslash{}\ C\textbackslash{}left{[}u\_\{xx\}\textbackslash{}right{]}\ \&=\ -\textbackslash{}left{[}n\textbackslash{}pi/L\textbackslash{}right{]}\^{}\{2\}C\textbackslash{}left{[}u\textbackslash{}right{]}\ -\ \textbackslash{}frac\{2\}\{L\}\textbackslash{}left{[}u\_\{x\}\textbackslash{}left(0,\textbackslash{}\ t\textbackslash{}right)\ +\ (-1)\^{}\{n+1\}u\_\{x\}\textbackslash{}left(L,\textbackslash{}\ t\textbackslash{}right)\textbackslash{}right{]}.\ \textbackslash{}end\{split\}}

\begin{verbatim}
Решение неоднородной краевой задачи с помощью синус-преобразования

Рассмотрим неоднородное волновое уравнение
\end{verbatim}

\texttt{(\textbackslash{}text\{UBP\})\ \textbackslash{}quad\ u\_\{tt\}\ =\ u\_\{xx\}\ +\ \textbackslash{}sin(\textbackslash{}pi\ x),\ \textbackslash{}quad\ 0\ \textless{}\ x\ \textless{}\ 1,\ \textbackslash{}quad\ 0\ \textless{}\ t\ \textless{}\ \textbackslash{}infty,}
\texttt{(\textbackslash{}text\{GU\})\ \textbackslash{}quad\ \textbackslash{}begin\{cases\}\ u(0,\ t)\ =\ 0,\ \textbackslash{}\textbackslash{}\ u(1,\ t)\ =\ 0,\ \textbackslash{}end\{cases\}\ \textbackslash{}quad\ 0\ \textless{}\ t\ \textless{}\ \textbackslash{}infty,}
\texttt{(\textbackslash{}text\{HV\})\ \textbackslash{}quad\ \textbackslash{}begin\{cases\}\ u(x,\ 0)\ =\ 1,\ \textbackslash{}\textbackslash{}\ u\_t(x,\ 0)\ =\ 0,\ \textbackslash{}end\{cases\}\ \textbackslash{}quad\ 0\ \textbackslash{}leqslant\ x\ \textbackslash{}leqslant\ 1.}

\begin{verbatim}
Решение задачи выполним за три следующих шага:

ШАГ І (Определение подходящего преобразования.)

Поскольку переменная x меняется от 0 до 1, мы воспользуемся конечным преобразованием, а именно синус-преобразованием; позже станет ясно почему. Мы могли бы решить эту задачу с помощью преобразования Лапласа по переменной t (по трудоемкости этот метод не отличается от конечного синус-преобразования).

ШАГ 2 (Выполнение преобразования).

Для удобства будем использовать обозначение  S_n(t) = S[u] . Применим синус-преобразование к исходному уравнению
\end{verbatim}

\texttt{S{[}u\_\{tt\}{]}\ =\ S{[}u\_\{xx\}{]}\ +\ S{[}\textbackslash{}sin(\textbackslash{}pi\ x){]}}

\begin{verbatim}
и воспользуемся тождеством для синус-преобразования, получим
\end{verbatim}

\texttt{\textbackslash{}begin\{split\}\ \textbackslash{}frac\{d\^{}2S\_n\textbackslash{}left(t\textbackslash{}right)\}\{dt\^{}2\}\ \&=\ -\textbackslash{}left(\textbackslash{}pi\ n\textbackslash{}right)\^{}2S\_n\textbackslash{}left(t\textbackslash{}right)\ +\ 2n\textbackslash{}pi\textbackslash{}left{[}u\textbackslash{}left(0,\textbackslash{}\ t\textbackslash{}right)\ +\ \textbackslash{}left(-1\textbackslash{}right)\^{}\{n+1\}u\textbackslash{}left(1,\textbackslash{}\ t\textbackslash{}right)\textbackslash{}right{]}\ +\ D\_n\textbackslash{}left(t\textbackslash{}right)\ =\ \textbackslash{}\textbackslash{}\ \&=\ -\textbackslash{}left(n\textbackslash{}pi\textbackslash{}right)\^{}2S\_n\textbackslash{}left(t\textbackslash{}right)\ +\ D\_n\textbackslash{}left(t\textbackslash{}right),\ \ \textbackslash{}end\{split\}}

\begin{verbatim}
где
\end{verbatim}

\texttt{D\_n(t)\ =\ S{[}\textbackslash{}sin(\textbackslash{}pi\ x){]}\ =\ \textbackslash{}begin\{cases\}\ 1,\ \&\ n\ =\ 1\ \textbackslash{}\ \textbackslash{}text\{(це\ коефіцієнти\ розкладу\}\ \textbackslash{}\textbackslash{}0,\ \&n\ =\ 2,\ 3,\ \textbackslash{}dots\ \textbackslash{}end\{cases\}}

\begin{verbatim}
Если теперь преобразовать начальные условия краевой задачи, то получим начальные условия для обыкновенного дифференциального уравнения
\end{verbatim}

\texttt{S{[}u(x,\ 0){]}\ =\ S\_n(0)\ =\ \textbackslash{}begin\{cases\}\ 4/n\textbackslash{}pi,\ \&\ n\ =\ 1,\ 3,\ \textbackslash{}dots,\ \textbackslash{}\textbackslash{}\ 0,\ \&n\ =\ 2,\ 4,\ \textbackslash{}dots,\ \textbackslash{}end\{cases\}}
\texttt{S{[}u\_t(x,\ 0){]}\ =\ \textbackslash{}frac\{dS\_n(0)\}\{dt\}\ =\ 0.}

\begin{verbatim}
Итак, решим теперь новое семейство задач Кощи:

(ОДУ)
\end{verbatim}

\texttt{\textbackslash{}frac\{d\^{}2S\_n\}\{dt\^{}2\}\ +\ (n\textbackslash{}pi)\^{}2\ S\_n\ =\ \textbackslash{}begin\{cases\}\ 1,\ \&n\ =\ 1,\ \textbackslash{}\textbackslash{}\ 0,\ \&n\ =\ 2,\ 3,\ \textbackslash{}dots,\ \textbackslash{}end\{cases\}}
\texttt{S\_n(0)\ =\ \textbackslash{}begin\{cases\}}
\texttt{4/n\textbackslash{}pi,\ \&\ n\ =\ 1,\ 3,\ \textbackslash{}dots,\ \textbackslash{}\textbackslash{}}
\texttt{0,\ \&\ n\ =\ 2,\ 4,\ \textbackslash{}dots,}
\texttt{\textbackslash{}end\{cases\}}
\texttt{\textbackslash{}frac\{dS\_n(0)\}\{dt\}\ =\ 0.}

\begin{verbatim}
В результате получаем
\end{verbatim}

\texttt{S\_1(t)\ =\ A\textbackslash{}cos(\textbackslash{}pi\ t)\ +\ (1/\textbackslash{}pi)\^{}2}

,

где

\begin{verbatim}
A = \frac{4}{\pi} - \frac{1}{\pi^2} = 1,17,
S_n(t) = \begin{cases} 0, &n = 2, 4, \dots, \\ (4/n\pi)\cos(n\pi t), &n = 3, 5, 7, \dots \end{cases}
\end{verbatim}

Следовательно, решение нашей задачи можно представить в виде

\begin{verbatim}
u(x, t) = [A\cos(\pi t) + (1/\pi)^2] \sin(\pi x) + \frac{4}{\pi} \sum_{n=1}^{\infty} \frac{1}{2n+1} \cos[(2n+1)\pi t] \sin[(2n+1)\pi x].
\end{verbatim}

\paragraph{ЗАМЕЧ
АНИЯ}\label{ux437ux430ux43cux435ux447-ux430ux43dux438ux44f}

\begin{enumerate}
\def\labelenumi{\arabic{enumi}.}
\tightlist
\item
  Для решения задачи методом конечного синус- или косинуспреобразования
  краевые условия при x=0 и x=L должны иметь вид
\end{enumerate}

\begin{verbatim}
\begin{cases} u(0, t) = f(t) \\ u(L, t) = g(t) \end{cases}
\end{verbatim}

(используется синус-преобразование),

\begin{verbatim}
\begin{cases} u_x(0, t) = f(t) \\ u_x(L, t) = g(x) \end{cases}
\end{verbatim}

(используется косинус-преобразование).

Другими словами, если заданы граничные условия \texttt{u(0,\ t)} =
\texttt{f(t)} и \texttt{u\_x(L,\ t)\ =\ g(t)} , то эти преобразования
неприменимы. Неприменимы они и при граничных условиях вида
\texttt{u\_x(0,\ t)\ +\ hu(0,\ t)\ =\ 0} . В этих случаях используются
другие преобразования. С обобщенными синус- и косинус-преобразованиями
можно познакомиться по литературе, указанной в конце книги.

\begin{enumerate}
\def\labelenumi{\arabic{enumi}.}
\setcounter{enumi}{1}
\tightlist
\item
  Для того чтобы можно было применять конечные синус- или
  косинус-преобразования, исходные уравнения не должны содержать первых
  производных по x (поскольку синус-преобра-
\end{enumerate}

нию второй и наоборот).

\begin{enumerate}
\def\labelenumi{\arabic{enumi}.}
\setcounter{enumi}{2}
\tightlist
\item
  В методах конечных синус- и косинус-преобразований существенно то, что
  все функции задачи, включая \texttt{u\_\{ti\}} , \texttt{u\_\{xx\}} ,
  начальные условия, граничные условия, разлагаются в ряд Фурье по
  синусам или косинусам, решается последовательность ОДУ для
  коэффициентов Фурье, а затем все результаты суммируются.
\end{enumerate}

\subsection{ЗАДАЧИ}\label{ux437ux430ux434ux430ux447ux438}

\begin{enumerate}
\def\labelenumi{\arabic{enumi}.}
\tightlist
\item
  Решите диффузионную задачу с изолированными границами, т. е.
\end{enumerate}

\begin{verbatim}
\begin{array}{ll} (\mathrm{Y} \mathrm{U} \Pi) & u_t = u_{xx}, \quad 0 < x < 1, \quad 0 < t < \infty, \\ (\mathrm{\Gamma} \mathrm{Y} & \left\{ \begin{array}{ll u_x} �  \left( 0, \ t \right) = 0 \\ u_x \left( 1, \ t \right) = 0, \end{array} \right. & 0 < t < \infty, \\ (\mathrm{H} \mathrm{Y}) & u \left( x, \ 0 \right) = 1 + \cos \left( \pi x \right) + 0.. 5 \cos \left( 3\pi x \right),  \quad 0 \leqslant x \leqslant 1. \end{array}
\end{verbatim}

\begin{enumerate}
\def\labelenumi{\arabic{enumi}.}
\setcounter{enumi}{1}
\tightlist
\item
  Решите в общем виде задачу
\end{enumerate}

\begin{verbatim}
\begin{array}{ll} (\text{PHP}) & u_t = \alpha^2 u_{xx} + bu + f(x, t), \quad 0 < x < 1, \quad 0 < t < \infty, \\ (\text{GU}) & \left\{ \begin{array}{ll} u(0, t) = 0, \\ u(1, t) = 0, \end{array} \right. & 0 < t < \infty, \\ (\text{NU}) & u(x, 0) = 0, \quad 0 \leqslant x \leqslant 1. \end{array}
\end{verbatim}

\begin{enumerate}
\def\labelenumi{\arabic{enumi}.}
\setcounter{enumi}{2}
\tightlist
\item
  Выведите основные правила вычисления \texttt{S{[}u\_\{xx\}{]}} и
  \texttt{C{[}u\_\{xx\}{]}} , которые были приведены в замечаниях.
  Понимаете ли вы, почему трудно решить дифференциальное уравнение,
  содержащее первую производную?
\end{enumerate}

\texttt{**4.**\ Чему\ равно\ конечное\ синус-преобразование\ функции\ \ \textquotesingle{}f(x)\ =\ \textbackslash{}sin\{(\textbackslash{}pi\ x)\}\ +\ 0.5\ (3\textbackslash{}pi\ x)\textquotesingle{}\ ?\ Постройте\ график\ преобразования\ при}

\begin{verbatim}
L=1.
\end{verbatim}

\textbf{5.} Найдите косинус-преобразование функции \texttt{f(x)} = x,
\texttt{0\ \textbackslash{}le\ x\ \textbackslash{}le\ 1} . Изобразите
график обратного преобразования при всех значениях x. (Конечно, внутри
отрезка \texttt{0\ \textbackslash{}le\ x\ \textbackslash{}le\ 1}
восстановится функция \texttt{f(x)} = x, а что будет вне отрезка {[}0,
1{]}?)

\begin{enumerate}
\def\labelenumi{\arabic{enumi}.}
\setcounter{enumi}{5}
\tightlist
\item
  Решите задачу
\end{enumerate}

\begin{verbatim}
\begin{array}{llll} (\text{PHP}) & u_t = u_{xx} + \sin{(3\pi x)}, & 0 < x < 1, & 0 < t < \infty, \\ (\text{GU}) & \begin{cases} u(0, \ t) = 0, \\ u(1, \ t) = 1, \end{cases} & 0 < t < \infty, \\ (\text{NU}) & u(x, \ 0) = \sin{(\pi x)}, & 0 \leqslant x \leqslant 1. \end{array}
\end{verbatim}

\section{Лекція 26. ПРИНЦИП СУПЕРПОЗИЦІЇ --- ОСНОВА ТЕОРІЇ ЛІНІЙНИХ
СИСТЕМ}\label{ux43bux435ux43aux446ux456ux44f-26.-ux43fux440ux438ux43dux446ux438ux43f-ux441ux443ux43fux435ux440ux43fux43eux437ux438ux446ux456ux457-ux43eux441ux43dux43eux432ux430-ux442ux435ux43eux440ux456ux457-ux43bux456ux43dux456ux439ux43dux438ux445-ux441ux438ux441ux442ux435ux43c}

ЦЕЛЬ ЛЕКЦИИ: Познакомить читателя с принципом суперпозиции и показать,
как с его помощью исходную задачу можно разложить на подзадачи, затем
решить сразу все подзадачи и, сложив полученные результаты, найти
решение исходной задачи. Мы хотим также показать, что два основных
метода решения линейных уравнений---метод разделения переменных и метод
интегральных преобразований---используют принцип суперпозиции.

Предположим, что инженеру нужно найти отклик линейной системы и на
входное воздействие f.~Общий подход к решению задачи таков:

\begin{itemize}
\tightlist
\item
  \begin{enumerate}
  \def\labelenumi{\arabic{enumi})}
  \tightlist
  \item
    разложить f на элементарные воздействия
    \texttt{f\ =\ \textbackslash{}sum\ f\_k} ;
  \end{enumerate}
\item
  \begin{enumerate}
  \def\labelenumi{\arabic{enumi})}
  \setcounter{enumi}{1}
  \tightlist
  \item
    найти отклик системы \texttt{u\_k} на воздействие \texttt{f\_k} ;
  \end{enumerate}
\end{itemize}

\begin{enumerate}
\def\labelenumi{\arabic{enumi})}
\setcounter{enumi}{2}
\tightlist
\item
  сложить отклики
  \texttt{u\ =\ \textbackslash{}sum\_\{i=1\}\^{}\{n\}\ u\_\{i\}} .
\end{enumerate}

Если система линейная, то полученная сумма и является откликом системы
на входное воздействие f.~В этом состоит принцип суперпоэиции (см. рнс.
26.1).

\pandocbounded{\includegraphics[keepaspectratio]{images/_page_189_Picture_10.jpeg}}

РИС, 26.1. Основная идея принципа суперпоэнции; L---дифференциальный
оператор; f---вход; u---выход.

Эту основную идею можно использовать при решении сложных смешанных
задач, разбивая их на простые подзадачи, ре-

шая каждую из них и затем суммируя полученные результаты (конечно,
дифференциальное уравнение и граничные условия должны быть линейны).

\subsection{Разложение смешанной задачи на две более
простые}\label{ux440ux430ux437ux43bux43eux436ux435ux43dux438ux435-ux441ux43cux435ux448ux430ux43dux43dux43eux439-ux437ux430ux434ux430ux447ux438-ux43dux430-ux434ux432ux435-ux431ux43eux43bux435ux435-ux43fux440ux43eux441ux442ux44bux435}

Предположим, нам необходимо решить линейную задачу (обозначим ее P)

(УЧП)

\begin{verbatim}
u_t = u_{xx} + \sin(\pi x), \quad 0 < x < 1, \quad 0 < t < \infty,
\end{verbatim}

(Р) (ГУ)

\begin{verbatim}
\begin{cases} u(0, t) = 0, \\ u(1, t) = 0, \\ u(x, 0) = \sin(2\pi x), \quad 0 \le x \le \le 1. \end{cases}
\end{verbatim}

Перед нами неоднородное уравнение теплопроводности, так что метод
разделения переменных неприменим. Можпо, конечно, воспользоваться
конечным синус-преобразованием по переменной x или преобразованием
Лапласа по переменной t, но вместо этого мы рассмотрим две подзадачи

\section{Лекція 27. РІВНЯННЯ ПЕРШОГО ПОРЯДКУ (МЕТОД
ХАРАКТЕРИСТИК)}\label{ux43bux435ux43aux446ux456ux44f-27.-ux440ux456ux432ux43dux44fux43dux43dux44f-ux43fux435ux440ux448ux43eux433ux43e-ux43fux43eux440ux44fux434ux43aux443-ux43cux435ux442ux43eux434-ux445ux430ux440ux430ux43aux442ux435ux440ux438ux441ux442ux438ux43a}

МЕТА ЛЕКЦІЇ: Ввести поняття диференціального рівняння в частинних
похідних першого порядку і ознайомитися з методом характеристик для
задач із початковими умовами.

Розглянемо типову задачу Коші для лінійного рівняння першого порядку:

\begin{itemize}
\tightlist
\item
  (УЧП): \texttt{a(x,t)u\_x\ +\ b(x,t)u\_t\ =\ c(x,t)},
  \texttt{-\textbackslash{}infty\ \textless{}\ x\ \textless{}\ \textbackslash{}infty},
  \texttt{0\ \textless{}\ t\ \textless{}\ \textbackslash{}infty}.
\item
  (НУ): \texttt{u(x,0)\ =\ \textbackslash{}varphi(x)},
  \texttt{-\textbackslash{}infty\ \textless{}\ x\ \textless{}\ \textbackslash{}infty}.
\end{itemize}

На відміну від рівняння теплопровідності, де початкове збурення миттєво
впливає на весь простір, для рівнянь першого порядку інформація
поширюється вздовж спеціальних кривих у площині \texttt{(x,t)}. Такі
криві називаються характеристиками.

Ідея методу характеристик така: замість того, щоб розв'язувати рівняння
одразу в координатах \texttt{(x,t)}, ми переходимо до параметра
\texttt{s}, який біжить уздовж характеристики. Уздовж правильно вибраної
характеристики частинні похідні поєднуються так, що задача зводиться до
звичайного диференціального рівняння.

Якщо характеристику записати у вигляді \texttt{x=x(s)}, \texttt{t=t(s)},
то природно вимагати, щоб вона задовольняла співвідношення:

\texttt{-\ \textquotesingle{}\textbackslash{}dfrac\{dx\}\{ds\}\ =\ a(x,t)\textquotesingle{},}
\texttt{-\ \textquotesingle{}\textbackslash{}dfrac\{dt\}\{ds\}\ =\ b(x,t)\textquotesingle{}.}

Після цього значення \texttt{u(x,t)} уздовж характеристики визначається
ще одним звичайним диференціальним рівнянням. Саме тому метод
характеристик перетворює складну задачу для УЧП на послідовність
простіших задач для звичайних рівнянь.

Геометрично це означає, що початкові дані з лінії \texttt{t=0} не
розмазуються по всій площині, а переносяться вздовж характеристичних
кривих. У наступному фрагменті лекції ці криві буде побудовано явно, а
тоді з їхньою допомогою буде отримано розв'язок задачі Коші.

\section{Лекція 28. НЕЛІНІЙНІ РІВНЯННЯ ПЕРШОГО ПОРЯДКУ (ЗАКОНИ
ЗБЕРЕЖЕННЯ)}\label{ux43bux435ux43aux446ux456ux44f-28.-ux43dux435ux43bux456ux43dux456ux439ux43dux456-ux440ux456ux432ux43dux44fux43dux43dux44f-ux43fux435ux440ux448ux43eux433ux43e-ux43fux43eux440ux44fux434ux43aux443-ux437ux430ux43aux43eux43dux438-ux437ux431ux435ux440ux435ux436ux435ux43dux43dux44f}

МЕТА ЛЕКЦІЇ: Показати, як закони збереження приводять до нелінійних
рівнянь першого порядку і як метод характеристик описує появу ударних
хвиль у задачі дорожнього руху.

\pandocbounded{\includegraphics[keepaspectratio]{images/_page_205_Figure_5.jpeg}}

Рис. 28.1. Типова залежність потоку від густини.

Тепер ми готові застосувати закон збереження до задачі дорожнього руху.
Якщо потік \texttt{f(u)} залежить від густини \texttt{u}, то

\begin{verbatim}
f_x=\frac{df}{du}\,u_x.
\end{verbatim}

Отже, рівняння збереження

\begin{verbatim}
u_t + f_x = 0
\end{verbatim}

можна переписати у вигляді

\begin{verbatim}
u_t + \frac{df}{du}u_x = u_t + g(u)u_x = 0,
\end{verbatim}

де \texttt{g(u)=f\textquotesingle{}(u)}.

Наприклад, якщо залежність потоку від густини квадратична, тобто
\texttt{f(u)=u\^{}2}, то закон збереження набуває форми

\begin{verbatim}
u_t + 2u\,u_x = 0.
\end{verbatim}

Якщо початкова густина розподілу автомобілів задається функцією
\texttt{u(x,0)=\textbackslash{}varphi(x)}, то потрібно розв'язати задачу
Коші

\begin{verbatim}
\begin{aligned}
(\text{УЧП})\quad & u_t + 2u\,u_x = 0, && -\infty < x < \infty,\ 0 < t < \infty, \\
(\text{НУ})\quad & u(x,0)=\varphi(x), && -\infty < x < \infty.
\end{aligned}
\end{verbatim}

Тепер розглянемо загальну нелінійну задачу Коші.

\section{Нелінійна задача з початковою умовою (задача
Коші)}\label{ux43dux435ux43bux456ux43dux456ux439ux43dux430-ux437ux430ux434ux430ux447ux430-ux437-ux43fux43eux447ux430ux442ux43aux43eux432ux43eux44e-ux443ux43cux43eux432ux43eux44e-ux437ux430ux434ux430ux447ux430-ux43aux43eux448ux456}

Розглянемо задачу

\begin{verbatim}
\begin{aligned}
(\text{УЧП})\quad & u_t + g(u)u_x = 0, && -\infty < x < \infty,\ 0 < t < \infty, \\
(\text{НУ})\quad & u(x,0)=\varphi(x), && -\infty < x < \infty.
\end{aligned}
\end{verbatim}

Нагадаємо, що для рівняння переносу

\begin{verbatim}
u_t + v u_x = 0
\end{verbatim}

розв'язок є сталим уздовж прямих \texttt{x=x\_0+vt}. У нелінійному
випадку відбувається те саме, але швидкість тепер залежить від самого
розв'язку. Якщо частинка середовища стартує з точки \texttt{(x\_0,0)} і
рухається зі швидкістю \texttt{g(u)}, то її траєкторія має вигляд

\begin{verbatim}
x = x_0 + g(u)\,t.
\end{verbatim}

Це і є рівняння характеристики. Вздовж характеристики значення
\texttt{u(x,t)} залишається сталим і дорівнює початковому значенню
\texttt{u(x\_0,0)}.

\pandocbounded{\includegraphics[keepaspectratio]{images/_page_206_Figure_8.jpeg}}

Фіг. 28.2. Характеристики рівняння \texttt{u\_t\ +\ g(u)u\_x\ =\ 0};
вздовж кожної характеристики розв'язок залишається сталим і дорівнює
значенню в точці \texttt{(x\_0,0)}, з якої ця характеристика виходить.

Отже, характеристика, що починається в точці \texttt{(x\_0,0)},
задається формулою

\begin{verbatim}
x = x_0 + g\!\bigl[u(x_0,0)\bigr]\,t.
\end{verbatim}

Наприклад, для рівняння

\begin{verbatim}
u_t + 3u\,u_x = 0
\end{verbatim}

якщо \texttt{u(2,0)=1}, то характеристика, що виходить із точки
\texttt{(2,0)}, визначається співвідношенням

\begin{verbatim}
x = 2 + g[1]\,t = 2 + 3t.
\end{verbatim}

Отже, на прямій \texttt{x=2+3t} розв'язок зберігає значення
\texttt{u=1}. Саме так читається зміст характеристик у нелінійній задачі
Коші.

\texttt{Тепер\ зрозуміло,\ що,\ знаючи\ характеристики\ в\ кожній\ точці\ і\ знаючи,\ що\ розв\textquotesingle{}язок\ за\ характеристиками\ не\ змінюється,\ можна\ отримати\ розв\textquotesingle{}язок\ \textquotesingle{}u\textbackslash{}left(x,\textbackslash{}\ t\textbackslash{}right)\textquotesingle{}\ у\ будь-який\ момент\ часу\ t.\ Ми\ не\ отримали\ явного\ виразу\ для\ \textquotesingle{}u\textbackslash{}left(x,\textbackslash{}\ t\textbackslash{}right)\textquotesingle{}\ через\ змінні\ x\ і\ t,\ але\ знання\ характеристик\ рівняння\ дозволяє\ розв\textquotesingle{}язати\ деякі\ цікаві\ задачі.}

\section{Задача дорожнього
руху}\label{ux437ux430ux434ux430ux447ux430-ux434ux43eux440ux43eux436ux43dux44cux43eux433ux43e-ux440ux443ux445ux443}

Розглянемо модель руху, у якій потік задається формулою
\texttt{f(u)=u\^{}2}, а початкова густина має вигляд, показаний на рис.
28.3.

\pandocbounded{\includegraphics[keepaspectratio]{images/_page_207_Figure_8.jpeg}}

. 28.3. Початкова густина автомобілів, що рухаються зліва направо.

Іншими словами, нам потрібно розв'язати задачу

\begin{verbatim}
\begin{aligned}
(\text{УЧП})\quad & u_t + 2u\,u_x = 0, && -\infty < x < \infty,\ 0 < t < \infty, \\
(\text{НУ})\quad &
u(x,0)=
\begin{cases}
1, & x \le 0,\\
1-x, & 0 < x < 1,\\
0, & 1 \le x,
\end{cases}
&& -\infty < x < \infty.
\end{aligned}
\end{verbatim}

Почнемо з визначення характеристик, що виходять із точок
\texttt{(x\_0,0)}.

Якщо \texttt{x\_0\textless{}0}, то \texttt{u(x\_0,0)=1}, отже

\begin{verbatim}
x = x_0 + g[1]\,t = x_0 + 2t,
\end{verbatim}

або, що те саме,

\begin{verbatim}
t = \frac{x-x_0}{2}.
\end{verbatim}

Тепер розглянемо точки \texttt{0\textless{}x\_0\textless{}1}. Тут
\texttt{u(x\_0,0)=1-x\_0}, тому

\begin{verbatim}
x = x_0 + g(1-x_0)t = x_0 + 2(1-x_0)t,
\end{verbatim}

звідки

\begin{verbatim}
t = \frac{x-x_0}{2(1-x_0)}.
\end{verbatim}

\texttt{Нарешті,\ якщо\ \textquotesingle{}1\textbackslash{}le\ x\_0\textless{}\textbackslash{}infty\textquotesingle{},\ то\ \textquotesingle{}u(x\_0,0)=0\textquotesingle{},\ і\ характеристики\ є\ вертикальними:}

\begin{verbatim}
x = x_0.
\end{verbatim}

Усі ці характеристики показані на рис. 28.4.

\pandocbounded{\includegraphics[keepaspectratio]{images/_page_208_Figure_10.jpeg}}

Фіг. 28.4. Характеристики рівняння \texttt{u\_t\ +\ 2uu\_x\ =\ 0} .

Ситуація на дорозі при \texttt{0\textless{}t\textless{}1/2} досить
зрозуміла. Рис. 28.5 показує розв'язок у різні моменти часу. Зверніть
увагу, що характеристики сходяться в одній точці при \texttt{t=1/2}.
Отже, для \texttt{t\textgreater{}1/2} виникає ударна хвиля, тобто розрив
розв'язку. Потрібно визначити швидкість руху її фронту.

Виявляється, що швидкість поширення розриву задається формулою
\texttt{S=\textbackslash{}dfrac\{f(u\_R)-f(u\_L)\}\{u\_R-u\_L\}}, де
\texttt{u\_R} і \texttt{u\_L} --- значення розв'язку праворуч і ліворуч
від фронту хвилі, а \texttt{f(u\_R)} і \texttt{f(u\_L)} --- значення
потоку при цих густинах.

\pandocbounded{\includegraphics[keepaspectratio]{images/_page_209_Figure_5.jpeg}}

Фіг. 28.5. Щільність руху в різний час.

У нашому прикладі швидкість поширення щілини виявляється рівною одиниці.
Справді, для \texttt{f(u)=u\^{}2} маємо

\begin{verbatim}
S=\frac{0-1}{0-1}=1.
\end{verbatim}

Це означає, що при \texttt{t\textgreater{}1/2} фронт хвилі рухатиметься
зліва направо зі швидкістю одиниці. Повне розв'язання цієї задачі
показано на рис. 28.5.

\section{ЗАУВАЖЕННЯ}\label{ux437ux430ux443ux432ux430ux436ux435ux43dux43dux44f-6}

\begin{enumerate}
\def\labelenumi{\arabic{enumi}.}
\item
  Ударна хвиля в нашому прикладі виникає через те, що потік швидко
  зростає зі збільшенням густини \texttt{u}. Якщо б потік задавався
  формулою \texttt{f(u)=u}, то рівняння мало б вигляд
  \texttt{u\_t+u\_x=0}, а розв'язок був би просто хвилею, що рухається
  вправо без спотворення форми.
\item
  Прямою підстановкою розв'язку вигляду
\end{enumerate}

\begin{verbatim}
u=\varphi(x-g(u)t)
\end{verbatim}

можна перевірити, що ця формула неявно задає розв'язок задачі

\begin{verbatim}
\begin{aligned}
(\text{РЧП})\quad & u_t + g(u)u_x = 0, && -\infty < x < \infty,\quad 0 < t < \infty,\\
(\text{ПУ})\quad & u(x,0)=\varphi(x), && -\infty < x < \infty.
\end{aligned}
\end{verbatim}

Наприклад, для задачі Коші

\begin{verbatim}
\begin{aligned}
(\text{РЧП})\quad & u_t + u\,u_x = 0, && -\infty < x < \infty,\quad 0 < t < \infty,\\
(\text{ПУ})\quad & u(x,0)=x, && -\infty < x < \infty,
\end{aligned}
\end{verbatim}

неявний розв'язок задається співвідношенням

\begin{verbatim}
u=\varphi[x-g(u)t]=x-ut.
\end{verbatim}

У нашому конкретному випадку можливо отримати рішення у явній формі (у
багатьох інших випадках це неможливо):

\begin{verbatim}
u(x,t)=\frac{x}{1+t}.
\end{verbatim}

Переконайтеся, що ця функція справді є розв'язком нашої задачі.

\section{ЗАВДАННЯ}\label{ux437ux430ux432ux434ux430ux43dux43dux44f-6}

\begin{enumerate}
\def\labelenumi{\arabic{enumi}.}
\tightlist
\item
  Розв'яжіть задачу Коші
\end{enumerate}

\begin{verbatim}
\begin{aligned}
(\text{УЧП})\quad & u_t + u\,u_x = 0, && -\infty < x < \infty,\quad 0 < t < \infty, \\
(\text{НУ})\quad &
u(x,0)=
\begin{cases}
0, & x < 0,\\
x, & 0 \le x.
\end{cases}
\end{aligned}
\end{verbatim}

Побудуйте графіки розв'язку для різних моментів часу і поясніть їх
фізичний зміст.

\begin{enumerate}
\def\labelenumi{\arabic{enumi}.}
\setcounter{enumi}{1}
\tightlist
\item
  Розв'яжіть задачу
\end{enumerate}

\begin{verbatim}
\begin{aligned}
(\text{УЧП})\quad & u_t + u^2 u_x = 0, && -\infty < x < \infty,\quad 0 < t < \infty, \\
(\text{НУ})\quad &
u(x,0)=
\begin{cases}
0, & x < 0,\\
x, & 0 \le x.
\end{cases}
\end{aligned}
\end{verbatim}

Порівняйте поведінку цього розв'язку з результатом задачі 1.

\begin{enumerate}
\def\labelenumi{\arabic{enumi}.}
\setcounter{enumi}{2}
\tightlist
\item
  Припустимо, що нев'язка рідина тече по трубі і водночас просочується
  крізь стінки згідно із законом \texttt{F(u)=ku}. Тоді закон збереження
  набуває вигляду
\end{enumerate}

\begin{verbatim}
u_t + f_x = -F(u).
\end{verbatim}

Знайдіть розв'язок, якщо \texttt{f(u)=u}, а початковий розподіл
задається довільною функцією
\texttt{\textquotesingle{}u(x,0)=\textbackslash{}varphi(x)\textquotesingle{}.}

\begin{enumerate}
\def\labelenumi{\arabic{enumi}.}
\setcounter{enumi}{3}
\item
  Розв'яжіть попередню задачу, якщо втрати задаються функцією
  \texttt{F(x,t)=1/x}.
\item
  Перевірте, що неявно задана функція
\end{enumerate}

\begin{verbatim}
u = \varphi[x-g(u)t]
\end{verbatim}

є розв'язком нелінійної задачі

\begin{verbatim}
\begin{aligned}
(\text{УЧП})\quad & u_t + g(u)u_x = 0,\\
(\text{НУ})\quad & u(x,0)=\varphi(x).
\end{aligned}
\end{verbatim}

\section{Лекція 29. СИСТЕМИ РІВНЯНЬ З ЧАСТКОВИМИ
ПОХІДНИМИ}\label{ux43bux435ux43aux446ux456ux44f-29.-ux441ux438ux441ux442ux435ux43cux438-ux440ux456ux432ux43dux44fux43dux44c-ux437-ux447ux430ux441ux442ux43aux43eux432ux438ux43cux438-ux43fux43eux445ux456ux434ux43dux438ux43cux438}

ЦЕЛЬ ЛЕКЦИИ: Познакомить читателя с тем фактом, что многие физические
системы (особенно в динамике жидкости) нельзя описать одним уравнением.
Для описания таких систем требуются системы связанных уравнений. В эти
системы входят некоторые неизвестные функции, такие, как давление
\texttt{p(x,\ y,\ z,\ t)}, плотность
\texttt{\textbackslash{}rho(x,\ y,\ z,\ t)} , температура
\texttt{T(x,\ y,\ z,\ t)} и их частные производные. Связи между этими
функциями определяются физическими законами и наша задача --- найти все
эти функции одновременно.

В лекции будет показано, как можно решить линейную систему уравнений

\begin{verbatim}
u_t + Au_x = 0
\end{verbatim}

если преобразовать ее в новую систему независимых уравнений

\begin{verbatim}
v_1 + \Lambda v_x = 0,
\end{verbatim}

а затем независимо решить каждое уравнение этой системы.

Во многих областях науки встречаются системы величин, которые связаны
между собой не одним, а сразу несколькими соотношениями. Например, в
динамике жидкостей---это четыре уравнения

\begin{verbatim}
u_t + uu_x + vu_y + \frac{1}{\rho} p_x = 0
\end{verbatim}

\texttt{(сохранение\ импульса\ вдоль\ оси\ \ \textquotesingle{}x\textquotesingle{}\ ),\ (29.1)\ \ \textquotesingle{}v\_t\ +\ uv\_x\ +\ vv\_y\ +\ \textbackslash{}frac\{1\}\{\textbackslash{}rho\}\ p\_y\ =\ 0\textquotesingle{}\ \ (сохранение\ импульса\ вдоль\ оси\ \ \textquotesingle{}y\textquotesingle{}\ ),\ (сохранение\ массы),\ (сохранение\ массы),\ (сохранение\ массы),\ \ \textquotesingle{}\textbackslash{}left(\textbackslash{}frac\{p\}\{\textbackslash{}rho\^{}\{\textbackslash{}gamma\}\}\textbackslash{}right)\_t\ +\ u\textbackslash{}left(\textbackslash{}frac\{p\}\{\textbackslash{}rho\^{}\{\textbackslash{}gamma\}\}\textbackslash{}right)\_x\ +\ v\textbackslash{}left(\textbackslash{}frac\{p\}\{\textbackslash{}rho\^{}\{\textbackslash{}gamma\}\}\textbackslash{}right)\_y\ =\ 0\textquotesingle{}\ \ (сохранение\ энергии),\ \ \textquotesingle{}p(x,\ y,\ z,\ t)\textquotesingle{}\ \ —\ давление\ в\ жидкости,\ \ \textquotesingle{}u(x,\ y,\ z,\ t)\textquotesingle{}\ \ —\ \ \textquotesingle{}x\textquotesingle{}\ -компонента\ скорости,\ \ \textquotesingle{}v(x,\ y,\ z,\ t)\textquotesingle{}\ \ —\ \ \textquotesingle{}y\textquotesingle{}\ -компонента\ скорости,\ \ \textquotesingle{}\textbackslash{}rho(x,\ y,\ z,\ t)\textquotesingle{}\ \ —\ плотность\ жидкости.}

Приведенная нелинейная система уравнений называется уравнениями Эйлера
движения жидкости. Задача состоит в одновременном отыскании всех
неизвестных функций р, и, v и р, которые удовлетворяли бы всем четырем
уравнениям (вместе с начальными и граничными условиями).

Есть и другие причины для изучения систем уравнений. В теории
обыкновенных дифференциальных уравнений (если читатель помнит) показано,
что одно уравнение второго порядка можно представить в виде системы двух
уравнений первого порядка. И хотя в теории уравнений с частными
производными все не так просто, очень часто одно уравнение с частными
производными высокого порядка можно свести к системе уравнений с
частными производными первого порядка. Например, телеграфное уравнение

\begin{verbatim}
u_{tt} = c^2 u_{xx} + a u_t + b u
\end{verbatim}

путем введения трех новых переменных

\begin{verbatim}
u_1 = u_1
u_2 = u_2
u_3 = u_4
\end{verbatim}

можно свести к системе уравнений первого порядка

(29.2)

\begin{verbatim}
\begin{cases} \frac{\partial u_1}{\partial x} = u_2, \\ \frac{\partial a_1}{\partial t} = u_3, \\ \frac{\partial u_3}{\partial t} = c^2 \frac{\partial u_2}{\partial x} + au_3 + bu_4. \end{cases}
\end{verbatim}

Для того чтобы решать системы типа (29.2), необходимо хорошо знать такие
понятия, как собственные значения и собственные векторы матрицы. По эфой
причине мы предлагаем вниманию читателя краткий обзор линейной алгебры.

\subsubsection{Обзор линейной
алгебры}\label{ux43eux431ux437ux43eux440-ux43bux438ux43dux435ux439ux43dux43eux439-ux430ux43bux433ux435ux431ux440ux44b}

Общий случай

Частный случай

Определение. Если А---квадратная матрица порядка n, то собственными
значениями матрицы называются n корней характе-

Пусть

\begin{verbatim}
A = \begin{bmatrix} 1 & 1 \\ 4 & 1 \end{bmatrix}
\end{verbatim}

\begin{verbatim}
.

ристического уравнения
\end{verbatim}

\texttt{\textbackslash{}det\textbackslash{}left(A-\textbackslash{}lambda\ I\textbackslash{}right)=0,}

\begin{verbatim}
где  \det(A-\lambda I) — определитель матрицы  A-\lambda I .

Некоторые собственные значения могут совпадать.

Тогда собственные значения являются корнями уравнения
\end{verbatim}

\texttt{\textbackslash{}det\ \textbackslash{}begin\{bmatrix\}\ 1-\textbackslash{}lambda\ \&\ 1\textbackslash{}\textbackslash{}\ 4\ \&\ 1-\textbackslash{}lambda\ \textbackslash{}end\{bmatrix\}\ =\ (1-\textbackslash{}lambda)\^{}2\ -\ 4\ =}
\texttt{=\ \textbackslash{}lambda\^{}2\ -\ 2\textbackslash{}lambda\ -\ 3\ =\ 0.}

\begin{verbatim}
Откуда
\end{verbatim}

\texttt{\textbackslash{}lambda\_1\ =\ -1}

\texttt{,\ \ \textquotesingle{}\textbackslash{}lambda\_2\ =\ 3\textquotesingle{}\ .}

Определение. Если \texttt{\textbackslash{}lambda} ---собственное
значение матрицы A, то соответствующим ему собственным вектором
называется такой ненулевой вектор, который удовлетворяет соотношению

\begin{verbatim}
Ax = \lambda x
\end{verbatim}

\begin{verbatim}
.

Пример. Вектор
\end{verbatim}

\texttt{X\ =\ \textbackslash{}begin\{bmatrix\}\ 1\ \textbackslash{}\textbackslash{}\ 2\ \textbackslash{}end\{bmatrix\}}

\begin{verbatim}
является собственным вектором, соответствующим собственному значению  \lambda_2 = 3 , так как
\end{verbatim}

\texttt{\textbackslash{}begin\{bmatrix\}\ 1\ \&\ 1\ \textbackslash{}\textbackslash{}\ 4\ \&\ 1\ \textbackslash{}end\{bmatrix\}\ \textbackslash{}begin\{bmatrix\}\ 1\ \textbackslash{}\textbackslash{}\ 2\ \textbackslash{}end\{bmatrix\}\ =\ 3\ \textbackslash{}begin\{bmatrix\}\ 1\ \textbackslash{}\textbackslash{}\ 2\ \textbackslash{}end\{bmatrix\}.}

\begin{verbatim}
Определение. Матрица  A^{-1}  называется обратной к матрице A, если
\end{verbatim}

AA\^{}\{-1\} = A\^{}\{-1\}A = I\_1

\begin{verbatim}
Пример, Обратной к
\end{verbatim}

\texttt{A\ =\ \textbackslash{}begin\{bmatrix\}\ 2\ \&\ -2\ \textbackslash{}\textbackslash{}\ 1\ \&\ 1\ \textbackslash{}end\{bmatrix\}}

\begin{verbatim}
является матрица
\end{verbatim}

\texttt{A\^{}\{-1\}\ =\ \textbackslash{}begin\{bmatrix\}\ 0.25\ \&\ 0.5\ \textbackslash{}\textbackslash{}\ -0.25\ \&\ 0.5\ \textbackslash{}end\{bmatrix\},}

\begin{verbatim}
где 1-единичная матрица.

поскольку
\end{verbatim}

\texttt{AA\^{}\{-1\}\ =\ A\^{}\{-1\}A\ =\ \textbackslash{}begin\{bmatrix\}\ 1\ \&\ 0\ \textbackslash{}\textbackslash{}\ 0\ \&\ 1\ \textbackslash{}end\{bmatrix\}.}

\begin{verbatim}
# Диагонализация матрицы

Пусть A — квадратная матрина порядка n и пусть все ее собственные значения  \lambda_1 ,  \lambda_2  ...  \lambda_n  различны. Составим матрицу P так, чтобы у нее *столбцами* были координаты собственных векторов (в k-м столбце стоит собственный вектор  X_k , соответствующий собственному значению  \lambda_k ), т. е.
\end{verbatim}

\texttt{P\ =\ \textbackslash{}begin\{bmatrix\}\ X\_1\ \&\ X\_\{\textbackslash{}tilde\{2\}\}\ \&\ \textbackslash{}dots\ \&\ X\_n\ \textbackslash{}\textbackslash{}\ X\_\{1\}\ \&\ X\_\{2\}\ \&\ \textbackslash{}dots\ \&\ X\_n\ \textbackslash{}end\{bmatrix\}\ =\ \textbackslash{}begin\{bmatrix\}\ x\_\{11\}\ \&\ x\_\{1\textbackslash{}tilde\{2\}\}\ \&\ \textbackslash{}dots\ \&\ x\_\{1n\}\ \textbackslash{}\textbackslash{}\ x\_\{21\}\ \&\ x\_\{22\}\ \&\ \textbackslash{}dots\ \&\ x\_\{2n\}\ \textbackslash{}\textbackslash{}\ \textbackslash{}vdots\ \&\ \textbackslash{}vdots\ \&\ \textbackslash{}ddots\ \&\ \textbackslash{}vdots\ \textbackslash{}\textbackslash{}\ x\_\{n1\}\ \&\ x\_\{n\textbackslash{}tilde\{3\}\}\ \&\ \textbackslash{}dots\ \&\ x\_\{nn\}\ \textbackslash{}end\{bmatrix\}.}

\begin{verbatim}
Тогда матрица Л, определяемая соотношением
\end{verbatim}

\texttt{\textbackslash{}Lambda\ =\ P\^{}\{-1\}AP,}

\begin{verbatim}
где  P^{-1}  — обратная к матрице P, будет диагональной матрицей вида

~~
\Лямбда = \begin{bmatrix} \lambda_1 & & & \\ & \lambda_2 & 0 & \\ & & \ddots & \\ & & & & \lambda_n \end{bmatrix}.
~~

Например, у матрицы
\end{verbatim}

\texttt{A\ =\ \textbackslash{}begin\{bmatrix\}\ 0\ \&\ 8\ \textbackslash{}\textbackslash{}\ 2\ \&\ 0\ \textbackslash{}end\{bmatrix\}}

\begin{verbatim}
собственные числа равны  \lambda_1 = 4  и  \lambda_2 = -4 , а собственные векторы соответственно
\end{verbatim}

\texttt{X\_1\ =\ \textbackslash{}begin\{bmatrix\}\ 2\ \textbackslash{}\textbackslash{}\ 1\ \textbackslash{}end\{bmatrix\}}

\begin{verbatim}
in
\end{verbatim}

\texttt{X\_2\ =\ \textbackslash{}begin\{bmatrix\}\ -2\ \textbackslash{}\textbackslash{}\ 1\ \textbackslash{}end\{bmatrix\}}

\begin{verbatim}
.

Значит, произведение

~~
P^{-1}AP = \begin{bmatrix} 2 & -2 \\ 1 & 1 \end{bmatrix}^{-1} \begin{bmatrix} 0 & 8 \\ 2 & 0 \end{bmatrix} \begin{bmatrix} \begin{bmatrix} 2 & -2 \\ 1 & 1 \end{bmatrix} = \\ = \begin{bmatrix} 1/4 & 1/2 \\ -1/4 & 1/2 \end{bmatrix} \begin{bmatrix} 0 & 8 \\ 2 & 0 \end{bmatrix} \begin{bmatrix} 2 & -2 \\ 1 & 1 \end{bmatrix}
~~

равно диагональной матрице \end{bmatrix}
\end{verbatim}

\texttt{\textbackslash{}Lambda\ =\ \textbackslash{}begin\{bmatrix\}\ 4\ \&\ 0\ \textbackslash{}\textbackslash{}\ 0\ \&\ -4\ \textbackslash{}end\{bmatrix\}.}

\begin{verbatim}
(Читатель может проверить это самостоятельно.)

Теперь можно решить простую систему двух уравнений с частными производными (с соответствующими начальными условиями).

# Решение линейной системы u_t + Au_x = 0 .

Рассмотрим задачу Коши для системы, содержащей два уравнения и два начальных условия:

(УЧП 1)
\end{verbatim}

\texttt{\textbackslash{}frac\{\textbackslash{}partial\ u\_1\}\{\textbackslash{}partial\ t\}\ +\ 8\ \textbackslash{}frac\{\textbackslash{}partial\ u\_2\}\{\textbackslash{}partial\ x\}\ =\ 0}

,\\
(29.3) (УЧП 2)
\texttt{\textbackslash{}frac\{\textbackslash{}partial\ u\_2\}\{\textbackslash{}partial\ t\}\ +\ 2\ \textbackslash{}frac\{\textbackslash{}partial\ u\_1\}\{\textbackslash{}partial\ x\}\ =\ 0}
,
\texttt{-\textbackslash{}infty\ \textless{}\ x\ \textless{}\ \textbackslash{}infty}
, \texttt{0\ \textless{}\ t\ \textless{}\ \textbackslash{}infty} ,\\
(HУ 1) \texttt{u\_1(x,\ 0)\ =\ f(x)} ,\\
\texttt{(HУ\ 2)\ \ \textquotesingle{}u\_2(x,\ 0)\ =\ g(x)\textquotesingle{}\ ,\ \ \textquotesingle{}-\textbackslash{}infty\ \textless{}\ x\ \textless{}\ \textbackslash{}infty\textquotesingle{}\ .}

Эта задача может соответствовать определению давления
\texttt{u\_1(x,\ t)} и плотности \texttt{u\_2(x,\ t)} как функции
пространственной координаты x и времени t по известным распределениям
этих величин в начальный момент.

Запишем систему уравнений в матричной форме

\begin{verbatim}
\begin{bmatrix} \frac{\partial u_1}{\partial t} \\ \frac{\partial u_2}{\partial t} \end{bmatrix} + \begin{bmatrix} 0 & 8 \\ 2 & 0 \end{bmatrix} \begin{bmatrix} \frac{\partial u_1}{\partial x} \\ \frac{\partial u_2}{\partial x} \end{bmatrix} = \begin{bmatrix} 0 \\ 0 \end{bmatrix}
\end{verbatim}

или \texttt{(29\ \textbackslash{}\ 4)}

u\_t + Au\_s = 0

\begin{verbatim}
где
\end{verbatim}

\texttt{A\ =\ \textbackslash{}begin\{bmatrix\}\ 0\ \&\ 8\ \textbackslash{}\textbackslash{}\ 2\ \&\ 0\ \textbackslash{}end\{bmatrix\},\ \textbackslash{}quad\ u\_1\ =\ \textbackslash{}begin\{bmatrix\}\ \textbackslash{}frac\{\textbackslash{}partial\ u\_1\}\{\textbackslash{}partial\ t\}\ \textbackslash{}\textbackslash{}\ \textbackslash{}frac\{\textbackslash{}partial\ u\_2\}\{\textbackslash{}partial\ t\}\ \textbackslash{}end\{bmatrix\},\ \textbackslash{}quad\ u\_x\ =\ \textbackslash{}begin\{bmatrix\}\ \textbackslash{}frac\{\textbackslash{}partial\ u\_1\}\{\textbackslash{}partial\ x\}\ \textbackslash{}\textbackslash{}\ \textbackslash{}frac\{\textbackslash{}partial\ u\_2\}\{\textbackslash{}partial\ x\}\ \textbackslash{}end\{bmatrix\},\ \textbackslash{}quad\ 0\ =\ \textbackslash{}begin\{bmatrix\}\ 0\ \textbackslash{}\textbackslash{}\ 0\ \textbackslash{}end\{bmatrix\}.}

\begin{verbatim}
Введем новые неизвестные функции
\end{verbatim}

\texttt{v\ =\ \textbackslash{}begin\{bmatrix\}\ v\_1\ \textbackslash{}\textbackslash{}\ v\_2\ \textbackslash{}end\{bmatrix\}}

\begin{verbatim}
с помощью преобразования
\end{verbatim}

u = Pv

.

где P---матрица, по столбцам которой стоят собственные векторы матрицы A
(матрица P нам уже известна). Оказывается, что после такого
преобразования для определения v получится очень простая система (два
уравнения относительно новых неизвестных \texttt{v\_1} и \texttt{v\_2}
оказываются независимыми). Значит, \texttt{v\_1} и \texttt{v\_2} легко
находятся. После того как \texttt{v\_1} и \texttt{v\_2} найдены, по
формуле

\begin{verbatim}
u = Pv
\end{verbatim}

находятся искомые функции и, и и2.

Сначала, однако, выясним, как выглядит система для определения v. Чтобы
разобраться в этом, читатель должен обратиться к матричному исчислению.
Продифференцировав обе части соотношения u = Pv, получаем

(29 5)

\begin{verbatim}
\frac{\partial u}{\partial t} = P \frac{\partial v}{\partial t},
\frac{\partial u}{\partial x} = P \frac{\partial v}{\partial x}.
\end{verbatim}

(Чтобы лучше понять дальнейшее изложение, читателю рекомендуется
записать эти уравнен я в развернутой форме.)

Теперь подставим соотношения (29.5) в систему

\begin{verbatim}
u_t + Au_x = 0
\end{verbatim}

и получим

\begin{verbatim}
Pv_t + APv_x = 0.
\end{verbatim}

Если умножить теперь обе части последнего уравнения на
\texttt{P\^{}\{-1\}} , то получим

\begin{verbatim}
v_t + P^{-1}APv_x = 0,
\end{verbatim}

или

(29,6) \texttt{v\_t} + \Лямбда \texttt{v\_x\ =\ 0}.

Раньше мы уже видели, что для нашей матрицы

\begin{verbatim}
\Lambda = \begin{bmatrix} 4 & 0 \\ 0 & -4 \end{bmatrix}.
\end{verbatim}

Запишем (29.6) в развернутой форме

(29.7)

\begin{verbatim}
\frac{\frac{\partial v_1}{\partial t} + 4 \frac{\partial v_1}{\partial x} = 0, \\ \frac{\partial v_2}{\partial t} - 4 \frac{\partial v_2}{\partial x} = 0.
\end{verbatim}

Получилась система из двух несвязанных уравнений, которые можно решать
независимо. Читателю должно быть известно, что решениями этих уравнений
являются бегущие волны

\begin{verbatim}
v_1(x, t) = \varphi(x-4t),
\end{verbatim}

\texttt{v\_2(x,\ t)\ =\ \textbackslash{}psi(x+4t),} где ф и
ф-произвольные дифференцируемые функции.

Таким образом, мы нашли функции \texttt{v\_1} и \texttt{v\_2} . Для
получения общего решения u нужно вычислить

\begin{verbatim}
u = Pv = \begin{bmatrix} 2 & -2 \\ 1 & 1 \end{bmatrix} \begin{bmatrix} v_1 \\ v_2 \end{bmatrix} =
= \begin{bmatrix} 2 & -2 \\ 1 & 1 \end{bmatrix} \begin{bmatrix} \varphi(x - 4t) \\ \psi(x + 4t) \end{bmatrix} =
= \begin{bmatrix} 2\varphi(x - 4t) - 2\psi(x + 4t) \\ \varphi(x - 4t) + \psi(x + 4t) \end{bmatrix}.
\end{verbatim}

Другими словами,

(29.8)

\begin{verbatim}
u_1(x, t) = 2\varphi(x-4t) - 2\psi(x+4t), u_2(x, t) = \varphi(x-4t) + \psi(x+4t).
\end{verbatim}

Следовательно, мы нашли общее решение системы двух уравнений (29.3).

Типичным решением (а их бесконечное число) могут быть, например, функции

\begin{verbatim}
\phi(\xi) = \sin \xi, \ \psi(\xi) = \xi^2
\end{verbatim}

(две произвольные функции)

и, следовательно,

\begin{verbatim}
u_1(x, t) = 2 \sin(x-4t) - 2(x+4t)^2,
\end{verbatim}

\texttt{u\_2(x,\ t)\ =\ \textbackslash{}sin(x-4t)\ +\ (x+4t)\^{}2.}
Подставим теперь общее решение (29.8) в начальные условия

\begin{verbatim}
u_1(x, 0) = f(x),
\end{verbatim}

\texttt{u\_2(x,\ 0)\ =\ g(x),} в результате получаем

\begin{verbatim}
2\varphi(x) - 2\psi(x) = f(x),
\varphi(x) + \psi(x) = g(x),
\end{verbatim}

откуда функци і ф и ф легко находятся:

\begin{verbatim}
\varphi(x) = \frac{1}{4} \left[ f(x) + 2g(x) \right],
\psi(x) = \frac{1}{4} \left[ 2g(x) - f(x) \right].
\end{verbatim}

\textbf{Тепе}рь можно написать решение задачи (29.3) в окончательном
виде

\begin{verbatim}
u_{1}(x, t) = 2\varphi(x-4t) - 2\psi(x+4t) =
= \frac{1}{2} \left[ f(x-4t) + 2g(x-4t) \right] - \frac{1}{2} \left[ 2g(x+4t) - f(x+4t) \right],
\end{verbatim}

(29.9)

\begin{verbatim}
u_{2}(x, t) = \varphi(x-4t) + \psi(x+4t) =
= \frac{1}{4} \left[ f(x-4t) + 2g(x-4t) \right] + \frac{1}{4} \left[ 2g(x+4t) - f(x+4t) \right].
\end{verbatim}

\paragraph{ЗАМЕЧАНИЯ}\label{ux437ux430ux43cux435ux447ux430ux43dux438ux44f-1}

\begin{enumerate}
\def\labelenumi{\arabic{enumi}.}
\item
  Многие численные методы ориентированы на решение систем уравнений и,
  значит, многне программы для ЭВМ написаны так, чтобы решать системы
  уравнений первого порядка. По этой причине, прежде чем воспользоваться
  программои, приходится преобразовать свое уравнение высокого порядка в
  систему уравнений первого порядка.
\item
  Линейную систему
\end{enumerate}

\begin{verbatim}
u_t + A(x, t) u_x + B(x, t) u = 0
\end{verbatim}

можно решать тем же методом, что и систему \texttt{u\_t\ +\ Au\_x\ =\ 0}
. Разница в том, что в этом случае матрица P из собственных векторов
матрицы \texttt{A(x,\ t)} будет функцией x и t.

\paragraph{ЗАДАЧИ}\label{ux437ux430ux434ux430ux447ux438-1}

\begin{enumerate}
\def\labelenumi{\arabic{enumi}.}
\tightlist
\item
  Запишите систему уравнений (29.2) в матричной форме
\end{enumerate}

\begin{verbatim}
Au_t + Bu_x + Cu = 0,
\end{verbatim}

где A, B и C-матрицы третьего порядка.

\begin{enumerate}
\def\labelenumi{\arabic{enumi}.}
\setcounter{enumi}{1}
\tightlist
\item
  Найдите собственные векторы и собственные значения матрицы
\end{enumerate}

\begin{verbatim}
A = \begin{bmatrix} 1 & 1 \\ 4 & 1 \end{bmatrix}.
\end{verbatim}

\begin{enumerate}
\def\labelenumi{\arabic{enumi}.}
\setcounter{enumi}{2}
\tightlist
\item
  Воспользовавшись результатами задачи 2, решите систему
\end{enumerate}

\begin{verbatim}
\begin{cases} \frac{\partial u_1}{\partial t} + \frac{\partial u_1}{\partial x} + \frac{\partial u_2}{\partial x} = 0, \\ \frac{\partial u_2}{\partial t} + 4 \frac{\partial u_1}{\partial x} + \frac{\partial u_2}{\partial x} = 0. \end{cases}
\end{verbatim}

УКАЗАНИЕ. Запишите сначала систему в матричной форме.

\begin{enumerate}
\def\labelenumi{\arabic{enumi}.}
\setcounter{enumi}{3}
\item
  Проверьте формулы (29.5), записав их в скалярной форме.
\item
  Проверьте, что функции \texttt{u\_1} и \texttt{u\_2} из (29.8)
  удовлетворяют обоим уравнениям (29.3).
\end{enumerate}

\section{Лекція 30. КОЛИВАННЯ МЕМБРАНИ (ХВИЛЬОВЕ РІВНЯННЯ В ПОЛЯРНИХ
КООРДИНАТАХ)}\label{ux43bux435ux43aux446ux456ux44f-30.-ux43aux43eux43bux438ux432ux430ux43dux43dux44f-ux43cux435ux43cux431ux440ux430ux43dux438-ux445ux432ux438ux43bux44cux43eux432ux435-ux440ux456ux432ux43dux44fux43dux43dux44f-ux432-ux43fux43eux43bux44fux440ux43dux438ux445-ux43aux43eux43eux440ux434ux438ux43dux430ux442ux430ux445}

ЦЕЛЬ ЛЕКЦИИ: Показать, что колебания мембраны можно описать волновым
уравнением

\begin{verbatim}
u_{tt} = c^2 \left( u_{rr} + \frac{1}{r} u_r + \frac{1}{r^2} u_{\theta\theta} \right)
\end{verbatim}

\begin{verbatim}
.

Если искать решение этого уравнения в виде стоячих волн:
\end{verbatim}

\texttt{u(r,\ \textbackslash{}theta,\ t)\ =\ P(r)\ \textbackslash{}Theta(\textbackslash{}theta)\ T(t),}

\begin{verbatim}
то оно распадается на три обыкновенных дифференциальных уравнения:
\end{verbatim}

\texttt{T\textquotesingle{}\textquotesingle{}\ +\ \textbackslash{}lambda\^{}2\ c\^{}2\ T\ =\ 0}

\begin{verbatim}
(уравнение гармонических колебаний),  r^2 R'' + rR' + (\lambda^2 r^2 - n^2) R = 0  (уравнение Бесселя), (уравнение гармонических колебаний).

Все произведения  R(r) \Theta(\theta) T(t)  решений этих трех обыкновенных дифференциальных уравнений описывают основные колебания мембраны, а функции  R(r) \Theta(\theta)  описывают форму мембраны. Чтобы найти колебания мембраны при произвольных начальных условиях, мы составим такую комбинацию фундаментальных решений, чтобы начальные условия удовлетворялись.

Цель этой лекции ясна—найти колебания круглой мембраны с заданными начальными и граничными условиями. Для простоты будем считать, что радиус мембраны равен единице, а смещение по границе равно нулю. Пусть  u\left(r,\theta,t\right)  обозначает возвышение точек мембраны под плоскостью равновесия. Нам необходимо решить задачу (см. рис. 30.1)
\end{verbatim}

\texttt{\textbackslash{}begin\{array\}\{lll\}\ (\textbackslash{}text\{UPP\})\ \&\ u\_\{tt\}\ =\ c\^{}2\ \textbackslash{}left(\ u\_\{rr\}\ +\ \textbackslash{}frac\{1\}\{r\}\ u\_r\ +\ \textbackslash{}frac\{1\}\{r\^{}2\}\ u\_\{\textbackslash{}theta\textbackslash{}theta\}\ \textbackslash{}right),\ \&\ 0\ \textless{}\ r\ \textless{}\ 1,\ \textbackslash{}\ 0\ \textless{}t\ \textless{}\ \textbackslash{}infty,\ \textbackslash{}\textbackslash{}\ (\textbackslash{}text\{GU\})\ \&\ u\ =\ 0\ \&\ \textbackslash{}text\{when\}\ r\ =\ 1,\ \&\ 0\ \textless{}\ t\ \textless{}\ \textbackslash{}infty,\ \textbackslash{}\textbackslash{}\ (\textbackslash{}text\{NU\})\ \&\ \textbackslash{}begin\{cases\}\ u\ =\ f(r,\ \textbackslash{}\ \textbackslash{}theta)\ \textbackslash{}\textbackslash{}\ u\_t\ =\ g(r,\ \textbackslash{}\ \textbackslash{}theta)\ \textbackslash{}end\{cases\}\ \&\ \textbackslash{}text\{when\}\ t\ =\ 0.\ \textbackslash{}end\{array\}}

\begin{verbatim}
Чтобы решить эту задачу, вспомним задачу о колебаниях скрипичной струны из лекции 20, которая решалась суперпозицией бесконечного числа простых колебаний. Если воспользовать-

ся этим же подходом, то решения следует искать в виде
\end{verbatim}

\texttt{u(r,\ \textbackslash{}theta,\ t)\ =\ U(r,\ \textbackslash{}theta)\ T(t)}

.

Форма этих колебаний определяется функцией
\texttt{U(r,\ \textbackslash{}theta)} , а характер осцилляций ---
множителем \texttt{T(t)}.

Подставив это представление в волновое уравнение (см. задачу 1), получим
два уравнения

\begin{verbatim}
\Delta U + \lambda^2 U = 0
\end{verbatim}

(уравнение Гельм-гольца),

\begin{verbatim}
T'' + \lambda^2 c^2 T = 0
\end{verbatim}

(уравнение гармонических колебаний),

\pandocbounded{\includegraphics[keepaspectratio]{images/_page_221_Picture_9.jpeg}}

РИС. 30.1. Қолеблющаяся мембрана (гиперболическая смешанная задача).

где

\begin{verbatim}
\Delta U = U_{rr} + \frac{1}{r} U_r + \frac{1}{r^2} U_{\theta\theta}.
\end{verbatim}

Отметим, что здесь мы требуем, чтобы константа разделения была
положительна (поэтому мы обозначим ее
\texttt{\textbackslash{}lambda\^{}\{2-1\}} )). Только в этом случае
функции \texttt{T(t)} будут периодическими.

На следующем шаге мы решим уравнение Гельмгольца, но сначала разберемся
с граничным условием для него. Чтобы найти его, подставим
\texttt{u(r,\ \textbackslash{}theta,\ t)\ =\ U(r,\ \textbackslash{}theta)\ T(t)}
в граничное условие для мембраны

\begin{verbatim}
u(1, \theta, t) = U(1, \theta) T(t) = 0, \quad 0 < t < \infty,
\end{verbatim}

или

\begin{verbatim}
U(1, \theta) = 0.
\end{verbatim}

Следовательно, для того чтобы найти форму фундаментальных колебаний
мембраны, необходимо решить задачу

\begin{verbatim}
\Delta U + \lambda^2 U = 0,
U(1, \theta) = 0.
\begin{cases} \Delta U + \mu u = 0 & \text{у колі,} \\ U = 0 & \text{на межі кола} \end{cases}
\end{verbatim}

при \texttt{\textbackslash{}mu\ \textbackslash{}le\ 0} имеет только
тривиальное решение U=0. --- Прим. ред.

1) Поскольку краевая задача

\texttt{Это\ хорошо\ известная\ эллиптическая\ задача\ на\ собственные\ значения,\ и\ нам\ надо\ найти\ все\ \ \textquotesingle{}\textbackslash{}lambda\textquotesingle{}\ ,\ при\ которых\ она\ имеет\ ненуленые\ решения.\ Решения\ \ \textquotesingle{}U\textbackslash{}left(r,\textbackslash{}theta\textbackslash{}right)\textquotesingle{}\ \ описывают\ форму\ фундаментальных\ колебаний\ мембраны,\ а\ собственные\ числа\ \ \textquotesingle{}\textbackslash{}lambda\textquotesingle{}\ \ являются\ нулями\ некоторых\ функций\ Бесселя\ и\ пропорциональны\ частотам\ этих\ колебаний.}

Итак, нам предстоит теперь решить задачу на собственные значения для
уравнения Гельмгольца (эта задача очень важна и сама по себе). Будем
решать ее тем же методом, что и другие линейные однородные уравнения с
нулевыми граничными условиями, а именно методом разделения переменных.

\subsubsection{Решение задачи на собственные значения для уравнения
Гельмгольца}\label{ux440ux435ux448ux435ux43dux438ux435-ux437ux430ux434ux430ux447ux438-ux43dux430-ux441ux43eux431ux441ux442ux432ux435ux43dux43dux44bux435-ux437ux43dux430ux447ux435ux43dux438ux44f-ux434ux43bux44f-ux443ux440ux430ux432ux43dux435ux43dux438ux44f-ux433ux435ux43bux44cux43cux433ux43eux43bux44cux446ux430}

Чтобы решить задачу

\begin{verbatim}
\begin{array}{ll} (\mathsf{Y}\mathsf{H}\mathsf{\Pi}) & \Delta U + \lambda^2 U = 0, \\ (\mathsf{\Gamma}\mathsf{Y}) & U(1, \theta) = 0, \end{array}
\end{verbatim}

подставим

\begin{verbatim}
U(r, \theta) = R(r) \Theta(\theta)
\end{verbatim}

в задачу Гельмгольца. Сделав это, получаем

\begin{verbatim}
r^2R'' + rR' + (\lambda^2r^2 - n^2)R = 0
\end{verbatim}

\texttt{(уравнение\ Бесселя),\ \ \textquotesingle{}R(1)\ =\ 0\textquotesingle{}\ ,\ \ \textquotesingle{}\textbackslash{}Theta\textquotesingle{}\textquotesingle{}\ +\ n\^{}2\textbackslash{}Theta\ =\ 0\textquotesingle{}}

\texttt{(проделайте\ все\ это\ сами).\ Отметим,\ что\ мы\ выбрали\ константу\ разделения\ в\ виде\ \ \textquotesingle{}n\^{}2\textquotesingle{}\ ,\ \ \textquotesingle{}n=0,\textbackslash{}\ 1,\textbackslash{}\ 2,\textbackslash{}\ \textbackslash{}ldots\textquotesingle{}\ ,\ так\ как\ хотим,\ чтобы\ функции\ \ \textquotesingle{}\textbackslash{}Theta\textbackslash{}left(\textbackslash{}theta\textbackslash{}right)\textquotesingle{}\ \ были\ периодическими\ с\ периодом\ \ \textquotesingle{}2\textbackslash{}pi\textquotesingle{}\ .\ Очевидно,\ что\ форма\ мембраны\ является\ периодической\ функцией\ по\ \ \textquotesingle{}\textbackslash{}theta\textquotesingle{}\ .\ Итак,\ для\ того\ чтобы\ решить\ задачу\ Гельмгольца,\ мы\ должны\ решить\ два\ обыкновенных\ дифференциальных\ уравнения}

\begin{verbatim}
r^2R'' + rR' + (\lambda^2r^2 - n^2) R = 0
\end{verbatim}

\begin{verbatim}
,  0 < r < 1 ,  R(0) < \infty  (физическое ограничение),  R(1) = 0 ,  \Theta'' + n^2\Theta = 0 .

# Уравнение Бесселя

Уравнение
\end{verbatim}

\texttt{r\^{}2R\textquotesingle{}\textquotesingle{}\ +\ rR\textquotesingle{}\ +\ (\textbackslash{}lambda\^{}2r\^{}2\ -\ n\^{}2)R\ =\ 0}

\begin{verbatim}
жорошо известно в теории обыкновенных дифференциальных уравнений. Это так называемое уравнение Бесселя. Оно имеет два

линейно-независимых решения

 R_{1}(r) = A J_{n}(\lambda r)  функция Бесселя *первого* рода *п*-го порядка,

 R_{2}(r) = BY_{n}(\lambda r)  функция Бесселя второго рода n-го порядка,

и, следовательно, общее решение этого уравнения представимо в виде
\end{verbatim}

\texttt{R(r)\ =\ AJ\_n(\textbackslash{}lambda\ r)\ +\ BY\_n(\textbackslash{}lambda\ r).}

\begin{verbatim}
Отметим, что рещения зависят от параметров n и  \lambda , входящих в уравнение. Некоторые графики функции Бесселя приведены на рис. 30.2. Для определения функций  J_n(\lambda r)  и  Y_n(\lambda r)  можно

![](images/_page_223_Figure_8.jpeg)

![](images/_page_223_Figure_9.jpeg)

Рис. 30 2 Функции Бесселя. a — функции Бесселя первого рода;  \pmb{\delta}  — функции Бесселя второго рода.

воспользоваться методом Фробениуса, который позволяет найти R(r) в виде степенных рядов. Этим методом находятся два линейно-независимых степенных ряда  J_n(\lambda r)  и  Y_n(\lambda r) . Поскольку функция  Y_n(\lambda r)  не ограничена при r=0, мы выбираем решение в виде
\end{verbatim}

\texttt{R(r)\ =\ AJ\_n(\textbackslash{}lambda\ r).}

\begin{verbatim}
Последний шаг в определении  R\left( r\right)  связан с использованием граничного условия  R\left( \mathbf{l}\right) =0  для определения всех  \lambda  (в данный

Таблица 30.1

| Корни | уравнений. | J_n(r) = 0 |
|-------|------------|--------------|
|-------|------------|--------------|

| Γ | 1 | n     |       |       |       |       |  |  |
|---|---|-------|-------|-------|-------|-------|--|--|
|   |   | 0     | 1     | 2     | 3     | 4     |  |  |
|   | 1 | 2,40  | 3,83  | 5,13  | 6,38  | 7,59  |  |  |
| ı | 2 | 5,52  | 7,02  | 8,42  | 9,76  | 11,06 |  |  |
| ľ | 3 | 8,65  | 10,17 | 11,62 | 13,02 | 14,37 |  |  |
|   | 4 | 11,79 | 13,32 | 14,80 | 16,22 | 17,62 |  |  |
|   | 5 | 14,93 | 16,47 | 17,96 | 19,41 | 20,83 |  |  |
|   | . | •     |       | •     |       | •     |  |  |
|   |   |       | •     |       | •     | •     |  |  |
|   | • | •     | •     | •     | •     | •     |  |  |

момент нас не интересует коэффициент A). Подставляя R (1) = 0, в R (r) =  J_n  ( \lambda r ), получаем  J_n  ( \lambda ) = 0. Другими словами, для того чтобы функция R (r) обращалась в нуль на границе круга, мы должны выбрать константу разделения так, чтобы она была корнем уравнения  J_n  (r) = 0,  \tau . e.
\end{verbatim}

\texttt{\textbackslash{}lambda\ =\ k\_\{nm\}}

\begin{verbatim}
где  k_{nm}  — m-й корень уравнения  J_n(r) = 0 .

Хорошо известны таблицы этих корней, но их можно вычислять на ЭВМ. Значения нескольких первых корней представлены в табл. 30.1. Зная эти корни, мы сразу можем построить решение краевой задачи для уравнения Гельмгольца. Если собственные значения  \lambda = k_{nm} , то собственные функции  U_{nm}(r)  определяются формулой
\end{verbatim}

\texttt{U\_\{nm\}(r,\ \textbackslash{}theta)\ =\ J\_n(k\_\{nm\}r)\ {[}A\ \textbackslash{}sin(n\textbackslash{}theta)\ +\ B\ \textbackslash{}cos(n\textbackslash{}theta){]}.}

\begin{verbatim}
На рис. 30.3 изображены некоторые из этих функций для различных значений n и m. Общая форма функции  U_{nm}(r,\theta)  не зависит от значений констант A и B. Эти константы влияют только на амплитуду колебаний и расположение линии узлов относительно  \theta=0 .

Каждая функция представляет собой фундаментальное колебание круглой мембраны с частотой
\end{verbatim}

\texttt{f\_\{nm\}\ =\ k\_\{nm\}c/2\textbackslash{}pi.}

\begin{verbatim}
Эти частоты мы нашли из решений
\end{verbatim}

\texttt{T\_\{nm\}(t)\ =\ A\ \textbackslash{}sin(k\_\{nm\}ct)\ +\ B\ \textbackslash{}cos(k\_\{nm\}ct)}

\begin{verbatim}
#### уравнения
\end{verbatim}

T'\,' + k\_\{nm\}\^{}2 c\^{}2 T = 0.

\begin{verbatim}
Интересно отметить, что отношение частоты колебания  U_{nm}(\mathbf{r}, \mathbf{\theta})  к частоте основного колебания  U_{nt}(\mathbf{r}, \mathbf{\theta}) , т.е.
\end{verbatim}

\texttt{\textbackslash{}frac\{f\_\{nm\}\}\{f\_\{01\}\}\ =\ \textbackslash{}frac\{k\_\{nm\}\}\{k\_\{01\}\}\ ,}

\begin{verbatim}
в отличие от одномерного случая не является целым числом. Другими словами, высшие колебания мембраны не являются чистыми обертонами основной частоты. Отметим также, что число m совпадает с числом круговых узловых линий у функции  U_{um}(r,\theta)

![](images/_page_225_Figure_7.jpeg)

РИС. 30.3. Линин узлов для основных гармоник круглой мембраны.

(амплитуда колебаний в точках узловых линий равна нулю). Радиусы этих узловых линий определяются отношениями
\end{verbatim}

\texttt{\textbackslash{}frac\{k\_\{01\}\}\{k\_\{0m\}\}}

\texttt{,\ \ \textquotesingle{}\textbackslash{}frac\{k\_\{02\}\}\{k\_\{0m\}\}\textquotesingle{}\ ,\ \ \textquotesingle{}\textbackslash{}frac\{k\_\{03\}\}\{k\_\{0m\}\}\textquotesingle{}\ ,\ ...,\ 1.}

Теперь, когда найдены собственные функции задачи Гельмгольца, их следует
умножить на временной множитель

\begin{verbatim}
T_{nm}(t) = A \sin(k_{nm}ct) + B \cos(k_{nm}ct)
\end{verbatim}

и найти такую линейную комбинацию получившихся функций, которая бы
удовлетворяла начальным условиям. Это значит, что решение исходной
задачи мы будем представлять в виде

\begin{verbatim}
u(r, \theta, t) = \sum_{n=0}^{\infty} \sum_{m=1}^{\infty} J_n(k_{nm}r) \cos(n\theta) [A_{nm}\sin(k_{nm}ct) + B_{nm}\cos(k_{nm}ct)].
\end{verbatim}

\texttt{Отметим,\ что,\ заменив\ величину\ \ \textquotesingle{}A\textbackslash{}sin\textbackslash{}left(n\textbackslash{}theta\textbackslash{}right)+B\textbackslash{}cos\textbackslash{}left(n\textbackslash{}theta\textbackslash{}right)\textquotesingle{}\ \ величиной\ \ \textquotesingle{}\textbackslash{}cos\textbackslash{}left(n\textbackslash{}theta\textbackslash{}right)\textquotesingle{}\ ,\ мы\ только\ по-другому\ выбрали\ начало\ отсчета\ угла\ \ \textquotesingle{}\textbackslash{}theta\textquotesingle{}\ .\ Константы\ \ \textquotesingle{}A\_\{nm\}\textquotesingle{}\ \ и\ \ \textquotesingle{}B\_\{nm\}\textquotesingle{}\ \ тоже\ можно\ было\ объединить\ в\ одну,\ но\ это\ уже\ несущественно.}

Вместо того чтобы строить сложную процедуру отыскания коэффициентов
\texttt{A\_\{nm\}} и \texttt{B\_\{nm\}} в общем случае, мы рассмотрим
случай, когда u не зависит от \texttt{\textbackslash{}theta} (тоже
довольно общий). Другими словами, предположим, что начальное смещение
мембраны зависит только от r; т. е.

\begin{verbatim}
u(r, \theta, 0) = f(r),
\end{verbatim}

\texttt{u\_i(r,\ \textbackslash{}theta,\ 0)\ =\ 0.} (Так же легко можно
рассмотреть случай \texttt{u\_t\ \textbackslash{}neq\ 0} .) При
сделанных предположениях решение следует искать в виде

(30.1)

\begin{verbatim}
u(r, t) = \sum_{m=1}^{\infty} A_m J_0(k_{nm}r) \cos(k_{nm}ct).
\end{verbatim}

Наша задача --- подобрать коэффициенты \texttt{A\_m} таким образом,
чтобы выполнялось начальное условие

(30.2)

\begin{verbatim}
F(R) = \sum_{M=1}^{\infty} A_m J_0(k_{0M}R)
\end{verbatim}

\texttt{(условие\ \ \textquotesingle{}u\_t=0\textquotesingle{}\ \ выполняется\ при\ этом\ автоматически).\ Для\ определения\ коэффициентов\ \ \textquotesingle{}A\_m\textquotesingle{}\ \ воспользуемся\ ортогональностью\ функций\ Бесселя\ \ \textquotesingle{}\textbackslash{}\{J\_0(k\_\{0m\}r)\textbackslash{}colon\ m=1,\textbackslash{}\ 2,\textbackslash{}\ \textbackslash{}dots\textbackslash{}\}\textquotesingle{}}

\begin{verbatim}
\int_{0}^{1} r J_{0}(k_{0i}r) J_{0}(k_{0f}r) dr = \begin{cases} 0, & i \neq j, \\ \frac{1}{2} J_{1}^{2}(k_{0i}), & i = j, \end{cases}
\end{verbatim}

\texttt{(метод\ вычисления\ этого\ интеграла\ можно\ найти\ в\ книгах\ по\ функциям\ Бесселя).\ Умножим\ обе\ части\ уравнения\ (30.2)\ на\ \ \textquotesingle{}rJ\_0\textbackslash{}left(k\_\{bf\}r\textbackslash{}right)\textquotesingle{}\ \ и\ проинтегрируем\ от\ 0\ до\ 1:}

\begin{verbatim}
A_{f}\int_{0}^{1}rJ_{\theta}^{2}(k_{\theta f}r)dr=\int_{0}^{1}rf(r)J_{\theta}(k_{\theta f}r)dr.
\end{verbatim}

Из последнего соотношения определяем коэффициенты

(30.3)

\begin{verbatim}
A_{j} = 2 \int_{0}^{1} r f(r) J_{0}(k_{0j}r) dr / J_{1}^{2}(k_{0j}), \quad j = 1, 2, \dots
\end{verbatim}

Итак, не зависящие от \texttt{\textbackslash{}theta} колебания мембраны
описываются функцией вида (30.1), коэффициенты которой вычисляются по
формулам (30.3). Это решение не так сложно, как могло бы показаться
читателю. Его можно интерпретировать следующим образом. Разложим
начальное условие \texttt{f(r)} в сумму

\begin{verbatim}
f(r) = A_1 J_0(k_{01}r) + A_2 J_0(k_{02}r) + A_3 J_0(k_{03}r) + \dots,
\end{verbatim}

а затем умножим каждое слагаемое на соответствующий осцилляторный
множитель \texttt{\textbackslash{}cos\{(k\_\{0m\}ct)\}} , т. е.

\begin{verbatim}
u(r, t) = A_1 J_0(k_{01}r) \cos(k_{01}ct) + A_2 J_0(k_{02}r) \cos(k_{02}ct) + \dots
\end{verbatim}

В результате получили решение исходной задачи. Например, решение задачи
о колебаниях мембраны с начальными условиями

\begin{verbatim}
u(r, \theta, 0) = J_0(2.4 \ r) + 0.5 \ J_0(8.65 \ r),
\end{verbatim}

\texttt{u\_t(r,\ \textbackslash{}theta,\ 0)\ =\ 0} имеет вид

\begin{verbatim}
u(r, t) = J_0(2, 4 r) \cos(2, 4 ct) + 0,5 J_0(8,65 r) \cos(8,65 ct).
\end{verbatim}

\texttt{(Поскольку\ \ \textquotesingle{}k\_\{01\}=2,4048\textbackslash{}ldots=2,40\textquotesingle{}\ ,\ \ \textquotesingle{}k\_\{03\}=8,6537\textbackslash{}ldots\textbackslash{}neq\ 8,65\textquotesingle{}\ ,\ то\ это\ решение\ не\ является\ точным.)}

\section{\texorpdfstring{Графики функций \texttt{J\_\{0\}(k\_\{01\}r)} ,
\texttt{J\_\{0\}(k\_\{02\}r)} , . .
.}{Графики функций J\_\{0\}(k\_\{01\}r) , J\_\{0\}(k\_\{02\}r) , . . .}}\label{ux433ux440ux430ux444ux438ux43aux438-ux444ux443ux43dux43aux446ux438ux439-j_0k_01r-j_0k_02r-.-.-.}

Читателю будет полезно иметь наглядное представление о графиках функций
\texttt{J\_\{\textbackslash{}mathfrak\{g\}\}(k\_\{\textbackslash{}mathfrak\{g\}\_1\}r)}
,
\texttt{J\_\{\textbackslash{}mathfrak\{g\}\}(k\_\{\textbackslash{}mathfrak\{g\}\_2\}r)}
\ldots, так как решение выражается через эти функции. На рис. 30.4
изображен график \texttt{J\_\{\textbackslash{}mathfrak\{g\}\}(r)} .

\pandocbounded{\includegraphics[keepaspectratio]{images/_page_227_Figure_15.jpeg}}

РИС. 30.4. Функция Бесселя нулевого порядка
\texttt{J\_\{\textbackslash{}phi\}(r)} .

Теперь, для того чтобы построить графики функций
\texttt{J\_0(k\_\{01\}r)} , \texttt{J\_0(k\_\{02\}r)} , \ldots,
\texttt{J\_0(k\_\{0m\}r)} , сожмем графики вдоль оси x так, чтобы m-й
нуль попал в точку r=1 (рис. 30.5).

\pandocbounded{\includegraphics[keepaspectratio]{images/_page_228_Figure_3.jpeg}}

\texttt{РИС.\ 30\ 5.\ Графики\ функций\ \ \textquotesingle{}J\_0\textbackslash{}left(k\_\{\textbackslash{}theta\ m\}r\textbackslash{}right)\textquotesingle{}\ \ (основные\ гармоники\ колебаний\ мембраны).}

\paragraph{ЗАМЕЧАНИЯ}\label{ux437ux430ux43cux435ux447ux430ux43dux438ux44f-2}

\begin{enumerate}
\def\labelenumi{\arabic{enumi}.}
\tightlist
\item
  Предположим, что начальное смещение мембраны, у которой c=1, задается
  формулами
\end{enumerate}

(НУ)

\begin{verbatim}
\begin{cases} u\left(r,\;\theta,\; 0\right) = J_{\phi}\left(2,4\; r\right) + 0.10\;\; J_{\phi}\left(5,52\; r\right),\\ u_{t}\left(r,\;\theta,\; 0\right) = 0 \qquad \text{(нульова початкова швидкість)}. \end{cases}
\end{verbatim}

\texttt{В\ соответствии\ с\ общим\ методом\ мы\ должны\ разложить\ начальное\ смещение\ мембраны\ в\ сумму\ функций\ вида\ \ \textquotesingle{}J\_\{0\}\textbackslash{}left(k\_\{0m\}r\textbackslash{}right)\textquotesingle{}\ ,\ найти\ колебания\ каждого\ слагаемого,\ а\ затем\ все\ колебания\ сложить.\ Основная\ особенность\ колебаний\ мембраны\ Р\ И}

\pandocbounded{\includegraphics[keepaspectratio]{images/_page_228_Figure_9.jpeg}}

С. 30 б. Начальные условия разлагаются по основным гармоникам.

в том, что каждое из смещений вида
\texttt{J\_\{\textbackslash{}mathrm\{\ \}0\}} (2,4 r),
\texttt{J\_\{\textbackslash{}mathrm\{\ \}0\}} (5,52 r), \ldots{}
\texttt{J\_\{\textbackslash{}mathrm\{\ \}0\}} (
\texttt{k\_\{\textbackslash{}mathrm\{\ \}0m\}} r) совершает простое
гармоническое движение (т. е. все точки мембраны колеблются с одной и
той же частотой). В нашем случае начальное смещение является суммой двух
простейших и, значит, мембрана будет участвовать в двух

гармонических колебаннях (см. рис. 30.6)

\begin{verbatim}
u(r, t) = J_0(2.4 r) \cos(2.4 t) + 0.10 J_0(5.52 r) \cos(5.52 t).
\end{verbatim}

В результате сложения этих простейших колебаний получается
результирующее сложное движение мембраны. Отметим еще, что частота
высокочастотного колебания не кратна частоте низкочастотного.

\begin{enumerate}
\def\labelenumi{\arabic{enumi}.}
\setcounter{enumi}{1}
\tightlist
\item
  Если начальные условия заданы в виде u=\texttt{f(r)}, \texttt{u\_t=0}
  , то мы можем разложить начальное смещение \texttt{f(r)} в сумму
  элементарных «кирпичиков» \texttt{A\_m\ J\_a(k\_\{0m\}r)} , каждый из
  которых совершает колебания вида
  \texttt{\textbackslash{}cos(k\_\{0m\}ct)} на своей собственной
  частоте. Сложив фундаментальные колебания
\end{enumerate}

\begin{verbatim}
A_m J_0(k_{0m}r)\cos(k_{0m}ct)
\end{verbatim}

\begin{verbatim}
,

мы получим результирующие колебания, вызванные начальным смещением f(r).

#### ЗАДАЧИ

1. Подставьте  u(r, \theta, t) = U(r, \theta) T(t)  в волновое уравнение
\end{verbatim}

\texttt{u\_\{tt\}\ =\ c\^{}2\ \textbackslash{}left(\ u\_\{rr\}\ +\ \textbackslash{}frac\{1\}\{r\}\ u\_r\ +\ \textbackslash{}frac\{1\}\{r\^{}2\}\ u\_\{\textbackslash{}theta\textbackslash{}theta\}\ \textbackslash{}right)}

,

чтобы разделить его на два уравнения

\begin{verbatim}
\Delta U + \lambda^2 U = 0,
T'' + \lambda^2 c^2 T = 0.
\end{verbatim}

\texttt{2.\ Подстановкой\ \ \textquotesingle{}U\textbackslash{}left(r,\textbackslash{};\textbackslash{}theta\textbackslash{}right)=R\textbackslash{}left(r\textbackslash{}right)\textbackslash{}Theta\textbackslash{}left(\textbackslash{}theta\textbackslash{}right)\textquotesingle{}\ \ приведите\ задачу\ Гельмгольца}

\begin{verbatim}
\Delta U + \lambda^2 U = 0,
U(1, \theta) = 0
\end{verbatim}

к виду

\begin{verbatim}
r^2R'' + rR' + (\lambda^2r^2 - n^2)R = 0,
\end{verbatim}

\texttt{R(1)\ =\ 0,}
\texttt{\textbackslash{}Theta\textquotesingle{}\textquotesingle{}\ +\ n\^{}2\textbackslash{}Theta\ =\ 0.}

Константу разделения мы обозначили \texttt{n\^{}2} . Почему n может
принимать только значения \texttt{0,\ 1,\ 2,\ \textbackslash{}ldots} ?

\begin{enumerate}
\def\labelenumi{\arabic{enumi}.}
\setcounter{enumi}{2}
\tightlist
\item
  Решите задачу
\end{enumerate}

\begin{verbatim}
\begin{array}{llll} \text{(GU)} & u_{tt} = \Delta u, & 0 < r < 1, & 0 < t < \infty, \\ \text{(GU)} & u\left(1, \; \theta, \; t\right) = 0, & 0 < t < \infty, \\ \text{(HY)} & \begin{cases} u\left(r, \; \theta, \; \theta\right) = 1 - r^2, \\ u_t\left(r, \; \theta, \; \theta\right) = 0. \end{cases} & 0 \leqslant r \leqslant 1, \\ \end{array}
\end{verbatim}

\begin{enumerate}
\def\labelenumi{\arabic{enumi}.}
\setcounter{enumi}{3}
\tightlist
\item
  Решите задачу
\end{enumerate}

\begin{verbatim}
\begin{array}{lll} \text{(GU)} & u_{tt} = \Delta u, & 0 < r < 1, & 0 < t < \infty, \\ \text{(GU)} & u\left(1, & \theta, & t\right) = 0, & 0 <t < \infty, \\ \text{(NU)}} & \begin{cases} u\left(r, & \theta, & 0\right) = J_0\left(2, 4 \r\right), \\ u_t\left(r, & \theta, & 0\right) = 0. \end{cases} & 0 \leqslant r \leqslant 1, \\ \end{array}
\end{verbatim}

\begin{itemize}
\tightlist
\item
  \begin{enumerate}
  \def\labelenumi{\arabic{enumi}.}
  \setcounter{enumi}{4}
  \tightlist
  \item
    Решите задачу 4, если начальное смещение мембраны задано в виде
    \texttt{u(r,\ \textbackslash{}theta,\ 0)\ =\ J\_\{\textbackslash{}mathrm\{\ \}o\}(2,4r)\ 0.5\textbackslash{}\ J\_\{\textbackslash{}mathrm\{\ \}o\}(8,65r)\ +\ 0.25\textbackslash{}\ J\_\{\textbackslash{}mathrm\{\ \}o\}(14,93r)}
    . Какова наивысшая частота колебаний мембраны в этом случае?
  \end{enumerate}
\item
  \begin{enumerate}
  \def\labelenumi{\arabic{enumi}.}
  \setcounter{enumi}{5}
  \tightlist
  \item
    Постройте графики следующих функций на отрезке 0≤r≤1
  \end{enumerate}
\item
  \begin{enumerate}
  \def\labelenumi{\alph{enumi})}
  \tightlist
  \item
    \texttt{J\_\{o\}(5,52\ r)} ,
  \end{enumerate}
\item
  \begin{enumerate}
  \def\labelenumi{\arabic{enumi})}
  \setcounter{enumi}{5}
  \tightlist
  \item
    \texttt{J\_0(14,93\ \textbackslash{}\ r)} .
  \end{enumerate}
\end{itemize}

\section{Лекція 31. ЛАПЛАСІАН (ІНТУЇТИВНИЙ
ОПИС)}\label{ux43bux435ux43aux446ux456ux44f-31.-ux43bux430ux43fux43bux430ux441ux456ux430ux43d-ux456ux43dux442ux443ux457ux442ux438ux432ux43dux438ux439-ux43eux43fux438ux441}

ЦЕЛЬ ЛЕКЦИИ: Дать интуитивное описание лапласиана и привести его
выражения в различных системах координат. В наиболее важных системах
координат двух-и трехмерного пространства лапласиан записывается в виде

трехмерного пространства лапласнан записывается в виде

\begin{verbatim}
\Delta u = u_{xx} + u_{yy} \qquad \qquad \qquad \qquad \qquad \qquad \qquad \qquad \qquad
\end{verbatim}

Поскольку преобразования координат вызывают затруднения у некогорых
студентов, мы обсудим правила, которыми надлежит пользоваться при
выполнении таких преобразований.

Оператор Лапласа

(31.1)

\begin{verbatim}
\Delta = \frac{\partial^2}{\partial x^2} + \frac{\partial^2}{\partial y^2} + \frac{\partial^2}{\partial z^2},
\end{verbatim}

вероятно, самый важный оператор математической физики. Каков смысл этого
оператора и какое отношение имеет сумма трех частных производных второго
порядка к законам природы? Ответы на эти вопросы связаны с тем фактом,
что лапласиан функции позволяет оценить значение функции в точке через
значения функции в соседних точках. Кроме того, лапласиан можно считать
обобщением второй производной функции одной переменной на многомерный
случай.

Теперь мы переходим к основному свойству лапласиана, благодаря которому
он широко используется.

\section{Основное свойство двумерного оператора
А}\label{ux43eux441ux43dux43eux432ux43dux43eux435-ux441ux432ux43eux439ux441ux442ux432ux43e-ux434ux432ux443ux43cux435ux440ux43dux43eux433ux43e-ux43eux43fux435ux440ux430ux442ux43eux440ux430-ux430}

Если

\texttt{1.\ \ \textquotesingle{}\textbackslash{}Delta\ u\ \textgreater{}\ 0\textquotesingle{}\ \ в\ точке\ (x,\ y),\ то\ \textquotesingle{}u(x,\ y)\textquotesingle{}\ меньше\ «среднего\ значения\ функции\ в\ соседних\ точках»\ 1)\ (например,\ на\ окружности\ с\ центром\ в\ точке\ (x,\ y));}

\texttt{2.\ \ \textquotesingle{}\textbackslash{}Delta\ u\ =\ 0\textquotesingle{}\ \ в\ точке\ (x,\ y),\ то\ \textquotesingle{}u(x,\ y)\textquotesingle{}\ равна\ «среднему\ значению}

функции в соседних точках»;

\texttt{3.\ \ \textquotesingle{}\textbackslash{}Delta\ u\ \textless{}\ 0\textquotesingle{}\ \ в\ точке\ (x,\ y),\ то\ \textquotesingle{}u(x,\ y)\textquotesingle{}\ больше\ «среднего\ значения\ функции\ в\ соседних\ точках».}

Используя этот факт 2), дадим интерпретацию основных уравнений
математической физики.

\subsubsection{Интуитивный смысл некоторых законов
физики}\label{ux438ux43dux442ux443ux438ux442ux438ux432ux43dux44bux439-ux441ux43cux44bux441ux43b-ux43dux435ux43aux43eux442ux43eux440ux44bux445-ux437ux430ux43aux43eux43dux43eux432-ux444ux438ux437ux438ux43aux438}

\texttt{Согласно\ уравнению\ теплопроводности\ \ \textquotesingle{}u\_t\ =\ \textbackslash{}alpha\^{}2\ \textbackslash{}Delta\ u\textquotesingle{}\ \ для\ температуры\ (или\ концентрации)\ и,\ скорость\ изменения\ температуры\ \ \textquotesingle{}u\_t\textquotesingle{}\ \ пропорциональна\ величине\ \ \textquotesingle{}\textbackslash{}Delta\ u\textquotesingle{}\ .\ Значит,\ если\ температура\ в\ точке\ температура\ на\ окружности,\ окружающей\ данную\ точку,\ то\ температура\ в\ данной\ точке\ будет\ возрастать.}

\pandocbounded{\includegraphics[keepaspectratio]{images/_page_232_Picture_11.jpeg}}

\pandocbounded{\includegraphics[keepaspectratio]{images/_page_232_Picture_12.jpeg}}

\texttt{РИС.\ 31.1.\ Интуитивное\ описание\ лапласиана\ \ \textquotesingle{}\textbackslash{}Delta\ u\textquotesingle{}\ :\ a—\ среднее\ значение\ функции\ u\ на\ окружности\ с\ центром\ в\ точке\ p\ равно\ \textquotesingle{}u\ (p)\textquotesingle{},\ \ \textquotesingle{}\textbackslash{}Delta\ u\textquotesingle{}\ \ (p)\ =\ 0;\ \ \textquotesingle{}\textbackslash{}delta\textquotesingle{}\ —\ среднее\ значение\ функции\ u\ на\ окружности\ с\ центром\ в\ точке\ p\ больше,\ чем\ \textquotesingle{}u\ (p)\textquotesingle{},\ \ \textquotesingle{}\textbackslash{}Delta\ u\textquotesingle{}\ \ (p)\ \textgreater{}\ 0;\ \ \textquotesingle{}\textbackslash{}delta\textquotesingle{}\ —\ среднее\ значение\ функции\ u\ на\ окружности\ с\ центром\ в\ точке\ p\ меньше,\ чем\ \textquotesingle{}u\ (p)\textquotesingle{},\ \ \textquotesingle{}\textbackslash{}Delta\ u\textquotesingle{}\ \ (p)\ \textless{}\ 0.}

\texttt{Согласно\ волновому\ уравнению\ \ \textquotesingle{}u\_\{tt\}\ =\ \textbackslash{}alpha\^{}2\ \textbackslash{}Delta\ u\textquotesingle{}\ \ для\ смещения\ мембраны,\ ускорение\ точки\ мембраны\ \ \textquotesingle{}u\_\{tt\}\textquotesingle{}\ \ (или\ сила,\ действующая\ на\ точку)\ пропорционально\ величине\ \ \textquotesingle{}\textbackslash{}Delta\ u\textquotesingle{}\ .\ Значит,\ точка\ мембраны}

2) Этот факт (теорема о среднем) не зависит от размерности пространства.
--- Прим. ред.

1) Под «средним значением функции в соседних точках» понимается среднее
значение функции либо по окружности, либо по кругу с центром в точке (x,
y). --- Прим. ред.

ускоряется вверх (сила направлена вверх), если ее смещение (по высоте)
меньше, чем среднее смещение соседних точек.

\texttt{Если\ функция\ удовлетворяет\ уравнению\ Лапласа\ \ \textquotesingle{}\textbackslash{}Delta\ u=0\textquotesingle{}\ ,\ то\ ее\ вначение\ всегда\ совпадает\ со\ средним\ значением.\ Например,\ натянутая\ неподвижная\ резиновая\ мембрана\ удовлетворяет\ уравнению\ Лапласа,\ следовательно,\ смещение\ мембраны\ в\ любой\ точке\ равно\ среднему\ смещению\ мембраны\ на\ окружности\ с\ центром\ в\ этой\ точке.\ Рис.\ 31.1\ иллюстрирует\ это\ свойство\ лапласиана.}

\texttt{У\ равнение\ Пуассона\ \ \textquotesingle{}\textbackslash{}Delta\ u\ =\ f\textquotesingle{}\ ,\ где\ f\ —\ функция,\ зависящая\ от\ пространственных\ переменных\ (может\ быть\ и\ константой),\ описывает\ значительное\ число\ различных\ явлений.}

\texttt{-\ 1.\ \ \textquotesingle{}\textbackslash{}Delta\ u\ =\ -\textbackslash{}rho\textquotesingle{}\ \ описывает\ потенциал\ электростатического\ поля,\ если\ \ \textquotesingle{}\textbackslash{}rho\textquotesingle{}\ —плотность\ статических\ зарядов.\ Зная\ основное\ свойство\ лапласиана,\ что\ бы\ вы\ могли\ сказать\ о\ свойствах\ потенциального\ поля?}
\texttt{-\ 2.\ \ \textquotesingle{}\textbackslash{}Delta\ u\ =\ -g(x,\ y)\textquotesingle{}\ \ описывает\ стационарное\ распределение\ температуры\ от\ теплового\ источника\ \textquotesingle{}g(x,\ y)\textquotesingle{}.\ Если\ \textquotesingle{}g(x,\ y)\textquotesingle{}\ положительна\ в\ точке,\ значит,\ в\ этой\ точке\ тепло\ выделяется,\ если\ же\ \textquotesingle{}g(x,\ y)\textquotesingle{}\ отрицательна,\ то\ тепло\ в\ этой\ точке\ поглощается.}
- 3. Δu + λu = 0 это уравнение Гельмгольца (приведенное волновое
уравнение), которое описывает (помимо многого другого) пространственную
часть собственных колебаний мембраны. Это уравнение возникает при
разделении пространственных и временной переменных в волновом уравнении
и уравнении теплопроводности.

Пока что мы говорили о физическом смысле лапласиана.

\begin{verbatim}
\Delta u = u_{xx} + u_{yy}
\end{verbatim}

\texttt{(двумерный\ случай),\ \ \textquotesingle{}\textbackslash{}Delta\ u\ =\ u\_\{xx\}\ +\ u\_\{yy\}\ +\ u\_\{zz\}\textquotesingle{}\ \ (трехмерный\ случай).}

Однако во многих задачах возникает необходимость знать выражение для
лапласнана в других координатах. Если, например, граница области ---
окружность, то естественно перейти к полярным координатам
\texttt{(r,\ \textbackslash{}theta)} , а если область трехмерна, а ее
граница --- сфера, то следует перейти к сферическим координатам
\texttt{(r,\ \textbackslash{}theta,\ \textbackslash{}phi)} . Итак, нам
предстоит ответить на вопрос: как выглядит лапласиан в различных
системах координат?

Прежде чем приступить к ответу на этот вопрос, напомним кратко о пяти
основных системах координат в двумерном и трехмерном пространствах:

декартова прямоугольная система в двумерном пространстве, декартова
прямоугольная система в трехмерном пространстве, полярные координаты в
двумерном пространстве,

цилиндрические координаты в трехмерном пространстве, сферические
координаты в трехмерном пространстве.

Полярные координаты определяются соотношениями

\begin{verbatim}
\begin{array}{|c|c|c|c|}
\end{verbatim}

\begin{verbatim}
\hлінія r^2 = x^2 + y^2, & x = r \cos \theta, \\
\end{verbatim}

\begin{verbatim}
\text{tg } \theta = y/x, & y = r \sin \theta.
\end{array}
\end{verbatim}

\pandocbounded{\includegraphics[keepaspectratio]{images/_page_234_Picture_4.jpeg}}

Цилиндрические координаты связаны с декартовыми следующими
соотношениями:

\begin{verbatim}
r^2 = x^2 - y^2
\end{verbatim}

\begin{verbatim}
,  x = r \cos \theta ,  
 tg \theta = y/x ,  tg \theta = y/x ,  tg \theta = z ,  tg \theta ,  tg \theta = z

![](images/_page_234_Picture_7.jpeg)

Сферические координаты определяются следующим образом:
\end{verbatim}

r\^{}2 = x\^{}2 + y\^{}2 + z\^{}2

,
\texttt{x\ =\ r\ \textbackslash{}sin\ \textbackslash{}varphi\ \textbackslash{}cos\ \textbackslash{}theta}
, \texttt{\textbackslash{}cos\ \textbackslash{}varphi\ =\ z/r} , или
\texttt{y\ =\ r\ \textbackslash{}sin\ \textbackslash{}varphi\ \textbackslash{}sin\ \textbackslash{}theta}
, \texttt{tg\ \textbackslash{}theta\ =\ y/x} ,
\texttt{z\ =\ r\ \textbackslash{}cos\ \textbackslash{}varphi} .

\pandocbounded{\includegraphics[keepaspectratio]{images/_page_235_Picture_2.jpeg}}

\paragraph{Преобразование
координат}\label{ux43fux440ux435ux43eux431ux440ux430ux437ux43eux432ux430ux43dux438ux435-ux43aux43eux43eux440ux434ux438ux43dux430ux442}

Теперь мы подготовлены к преобразованию координат. В качестве примера
посмотрим, как преобразуется лапласиан при переходе к полярным
координатам. Если мы разберемся в этом, то Р И

\pandocbounded{\includegraphics[keepaspectratio]{images/_page_235_Picture_5.jpeg}}

С 312 Функционатьная зависимость функции u от переменных
\texttt{(r,\ \textbackslash{}theta)} и (x, y)

\pandocbounded{\includegraphics[keepaspectratio]{images/_page_235_Picture_7.jpeg}}

РИС. 31 3 Функциональная зависимость производной \texttt{u\_x} от
переменных (r, 0) и (x, y).

сможем легко переходить от одной системы координат к другой. Итак, у нас
есть лапласиан

\begin{verbatim}
\Delta u = u_{xx} + u_{yy}.
\end{verbatim}

Чтобы найти лапласиан в новой системе координат, обратимся к диаграмме,
изображенной на рис. 31.2. Цель диаграммы --- точно показать, что от
чего зависит. Например, из этой диаграммы видно, что функция u зависит
от новых переменных r и \texttt{\textbackslash{}theta} .

Каждая из этих переменных в свою очередь зависит от x и y (соединяющие
линии указывают на эту зависимость).

Мы могли бы нарисовать диаграмму, в основании которой находились бы r и
\texttt{\textbackslash{}theta} , но поскольку нам предстоит вычислять
\texttt{u\_\{xx\}} и \texttt{u\_\{yy\}} , то больше подходит именно эта
диаграмма. Как же теперь выразить \texttt{u\_\{xx\}} и
\texttt{u\_\{yy\}} через r и \texttt{\textbackslash{}theta} ?

На самом деле сделать это очень легко. Сначала найдем \texttt{u\_x} и
\texttt{u\_y} . Мы найдем \texttt{u\_x} , складывая пути на диаграмме,
ведущие от u к x. В данном случае существуют только два таких пути, так
что

\begin{verbatim}
u_x = u_r r_x + u_\theta \theta_x = u_r (\cos \theta) - u_\theta (\sin \theta/r).
\end{verbatim}

По правилам дифференцирования сложных функций следует перемножить
производные, написанные у ребер каждого пути, а затем сложить их.
Аналогично

\begin{verbatim}
u_y = u_r r_y + u_\theta \theta_y = u_r \sin \theta + u_\theta (\cos \theta/r).
\end{verbatim}

Обратимся теперь ко вторым производным. Найдем \texttt{u\_\{xx\}} (в
качестве упражнения предоставим читателю возможность найти
\texttt{u\_\{yy\}} ).

``Сначала, однако, построим другую диаграмму (см. рис. 31.3). Читателю,
вероятно, ясно, что если u зависит от r и \texttt{\textbackslash{}theta}
(которые в свою очередь зависят от x и y), то \texttt{u\_x} так же
(вообще говоря) зависит от тех же переменных. Таким образом,

\begin{verbatim}
u_{xx} = (u_x)_x =
\end{verbatim}

= (\texttt{u\_x})\emph{r r}\лямбда + (\texttt{u\_x})\_\theta \theta\_x =

\begin{verbatim}
= (u_r \cos \theta - u_\theta \sin \theta/r)_r \cos \theta + (u_r \cos \theta - u_\theta \sin \theta/r)_\theta (\sin \theta/r) =
= ((u_{rr} \cos \theta - u_{r\theta} \sin \theta/r + u_\theta \sin \theta/r^2) \cos \theta +
+ (u_{r\theta} \cos \theta - u_r \sin \theta - u_{\theta\theta} \sin \theta/r - u_\theta \cos \theta/r) (-\sin \theta/r).
\end{verbatim}

Аналогично

\begin{verbatim}
u_{yy} = (u_{rr}\sin\theta + u_{r\theta}\cos\theta/r - \cos\theta/r^2)\sin\theta + + (u_{r\theta}\sin\theta + u_{r}\cos\theta + u_{\theta\theta}\cos\theta/r - u_{\theta}\sin\theta/r)(\cos\theta/r).
\end{verbatim}

Складывая \texttt{u\_\{xx\}} и \texttt{u\_\{yy\}} , получаем
окончательный результат

\begin{verbatim}
\Delta u = u_{xx} + u_{yy} = u_{rr} + \frac{1}{r} u_r + \frac{1}{r^2} u_{\theta\theta},
\end{verbatim}

т. е. выражение для лапласиана в полярных координатах. Его смысл такой
же, как и у лапласиана в декартовых координатах, но форма записи другая.
К сожалению, теперь это оператор с переменными коэффициентами, так что
уравнение, содержащее лапласианы в полярных координатах, решать труднее.

Аналогично можно преобразовать и трехмерный лапласиан

\begin{verbatim}
\Delta u = u_{xx} + u_{yy} + u_{zz}.
\end{verbatim}

В результате получаем в цилиндрических координатах

\begin{verbatim}
\Delta u = u_{rr} + \frac{1}{r} u_r + \frac{1}{r^2} u_{\theta\theta} + u_{zz},
\end{verbatim}

в сферических координатах

\begin{verbatim}
\Delta u = u_{rr} + \frac{2}{r} u_r + \frac{1}{r^2} u_{\varphi\varphi} + \frac{\operatorname{ctg} \varphi}{r^2} u_{\varphi} + \frac{1}{r^2 \sin^2 \varphi} u_{\theta\theta}.
\end{verbatim}

\paragraph{ЗАМЕЧАНИЯ}\label{ux437ux430ux43cux435ux447ux430ux43dux438ux44f-3}

\begin{enumerate}
\def\labelenumi{\arabic{enumi}.}
\item
  Только в декартовой системе координат лапласиан обладает постоянными
  коэффициентами. Этим объясняется, почему задачи в других системах
  координат решать труднее. Тем не менее к этим уравнениям с переменными
  коэффициентами все еще применим метод разделения переменных, правда,
  получающиеся при этом обыкновенные дифференциальные уравнения тоже
  будут с переменными коэффициентами. Так получаются уравнение Бесселя,
  уравнение Лежандра и другие, так называемые классические уравнения
  математической физики. Такие уравнения решаются обычно методом
  разложения в степенные ряды.
\item
  Кроме упомянутых здесь, существует много др гих координатных систем
  1). Например, задача о колебаниях эллиптической мембраны требует
  представления лапласиана в эллиптической системе координат. Чтобы
  сделать это, требуется несколько иной подход: преобразовать эллипс в
  круг, решить получившуюся задачу в новых (полярных) координатах, и
  затем полученное решение преобразовать к исходным координатам 2).
\end{enumerate}

2) См. Тихонов А. Н., Самарский А. А. Уравнения математической физики.
--- М.: Наука, 1972 и др. издания. --- Прим. ред.

2) Целесообразнее непосредственно перейти к эллиптическим координатам.
--- Прим. ред.

\paragraph{ЗАЛАЧИ}\label{ux437ux430ux43bux430ux447ux438}

\texttt{-\ 1.\ Как\ записать\ волновое\ уравнение\ \ \textquotesingle{}u\_\{tt\}\ =\ \textbackslash{}alpha\^{}2\ \textbackslash{}Delta\ u\textquotesingle{}\ \ в\ сферических\ координатах,\ если\ вам\ известно,\ что\ решение\ зависит\ только\ OT\ r\ H\ t?}
\texttt{-\ 2.\ Как\ записать\ волновое\ уравнение\ \ \textquotesingle{}u\_\{tt\}\ =\ \textbackslash{}alpha\^{}2\ \textbackslash{}Delta\ u\textquotesingle{}\ \ в\ полярных\ координатах,\ если\ известно,\ что\ решение\ зависит\ только\ от\ г\ -\ 3.\ Как\ записать\ уравнение\ Лапласа\ \ \textquotesingle{}\textbackslash{}Delta\ u\ =\ 0\textquotesingle{}\ \ в\ полярных\ координатах,\ если\ известно,\ что\ u\ зависит\ только\ от\ r?\ Каково\ решение\ этого\ уравнения?\ Эти\ функции\ называются\ потенциалами\ с\textquotesingle{}криговой\ симметрией.}
- 4. Как записать уравнение Лапласа в сферических координатах, если
решение и зависит только от г? Можете ли вы найти решения этого
уравнения? Эти функции называются сферически-симметричными потенциалами,
- 5. Что вы можете сказать о поверхности \texttt{u(x,\ y)\ =\ xy\^{}2}
\texttt{-\ 6.\ Преобразуйте\ трехмерный\ лапласиан\ \ \textquotesingle{}\textbackslash{}Delta\ u\ =\ u\_\{xx\}\ +\ u\_\{yy\}\ +\ u\_\{zz\}\textquotesingle{}\ \ к\ сферическим\ координатам.}

\section{Лекція 32. ЗАГАЛЬНІ ВЛАСТИВОСТІ КРАЙОВИХ
ЗАДАЧ}\label{ux43bux435ux43aux446ux456ux44f-32.-ux437ux430ux433ux430ux43bux44cux43dux456-ux432ux43bux430ux441ux442ux438ux432ux43eux441ux442ux456-ux43aux440ux430ux439ux43eux432ux438ux445-ux437ux430ux434ux430ux447}

ЦЕЛЬ ЛЕКЦИИ: Объяснить, как возникают уравнения с частными производными,
не содержащие производных по времени. В отличие от гиперболического
волнового уравнения или параболического уравнения теплопроводности эти
уравнения не требуют начальных условий. Для них нужны только граничные
условия. По этой причине задачи для таких уравнений называются краевыми
задачами

Укажем три наиболее важных типа граничных условий.

\begin{itemize}
\tightlist
\item
  Граничное условие первого рода (условие Дирихле). Граничное условие
  второго рода (условие Неймана).
\item
  \begin{enumerate}
  \def\labelenumi{\arabic{enumi}.}
  \setcounter{enumi}{2}
  \tightlist
  \item
    Граничное условие третьего рода (условие Робена).
  \end{enumerate}
\item
  В лекции проводятся примеры и пояснения этих граничных условий.
\end{itemize}

До сих пор мы изучали явления, которые развивались в пространстве и во
времени. Существует, однако, много явлений, которые не изменяются с
течением времени. Эти явления в большинстве случаев описываются краевыми
эллиптическими задачами, и цель нашей лекции---познакомить читателя с
такого типа задачами.

Можно указать два случая, когда в уравнениях математической физики не
содержится время

\begin{enumerate}
\def\labelenumi{\arabic{enumi}.}
\tightlist
\item
  Стационарные задачи
\end{enumerate}

2 Задачи, в которых зависимость от времени входит в виде временного
множителя.

Сначала давайте рассмотрим стационарные задачи.

\subsubsection{Стационарные
задачи}\label{ux441ux442ux430ux446ux438ux43eux43dux430ux440ux43dux44bux435-ux437ux430ux434ux430ux447ux438}

Предположим, что нас интересует стационарное (при
\texttt{t\ \textbackslash{}to\ \textbackslash{}infty} ) решение
уравнения теплопроводности

\begin{verbatim}
u_i = \alpha^2 \Delta u
\end{verbatim}

\begin{verbatim}
.

Очевидно, если решение не изменяется во времени, то  u_t = 0  и равнение теплопроводности переходит в уравнение Лапласа
\end{verbatim}

\texttt{\textbackslash{}Delta\ u\ =\ 0}

\begin{verbatim}
Проиллюстрируем подробно идею стационарности на примере. Рассмотрим задачу
\end{verbatim}

\texttt{\textbackslash{}begin\{array\}\{ll\}\ (\textbackslash{}text\{UBP\})\ \&\ u\_t\ =\ u\_\{xx\}\ +\ \textbackslash{}sin\{(\textbackslash{}pi\ x)\},\ \&\ 0\ \textless{}\ x\ \textless{}\ 1,\ \&\ 0\ \textless{}\ t\ \textless{}\ \textbackslash{}infty,\ \textbackslash{}\textbackslash{}\ (32.1)\ \&\ (\textbackslash{}Gamma\ \textbackslash{}text\{Y\})\ \&\ \textbackslash{}begin\{cases\}\ u(0,\ t)\ =\ 0,\ \textbackslash{}\textbackslash{}\ u(1,\ t)\ =\ 0,\ \textbackslash{}\textbackslash{}\ u(x,\ 0)\ =\ \textbackslash{}sin\{(3\textbackslash{}pi\ x)\},\ \&\ 0\ \textbackslash{}le\ x\ \textbackslash{}le\ \textbackslash{}le\ 1.\ \textbackslash{}end\{cases\}\ \textbackslash{}\textbackslash{}\ \textbackslash{}end\{array\}}

\begin{verbatim}
Чтобы найти стационарное решение  u\left(x,\infty\right)  (если оно существует), положим  u_{t}\!=\!0  и решим краевую задачу
\end{verbatim}

\texttt{\textbackslash{}begin\{cases\}\ \textbackslash{}frac\{d\^{}\{2\}u\}\{dx\^{}\{2\}\}\ =\ -\textbackslash{}sin(\textbackslash{}pi\ x),\ \&\ 0\ \textless{}\ x\ \textless{}\ 1,\ \textbackslash{}\textbackslash{}\ u(0)\ =\ 0,\ \textbackslash{}\textbackslash{}\ u(1)\ =\ 0.\ \textbackslash{}end\{cases\}}

\begin{verbatim}
Решение этой задачи имеет вид
\end{verbatim}

\texttt{u(x,\ \textbackslash{}infty)\ =\ \textbackslash{}frac\{1\}\{\textbackslash{}pi\^{}2\}\ \textbackslash{}sin\{(\textbackslash{}pi\ x)\}.}

\begin{verbatim}
Если решить задачу (32.1), то ее решение в начальный момент будет совпадать с  \sin{(3\pi x)} , а затем все в большей и большей степен походить на
\end{verbatim}

\texttt{\textbackslash{}frac\{1\}\{\textbackslash{}pi\^{}2\}\textbackslash{}sin\{(\textbackslash{}pi\ x)\}.}

\begin{verbatim}
Для некоторых задач может не существовать стационарное решение (не зависящее от времени), а для других стационарным решением может оказаться временная синусоида. Следовательно, не всегда справедливо, что производные  u_t, u_{tt} \dots  равны нулю. Поэтому мы, как правило, должны иметь некоторую информацию о физическом смысле задачи.

# Выделение временной компоненты в гиперболических и параболических задачах

В задаче о колебаниях мембраны

(УЧП)
\end{verbatim}

\texttt{u\_\{tt\}\ =\ \textbackslash{}Delta\ u,\ \textbackslash{}quad\ 0\ \textbackslash{}le\ r\ \textless{}\ 1,}

\begin{verbatim}
(ГУ)  u = 0, \quad \text{на окружности,}   
(НУ)
\end{verbatim}

\texttt{\textbackslash{}begin\{cases\}\ u(r,\ \textbackslash{}theta,\ 0)\ =\ f(r,\ \textbackslash{}theta),\ \textbackslash{}\textbackslash{}\ u\_t(r,\ \textbackslash{}theta,\ 0)\ =\ g(r,\ \textbackslash{}theta)\ \textbackslash{}end\{cases\}}

\begin{verbatim}
мы искали решение в виде  u\left(r,\,\theta,\,t\right)=U\left(r,\,\theta\right)T\left(t\right) . При этом функция  U\left(r,\,\theta\right)  должна быть решением краевой задачи для уравнения Гельмгольца
\end{verbatim}

\texttt{\textbackslash{}begin\{array\}\{ll\}\ (\textbackslash{}mathrm\{Y\}\textbackslash{}mathrm\{H\}\textbackslash{}mathrm\{\textbackslash{}Pi\})\ \&\ \textbackslash{}Delta\ U\ +\ \textbackslash{}lambda\^{}2\ U\ =\ 0,\ \textbackslash{}\textbackslash{}\ (\textbackslash{}mathrm\{H\}\textbackslash{}mathrm\{Y\})\ \&\ U\textbackslash{}left(1,\textbackslash{},\textbackslash{}theta\textbackslash{}right)\ =\ 0.\ \textbackslash{}end\{array\}}

\begin{verbatim}
Такая ситуация характерна для теории уравнений с частными производными, где решение зачастую представляет собой произведение пространственной части  U\left(r,\theta\right)  на временную часть  T\left(t\right) . Фактически мы получим то же самое уравнение Гельмгольца, если попытаемся выделить временную компоненту в уравнении теплопроводности
\end{verbatim}

\texttt{u\_i\ =\ \textbackslash{}Delta\ u}

.

При решении краевых задач чаще всего используются три типа граничных
условий. На них мы сейчас остановимся.

\subsubsection{Три основных типа граничных условий в краевых
задачах}\label{ux442ux440ux438-ux43eux441ux43dux43eux432ux43dux44bux445-ux442ux438ux43fux430-ux433ux440ux430ux43dux438ux447ux43dux44bux445-ux443ux441ux43bux43eux432ux438ux439-ux432-ux43aux440ux430ux435ux432ux44bux445-ux437ux430ux434ux430ux447ux430ux445}

Краевая задача с граничными условиями первого рода (задача Дирихле)

Требуется найти решение уравнения в некоторой области пространства,
которое принимает на границе области заданные значения. В качестве
примера можно привести задачу о нахождении стационарного распределения
температуры внутри области,

если задана температура на границе этой области. Другой пример: найти
распределение потенциала внутри области, если известен потенциал на
границе. Ниже приведен еще ряд примеров.

\paragraph{ПРИМЕРЫ ЗАДАЧ
ДИРИХЛЕ}\label{ux43fux440ux438ux43cux435ux440ux44b-ux437ux430ux434ux430ux447-ux434ux438ux440ux438ux445ux43bux435}

Рассмотрим уравнение Лапласа внутри круга с заданными значениями решения
на границе (см. рис. 32.1):

(УЧП)

\begin{verbatim}
u_{rr} + \frac{1}{r}u_r + \frac{1}{r^2}u_{\theta\theta} = 0, \quad 0 < r < 1,
\end{verbatim}

\texttt{(ГУ)\ \ \textquotesingle{}u(1,\ \textbackslash{}theta)\ =\ \textbackslash{}sin\ \textbackslash{}theta,\ \textbackslash{}quad\ 0\ \textbackslash{}le\ \textbackslash{}theta\ \textless{}\ 2\textbackslash{}pi.\textquotesingle{}}

В качестве еще одного примера можно указать внешнюю задачу Дирихле, для
которой решение уравнения Лапласа пишется вне круга, а граничные
значения заданы по окружности (см. рис. 32.2).

\pandocbounded{\includegraphics[keepaspectratio]{images/_page_241_Picture_7.jpeg}}

РИС, 32.1. Внутренняя задача Дирихле.

\pandocbounded{\includegraphics[keepaspectratio]{images/_page_241_Picture_9.jpeg}}

РИС. 32.2. Внешняя задача Дирихле.

\section{Краевые задачи с граничным условием второго рода (задача
Неймана)}\label{ux43aux440ux430ux435ux432ux44bux435-ux437ux430ux434ux430ux447ux438-ux441-ux433ux440ux430ux43dux438ux447ux43dux44bux43c-ux443ux441ux43bux43eux432ux438ux435ux43c-ux432ux442ux43eux440ux43eux433ux43e-ux440ux43eux434ux430-ux437ux430ux434ux430ux447ux430-ux43dux435ux439ux43cux430ux43dux430}

Требуется найти решение уравнения в некоторой области пространства, на
границе которой задана внешняя нормальная производная ди/дп (которая
пропорциональна втекающему потоку). Это общая задача и для стационарной
теплопроводности, и для электростатики, если на границе задан поток
(тепла, электронов и т. д.).

Например, предположим, что тепловой поток на границе круга

изменяется по закону

\begin{verbatim}
\frac{\partial u}{\partial r} = \sin \theta.
\end{verbatim}

Тогда стационарное распределение температуры внутри круга является
решением краевой задачи

\begin{verbatim}
\begin{cases} \Delta u = 0, & 0 < r < 1, \\ \frac{\partial u}{\partial r} = \sin \theta, & r = 1, & 0 \le \theta < 2\pi. \end{cases}
\end{verbatim}

\texttt{В\ этой\ задаче\ *поток*\ тепла\ (кал/см·с)\ через\ границу\ направлен\ енутрь\ при\ \ \textquotesingle{}0\ \textbackslash{}leqslant\ \textbackslash{}theta\ \textless{}\ \textbackslash{}pi\textquotesingle{}\ \ и\ наружу\ при\ \ \textquotesingle{}\textbackslash{}pi\ \textbackslash{}leqslant\ \textbackslash{}theta\ \textless{}\ 2\textbackslash{}pi\textquotesingle{}\ .}

Однако, поскольку полный поток

\begin{verbatim}
\int_{0}^{2\pi} \frac{\partial u}{\partial r} d\theta = \int_{0}^{2\pi} \sin\theta \, d\theta = 0
\end{verbatim}

(условие, которое всегда должно выполняться в задаче Неймана), можно
утверждать, что температура каждой точки внутри круга не изменяется во
времени.

\pandocbounded{\includegraphics[keepaspectratio]{images/_page_242_Picture_9.jpeg}}

\textbf{РИС.} 32.3. Поток тепла в задаче Неймана. Температура в каждой
точке \textbf{мостоянна}, поскольку полное количество тепла, втекающего
в точку и вытежающего из него, равно нулю.

Другими словами, задача Неймана (для уравнения Лапласа) имеет смысл
только в том случае, когда полный поток тепла через границу равен нулю.
Математически это значит, что на границе должно выполняться соотношение

\begin{verbatim}
\int_{C} \frac{\partial u}{\partial n} = 0,
\end{verbatim}

иначе задача не будет иметь решения. Например, внутренняя задача Неймана

\begin{verbatim}
\begin{cases} \Delta u = 0, & 0 < r < 1, \\ \frac{\partial u}{\partial r}(1, \theta) = 1, & r = 1, & 0 \le \theta < 2\pi, \end{cases}
\end{verbatim}

не имеет физического смысла, поскольку постоянный единичный поток внутрь
области не может обеспечить стационарность решения.

Еще одна особенность задачи Неймана, которая отличает ее от других
граничных задач, это неединственность решения. Другими словами, задача
Неймана

\begin{verbatim}
\begin{cases} \Delta u = 0, & 0 < r < 1, \\ \frac{\partial u}{\partial r}(1, \theta) = \cos(2\theta), & r = 1, & 0 \le \theta < 2\pi, \end{cases}
\end{verbatim}

имеет бесконечное множество решений \texttt{u(r,\textbackslash{}theta)}
(заметим, что общий поток через границу равен нулю). Если у нас есть
одно решение, то, добавляя к нему константы, можно пайти и все другие
решения. Например, одно из решений нашей задачи Неймана имеет вид

\begin{verbatim}
u(r,\theta)=r^2\cos{(2\theta)}.
\end{verbatim}

Если прибавить к нему произвольную константу, то получится другое
решение. По этой причине для выделения единственного решения нужно иметь
дополнительную информацию (например, знать значение решения в одной
точке).

\section{\texorpdfstring{\textbf{Краевые задачи с граничным условием}
третьего
рода}{Краевые задачи с граничным условием третьего рода}}\label{ux43aux440ux430ux435ux432ux44bux435-ux437ux430ux434ux430ux447ux438-ux441-ux433ux440ux430ux43dux438ux447ux43dux44bux43c-ux443ux441ux43bux43eux432ux438ux435ux43c-ux442ux440ux435ux442ux44cux435ux433ux43e-ux440ux43eux434ux430}

Требуется найти такое решение уравнения в некоторой области
пространства, которое удовлетворяет на границе условию вида

\begin{verbatim}
\frac{\partial u}{\partial n} + h (u - g) = 0,
\end{verbatim}

где h---заданная константа, а g---заданная функция, которая, вообще
говоря, меняется вдоль границы. Обратим внимание на то, что новое
граничное условие является как бы смесью граничных условий первых двух
типов. Более привычная форма

этого условия имеет вид

\begin{verbatim}
\frac{\partial u}{\partial n} = -h (u - g).
\end{verbatim}

Согласно этому граничному условию, поток, втекающий в область через
границу, пропорционален разности между температурой u и некоторой
заданной температурой g. Это означает, что (при h\textgreater0)

\begin{enumerate}
\def\labelenumi{\arabic{enumi})}
\item
  если температура \emph{и} на границе выше температуры окружающей среды
  на границе \emph{g}, то тепло вытекает из области;
\item
  если и меньше температуры g, то тепло втекает в область. Конечно, это
  не что иное, как закон теплообмена Ньютона, а граничные условия всех
  трех типов совершенно естественны для стационарных задач
  теплопроводности.
\end{enumerate}

\paragraph{ПРИМЕР}\label{ux43fux440ux438ux43cux435ux440}

Предположим, что вне единичного круга температура распределена по закону
\texttt{g(\textbackslash{}theta)\ =\ \textbackslash{}sin\ \textbackslash{}theta}
. В этом случае поток тепла через границу должен задаваться в виде

\begin{verbatim}
\frac{\partial u}{\partial r} = -h(u - \sin \theta), \quad r = 1, \quad 0 \le \theta < 2\pi.
\end{verbatim}

Следовательно, стационарное распределение температуры внутри круга
является решением краевой задачи

(32.2)

\begin{verbatim}
\begin{array}{ll} (\text{UBP}) & \Delta u = 0, \quad 0 < r < 1, \\ (\Gamma \text{Y}) & \frac{\partial u}{\partial r} + h \left( u - \sin \theta \right) = 0, \quad 0 \leqslant \theta < 2\pi. \end{array}
\end{verbatim}

Константа h является физическим параметром и численно равна количеству
тепла, протекающему через границу при разности температур в один градус.
Если h велика, то при разности температур в один градус тепловой поток
велик и решение становится похожим на решение задачи Дирихле с граничным
условием f = \texttt{\textbackslash{}sin\ \textbackslash{}theta} . С
другой стороны, если f = 0, то граничное условне переходит в условие f f
f f f f f f f f

\begin{verbatim}
\frac{\partial u}{\partial r} = 0.
\end{verbatim}

Читатель может представить, как изменяются решения задачи (32.2), когда
параметр h пробегает диапазон от 0 до \texttt{+\textbackslash{}infty} .
Если h=0, то граница теплоизолирована и, следовательно, решение
\texttt{u(r,\textbackslash{}theta)} является константой (решение в
данном случае не единственное, так что любая константа является
решением). Если h становится все больше и больше, то решение становится
все всявие и больше похожим на решение задачи Дирихле с гранич-

\texttt{ным\ условием\ \ \textquotesingle{}u=\textbackslash{}sin\textbackslash{}theta\textquotesingle{}\ .\ Если\ h\textgreater{}0\ и\ стремится\ к\ нулю,\ то\ и\ решение\ также\ будет\ стремиться\ к\ тождественному\ нулю\ (поскольку\ нуль—это\ среднее\ значение\ \ \textquotesingle{}\textbackslash{}sin\textbackslash{}theta\textquotesingle{}\ \ на\ границе).}

\paragraph{ЗАДАЧИ}\label{ux437ux430ux434ux430ux447ux438-2}

\begin{enumerate}
\def\labelenumi{\arabic{enumi}.}
\tightlist
\item
  Можете ли вы из интуптивных соображений решить задачу Дирихле?
\end{enumerate}

\begin{verbatim}
\begin{array}{lll} (\text{UCH}\Pi) & \Delta u = 0, & 0 < r < 1, \\ (\text{GU}) & u(1, \theta) = \sin \theta, & 0 \leqslant \theta < 2\pi. \end{array}
\end{verbatim}

\begin{enumerate}
\def\labelenumi{\arabic{enumi}.}
\setcounter{enumi}{1}
\tightlist
\item
  Имеет ли задача Неймана
\end{enumerate}

\begin{verbatim}
(\forall \Pi)
\end{verbatim}

\texttt{\textbackslash{}Delta\ u\ =\ 0} ,
\texttt{0\ \textless{}\ r\ \textless{}\ 1} ,
\texttt{(\textbackslash{}Gamma\ \textbackslash{}forall)}
\texttt{\textbackslash{}frac\{\textbackslash{}partial\ u\}\{\textbackslash{}partial\ r\}\ =\ \textbackslash{}sin\^{}2\ \textbackslash{}theta}

решение внутри круга?

\begin{enumerate}
\def\labelenumi{\arabic{enumi}.}
\setcounter{enumi}{2}
\tightlist
\item
  Попытайтесь представить, как выглядит решение задачи
\end{enumerate}

(УЧП)

\begin{verbatim}
\Delta u = 0
\end{verbatim}

\begin{verbatim}
,  0 < r < 1 ,  
(ГУ)  \frac{\partial u}{\partial r} + h(u - \sin \theta) = 0 

при различных значениях параметра h.

4. Какую краевую задачу нужно решить, чтобы найти стационарные решения следующей задачи:

~~
\begin{array}{ll} (\text{PHP}) & u_{tt} = u_{xx} - u_t + u, \quad 0 < x < 1, \quad 0 < t < \infty, \\ (\text{GU}) & \left\{ \begin{array}{ll} u\left(0, \t\right) = 0, \\ u\left(1, \t\right) = 0, \end{array} \right. \quad 0 < t < \infty, \\ (\text{NU}) & \left\{ \begin{array}{ll} u\left(x, \ 0\right) = \sin\left(3\pi x\right), \\ u_t\ left(x, \ 0\right) = 0? \end{array} \right. \quad 0 \leqslant x \leqslant 1, \end{array}
~~

 Исходя из физической интерпретации лапласиана, что вы можете сказать об общих свойствах решений краевой задачи для уравнения Гельмгольца

(УЧП)
\end{verbatim}

\texttt{\textbackslash{}Delta\ u\ =\ -\textbackslash{}lambda\^{}2\ u}

\texttt{,\ \ \textquotesingle{}0\ \textless{}\ r\ \textless{}\ 1\textquotesingle{}\ ,\ \ \textquotesingle{}(\textbackslash{}Gamma\ Y)\textquotesingle{}\ \ \ \textquotesingle{}u(1,\ \textbackslash{}theta)\ =\ 0\textquotesingle{}\ ,\ \ \textquotesingle{}0\ \textbackslash{}le\ \textbackslash{}theta\ \textless{}\ 2\textbackslash{}pi\textquotesingle{}\ ?}

\begin{enumerate}
\def\labelenumi{\arabic{enumi}.}
\setcounter{enumi}{5}
\tightlist
\item
  Дайте физическую интерпретацию краевой задачи с граничными
\end{enumerate}

условиями смешанного типа внутри квадрата

(YUII)

\begin{verbatim}
u_{xx} + u_{yy} = 0, \quad 0 < x < 1, \quad 0 < y < 1,
\begin{cases} u_y(x, 0) - h[u(x, 0) - 2] = 0, & 0 < x < 1, \\ u(x, 1) = 1, & 0 < x < 1, \\ u_x(0, y) = 0, & 0 < y < 1. \end{cases}
\end{verbatim}

(FY)

\begin{verbatim}
\begin{cases} u_y(x, 0) - h[u(x, 0) - 2] = 0, & 0 < x < 1, \\ u_x(0, y) = 0, & 0 < y < 1. \end{cases}
\end{verbatim}

\section{Лекція 33. ВНУТРІШНЯ ЗАДАЧА ДІРІХЛЕ ДЛЯ
КРУГА}\label{ux43bux435ux43aux446ux456ux44f-33.-ux432ux43dux443ux442ux440ux456ux448ux43dux44f-ux437ux430ux434ux430ux447ux430-ux434ux456ux440ux456ux445ux43bux435-ux434ux43bux44f-ux43aux440ux443ux433ux430}

ЦЕЛЬ ЛЕКЦИИ: Показать, как методом разделения переменных решигь
внутреннюю задачу Дирихле для круга

(УЧП)

\begin{verbatim}
u_{rr} + \frac{1}{r} u_r + \frac{1}{r^2} u_{\theta\theta} = 0, \quad 0 < r < 1,
\end{verbatim}

\texttt{(ГУ)\ \ \textquotesingle{}u(1,\ \textbackslash{}theta)\ =\ g(\textbackslash{}theta),\ \textbackslash{}quad\ 0\ \textbackslash{}le\ \textbackslash{}theta\ \textless{}\ 2\textbackslash{}pi.\textquotesingle{}}

Сначала граничное условие разлагается в ряд

\begin{verbatim}
g(\theta) = \sum_{n=0}^{\infty} [a_n \cos(n\theta) + b_n \sin(n\theta)],
\end{verbatim}

а затем ищется решение каждой из задач

\begin{verbatim}
\Delta u = 0,
\end{verbatim}

\texttt{\textbackslash{}Delta\ u\ =\ 0,}
\texttt{u(1,\ \textbackslash{}theta)\ =\ \textbackslash{}sin(n\textbackslash{}theta),}
\texttt{u(1,\ \textbackslash{}theta)\ =\ \textbackslash{}cos\ n\textbackslash{}theta.}

Решения этих задач записываются в виде

\begin{verbatim}
u(r, \theta) = r^n \sin(n\theta), \quad u(r, \theta) = r^n \cos(n\theta),
\end{verbatim}

откуда получается решение задачи Дирихле

\begin{verbatim}
u(r, \theta) = \sum_{n=0}^{\infty} r^n [a_n \cos(n\theta) + b_n \sin(n\theta)].
\end{verbatim}

После того как решение в виде ряда получено, путем алгебраических
преобразований решению можно придать другую, интегральную форму, которая
называется интегральной формулой Пуассона. В этой лекции предлагается
множество новых идей, как решить задачу Дирихле

(33.1)

\begin{verbatim}
\begin{aligned} (\text{UCH}\Pi) & u_{rr} + \frac{1}{r} u_r + \frac{1}{r^2} u_{\theta\theta} = 0, \quad 0 < r < 1, \\ (\text{GU}) & u(1, \theta) = g(\theta), \quad 0 \leqslant \theta < 2\pi. \end{aligned}
\end{verbatim}

Сначала мы будем решать эту задачу методом разделения переменных, а
затем, когда решение в виде ряда будет найдено, мы преобразуем его к
другому виду (интегральная формула Пуассона).

\pandocbounded{\includegraphics[keepaspectratio]{images/_page_247_Figure_5.jpeg}}

РИС. 33.1. Схема решения внутренней задачи Дирихле.

Задача (33.1) очень важна в физических приложениях. Ее можно
интерпретировать как задачу нахождения электростатического потенциала
внутри круга по известному распределению потенциала на границе. Другая
интерпретация --- модель мыльной пленки. Если сделать из проволоки
кольцо и изотнуть его так, чтобы изгиб задавался функцией
\texttt{g(\textbackslash{}theta)} , а затем погрузить в мыльный раствор,
то мыльная пленка натянется в соответствии с формой проволоки.
Возвышения точек пленки описываются решением задачи (33.1), если
смещения \texttt{g(\textbackslash{}theta)} малы.

Читатель должен хорошо разобраться в методе разделения переменных, в
общих чертах изображенном на рис. 33.1, и тщательно проработать его.

Дадим несколько пояснений к рис. 33.1.

\begin{itemize}
\tightlist
\item
  \begin{enumerate}
  \def\labelenumi{\arabic{enumi}.}
  \tightlist
  \item
    Читатель должен проверить, что константа разделения должна быть
    неотрицательной (вот почему она обозначена
    \texttt{\textbackslash{}lambda\^{}2} ). Если бы константа разделения
    была отрицательной, то функции
    \texttt{\textbackslash{}Theta(\textbackslash{}theta)} не были бы
    периодическими. Если же константа разделения нулевая, то в решение
    необходимо ввести логарифмический член, т. е.
    \texttt{R(r)\ =\ a\ +\ \textbackslash{}ln\ r} .
  \end{enumerate}
\item
  \begin{enumerate}
  \def\labelenumi{\arabic{enumi}.}
  \setcounter{enumi}{1}
  \tightlist
  \item
    Читатель должен знать, как получить константы \texttt{a\_n} и
    \texttt{b\_n} . Итак, решение внутренней задачи Дирихле записывается
    в виде
  \end{enumerate}
\end{itemize}

(33.2)

\begin{verbatim}
u(r, \theta) = \sum_{n=0}^{\infty} r^n \left[ a_n \cos(n\theta) + b_n \sin(n\theta) \right].
\end{verbatim}

Прежде чем перейти к другой форме представления решения, проведем анализ
формулы (33.2).

\section{Анализ решения задачи
Дирихле}\label{ux430ux43dux430ux43bux438ux437-ux440ux435ux448ux435ux43dux438ux44f-ux437ux430ux434ux430ux447ux438-ux434ux438ux440ux438ux445ux43bux435}

\begin{enumerate}
\def\labelenumi{\arabic{enumi}.}
\tightlist
\item
  Остановимся сначала на структуре решения (33.2). Согласно этому
  решению, мы в первую очередь должны разложить граничную функцию в ряд
  Фурье
\end{enumerate}

\begin{verbatim}
g(\theta) = \sum_{n=0}^{\infty} [a_n \cos(n\theta) + b_n \sin(n\theta)],
\end{verbatim}

\texttt{а\ затем\ решить\ задачу\ для\ каждого\ синуса\ и\ косинуса,\ входящих\ в\ ряд.\ Каждое\ из\ этих\ слагаемых\ порождает\ решения\ вида\ \ \textquotesingle{}r\^{}n\ \textbackslash{}sin\{(n\textbackslash{}theta)\}\textquotesingle{}\ \ и\ \ \textquotesingle{}r\^{}n\ \textbackslash{}cos\{(n\textbackslash{}theta)\}\textquotesingle{}\ ,\ так\ что\ по\ принцицу\ суперпозиции}

\begin{verbatim}
u(r, \theta) = \sum_{n=0}^{\infty} r^n [a_n \cos(n\theta) + b_n \sin(n\theta)].
\end{verbatim}

\begin{enumerate}
\def\labelenumi{\arabic{enumi}.}
\setcounter{enumi}{1}
\tightlist
\item
  Решение, например, задачи
\end{enumerate}

(УЧП)

\begin{verbatim}
\Delta u = 0
\end{verbatim}

\begin{verbatim}
,  0 < r < 1 ,  
(ГУ)  u(1, \theta) = 1 + \sin \theta + \frac{1}{2} \sin 3\theta + \cos (4\theta) 

должно иметь вид
\end{verbatim}

\texttt{u(r,\ \textbackslash{}theta)\ =\ 1\ +\ r\ \textbackslash{}sin\ \textbackslash{}theta\ +\ \textbackslash{}frac\{r\^{}3\}\{2\}\ \textbackslash{}sin\ (3\textbackslash{}theta)\ +\ r\^{}4\ \textbackslash{}cos\ (4\textbackslash{}theta)}

\begin{verbatim}
поскольку  g(\theta)  в данном случае уже разложена в ряд Фурье с коэффициентами  a_0=1,\ a_4=1,\ b_1=1,\ b_3=0,5,  все остальные значения  a_n  и  b_n  равны нулю, поэтому для их вычисления нет необходимости пользоваться формулами.

3. Если радиус круга произволен (скажем, R), то решение запишется в виде
\end{verbatim}

\texttt{u(r,\ \textbackslash{}theta)\ =\ \textbackslash{}sum\_\{n=0\}\^{}\{\textbackslash{}infty\}\ (r/R)\^{}n\ {[}a\_n\ \textbackslash{}cos(n\textbackslash{}theta)\ +\ b\_n\ \textbackslash{}sin(n\textbackslash{}theta){]}.}

\begin{verbatim}
4. Заметим, что константа  a_{\theta}  в решении (33.2) представляет собой среднее значение функции  g\left(\theta\right)
\end{verbatim}

\texttt{a\_0\ =\ \textbackslash{}frac\{1\}\{2\textbackslash{}pi\}\ \textbackslash{}int\_0\^{}\{2\textbackslash{}pi\}\ g(\textbackslash{}theta)\ d\textbackslash{}theta.}

\begin{verbatim}
Этим завершается решение задачи методом разделения переменных. Перейдем теперь к весьма интересной интегральной формуле Пуассона.

# Интегральная формула Пуассона

Возьмем решение задачи Дирихле, полученное методом разделения переменных
\end{verbatim}

\texttt{u(r,\ \textbackslash{}theta)\ =\ \textbackslash{}sum\_\{n=0\}\^{}\{\textbackslash{}infty\}\ (r/R)\^{}n\ \textbackslash{}left{[}\ a\_n\ \textbackslash{}cos(n\textbackslash{}theta)\ +\ b\_n\ \textbackslash{}sin(n\textbackslash{}theta)\ \textbackslash{}right{]},}

\begin{verbatim}
(мы взяли случай произвольного радиуса круга) и подставим в него формулы для вычисления коэффициентов  a_n  и  b_n . После ряда выкладок с применением алгебры, интегрального исчисления и тригонометрии получаем
\end{verbatim}

\texttt{u(r,\ \textbackslash{}theta)\ =\ \textbackslash{}frac\{1\}\{2\textbackslash{}pi\}\ \textbackslash{}int\_\{0\}\^{}\{2\textbackslash{}pi\}\ g(\textbackslash{}theta)\ d\textbackslash{}theta\ +\ \textbackslash{}frac\{1\}\{\textbackslash{}pi\}\ \textbackslash{}sum\_\{n=1\}\^{}\{\textbackslash{}infty\}\ (r/R)\^{}\{n\}\ \textbackslash{}int\_\{0\}\^{}\{2\textbackslash{}pi\}\ g(\textbackslash{}alpha)\ \textbackslash{}times\ \textbackslash{}\textbackslash{}\ \textbackslash{}times\ {[}\textbackslash{}cos(n\textbackslash{}alpha)\textbackslash{}cos(n\textbackslash{}theta)\ +\ \textbackslash{}sin.\ (n\textbackslash{}alpha)\textbackslash{}sin(n\textbackslash{}theta){]}\ d\textbackslash{}alpha\ =\ \textbackslash{}\textbackslash{}\ =\ \textbackslash{}frac\{1\}\{2\textbackslash{}pi\}\ \textbackslash{}int\_\{0\}\^{}\{2\textbackslash{}pi\}\ \textbackslash{}left\textbackslash{}\{\ 1\ +\ 2\ \textbackslash{}sum\_\{n=1\}\^{}\{\textbackslash{}infty\}\ (r/R)\^{}\{n\}\ \textbackslash{}cos{[}n(\textbackslash{}theta\ -\ \textbackslash{}alpha){]}\ \textbackslash{}right\textbackslash{}\}\ g(\textbackslash{}alpha)\ d\textbackslash{}alpha\ =\ \textbackslash{}\textbackslash{}\ =\ \textbackslash{}frac\{\ 1\}\{2\textbackslash{}pi\}\ \textbackslash{}int\_\{0\}\^{}\{2\textbackslash{}pi\}\ \textbackslash{}left\textbackslash{}\{\ 1\ +\ \textbackslash{}sum\_\{n=1\}\^{}\{\textbackslash{}infty\}\ (r/R)\^{}\{n\}\ \textbackslash{}left{[}\ e\^{}\{in(\textbackslash{}theta\ -\ \textbackslash{}alpha)\}\ +\ e\^{}\{-in(\textbackslash{}theta\ -\ \textbackslash{}alpha)\}\ \textbackslash{}right{]}\ \textbackslash{}right\textbackslash{}\}\ g(\textbackslash{}alpha)\ d\textbackslash{}alpha\ =\ \textbackslash{}\textbackslash{}\ =\ \textbackslash{}frac\{1\}\{2\textbackslash{}pi\}\ \textbackslash{}int\_\{\ 0\}\^{}\{2\textbackslash{}pi\}\ \textbackslash{}\ left\textbackslash{}\{\ 1\ +\ \textbackslash{}frac\{re\^{}\{i(\textbackslash{}theta\ -\ \textbackslash{}alpha)\}\}\{\textbackslash{}left{[}R\ -\ re\^{}\{i(\textbackslash{}theta\ -\ \textbackslash{}alpha)\}\textbackslash{}right{]}\}\ +\ \textbackslash{}frac\{re\^{}\{-i(\textbackslash{}theta\ -\ \textbackslash{}alpha)\}\}\{\textbackslash{}left{[}R\ -\ re\^{}\{-i(\textbackslash{}theta\ -\ \textbackslash{}alpha)\}\textbackslash{}right{]}\}\ \textbackslash{}right\textbackslash{}\}\ g(\textbackslash{}alpha)\ d\textbackslash{}alpha\ =\ \textbackslash{}\textbackslash{}\ =\ \textbackslash{}frac\{1\}\{2\textbackslash{}pi\}\ \textbackslash{}int\_\{0\}\^{}\{2\textbackslash{}pi\}\ \textbackslash{}left{[}\ \textbackslash{}frac\{R\^{}\{2\}\ -\ r\^{}\{2\}\}\{R\^{}\{2\}\ -\ 2rR\textbackslash{}cos(\textbackslash{}theta\ -\ \textbackslash{}alpha)\ +\ r\^{}\{2\}\}\ \textbackslash{}right{]}\ g(\textbackslash{}alpha)\ d\textbackslash{}alpha.}

\begin{verbatim}
Последнее выражение и называется интегральной формулой Пуассона. Итак, мы получили еще одну форму представления реше-

ния внутренией задачи Дирихле

(33.3)
\end{verbatim}

\texttt{U(R,\ \textbackslash{}theta)\ =\ \textbackslash{}frac\{1\}\{2\textbackslash{}pi\}\ \textbackslash{}int\_\{0\}\^{}\{2\textbackslash{}pi\}\ \textbackslash{}left{[}\ \textbackslash{}frac\{R\^{}\{2\}\ -\ r\^{}\{2\}\}\{R\^{}\{2\}\ -\ 2rR\textbackslash{}cos((0\ -\ \textbackslash{}alpha)\ +\ r\^{}\{2\})\}\ \textbackslash{}right{]}\ G(\textbackslash{}alpha)\ D\textbackslash{}alpha.}

\begin{verbatim}
# Анализ интегральной формы решения

1. Можно считать, что по формуле Пуассона (33.3) потенциал u в точке  (r,\theta)  является взвещенным средним граничного потенциала. Весовая функция называется ядром Пуассона и имеет вид

Ядро Пуассона =
\end{verbatim}

\texttt{\textbackslash{}frac\{R\^{}2-r\^{}2\}\{R\^{}2-2rR\textbackslash{}cos(\textbackslash{}theta-\textbackslash{}alpha)+r\^{}2\}.}

\begin{verbatim}
Это говорит нам кое-что о физике системы, а именно: потенциал в точке является взвешенным средним потенциалов соседних точек.

![](images/_page_250_Picture_7.jpeg)

РИС 33.2. Решение  a\left(r,\theta\right)  — взвешенная сумма граничных потенциалов. a — Квадрат этон длины является знаменателем ядра Пуассона.

Ядро Пуассона показывает только, каков вклад каждой точки в общий потенциал (рис. 33.2).

Для граничных значений  (R, \alpha) , близких к  (r, \theta) , ядро Пуассона становится большим, так как значенатель ядра равен квадрату расстояния между точками  (r, \theta)  и  (R, \alpha)  и, значит, мал (рис. 33.2) Благодаря этому больший вклад в интеграл дают значения  g(\alpha) , для которых  (R, \alpha)  близки к  (r, \theta) . К сожалению, если точка  (r, \theta)  как угодно близко подходит к границе r = R, то ядро Пуассона также неограниченно возрастает. Поэтому

численные вначения решения вблизи границы удобнее вычислять с помощью рядов.

2. Если вычислить потенциал в центре круга по формуле Пуассона, то получим
\end{verbatim}

\texttt{u(0,\ 0)\ =\ \textbackslash{}frac\{1\}\{2\textbackslash{}pi\}\ \textbackslash{}int\_\{0\}\^{}\{2\textbackslash{}pi\}\ g(\textbackslash{}alpha)\ d\textbackslash{}alpha.}

\begin{verbatim}
Другими словами, потенциал в центре равен среднему значению потенциала на окружности.

Этим мы завершаем обсуждение задачи Дирихле. Читателю может показаться, что на самом деле эта задача и не так уж важна, поскольку область, в которой она решается, так проста. Действительно, уравнение Лапласа упрощается, если область, в которой ищется решение, представляет собой круг, квадрат, полуплоскость и т. д. Однако здесь существенны два момента.

Во многих случаях экспериментатор, разрабатывая свою физическую аппаратуру, может выбирать форму границы по своему

усмотрению (в том числе и в виде окружности).

2. Далее мы будем изучать конформные отображения, которые позволяют преобразовывать сложные области в области простой формы (например, круга). Следовательно, чтобы решить задачу Дирихле в произвольной области, можно конформно отобразить ее на круг, воспользоваться решением, полученным в этой лекции, и преобразовать затем его к исходным координатам.

#### ЗАМЕЧАНИЯ

1. Краевую задачу для неоднородного уравнения

(PPP)

  \Delta u = f , внутри  D ,  (ГУ)   u = 0 , на границе  D ,

всегда можно рещить следующим приемом.

Найдем частное решение V уравнения  \Delta V = f  и решим краевую задачу
\end{verbatim}

\texttt{\textbackslash{}Delta\ W\ =\ 0}

, внутри \texttt{D} , \texttt{W\ =\ V} , на границе \texttt{D} .

Тогда решение исходной задачи находится по формуле u=V-W. Другими
словами, неоднородность уравнения можно перенести в граничное условие.

\begin{enumerate}
\def\labelenumi{\arabic{enumi}.}
\setcounter{enumi}{1}
\tightlist
\item
  Краевую вадачу с неоднородными граничными условиями
\end{enumerate}

(PPP)

\texttt{\textquotesingle{}\textbackslash{}Delta\ u\ =\ 0\textquotesingle{}\ ,\ внутри\ \ \textquotesingle{}D\textquotesingle{}\ ,\ \ \textquotesingle{}(ГУ)\textquotesingle{}\ \ \ \textquotesingle{}u\ =\ f\textquotesingle{}\ ,\ на\ границе\ \ \textquotesingle{}D\textquotesingle{}\ ,}

можно решить следующим способом:

\begin{itemize}
\tightlist
\item
  а) найти любую функцию V, которая удовлетворяет на границе D условию
  V=f;
\item
  б) решить новую краевую задачу
\end{itemize}

\texttt{\textbackslash{}Дельта\ W\ =\ \textbackslash{}Дельта\ V}

\begin{verbatim}
, внутри  D ,  W = 0 , на границе  D ;

в) найти решение исходной задачи по формуле u = V - W. Другими словами, неоднородность из граничного условия можно перенести в уравнение.

#### ЗАДАЧИ

1. Тщательно разберитесь во всех деталях применения метода разделения переменных при решении задачи Дирихле (33.1). Это важно потому, что, когда вы будете решать следующие задачи, возникнут интересные отличия. Вы должны особенно хорошо понимать, почему константа разделения не может быть отрицательной, и при  \lambda=0  отбрасывается очень важное решение. Почему? Посмотрите внимательно на рис. 33.1.

2. Найдите решение внутренней задачи Дирихле

(УЧП)
\end{verbatim}

\texttt{u\_\{rr\}\ +\ \textbackslash{}frac\{1\}\{r\}\ u\_r\ +\ \textbackslash{}frac\{1\}\{r\^{}3\}\ u\_\{\textbackslash{}theta\textbackslash{}theta\}\ =\ 0}

, \texttt{0\ \textless{}\ r\ \textless{}\ 1} ,

при следующих граничных условиях:

\begin{enumerate}
\def\labelenumi{\alph{enumi})}
\tightlist
\item
\end{enumerate}

\begin{verbatim}
u(1, \theta) = 1 + \sin \theta + \frac{1}{2} \cos \theta
\end{verbatim}

\begin{verbatim}
,

- 6)  u(1, \theta) = 2 ,
- B)  u(1, \theta) = \sin \theta ,
- r)  u(1, \theta) = \sin 3\theta .

Как выглядят ваши решения? Удовлетворяют ли они уравнению Лапласа?

3. Найдите решение внутренней задачи Дирихле в круге радиусом R=2.

(УЧП)
\end{verbatim}

\texttt{\textbackslash{}Delta\ u\ =\ 0}

, \texttt{0\ \textless{}\ r\ \textless{}\ 2} ,\\
\texttt{(ГУ)\ \ \textquotesingle{}u(2,\ \textbackslash{}theta)\ =\ \textbackslash{}sin\ \textbackslash{}theta\textquotesingle{}\ ,\ \ \textquotesingle{}0\ \textbackslash{}le\ \textbackslash{}theta\ \textbackslash{}le\ 2\textbackslash{}pi\textquotesingle{}\ .}

Как выглядит график этого решения?

\texttt{-\ 4.\ Как\ изменится\ решение\ задачи\ 3,\ если\ граничное\ условие\ заменить\ на\ \ \textquotesingle{}u(2,\ \textbackslash{}theta)\ =\ \textbackslash{}sin(2\textbackslash{}theta)\textquotesingle{}\ ?\ Как\ выглядит\ график\ этого\ решения?}
- 5. Решите задачу

\begin{verbatim}
(\mathsf{Y}\mathsf{Y}\mathsf{\Pi}) \quad \Delta u = 0, \quad 0 < r < 1,
(\mathsf{\Gamma}\mathsf{Y}) \quad u(1, \, \theta) = \begin{cases} \sin \theta, & 0 \le \theta < \pi, \\ 0, & \pi \le \theta < 2\pi. \end{cases}
\end{verbatim}

Прикиньте, как выглядит график решения?

\texttt{-\ 6.\ Пусть\ в\ ядре\ Пуассона\ r=3R/4,\ \ \textquotesingle{}\textbackslash{}theta=\textbackslash{}pi/2\textquotesingle{}\ .\ Постройте\ график\ зависимости\ ядра\ Пуассона\ от\ угла\ \ \textquotesingle{}\textbackslash{}alpha\textquotesingle{}\ \ в\ диапазоне\ \ \textquotesingle{}0\ \textbackslash{}le\ \textbackslash{}alpha\ \textless{}\ 2\textbackslash{}pi\textquotesingle{}\ .}
- 7. Убедитесь в справедливости замечания 1 из этой лекции. - 8.
Убедитесь в справедливости замечания 2 из этой лекции.

\section{Лекція 34. ЗАДАЧА ДІРІХЛЕ В
КІЛЬЦІ}\label{ux43bux435ux43aux446ux456ux44f-34.-ux437ux430ux434ux430ux447ux430-ux434ux456ux440ux456ux445ux43bux435-ux432-ux43aux456ux43bux44cux446ux456}

ЦЕЛЬ ЛЕКЦИИ: Решить задачу Дирихле в кольце

(УЧП)

\begin{verbatim}
u_{rr} + \frac{1}{r} u_r + \frac{1}{r^2} u_{\theta\theta} = 0, \quad R_1 < r < R_2,
\end{verbatim}

(ГУ)

\begin{verbatim}
\begin{cases} u(R_1, \theta) = g_1(\theta), \\ u(R_2, \theta) = g_2(\theta). \end{cases}
\end{verbatim}

\texttt{0\ \textbackslash{}le\ \textbackslash{}theta\ \textbackslash{}le\ 2\textbackslash{}pi,}

Здесь мы опять воспользуемся методом разделения переменных, но в отличие
от внутренней зодачи Дирихле мы не станем отбрасывать решения вида

\begin{verbatim}
\frac{1}{r^n}\sin{(n\theta)}, \quad \frac{1}{r^n}\cos{(n\theta)}, \quad \ln r.
\end{verbatim}

Следовательно, решение должно записываться в виде
\texttt{u(r,\ \textbackslash{}theta)\ =\ a\_0\ +\ b\_0\ \textbackslash{}ln\ r\ +}

\begin{verbatim}
+\sum_{n=1}^{\infty} \left[ (a_n r^n + b_n r^{-n}) \cos(n\theta) + (c_n r^n + d_n r^{-n}) \sin(n\theta) \right].
\end{verbatim}

Мы также кратко обсудим решение внешней задачи Дирихле. В этом случае
нам придется отбросить слагаемые, неограниченные при
\texttt{r=\textbackslash{}infty} . Следовательно, решение внешней задачи
Дирихле

\begin{verbatim}
\begin{array}{ll} (\text{PHP}) & \Delta u = 0, \quad 1 < r < \infty, \\ (\text{GU}) & u\left(1, \; \theta\right) = g_1\left(\theta\right), \quad 0 \leqslant \theta < 2\pi, \end{array}
\end{verbatim}

следует искать в виде

\begin{verbatim}
u(r, \theta) = \sum_{n=0}^{\infty} r^{-n} [a_n \cos(n\theta) + b_n \sin(n\theta)].
\end{verbatim}

Существует множество разнообразных по форме областей, для которых можно
в явном виде выписать решение задачи Дирихле. Мы перечислим только
некоторые из них:

\begin{itemize}
\tightlist
\item
  а) в круге (лекция 33),
\item
  б) в кольце (данная лекция),
\item
  в) вне круга (данная лекция),
\item
  г) в шаре (в последующих лекциях),
\item
  д) в сферическом слое (в последующих лекциях),
\item
  е) в полосе (двумерный случай),
\item
  ж) в слое между плоскостями (трехмерный случай).
\end{itemize}

Эгот список можно продолжать до бесконечности. Мы хотим решить простую,
но достаточно типичную задачу Дприхле, так чтобы читатель освоил общие
принципы и мог при необходимости решить свою задачу.

\pandocbounded{\includegraphics[keepaspectratio]{images/_page_254_Picture_3.jpeg}}

РИС. 34.1. Уравнение Лапласа в кольце.

Здесь наша цель---найти форму мыльной пленки, натинутой между двумя
проволочными ободками. Вероятно, в этой задаче интунция уже не играет
такой роли, как во внутренней задаче Дирихле. Математическую модель
задачи запишем следующим образом:

(34.1)

\begin{verbatim}
\begin{aligned} & (\text{UBP}) \quad u_{rr} + \frac{1}{r} u_r + \frac{1}{r^2} u_{\theta\theta} = 0, \quad R_1 < r < R_2, \\ & (\text{GU}) \quad \begin{cases} u(R_1, \ \theta) = g_1(\theta), \\ u(R_2, \ \theta) = g_2(\theta). \end{cases} \quad 0 \leqslant \theta \leqslant 2\pi, \end{aligned}
\end{verbatim}

Общая схема этой задачи изображена на рис. 34.1.

Задача состоит в том, чтобы найти решение
\texttt{u(r,\ \textbackslash{}theta)} между окружностями
\texttt{r\ =\ R\_1} и \texttt{r\ =\ R\_2} , совпадающее на окружностях с
\texttt{g\_1(\textbackslash{}theta)} и соответственно с
\texttt{g\_2(\textbackslash{}theta)} . Будем искать решения вида

\begin{verbatim}
u(r, \theta) = R(r) \Theta(\theta)
\end{verbatim}

\begin{verbatim}
.

Подставляя их в уравнение Лапласа, приходим к двум обыкновенным дифференциальным уравнениям для функций R(r) и  \Theta(\theta)
\end{verbatim}

\texttt{r\^{}2R\textquotesingle{}\textquotesingle{}\ +\ rR\textquotesingle{}\ -\ \textbackslash{}lambda\^{}2R\ =\ 0}

\begin{verbatim}
(уравнение Эйлера),  \Theta'' + \lambda^2\Theta = 0 .

Отметим, что константа разделения в этих уравнениях должна быть больше или равна нулю (мы обозначили ее  \lambda^2 ), иначе реше-

ние для  \Theta  ( \theta ) не будет периодическим. Далее мы решим оба этих уравнения, а затем перемножим эти решения, в результате получим все решения уравнения Лапласа, которые представимы в виде R (r)  \Theta  ( \theta ).

Из этих двух уравнений наибольший интерес представляет уравнение Эйлера. К счастью, это одно из немногих обыкновенных дифференциальных уравнений с переменными коэффициентами, которые решаются сравнительно легко.

#### РЕШЕНИЕ УРАВНЕНИЯ ЭЙЛЕРА
\end{verbatim}

\texttt{(r\^{}2R\textquotesingle{}\textquotesingle{}\ +\ rR\textquotesingle{}\ -\ \textbackslash{}lambda\^{}2R\ =\ 0)}

\begin{verbatim}
СЛУЧАЙ I (\lambda = 0). В этом случае уравнение Эйлера принимает вид
\end{verbatim}

r\^{}2R'\,' + rR' = 0,

\begin{verbatim}
общее решение которого легко находится:
\end{verbatim}

\texttt{R(r)\ =\ a\ +\ b\ \textbackslash{}ln\ r.}

\begin{verbatim}
В самом деле, сделав подстановку  V\left(r\right) \Rightarrow R'\left(r\right) , получаем для V уравнение с разделяющимися переменными
\end{verbatim}

rV' + V = 0

.

\texttt{Решив\ это\ уравнение,\ находим\ \ \textquotesingle{}V\textbackslash{}left(r\textbackslash{}right)=c\_\{1\}/r\textquotesingle{}\ ,\ откуда}

\begin{verbatim}
R(r) = c_1 \ln r + c_2
\end{verbatim}

\begin{verbatim}
.

СЛУЧАЙ 2. ( \lambda > 0 ). В этом случае уравнение Эйлера имеет вид
\end{verbatim}

\texttt{r\^{}2R\textquotesingle{}\textquotesingle{}\ +\ rR\textquotesingle{}\ -\ \textbackslash{}lambda\^{}2R\ =\ 0.}

\begin{verbatim}
Будем искать решение этого уравнения в виде  R(r)=r^{\alpha} . Наша цель—найти два значения  \alpha  (скажем,  \alpha_1  и  \alpha_2 ), таких, чтобы функция
\end{verbatim}

\texttt{R(r)\ =\ ar\^{}\{\textbackslash{}alpha\_1\}\ +\ br\^{}\{\textbackslash{}alpha\}}

\begin{verbatim}
являлась общим решением уравнения. Подстановка  R(r) = r^{\alpha}  в уравнение Эйлера приводит к  \alpha = \lambda  и  \alpha = -\lambda . Следовательно,
\end{verbatim}

\texttt{R(r)\ =\ ar\textquotesingle{}\ +\ br\^{}\{-\textbackslash{}lambda\}.}

\begin{verbatim}
Итак, используя решение уравнения Эйлера, мы получили
\end{verbatim}

\texttt{\textbackslash{}lambda\ =\ 0\ \textbackslash{}begin\{cases\}\ R(r)\ =\ a\ +\ b\ \textbackslash{}ln\ r,\ \textbackslash{}\textbackslash{}\ \textbackslash{}Theta(\textbackslash{}theta)\ =\ c\ +\ d\textbackslash{}theta,\ \textbackslash{}\textbackslash{}\ \textbackslash{}lambda\ \textgreater{}\ 0\ \textbackslash{}end\{cases\}}
\texttt{\textbackslash{}begin\{cases\}\ R(r)\ =\ ar\^{}\{\textbackslash{}lambda\}\ +\ br\^{}\{-\textbackslash{}lambda\},\ \textbackslash{}\textbackslash{}\ \textbackslash{}Theta(\textbackslash{}theta)\ =\ c\ \textbackslash{}cos(\textbackslash{}lambda\ \textbackslash{}theta)\ +\ d\ \textbackslash{}sin(\textbackslash{}lambda\ \textbackslash{}theta).\ \textbackslash{}end\{cases\}}

\begin{verbatim}
Функция  \Theta(\theta)  должна быть *периодической* с периодом  2\pi . Значит, величина  \lambda  должна принимать значения  0, 1, 2, 3, \ldots . Теперь мы можем с помощью произведения составить следующие решения.

# Решения уравнения Лапласа
\end{verbatim}

\texttt{c\ -\ \textbackslash{}text\{constant\},}

\begin{verbatim}
c \ln r,  cr^n \cos(n\theta) \ (n = 1, 2, ...),    cr^n \sin(n\theta) \ (n = 1, 2, ...),  cr^{-n} \cos(n\theta) \ (n = 1, 2, ...),    cr^{-n} \sin(n\theta) \ (n = 1, 2, ...). Поскольку сумма этих решений также является решением, то делаем вывод, что общее решение уравнения Лапласа в нашем случае должно иметь вид

(34.2)
\end{verbatim}

\texttt{u(r,\ \textbackslash{}theta)\ =\ a\_0\ +\ b\_0\ \textbackslash{}ln\ r\ +\ \textbackslash{}sum\_\{n=1\}\^{}\{\textbackslash{}infty\}\ \textbackslash{}left{[}\ (a\_n\ r\^{}n\ +\ b\_n\ r\^{}\{-n\})\ \textbackslash{}cos(n\textbackslash{}theta)\ +\ (c\_n\ r\^{}n\ +\ d\_n\ r\^{}\{-n\})\ \textbackslash{}sin(n\textbackslash{}theta)\ \textbackslash{}right{]}.}

\begin{verbatim}
Осталось только найти все коэффициенты в сумме (34.2) так, чтобы удовлетворить граничным условиям
\end{verbatim}

\texttt{\textbackslash{}begin\{cases\}\ u(R\_1,\ \textbackslash{}theta)\ =\ g\_1(\textbackslash{}theta),\ \textbackslash{}\textbackslash{}\ u(R\_2,\ \textbackslash{}theta)\ =\ g\_2(\textbackslash{}theta).\ \textbackslash{}end\{cases\}}

\begin{verbatim}
Подстановка (34.2) в граничные условия приводит к следующим уравнениям:
\end{verbatim}

\texttt{\textbackslash{}begin\{cases\}\ a\_0\ +\ b\_0\ \textbackslash{}ln\ R\_1\ =\ \textbackslash{}frac\{1\}\{2\textbackslash{}pi\}\ \textbackslash{}int\textbackslash{}limits\_0\^{}\{2\textbackslash{}pi\}\ g\_1(s)\ \textbackslash{},\ ds\ \textbackslash{}\textbackslash{}\ a\_0\ +\ b\_0\ \textbackslash{}ln\ R\_2\ =\ \textbackslash{}frac\{1\}\{2\textbackslash{}pi\}\ \textbackslash{}int\textbackslash{}limits\_0\^{}\{2\textbackslash{}pi\}\ g\_2(s)\ \textbackslash{},\ ds\ \textbackslash{}end\{cases\}}

\begin{verbatim}
(решаются относительно  a_0  и  b_0 ),
\end{verbatim}

\texttt{\textbackslash{}begin\{cases\}\ a\_n\ R\_1\^{}n\ +\ b\_n\ R\_1\^{}\{-n\}\ =\ \textbackslash{}frac\{1\}\{\textbackslash{}pi\}\ \textbackslash{}int\textbackslash{}limits\_0\^{}\{2\textbackslash{}pi\}\ g\_1(s)\ \textbackslash{}cos(ns)\ \textbackslash{},\ ds\ \textbackslash{}\textbackslash{}\ a\_n\ R\_2\^{}n\ +\ b\_n\ R\_2\^{}\{-n\}\ =\ \textbackslash{}frac\{1\}\{\textbackslash{}pi\}\ \textbackslash{}int\textbackslash{}limits\_0\^{}\{2\textbackslash{}pi\}\ g\_2(s)\ \textbackslash{}cos(ns)\ \textbackslash{},\ ds\ \textbackslash{}end\{cases\}}

\begin{verbatim}
(решаются относительно  a_n  и  b_n ),
\end{verbatim}

\texttt{\textbackslash{}begin\{cases\}\ c\_n\ R\_1\^{}n\ +\ d\_n\ R\_1\^{}\{-n\}\ =\ \textbackslash{}frac\{1\}\{\textbackslash{}pi\}\ \textbackslash{}int\textbackslash{}limits\_0\^{}\{2\textbackslash{}pi\}\ g\_1(s)\ \textbackslash{}sin\{(ns)\}\ \textbackslash{},\ ds\ \&\ \textbackslash{}text\{(потрібно\ розв\textquotesingle{}язати\ відносно\ \}\ c\_n\ \textbackslash{}\textbackslash{}\ c\_n\ R\_2\^{}n\ +\ d\_n\ R\_2\^{}\{-n\}\ =\ \textbackslash{}frac\{1\}\{\textbackslash{}pi\}\ \textbackslash{}int\textbackslash{}limits\_0\^{}\{2\textbackslash{}pi\}\ g\_2(s)\ \textbackslash{}sin\{(ns)\}\ \textbackslash{},\ ds\ \&\ \textbackslash{}text\{and\ \}\ d\_n\ \textbackslash{}text\{)\}\ (s)\ \textbackslash{}sin\{(ns)\}\ \textbackslash{},\ ds\ \&\ \textbackslash{}text\{and\ \}\ \ \textbackslash{}text\{)\}.\ \textbackslash{}end\{cases\}}

\begin{verbatim}
Тем самым из этих уравнений находятся все неизвестные коэффициенты  a_0 ,  b_0 ,  a_n ,  b_n ,  c_n  и  d_n . Теперь задача (34.1) полностью решена. Решение дается выражением (34.2), коэффициенты которого определяются по формулам (34.3).

Чтобы лучше понять полученные результаты, рассмотрим

несколько примеров.

# Некоторые варианты задачи Дирихле в кольце

ПРИМЕР 1.

Предположим, что потенциал на внутренней окружности равен нулю, а потенциал на внешней равен  \sin\theta . Тогда

(УЧП)
\end{verbatim}

\texttt{\textbackslash{}Delta\ u\ =\ 0}

, \texttt{1\ \textless{}\ r\ \textless{}\ 2} ,\\
(ГУ)

\begin{verbatim}
\begin{cases} u(1, \theta) = 0, \\ u(2, \theta) = \sin \theta. \end{cases}
\end{verbatim}

\texttt{\textquotesingle{}0\ \textbackslash{}le\ \textbackslash{}theta\ \textbackslash{}le\ 2\textbackslash{}pi\textquotesingle{}\ ,}

Первое, что мы должны сделать---это вычислить все интегралы в (34.3) и
решить соответствующие уравнения для \texttt{a\_0} , \texttt{b\_0} ,
\texttt{a\_n} , \texttt{b\_n} , \texttt{c\_n} и \texttt{d\_n} .
Целесообразно сначала попытаться подобрать такие частные решения (из
выписанных выше), линейная комбинация которых удовлетворяет граничным
условням (здесь
\texttt{u\ =\ Ar\ \textbackslash{}sin\ \textbackslash{}theta\ +\ Br\^{}\{-1\}\ \textbackslash{}sin\ \textbackslash{}theta}
). Затем из граничных условий найти коэффициенты (здесь A = 2/3, B =
-2/3). В результате этих простых вычислений получаем

\begin{verbatim}
a_0=0, \quad b_0=0, \quad a_n=0, \quad b_n=0 \quad (n=1, 2, \ldots)
c_n=\begin{cases}
2/3, & n=1, \\
0 & \text{для всіх інших } n,
\end{cases}
\qquad d_n=\begin{cases}
-2/3, & n=1, \\
0 & \text{для всіх інших } n.
\end{cases}
\end{verbatim}

Зная все коэффициенты, выпишем решение задачи

\begin{verbatim}
u(r, \theta) = \frac{2}{3} \left(r - \frac{1}{r}\right) \sin \theta.
\end{verbatim}

Легко проверить, что эта функция удовлетворяет обоим граничным условиям.
Очевидно и то, что она удовлетворяет уравнению Лапласа, так как
принадлежит классу функций вида (34.2).

\paragraph{ПРИМЕР 2.}\label{ux43fux440ux438ux43cux435ux440-2.}

Рассмотрим задачу с \emph{постоянными потенциалами} на гранинах кольца

(УЧП)

\begin{verbatim}
\Delta u = 0
\end{verbatim}

\begin{verbatim}
,  1 < r < 2 .  
(ГУ)
\end{verbatim}

\texttt{\textbackslash{}begin\{cases\}\ u(1,\ \textbackslash{}theta)\ =\ 3,\ \textbackslash{}\textbackslash{}\ u(2,\ \textbackslash{}theta)\ =\ 5.\ \textbackslash{}end\{cases\}\ 0\ \textbackslash{}le\ \textbackslash{}theta\ \textbackslash{}le\ 2\textbackslash{}pi,}

\begin{verbatim}
В этом случае можно сэкономить массу времени, поскольку очевидно, что решение не зависит от  \theta  (так как граничные условия не зависят от  \theta ). Другими словами, наше решение следует искать среди функций вида  a_0 + b_0 \ln r . Подставляя эту функцию в граничные условия, получаем
\end{verbatim}

\texttt{a\_0\ +\ b\_0\ \textbackslash{}ln\ 1\ =\ 3}

,\\
\texttt{\textquotesingle{}a\_0\ +\ b\_0\ \textbackslash{}ln\ 2\ =\ 5\textquotesingle{}\ ,}

откуда \texttt{a\_0\ =\ 3} ,
\texttt{b\_0\ =\ 2/\textbackslash{}ln\ 2\ =\ 2.9} . Следовательно,
решением будет функция

\begin{verbatim}
u(r, \theta) = 3 + 2.9 \ln r
\end{verbatim}

\begin{verbatim}
.

Ее график изображен на рис. 342.

![](images/_page_258_Figure_10.jpeg)

РИС. 34.2. Зависимость потенциала от радиуса внутри кольца.

#### ПРИМЕР 3

Решим еще одну интересную задачу

(УЧП)
\end{verbatim}

\texttt{\textbackslash{}Delta\ u\ =\ 0}

, \texttt{1\ \textless{}\ r\ \textless{}\ 2} ,\\
(ГУ)

\begin{verbatim}
\begin{cases} u(1, \theta) = \sin \theta, \\ u(2, \theta) = \sin \theta. \end{cases} \quad 0 \le \theta \le 2\pi,
\end{verbatim}

Можно быстро проверить, что все коэффициенты \texttt{a\_0} ,
\texttt{b\_0} , \texttt{a\_n} , \texttt{b\_n} , \texttt{c\_n} и
\texttt{d\_n} равны нулю, а коэффициенты \texttt{c\_i} и \texttt{d\_i}
определяются из системы уравнений

\begin{verbatim}
c_1 + d_1 = \frac{1}{\pi} \int_0^{2\pi} \sin^2 \theta \ d\theta = 1,
2c_1 + d_1/2 = \frac{1}{\pi} \int_0^{2\pi} \sin^2 \theta \ d\theta = 1.
\end{verbatim}

Решая эту систему, получаем \texttt{c\_1=1/3} и \texttt{d\_1=2/3} .
Следовательно, решением задачи является функция

\begin{verbatim}
u(r, \theta) = \left(\frac{1}{3}r + \frac{2}{3r}\right)\sin\theta.
\end{verbatim}

Форма кривых при различных значениях \texttt{\textbackslash{}theta}
показана на рис. 34.3.

\pandocbounded{\includegraphics[keepaspectratio]{images/_page_259_Figure_7.jpeg}}

\texttt{РИС.\ 34.3\ Мыльная\ оченка\ между\ \ \textquotesingle{}u\textbackslash{}left(1,\textbackslash{};\textbackslash{}theta\textbackslash{}right)=\textbackslash{}sin\textbackslash{}theta\textquotesingle{}\ \ и\ \ \textquotesingle{}u\textbackslash{}left(2,\textbackslash{};\textbackslash{}theta\textbackslash{}right)=\textbackslash{}sin\textbackslash{}theta\textquotesingle{}\ .}

Наша лекция завершается кратким обсуждением внешней за-дачи Дирихле для
круга.

\section{Внешняя задача
Дирихле}\label{ux432ux43dux435ux448ux43dux44fux44f-ux437ux430ux434ux430ux447ux430-ux434ux438ux440ux438ux445ux43bux435}

Внешняя задача Дирихле

\begin{verbatim}
\begin{array}{ll} (\text{UBP}) & u_{rr} + \frac{1}{r} u_r + \frac{1}{r^2} u_{\theta\theta} = 0, \quad 1 < r < \infty, \\ (\text{GU}) & u(1, \theta) = g(\theta), \quad 0 \leqslant \theta < 2\pi, \end{array}
\end{verbatim}

решается точно так же, как решалась внутренняя задача в лекции 33. За
исключением того, что теперь необходимо отбросить те решения, которые не
ограничены при стремлении г к бесконечности, т. е.

\begin{verbatim}
\ln r
\end{verbatim}

\begin{verbatim}
,  r^n \cos(n\theta)   H   r^n \sin(n\theta) .

Следовательно, в качестве решения остается взять функцию

(34.4)
\end{verbatim}

\texttt{u(r,\ \textbackslash{}theta)\ =\ \textbackslash{}sum\_\{n=0\}\^{}\{\textbackslash{}infty\}\ r\^{}\{-n\}\ {[}a\_n\ \textbackslash{}cos(n\textbackslash{}theta)\ +\ b\_n\ \textbackslash{}sin(n\textbackslash{}theta){]},}

\begin{verbatim}
где коэффициенты  a_n  и  b_n  вычисляются по тем же формулам, что и раньше:
\end{verbatim}

\texttt{a\_0\ =\ \textbackslash{}frac\{1\}\{2\textbackslash{}pi\}\ \textbackslash{}int\_0\^{}\{2\textbackslash{}pi\}\ g(\textbackslash{}theta)\ d\textbackslash{}theta,}
\texttt{a\_n\ =\ \textbackslash{}frac\{1\}\{\textbackslash{}pi\}\ \textbackslash{}int\_0\^{}\{2\textbackslash{}pi\}\ g(\textbackslash{}theta)\ \textbackslash{}cos(n\textbackslash{}theta)\ d\textbackslash{}theta,\ \textbackslash{}quad\ b\_n\ =\ \textbackslash{}frac\{1\}\{\textbackslash{}pi\}\ \textbackslash{}int\_0\^{}\{2\textbackslash{}pi\}\ g(\textbackslash{}theta)\ \textbackslash{}sin(n\textbackslash{}theta)\ d\textbackslash{}theta.}

\begin{verbatim}
Другими словами, мы просто разлагаем функцию  u\left(1,\;\theta\right)=g\left(\theta\right)  в ряд Фурье
\end{verbatim}

\texttt{g(\textbackslash{}theta)\ =\ \textbackslash{}sum\_\{n=0\}\^{}\{\textbackslash{}infty\}\ {[}a\_n\ \textbackslash{}cos(n\textbackslash{}theta)\ +\ b\_n\ \textbackslash{}sin(n\textbackslash{}theta){]},}

\begin{verbatim}
а затем каждый член этого ряда домножаем на коэффициент  r^{-n} . Чтобы теснее познакомиться с этим решением, приведем два примера.

Примеры внешней задачи Дирихле

Внешняя задача

(УЧП)
\end{verbatim}

\texttt{\textbackslash{}Delta\ u\ =\ 0}

\texttt{,\ \ \textquotesingle{}1\ \textless{}\ r\ \textless{}\ \textbackslash{}infty\textquotesingle{}\ ,}
\texttt{(ГУ)\ \ \textquotesingle{}u(1,\ \textbackslash{}theta)\ =\ 1\ +\ \textbackslash{}sin\ \textbackslash{}theta\ +\ \textbackslash{}cos\ (3\textbackslash{}theta)\textquotesingle{}\ ,\ \ \textquotesingle{}0\ \textbackslash{}le\ \textbackslash{}theta\ \textbackslash{}le\ 2\textbackslash{}pi\textquotesingle{}\ ,}

имеет решение

\begin{verbatim}
u(r, \theta) = 1 + \frac{1}{r} \sin \theta + \frac{1}{r^3} \cos (3\theta).
\end{verbatim}

ПРИМЕР 2

Внешняя задача

(УЧП)

\begin{verbatim}
\Delta u = 0
\end{verbatim}

\begin{verbatim}
,  1 < r < \infty ,  
(ГУ)  u(r, \theta) = \cos(4\theta) ,  0 \le \theta \le 2\pi ,

имеет решение
\end{verbatim}

\texttt{u(r,\ \textbackslash{}theta)\ =\ \textbackslash{}frac\{1\}\{r\^{}4\}\ \textbackslash{}cos(4\textbackslash{}theta)}

\begin{verbatim}
Читатель должен попытаться представить, как выглядит это решение

#### ЗАМЕЧАНИЯ

Решение внешней задачи Дирихле при произвольном радиусе R
\end{verbatim}

\texttt{\textbackslash{}begin\{array\}\{ll\}\ (\textbackslash{}text\{YHI\})\ \&\ \textbackslash{}Delta\ u\ =\ 0,\ \&\ R\ \textless{}\ r\ \textless{}\ \textbackslash{}infty,\ \textbackslash{}\textbackslash{}\ (\textbackslash{}text\{\textbackslash{}GammaY\})\ \&\ u\ (R,\ \textbackslash{}theta)\ =\ g\ (\textbackslash{}theta),\ \&\ 0\ \textbackslash{}leqslant\ \textbackslash{}theta\ \textbackslash{}leqslant\ 2\textbackslash{}pi,\ \textbackslash{}end\{array\}}

\begin{verbatim}
имеет вид
\end{verbatim}

\texttt{u(r,\ \textbackslash{}theta)\ =\ \textbackslash{}sum\_\{n=0\}\^{}\{\textbackslash{}infty\}\ (r/R)\^{}\{-n\}\ {[}a\_n\ \textbackslash{}cos(n\textbackslash{}theta)\ +\ b\_n\ \textbackslash{}sin(n\textbackslash{}theta){]}.}

\begin{verbatim}
2 Существует только два решения двумерного уравнения Лапласа, которые не зависят от  \theta , а зависят только от r. Это константа c и  \ln r  Потенциал  \ln r  играет большую роль в теорин и называется логарифмическим потенциалом. Ниже мы подробнее остановимся на его свойствах

#### ЗАДАЧИ

1 Решите задачу Дирихле

(УЧП)
\end{verbatim}

\texttt{\textbackslash{}Delta\ u\ =\ 0}

, \texttt{1\ \textless{}\ r\ \textless{}\ 2} ,\\
(ГУ)

\begin{verbatim}
\begin{cases} u(1, \theta) = \cos \theta, \\ u(2, \theta) = \sin \theta. \end{cases}
\end{verbatim}

\begin{enumerate}
\def\labelenumi{\arabic{enumi}.}
\setcounter{enumi}{1}
\tightlist
\item
  Найдите решение внешней задачи Дирихле
\end{enumerate}

\begin{verbatim}
(\forall \Pi) \quad \Delta u = 0, \quad 1 < r < \infty,
\end{verbatim}

при следующих граничных условиях:

\begin{enumerate}
\def\labelenumi{\alph{enumi})}
\tightlist
\item
\end{enumerate}

\begin{verbatim}
u(1, \theta) = 1
\end{verbatim}

\begin{verbatim}
,

6)
\end{verbatim}

\texttt{u(1,\ \textbackslash{}theta)\ =\ 1\ +\ \textbackslash{}cos(3\textbackslash{}theta)}

,

\begin{enumerate}
\def\labelenumi{\Alph{enumi})}
\setcounter{enumi}{1}
\tightlist
\item
\end{enumerate}

\begin{verbatim}
u(1, \theta) = \sin(\theta) + \cos(3\theta)
\end{verbatim}

\begin{verbatim}
,

B)
\end{verbatim}

\texttt{u(1,\ \textbackslash{}theta)\ =\ \textbackslash{}sin(\textbackslash{}theta)\ +\ \textbackslash{}cos(3\textbackslash{}theta)}

,\\
r)

\begin{verbatim}
u(1, \theta) =\begin{cases} 1, & 0 \le \theta < \pi, \\ 0, & \pi \le \theta < 2\pi. \end{cases}
\end{verbatim}

3 Рещение внешней задачи Неймана

\begin{verbatim}
\begin{array}{ll} (\forall \Pi \Pi) & \Delta u = 0, \quad 1 < r < \infty, \\ (\Gamma \forall) & \frac{\partial u}{\partial r} (1, \theta) = g(\theta), \quad 0 \leqslant \theta \leqslant 2\pi, \end{array}
\end{verbatim}

имеет тот же вид, что и решение задачи Дирихле, т. е.

\begin{verbatim}
u(r, \theta) = \sum_{n=0}^{\infty} r^{-n} [a_n \cos(n\theta) + b_n \sin(n\theta)],
\end{verbatim}

но теперь коэффициенты \texttt{a\_n} и \texttt{b\_n} должны
удовлетворять новым граничным условиям. Подставьте это решение в
граничное условие

\begin{verbatim}
\frac{\partial u}{\partial r}(1, \theta) = \sin \theta
\end{verbatim}

и получите решение задачи

\begin{verbatim}
\begin{cases} \Delta u = 0, & 1 < r < \infty, \\ \frac{\partial u}{\partial r}(1, \theta) = \sin \theta. \end{cases}
\end{verbatim}

Проверьте ваше решение. Конечно, если вы прибавите к своему решению
какую-нибудь константу, то снова получите решение.

\begin{enumerate}
\def\labelenumi{\arabic{enumi}.}
\setcounter{enumi}{3}
\tightlist
\item
  Подставьте общее решение (34.2) в граничные условия
\end{enumerate}

\begin{verbatim}
\begin{cases} u(R_1, \theta) = g_1(\theta), \\ u(R_2, \theta) = g_2(\theta) \end{cases}
\end{verbatim}

и получите уравнения (34.3).

\section{Лекція 35. РІВНЯННЯ ЛАПЛАСА В СФЕРИЧНИХ КООРДИНАТАХ (СФЕРИЧНІ
ГАРМОНІКИ)}\label{ux43bux435ux43aux446ux456ux44f-35.-ux440ux456ux432ux43dux44fux43dux43dux44f-ux43bux430ux43fux43bux430ux441ux430-ux432-ux441ux444ux435ux440ux438ux447ux43dux438ux445-ux43aux43eux43eux440ux434ux438ux43dux430ux442ux430ux445-ux441ux444ux435ux440ux438ux447ux43dux456-ux433ux430ux440ux43cux43eux43dux456ux43aux438}

ЦЕЛЬ ЛЕКЦИИ: Найти такие частные решения уравнения Лапласа в сферических
координатах, чтобы из них можно было строить решения различных задач
(типа задач Дирихле, Неймана). Кроме того, мы решим внутреннюю задачу
Дирихле

(УЧП)

\begin{verbatim}
(r^2u_r)_r + \frac{1}{\sin\varphi} \left[\sin\varphi u_q\right]_{\varphi} + \frac{1}{\sin^2\varphi} u_{\theta\theta} = 0^{-1}
\end{verbatim}

\begin{verbatim}
),

(ГУ)
\end{verbatim}

\texttt{u(1,\ \textbackslash{}theta,\ \textbackslash{}varphi)\ =\ g(\textbackslash{}varphi),\ \textbackslash{}quad\ 0\ \textbackslash{}le\ \textbackslash{}varphi\ \textbackslash{}le\ \textbackslash{}pi,}

\begin{verbatim}
для одного частного случая, когда потенциал на границе зависит только от  \phi .

<sup>1)</sup> Здесь левая часть уравнения есть  r^2 \Lambda u  в сферических координатах.—  \Pi \rho u m .  \rho e \partial_v 

Для решения задачи нам придется разложить потенциал на границе  g(\varphi)  в ряд по сферическим гармоникам
\end{verbatim}

\texttt{g(\textbackslash{}varphi)\ =\ \textbackslash{}sum\_\{n=0\}\^{}\{\textbackslash{}infty\}\ a\_n\ P\_n(\textbackslash{}cos\ \textbackslash{}varphi),}

\begin{verbatim}
где частные решения уравнения Лапласа, так называемые сферические гармоники  P_n  ( \cos \varphi ), являются полиномами Лежандра от  \cos \varphi . После того как разложение произведено, решение находится сразу:
\end{verbatim}

\texttt{U(R,\ \textbackslash{}varphi)\ =\ \textbackslash{}sum\_\{n=0\}\^{}\{\textbackslash{}infty\}\ a\_n\ r\^{}n\ P\_n(\textbackslash{}cos\ \textbackslash{}varphi).}

\begin{verbatim}
Аналогично находится решение внешней задачи Дирихле:
\end{verbatim}

\texttt{U(R,\ \textbackslash{}varphi)\ =\ \textbackslash{}sum\_\{n=0\}\^{}\{\textbackslash{}infty\}\ \textbackslash{}frac\{a\_n\}\{r\^{}\{n+1\}\}\ P\_n(\textbackslash{}cos\ \textbackslash{}varphi).}

\begin{verbatim}
Одной из важных задач физики является задача определения потенциала внутри или вие сферы, если задан потенциал на поверхности сферы. Для внутренней задачи мы должны найти функцию, которая является решением задачи

(35 1)
\end{verbatim}

\texttt{(Y\ \textbackslash{}Psi\ \textbackslash{}Pi)\ \textbackslash{}quad\ (r\^{}2\ u\_r)\_r\ +\ \textbackslash{}frac\{1\}\{\textbackslash{}sin\ \textbackslash{}varphi\}\ {[}\textbackslash{}sin\ \textbackslash{}varphi\ \textbackslash{},\ u\_\textbackslash{}varphi{]}\_\textbackslash{}varphi\ +\ \textbackslash{}frac\{1\}\{\textbackslash{}sin\^{}2\ \textbackslash{}varphi\}\ \textbackslash{},\ u\_\{\textbackslash{}theta\textbackslash{}theta\}\ =\ 0,}
\texttt{(\textbackslash{}Gamma\ Y)\ \textbackslash{}quad\ u(1,\ 0,\ \textbackslash{}varphi)\ =\ g(\textbackslash{}theta,\ \textbackslash{}varphi),\ \textbackslash{}quad\ -\textbackslash{}pi\ \textbackslash{}leqslant\ \textbackslash{}theta\ \textbackslash{}leqslant\ \textbackslash{}pi,\ \textbackslash{}quad\ 0\ \textbackslash{}leqslant\ \textbackslash{}varphi\ \textbackslash{}leqslant\ \textbackslash{}pi.}

\begin{verbatim}
Обратим внимание на то, что сферический лапласнан записан не так, как раньше. Эта форма записи более компактна и более удобна для использования. Другая типичная интерпретация этой задачи такова: найти распределение температуры внутри сферы, если температура на границе сферы задана. Очень часто функция  g(\theta, \phi)  имеет специальный вид, так что нет необходимости решать задачу в самом общем виде.

В этой лекции мы остановимся на двух важных частных случаях. В одном случае функция  g(\theta, \phi)  будет константой, а в другом — будет зависеть только от угла  \phi  (этот угол отсчитывается от северного полюса сферы).

#### Частные случаи задачи Дирихле

Частный случай  I(g(\theta, \phi) = \text{const}) 

Ясно, что в этом случае решение не зависит от углов ф и 0 и уравнение Лапласа переходит в обыкновенное дифференциальное уравнение
\end{verbatim}

(r\^{}2u\_r)\_r=0.

\begin{verbatim}
Это очень простое уравнение, и читатель может легко решить его. Общее решение этого уравнения имеет вид

(35.2)
\end{verbatim}

\texttt{u(r)\ =\ \textbackslash{}frac\{a\}\{r\}\ +\ b}

.

Другими словами, среди всех потенциалов существует только два таких,
которые не зависят от углов, а зависят только от r. Эти потенциалы---
константа и c/r. Потенциал 1/r играет очень важную роль в физике и
называется ньютоновским потенциалом.

\pandocbounded{\includegraphics[keepaspectratio]{images/_page_264_Picture_9.jpeg}}

РИС. 35.1. Задача Дирихле в щаре

Теперь мы решим две задачи, в каждой из которых потенциал зависит только
от r.

ЗАДАЧА 1 (Потенциал внутри сферы)

(УЧП)

\begin{verbatim}
\Delta u = 0
\end{verbatim}

\begin{verbatim}
,  0 < r < 1 , (ГУ)  u(1, \theta, \varphi) = 3 .

Согласно (35.2), получаем ограниченное решение  u(r, \theta, \phi) = 3 .  3A \square A \square A A = 2  (Потенциал между двумя сферами, каждая из которых обладает постоянным потенциалом)

Предположим, что мы хотим найти стационарное распределение температуры между двумя сферами, которые поддерживаются при различных температурах:
\end{verbatim}

\texttt{(\textbackslash{}text\{YYII\})\ \textbackslash{}quad\ \textbackslash{}Delta\ u\ =\ 0,\ \textbackslash{}quad\ R\_1\ \textless{}\ r\ \textless{}\ R\_2,}
\texttt{(\textbackslash{}text\{\textbackslash{}GammaY\})\ \textbackslash{}quad\ \textbackslash{}begin\{cases\}\ u(R\_1,\ \textbackslash{}theta,\ \textbackslash{}varphi)\ =\ A,\ \textbackslash{}\textbackslash{}\ u(R\_2,\ \textbackslash{}theta,\ \textbackslash{}varphi)\ =\ B.\ \textbackslash{}end\{cases\}}

\begin{verbatim}
Поскольку мы знаем, что общая форма потенциала, не зависящего от  \theta  и  \varphi ,
\end{verbatim}

\texttt{u(r)\ =\ \textbackslash{}frac\{a\}\{r\}\ +\ b,}

\begin{verbatim}
то подставим его в граничное условие и решим получившуюся систему относительно а и b. В результате получаем
\end{verbatim}

\texttt{a\ =\ \textbackslash{}frac\{(A-B)R\_1R\_2\}\{R\_2-R\_1\}}

\texttt{,\ \ \textquotesingle{}b\ =\ \textbackslash{}frac\{R\_2B-R\_1A\}\{R\_2-R\_1\}\textquotesingle{}\ .}

Графики этого потенциала при различных значениях потенциалов A и B
изображены на рис. 35.2.

\pandocbounded{\includegraphics[keepaspectratio]{images/_page_265_Figure_9.jpeg}}

РИС 35 2. Потенциал между двумя концентрическими сферами, одна из
которых обладает потенциалом A, а другая---потенциалом B.

Частный случай 2 (
\texttt{g(\textbackslash{}theta,\ \textbackslash{}varphi)} зависит
только от \texttt{\textbackslash{}varphi} )

В этом случае задача Дирихле записывается в виде

(УЧП)

\begin{verbatim}
(r^2u_r)_r + \frac{1}{\sin\varphi} [\sin\varphi u_\varphi]_\varphi = 0
\end{verbatim}

\begin{verbatim}
,  0 < r < 1 ,  0 \le \varphi \le \pi , (ГУ)  u(1, \theta, \varphi) = g(\varphi) ,  0 \le \varphi \le \pi .

Воспользуемся методом разделення переменных и найдем решения вида
\end{verbatim}

\texttt{u(r,\ \textbackslash{}varphi)\ =\ R(r)\ \textbackslash{}Phi(\textbackslash{}varphi).}

\begin{verbatim}
В результате приходим к двум обыкновенным дифференциальным

уравнениям:
\end{verbatim}

r\^{}2R'\,' + 2rR' - n(n+1)R = 0

\begin{verbatim}
(уравнение Эйлера), [ \sin \phi \Phi' ]' +  n(n+1)\sin \phi \Phi = 0  (уравнение Лежандра).

Немного позже читатель увидит, почему константа разделения выбрана в виде n(n+1).

Уравнение Эйлера мы решим, подставив в него  R(r) = r^x  и определив возможные значения  \alpha . Сделав это, получим
\end{verbatim}

\texttt{\textbackslash{}alpha\ =\ \textbackslash{}begin\{cases\}\ n\ \textbackslash{}\textbackslash{}\ -(n+1).\ \textbackslash{}end\{cases\}}

\begin{verbatim}
Следовательно, общее решение уравнения Эйлера имеет вид  R(r) = ar^n + br^{-(n+1)} .

Уравнение Лежандра решить не так легко. Общий метод решения основан на подстановке
\end{verbatim}

\texttt{x\ =\ \textbackslash{}cos\ \textbackslash{}varphi}

.

Проделав эту замену переменной, получаем новую форму уравнения Лежандра

\begin{verbatim}
(1-x^2)\frac{d^2\Phi}{dx^2} - 2x\frac{d\Phi}{dx} + (n+1)n\Phi = 0, \quad -1 < x < 1.
\end{verbatim}

\texttt{Теперь\ нужно\ решить\ это\ уравнение\ и\ подставить\ в\ решение\ \ \textquotesingle{}x\ =\ \textbackslash{}cos\ \textbackslash{}varphi\textquotesingle{}\ .\ Уравнение\ Лежандра\ является\ линейным\ уравнением\ вгорого\ порядка\ с\ переменными\ коэффициентами.\ Одна\ из\ особенностей\ этого\ уравнения\ в\ том,\ что\ коэффициент\ \ \textquotesingle{}(1-x\^{}2)\textquotesingle{}\ \ при\ \ \textquotesingle{}d\^{}2\textbackslash{}Phi/dx\^{}2\textquotesingle{}\ \ обращается\ в\ нуль\ на\ концах\ отрезка\ \ \textquotesingle{}-1\ \textbackslash{}leqslant\ x\ \textbackslash{}leqslant\ 1\textquotesingle{}\ .\ Уравнения\ этого\ типа\ называются\ сингулярными\ дифференциальными\ уравнениями\ и\ очень\ часто\ решаются\ методом\ Фробеннуса\ Не\ вдаваясь\ в\ подробности\ этого\ метода,\ мы\ приведем\ только\ наиболее\ важный\ результат.\ Ограниченные\ решения\ уравнения\ Лежандра\ могут\ существовать,\ только\ если\ \ \textquotesingle{}n=0,1,2,\textbackslash{}ldots\textquotesingle{}\ \ При\ этом\ оказывается,\ что\ они\ являются\ полиномами\ \ \textquotesingle{}P\_n(x)\textquotesingle{}\ \ (полиномами\ Лежандра)}

\begin{verbatim}
\begin{array}{llll} n=0, & P_{0}(x)=1, \\ n=1, & P_{1}(x)=x, \\ n=2, & P_{2}(x)=\frac{1}{2}(3x^{2}-1), \\ n=3, & P_{3}(x)=\frac{1}{2}(5x^{3}-3x), & -1\leqslant x\leqslant 1, \\ & \vdots & \vdots & \vdots \\ n, & P_{n}(x)=\frac{1}{2^{n}n!}\frac{d^{n}}{dx^{n}}[(x^{2}-1)^{n}] & \left(\frac{\text{formula}}{\text{Rodriga}}\right). \end{array}
\end{verbatim}

Графики нескольких полиномов Лежандра изображены на рис. 35.3. Теперь
ясно, что ограниченные решения уравнений

\begin{verbatim}
r^2R'' + 2rR' - n(n+1)R = 0,
\end{verbatim}

\texttt{0\ \textless{}\ r\ \textless{}\ 1,}
\texttt{0\ \textbackslash{}le\ \textbackslash{}phi\ \textbackslash{}le\ \textbackslash{}pi,}

имеют вид

\begin{verbatim}
R(r) = ar^n
\end{verbatim}

\begin{verbatim}
,  
 \Phi(\varphi) = aP_n(\cos \varphi) .

Функция  P_n(\cos\phi) —эго полином Лежандра n-го порядка, в котором x заменен на  \cos\phi . И наконец, ищем решение исходной задачи в виде ряда

(35.3)
\end{verbatim}

\texttt{U(R,\ \textbackslash{}varphi)\ =\ \textbackslash{}sum\_\{N=0\}\^{}\{\textbackslash{}infty\}\ a\_n\ r\^{}n\ P\_n(\textbackslash{}cos\ \textbackslash{}varphi),}

\begin{verbatim}
коэффициенты которого выбираем так, чтобы удовлетворить граничному условию  u\left(1,\,\varphi\right)=g\left(\varphi\right) . Подстановка решения (35-3)

![](images/_page_267_Figure_9.jpeg)

РИС 35 3. Полиномы Лежандра  P_{n}\left(x\right) .

в граничные условия дает
\end{verbatim}

\texttt{\textbackslash{}sum\_\{n=0\}\^{}\{\textbackslash{}infty\}\ a\_n\ P\_n\ (\textbackslash{}cos\ \textbackslash{}varphi)\ =\ g(\textbackslash{}varphi).}

\begin{verbatim}
Если умножить обе части этого соотношения на  P_m (\cos \varphi) \sin \varphi  и проинтегрировать по  \varphi  от 0 до  \pi , то получим
\end{verbatim}

\texttt{\textbackslash{}int\_\{0\}\^{}\{\textbackslash{}pi\}\ g(\textbackslash{}varphi)\ P\_\{m\}(\textbackslash{}cos\ \textbackslash{}varphi)\ \textbackslash{}sin\ \textbackslash{}varphi\ \textbackslash{},\ d\textbackslash{}varphi\ =\ \textbackslash{}sum\_\{n=0\}\^{}\{\textbackslash{}infty\}\ a\_\{n\}\ \textbackslash{}int\_\{0\}\^{}\{\textbackslash{}pi\}\ P\_\{n\}(\textbackslash{}cos\ \textbackslash{}varphi)\ P\_\{m\}(\textbackslash{}cos\ \textbackslash{}varphi)\ \textbackslash{}sin\ \textbackslash{}varphi\ \textbackslash{},\ d\textbackslash{}varphi\ =}
\texttt{=\ \textbackslash{}sum\_\{n=0\}\^{}\{\textbackslash{}infty\}\ a\_\{n\}\ \textbackslash{}int\_\{-1\}\^{}\{1\}\ P\_\{n\}(x)\ P\_\{m\}(x)\ \textbackslash{},\ dx\ =}
\texttt{=\ \textbackslash{}begin\{cases\}\ 0,\ \&\ n\ \textbackslash{}neq\ m,\ \textbackslash{}\textbackslash{}\ \textbackslash{}frac\{2a\_\{m\}\}\{2m+1\},\ \&\ n\ =\ m.\ \textbackslash{}end\{cases\}}

\begin{verbatim}
Можно убедиться в том, что полиномы Лежандра ортогональны на отрезке [-1, 1]. Следовательно,

(35.4)
\end{verbatim}

\texttt{a\_m\ =\ \textbackslash{}frac\{2m+1\}\{2\}\ \textbackslash{}int\_0\^{}\{\textbackslash{}pi\}\ g(\textbackslash{}varphi)\ P\_m(\textbackslash{}cos\ \textbackslash{}varphi)\ \textbackslash{}sin\ \textbackslash{}varphi\ \textbackslash{},\ d\textbackslash{}varphi\ \textbackslash{}quad\ (m=0,\ 1,\ 2,\ \textbackslash{}ldots).}

\begin{verbatim}
Итак, решение задачи Дирихле (35.1) дается выражением

(35.5)
\end{verbatim}

\texttt{U(R,\ \textbackslash{}varphi)\ =\ \textbackslash{}sum\_\{N=0\}\^{}\{\textbackslash{}infty\}\ a\_n\ r\^{}n\ P\_n(\textbackslash{}cos\ \textbackslash{}varphi),}

\begin{verbatim}
где коэффициенты  a_n  вычисляются по формулам (35.4).

Приведем теперь пример цилиндрически-симметричного потенциала.

Цилиндрически-симметричный потенциал (не зависящий от  \theta )

Предположим, что температура поверхности сферы определяется формулой
\end{verbatim}

\texttt{g(\textbackslash{}varphi)\ =\ 1\ -\ \textbackslash{}cos(2\textbackslash{}varphi),\ \textbackslash{}quad\ 0\ \textbackslash{}leqslant\ \textbackslash{}varphi\ \textbackslash{}leqslant\ \textbackslash{}pi,}

\begin{verbatim}
и предположим, что мы хотим найти температуру внутри сферы. В нашей задаче температура постоянная на линиях постоянной широты (например, температура равна 2 на экваторе:  \varphi = \pi/2 , и Р И

![](images/_page_268_Figure_12.jpeg)

С. 35.4. Зависимость температуры от широты в радианах (цачало отсчета на северном жолюсе).

нулю—на полюсах:  \phi = 0 ,  \phi < \pi ). Чтобы найти u, нужно решить задачу

(УЧП)
\end{verbatim}

\texttt{\textbackslash{}Delta\ u\ =\ 0}

, \texttt{0\ \textless{}\ r\ \textless{}\ 1} ,\\
\texttt{(ГУ)\ \ \textquotesingle{}u(1,\ \textbackslash{}theta,\ \textbackslash{}varphi)\ =\ 1\ -\ \textbackslash{}cos(2\textbackslash{}varphi)\textquotesingle{}\ ,\ \ \textquotesingle{}0\ \textbackslash{}le\ \textbackslash{}varphi\ \textbackslash{}le\ \textbackslash{}pi\textquotesingle{}\ .}

График температуры на поверхности изображен на рис. 35.4. Нам необходимо
найти коэффициенты \texttt{a\_n} для решения (35.5). Чтобы сделать это,
мы должны

\begin{itemize}
\tightlist
\item
  а) либо воспользоваться формулой (35.4),
\item
  б) либо воспользоваться программой, существующей на вашем
  вычислительном центре и предназначенной для вычисления коэффициентов
  разложения по полиномам Лежандра,
\item
  в) либо проявить немного здравого смысла.
\end{itemize}

\pandocbounded{\includegraphics[keepaspectratio]{images/_page_269_Figure_5.jpeg}}

РИС. 35.5. Зависимость температуры от расстояния до центра шара.

Попытаемся воспользоваться последним рецептом. Рассмотрим
тригонометрическое тождество

\begin{verbatim}
\cos(2\varphi) = 2\cos^2\varphi - 1,
\end{verbatim}

которое позволяет переписать температуру на границе в виде

\begin{verbatim}
1 - \cos(2\varphi) = 1 - [2\cos^2\varphi - 1] =
= 1 - \frac{2}{3}[3\cos^2\varphi - 1] + \frac{1}{3} =
= \frac{4}{3}P_{\theta}(\cos\varphi) - \frac{4}{3}P_{\theta}(\cos\varphi).
\end{verbatim}

Мы получили разложение функции \texttt{g(\textbackslash{}phi)} в ряд по
полиномам Лежандра. Теперь решение задачи легко находится

\section{Лекція 36. НЕОДНОРІДНА ЗАДАЧА ДІРІХЛЕ (ФУНКЦІЯ
ГРІНА)}\label{ux43bux435ux43aux446ux456ux44f-36.-ux43dux435ux43eux434ux43dux43eux440ux456ux434ux43dux430-ux437ux430ux434ux430ux447ux430-ux434ux456ux440ux456ux445ux43bux435-ux444ux443ux43dux43aux446ux456ux44f-ux433ux440ux456ux43dux430}

МЕТА ЛЕКЦІЇ: Показати, як гетерогенну задачу Діріхле можна розв'язати за
допомогою функції Гріна (вихідної функції). У цьому важливому методі
права частина рівняння розглядається як певна вхідна дія і розкладається
на неперервний набір дельтаподібних джерел, розподілених по заданій
площі. Потім визначається відповідь системи на кожне таке джерело
(функція Гріна), і всі відповіді підсумовуються (інтегровані). У
результаті ми отримуємо повне розв'язання проблеми.

Досить загальною задачею прикладної математики є визначення потенціалу в
деякій області простору, якщо всередині цієї області задано розподіл
джерел \texttt{f(x,y)}. В електростатиці потенціал у області \texttt{D}
визначається розподілом зарядів із щільністю \texttt{f(x,y)}, а в теорії
теплопровідності аналогічно задається джерелами тепла. Типова двовимірна
задача в крузі приводить до рівняння Пуассона.

\begin{verbatim}
\begin{aligned}(36.1)\quad (\text{УЧП})\quad &u_{rr} + \frac{1}{r}u_r + \frac{1}{r^2}u_{\theta\theta} = f(r,\theta),&& 0 \le r < 1,\quad 0 \le \theta < 2\pi, \\(\text{ГУ})\quad &u(1,\theta)=0,&& 0 \le \theta \le 2\pi.\end{aligned}
\end{verbatim}

Ми не випадково обрали нульові граничні умови. Якщо б потрібно було
розв'язати задачу в загальному випадку, з ненульовими граничними умовами
та неоднорідним рівнянням, то внесок граничних даних можна було б знайти
за допомогою інтеграла Пуассона з лекції 33.

Щоб отримати наочне уявлення про неоднорідні еліптичні рівняння,
розглянемо задачу

\begin{verbatim}
\begin{aligned}(\text{УЧП})\quad & \Delta u = -q, && 0 < r < 1,\quad 0 \le \theta \le 2\pi, \\(\text{ГУ})\quad & u(1,\theta)=0, && 0 \le \theta \le 2\pi,\end{aligned}
\end{verbatim}

де \texttt{q\textgreater{}0} --- стала. На межі потенціал підтримується
рівним нулю, а лапласіан є сталим у всій області. Рівняння Пуассона
означає, що графік \texttt{u(r,\textbackslash{}theta)} має всюди
однакову опуклість. Геометрично його можна уявляти як тонку мембрану з
нерухомим краєм, яку прогинає зовнішня сила.

Тепер перейдемо до головної частини лекції: визначення функції Гріна і
побудови розв'язку рівняння (36.1). Але спершу введемо потенціали
точкових джерел і стоків.

\section{Потенціали точкових джерел і
стоків}\label{ux43fux43eux442ux435ux43dux446ux456ux430ux43bux438-ux442ux43eux447ux43aux43eux432ux438ux445-ux434ux436ux435ux440ux435ux43b-ux456-ux441ux442ux43eux43aux456ux432}

Щоб розв'язати неоднорідне лінійне рівняння, достатньо зрозуміти
відповідь системи на одне точкове джерело. Тоді загальний розв'язок
можна отримати підсумовуванням таких елементарних внесків. У
теплопровідності точкове джерело є місцем виділення тепла, а в
електростатиці --- точковим зарядом.

Припустимо, що джерело інтенсивності \texttt{+q} розміщене в початку
координат. Тоді потік є радіальним, а повний потік через коло радіуса
\texttt{r} дорівнює

\begin{verbatim}
-\int_0^{2\pi} u_r(r)\,r\,d\theta = -2\pi r\,u_r(r).
\end{verbatim}

За законом збереження цей потік повинен дорівнювати кількості тепла, що
генерується всередині кола, тобто

\begin{verbatim}
-2\pi r\,u_r(r)=q.
\end{verbatim}

Розв'язуючи це рівняння відносно функції \texttt{u(r)}, отримуємо
двовимірний потенціал точкового джерела.

\begin{figure}
\centering
\pandocbounded{\includegraphics[keepaspectratio,alt={Рис. 36.1. Радіальний потік тепла з точкового джерела.}]{images/_page_273_Picture_10.jpeg}}
\caption{Рис. 36.1. Радіальний потік тепла з точкового джерела.}
\end{figure}

Рис. 36.1. Радіальний потік тепла з точкового джерела; A --- тепло
виділяється в початку.

\begin{verbatim}
u(r) = \frac{-q}{2\pi} \ln r = \frac{q}{2\pi} \ln \frac{1}{r}.
\end{verbatim}

В електростатиці різниця потужності \texttt{u(B)} до \texttt{u(A)}
дорівнює роботі, яку потрібно виконати для переміщення одного
позитивного заряду з точки A в точку B.

\begin{figure}
\centering
\pandocbounded{\includegraphics[keepaspectratio,alt={Рис. 36.2. Потенціал точкового заряду у двовимірному випадку.}]{images/_page_273_Figure_15.jpeg}}
\caption{Рис. 36.2. Потенціал точкового заряду у двовимірному випадку.}
\end{figure}

Рис. 36.2. Потенціал точкового заряду у двовимірному випадку. Тут
\texttt{\textbackslash{}alpha} --- робота, потрібна для переміщення
одиничного заряду з точки A до точки B.

Стік можна інтерпретувати як від'ємне джерело, так що стік \texttt{-q}
створює потенціал

\begin{verbatim}
u\left(r\right)=-\frac{q}{2\pi}\ln\frac{1}{r}.
\end{verbatim}

На цьому завершується обговорення потенціалів точкових зарядів. Тепер ми
готові розв'язати неоднорідне рівняння за допомогою методу функції
Гріна.

\section{Рівняння Пуассона в
крузі}\label{ux440ux456ux432ux43dux44fux43dux43dux44f-ux43fux443ux430ux441ux441ux43eux43dux430-ux432-ux43aux440ux443ux437ux456}

Розглянемо важливу задачу.

\begin{verbatim}
(36.2)\qquad\begin{aligned}(\text{УЧП})\quad & u_{rr} + \frac{1}{r} u_r + \frac{1}{r^2} u_{\theta\theta} = f(r, \theta), && 0 < r < 1,\ 0 \le \theta \le 2\pi, \\(\text{ГУ})\quad & u(1, \theta) = 0, && 0 \le \theta \le 2\pi.\end{aligned}
\end{verbatim}

Метод функції Гріна складається з двох кроків.

\begin{enumerate}
\def\labelenumi{\arabic{enumi}.}
\item
  Встановіть заряд точкової одиниці в точці
  \texttt{(\textbackslash{}rho,\ \textbackslash{}varphi)} і знайдіть
  \texttt{G(r,\ \textbackslash{}theta,\ \textbackslash{}rho,\ \textbackslash{}varphi)}
  потенціал, створений цим зарядом і який на межі дорівнює нулю.
\item
  Підсумуйте окремі відповіді,
  \texttt{G(r,\ \textbackslash{}theta,\ \textbackslash{}rho,\ \textbackslash{}phi)},
  зважені правильною частиною (щільністю заряду)
  \texttt{f(r,\ \textbackslash{}theta)}, вздовж усього кола. У
  результаті ми отримуємо те рішення, яке шукаємо
\end{enumerate}

\begin{verbatim}
u\left(r,\theta\right) = \int_{0}^{2\pi} \int_{0}^{1} G\left(r,\theta,\rho,\varphi\right) f\left(\rho,\varphi\right)\rho\, d\rho\, d\varphi.
\end{verbatim}

Тепер знайдемо вихідну функцію для нашого завдання.

\section{\texorpdfstring{Означення функції джерела
\texttt{G(r,\ \textbackslash{}theta,\ \textbackslash{}rho,\ \textbackslash{}phi)}}{Означення функції джерела G(r, \textbackslash theta, \textbackslash rho, \textbackslash phi)}}\label{ux43eux437ux43dux430ux447ux435ux43dux43dux44f-ux444ux443ux43dux43aux446ux456ux457-ux434ux436ux435ux440ux435ux43bux430-gr-theta-rho-phi}

Замінимо праву частину \texttt{f(r,\textbackslash{}theta)} рівняння
точковим джерелом одиничної величини, розташованим у точці
\texttt{(\textbackslash{}rho,\ \textbackslash{}varphi)}. Математично
таке джерело описується дельта-функцією. Наше завдання --- знайти
потенціал цього джерела за умови, що на межі кола потенціал дорівнює
нулю. Така функція і називається функцією Гріна, або функцією джерела.

Потенціал точкового заряду дорівнює

\begin{verbatim}
\frac{1}{2\pi} \ln \left( \frac{1}{r} \right),
\end{verbatim}

де \texttt{r} --- відстань від заряду до точки спостереження.

Шуканий потенціал має такий фізичний зміст:

\begin{enumerate}
\def\labelenumi{\arabic{enumi}.}
\tightlist
\item
  У теорії теплопровідності це стаціонарна температура всередині круга,
  якщо джерело тепла розташоване в точці
  \texttt{(\textbackslash{}rho,\textbackslash{}varphi)}, а на межі
  підтримується нульова температура.
\item
  Для мембрани це відхилення від положення рівноваги, якщо край
  закріплено на нулі, а в точці
  \texttt{(\textbackslash{}rho,\textbackslash{}varphi)} прикладено
  зосереджене навантаження.
\item
  В електростатиці це розподіл потенціалу всередині круга, якщо точковий
  позитивний заряд розташований у точці
  \texttt{(\textbackslash{}rho,\textbackslash{}varphi)}, а коло
  заземлено.
\end{enumerate}

\begin{figure}
\centering
\pandocbounded{\includegraphics[keepaspectratio,alt={Рис. 36.3. Функція Гріна для точкового джерела.}]{images/_page_275_Picture_11.jpeg}}
\caption{Рис. 36.3. Функція Гріна для точкового джерела.}
\end{figure}

Рис. 36.3. Функція Гріна для точкового джерела, розташованого в точці
\texttt{(\textbackslash{}rho,\ \textbackslash{}phi)}, всередині
одиничного кола.

Тепер побудуємо функцію Гріна; за формою вона буде подібна до графіка на
рис. 36.3.

\section{Побудова
розв'язку}\label{ux43fux43eux431ux443ux434ux43eux432ux430-ux440ux43eux437ux432ux44fux437ux43aux443}

\textbf{КРОК 1.} Оскільки

\begin{verbatim}
\frac{1}{2\pi} \ln \frac{1}{R}
\end{verbatim}

--- це потенціал поля в точці \texttt{P=(r,\textbackslash{}theta)}, який
створюється одиничним точковим зарядом, розміщеним у точці
\texttt{Q=(\textbackslash{}rho,\textbackslash{}phi)}. Залишається лише
модифікувати цю функцію так, щоб на межі кола вона перетворювалася на
нуль.

\textbf{КРОК 2.} З експерименту відомо, що рівнопотенціальні лінії
точкових зарядів є колами (рис. 36.4).

Основна ідея побудови функції Гріна така: потрібно розмістити інший,
від'ємний, заряд поза колом так, щоб сумарний потенціал на колі
\texttt{r=1} був сталим. Потім цю константу можна відняти й отримати
нульовий потенціал на межі.

\begin{figure}
\centering
\pandocbounded{\includegraphics[keepaspectratio,alt={Рис. 36.4. Потенціал поля, утвореного двома протилежно зарядженими частинками.}]{images/_page_276_Picture_3.jpeg}}
\caption{Рис. 36.4. Потенціал поля, утвореного двома протилежно
зарядженими частинками.}
\end{figure}

Рис. 36.4. Потенціал поля, утвореного двома протилежно зарядженими
частинками; A і B --- лінії сталого потенціалу.

\begin{figure}
\centering
\pandocbounded{\includegraphics[keepaspectratio,alt={Рис. 36.5. Заряди Q і Q̄ створюють сталий потенціал на колі r=1.}]{images/_page_276_Picture_5.jpeg}}
\caption{Рис. 36.5. Заряди Q і Q̄ створюють сталий потенціал на колі
r=1.}
\end{figure}

Рис. 36.5. Заряди \texttt{Q} і \texttt{\textbackslash{}overline\{Q\}}
створюють сталий потенціал на колі \texttt{r=1}.

Легко побачити, що якщо негативний заряд розміщений у точці
\texttt{\textbackslash{}bar\{Q\}\ =\ (\textbackslash{}bar\{\textbackslash{}rho\},\ \textbackslash{}bar\{\textbackslash{}phi\})\ =\ (1/\textbackslash{}rho,\ \textbackslash{}phi)},
то повний потенціал

\begin{verbatim}
u(r, \theta) = \frac{1}{2\pi} \ln \frac{1}{R} - \frac{1}{2\pi} \ln \frac{1}{\overline{R}}
\end{verbatim}

буде сталою на колі \texttt{r=1}. Тут відстані від зарядів до точки
спостереження позначено через \texttt{R} і
\texttt{\textbackslash{}overline\{R\}} (див. рис. 36.5). Можна показати,
що значення потенціалу на колі \texttt{r=1} дорівнює

\begin{verbatim}
-\frac{1}{2\pi}\ln\rho
\end{verbatim}

(додатна константа).

З урахуванням цього результату будуємо функцію Гріна

\begin{verbatim}
(36.3)\qquad G(r, \theta, \rho, \phi) =\frac{1}{2\pi} \ln \frac{1}{R}- \frac{1}{2\pi} \ln \frac{1}{\overline{R}}+ \frac{1}{2\pi} \ln \rho
\end{verbatim}

Отже, функція Гріна складається з потенціалу додатного заряду в точці
\texttt{Q}, потенціалу від'ємного заряду в точці
\texttt{\textbackslash{}overline\{Q\}} та коригувальної константи, яка
забезпечує нульове граничне значення.

де

\begin{verbatim}
\begin{aligned}R &= \sqrt{r^2 - 2\rho r \cos(\theta - \phi) + \rho^2},\\\overline{R} &= \sqrt{r^2 - 2\frac{r}{\rho}\cos(\theta - \phi) + \frac{1}{\rho^2}}.\end{aligned}
\end{verbatim}

Це формули для відстаней між відповідними точками в полярних
координатах. Щоб знайти розв'язок початкової задачі, потрібно
просумувати внесок усіх таких елементарних джерел. Переходимо до
останнього кроку.

\textbf{КРОК 3.} Суперпозиція імпульсних характеристик. Цей крок досить
простий. З відношення

\begin{verbatim}
u(r, \theta) = \int_{0}^{2\pi} \int_{0}^{1} G(r, \theta, \rho, \varphi) f(\rho, \varphi) \rho\, d\rho\, d\varphi
\end{verbatim}

отримуємо

\begin{verbatim}
(36.4)\qquad u(r, \theta) = \frac{1}{2\pi} \int_{0}^{2\pi} \int_{0}^{1}\ln \left(\frac{\rho \overline{R}}{R}\right) f(\rho, \varphi)\,\rho\, d\rho\, d\varphi.
\end{verbatim}

Формула (36.4) виражає розв'язок задачі Діріхле для рівняння Пуассона
всередині кола через функцію Гріна цієї задачі. Якщо щільність заряду
\texttt{f(r,\ \textbackslash{}theta)} відома, то інтеграл можна
обчислити, наприклад, чисельно.

\section{ЗАУВАЖЕННЯ}\label{ux437ux430ux443ux432ux430ux436ux435ux43dux43dux44f-7}

\textbf{Зауваження 1.} Задачу

\begin{verbatim}
\begin{aligned}(\text{УЧП}) \quad & \Delta u = 0, \quad 0 < r < 1, \quad 0 \leqslant \theta < 2\pi, \\(\text{ГУ}) \quad & u(1, \theta) = g(\theta), \quad 0 \leqslant \theta \leqslant 2\pi.\end{aligned}
\end{verbatim}

також можна розв'язати за допомогою методу \emph{функції Гріна}. У цьому
випадку розв'язок має вигляд

\begin{verbatim}
u(r, \theta) = \int_{0}^{2\pi} \frac{\partial G}{\partial r}(r, \theta, 1, \varphi)\, g(\varphi)\, d\varphi.
\end{verbatim}

Останнє співвідношення можна переписати в більш обчислювально зручному
вигляді, явно підставивши вираз для
\texttt{\textbackslash{}partial\ G/\textbackslash{}partial\ r}. У
результаті отримуємо інтегральну формулу Пуассона, виведену раніше в
лекції 33.

\begin{verbatim}
(36.5)\qquad u(r, \theta) = \frac{1}{2\pi} \int_{0}^{2\pi}\left[ \frac{1 - r^2}{1 - 2r\cos(\theta - \varphi) + r^2} \right] g(\varphi)\, d\varphi.
\end{verbatim}

\textbf{Зауваження 2.} Розв'язок загальної задачі Діріхле

\begin{verbatim}
\begin{aligned}(\text{УЧП}) \quad & \Delta u = f(r, \theta), \quad 0 < r < 1, \quad 0 \le \theta < 2\pi, \\(\text{ГУ}) \quad & u(1, \theta) = g(\theta), \quad 0 \le \theta \le 2\pi.\end{aligned}
\end{verbatim}

знаходиться як сума розв'язків (36.4) і (36.5).

\textbf{Зауваження 3.} Метод функцій Гріна дозволяє отримати розв'язки
багатьох задач в областях різної форми. Однак для кожної області, а
точніше для кожного оператора та граничної умови, потрібно знайти власну
функцію Гріна, і це не завжди легко зробити.

\textbf{Зауваження 4.} Щоб дійсно знайти розв'язок за формулою (36.4), у
більшості випадків необхідно обчислити інтеграли на комп'ютері.

\section{ЗАВДАННЯ}\label{ux437ux430ux432ux434ux430ux43dux43dux44f-7}

\begin{enumerate}
\def\labelenumi{\arabic{enumi}.}
\item
  Знайти потенціал точкового джерела у тривимірному просторі.
\item
  Знайти функцію Гріна
  \texttt{G(x,y,\textbackslash{}xi,\textbackslash{}eta)} для рівняння
  Лапласа у верхній півплощині \texttt{y\textgreater{}0}. Іншими
  словами, знайти потенціал у точці \texttt{(x,y)}, якщо точковий заряд
  розташований у \texttt{(\textbackslash{}xi,\textbackslash{}eta)}, а
  потенціал на межі \texttt{y=0} дорівнює нулю.
\end{enumerate}

\begin{figure}
\centering
\pandocbounded{\includegraphics[keepaspectratio,alt={Примітка до задачі 2.}]{images/_page_278_Figure_11.jpeg}}
\caption{Примітка до задачі 2.}
\end{figure}

ПРИМІТКА. Якщо від'ємний заряд розмістити в точці
\texttt{\textbackslash{}overline\{Q\}=(\textbackslash{}xi,-\textbackslash{}eta)},
то сумарний потенціал на прямій \texttt{y=0} зануляється. Це і є основа
методу дзеркальних відображень.

\texttt{3.\ Використовуючи\ результати\ задачі\ 2,\ знайдіть\ розв\textquotesingle{}язок\ рівняння\ Пуассона\ \textquotesingle{}\textbackslash{}Delta\ u=-k\textquotesingle{}\ у\ верхній\ півплощині\ з\ нульовою\ граничною\ умовою.}

\begin{enumerate}
\def\labelenumi{\arabic{enumi}.}
\setcounter{enumi}{3}
\item
  Як би ви побудували функцію Гріна для першого квадранта
  \texttt{x\textgreater{}0}, \texttt{y\textgreater{}0}?
\item
  Інший підхід до розв'язання рівняння Пуассона такий. Припустимо, що
  потрібно розв'язати задачу
\end{enumerate}

\begin{verbatim}
\begin{aligned}(\text{УЧП})\quad & \Delta u = 1, && 0 < r < 1,\quad 0 \le \theta < 2\pi,\\(\text{ГУ})\quad & u(1,\theta)=\sin\theta, && 0 \le \theta \le 2\pi.\end{aligned}
\end{verbatim}

Спочатку спробуйте знайти частинний розв'язок підстановкою

\begin{verbatim}
u_p(r,\theta)=Ar^2.
\end{verbatim}

Після цього підставте вираз \texttt{u\_p} в рівняння Пуассона і знайдіть
константу \texttt{A}. Далі подайте шуканий розв'язок у вигляді
\texttt{u=w+u\_p}, де функція \texttt{w(r,\textbackslash{}theta)}
задовольняє відповідну однорідну задачу з модифікованими граничними
умовами. Знайдіть \texttt{w(r,\textbackslash{}theta)} і запишіть
остаточний розв'язок \texttt{u(r,\textbackslash{}theta)}. Ретельно
проаналізуйте отриману відповідь.

\section{Частина 5. ЧИСЕЛЬНІ ТА ПРИБЛИЗНІ
МЕТОДИ}\label{ux447ux430ux441ux442ux438ux43dux430-5.-ux447ux438ux441ux435ux43bux44cux43dux456-ux442ux430-ux43fux440ux438ux431ux43bux438ux437ux43dux456-ux43cux435ux442ux43eux434ux438}

\section{Лекція 37. ЧИСЕЛЬНІ РОЗВ'ЯЗКИ (ЕЛІПТИЧНІ
ЗАДАЧІ)}\label{ux43bux435ux43aux446ux456ux44f-37.-ux447ux438ux441ux435ux43bux44cux43dux456-ux440ux43eux437ux432ux44fux437ux43aux438-ux435ux43bux456ux43fux442ux438ux447ux43dux456-ux437ux430ux434ux430ux447ux456}

МЕТА ЛЕКЦІЇ: Показати, як рівняння в частинних похідних можна звести до
систем алгебраїчних рівнянь, якщо замінити похідні наближеннями
скінченних різниць. Розв'язуючи ці системи ітеративними методами, можна
отримати наближені розв'язки еліптичних задач.

На цей момент ми вже бачили кілька точних методів розв'язання лінійних
рівнянь у частинних похідних. Проте більшість із них працює лише для
простих рівнянь і простих областей. У реальних прикладних задачах цього
часто недостатньо, тому доводиться вдаватися до чисельних методів.

У цій та наступних двох лекціях ми покажемо, як розв'язувати еліптичні,
гіперболічні та параболічні рівняння методом скінченних різниць. Спершу
ознайомимося з самим поняттям скінченних різниць, а потім покажемо, як
застосовувати їх до задачі Діріхле в квадраті.

\section{Апроксимації скінченних
різниць}\label{ux430ux43fux440ux43eux43aux441ux438ux43cux430ux446ux456ux457-ux441ux43aux456ux43dux447ux435ux43dux43dux438ux445-ux440ux456ux437ux43dux438ux446ux44c}

Згадаємо ряд Тейлора для функції \texttt{f(x)}:

\begin{verbatim}
f(x+h)=f(x)+f'(x)h+\frac{f''(x)}{2!}h^2+\dots
\end{verbatim}

Якщо обмежитися першими двома членами, отримаємо

\begin{verbatim}
f(x+h)\cong f(x)+f'(x)h,
\end{verbatim}

звідки випливає права різницева похідна

\begin{verbatim}
(37.1)\qquad f'(x)\cong \frac{f(x+h)-f(x)}{h}.
\end{verbatim}

У розкладі Тейлора можна замінити \texttt{h} на \texttt{-h} і отримати
ліву різницеву похідну:

\begin{verbatim}
(37.2)\qquad f'(x)\cong \frac{f(x)-f(x-h)}{h}.
\end{verbatim}

Розкладаючи \texttt{f(x-h)} і \texttt{f(x+h)} в ряд Тейлора та
відкидаючи члени вищого порядку, отримуємо центральну різницеву похідну:

\begin{verbatim}
(37.3)\qquad f'(x)\simeq \frac{f(x+h)-f(x-h)}{2h}.
\end{verbatim}

Якщо в ряду Тейлора залишити ще один член, то для другої похідної маємо

\begin{verbatim}
(37.4)\qquad f''(x)\cong \frac{f(x+h)-2f(x)+f(x-h)}{h^2}.
\end{verbatim}

Тепер поширимо поняття скінченних різниць на часткові похідні. Виходячи
з розкладів Тейлора для функції двох змінних,

\begin{verbatim}
\begin{aligned}u(x+h,y)&=u(x,y)+u_x(x,y)h+u_{xx}(x,y)\frac{h^2}{2!}+\dots,\\u(x-h,y)&=u(x,y)-u_x(x,y)h+u_{xx}(x,y)\frac{h^2}{2!}-\dots.\end{aligned}
\end{verbatim}

отримуємо такі наближення:

\begin{verbatim}
\begin{aligned}u_x(x,y) &\cong \frac{u(x+h,y)-u(x,y)}{h},\\u_{xx}(x,y) &\cong \frac{u(x+h,y)-2u(x,y)+u(x-h,y)}{h^2},\\u_y(x,y) &\cong \frac{u(x,y+k)-u(x,y)}{k},\\u_{yy}(x,y) &\cong \frac{u(x,y+k)-2u(x,y)+u(x,y-k)}{k^2}.\end{aligned}
\end{verbatim}

У цих формулах часткові похідні апроксимуються правою, центральною та
лівою різницею, але в цій лекції ми будемо використовувати лише
центральні наближення.

\begin{figure}
\centering
\pandocbounded{\includegraphics[keepaspectratio,alt={Фіг. 37.1. Сітка для задачі Діріхле всередині квадрата.}]{images/_page_282_Figure_5.jpeg}}
\caption{Фіг. 37.1. Сітка для задачі Діріхле всередині квадрата.}
\end{figure}

Фіг. 37.1. Сітка для задачі Діріхле всередині квадрата.

Щоб ознайомитися з основними правилами використання апроксимації
скінченних різниць, розв'яжемо просту задачу Діріхле.

\subsection{Розв'язання задачі Діріхле методом скінченної
різниці}\label{ux440ux43eux437ux432ux44fux437ux430ux43dux43dux44f-ux437ux430ux434ux430ux447ux456-ux434ux456ux440ux456ux445ux43bux435-ux43cux435ux442ux43eux434ux43eux43c-ux441ux43aux456ux43dux447ux435ux43dux43dux43eux457-ux440ux456ux437ux43dux438ux446ux456}

Нехай потрібно розв'язати задачу Діріхле

\begin{verbatim}
\begin{aligned}(\text{УЧП})\quad & u_{xx}+u_{yy}=0, && 0<x<1,\quad 0<y<1,\\(\text{ГУ})\quad & u=0, && \text{на верхній та бічних сторонах квадрата},\\& u(x,0)=\sin(\pi x), && 0\le x\le 1.\end{aligned}
\end{verbatim}

Намалюємо сітку на площині \texttt{(x,y)}, як показано на рис. 37.1.
Зручно (особливо якщо ми збираємося користуватися комп'ютером)
використовувати такі позначення:

\begin{verbatim}
\begin{aligned}u(x,y)&=u_{ij},\\u(x,y+k)&=u_{i+1,j},\\u(x,y-k)&=u_{i-1,j},\\u(x-h,y)&=u_{i,j-1},\\u(x+h,y)&=u_{i,j+1},\\u_x(x,y)&\cong \frac{u_{i,j+1}-u_{i,j-1}}{2h},\\u_y(x,y)&\cong \frac{u_{i+1,j}-u_{i-1,j}}{2k},\\u_{xx}(x,y)&\cong \frac{u_{i,j+1}-2u_{i,j}+u_{i,j-1}}{h^2},\\u_{yy}(x,y)&\cong \frac{u_{i+1,j}-2u_{i,j}+u_{i-1,j}}{k^2}.\end{aligned}
\end{verbatim}

Основна ідея розв'язання задачі Діріхле базується на заміні часткових
похідних у рівнянні Лапласа їхніми скінченнорізницевими апроксимаціями.
Зробивши це і використавши компактну нотацію, отримуємо різницеве
рівняння

\begin{verbatim}
\Delta u \cong\frac{u_{i,j+1}-2u_{i,j}+u_{i,j-1}}{h^2}+\frac{u_{i+1,j}-2u_{i,j}+u_{i-1,j}}{k^2}=0.
\end{verbatim}

У випадку \texttt{h=k} рівняння Лапласа приводиться до вигляду

\begin{verbatim}
u_{i+1,j}+u_{i-1,j}+u_{i,j+1}+u_{i,j-1}-4u_{ij}=0.
\end{verbatim}

Розв'язавши його відносно \texttt{u\_\{ij\}}, отримуємо

\begin{verbatim}
u_{ij}=\frac{1}{4}\left(u_{i+1,j}+u_{i-1,j}+u_{i,j+1}+u_{i,j-1}\right).
\end{verbatim}

Отже, значення розв'язку у внутрішньому вузлі апроксимується середнім
значенням розв'язку в чотирьох сусідніх вузлах. На цьому і ґрунтується
метод Лібмана.

\subsection{Алгоритм чисельного розв'язання задачі Діріхле (метод
Лібмана)}\label{ux430ux43bux433ux43eux440ux438ux442ux43c-ux447ux438ux441ux435ux43bux44cux43dux43eux433ux43e-ux440ux43eux437ux432ux44fux437ux430ux43dux43dux44f-ux437ux430ux434ux430ux447ux456-ux434ux456ux440ux456ux445ux43bux435-ux43cux435ux442ux43eux434-ux43bux456ux431ux43cux430ux43dux430}

\textbf{КРОК 1.} Присвоїмо значенням \texttt{u\_\{ij\}} у внутрішніх
вузлах початкове наближення, наприклад середнє значення всіх граничних
умов.

\textbf{КРОК 2.} Перераховуємо значення у всіх внутрішніх вузлах,
замінюючи старе значення середнім значенням чотирьох сусідніх точок.
Після кількох ітерацій процес збігається до наближеного розв'язку
задачі.

\section{ЗАУВАЖЕННЯ}\label{ux437ux430ux443ux432ux430ux436ux435ux43dux43dux44f-8}

\begin{enumerate}
\def\labelenumi{\arabic{enumi}.}
\tightlist
\item
  Якщо взяти сітку з чотирма внутрішніми точками \texttt{(m=n=4)}, то
  система набуває вигляду
\end{enumerate}

\begin{verbatim}
\begin{aligned}-4u_{22} + u_{23} + u_{32} &= -\sin(\pi/3),\\u_{22} - 4u_{23} + u_{33} &= -\sin(2\pi/3),\\u_{22} - 4u_{32} + u_{33} &= 0,\\u_{23} + u_{32} - 4u_{33} &= 0.\end{aligned}
\end{verbatim}

Цю систему потрібно розв'язувати відносно \texttt{u\_\{22\}},
\texttt{u\_\{23\}}, \texttt{u\_\{32\}} та \texttt{u\_\{33\}}. Розв'язок
такої системи можна знайти ітеративним методом. Метод Лібмана також є
одним з ітеративних методів.

\begin{enumerate}
\def\labelenumi{\arabic{enumi}.}
\setcounter{enumi}{1}
\item
  Якщо зменшити кроки сітки \texttt{h} і \texttt{k}, то отримаємо
  систему рівнянь того самого типу, але більшого розміру.
\item
  Систему рівнянь можна записати у матричній формі:
\end{enumerate}

\begin{verbatim}
\begin{bmatrix}-4 & 1 & 1 & 0 \\1 & -4 & 0 & 1 \\1 & 0 & -4 & 1 \\0 & 1 & 1 & -4\end{bmatrix}\begin{bmatrix}u_{22}\\u_{23}\\u_{32}\\u_{33}\end{bmatrix}=\begin{bmatrix}-\sin(\pi/3)\\-\sin(2\pi/3)\\0\\0\end{bmatrix}\approx\begin{bmatrix}-0.86\\-0.86\\0\\0\end{bmatrix}.
\end{verbatim}

Зі збільшенням кількості рівнянь матриця коефіцієнтів стає розрідженою,
тобто містить багато нулів. Такі системи розв'язуються спеціальними
методами, серед яких особливо важливі метод Якобі, метод Гаусса-Зайделя
та послідовна верхня релаксація.

\begin{enumerate}
\def\labelenumi{\arabic{enumi}.}
\setcounter{enumi}{3}
\item
  При розв'язанні задачі Неймана ті \emph{похідні}, що входять до
  граничної умови, також слід замінити різницею апроксимаціями.
\item
  Таким чином ви зможете розв'язувати задачі такого типу:
\end{enumerate}

а) \texttt{u\_\{xx\}\ +\ u\_\{yy\}\ =\ f(x,\ y)}

\texttt{б)\ \textquotesingle{}x\ u\_\{xx\}\ +\ u\_\{yy\}\ +\ 2u\ =\ \textbackslash{}sin(x\ -\ y)\textquotesingle{}}

\texttt{в)\ \textquotesingle{}\textbackslash{}sin\ x\textbackslash{},u\_\{xx\}\ +\ u\_\{xy\}\ +\ 3u\ =\ 0\textquotesingle{}}

(гетерогенні рівняння зі змінними коефіцієнтами).

\begin{enumerate}
\def\labelenumi{\arabic{enumi}.}
\setcounter{enumi}{5}
\tightlist
\item
  Якщо область, у якій розв'язується задача, має неправильну форму, її
  можна покрити сіткою і наближати розв'язок у вузлах, найближчих до
  межі, інтерполюючи граничні умови. Після цього задачу розв'язують
  звичайним способом.
\end{enumerate}

\begin{figure}
\centering
\pandocbounded{\includegraphics[keepaspectratio,alt={Фіг. 37.2. Поблизу криволінійної межі нові граничні умови отримують інтерполяцією.}]{images/_page_285_Figure_8.jpeg}}
\caption{Фіг. 37.2. Поблизу криволінійної межі нові граничні умови
отримують інтерполяцією.}
\end{figure}

Фіг. 37.2. Поблизу криволінійної межі нові граничні умови отримують
інтерполяцією.

\begin{enumerate}
\def\labelenumi{\arabic{enumi}.}
\setcounter{enumi}{6}
\tightlist
\item
  Існують спеціальні програмні пакети для чисельного розв'язання
  еліптичних задач. Один з відомих прикладів --- пакет ELLPACK, який
  дозволяє працювати з дво- та тривимірними задачами в різних
  координатних системах і для досить широкого класу граничних умов.
\end{enumerate}

\section{ЗАВДАННЯ}\label{ux437ux430ux432ux434ux430ux43dux43dux44f-8}

\begin{enumerate}
\def\labelenumi{\arabic{enumi}.}
\tightlist
\item
  Отримайте наближення для другої похідної
\end{enumerate}

\begin{verbatim}
f''(x)\approx \frac{f(x+h)-2f(x)+f(x-h)}{h^2}.
\end{verbatim}

\begin{enumerate}
\def\labelenumi{\arabic{enumi}.}
\setcounter{enumi}{1}
\item
  Виконайте дві ітерації для задачі Діріхле (37.5) методом Лібмана.
\item
  До якої алгебраїчної системи зводиться задача Діріхле для рівняння
  Пуассона в квадраті
\end{enumerate}

\begin{verbatim}
\begin{aligned}(\text{УЧП})\quad & u_{xx}+u_{yy}=f(x,y), && 0<x<1,\quad 0<y<1,\\(\text{ГУ})\quad & u(x,y)=g(x,y), && \text{на межі квадрата},\end{aligned}
\end{verbatim}

якщо розв'язувати її методом скінченних різниць?

\begin{enumerate}
\def\labelenumi{\arabic{enumi}.}
\setcounter{enumi}{3}
\tightlist
\item
  Розв'яжіть задачу 3 для рівняння
\end{enumerate}

\begin{verbatim}
u_{xx}+u_{yy}+2u=0,\qquad 0<x<1,\quad 0<y<1,
\end{verbatim}

за умови \texttt{u(x,y)=g(x,y)} на межі квадрата.

\begin{enumerate}
\def\labelenumi{\arabic{enumi}.}
\setcounter{enumi}{4}
\tightlist
\item
  Як би ви розв'язували задачу Неймана в квадраті
\end{enumerate}

\begin{verbatim}
\begin{aligned}(\text{УЧП})\quad & u_{xx}+u_{yy}=0, && 0<x<1,\quad 0<y<1,\\(\text{ГУ})\quad & u=0 && \text{на трьох сторонах квадрата},\\& \frac{\partial u}{\partial x}(1,y)=1, && 0\le y\le 1,\end{aligned}
\end{verbatim}

методом скінченних різниць?

\begin{enumerate}
\def\labelenumi{\arabic{enumi}.}
\setcounter{enumi}{5}
\tightlist
\item
  Побудуйте блок-схему алгоритму розв'язання задачі Діріхле в квадраті
\end{enumerate}

\begin{verbatim}
\begin{aligned}(\text{УЧП})\quad & u_{xx}+u_{yy}=f(x,y), && 0<x<1,\quad 0<y<1,\\(\text{ГУ})\quad & u(x,y)=g(x,y), && \text{на межі квадрата},\end{aligned}
\end{verbatim}

для довільної кількості вузлів сітки. Якщо знаєте мову програмування,
напишіть програму для виконання цих обчислень.

\section{Лекція 38. ЯВНІ РІЗНИЦЕВІ
СХЕМИ}\label{ux43bux435ux43aux446ux456ux44f-38.-ux44fux432ux43dux456-ux440ux456ux437ux43dux438ux446ux435ux432ux456-ux441ux445ux435ux43cux438}

МЕТА ЛЕКЦІЇ: Ввести поняття явної різницевої схеми і показати, як її
можна застосовувати до задач теплопровідності та хвильових задач.
Основна ідея полягає в тому, щоб замінити похідні скінченними різницями
і отримати формули, які безпосередньо виражають розв'язок на новому
часовому шарі через значення на попередньому шарі.

\texttt{На\ відміну\ від\ стаціонарних\ еліптичних\ задач,\ де\ значення\ у\ внутрішніх\ вузлах\ сітки\ визначаються\ одночасно,\ для\ задач,\ що\ залежать\ від\ часу,\ можна\ обчислювати\ розв\textquotesingle{}язок\ послідовно:\ \textquotesingle{}t=\textbackslash{}Delta\ t,\textbackslash{}\ 2\textbackslash{}Delta\ t,\textbackslash{}\ 3\textbackslash{}Delta\ t,\textbackslash{}dots\textquotesingle{}}

\subsection{Явна схема для рівняння
теплопровідності}\label{ux44fux432ux43dux430-ux441ux445ux435ux43cux430-ux434ux43bux44f-ux440ux456ux432ux43dux44fux43dux43dux44f-ux442ux435ux43fux43bux43eux43fux440ux43eux432ux456ux434ux43dux43eux441ux442ux456}

Розглянемо задачу теплопровідності для стрижня, початкова температура
якого дорівнює нулю. Температура лівого кінця фіксована, а на правому
кінці виконується умова теплообміну з навколишнім середовищем.

\begin{figure}
\centering
\pandocbounded{\includegraphics[keepaspectratio,alt={Фіг. 38.1. Сітка для рівняння теплопровідності.}]{images/_page_288_Figure_3.jpeg}}
\caption{Фіг. 38.1. Сітка для рівняння теплопровідності.}
\end{figure}

Фіг. 38.1. Сітка для рівняння теплопровідності.

Нехай температура середовища задається функцією \texttt{g(t)}. Тоді
маємо задачу

\begin{verbatim}
\begin{aligned}(38.1)\quad (\text{УЧП})\quad & u_t=u_{xx}, && 0<x<1,\quad 0<t<\infty,\\(\text{ГУ})\quad & u(0,t)=1,\quad u_x(1,t)=-\left[u(1,t)-g(t)\right], && 0<t<\infty,\\(\text{НУ})\quad & u(x,0)=0, && 0\le x\le 1.\end{aligned}
\end{verbatim}

Побудуємо прямокутну сітку з вузлами

\begin{verbatim}
\begin{aligned}x_j&=jh,\qquad j=0,1,2,\dots,n,\\t_i&=ik,\qquad i=0,1,2,\dots,m.\end{aligned}
\end{verbatim}

Позначимо через \texttt{u\_\{i,j\}} наближене значення функції
\texttt{u(x\_j,t\_i)}. Значення на лівій межі та на початковому шарі
відомі з умов задачі. Щоб знайти решту значень, замінимо похідні
наближеннями

\begin{verbatim}
\begin{aligned}u_t(x_j,t_i)&\approx \frac{u_{i+1,j}-u_{i,j}}{k},\\u_{xx}(x_j,t_i)&\approx \frac{u_{i,j+1}-2u_{i,j}+u_{i,j-1}}{h^2}.\end{aligned}
\end{verbatim}

Підставляючи ці вирази в рівняння теплопровідності, дістаємо явну
формулу

\begin{verbatim}
(38.2)\qquad u_{i+1,j}=u_{i,j}+\frac{k}{h^2}\left[u_{i,j+1}-2u_{i,j}+u_{i,j-1}\right].
\end{verbatim}

Щоб врахувати праву граничну умову, замінимо похідну \texttt{u\_x(1,t)}
лівою різницею:

\begin{verbatim}
\frac{u_{i,n}-u_{i,n-1}}{h}=-\left[u_{i,n}-g_i\right],\qquad g_i=g(ik).
\end{verbatim}

Звідси отримуємо

\begin{verbatim}
(38.4)\qquad u_{i,n}=\frac{u_{i,n-1}+h g_i}{1+h}.
\end{verbatim}

Формули (38.2) і (38.4) дають повну явну схему для послідовного переходу
від одного часового шару до наступного.

\subsection{Алгоритм обчислень за явною
схемою}\label{ux430ux43bux433ux43eux440ux438ux442ux43c-ux43eux431ux447ux438ux441ux43bux435ux43dux44c-ux437ux430-ux44fux432ux43dux43eux44e-ux441ux445ux435ux43cux43eux44e}

\texttt{**КРОК\ 1.**\ Обчислити\ значення\ розв\textquotesingle{}язку\ на\ шарі\ \textquotesingle{}t=\textbackslash{}Delta\ t\textquotesingle{}\ за\ формулою}

\begin{verbatim}
u_{2,j}=u_{1,j}+\frac{k}{h^2}\left[u_{1,j+1}-2u_{1,j}+u_{1,j-1}\right],\qquad j=2,3,\dots,n-1.
\end{verbatim}

\textbf{КРОК 2.} Значення на правій межі знайти з (38.4):

\begin{verbatim}
u_{2,n}=\frac{u_{2,n-1}+h g_2}{1+h}.
\end{verbatim}

\texttt{**КРОК\ 3.**\ Для\ переходу\ до\ шарів\ \textquotesingle{}t=2\textbackslash{}Delta\ t,\textbackslash{}\ 3\textbackslash{}Delta\ t,\textbackslash{}dots\textquotesingle{}\ повторювати\ кроки\ 1\ і\ 2,\ щоразу\ використовуючи\ значення\ з\ попереднього\ часового\ шару.}

\begin{figure}
\centering
\pandocbounded{\includegraphics[keepaspectratio,alt={Фіг. 38.2. Ілюстрація явної різницевої схеми.}]{images/_page_290_Figure_3.jpeg}}
\caption{Фіг. 38.2. Ілюстрація явної різницевої схеми.}
\end{figure}

Фіг. 38.2. Ілюстрація явної різницевої схеми.

\section{ЗАУВАЖЕННЯ}\label{ux437ux430ux443ux432ux430ux436ux435ux43dux43dux44f-9}

\begin{enumerate}
\def\labelenumi{\arabic{enumi}.}
\tightlist
\item
  Явна схема має істотний недолік: якщо часовий крок занадто великий
  порівняно з кроком за просторовою координатою, похибки округлення
  можуть швидко зрости і зруйнувати чисельний розв'язок. Для рівняння
  теплопровідності зазвичай потрібно, щоб
\end{enumerate}

\begin{verbatim}
\frac{k}{h^2}\le 0.5.
\end{verbatim}

\texttt{2.\ Зменшення\ кроків\ \textquotesingle{}\textbackslash{}Delta\ t\textquotesingle{}\ і\ \textquotesingle{}\textbackslash{}Delta\ x\textquotesingle{}\ зазвичай\ підвищує\ точність,\ але\ водночас\ збільшує\ обсяг\ обчислень.\ Залежність\ повної\ похибки\ від\ кроку\ сітки\ показано\ на\ рис.\ 38.5.}

\begin{enumerate}
\def\labelenumi{\arabic{enumi}.}
\setcounter{enumi}{2}
\tightlist
\item
  Для гіперболічної задачі
\end{enumerate}

\begin{verbatim}
\begin{aligned}(\text{УЧП})\quad & u_{tt}=u_{xx}, && 0<x<1,\quad 0<t<\infty,\\(\text{ГУ})\quad & u(0,t)=g_1(t),\quad u(1,t)=g_2(t), && 0<t<\infty,\\(\text{НУ})\quad & u(x,0)=\varphi(x),\quad u_t(x,0)=\psi(x), && 0\le x\le 1\end{aligned}
\end{verbatim}

також можна побудувати явну різницеву схему. Для цього використовують
центральні різниці

\begin{verbatim}
\begin{aligned}u_{tt}(x,t)&\approx \frac{u(x,t+k)-2u(x,t)+u(x,t-k)}{k^2},\\u_{xx}(x,t)&\approx \frac{u(x+h,t)-2u(x,t)+u(x-h,t)}{h^2}.\end{aligned}
\end{verbatim}

а початкову умову для швидкості наближують рівністю

\begin{verbatim}
u_t(x,0)\approx \frac{u(x,k)-u(x,0)}{k}=\frac{u(x,k)-\varphi(x)}{k}.
\end{verbatim}

Звідси отримуємо явну схему

\begin{verbatim}
(38.5)\qquad u(x,t+k)=2u(x,t)-u(x,t-k)+\left(\frac{k}{h}\right)^2\left[u(x+h,t)-2u(x,t)+u(x-h,t)\right].
\end{verbatim}

\texttt{Щоб\ почати\ обчислення,\ потрібно\ спочатку\ знайти\ значення\ на\ шарі\ \textquotesingle{}t=\textbackslash{}Delta\ t\textquotesingle{}:}

\begin{verbatim}
u(x,k)=\varphi(x)+k\psi(x).
\end{verbatim}

\begin{figure}
\centering
\pandocbounded{\includegraphics[keepaspectratio,alt={Фіг. 38.3. Алгоритм розв'язання гіперболічної задачі за явною схемою.}]{images/_page_291_Figure_4.jpeg}}
\caption{Фіг. 38.3. Алгоритм розв'язання гіперболічної задачі за явною
схемою.}
\end{figure}

Фіг. 38.3. Алгоритм розв'язання гіперболічної задачі за явною схемою.

\begin{figure}
\centering
\pandocbounded{\includegraphics[keepaspectratio,alt={Фіг. 38.4. Використання масивів при реалізації явної схеми.}]{images/_page_292_Figure_6.jpeg}}
\caption{Фіг. 38.4. Використання масивів при реалізації явної схеми.}
\end{figure}

Фіг. 38.4. Використання масивів при реалізації явної схеми.

\section{ЗАВДАННЯ}\label{ux437ux430ux432ux434ux430ux43dux43dux44f-9}

\begin{enumerate}
\def\labelenumi{\arabic{enumi}.}
\tightlist
\item
  Побудуйте явну скінченнорізницеву схему для задачі
\end{enumerate}

\begin{verbatim}
\begin{aligned}(\text{УЧП})\quad & u_t=u_{xx}, && 0<x<1,\quad 0<t<\infty,\\(\text{ГУ})\quad & u(0,t)=0,\quad u(1,t)=0, && 0<t<\infty,\\(\text{НУ})\quad & u(x,0)=\sin(\pi x), && 0\le x\le 1.\end{aligned}
\end{verbatim}

\texttt{Вважаючи\ \textquotesingle{}h=\textbackslash{}Delta\ x=0.1\textquotesingle{},\ знайдіть\ чисельний\ розв\textquotesingle{}язок\ при\ \textquotesingle{}t\_1=0.005\textquotesingle{},\ \textquotesingle{}t\_2=0.010\textquotesingle{},\ \textquotesingle{}t\_3=0.015\textquotesingle{}\ і\ побудуйте\ графік\ розв\textquotesingle{}язку\ для\ \textquotesingle{}t=0.015\textquotesingle{}.}

\begin{figure}
\centering
\pandocbounded{\includegraphics[keepaspectratio,alt={Фіг. 38.5. Залежність повної похибки від кроку сітки.}]{images/_page_293_Figure_6.jpeg}}
\caption{Фіг. 38.5. Залежність повної похибки від кроку сітки.}
\end{figure}

Фіг. 38.5. Залежність повної похибки від кроку сітки.

\begin{enumerate}
\def\labelenumi{\arabic{enumi}.}
\setcounter{enumi}{1}
\item
  Розв'яжіть задачу 1 аналітично методом розділення змінних і порівняйте
  значення у вузлах сітки з чисельним розв'язком.
\item
  Побудуйте блок-схему для чисельного розв'язання наведеної вище
  гіперболічної задачі.
\item
  Розв'яжіть задачу 1 з новою граничною умовою при \texttt{x=1}:
\end{enumerate}

\begin{verbatim}
u_x(1,t)=-\left[u(1,t)-1\right].
\end{verbatim}

\section{Лекція 39. СХЕМИ НЕЯВНИХ РІЗНИЦЬ (СХЕМА
КРАНКА-НІКОЛСОНА)}\label{ux43bux435ux43aux446ux456ux44f-39.-ux441ux445ux435ux43cux438-ux43dux435ux44fux432ux43dux438ux445-ux440ux456ux437ux43dux438ux446ux44c-ux441ux445ux435ux43cux430-ux43aux440ux430ux43dux43aux430-ux43dux456ux43aux43eux43bux441ux43eux43dux430}

МЕТА ЛЕКЦІЇ: Показати, як задачі, що залежать від часу, можна
розв'язувати за допомогою неявної різницевої схеми. У такій схемі
значення на новому часовому шарі вже не виражаються явно через
попередній шар; замість цього на кожному кроці потрібно розв'язувати
систему алгебраїчних рівнянь.

Перевага неявних схем над явними полягає в тому, що вони дозволяють
брати значно більший часовий крок без руйнування розв'язку через
накопичення похибок округлення. Для параболічних задач ми використаємо
відому схему Кранка-Ніколсона.

\texttt{Як\ уже\ зазначалося\ в\ попередній\ лекції,\ явна\ схема\ для\ рівняння\ теплопровідності\ працює\ надійно\ лише\ тоді,\ коли\ кроки\ сітки\ \textquotesingle{}\textbackslash{}Delta\ t\textquotesingle{}\ і\ \textquotesingle{}\textbackslash{}Delta\ x\textquotesingle{}\ задовольняють\ умову}

\begin{verbatim}
\frac{\Delta t}{(\Delta x)^2}\le 0.5.
\end{verbatim}

\texttt{Інакше\ схема\ стає\ чисельно\ нестабільною.\ Наприклад,\ якщо\ \textquotesingle{}\textbackslash{}Delta\ x=0.1\textquotesingle{},\ то\ часовий\ крок\ не\ повинен\ перевищувати\ \textquotesingle{}0.5\textbackslash{}Delta\ x\^{}2=0.005\textquotesingle{},\ і\ для\ переходу\ від\ \textquotesingle{}t=0\textquotesingle{}\ до\ \textquotesingle{}t=1\textquotesingle{}\ потрібно\ зробити\ 200\ кроків.}

Неявні схеми дають можливість працювати з більшими часовими кроками.
Ціна за це --- необхідність на кожному часовому кроці розв'язувати
систему алгебраїчних рівнянь.

\subsection{Неявна схема для рівняння
теплопровідності}\label{ux43dux435ux44fux432ux43dux430-ux441ux445ux435ux43cux430-ux434ux43bux44f-ux440ux456ux432ux43dux44fux43dux43dux44f-ux442ux435ux43fux43bux43eux43fux440ux43eux432ux456ux434ux43dux43eux441ux442ux456}

Розглянемо задачу

\begin{verbatim}
\begin{aligned}(39.2)\quad (\text{УЧП})\quad & u_t=u_{xx}, && 0<x<1,\quad 0<t<\infty,\\(\text{ГУ})\quad & u(0,t)=0,\quad u(1,t)=0, && 0<t<\infty,\\(\text{НУ})\quad & u(x,0)=1, && 0<x<1.\end{aligned}
\end{verbatim}

Використаємо такі наближення для часткових похідних:

\begin{verbatim}
\begin{aligned}u_t(x,t) &\cong \frac{u(x,t+k)-u(x,t)}{k},\\u_{xx}(x,t) &\cong\frac{\lambda}{h^2}\bigl[u(x+h,t+k)-2u(x,t+k)+u(x-h,t+k)\bigr]\\&\qquad +\frac{1-\lambda}{h^2}\bigl[u(x+h,t)-2u(x,t)+u(x-h,t)\bigr].\end{aligned}
\end{verbatim}

Тут параметр \texttt{\textbackslash{}lambda} належить відрізку
\texttt{{[}0,1{]}}. При \texttt{\textbackslash{}lambda=0.5} отримуємо
схему Кранка-Ніколсона, а при \texttt{\textbackslash{}lambda=0} ---
звичайну явну схему.

Після підстановки цих апроксимацій у рівняння теплопровідності отримуємо

\begin{verbatim}
\frac{u_{i+1,j}-u_{i,j}}{k}=\frac{\lambda}{h^2}\left(u_{i+1,j+1}-2u_{i+1,j}+u_{i+1,j-1}\right)+\frac{1-\lambda}{h^2}\left(u_{i,j+1}-2u_{i,j}+u_{i,j-1}\right).
\end{verbatim}

Позначаючи \texttt{r=k/h\^{}2} і переносячи всі невідомі з верхнього
часового шару в ліву частину, одержуємо різницеву схему

\begin{verbatim}
(39.4)\qquad-\lambda r\,u_{i+1,j+1}+(1+2\lambda r)u_{i+1,j}-\lambda r\,u_{i+1,j-1}=r(1-\lambda)u_{i,j+1}+[1-2r(1-\lambda)]u_{i,j}+r(1-\lambda)u_{i,j-1}.
\end{verbatim}

Якщо індекс \texttt{i} зафіксовано, а \texttt{j} пробігає значення
\texttt{2,3,\textbackslash{}dots,n-1}, то співвідношення (39.4) задають
систему з \texttt{n-2} рівнянь для \texttt{n-2} невідомих
\texttt{u\_\{i+1,2\},u\_\{i+1,3\},\textbackslash{}dots,u\_\{i+1,n-1\}}.

\begin{figure}
\centering
\pandocbounded{\includegraphics[keepaspectratio,alt={Фіг. 39.1. Сітка для схеми неявних різниць.}]{images/_page_296_Figure_5.jpeg}}
\caption{Фіг. 39.1. Сітка для схеми неявних різниць.}
\end{figure}

Фіг. 39.1. Сітка для схеми неявних різниць.

\begin{figure}
\centering
\pandocbounded{\includegraphics[keepaspectratio,alt={Фіг. 39.2. Шаблон неявної схеми.}]{images/_page_296_Figure_7.jpeg}}
\caption{Фіг. 39.2. Шаблон неявної схеми.}
\end{figure}

Фіг. 39.2. Шаблон неявної схеми
\texttt{u\_\{i+1,4\},\textbackslash{}ldots,u\_\{i+1,n-1\}}.

Фіг. 39.2 дає візуальне уявлення про структуру кожного рівняння системи
(39.4).

\subsection{Алгоритм розв'язання задачі
(39.2)}\label{ux430ux43bux433ux43eux440ux438ux442ux43c-ux440ux43eux437ux432ux44fux437ux430ux43dux43dux44f-ux437ux430ux434ux430ux447ux456-39.2}

\texttt{**КРОК\ 1.**\ Виберіть\ значення\ параметра\ \textquotesingle{}\textbackslash{}lambda\textquotesingle{},\ де\ \textquotesingle{}0\textbackslash{}le\textbackslash{}lambda\textbackslash{}le1\textquotesingle{}.}

\texttt{**КРОК\ 2.**\ Нехай,\ наприклад,\ \textquotesingle{}h=\textbackslash{}Delta\ x=0.2\textquotesingle{}\ і\ \textquotesingle{}k=\textbackslash{}Delta\ t=0.08\textquotesingle{}.\ Тоді\ \textquotesingle{}r=k/h\^{}2=2\textquotesingle{}.\ Візьмемо\ \textquotesingle{}\textbackslash{}lambda=0.5\textquotesingle{};\ це\ і\ є\ схема\ Кранка-Ніколсона.\ Для\ перших\ чотирьох\ внутрішніх\ вузлів\ нового\ часового\ шару\ одержуємо\ систему}

\begin{verbatim}
\begin{aligned}3u_{22}-u_{23} &= 1,\\-u_{22}+3u_{23}-u_{24} &= 1,\\-u_{23}+3u_{24}-u_{25} &= 1,\\-u_{24}+3u_{25} &= 1.\end{aligned}
\end{verbatim}

У матричній формі це записується так:

\begin{verbatim}
(39.5)\qquad\begin{bmatrix}3 & -1 & 0 & 0\\-1 & 3 & -1 & 0\\0 & -1 & 3 & -1\\0 & 0 & -1 & 3\end{bmatrix}\begin{bmatrix}u_{22}\\u_{23}\\u_{24}\\u_{25}\end{bmatrix}=\begin{bmatrix}1\\1\\1\\1\end{bmatrix}.
\end{verbatim}

Матриця цієї системи є тридіагональною. Такі системи зручно розв'язувати
методом прогонки. Для загальної системи

\begin{verbatim}
\begin{bmatrix}b_1 & c_1 & 0 & \cdots & 0\\a_1 & b_2 & c_2 & \cdots & 0\\0 & a_2 & b_3 & \ddots & \vdots\\\vdots & \ddots & \ddots & \ddots & c_{n-1}\\0 & \cdots & 0 & a_{n-1} & b_n\end{bmatrix}\begin{bmatrix}x_1\\x_2\\x_3\\\vdots\\x_n\end{bmatrix}=\begin{bmatrix}d_1\\d_2\\d_3\\\vdots\\d_n\end{bmatrix}
\end{verbatim}

вводимо коефіцієнти

\begin{verbatim}
c_1^*=\frac{c_1}{b_1},\qquad d_1^*=\frac{d_1}{b_1},
\end{verbatim}

а далі рекурентно обчислюємо

\begin{verbatim}
c_j^*=\frac{c_j}{b_j-a_{j-1}c_{j-1}^*},\qquad d_j^*=\frac{d_j-a_{j-1}d_{j-1}^*}{b_j-a_{j-1}c_{j-1}^*},\qquad j=2,3,\dots,n.
\end{verbatim}

Після цього розв'язок знаходиться зворотною підстановкою:

\begin{verbatim}
x_n=d_n^*,\qquad x_j=d_j^*-c_j^*x_{j+1},\qquad j=n-1,n-2,\dots,1.
\end{verbatim}

Застосувавши цей метод до системи (39.5), отримуємо

\begin{verbatim}
u_{22}=0.60,\qquad u_{23}=0.80,\qquad u_{24}=0.80,\qquad u_{25}=0.60.
\end{verbatim}

\texttt{Це\ наближені\ значення\ розв\textquotesingle{}язку\ у\ внутрішніх\ точках\ сітки\ на\ шарі\ \textquotesingle{}t=\textbackslash{}Delta\ t\textquotesingle{}.\ Щоб\ перейти\ до\ наступного\ моменту\ часу,\ потрібно\ знову\ розв\textquotesingle{}язати\ аналогічну\ тридіагональну\ систему.}

У неявній схемі обсяг обчислень на кожному кроці більший, ніж у явній,
зате хорошу точність можна отримати навіть при значно більшому часовому
кроці.

\section{ЗАВДАННЯ}\label{ux437ux430ux432ux434ux430ux43dux43dux44f-10}

\begin{enumerate}
\def\labelenumi{\arabic{enumi}.}
\item
  Виведіть рівняння (39.4) з різницевої схеми, що передує йому.
\item
  Побудуйте неявну скінченнорізницеву схему для задачі
\end{enumerate}

\begin{verbatim}
\begin{aligned}(\text{УЧП})\quad & u_t=u_{xx}, && 0<x<1,\quad 0<t<\infty,\\(\text{ГУ})\quad & u(0,t)=1,\quad u_x(1,t)+u(1,t)=g(t), && 0<t<\infty,\\(\text{НУ})\quad & u(x,0)=0, && 0\le x\le 1.\end{aligned}
\end{verbatim}

\begin{enumerate}
\def\labelenumi{\arabic{enumi}.}
\setcounter{enumi}{2}
\tightlist
\item
  Як за допомогою неявної схеми розв'язати задачу
\end{enumerate}

\begin{verbatim}
\begin{aligned}(\text{УЧП})\quad & u_t=u_{xx}+u, && 0<x<1,\quad 0<t<\infty,\\(\text{ГУ})\quad & u(0,t)=0,\quad u(1,t)=0, && 0<t<\infty,\\(\text{НУ})\quad & u(x,0)=\sin(\pi x), && 0\le x\le 1?\end{aligned}
\end{verbatim}

\begin{enumerate}
\def\labelenumi{\arabic{enumi}.}
\setcounter{enumi}{3}
\item
  Як виглядає шаблон рівняння (39.4), коли
  \texttt{\textbackslash{}lambda=1}?
\item
  Побудуйте блок-схему для розв'язання рівняння теплопровідності (39.2).
  Якщо можливо, напишіть комп'ютерну програму. Корисно також розв'язати
  цю задачу чисельно для різних значень \texttt{\textbackslash{}lambda}
  і порівняти результати з аналітичним розв'язком.
\end{enumerate}

\section{Лекція 40. АНАЛІТИЧНІ РОЗВ'ЯЗКИ VERSUS ЧИСЛОВІ РІШЕННЯ
Р}\label{ux43bux435ux43aux446ux456ux44f-40.-ux430ux43dux430ux43bux456ux442ux438ux447ux43dux456-ux440ux43eux437ux432ux44fux437ux43aux438-versus-ux447ux438ux441ux43bux43eux432ux456-ux440ux456ux448ux435ux43dux43dux44f-ux440}

\begin{figure}
\centering
\pandocbounded{\includegraphics[keepaspectratio,alt={Рис. 40.1. Розв'язок рівняння теплопровідності у різні моменти часу.}]{images/_page_301_Figure_4.jpeg}}
\caption{Рис. 40.1. Розв'язок рівняння теплопровідності у різні моменти
часу.}
\end{figure}

ис. 40.1. Розв'язок рівняння теплопровідності (40.1) у різні моменти
часу.

\subsection{Переваги аналітичних
розв'язків}\label{ux43fux435ux440ux435ux432ux430ux433ux438-ux430ux43dux430ux43bux456ux442ux438ux447ux43dux438ux445-ux440ux43eux437ux432ux44fux437ux43aux456ux432}

\begin{enumerate}
\def\labelenumi{\arabic{enumi}.}
\item
  Аналітичний розв'язок дозволяє обчислювати точне значення функції в
  будь-якій точці, а не лише чисельно наближене.
\item
  Аналітичне розв'язання завжди дозволяє обчислити значення розв'язку в
  одній точці \texttt{(x,\ t)} без необхідності обчислювати значення
  розв'язків в інших точках, як це відбувається при розв'язанні задач із
  явними або неявними різницевими схемами.
\item
  Аналітичне рішення дозволяє визначити розв'язок у будь-якій точці, а
  не лише на вузлах сітки.
\item
  Найважливішою перевагою аналітичного розв'язку для нас є можливість
  простежити вплив фізичних параметрів, початкових і граничних умов на
  природу розв'язку.
\end{enumerate}

Чисельні методи не виявляють цих регулярностей, оскільки дозволяють
знаходити розв'язок лише за заданих параметрів --- початкових і
граничних умовах. Іноді важливо знати взаємозв'язок між параметрами
моделі та розв'язком, особливо коли йдеться про оцінку фізичних
параметрів за типом розв'язку. Наприклад, припустимо, що розв'язок
\texttt{u} визначається експериментально, і ми знаємо аналітичний
розв'язок:

\begin{verbatim}
u = \text{функція від параметрів}.
\end{verbatim}

Тоді можна поставити питання про визначення параметрів як функцій
вхідних даних за схемою:

\begin{verbatim}
\text{Параметри} = \text{функція від } u = \text{функція вхідних даних}.
\end{verbatim}

Така задача називається \textbf{параметричною ідентифікацією}. Для
параметричної ідентифікації необхідно вміти розв'язувати рівняння з
частковими похідними. Трохи пізніше в цій лекції наведемо приклад
параметричної ідентифікації в біології, але спочатку розглянемо переваги
числових розв'язань.

\subsection{Переваги числових
розв'язань}\label{ux43fux435ux440ux435ux432ux430ux433ux438-ux447ux438ux441ux43bux43eux432ux438ux445-ux440ux43eux437ux432ux44fux437ux430ux43dux44c}

Головна перевага числових розв'язків полягає в тому, що їх можна
отримати навіть якщо аналітичне рішення неможливо. Майже всі нелінійні
диференціальні рівняння в частинних похідних потрібно розв'язувати
числовими методами, і більшість реальних моделей фізики, хімії, біології
тощо є нелінійними. Як правило, лінійні моделі наближаються до
нелінійних лише якщо в них відкинути нелінійні члени. Наприклад, ось
кілька нелінійних рівнянь:

\begin{enumerate}
\def\labelenumi{\arabic{enumi}.}
\tightlist
\item
  нелінійне хвильове рівняння \texttt{u\_\{tt\}\ =\ u\_\{xx\}\ +\ f(u)};
\item
  рівняння реакції з нелінійною дифузією
  \texttt{u\_t\ =\ u\_\{xx\}\ +\ f(u)};
\item
  система Ходжкіна--Хакслі (складна нелінійна система, що описує
  генерацію нервових імпульсів).
\end{enumerate}

Аналітичні розв'язки для цих рівнянь невідомі ні при яких нелінійних
функціях \texttt{f} і \texttt{g}. Тому загальний підхід до розв'язання
нелінійних (а в ряді випадків і лінійних) задач зазвичай базується на
числових розв'язаннях.

Тепер давайте розглянемо приклад того, як можна використати аналітичне
рішення для визначення фізичних параметрів.

\subsection{Параметрична ідентифікація (в
біології)}\label{ux43fux430ux440ux430ux43cux435ux442ux440ux438ux447ux43dux430-ux456ux434ux435ux43dux442ux438ux444ux456ux43aux430ux446ux456ux44f-ux432-ux431ux456ux43eux43bux43eux433ux456ux457}

Припустимо, що біолог намагається визначити швидкість, з якою іони калію
\texttt{(K\^{}+)} дифундують у екзоплазмовому розчині. Якщо знати цей
коефіцієнт дифузії, можна багато сказати про те, як нервові імпульси
передаються вздовж аксонів. Проблема в тому, що це співвідношення майже
неможливо знайти прямим вимірюванням. Однак можливо знайти математичний
зв'язок між концентрацією поташу \texttt{u(x,t)} і коефіцієнтом дифузії
\texttt{D}, а потім за вимірюваннями величини \texttt{u(x,t)} визначити
величину \texttt{D}.

Біологи Ходжкін і Кейз виявили, що після розміщення гігантського аксону
у спеціальному соляному розчині концентрація радіоактивного калію
(\texttt{\{\}\^{}\{42\}K}) вздовж аксону приблизно описується кривою в
початковий момент:

\begin{verbatim}
u(x, 0) = A \, e^{-x^2/a},
\end{verbatim}

де \texttt{A} і \texttt{a} --- параметри, що підбираються так, щоб крива
відповідала експериментальним даним.

\begin{figure}
\centering
\pandocbounded{\includegraphics[keepaspectratio,alt={Рис. 40.2. Початкова концентрація іонів.}]{images/_page_303_Figure_6.jpeg}}
\caption{Рис. 40.2. Початкова концентрація іонів.}
\end{figure}

Рис. 40.2. Початкова концентрація іонів (\texttt{\{\}\^{}\{42\}K}): а
--- теоретична крива; б --- експериментальні точки.

З часом іони поташу під дією дифузії та конвекції розтікаються вздовж
аксону. Ходжкін і Кейз припустили, що концентрація \texttt{u} є
розв'язком наступної задачі з початковими даними для рівняння дифузії:

\begin{verbatim}
(40.3)\qquad\begin{cases}u_t = D u_{xx} - V u_x, \quad -\infty < x < \infty,\ t > 0, \\u(x, 0) = A\, e^{-x^2/a}.\end{cases}
\end{verbatim}

Цю проблему вирішили в лекції 15, переключившись на рухому систему
координат. Загальна схема методу виглядає так: спочатку відкидаємо
конвективний член і отримуємо розв'язок суто дифузійної задачі (при
\texttt{V\ =\ 0}):

\begin{verbatim}
u_0(x, t) = \frac{A\sqrt{a}}{\sqrt{a+4Dt}}\, e^{-x^2/(a+4Dt)},
\end{verbatim}

а потім враховуємо конвективний член і отримуємо:

\begin{verbatim}
u(x, t) = \frac{A\sqrt{a}}{\sqrt{a+4Dt}}\, e^{-[(x-Vt)^{2}/(a+4Dt)]}.
\end{verbatim}

Параметри \texttt{A}, \texttt{a}, \texttt{D} і \texttt{V} визначаються з
умови найкращого узгодження з експериментом (наприклад, методом
найменших квадратів):

\begin{verbatim}
SS = \sum_{i=1}^{n} [y_i - (a + bx_i)]^2.
\end{verbatim}

\section{Лекція 41. КЛАСИФІКАЦІЯ РІВНЯНЬ З ЧАСТКОВИМИ ПОХІДНИМИ ДРУГОГО
ПОРЯДКУ ТА ЗВЕДЕННЯ ДО КАНОНІЧНОЇ
ФОРМИ}\label{ux43bux435ux43aux446ux456ux44f-41.-ux43aux43bux430ux441ux438ux444ux456ux43aux430ux446ux456ux44f-ux440ux456ux432ux43dux44fux43dux44c-ux437-ux447ux430ux441ux442ux43aux43eux432ux438ux43cux438-ux43fux43eux445ux456ux434ux43dux438ux43cux438-ux434ux440ux443ux433ux43eux433ux43e-ux43fux43eux440ux44fux434ux43aux443-ux442ux430-ux437ux432ux435ux434ux435ux43dux43dux44f-ux434ux43e-ux43aux430ux43dux43eux43dux456ux447ux43dux43eux457-ux444ux43eux440ux43cux438}

Загальне лінійне рівняння другого порядку з двома незалежними змінними
\texttt{x} і \texttt{y} має вигляд:

\begin{verbatim}
A u_{xx} + B u_{xy} + C u_{yy} + D u_x + E u_y + F u = G,
\end{verbatim}

де коефіцієнти \texttt{A,\ B,\ C,\ D,\ E,\ F,\ G} можуть залежати від
\texttt{x} і \texttt{y}. Класифікація рівняння визначається знаком
дискримінанта \texttt{B\^{}2\ -\ 4AC}:

\begin{itemize}
\tightlist
\item
  \texttt{B\^{}2\ -\ 4AC\ \textless{}\ 0} → \textbf{еліптичне}
  (наприклад, рівняння Лапласа);
\item
  \texttt{B\^{}2\ -\ 4AC\ =\ 0} → \textbf{параболічне} (наприклад,
  рівняння теплопровідності);
\item
  \texttt{B\^{}2\ -\ 4AC\ \textgreater{}\ 0} → \textbf{гіперболічне}
  (наприклад, хвильове рівняння).
\end{itemize}

\subsection{Зведення до канонічної форми (загальна
схема)}\label{ux437ux432ux435ux434ux435ux43dux43dux44f-ux434ux43e-ux43aux430ux43dux43eux43dux456ux447ux43dux43eux457-ux444ux43eux440ux43cux438-ux437ux430ux433ux430ux43bux44cux43dux430-ux441ux445ux435ux43cux430}

Щоб привести рівняння до канонічної форми, вводять нові змінні
\texttt{\textbackslash{}xi\ =\ \textbackslash{}xi(x,\ y)},
\texttt{\textbackslash{}eta\ =\ \textbackslash{}eta(x,\ y)}. Після
заміни змінних рівняння набирає вигляду:

\begin{verbatim}
(41.2)\qquad\overline{A} u_{\xi\xi} + \overline{B} u_{\xi\eta} + \overline{C} u_{\eta\eta} + \overline{D} u_{\xi} + \overline{E} u_{\eta} + \overline{F} u = \overline{G},
\end{verbatim}

де нові коефіцієнти виражаються через старі та часткові похідні заміни:

\begin{verbatim}
\overline{A} = A\xi_x^2 + B\xi_x\xi_y + C\xi_y^2,
\overline{B} = 2A\xi_x\eta_x + B(\xi_x\eta_y + \xi_y\eta_x) + 2C\xi_y\eta_y,
\overline{C} = A\eta_x^2 + B\eta_x\eta_y + C\eta_y^2,
\overline{D} = A\xi_{xx} + B\xi_{xy} + C\xi_{yy} + D\xi_x + E\xi_y,
\overline{E} = A\eta_{xx} + B\eta_{xy} + C\eta_{yy} + D\eta_x + E\eta_y,
\end{verbatim}

\texttt{\textbackslash{}overline\{F\}\ =\ F,\ \textbackslash{}quad\ \textbackslash{}overline\{G\}\ =\ G.}

\begin{verbatim}
## Приклад зведення до канонічної форми (гіперболічний випадок)

Розглянемо рівняння
\end{verbatim}

u\_\{xx\} - x\^{}2 u\_\{yy\} = 0.

\begin{verbatim}
Знаходимо: A = 1, B = 0, C = -x^2, значить B^2 - 4AC = 4x^2 > 0 — рівняння гіперболічне. Характеристичні рівняння:
\end{verbatim}

\texttt{\textbackslash{}frac\{dy\}\{dx\}\ =\ \textbackslash{}pm\ x\ \textbackslash{}quad\ \textbackslash{}Rightarrow\ \textbackslash{}quad\ y\ =\ \textbackslash{}pm\ \textbackslash{}frac\{x\^{}2\}\{2\}\ +\ c.}

\begin{verbatim}
Нові координати:
\end{verbatim}

\texttt{\textbackslash{}xi\ =\ y\ -\ \textbackslash{}frac\{x\^{}2\}\{2\},\ \textbackslash{}quad\ \textbackslash{}eta\ =\ y\ +\ \textbackslash{}frac\{x\^{}2\}\{2\}.}

\begin{verbatim}
Після підстановки в формули для \overline{A}, \overline{B}, \overline{C}:
\end{verbatim}

\texttt{\textbackslash{}overline\{A\}\ =\ 0,\ \textbackslash{}quad\ \textbackslash{}overline\{B\}\ \textbackslash{}neq\ 0,\ \textbackslash{}quad\ \textbackslash{}overline\{C\}\ =\ 0,}

\begin{verbatim}
і рівняння набирає канонічної гіперболічної форми:
\end{verbatim}

\texttt{u\_\{\textbackslash{}xi\textbackslash{}eta\}\ =\ \textbackslash{}frac\{-(x\^{}2+y\^{}2)\textbackslash{},u\_\textbackslash{}xi\ +\ (y\^{}2-x\^{}2)\textbackslash{},u\_\textbackslash{}eta\}\{8x\^{}2\ y\^{}3\}.}

\begin{verbatim}
## Нова система координат для параболічного прикладу

Нова система координат \xi = y - x, \eta = y (рис. 41.1).

![Рис. 41.1. Нова система координат.](images/_page_309_Figure_3.jpeg)

Рис. 41.1. Нова система координат \xi = y - x, \eta = y.

**ШАГ 2.** Залишилось знайти канонічну форму рівняння. Підставляючи \xi і \eta у формули для коефіцієнтів \overline{A}, \overline{B}, \overline{C}, \overline{D}, \overline{E}, \overline{F} і \overline{G}, отримуємо:
\end{verbatim}

\texttt{\textbackslash{}overline\{A\}\ =\ 0\ \textbackslash{}quad\ (\textbackslash{}text\{мало\ бути\ нулем;\ ми\ знаходили\ \}\ \textbackslash{}xi\ \textbackslash{}text\{\ з\ умови\ рівності\ нулю\ цього\ коефіцієнта\}),}
\texttt{\textbackslash{}overline\{B\}\ =\ 0\ \textbackslash{}quad\ (\textbackslash{}text\{ми\ вже\ раніше\ показали,\ що\ цей\ коефіцієнт\ теж\ обертається\ в\ нуль\}),}
\texttt{\textbackslash{}overline\{C\}\ \textbackslash{}neq\ 0.}

\begin{verbatim}
Нові допоміжні коефіцієнти (41.5):
\end{verbatim}

\texttt{\textbackslash{}overline\{D\}\ =\ A\textbackslash{}xi\_\{xx\}\ +\ B\textbackslash{}xi\_\{xy\}\ +\ C\textbackslash{}eta\_\{yy\}\ +\ D\textbackslash{}xi\_x\ +\ E\textbackslash{}xi\_y\ =\ 0,}
\texttt{\textbackslash{}overline\{E\}\ =\ A\textbackslash{}eta\_\{xx\}\ +\ B\textbackslash{}eta\_\{xy\}\ +\ C\textbackslash{}eta\_\{yy\}\ +\ D\textbackslash{}eta\_x\ +\ E\textbackslash{}eta\_y\ =\ 2(y\^{}2\ -\ x\^{}2),}
\texttt{\textbackslash{}tilde\{F\}\ =\ F\ =\ 0,\ \textbackslash{}quad\ \textbackslash{}overline\{G\}\ =\ G\ =\ 0.}

\begin{verbatim}
Отже, нове рівняння
\end{verbatim}

\texttt{u\_\{\textbackslash{}xi\textbackslash{}eta\}\ =\ \textbackslash{}frac\{\textbackslash{}eta\textbackslash{},\ u\_\textbackslash{}xi\ -\ \textbackslash{}xi\textbackslash{},\ u\_\textbackslash{}eta\}\{2(\textbackslash{}xi\^{}2\ -\ \textbackslash{}eta\^{}2)\}}

\begin{verbatim}
є загальним розв'язком відповідного рівняння.

## Конвективна дифузія — приклад застосування канонічних координат

Задача (15.2):
\end{verbatim}

\texttt{\textbackslash{}begin\{cases\}u\_t\ =\ D\ u\_\{xx\}\ -\ V\ u\_x,\ \textbackslash{}quad\ -\textbackslash{}infty\ \textless{}\ x\ \textless{}\ \textbackslash{}infty,\textbackslash{}\ t\ \textgreater{}\ 0,\ \textbackslash{}\textbackslash{}u(x,\ 0)\ =\ H(x),\textbackslash{}end\{cases\}}

\begin{verbatim}
де H(x) — функція Хевісайда. Замінимо змінні:
\end{verbatim}

\texttt{\textbackslash{}xi\ =\ x\ -\ Vt,\ \textbackslash{}quad\ \textbackslash{}eta\ =\ t.}

\begin{verbatim}
Розв'язок задачі конвективної дифузії:
\end{verbatim}

\texttt{u(x,\ t)\ =\ \textbackslash{}begin\{cases\}\ 0,\ \&\ t\ \textless{}\ x/V,\ \textbackslash{}\textbackslash{}\ P,\ \&\ t\ \textbackslash{}ge\ x/V.\ \textbackslash{}end\{cases\}}

\begin{verbatim}
![](images/_page_319_Figure_14.jpeg)

# Лекція 42. ОБЧИСЛЕННЯ ЗА МЕТОДОМ МОНТЕ-КАРЛО

**МЕТА ЛЕКЦІЇ:** Показати, що деякі інтеграли можна оцінити імовірнісним шляхом — методом Монте-Карло — і що ця ідея переноситься на розв'язок рівнянь в частинних похідних.

![Фіг. 42.2. Обчислення інтеграції Монте-Карло.](images/_page_315_Figure_7.jpeg)

Фіг. 42.2. Обчислення інтеграції Монте-Карло.

Очевидно, якщо випадковим чином кинути близько 100 дротиків у прямокутник R, що містить графік функції, можна отримати оцінку значення інтеграла, помноживши площу прямокутника R на відносну кількість дротиків під графіком функції. Результат нашої гри:
\end{verbatim}

\texttt{I\ \textbackslash{}approx\ \textbackslash{}text\{(площа\ \}\ R)\ \textbackslash{}times\ \textbackslash{}frac\{\textbackslash{}text\{кількість\ дротиків\ під\ графіком\}\}\{\textbackslash{}text\{загальна\ кількість\ дротиків\}\}.}

\begin{verbatim}
Щоб виконати ці обчислення на комп'ютері, нам потрібно якимось чином отримати послідовність випадкових точок, тобто наказати комп'ютеру «кидати дротики». Блок-схема на рис. 42.3 показує, як комп'ютер вирішить цю проблему.

![РИС. 42.3. Блок-схема розрахунку інтегралу методом Монте-Карло.](images/_page_316_Figure_2.jpeg)

**РИС. 42.3.** Блок-схема розрахунку інтегралу \int_{a}^{b} f(x)\,dx методом Монте-Карло (100 кидків); M = \max_a f(x).

## Випадкові числа

Перш ніж перейти до методу Монте-Карло для рівнянь з частковими похідними, зупинимося на важливому питанні — випадкових числах. Щойно, обчислюючи інтеграл, нам потрібно було побудувати послідовність випадкових точок P_i = (x_i, y_i) всередині прямокутника R. Іншими словами, x-координата точки P має бути випадковим числом з інтервалу [a, b], а y_i — з інтервалу [0, M]. Щоб знайти випадкові числа з цих інтервалів, звертаємося до послідовності випадкових чисел, рівномірно розподілених на відрізку [0, 1]. Тоді числа, рівномірно розподілені на відрізку [a, b], можна обчислити за формулою
\end{verbatim}

x = a + (b - a),r\_i,

\begin{verbatim}
де r_i \in [0, 1] — рівномірно розподілене випадкове число.

Тепер, звичайно, виникає питання: як отримати послідовність випадкових чисел \{r_i: i=1,2,\ldots\}, рівномірно розподілених на [0,1]? Найпоширенішим методом сьогодні є **метод конгруентних відрахувань**. Використовуючи цей метод, ми отримуємо послідовність випадкових цілих чисел, наприклад: 2120,\ 1401,\ 177,\ 3013, \ldots, а потім ліворуч від кожного числа розміщуємо десяткову крапку: 0{,}212;\ 0{,}1401;\ 0{,}0177;\ 0{,}3013; \ldots

## Отримання випадкових чисел за допомогою методу відрахування

1. Як перше випадкове ціле число, оберіть будь-яке число між 0 і P (число P обирається наперед).
2. Помножте це випадкове число на деяке фіксоване число — попередньо обраний множник M.
3. Додайте до добутку деяке фіксоване ціле число K, обране заздалегідь.
4. Поділіть отриману суму на P і візьміть залишок — це і є нове випадкове число.

Метод відрахування записується у вигляді формули:
\end{verbatim}

\texttt{r\_\{i+1\}\ \textbackslash{}equiv\ (M\ r\_i\ +\ K)\ \textbackslash{}bmod\ P.}

\begin{verbatim}
Це означає, що для отримання наступного r_{i+1} потрібно взяти попередній r_i, помножити його на M, додати K, поділити суму на P і взяти решту ділення.

**Наприклад**, якщо P = 100, M = 37, K = 16, r_0 = 15:
\end{verbatim}

\texttt{r\_1\ \textbackslash{}equiv\ (37\ \textbackslash{}cdot\ 15\ +\ 16)\ \textbackslash{}bmod\ 100\ =\ 571\ \textbackslash{}bmod\ 100\ =\ 71,}

\begin{verbatim}
M r_1 + K = 37 \cdot 71 + 16 = 2643, \quad r_2 = 2643 \bmod 100 = 43, \quad \ldots
\end{verbatim}

`\#\# Приклади інтегралів для методу Монте-Карло

\begin{verbatim}
I = \int_{0}^{1} e^{\sin x}\, dx.
I = \int_{0}^{1}\int_{0}^{1}\int_{0}^{1} e^{-(x^{2}+y^{2}+z^{2})}\, dx\, dy\, dz.
I = \int_{0}^{1}\int_{0}^{1}\int_{0}^{1}\int_{0}^{1} e^{-(x^{2}+y^{2}+z^{2}+\omega^{2})}\, dx\, dy\, dz\, d\omega.
\end{verbatim}

\subsection{ЗАВДАННЯ}\label{ux437ux430ux432ux434ux430ux43dux43dux44f-11}

\begin{enumerate}
\def\labelenumi{\arabic{enumi}.}
\item
  Яку послідовність випадкових чисел дає
  формула\texttt{r\_\{i+1\}\ \textbackslash{}equiv\ (3r\_i\ +\ 4)\ \textbackslash{}bmod\ 7},
  де\texttt{r\_0\ =\ 0}?
\item
  Як отримати послідовність випадкових чисел, розподіл яких показаний на
  рисунку?
\end{enumerate}

\pandocbounded{\includegraphics[keepaspectratio]{images/_page_319_Figure_16.jpeg}}

Іншими словами, як отримати послідовність
чисел\texttt{\textbackslash{}\{0,\ 1,\ 2\textbackslash{}\}}, якщо
ймовірність виникнення\texttt{0}і\texttt{2}дорівнює\texttt{0\{,\}25}, а
ймовірність появи\texttt{1}---\texttt{0\{,\}5}?

\section{Лекція 43. РОЗВ'ЯЗАННЯ РІВНЯНЬ З ЧАСТКОВИМИ ПОХІДНИМИ МЕТОДОМ
МОНТЕ-КАРЛО}\label{ux43bux435ux43aux446ux456ux44f-43.-ux440ux43eux437ux432ux44fux437ux430ux43dux43dux44f-ux440ux456ux432ux43dux44fux43dux44c-ux437-ux447ux430ux441ux442ux43aux43eux432ux438ux43cux438-ux43fux43eux445ux456ux434ux43dux438ux43cux438-ux43cux435ux442ux43eux434ux43eux43c-ux43cux43eux43dux442ux435-ux43aux430ux440ux43bux43e}

\textbf{МЕТА ЛЕКЦІЇ:} Показати, як побудувати таку гру, результатом якої
є приблизне розв'язання диференціального рівняння. Одна з цих ігор
(модель випадкової ходьби) веде до скінченно-різницевої апроксимації
задачі Діріхле у квадраті. Цю модель можна узагальнити для отримання
розв'язків інших задач.

У попередній лекції було сказано, що можливо винайти такі ігри на удачу,
результатом яких є розв'язки (можливо, приблизні) задач для рівнянь з
частковими похідними. У цій лекції ми покажемо, як побудувати таку гру
для наближеного розв'язання задачі Діріхле:

\begin{verbatim}
\begin{cases}u_{xx} + u_{yy} = 0, & 0 < x < 1,\ 0 < y < 1, \\u(x,y) = g(x,y), & \text{на межі квадрата.}\end{cases}
\end{verbatim}

Після того, як цю конкретну задачу буде розв'язано методом Монте-Карло,
ми розглянемо, як розв'язувати більш загальні задачі за допомогою цього
методу.

Щоб проілюструвати метод Монте-Карло, розглянемо гру під назвою
\textbf{«Блукаючий п'яниця»}. Щоб грати в цю гру, потрібна таблиця з
сіткою ліній, як показано на рис. 43.1.

\subsection{Як грати в «Блукаючого
п'яницю»?}\label{ux44fux43a-ux433ux440ux430ux442ux438-ux432-ux431ux43bux443ux43aux430ux44eux447ux43eux433ux43e-ux43fux44fux43dux438ux446ux44e}

\textbf{КРОК 1.} Блукання п'яниці починаються з довільної точки сітки (у
нашому випадку --- точки\texttt{A}).

\begin{figure}
\centering
\pandocbounded{\includegraphics[keepaspectratio,alt={Фіг. 43.1. Ігрове поле для гри «Блукаючий п'яниця».}]{images/_page_321_Figure_4.jpeg}}
\caption{Фіг. 43.1. Ігрове поле для гри «Блукаючий п'яниця».}
\end{figure}

Фіг. 43.1. Ігрове поле для гри «Блукаючий п'яниця».

\textbf{КРОК 2.} У кожному ході гри п'яний випадково переходить до
однієї з чотирьох суміжних точок сітки (у нашому випадку сусідні точки
для\texttt{A}--- це точки\texttt{B},\texttt{C},\texttt{D}і\texttt{E}, і
ймовірність досягти кожної з цих точок становить\texttt{0\{,\}25}).

{\def\LTcaptype{none} % do not increment counter
\begin{longtable}[]{@{}
  >{\raggedright\arraybackslash}p{(\linewidth - 4\tabcolsep) * \real{0.3333}}
  >{\raggedright\arraybackslash}p{(\linewidth - 4\tabcolsep) * \real{0.3333}}
  >{\raggedright\arraybackslash}p{(\linewidth - 4\tabcolsep) * \real{0.3333}}@{}}
\toprule\noalign{}
\begin{minipage}[b]{\linewidth}\raggedright
\textbf{Таблиця 43.1.} Ймовірність закінчення випадкового блукання в
граничній точці\texttt{p\_i}і величина винагороди\texttt{g\_i}
\end{minipage} & \begin{minipage}[b]{\linewidth}\raggedright
Гранична точка\texttt{p\_i}
\end{minipage} & \begin{minipage}[b]{\linewidth}\raggedright
\texttt{P\_A(p\_i)}
\end{minipage} \\
\midrule\noalign{}
\endhead
\bottomrule\noalign{}
\endlastfoot
1 & 0,04 & 1 \\
2 & 0,15 & 1 \\
3 & 0,03 & 1 \\
4 & 0,06 & 0 \\
5 & 0,17 & 0 \\
6 & 0,05 & 0 \\
7 & 0,06 & 0 \\
8 & 0,15 & 0 \\
9 & 0,03 & 0 \\
10 & 0,06 & 0 \\
11 & 0,16 & 0 \\
12 & 0,04 & 0 \\
\end{longtable}
}

\textbf{КРОК 3.} Після переходу до наступної точки процес блукання
відновлюється. Таким чином, п'яниця ходить від точки до точки, поки
випадково не опиниться на межі\texttt{p\_t}. Тут він зупиняється, і ми
записуємо номер точки\texttt{p\_t}. Це кінець однієї випадкової
прогулянки.

\textbf{КРОК 4.} Повторимо кроки 1--3 достатньо разів і знайдемо
відносну кількість прибуттів п'яниці в кожній з граничних
точок\texttt{p\_i}. Таблиця 43.1 показує типові результати, отримані
внаслідок 100 випадкових прогулянок.

\textbf{КРОК 5.} Припустимо, що п'яниця отримує
винагороду\texttt{g\_i\ =\ g(p\_i)}(значення граничної умови в
точці\texttt{p\_i}, якщо він досягає точки\texttt{p\_i}), і припустимо,
що мета гри --- обчислити середню\texttt{R(A)}винагороду за всі
випадкові прогулянки, починаючи з точки\texttt{A}. Цей середній виграш
визначається формулою:

\begin{verbatim}
R(A) = g_1 P_A(p_1) + g_2 P_A(p_2) + \dots + g_{12} P_A(p_{12}).
\end{verbatim}

Гра закінчується, коли знайдено значення\texttt{R(A)}. Відповідно до
даних таблиці 43.1:

\begin{verbatim}
R(A) = 1\cdot(0{,}04) + 1\cdot(0{,}15) + 1\cdot(0{,}03) + 0\cdot(0{,}06) + \dots + 0\cdot(0{,}04).
\end{verbatim}

\subsection{Для чого грати в «Блукаючого
п'яницю»?}\label{ux434ux43bux44f-ux447ux43eux433ux43e-ux433ux440ux430ux442ux438-ux432-ux431ux43bux443ux43aux430ux44eux447ux43eux433ux43e-ux43fux44fux43dux438ux446ux44e}

Виявляється, що середня винагорода\texttt{R(A)}є приблизним розв'язком
задачі Діріхле в точці\texttt{A}. Це цікаве спостереження базується на
двох фактах.

\textbf{1.} Припустимо, що п'яниця почав свій шлях з точки, що лежить на
межі. Кожна така прогулянка одразу закінчується в одному й тому ж місці,
і він отримує винагороду\texttt{g\_i}. Отже, середня винагорода за кожну
граничну точку дорівнює\texttt{g\_i}.

\textbf{2.} Тепер припустимо, що прогулянка починається з внутрішньої
точки. Тоді середня винагорода\texttt{R(A)}буде арифметичним середнім
середніх нагород за чотири сусідні бали:

\begin{verbatim}
R(A) = \frac{1}{4}\bigl[R(B) + R(C) + R(D) + R(E)\bigr].
\end{verbatim}

Давайте запитаємо себе: чому середнє значення
винагороди\texttt{R(A)}наближається до розв'язання задачі Діріхле в
точці\texttt{A}? Ми бачили, що значення\texttt{R(A)}задовольняє двом
рівнянням:

\begin{verbatim}
R(A) = \frac{1}{4}\bigl[R(B) + R(C) + R(D) + R(E)\bigr] \quad \text{(у внутрішніх точках)},
~~~`R(A)`=`g_i`\quad \text{(у граничних точках)}.
\end{verbatim}

Якщо g\_i --- значення функції g(x, y) з граничної умови у граничних
точках p\_i, то наші два рівняння точно збігаються з двома рівняннями,
отриманими при розв'язанні задачі Діріхле методом скінченних різниць:
\textsubscript{}\textasciitilde{}

\texttt{u\_\{i,j\}\ =\ \textbackslash{}frac\{1\}\{4\}\textbackslash{}bigl{[}u\_\{i-1,j\}\ +\ u\_\{i+1,j\}\ +\ u\_\{i,j-1\}\ +\ u\_\{i,j+1\}\textbackslash{}bigr{]},\ \textbackslash{}quad\ (i,j)\ \textbackslash{}text\{\ —\ внутрішня\ точка\},}

\begin{verbatim}
u_{ij} = g_{ij}, \quad (i,j) \text{— гранична точка}.
\end{verbatim}

Отже, значення\texttt{R(A)}дійсно наближає розв'язок рівняння Лапласа в
точці\texttt{A}.

\subsection{Розв'язання рівняння Лапласа за методом
Монте-Карло}\label{ux440ux43eux437ux432ux44fux437ux430ux43dux43dux44f-ux440ux456ux432ux43dux44fux43dux43dux44f-ux43bux430ux43fux43bux430ux441ux430-ux437ux430-ux43cux435ux442ux43eux434ux43eux43c-ux43cux43eux43dux442ux435-ux43aux430ux440ux43bux43e}

Наступні правила дозволяють отримати розв'язок в одній внутрішній точці
квадрата.

\textbf{КРОК 1.} Зробимо кілька випадкових прогулянок, починаючи з
точки\texttt{A}і закінчуючи однією з граничних точок. Відстежуємо,
скільки разів прогулянки закінчувалися на кожній межі.

\textbf{КРОК 2.} Після завершення всіх прогулянок для кожної граничної
точки обчислюємо відносну кількість прогулянок, які закінчилися в цій
точці. Позначимо ці значення\texttt{P\_A(p\_i)}.

\textbf{КРОК 3.} Обчислюємо наближений розв'язок\texttt{u(A)}за
формулою:

\begin{verbatim}
u(A) = g_1 P_A(p_1) + g_2 P_A(p_2) + \dots + g_N P_A(p_N),
\end{verbatim}

де\texttt{g\_i}--- значення функції\texttt{g}у точці\texttt{p\_i},
а\texttt{N}--- кількість граничних точок.

\subsection{Зауваження}\label{ux437ux430ux443ux432ux430ux436ux435ux43dux43dux44f-10}

\begin{enumerate}
\def\labelenumi{\arabic{enumi}.}
\item
  Зверніть увагу, що якщо відносні частоти\texttt{P\_A(p\_i)}відомі, то
  значення\texttt{u(A)}за інших граничних умов\texttt{g\_i}легко знайти.
  Для цього достатньо підставити нові\texttt{g\_i}у формулу --- не
  потрібно наново розіграти випадкові прогулянки.
\item
  У багатьох випадках дослідник зацікавлений у розв'язанні
  диференціального рівняння в певній точці. Якщо межа області є складною
  і рівняння залежить від трьох або чотирьох змінних, метод Монте-Карло
  може бути зручним. Методи Монте-Карло виникли для розв'язання дуже
  складних задач дифузії нейтронів. Ці задачі не можна розв'язати
  аналітично.
\end{enumerate}

\subsection{Розв'язання задачі Діріхле зі змінними
коефіцієнтами}\label{ux440ux43eux437ux432ux44fux437ux430ux43dux43dux44f-ux437ux430ux434ux430ux447ux456-ux434ux456ux440ux456ux445ux43bux435-ux437ux456-ux437ux43cux456ux43dux43dux438ux43cux438-ux43aux43eux435ux444ux456ux446ux456ux454ux43dux442ux430ux43cux438}

Розглянемо еліптичну крайову задачу в квадраті:

\begin{verbatim}
u_{xx} + (\sin x)\, u_{yy} = 0, \quad 0 < x < \pi,\ 0 < y < \pi,
~~~`u(x, y)`=`g(x, y)`\quad \text{на межі квадрата.}
\end{verbatim}

Щоб розв'язати цю задачу, замінимо u\_\{xx\}, u\_\{yy\} центральними
різницевими похідними: \textsubscript{}\textasciitilde{}

\texttt{u\_\{xx\}\ =\ u\_\{i,\textbackslash{},j+1\}\ -\ 2u\_\{i,\textbackslash{},j\}\ +\ u\_\{i,\textbackslash{},j-1\},\ \textbackslash{}quad\ u\_\{yy\}\ =\ u\_\{i+1,\textbackslash{},j\}\ -\ 2u\_\{i,\textbackslash{},j\}\ +\ u\_\{i-1,\textbackslash{},j\}.}

\begin{verbatim}
Підставляємо у рівняння і розв'язуємо за u_{i,j}:
\end{verbatim}

\texttt{u\_\{i,j\}\ =\ \textbackslash{}frac\{u\_\{i,j+1\}\ +\ u\_\{i,j-1\}\ +\ \textbackslash{}sin\ x\_j\textbackslash{},(u\_\{i+1,j\}\ +\ u\_\{i-1,j\})\}\{2(1+\textbackslash{}sin\ x\_j)\}.}

\begin{verbatim}
Коефіцієнти при u_{i+1,j}, u_{i,j+1}, u_{i,j-1}, u_{i-1,j} є додатними і в сумі дорівнюють одиниці. Отже, u_{i,j} є зваженим середнім значень у чотирьох сусідніх точках.

![Рис. 43.2. Вузли сітки для випадкового маршруту.](images/_page_324_Picture_9.jpeg)

Рис. 43.2. Вузли сітки для випадкового маршруту.

Можна **модифікувати Блукаючого П'яницю** так, щоб ймовірності переміщення до сусідніх точок дорівнювали коефіцієнтам відповідних членів. Якщо п'яниця знаходиться в точці (i, j), то він переходить у точку:
\end{verbatim}

\texttt{\textbackslash{}begin\{aligned\}\&(i,\ j+1)\ \textbackslash{}text\{\ з\ імовірністю\ \}\ \textbackslash{}frac\{1\}\{2(1+\textbackslash{}sin\ x\_j)\},\textbackslash{}\textbackslash{}\&(i,\ j-1)\ \textbackslash{}text\{\ з\ імовірністю\ \}\ \textbackslash{}frac\{1\}\{2(1+\textbackslash{}sin\ x\_j)\},\textbackslash{}\textbackslash{}\&(i+1,\ j)\ \textbackslash{}text\{\ з\ імовірністю\ \}\ \textbackslash{}frac\{\textbackslash{}sin\ x\_j\}\{2(1+\textbackslash{}sin\ x\_j)\},\textbackslash{}\textbackslash{}\&(i-1,\ j)\ \textbackslash{}text\{\ з\ імовірністю\ \}\ \textbackslash{}frac\{\textbackslash{}sin\ x\_j\}\{2(1+\textbackslash{}sin\ x\_j)\}.\textbackslash{}end\{aligned\}}

\begin{verbatim}
Інакше гра не змінюється. Модифікація для інших задач може бути більш винахідливою, але основна ідея залишається незмінною.

## ЗАВДАННЯ

1. Побудуйте блок-схему гри «Блукаючий п'яниця», щоб розв'язати задачу u_{xx} + u_{yy} = 0 у вузлах внутрішньої сітки.

2. Напишіть програму для розв'язання задачі відповідно до блок-схеми з завдання 1.

3. Як модифікувати гру для вирішення задачі:
\end{verbatim}

\texttt{\textbackslash{}begin\{cases\}u\_\{xx\}\ +\ u\_\{yy\}\ =\ 0,\ \&\ 0\ \textless{}\ x\ \textless{}\ 1,\textbackslash{}\ 0\ \textless{}\ y\ \textless{}\ 1,\ \textbackslash{}\textbackslash{}u(x,\ y)\ =\ g(x,\ y)\ \&\ \textbackslash{}text\{на\ межі?\}\textbackslash{}end\{cases\}}

\begin{verbatim}
Як виглядає в цій грі випадкова прогулянка?

4. Чи можна так модифікувати гру, щоб розв'язати задачу:
\end{verbatim}

\texttt{\textbackslash{}begin\{cases\}u\_\{xx\}\ +\ u\_\{yy\}\ =\ 0,\ \&\ 0\ \textless{}\ x\ \textless{}\ 1,\textbackslash{}\ 0\ \textless{}\ y\ \textless{}\ 1,\ \textbackslash{}\textbackslash{}u(x,\ 1)\ =\ 0,\ \&\ u(x,\ 0)\ =\ 0,\ \textbackslash{}\textbackslash{}u(0,\ y)\ =\ 1,\ \&\ \textbackslash{}tfrac\{\textbackslash{}partial\ u\}\{\textbackslash{}partial\ x\}(1,\ y)\ =\ 0?\textbackslash{}end\{cases\}}

\begin{verbatim}
5. Як модифікувати гру для розв'язання задачі:
\end{verbatim}

\texttt{\textbackslash{}begin\{aligned\}u\_t\ \&=\ \textbackslash{}alpha\^{}2\ u\_\{xx\}\ -\ \textbackslash{}beta\ u\_x\ -\ \textbackslash{}gamma\ u,\ \&\ 0\ \textless{}\ x\ \textless{}\ 1,\textbackslash{}\ 0\ \textless{}\ t\ \textless{}\ \textbackslash{}infty,\ \textbackslash{}\textbackslash{}u(0,t)\ \&=\ f(t),\textbackslash{}quad\ u(1,t)\ =\ g(t),\ \&\ 0\ \textless{}\ t\ \textless{}\ \textbackslash{}infty,\ \textbackslash{}\textbackslash{}u(x,0)\ \&=\ \textbackslash{}varphi(x),\ \&\ 0\ \textbackslash{}le\ x\ \textbackslash{}le\ 1.\textbackslash{}end\{aligned\}}

\begin{verbatim}
(ПРИМІТКА. Замінимо це рівняння скінченною різницею Кранка–Ніколсона, виразимо u_{i+1,j} через п'ять значень у суміжних точках.)

# Лекція 44. ВАРІАЦІЙНЕ ЧИСЛЕННЯ (РІВНЯННЯ ЕЙЛЕРА–ЛАГРАНЖА)

**МЕТА ЛЕКЦІЇ:** Ввести поняття функціоналу (функції від функції) і пояснити, як функціонали природно виникають у фізиці. Типовий функціонал є інтегралом
\end{verbatim}

\texttt{J{[}y{]}\ =\ \textbackslash{}int\_a\^{}b\ F(x,\textbackslash{},\ y,\textbackslash{},\ y\textquotesingle{})\textbackslash{},\ dx,}

\begin{verbatim}
що є функцією від функції y (припускається, що підінтегральна функція F(x, y, y') задана). Ми покажемо, як знайти функцію \overline{y}(x), яка мінімізує функціонал J[y]. Виявляється, що мінімізуюча функція \overline{y} має задовольняти так зване **рівняння Ейлера–Лагранжа**.

## Вступ. Задача брахістохрони

Варіаційне числення виникло одночасно з математичним аналізом у зв'язку з розв'язанням задач максимізації та мінімізації функцій функцій (функціоналів). Першою задачею варіаційного числення була задача брахістохрони, сформульована Йоганном Бернуллі у 1696 році.

Бернуллі показав, що час сходження записується у вигляді:
\end{verbatim}

\texttt{T\textbackslash{}\{y\textbackslash{}\}\ =\ \textbackslash{}int\_a\^{}b\ \textbackslash{}sqrt\{\textbackslash{}frac\{1+y\textquotesingle{}\^{}2\}\{2g\textbackslash{},y\}\}\textbackslash{},\ dx,}

\begin{verbatim}
і, отже, є функцією від функції. Оскільки багато функціоналів мають такий вигляд, ми зосередимося на вивченні загального функціоналу
\end{verbatim}

\texttt{(44.1)\textbackslash{}qquad\ J{[}y{]}\ =\ \textbackslash{}int\_a\^{}b\ F(x,\textbackslash{},\ y,\textbackslash{},\ y\textquotesingle{})\textbackslash{},\ dx.}

\begin{verbatim}
![Фіг. 44.1. Задача брахістохрони.](images/_page_329_Figure_5.jpeg)

Фіг. 44.1. Задача брахістохрони (час сходження треба мінімізувати).

## Мінімізація функціоналу

Розглянемо задачу пошуку функції y(x), яка мінімізує функціонал J[y] у класі гладких функцій, що задовольняють граничні умови y(a) = A, y(b) = B (рис. 44.2).

Нехай існує y = \overline{y}(x) — потрібна функція, і розглянемо малу варіацію функції \overline{y}, тобто функцію \overline{y} + \varepsilon\,\eta(x), де \varepsilon — мале число, а \eta(x) — гладка функція, яка задовольняє граничні умови \eta(a) = \eta(b) = 0 (рис. 44.2).

![Рис. 44.2. Варіація функції.](images/_page_329_Figure_7.jpeg)

Рис. 44.2. Варіація функції: а — варіація функції \overline{y} + \varepsilon\,\eta(x); б — мінімізуюча функція \overline{y}(x).

Оскільки \overline{y} мінімізує J, то
\end{verbatim}

\texttt{\textbackslash{}phi(\textbackslash{}varepsilon)\ =\ J{[}\textbackslash{}overline\{y\}\ +\ \textbackslash{}varepsilon\textbackslash{},\textbackslash{}eta{]}\ \textbackslash{}ge\ J{[}\textbackslash{}overline\{y\}{]}\ =\ \textbackslash{}phi(0)\ \textbackslash{}quad\ \textbackslash{}text\{для\ всіх\ \}\ \textbackslash{}varepsilon.}

\begin{verbatim}
Іншими словами, графік \phi(\varepsilon) = J[\overline{y}+\varepsilon\eta] як функція \varepsilon буде виглядати як рис. 44.3.

![Рис. 44.3. Графік функції в околиці мінімуму.](images/_page_330_Figure_6.jpeg)

Рис. 44.3. Графік функції J[\overline{y}+\varepsilon\eta] в околиці \varepsilon = 0.

З рис. 44.3 очевидно, що в точці мінімуму:
\end{verbatim}

\texttt{\textbackslash{}phi\textquotesingle{}(0)\ =\ \textbackslash{}frac\{d\}\{d\textbackslash{}varepsilon\}\ J{[}\textbackslash{}overline\{y\}\ +\ \textbackslash{}varepsilon\textbackslash{},\textbackslash{}eta{]}\ \textbackslash{}bigg\textbar{}\_\{\textbackslash{}varepsilon=0\}\ =\ 0.}

\begin{verbatim}
Диференціюючи під знаком інтеграла й інтегруючи за частинами, приходимо до висновку, що друга частина підінтегрального виразу має дорівнювати нулю.

~~
(44.2)\qquad\frac{\partial F}{\partial y} - \frac{d}{dx}\!\left(\frac{\partial F}{\partial y'}\right) = 0 \qquad \textbf{(Рівняння Ейлера–Лагранжа)}.
~~

Рівняння (44.2) називається **рівнянням Ейлера–Лагранжа**, і хоча загалом воно виглядає складно, коли в нього підставляється конкретна функція F(x, y, y'), воно стає звичайним диференціальним рівнянням другого порядку відносно невідомої функції \overline{y}(x).

## Теорема (необхідна умова мінімуму функціоналу)

Якщо функція y(x) мінімізує функціонал
\end{verbatim}

\texttt{J{[}y{]}\ =\ \textbackslash{}int\_a\^{}b\ F(x,\ y,\ y\textquotesingle{})\textbackslash{},\ dx}

\begin{verbatim}
у класі гладких функцій з граничними умовами y(a) = A, y(b) = B, то вона повинна задовольняти рівняння Ейлера–Лагранжа:
\end{verbatim}

\texttt{\textbackslash{}frac\{\textbackslash{}partial\ F\}\{\textbackslash{}partial\ y\}\ -\ \textbackslash{}frac\{d\}\{dx\}\textbackslash{}!\textbackslash{}left(\textbackslash{}frac\{\textbackslash{}partial\ F\}\{\textbackslash{}partial\ y\textquotesingle{}\}\textbackslash{}right)\ =\ 0,\ \textbackslash{}quad\ y(a)\ =\ A,\textbackslash{}\ y(b)\ =\ B.}

\begin{verbatim}
## Приклад: мінімальний функціонал J[y] = \int_{0}^{1}[y^2 + y'^2]\,dx

Спробуємо знайти функцію y(x), яка проходить через точки (0,0) і (1,1) і мінімізує функціонал:
\end{verbatim}

\texttt{J{[}y{]}\ =\ \textbackslash{}int\_\{0\}\^{}\{1\}{[}y\^{}2\ +\ y\textquotesingle{}\^{}2{]}\textbackslash{},\ dx,\ \textbackslash{}qquad\ y(0)\ =\ 0,\textbackslash{}\ y(1)\ =\ 1.}

\begin{verbatim}
Тут F(x, y, y') = y^2 + y'^2, отже:
\end{verbatim}

\texttt{\textbackslash{}frac\{\textbackslash{}partial\ F\}\{\textbackslash{}partial\ y\}\ =\ 2y,\ \textbackslash{}quad\ \textbackslash{}frac\{\textbackslash{}partial\ F\}\{\textbackslash{}partial\ y\textquotesingle{}\}\ =\ 2y\textquotesingle{}.}

\begin{verbatim}
Рівняння Ейлера–Лагранжа набирає вигляду:
\end{verbatim}

\texttt{2y\ -\ \textbackslash{}frac\{d\}\{dx\}(2y\textquotesingle{})\ =\ 0\ \textbackslash{}quad\ \textbackslash{}Rightarrow\ \textbackslash{}quad\ y\textquotesingle{}\textquotesingle{}\ -\ y\ =\ 0.}

\begin{verbatim}
Загальний розв'язок — y = C_1 e^x + C_2 e^{-x}. Розв'язуючи з граничними умовами y(0)=0 та y(1)=1, отримуємо:
\end{verbatim}

\texttt{\textbackslash{}overline\{y\}(x)\ =\ \textbackslash{}frac\{\textbackslash{}sinh\ x\}\{\textbackslash{}sinh\ 1\}\ =\ 0\{,\}42\textbackslash{},e\^{}x\ -\ 0\{,\}42\textbackslash{},e\^{}\{-x\}.}

\begin{verbatim}
![Фіг. 44.4. Допустимі криві та мінімізуюча функція.](images/_page_332_Figure_5.jpeg)

Фіг. 44.4. Усі допустимі гладкі криві задовольняють граничні умови y(0)=0 та y(1)=1: а — допустима крива; б — мінімізуюча функція y(x) = 0{,}42\,e^x - 0{,}42\,e^{-x}.

## Зауваження

1. Рівняння Ейлера–Лагранжа аналогічне рівності нульової похідної в диференціальному численні. Рівність нульової похідної є лише **необхідною**, але не достатньою умовою екстремуму. Те саме стосується рівняння Ейлера–Лагранжа. Однак дуже часто розв'язок рівняння Ейлера–Лагранжа фактично дає локальний або навіть глобальний мінімум.

2. Якщо y(x) мінімізує функціонал
\end{verbatim}

\texttt{J{[}y{]}\ =\ \textbackslash{}int\_a\^{}b\ 2\textbackslash{}pi\ y\ \textbackslash{}sqrt\{1\ +\ y\textquotesingle{}\^{}2\}\textbackslash{},\ dx,}

\begin{verbatim}
тоді площа поверхні, утворена обертанням кривої y(x) навколо осі x, буде мінімальною (рис. 44.5). Розв'язок рівняння Ейлера–Лагранжа в цьому випадку:

~~
y(x) = \alpha \cosh\!\left(\frac{x - \beta}{\alpha}\right) \quad \textbf{(ланцюгова лінія)},
~~

де константи \alpha і \beta визначаються з умов y(a) = A та y(b) = B.

3. У цій лекції ми виявили, що функція \Delta u = 0 доставляє мінімум функціоналу \Phi[u].

![Рис. 44.5. Мінімальна поверхня обертання.](images/_page_333_Picture_8.jpeg)

Рис. 44.5. Мінімальна поверхня обертання.

4. Основні закони фізики найчастіше формулюються мовою варіаційних принципів, а не диференціальних рівнянь. Прикладами є:

- **Принцип Ферма**: світло, поширюючись з однієї точки в іншу, обирає шлях, якому відповідає найкоротший час поширення.
- **Принцип Гамільтона**: у консервативному полі частинка рухається так, що інтеграл дії
\end{verbatim}

\texttt{\textbackslash{}int\_\{t\_1\}\^{}\{t\_2\}\ (E\_k\ -\ E\_p)\textbackslash{},\ dt\ \textbackslash{}quad\ \textbackslash{}text\{(кінетична\ енергія\ мінус\ потенціальна)\}}

\begin{verbatim}
буде стаціонарним. Звідси отримують диференціальне рівняння руху матеріальної точки.

Основні ідеї варіаційного числення можна розширити на функціонали, які залежать від функцій кількох незалежних змінних.

## ЗАВДАННЯ

1. Показати, що в класі гладких функцій, що задовольняють граничні умови y(0)=0 та y(\pi/2)=1, мінімальний функціонал має вигляд:
\end{verbatim}

\texttt{J{[}y{]}\ =\ \textbackslash{}int\_0\^{}\{\textbackslash{}pi/2\}\ {[}y\textquotesingle{}\^{}2\ -\ y\^{}2{]}\textbackslash{},\ dx.}

\begin{verbatim}
2. Якщо y(x) мінімізує функціонал J[y] = \int_a^b \sqrt{1+y'^2}\,dx, то крива y(x) між двома точками є відрізком прямої лінії.

3. Покажіть, що ланцюгова лінія y(x) = \alpha\cosh((x-\beta)/\alpha) задовольняє рівняння Ейлера–Лагранжа для функціоналу J[y] = \int_a^b y\sqrt{1+y'^2}\,dx.

4. Принцип Гамільтона стверджує, що матеріальна точка рухатиметься так, що інтеграл дії буде стаціонарним. Отримайте диференціальне рівняння руху матеріальної точки.

# Лекція 45. МЕТОД РІТЦА

**Мета лекції:** Показати, як задачі на мінімум функціоналів можна розв'язувати наближено, замінюючи нескінченновимірний простір допустимих функцій скінченновимірним підпростором.

**Наближений розв'язок** шукається у вигляді:
\end{verbatim}

\texttt{u\_n(x,y)\ =\ a\_1\textbackslash{}varphi\_1\ +\ a\_2\textbackslash{}varphi\_2\ +\ \textbackslash{}dots\ +\ a\_n\textbackslash{}varphi\_n,}

\begin{verbatim}
де \{\varphi_k\} — вибрані базисні функції (наприклад, тригонометричні поліноми або поліноми Лежандра). Коефіцієнти a_k обираються так, щоб значення функціоналу J[u_n] було мінімальним.

![Рис. 45.2. Блок-схема методу Рітца.](images/_page_343_Figure_2.jpeg)

Рис. 45.2. Блок-схема методу Рітца.

## ЗАВДАННЯ

1. Підставте розклад (46.8) у задачу (46.7) і отримайте послідовність задач P_0,\ P_1,\ P_2,\ldots

2. Покажіть, що нелінійну задачу можна наближено розв'язати методом збурень.

# Лекція 46. РОЗВ'ЯЗАННЯ РІВНЯНЬ З ЧАСТКОВИМИ ПОХІДНИМИ МЕТОДАМИ ТЕОРІЇ ЗБУРЕНЬ

**Мета лекції** — показати, як складні задачі з нелінійностями, змінними коефіцієнтами або нерегулярними областями можна розв'язувати як збурення простіших задач.

Наприклад, можна розглядати задачу Діріхле в деформованій області:
\end{verbatim}

\texttt{\textbackslash{}Delta\ u\ =\ 0,\ \textbackslash{}qquad\ 0\ \textless{}\ r\ \textless{}\ 1\ +\ \textbackslash{}tfrac\{1\}\{4\}\textbackslash{}sin\textbackslash{}theta,}

\begin{verbatim}
з граничною умовою:
\end{verbatim}

\texttt{u\textbackslash{}!\textbackslash{}left(1+\textbackslash{}tfrac\{1\}\{4\}\textbackslash{}sin\textbackslash{}theta,\textbackslash{},\textbackslash{}theta\textbackslash{}right)\ =\ \textbackslash{}cos\textbackslash{}theta,\ \textbackslash{}qquad\ 0\ \textbackslash{}le\ \textbackslash{}theta\ \textbackslash{}le\ 2\textbackslash{}pi.}

\begin{verbatim}
Дуже часто задачу, розв'язок якої невідомий, можна шляхом неперервної зміни деяких параметрів звести до близької, але легко розв'язуваної задачі. Тоді і розв'язок вихідної задачі буде близьким до розв'язку модифікованої задачі.

Наприклад, рівняння Лапласа можна неперервно модифікувати в нелінійне рівняння \Delta u + \varepsilon u^2 = 0 за допомогою сімейства задач, параметризованих числом \varepsilon, де 0 \le \varepsilon \le 1.

Розглядається область 0 < r < 1, 0 \le \theta < 2\pi з граничною умовою u(1,\theta) = g(\theta).

Розв'язок у вигляді ряду за степенями \varepsilon:
\end{verbatim}

\texttt{u\ =\ u\_0\ +\ \textbackslash{}varepsilon\ u\_1\ +\ \textbackslash{}varepsilon\^{}2\ u\_2\ +\ \textbackslash{}dots}

\begin{verbatim}
Підстановка цього розкладу дає послідовність лінійних задач P_0, P_1, P_2, \ldots, кожна з яких розв'язується аналітично.

Перше наближення:
\end{verbatim}

\texttt{u(r,\textbackslash{}theta)\ =\ u\_0(r,\textbackslash{}theta)\ +\ \textbackslash{}tfrac\{1\}\{4\}\textbackslash{},u\_1(r,\textbackslash{}theta),}

\begin{verbatim}
де:
\end{verbatim}

\texttt{u\_0(r,\textbackslash{}theta)\ =\ r\textbackslash{}cos\textbackslash{}theta,}

\begin{verbatim}
u = u_0 + u_1 = r\cos\theta -
\end{verbatim}

\texttt{\textbackslash{}frac\{r\^{}4-1\}\{32\}\ -\ \textbackslash{}frac\{(r\^{}4-r\^{}2)\}\{24\}\textbackslash{}cos(2\textbackslash{}theta).}

\begin{verbatim}
За деяких граничних умов (наприклад, g(\theta) = \cos\theta) задача має два різних матеріальні розв'язки, тоді як за інших (наприклад, g(\theta) = A\cos\theta з A > 20{,}65) задача не має реального розв'язку.

## ЗАВДАННЯ

1. Отримайте перший доданок розкладу для задачі:
\end{verbatim}

\texttt{\textbackslash{}Delta\ u\ =\ 0,\ \textbackslash{}qquad\ 0\ \textless{}\ r\ \textless{}\ 1\ +\ \textbackslash{}epsilon\textbackslash{}sin\textbackslash{}theta,}

\begin{verbatim}
u(1+\epsilon\sin\theta, \theta) = \cos\theta, \quad 0 \le \epsilon \le 1/4.
\end{verbatim}

\begin{enumerate}
\def\labelenumi{\arabic{enumi}.}
\setcounter{enumi}{1}
\item
  Розгляньте задачу методом збурень при
  малому\texttt{\textbackslash{}epsilon}і знайдіть
  наближення\texttt{u\_0}і\texttt{u\_1}.
\item
  Покажіть, що наближений розв'язок задачі збурення задовольняє формулу:
\end{enumerate}

\begin{verbatim}
f(x+h) = f(x) + f'(x)\,h + f''(x)\frac{h^2}{2!} + \dots
\end{verbatim}

(формула Тейлора, яку використовують для виведення умов збурень).

\begin{enumerate}
\def\labelenumi{\arabic{enumi}.}
\setcounter{enumi}{3}
\tightlist
\item
  Задача параметрів збурення для теплового рівняння:
\end{enumerate}

\begin{verbatim}
\begin{cases}u_t = \alpha^2 u_{xx}, \quad 0 < x < 1,\ 0 < t < \infty, \\u(0,t) = 0, \quad u(1,t) = 0, \quad 0 < t < \infty, \\u(x,0) = \sin(\pi x), \quad u_t(x,0) = 0, \quad 0 \le x \le 1.\end{cases}
\end{verbatim}

Дайте фізичну інтерпретацію цієї задачі.

\section{Лекція 47. РОЗВ'ЯЗАННЯ ДВОВИМІРНИХ ЗАДАЧ МЕТОДОМ КОНФОРМНИХ
ВІДОБРАЖЕНЬ}\label{ux43bux435ux43aux446ux456ux44f-47.-ux440ux43eux437ux432ux44fux437ux430ux43dux43dux44f-ux434ux432ux43eux432ux438ux43cux456ux440ux43dux438ux445-ux437ux430ux434ux430ux447-ux43cux435ux442ux43eux434ux43eux43c-ux43aux43eux43dux444ux43eux440ux43cux43dux438ux445-ux432ux456ux434ux43eux431ux440ux430ux436ux435ux43dux44c}

МЕТА ЛЕКЦІЇ: Показати, як деякі двовимірні крайові задачі для рівняння
Лапласа можна звести до простіших задач за допомогою конформного
відображення.

\pandocbounded{\includegraphics[keepaspectratio]{images/_page_353_Figure_9.jpeg}}

Ідея методу полягає в тому, що складну область на
площині\texttt{(x,y)}замінюють простішою областю на
площині\texttt{(u,v)}за допомогою комплексного
відображення\texttt{w=f(z)}, де

\begin{verbatim}
z = x + iy, \qquad w = u + iv.
\end{verbatim}

Для конформних відображень гармонічність зберігається:
якщо\texttt{\textbackslash{}varphi(x,y)}задовольняє рівняння Лапласа в
початковій області, то після переходу до нових координат ми знову
дістаємо рівняння Лапласа, але вже в геометрично простішій області.

У цій лекції розглянуто:

\begin{enumerate}
\def\labelenumi{\arabic{enumi}.}
\tightlist
\item
  основні властивості конформних відображень;
\item
  приклади побудови таких відображень;
\item
  застосування методу до задач Діріхле.
\end{enumerate}

\subsection{Комплексні функції та конформні
відображення}\label{ux43aux43eux43cux43fux43bux435ux43aux441ux43dux456-ux444ux443ux43dux43aux446ux456ux457-ux442ux430-ux43aux43eux43dux444ux43eux440ux43cux43dux456-ux432ux456ux434ux43eux431ux440ux430ux436ux435ux43dux43dux44f}

Комплексну функцію записують у вигляді

\begin{verbatim}
w = f(z).
\end{verbatim}

Вона переводить точки площини\texttt{z}у точки площини\texttt{w}.
Наприклад, відображення

\begin{verbatim}
w = z^2
\end{verbatim}

переносить перший квадрант площини\texttt{z}у верхню півплощину
площини\texttt{w}.

\begin{figure}
\centering
\pandocbounded{\includegraphics[keepaspectratio,alt={Рис. 47.1. Площина комплексної змінної z.}]{images/_page_355_Figure_5.jpeg}}
\caption{Рис. 47.1. Площина комплексної змінної z.}
\end{figure}

Рис. 47.1. Площина комплексної змінної\texttt{z}.

\begin{figure}
\centering
\pandocbounded{\includegraphics[keepaspectratio,alt={Рис. 47.2. Відображення першого квадранта площини z у верхню півплощину площини w.}]{images/_page_355_Figure_7.jpeg}}
\caption{Рис. 47.2. Відображення першого квадранта площини z у верхню
півплощину площини w.}
\end{figure}

Рис. 47.2. Відображення першого квадранта площини\texttt{z}у верхню
півплощину площини\texttt{w}.

Щоб побачити це в дійсних координатах, підставимо\texttt{z=x+iy}у
формулу\texttt{w=z\^{}2}:

\begin{verbatim}
w = (x+iy)^2 = x^2 - y^2 + 2ixy.
\end{verbatim}

Отже,

\begin{verbatim}
u = x^2 - y^2, \qquad v = 2xy.
\end{verbatim}

Це і є дійсна форма того самого відображення. Вона показує, як криві у
площині\texttt{(x,y)}переходять у криві в площині\texttt{(u,v)}.

\subsection{Означення конформного
відображення}\label{ux43eux437ux43dux430ux447ux435ux43dux43dux44f-ux43aux43eux43dux444ux43eux440ux43cux43dux43eux433ux43e-ux432ux456ux434ux43eux431ux440ux430ux436ux435ux43dux43dux44f}

\texttt{Відображення\textquotesingle{}w=f(z)\textquotesingle{}називається\ конформним\ у\ точці\textquotesingle{}z\_0\textquotesingle{},\ якщо\textquotesingle{}f\textquotesingle{}(z\_0)\textbackslash{}neq\ 0\textquotesingle{}.\ Воно\ називається\ конформним\ в\ області\textquotesingle{}D\textquotesingle{},\ якщо\ ця\ умова\ виконується\ в\ кожній\ точці\ області.}

Геометрично це означає, що відображення зберігає кути між кривими. Саме
ця властивість робить його корисним для розв'язування крайових задач:
складну межу можна перенести в простішу, не зруйнувавши структуру
задачі.

\subsection{Рівняння Лапласа у верхній
півплощині}\label{ux440ux456ux432ux43dux44fux43dux43dux44f-ux43bux430ux43fux43bux430ux441ux430-ux443-ux432ux435ux440ux445ux43dux456ux439-ux43fux456ux432ux43fux43bux43eux449ux438ux43dux456}

Розглянемо задачу Діріхле у верхній півплощині. Основна ідея полягає в
тому, щоб підібрати відображення, яке переведе верхню півплощину у
горизонтальну смугу. Для цього використовується класичне перетворення

\begin{verbatim}
w = \ln\left(\frac{z-1}{z+1}\right).
\end{verbatim}

Це відображення переносить верхню півплощину у смугу

\begin{verbatim}
-\infty < u < \infty, \qquad 0 < v < \pi.
\end{verbatim}

\begin{figure}
\centering
\pandocbounded{\includegraphics[keepaspectratio,alt={Рис. 47.3. Крайова задача у верхній півплощині.}]{images/_page_357_Figure_11.jpeg}}
\caption{Рис. 47.3. Крайова задача у верхній півплощині.}
\end{figure}

Рис. 47.3. Крайова задача у верхній півплощині.

\begin{figure}
\centering
\pandocbounded{\includegraphics[keepaspectratio,alt={Рис. 47.4. Перетворення складної задачі на простішу.}]{images/_page_358_Figure_2.jpeg}}
\caption{Рис. 47.4. Перетворення складної задачі на простішу.}
\end{figure}

Рис. 47.4. Перетворення складної задачі на простішу за допомогою
конформного відображення.

У смузі розв'язок уже знаходиться елементарно. Якщо на одній межі задано
нульовий потенціал, а на іншій одиничний, то

\begin{verbatim}
\varphi(u,v) = \frac{v}{\pi}.
\end{verbatim}

Далі потрібно виразити\texttt{v}через початкові
координати\texttt{x}і`y\$ та підставити цей вираз у знайдену формулу.

\subsection{Область між двома
колами}\label{ux43eux431ux43bux430ux441ux442ux44c-ux43cux456ux436-ux434ux432ux43eux43cux430-ux43aux43eux43bux430ux43cux438}

Другий типовий приклад полягає у відображенні області між двома
неконцентричними колами на кільце. Після такого перетворення задача
Діріхле зводиться до стандартної задачі в кільці, яку вже можна
розв'язувати у полярних координатах.

\begin{figure}
\centering
\pandocbounded{\includegraphics[keepaspectratio,alt={Рис. 47.5. Конформне відображення області на кільце.}]{images/_page_359_Figure_10.jpeg}}
\caption{Рис. 47.5. Конформне відображення області на кільце.}
\end{figure}

Рис. 47.5. Конформне відображення області на кільце.

Суть методу лишається тією самою:

\begin{enumerate}
\def\labelenumi{\arabic{enumi}.}
\tightlist
\item
  знайти зручне конформне відображення;
\item
  перенести задачу у простішу область;
\item
  розв'язати її в нових координатах;
\item
  повернутися до початкових змінних.
\end{enumerate}

Саме тому конформні відображення є одним із найсильніших інструментів
для двовимірних гармонічних задач.

\end{document}
